\documentclass[11pt,a4paper]{article}
\usepackage[utf8]{inputenc}
\usepackage[vietnamese,english]{babel}
\usepackage{amsmath,amssymb}
\usepackage{geometry}
\geometry{top=2cm,bottom=2cm,left=2cm,right=2cm}
\usepackage{tcolorbox}
\usepackage{tikz}
\usetikzlibrary{shapes,arrows,positioning}
\usepackage{enumitem}
\usepackage{hyperref}
\usepackage{booktabs}
\usepackage{tabularx}
\usepackage{graphicx}

\hypersetup{
    colorlinks=true,
    linkcolor=blue,
    urlcolor=blue,
    pdftitle={DBI202 - Database Systems - Ultra Edition},
    pdfauthor={DBI202 Team - FPT University}
}

% Định nghĩa môi trường cho câu hỏi
\newcounter{questioncount}
\newenvironment{question}[1]{%
  \par\vspace{1em}\noindent\textbf{\large Câu #1}\par\vspace{0.4em}
}{%
  \par\vspace{1em}\hrule\vspace{0.6em}
}

% Định nghĩa môi trường choices
\newenvironment{choices}{%
  \begin{enumerate}[label=\textbf{\Alph*.},leftmargin=2.2em,itemsep=2pt,topsep=3pt]
}{%
  \end{enumerate}
}
\newcommand{\choice}{\item}

% Định nghĩa môi trường codebox
\newenvironment{codebox}{%
  \begin{tcolorbox}[colback=gray!8!white,colframe=gray!50!black,arc=1.5mm,boxrule=0.6pt,left=3mm,right=3mm,top=2.5mm,bottom=2.5mm]
  \ttfamily\small
}{%
  \end{tcolorbox}
}

\begin{document}


\newpage
\section*{Chương 1: Introduction to Database Systems}

\begin{question}{1}
% [ID gốc: 1.1]
% [Nguồn: DBI202 Order 1 Q1]
A ____ is a logically coherent collection of data with some inherent meaning, representing some aspect of real world and being designed, built and populated with data for a specific purpose
\begin{choices}
\choice Database
\choice Database Instance
\choice Schema
\choice Schema Instance
\end{choices}
\textbf{Đáp án:} A
\textbf{Giải thích:} Đáp án theo dữ liệu đề: A — A. Database.

Nền tảng kiến thức: Chương này phân biệt dữ liệu, cơ sở dữ liệu, hệ quản trị CSDL (DBMS), mô hình dữ liệu, lược đồ và trạng thái dữ liệu. Cơ sở dữ liệu lưu dữ liệu có liên hệ; DBMS là phần mềm định nghĩa, truy vấn, cập nhật và bảo vệ dữ liệu đó.

Phân tích: Cơ sở dữ liệu là tập dữ liệu/thông tin có liên hệ, được tổ chức và duy trì để phục vụ một mục đích; DBMS là phần mềm quản lý tập dữ liệu ấy. Vì vậy, database không đồng nghĩa với một tệp đơn lẻ, một lần chụp dữ liệu (instance) hay phần mềm quản trị.

Đối chiếu phương án: B. Database Instance; C. Schema; D. Schema Instance. Các lựa chọn này không đáp ứng định nghĩa/điều kiện được hỏi; khi ôn, hãy chỉ ra chính xác điểm lệch của từng lựa chọn thay vì học thuộc mỗi chữ cái.
\end{question}

\begin{question}{2}
% [ID gốc: 1.2]
% [Nguồn: DBI202 Order 1 Q2]
Choose the most correct statement.
\begin{choices}
\choice Database is created and maintained by a DMBS
\choice All of the others
\choice Database is a collection of data that is managed by a DBMS
\choice Database is a collection of information that exists over a long period of time
\end{choices}
\textbf{Đáp án:} B
\textbf{Giải thích:} Đáp án theo dữ liệu đề: B — B. All of the others.

Nền tảng kiến thức: Chương này phân biệt dữ liệu, cơ sở dữ liệu, hệ quản trị CSDL (DBMS), mô hình dữ liệu, lược đồ và trạng thái dữ liệu. Cơ sở dữ liệu lưu dữ liệu có liên hệ; DBMS là phần mềm định nghĩa, truy vấn, cập nhật và bảo vệ dữ liệu đó.

Phân tích: Với lựa chọn tổng hợp, không chọn theo cảm tính: kiểm tra từng phát biểu độc lập. “All of the others/above” chỉ đúng khi mọi phát biểu còn lại đúng; “None” chỉ đúng khi không phát biểu nào trong nhóm thỏa yêu cầu. Đáp án lưu cần được đọc cùng các mệnh đề thành phần.

Đối chiếu phương án: A. Database is created and maintained by a DMBS; C. Database is a collection of data that is managed by a DBMS; D. Database is a collection of information that exists over a long period of time. Các lựa chọn này không đáp ứng định nghĩa/điều kiện được hỏi; khi ôn, hãy chỉ ra chính xác điểm lệch của từng lựa chọn thay vì học thuộc mỗi chữ cái.
\end{question}

\begin{question}{3}
% [ID gốc: 1.3]
% [Nguồn: DBI202 Order 1 Q3]
Which of following is never used as a data model?
\begin{choices}
\choice Hierarchical data model
\choice None of the others
\choice Graph-based data model
\choice Tree-based data model
\choice Relational data model
\end{choices}
\textbf{Đáp án:} B
\textbf{Giải thích:} Đáp án theo dữ liệu đề: B — B. None of the others.

Nền tảng kiến thức: Chương này phân biệt dữ liệu, cơ sở dữ liệu, hệ quản trị CSDL (DBMS), mô hình dữ liệu, lược đồ và trạng thái dữ liệu. Cơ sở dữ liệu lưu dữ liệu có liên hệ; DBMS là phần mềm định nghĩa, truy vấn, cập nhật và bảo vệ dữ liệu đó.

Phân tích: Data model là bộ khái niệm dùng để mô tả cấu trúc dữ liệu, quan hệ giữa dữ liệu và các ràng buộc. Relational, hierarchical, network và graph là các mô hình có thật; cần chọn mô hình khớp đặc điểm được hỏi, không nhầm mô hình với DBMS hay ngôn ngữ thao tác.

Đối chiếu phương án: A. Hierarchical data model; C. Graph-based data model; D. Tree-based data model; E. Relational data model. Các lựa chọn này không đáp ứng định nghĩa/điều kiện được hỏi; khi ôn, hãy chỉ ra chính xác điểm lệch của từng lựa chọn thay vì học thuộc mỗi chữ cái.
\end{question}

\begin{question}{4}
% [ID gốc: 1.4]
% [Nguồn: DBI202 Order 1 Q4]
A ___ is a notation for describing the structure of the data in a database, along with the constraints on that data
\begin{choices}
\choice data model
\choice database management system
\choice data operation
\choice data manipulation
\end{choices}
\textbf{Đáp án:} A
\textbf{Giải thích:} Đáp án theo dữ liệu đề: A — A. data model.

Nền tảng kiến thức: Chương này phân biệt dữ liệu, cơ sở dữ liệu, hệ quản trị CSDL (DBMS), mô hình dữ liệu, lược đồ và trạng thái dữ liệu. Cơ sở dữ liệu lưu dữ liệu có liên hệ; DBMS là phần mềm định nghĩa, truy vấn, cập nhật và bảo vệ dữ liệu đó.

Phân tích: Data model là bộ khái niệm dùng để mô tả cấu trúc dữ liệu, quan hệ giữa dữ liệu và các ràng buộc. Relational, hierarchical, network và graph là các mô hình có thật; cần chọn mô hình khớp đặc điểm được hỏi, không nhầm mô hình với DBMS hay ngôn ngữ thao tác.

Đối chiếu phương án: B. database management system; C. data operation; D. data manipulation. Các lựa chọn này không đáp ứng định nghĩa/điều kiện được hỏi; khi ôn, hãy chỉ ra chính xác điểm lệch của từng lựa chọn thay vì học thuộc mỗi chữ cái.
\end{question}

\begin{question}{5}
% [ID gốc: 1.5]
% [Nguồn: DBI202 Order 1 Q5]
What is a graph database?
\begin{choices}
\choice A graph database is a database that uses tree structures with nodes to represent and store information
\choice A graph database is a database that uses table structures to represent and store information
\choice A graph database is a database that uses graph structures with nodes, edges and properties to represent and store information
\choice None of the others
\end{choices}
\textbf{Đáp án:} C
\textbf{Giải thích:} Đáp án theo dữ liệu đề: C — C. A graph database is a database that uses graph structures with nodes, edges and properties to represent and store information.

Nền tảng kiến thức: Chương này phân biệt dữ liệu, cơ sở dữ liệu, hệ quản trị CSDL (DBMS), mô hình dữ liệu, lược đồ và trạng thái dữ liệu. Cơ sở dữ liệu lưu dữ liệu có liên hệ; DBMS là phần mềm định nghĩa, truy vấn, cập nhật và bảo vệ dữ liệu đó.

Phân tích: Mô hình đồ thị/mạng biểu diễn đối tượng thành nút và liên kết thành cạnh; mô hình phân cấp thường tổ chức quan hệ cha-con dạng cây. Hãy dựa vào cấu trúc liên kết được mô tả để phân biệt, không chỉ nhìn thấy từ “node”.

Đối chiếu phương án: A. A graph database is a database that uses tree structures with nodes to represent and store inf…; B. A graph database is a database that uses table structures to represent and store information; D. None of the others. Các lựa chọn này không đáp ứng định nghĩa/điều kiện được hỏi; khi ôn, hãy chỉ ra chính xác điểm lệch của từng lựa chọn thay vì học thuộc mỗi chữ cái.
\end{question}

\begin{question}{6}
% [ID gốc: 1.6]
% [Nguồn: DBI202 Order 1 Q6]
What is a database?
\begin{choices}
\choice A database is a collection of information that is organized so that it can easily be accessed, managed, and updated
\choice A database is a collection of data files
\choice A database is a collection of records
\choice A database is a collection of tuples
\end{choices}
\textbf{Đáp án:} A
\textbf{Giải thích:} Đáp án theo dữ liệu đề: A — A. A database is a collection of information that is organized so that it can easily be accessed, managed, and updated.

Nền tảng kiến thức: Chương này phân biệt dữ liệu, cơ sở dữ liệu, hệ quản trị CSDL (DBMS), mô hình dữ liệu, lược đồ và trạng thái dữ liệu. Cơ sở dữ liệu lưu dữ liệu có liên hệ; DBMS là phần mềm định nghĩa, truy vấn, cập nhật và bảo vệ dữ liệu đó.

Phân tích: Cơ sở dữ liệu là tập dữ liệu/thông tin có liên hệ, được tổ chức và duy trì để phục vụ một mục đích; DBMS là phần mềm quản lý tập dữ liệu ấy. Vì vậy, database không đồng nghĩa với một tệp đơn lẻ, một lần chụp dữ liệu (instance) hay phần mềm quản trị.

Đối chiếu phương án: B. A database is a collection of data files; C. A database is a collection of records; D. A database is a collection of tuples. Các lựa chọn này không đáp ứng định nghĩa/điều kiện được hỏi; khi ôn, hãy chỉ ra chính xác điểm lệch của từng lựa chọn thay vì học thuộc mỗi chữ cái.
\end{question}

\begin{question}{7}
% [ID gốc: 1.7]
% [Nguồn: DBI202 Order 1 Q7]
A ______ is a collection of programs that enables user to create and maintain databases. In other words, it is general-purpose software that provides the users with the processes of defining, constructing and manipulating the databases for various applications
\begin{choices}
\choice MSSQL
\choice DBMS
\choice DBS
\choice DB2
\end{choices}
\textbf{Đáp án:} B
\textbf{Giải thích:} Đáp án theo dữ liệu đề: B — B. DBMS.

Nền tảng kiến thức: Chương này phân biệt dữ liệu, cơ sở dữ liệu, hệ quản trị CSDL (DBMS), mô hình dữ liệu, lược đồ và trạng thái dữ liệu. Cơ sở dữ liệu lưu dữ liệu có liên hệ; DBMS là phần mềm định nghĩa, truy vấn, cập nhật và bảo vệ dữ liệu đó.

Phân tích: DBMS là viết tắt của Database Management System, tức hệ quản trị cơ sở dữ liệu. Nó cung cấp chức năng định nghĩa, lưu trữ, truy vấn, cập nhật và quản lý quyền trên CSDL; DBMS là phần mềm, còn database là dữ liệu được phần mềm quản lý.

Đối chiếu phương án: A. MSSQL; C. DBS; D. DB2. Các lựa chọn này không đáp ứng định nghĩa/điều kiện được hỏi; khi ôn, hãy chỉ ra chính xác điểm lệch của từng lựa chọn thay vì học thuộc mỗi chữ cái.
\end{question}

\begin{question}{8}
% [ID gốc: 1.9]
% [Nguồn: DBI202 Order 1 Q9]
Operations on data, constraint on data and structure of the data
A data model is a notation for describing data or information. And the description generally consist of:
\begin{choices}
\choice 2 parts
\choice 3 parts
\choice 4 parts
\choice 5 parts
\end{choices}
\textbf{Đáp án:} B
\textbf{Giải thích:} Đáp án theo dữ liệu đề: B — B. 3 parts.

Nền tảng kiến thức: Chương này phân biệt dữ liệu, cơ sở dữ liệu, hệ quản trị CSDL (DBMS), mô hình dữ liệu, lược đồ và trạng thái dữ liệu. Cơ sở dữ liệu lưu dữ liệu có liên hệ; DBMS là phần mềm định nghĩa, truy vấn, cập nhật và bảo vệ dữ liệu đó.

Phân tích: Data model là bộ khái niệm dùng để mô tả cấu trúc dữ liệu, quan hệ giữa dữ liệu và các ràng buộc. Relational, hierarchical, network và graph là các mô hình có thật; cần chọn mô hình khớp đặc điểm được hỏi, không nhầm mô hình với DBMS hay ngôn ngữ thao tác.

Đối chiếu phương án: A. 2 parts; C. 4 parts; D. 5 parts. Các lựa chọn này không đáp ứng định nghĩa/điều kiện được hỏi; khi ôn, hãy chỉ ra chính xác điểm lệch của từng lựa chọn thay vì học thuộc mỗi chữ cái.
\end{question}

\begin{question}{9}
% [ID gốc: 1.10]
% [Nguồn: DBI202 Order 1 Q10]
DBMS stands for what?
\begin{choices}
\choice Database Managing Systems
\choice Database Management System
\choice Database Manage System
\choice Database Managable System
\end{choices}
\textbf{Đáp án:} B
\textbf{Giải thích:} Đáp án theo dữ liệu đề: B — B. Database Management System.

Nền tảng kiến thức: Chương này phân biệt dữ liệu, cơ sở dữ liệu, hệ quản trị CSDL (DBMS), mô hình dữ liệu, lược đồ và trạng thái dữ liệu. Cơ sở dữ liệu lưu dữ liệu có liên hệ; DBMS là phần mềm định nghĩa, truy vấn, cập nhật và bảo vệ dữ liệu đó.

Phân tích: DBMS là viết tắt của Database Management System, tức hệ quản trị cơ sở dữ liệu. Nó cung cấp chức năng định nghĩa, lưu trữ, truy vấn, cập nhật và quản lý quyền trên CSDL; DBMS là phần mềm, còn database là dữ liệu được phần mềm quản lý.

Đối chiếu phương án: A. Database Managing Systems; C. Database Manage System; D. Database Managable System. Các lựa chọn này không đáp ứng định nghĩa/điều kiện được hỏi; khi ôn, hãy chỉ ra chính xác điểm lệch của từng lựa chọn thay vì học thuộc mỗi chữ cái.
\end{question}

\begin{question}{10}
% [ID gốc: 1.12]
% [Nguồn: DBI202 Order 1 Q12]
Choose an incorrect statement
\begin{choices}
\choice None of the others
\choice Database is created and maintained by a DMBS
IC]
Database is a collection of information that exists over a long opened of time
\choice Database is a collection of data that is managed by a DBMS
\end{choices}
\textbf{Đáp án:} A
\textbf{Giải thích:} Đáp án theo dữ liệu đề: A — A. None of the others.

Nền tảng kiến thức: Chương này phân biệt dữ liệu, cơ sở dữ liệu, hệ quản trị CSDL (DBMS), mô hình dữ liệu, lược đồ và trạng thái dữ liệu. Cơ sở dữ liệu lưu dữ liệu có liên hệ; DBMS là phần mềm định nghĩa, truy vấn, cập nhật và bảo vệ dữ liệu đó.

Phân tích: Với lựa chọn tổng hợp, không chọn theo cảm tính: kiểm tra từng phát biểu độc lập. “All of the others/above” chỉ đúng khi mọi phát biểu còn lại đúng; “None” chỉ đúng khi không phát biểu nào trong nhóm thỏa yêu cầu. Đáp án lưu cần được đọc cùng các mệnh đề thành phần.

Đối chiếu phương án: B. Database is created and maintained by a DMBS IC] Database is a collection of information that…; D. Database is a collection of data that is managed by a DBMS. Các lựa chọn này không đáp ứng định nghĩa/điều kiện được hỏi; khi ôn, hãy chỉ ra chính xác điểm lệch của từng lựa chọn thay vì học thuộc mỗi chữ cái.
\end{question}

\begin{question}{11}
% [ID gốc: 1.13]
% [Nguồn: DBI202 Order 1 Q13]
What is the oldest data model?
\begin{choices}
\choice Network Data Model
\choice Flat File Data Model
\choice Relational Data Model
\choice Hierachical Data Model
\end{choices}
\textbf{Đáp án:} B
\textbf{Giải thích:} Đáp án theo dữ liệu đề: A — A. Network Data Model.

Nền tảng kiến thức: Chương này phân biệt dữ liệu, cơ sở dữ liệu, hệ quản trị CSDL (DBMS), mô hình dữ liệu, lược đồ và trạng thái dữ liệu. Cơ sở dữ liệu lưu dữ liệu có liên hệ; DBMS là phần mềm định nghĩa, truy vấn, cập nhật và bảo vệ dữ liệu đó.

Phân tích: Data model là bộ khái niệm dùng để mô tả cấu trúc dữ liệu, quan hệ giữa dữ liệu và các ràng buộc. Relational, hierarchical, network và graph là các mô hình có thật; cần chọn mô hình khớp đặc điểm được hỏi, không nhầm mô hình với DBMS hay ngôn ngữ thao tác.

Đối chiếu phương án: B. Flat File Data Model; C. Relational Data Model; D. Hierachical Data Model. Các lựa chọn này không đáp ứng định nghĩa/điều kiện được hỏi; khi ôn, hãy chỉ ra chính xác điểm lệch của từng lựa chọn thay vì học thuộc mỗi chữ cái.
\end{question}

\begin{question}{12}
% [ID gốc: 1.14]
% [Nguồn: DBI202 Order 1 Q14]
How many data models mentioned in our database course?
\begin{choices}
\choice 3
\choice 4
\choice 5
\choice 6
\end{choices}
\textbf{Đáp án:} C
\textbf{Giải thích:} Đáp án theo dữ liệu đề: D — D. 6.

Nền tảng kiến thức: Chương này phân biệt dữ liệu, cơ sở dữ liệu, hệ quản trị CSDL (DBMS), mô hình dữ liệu, lược đồ và trạng thái dữ liệu. Cơ sở dữ liệu lưu dữ liệu có liên hệ; DBMS là phần mềm định nghĩa, truy vấn, cập nhật và bảo vệ dữ liệu đó.

Phân tích: Đây là câu đếm theo phạm vi giáo trình cụ thể, không có một con số phổ quát cho mọi môn CSDL. Đối chiếu danh sách mô hình đã được giảng/ghi trong đề cương của học phần; đáp án giữ nguyên theo khóa nguồn, còn phần giải thích không khẳng định ngoài phạm vi tài liệu môn.

Đối chiếu phương án: A. 3; B. 4; C. 5. Các lựa chọn này không đáp ứng định nghĩa/điều kiện được hỏi; khi ôn, hãy chỉ ra chính xác điểm lệch của từng lựa chọn thay vì học thuộc mỗi chữ cái.
\end{question}

\begin{question}{13}
% [ID gốc: 1.15]
% [Nguồn: DBI202 Order 1 Q15]
Database data models usually have a way to go describe limitations on what the data can be, is called
\begin{choices}
\choice Constraints on the data
\choice Operations on the data
\choice Structure of the data
\choice definition of data
\end{choices}
\textbf{Đáp án:} A
\textbf{Giải thích:} Đáp án theo dữ liệu đề: A — A. Constraints on the data.

Nền tảng kiến thức: Chương này phân biệt dữ liệu, cơ sở dữ liệu, hệ quản trị CSDL (DBMS), mô hình dữ liệu, lược đồ và trạng thái dữ liệu. Cơ sở dữ liệu lưu dữ liệu có liên hệ; DBMS là phần mềm định nghĩa, truy vấn, cập nhật và bảo vệ dữ liệu đó.

Phân tích: Constraint là quy tắc giới hạn giá trị/hình dạng dữ liệu hợp lệ, như miền giá trị, khóa hoặc tham chiếu. Operations mô tả thao tác trên dữ liệu, còn structure mô tả thành phần dữ liệu; constraint trả lời trực tiếp câu hỏi “giá trị nào được phép”.

Đối chiếu phương án: B. Operations on the data; C. Structure of the data; D. definition of data. Các lựa chọn này không đáp ứng định nghĩa/điều kiện được hỏi; khi ôn, hãy chỉ ra chính xác điểm lệch của từng lựa chọn thay vì học thuộc mỗi chữ cái.
\end{question}

\begin{question}{14}
% [ID gốc: 1.17]
% [Nguồn: DBI202 Order 1 Q17]
What are the disadvantages of network data model?
\begin{choices}
\choice Not support the high-level query language
\choice Not support the relationships between nodes
\choice Not support the database management system
\choice Not support the storage method of very large amounts of data
\end{choices}
\textbf{Đáp án:} A
\textbf{Giải thích:} Đáp án theo dữ liệu đề: A — A. Not support the high-level query language.

Nền tảng kiến thức: Chương này phân biệt dữ liệu, cơ sở dữ liệu, hệ quản trị CSDL (DBMS), mô hình dữ liệu, lược đồ và trạng thái dữ liệu. Cơ sở dữ liệu lưu dữ liệu có liên hệ; DBMS là phần mềm định nghĩa, truy vấn, cập nhật và bảo vệ dữ liệu đó.

Phân tích: Mô hình đồ thị/mạng biểu diễn đối tượng thành nút và liên kết thành cạnh; mô hình phân cấp thường tổ chức quan hệ cha-con dạng cây. Hãy dựa vào cấu trúc liên kết được mô tả để phân biệt, không chỉ nhìn thấy từ “node”.

Đối chiếu phương án: B. Not support the relationships between nodes; C. Not support the database management system; D. Not support the storage method of very large amounts of data. Các lựa chọn này không đáp ứng định nghĩa/điều kiện được hỏi; khi ôn, hãy chỉ ra chính xác điểm lệch của từng lựa chọn thay vì học thuộc mỗi chữ cái.
\end{question}

\begin{question}{15}
% [ID gốc: 1.18]
% [Nguồn: DBI202 Order 1 Q18]
What is information?
\begin{choices}
\choice A collection of unprocessed data
\choice A collection of processed data
\choice A collection of processed information
\choice A collection of unprocessed information
\end{choices}
\textbf{Đáp án:} B
\textbf{Giải thích:} Đáp án theo dữ liệu đề: B — B. A collection of processed data.

Nền tảng kiến thức: Chương này phân biệt dữ liệu, cơ sở dữ liệu, hệ quản trị CSDL (DBMS), mô hình dữ liệu, lược đồ và trạng thái dữ liệu. Cơ sở dữ liệu lưu dữ liệu có liên hệ; DBMS là phần mềm định nghĩa, truy vấn, cập nhật và bảo vệ dữ liệu đó.

Phân tích: Information là dữ liệu đã được xử lý/diễn giải để có ý nghĩa đối với mục tiêu sử dụng. Dữ liệu thô là các giá trị ghi nhận ban đầu; sau khi được tổng hợp hoặc đặt vào ngữ cảnh, chúng có thể trở thành thông tin.

Đối chiếu phương án: A. A collection of unprocessed data; C. A collection of processed information; D. A collection of unprocessed information. Các lựa chọn này không đáp ứng định nghĩa/điều kiện được hỏi; khi ôn, hãy chỉ ra chính xác điểm lệch của từng lựa chọn thay vì học thuộc mỗi chữ cái.
\end{question}

\begin{question}{16}
% [ID gốc: 1.22]
% [Nguồn: DBI202 Order 1 Q22]
Which of the following is not a feature of DBMS?
\begin{choices}
\choice Minimum duplication of data
\choice Redundancy of data
\choice Single-user access only
\choice Support ACID Property
\end{choices}
\textbf{Đáp án:} C
\textbf{Giải thích:} Đáp án theo dữ liệu đề: C — C. Single-user access only.

Nền tảng kiến thức: Chương này phân biệt dữ liệu, cơ sở dữ liệu, hệ quản trị CSDL (DBMS), mô hình dữ liệu, lược đồ và trạng thái dữ liệu. Cơ sở dữ liệu lưu dữ liệu có liên hệ; DBMS là phần mềm định nghĩa, truy vấn, cập nhật và bảo vệ dữ liệu đó.

Phân tích: Đối chiếu đúng thuộc tính ACID được hỏi: atomicity là tất-cả-hoặc-không; consistency giữ các ràng buộc; isolation kiểm soát ảnh hưởng giữa giao dịch đồng thời; durability bảo đảm giao dịch đã commit sống sót qua sự cố. COMMIT xác nhận, ROLLBACK hủy thay đổi chưa xác nhận.

Đối chiếu phương án: A. Minimum duplication of data; B. Redundancy of data; D. Support ACID Property. Các lựa chọn này không đáp ứng định nghĩa/điều kiện được hỏi; khi ôn, hãy chỉ ra chính xác điểm lệch của từng lựa chọn thay vì học thuộc mỗi chữ cái.
\end{question}

\begin{question}{17}
% [ID gốc: 1.24]
% [Nguồn: DBI202 Order 1 Q24]
What is the full form of DBMS?
\begin{choices}
\choice Data of Binary Management System
\choice Database Management System
\choice Database Management Service
\choice Data Backup Management System
\end{choices}
\textbf{Đáp án:} B
\textbf{Giải thích:} Đáp án theo dữ liệu đề: B — B. Database Management System.

Nền tảng kiến thức: Chương này phân biệt dữ liệu, cơ sở dữ liệu, hệ quản trị CSDL (DBMS), mô hình dữ liệu, lược đồ và trạng thái dữ liệu. Cơ sở dữ liệu lưu dữ liệu có liên hệ; DBMS là phần mềm định nghĩa, truy vấn, cập nhật và bảo vệ dữ liệu đó.

Phân tích: DBMS là viết tắt của Database Management System, tức hệ quản trị cơ sở dữ liệu. Nó cung cấp chức năng định nghĩa, lưu trữ, truy vấn, cập nhật và quản lý quyền trên CSDL; DBMS là phần mềm, còn database là dữ liệu được phần mềm quản lý.

Đối chiếu phương án: A. Data of Binary Management System; C. Database Management Service; D. Data Backup Management System. Các lựa chọn này không đáp ứng định nghĩa/điều kiện được hỏi; khi ôn, hãy chỉ ra chính xác điểm lệch của từng lựa chọn thay vì học thuộc mỗi chữ cái.
\end{question}

\begin{question}{18}
% [ID gốc: 1.26]
% [Nguồn: DBI202 Order 1 Q26]
A database administrator: a person or persons responsible for the struc-ture or schema of the database.
\begin{choices}
\choice True
\choice False
\end{choices}
\textbf{Đáp án:} A
\textbf{Giải thích:} Đáp án theo dữ liệu đề: A — A. True.

Nền tảng kiến thức: Chương này phân biệt dữ liệu, cơ sở dữ liệu, hệ quản trị CSDL (DBMS), mô hình dữ liệu, lược đồ và trạng thái dữ liệu. Cơ sở dữ liệu lưu dữ liệu có liên hệ; DBMS là phần mềm định nghĩa, truy vấn, cập nhật và bảo vệ dữ liệu đó.

Phân tích: Database administrator (DBA) chịu trách nhiệm quản trị và vận hành CSDL, bao gồm tổ chức schema, quyền truy cập, sao lưu/khôi phục và hiệu năng. DBA quản lý cấu trúc/hệ thống chứ không phải một bảng dữ liệu hay ngôn ngữ truy vấn.

Đối chiếu phương án: B. False. Các lựa chọn này không đáp ứng định nghĩa/điều kiện được hỏi; khi ôn, hãy chỉ ra chính xác điểm lệch của từng lựa chọn thay vì học thuộc mỗi chữ cái.
\end{question}

\begin{question}{19}
% [ID gốc: 1.27]
% [Nguồn: DBI202 Order 1 Q27]
What data structure will be represented for the hierarchical data model?
\begin{choices}
\choice Array structure
\choice Table structure
\choice Tree structure
\choice Pointers structure
\end{choices}
\textbf{Đáp án:} C
\textbf{Giải thích:} Đáp án theo dữ liệu đề: C — C. Tree structure.

Nền tảng kiến thức: Chương này phân biệt dữ liệu, cơ sở dữ liệu, hệ quản trị CSDL (DBMS), mô hình dữ liệu, lược đồ và trạng thái dữ liệu. Cơ sở dữ liệu lưu dữ liệu có liên hệ; DBMS là phần mềm định nghĩa, truy vấn, cập nhật và bảo vệ dữ liệu đó.

Phân tích: Mô hình đồ thị/mạng biểu diễn đối tượng thành nút và liên kết thành cạnh; mô hình phân cấp thường tổ chức quan hệ cha-con dạng cây. Hãy dựa vào cấu trúc liên kết được mô tả để phân biệt, không chỉ nhìn thấy từ “node”.

Đối chiếu phương án: A. Array structure; B. Table structure; D. Pointers structure. Các lựa chọn này không đáp ứng định nghĩa/điều kiện được hỏi; khi ôn, hãy chỉ ra chính xác điểm lệch của từng lựa chọn thay vì học thuộc mỗi chữ cái.
\end{question}

\begin{question}{20}
% [ID gốc: 1.29]
% [Nguồn: DBI202 Order 1 Q29]
In a network data model, a node can be linked to multiple other nodes. What is this relationship called?
\begin{choices}
\choice Many-to-one
\choice One-to-many
\choice Many-to-many
\choice One-to-one
\end{choices}
\textbf{Đáp án:} C
\textbf{Giải thích:} Đáp án theo dữ liệu đề: C — C. Many-to-many.

Lời giải gốc: Mô hình dữ liệu mạng (Network Data Model) cho phép một bản ghi con có nhiều bản ghi cha và ngược lại, tạo thành mối quan hệ nhiều-nhiều (Many-to-many).

Nền tảng kiến thức: Chương này phân biệt dữ liệu, cơ sở dữ liệu, hệ quản trị CSDL (DBMS), mô hình dữ liệu, lược đồ và trạng thái dữ liệu. Cơ sở dữ liệu lưu dữ liệu có liên hệ; DBMS là phần mềm định nghĩa, truy vấn, cập nhật và bảo vệ dữ liệu đó.

Phân tích: Mô hình đồ thị/mạng biểu diễn đối tượng thành nút và liên kết thành cạnh; mô hình phân cấp thường tổ chức quan hệ cha-con dạng cây. Hãy dựa vào cấu trúc liên kết được mô tả để phân biệt, không chỉ nhìn thấy từ “node”.

Đối chiếu phương án: A. Many-to-one; B. One-to-many; D. One-to-one. Các lựa chọn này không đáp ứng định nghĩa/điều kiện được hỏi; khi ôn, hãy chỉ ra chính xác điểm lệch của từng lựa chọn thay vì học thuộc mỗi chữ cái.
\end{question}

\begin{question}{21}
% [ID gốc: 1.32]
% [Nguồn: DBI202 Order 1 Q32]
Which of the following best defines a database?
\begin{choices}
\choice A collection of unorganized files
\choice A program used to analyze data
\choice A structured collection of data
\choice A spreadsheet
\end{choices}
\textbf{Đáp án:} C
\textbf{Giải thích:} Đáp án theo dữ liệu đề: C — C. A structured collection of data.

Lời giải gốc: Định nghĩa cơ bản nhất của Cơ sở dữ liệu (Database) là một tập hợp dữ liệu được tổ chức có cấu trúc, giúp việc lưu trữ, tìm kiếm và truy xuất thông tin hiệu quả.

Nền tảng kiến thức: Chương này phân biệt dữ liệu, cơ sở dữ liệu, hệ quản trị CSDL (DBMS), mô hình dữ liệu, lược đồ và trạng thái dữ liệu. Cơ sở dữ liệu lưu dữ liệu có liên hệ; DBMS là phần mềm định nghĩa, truy vấn, cập nhật và bảo vệ dữ liệu đó.

Phân tích: Cơ sở dữ liệu là tập dữ liệu/thông tin có liên hệ, được tổ chức và duy trì để phục vụ một mục đích; DBMS là phần mềm quản lý tập dữ liệu ấy. Vì vậy, database không đồng nghĩa với một tệp đơn lẻ, một lần chụp dữ liệu (instance) hay phần mềm quản trị.

Đối chiếu phương án: A. A collection of unorganized files; B. A program used to analyze data; D. A spreadsheet. Các lựa chọn này không đáp ứng định nghĩa/điều kiện được hỏi; khi ôn, hãy chỉ ra chính xác điểm lệch của từng lựa chọn thay vì học thuộc mỗi chữ cái.
\end{question}

\begin{question}{22}
% [ID gốc: 1.33]
% [Nguồn: DBI202 Order 1 Q33]
The people whose jobs require access to the database for querying, updating, and generating reports; the database primarily exists for their use. What type of user is this?
\begin{choices}
\choice DBA
\choice End users
\choice Database designer
\choice System Analysts and Application Programmers
\end{choices}
\textbf{Đáp án:} B
\textbf{Giải thích:} Đáp án theo dữ liệu đề: B — B. End users.

Lời giải gốc: Người dùng cuối (End users) là những người trực tiếp sử dụng cơ sở dữ liệu để truy vấn, cập nhật dữ liệu và tạo báo cáo phục vụ cho công việc hàng ngày của họ.

Nền tảng kiến thức: Chương này phân biệt dữ liệu, cơ sở dữ liệu, hệ quản trị CSDL (DBMS), mô hình dữ liệu, lược đồ và trạng thái dữ liệu. Cơ sở dữ liệu lưu dữ liệu có liên hệ; DBMS là phần mềm định nghĩa, truy vấn, cập nhật và bảo vệ dữ liệu đó.

Phân tích: OLTP ghi nhận giao dịch nghiệp vụ ngắn và cập nhật thường xuyên; decision support/OLAP tập trung truy vấn phân tích, tổng hợp và báo cáo. Hãy phân loại theo mục đích thao tác chứ không theo việc cả hai cùng đọc dữ liệu.

Đối chiếu phương án: A. DBA; C. Database designer; D. System Analysts and Application Programmers. Các lựa chọn này không đáp ứng định nghĩa/điều kiện được hỏi; khi ôn, hãy chỉ ra chính xác điểm lệch của từng lựa chọn thay vì học thuộc mỗi chữ cái.
\end{question}

\begin{question}{23}
% [ID gốc: 1.36]
% [Nguồn: DBI202 Order 1 Q36]
Which of the following is the appropriate characteristic of a database?
\begin{choices}
\choice Because a database is created to suit the format of the data, it cannot respond flexibly to data format changes.
\choice The procedure for making backups is complicated.
\choice It is difficult to share data between operations due to an exclusive control function.
\choice It can be accessed by multiple users at the same time due to an exclusive control function.
\end{choices}
\textbf{Đáp án:} D
\textbf{Giải thích:} Đáp án theo dữ liệu đề: D — D. It can be accessed by multiple users at the same time due to an exclusive control function..

Lời giải gốc: Đặc điểm quan trọng của Cơ sở dữ liệu (Database) là hỗ trợ tính đồng thời (Concurrency). Nhiều người dùng có thể truy cập cùng lúc nhờ vào các cơ chế khóa (locking) hoặc kiểm soát giao dịch (exclusive control) nhằm đảm bảo dữ liệu không bị sai lệch.

Nền tảng kiến thức: Chương này phân biệt dữ liệu, cơ sở dữ liệu, hệ quản trị CSDL (DBMS), mô hình dữ liệu, lược đồ và trạng thái dữ liệu. Cơ sở dữ liệu lưu dữ liệu có liên hệ; DBMS là phần mềm định nghĩa, truy vấn, cập nhật và bảo vệ dữ liệu đó.

Phân tích: Stored procedure là tập lệnh SQL có tên lưu trong CSDL, được gọi để thực thi một thủ tục; nó có thể nhận tham số và chứa nhiều câu lệnh. Đừng nhầm với trigger tự chạy khi có sự kiện hoặc function có quy tắc trả về giá trị.

Đối chiếu phương án: A. Because a database is created to suit the format of the data, it cannot respond flexibly to da…; B. The procedure for making backups is complicated.; C. It is difficult to share data between operations due to an exclusive control function.. Các lựa chọn này không đáp ứng định nghĩa/điều kiện được hỏi; khi ôn, hãy chỉ ra chính xác điểm lệch của từng lựa chọn thay vì học thuộc mỗi chữ cái.
\end{question}

\begin{question}{24}
% [ID gốc: 1.37]
% [Nguồn: DBI202 Order 1 Q37]
Which statement about the network data model is true?
\begin{choices}
\choice It supports high-level query languages like SQL.
\choice It only allows one parent per record.
\choice It organizes data using tables and relations.
\choice It allows each record to have multiple parents and children.
\end{choices}
\textbf{Đáp án:} D
\textbf{Giải thích:} Đáp án theo dữ liệu đề: D — D. It allows each record to have multiple parents and children..

Lời giải gốc: Đặc trưng của mô hình dữ liệu mạng là cho phép mỗi bản ghi có nhiều nút cha và nhiều nút con, phức tạp hơn mô hình phân cấp.

Nền tảng kiến thức: Chương này phân biệt dữ liệu, cơ sở dữ liệu, hệ quản trị CSDL (DBMS), mô hình dữ liệu, lược đồ và trạng thái dữ liệu. Cơ sở dữ liệu lưu dữ liệu có liên hệ; DBMS là phần mềm định nghĩa, truy vấn, cập nhật và bảo vệ dữ liệu đó.

Phân tích: Mô hình đồ thị/mạng biểu diễn đối tượng thành nút và liên kết thành cạnh; mô hình phân cấp thường tổ chức quan hệ cha-con dạng cây. Hãy dựa vào cấu trúc liên kết được mô tả để phân biệt, không chỉ nhìn thấy từ “node”.

Đối chiếu phương án: A. It supports high-level query languages like SQL.; B. It only allows one parent per record.; C. It organizes data using tables and relations.. Các lựa chọn này không đáp ứng định nghĩa/điều kiện được hỏi; khi ôn, hãy chỉ ra chính xác điểm lệch của từng lựa chọn thay vì học thuộc mỗi chữ cái.
\end{question}

\begin{question}{25}
% [ID gốc: 1.38]
% [Nguồn: DBI202 Order 1 Q38]
What is a disadvantage of the network data model?
\begin{choices}
\choice It does not support a high-level query language
\choice It does not support relationships between nodes
\choice It does not support the database management system
\choice It does not support storing very large amounts of data
\end{choices}
\textbf{Đáp án:} A
\textbf{Giải thích:} Đáp án theo dữ liệu đề: A — A. It does not support a high-level query language.

Lời giải gốc: Nhược điểm lớn của mô hình mạng là cấu trúc phức tạp và không hỗ trợ ngôn ngữ truy vấn cấp cao linh hoạt như SQL.

Nền tảng kiến thức: Chương này phân biệt dữ liệu, cơ sở dữ liệu, hệ quản trị CSDL (DBMS), mô hình dữ liệu, lược đồ và trạng thái dữ liệu. Cơ sở dữ liệu lưu dữ liệu có liên hệ; DBMS là phần mềm định nghĩa, truy vấn, cập nhật và bảo vệ dữ liệu đó.

Phân tích: Mô hình đồ thị/mạng biểu diễn đối tượng thành nút và liên kết thành cạnh; mô hình phân cấp thường tổ chức quan hệ cha-con dạng cây. Hãy dựa vào cấu trúc liên kết được mô tả để phân biệt, không chỉ nhìn thấy từ “node”.

Đối chiếu phương án: B. It does not support relationships between nodes; C. It does not support the database management system; D. It does not support storing very large amounts of data. Các lựa chọn này không đáp ứng định nghĩa/điều kiện được hỏi; khi ôn, hãy chỉ ra chính xác điểm lệch của từng lựa chọn thay vì học thuộc mỗi chữ cái.
\end{question}

\begin{question}{26}
% [ID gốc: 1.52]
% [Nguồn: DBI202 Order 1 Q52]
A database administrator is a person or persons responsible for the structure or schema of the database.
\begin{choices}
\choice True
\choice False
\end{choices}
\textbf{Đáp án:} A
\textbf{Giải thích:} Đáp án theo dữ liệu đề: A — A. True.

Lời giải gốc: Quản trị viên cơ sở dữ liệu (DBA) là người chịu trách nhiệm quản lý cấu trúc, bảo mật và lược đồ cơ sở dữ liệu.

Nền tảng kiến thức: Chương này phân biệt dữ liệu, cơ sở dữ liệu, hệ quản trị CSDL (DBMS), mô hình dữ liệu, lược đồ và trạng thái dữ liệu. Cơ sở dữ liệu lưu dữ liệu có liên hệ; DBMS là phần mềm định nghĩa, truy vấn, cập nhật và bảo vệ dữ liệu đó.

Phân tích: Database administrator (DBA) chịu trách nhiệm quản trị và vận hành CSDL, bao gồm tổ chức schema, quyền truy cập, sao lưu/khôi phục và hiệu năng. DBA quản lý cấu trúc/hệ thống chứ không phải một bảng dữ liệu hay ngôn ngữ truy vấn.

Đối chiếu phương án: B. False. Các lựa chọn này không đáp ứng định nghĩa/điều kiện được hỏi; khi ôn, hãy chỉ ra chính xác điểm lệch của từng lựa chọn thay vì học thuộc mỗi chữ cái.
\end{question}

\begin{question}{27}
% [ID gốc: 1.53]
% [Nguồn: DBI202 Order 1 Q53]
Which of the following is NOT a feature of a DBMS?
\begin{choices}
\choice Minimum duplication of data
\choice Redundancy of data
\choice Single-user access only
\choice Support for the ACID property
\end{choices}
\textbf{Đáp án:} C
\textbf{Giải thích:} Đáp án theo dữ liệu đề: C — C. Single-user access only.

Lời giải gốc: Hệ quản trị cơ sở dữ liệu được thiết kế để hỗ trợ đa người dùng truy cập đồng thời, do đó "Single-user access only" không phải là đặc điểm của DBMS.

Nền tảng kiến thức: Chương này phân biệt dữ liệu, cơ sở dữ liệu, hệ quản trị CSDL (DBMS), mô hình dữ liệu, lược đồ và trạng thái dữ liệu. Cơ sở dữ liệu lưu dữ liệu có liên hệ; DBMS là phần mềm định nghĩa, truy vấn, cập nhật và bảo vệ dữ liệu đó.

Phân tích: Đối chiếu đúng thuộc tính ACID được hỏi: atomicity là tất-cả-hoặc-không; consistency giữ các ràng buộc; isolation kiểm soát ảnh hưởng giữa giao dịch đồng thời; durability bảo đảm giao dịch đã commit sống sót qua sự cố. COMMIT xác nhận, ROLLBACK hủy thay đổi chưa xác nhận.

Đối chiếu phương án: A. Minimum duplication of data; B. Redundancy of data; D. Support for the ACID property. Các lựa chọn này không đáp ứng định nghĩa/điều kiện được hỏi; khi ôn, hãy chỉ ra chính xác điểm lệch của từng lựa chọn thay vì học thuộc mỗi chữ cái.
\end{question}

\begin{question}{28}
% [ID gốc: 1.59]
% [Nguồn: DBI202 Order 1 Q59]
In an employee table, to include the attributes whose value always have some value, which of the following constraints must be used?
\begin{choices}
\choice Null
\choice Not null
\choice Unique
\choice Distinct
\end{choices}
\textbf{Đáp án:} B
\textbf{Giải thích:} Đáp án theo dữ liệu đề: B — B. Not null.

Lời giải gốc: Ràng buộc NOT NULL bắt buộc một cột phải luôn chứa giá trị, không được để trống.

Nền tảng kiến thức: Chương này phân biệt dữ liệu, cơ sở dữ liệu, hệ quản trị CSDL (DBMS), mô hình dữ liệu, lược đồ và trạng thái dữ liệu. Cơ sở dữ liệu lưu dữ liệu có liên hệ; DBMS là phần mềm định nghĩa, truy vấn, cập nhật và bảo vệ dữ liệu đó.

Phân tích: NOT NULL là constraint cấm giá trị NULL ở cột, phù hợp khi mỗi hàng bắt buộc phải có giá trị. Nó không tự bảo đảm giá trị duy nhất (UNIQUE/PRIMARY KEY) và cũng không kiểm tra miền/range như CHECK.

Đối chiếu phương án: A. Null; C. Unique; D. Distinct. Các lựa chọn này không đáp ứng định nghĩa/điều kiện được hỏi; khi ôn, hãy chỉ ra chính xác điểm lệch của từng lựa chọn thay vì học thuộc mỗi chữ cái.
\end{question}

\begin{question}{29}
% [ID gốc: 1.60]
% [Nguồn: DBI202 Order 1 Q60]
What data structure represents the hierarchical data model?
\begin{choices}
\choice Array structure
\choice Table structure
\choice Tree structure
\choice Pointer structure
\end{choices}
\textbf{Đáp án:} C
\textbf{Giải thích:} Đáp án theo dữ liệu đề: C — C. Tree structure.

Lời giải gốc: Mô hình dữ liệu phân cấp sử dụng cấu trúc cây (tree structure) để tổ chức mối quan hệ cha-con giữa các dữ liệu.

Nền tảng kiến thức: Chương này phân biệt dữ liệu, cơ sở dữ liệu, hệ quản trị CSDL (DBMS), mô hình dữ liệu, lược đồ và trạng thái dữ liệu. Cơ sở dữ liệu lưu dữ liệu có liên hệ; DBMS là phần mềm định nghĩa, truy vấn, cập nhật và bảo vệ dữ liệu đó.

Phân tích: Mô hình đồ thị/mạng biểu diễn đối tượng thành nút và liên kết thành cạnh; mô hình phân cấp thường tổ chức quan hệ cha-con dạng cây. Hãy dựa vào cấu trúc liên kết được mô tả để phân biệt, không chỉ nhìn thấy từ “node”.

Đối chiếu phương án: A. Array structure; B. Table structure; D. Pointer structure. Các lựa chọn này không đáp ứng định nghĩa/điều kiện được hỏi; khi ôn, hãy chỉ ra chính xác điểm lệch của từng lựa chọn thay vì học thuộc mỗi chữ cái.
\end{question}

\begin{question}{30}
% [ID gốc: 1.61]
% [Nguồn: DBI202 Order 1 Q61]
Which of the following is an appropriate characteristic of a database?
\begin{choices}
\choice Because a database is created to suit the format of the data, it cannot respond flexibly to data format changes.
\choice The procedure for making backups is complicated.
\choice It is difficult to share data between operations due to an exclusive control function.
\choice It can be accessed by multiple users at the same time due to an exclusive control function.
\end{choices}
\textbf{Đáp án:} D
\textbf{Giải thích:} Đáp án theo dữ liệu đề: D — D. It can be accessed by multiple users at the same time due to an exclusive control function..

Lời giải gốc: Cơ sở dữ liệu cho phép nhiều người dùng truy cập đồng thời nhờ các cơ chế điều khiển tranh chấp và kiểm soát đồng thời.

Nền tảng kiến thức: Chương này phân biệt dữ liệu, cơ sở dữ liệu, hệ quản trị CSDL (DBMS), mô hình dữ liệu, lược đồ và trạng thái dữ liệu. Cơ sở dữ liệu lưu dữ liệu có liên hệ; DBMS là phần mềm định nghĩa, truy vấn, cập nhật và bảo vệ dữ liệu đó.

Phân tích: Stored procedure là tập lệnh SQL có tên lưu trong CSDL, được gọi để thực thi một thủ tục; nó có thể nhận tham số và chứa nhiều câu lệnh. Đừng nhầm với trigger tự chạy khi có sự kiện hoặc function có quy tắc trả về giá trị.

Đối chiếu phương án: A. Because a database is created to suit the format of the data, it cannot respond flexibly to da…; B. The procedure for making backups is complicated.; C. It is difficult to share data between operations due to an exclusive control function.. Các lựa chọn này không đáp ứng định nghĩa/điều kiện được hỏi; khi ôn, hãy chỉ ra chính xác điểm lệch của từng lựa chọn thay vì học thuộc mỗi chữ cái.
\end{question}

\begin{question}{31}
% [ID gốc: 1.62]
% [Nguồn: DBI202 Order 1 Q62]
A _ _ _ is a logically coherent collection of data with some inherent meaning, representing some aspect of the real world, and which is designed, built, and populated with data for a specific purpose.
\begin{choices}
\choice Database
\choice Schema Instance
\choice Schema
\choice Database Instance
\end{choices}
\textbf{Đáp án:} A
\textbf{Giải thích:} Đáp án theo dữ liệu đề: A — A. Database.

Lời giải gốc: Khái niệm này định nghĩa về Cơ sở dữ liệu (Database). ---

Nền tảng kiến thức: Chương này phân biệt dữ liệu, cơ sở dữ liệu, hệ quản trị CSDL (DBMS), mô hình dữ liệu, lược đồ và trạng thái dữ liệu. Cơ sở dữ liệu lưu dữ liệu có liên hệ; DBMS là phần mềm định nghĩa, truy vấn, cập nhật và bảo vệ dữ liệu đó.

Phân tích: Cơ sở dữ liệu là tập dữ liệu/thông tin có liên hệ, được tổ chức và duy trì để phục vụ một mục đích; DBMS là phần mềm quản lý tập dữ liệu ấy. Vì vậy, database không đồng nghĩa với một tệp đơn lẻ, một lần chụp dữ liệu (instance) hay phần mềm quản trị.

Đối chiếu phương án: B. Schema Instance; C. Schema; D. Database Instance. Các lựa chọn này không đáp ứng định nghĩa/điều kiện được hỏi; khi ôn, hãy chỉ ra chính xác điểm lệch của từng lựa chọn thay vì học thuộc mỗi chữ cái.
\end{question}

\begin{question}{32}
% [ID gốc: 1.63]
% [Nguồn: DBI202 Order 1 Q63]
A. person or persons responsible for the structure or schema of the database is _ _ _ _ _.
\begin{choices}
\choice A User
\choice A Coder
\choice A Customer
\choice A database administrator
\end{choices}
\textbf{Đáp án:} D
\textbf{Giải thích:} Đáp án theo dữ liệu đề: D — D. A database administrator.

Lời giải gốc: Người quản trị cơ sở dữ liệu (Database Administrator - DBA) chịu trách nhiệm chính trong việc thiết kế, quản lý cấu trúc, sơ đồ (schema) và vận hành cơ sở dữ liệu.

Nền tảng kiến thức: Chương này phân biệt dữ liệu, cơ sở dữ liệu, hệ quản trị CSDL (DBMS), mô hình dữ liệu, lược đồ và trạng thái dữ liệu. Cơ sở dữ liệu lưu dữ liệu có liên hệ; DBMS là phần mềm định nghĩa, truy vấn, cập nhật và bảo vệ dữ liệu đó.

Phân tích: Database administrator (DBA) chịu trách nhiệm quản trị và vận hành CSDL, bao gồm tổ chức schema, quyền truy cập, sao lưu/khôi phục và hiệu năng. DBA quản lý cấu trúc/hệ thống chứ không phải một bảng dữ liệu hay ngôn ngữ truy vấn.

Đối chiếu phương án: A. A User; B. A Coder; C. A Customer. Các lựa chọn này không đáp ứng định nghĩa/điều kiện được hỏi; khi ôn, hãy chỉ ra chính xác điểm lệch của từng lựa chọn thay vì học thuộc mỗi chữ cái.
\end{question}

\begin{question}{33}
% [ID gốc: 1.65]
% [Nguồn: DBI202 Order 1 Q65]
Which of the following is a function of a DBMS?
\begin{choices}
\choice Designing web pages
\choice Give users the ability to query the data
\choice Performing system backups
\choice Compiling source code
\end{choices}
\textbf{Đáp án:} B
\textbf{Giải thích:} Đáp án theo dữ liệu đề: B — B. Give users the ability to query the data.

Lời giải gốc: Một trong những chức năng cốt lõi của Hệ quản trị cơ sở dữ liệu (DBMS) là cung cấp ngôn ngữ và giao diện cho phép người dùng truy vấn (query) và quản lý dữ liệu.

Nền tảng kiến thức: Chương này phân biệt dữ liệu, cơ sở dữ liệu, hệ quản trị CSDL (DBMS), mô hình dữ liệu, lược đồ và trạng thái dữ liệu. Cơ sở dữ liệu lưu dữ liệu có liên hệ; DBMS là phần mềm định nghĩa, truy vấn, cập nhật và bảo vệ dữ liệu đó.

Phân tích: Đọc truy vấn theo các giai đoạn logic: FROM/JOIN tạo dữ liệu nguồn, WHERE lọc từng hàng, GROUP BY tạo nhóm, HAVING lọc nhóm, SELECT xác định biểu thức đầu ra và ORDER BY sắp kết quả. Với phép join, kiểm tra điều kiện ghép và các hàng không khớp trước khi tính kết quả.

Đối chiếu phương án: A. Designing web pages; C. Performing system backups; D. Compiling source code. Các lựa chọn này không đáp ứng định nghĩa/điều kiện được hỏi; khi ôn, hãy chỉ ra chính xác điểm lệch của từng lựa chọn thay vì học thuộc mỗi chữ cái.
\end{question}

\begin{question}{34}
% [ID gốc: 2.101]
% [Nguồn: DBI202 Order 2 Q101]
Assume that you are required to develop an application to manage courses at the university, which requires: Structured data (tables, rows, columns), strict consistency, and the need to use complex queries. The enforcement of ACID properties (Atomicity, Consistency, Isolation, Durability). In the context of databases, which most appropriate data model should you choose for developing this application?
\begin{choices}
\choice Relational data model (example: Microsoft SQL Server, MySQL, ..)
\choice Document-based data model (example: MongoDB, Cassandra, Couch DB)
\choice Key-value data model (example: Redis, DynamoDB, Riak)
\choice Hierarchical data model
\end{choices}
\textbf{Đáp án:} A
\textbf{Giải thích:} Đáp án theo dữ liệu đề: A — A. Relational data model (example: Microsoft SQL Server, MySQL, ..).

Lời giải gốc: Dữ liệu có cấu trúc dạng bảng, yêu cầu tính nhất quán cao, truy vấn phức tạp và tuân thủ tính chất ACID là đặc trưng chính của Mô hình dữ liệu quan hệ (Relational data model).

Nền tảng kiến thức: Mô hình ER mô tả yêu cầu ở mức khái niệm: thực thể là đối tượng cần lưu, thuộc tính mô tả thực thể hoặc mối quan hệ, còn lực lượng (cardinality) và participation diễn tả các ràng buộc kết hợp. Khóa của thực thể yếu thường ghép khóa bộ phận với khóa của thực thể chủ.

Phân tích: Data model là bộ khái niệm dùng để mô tả cấu trúc dữ liệu, quan hệ giữa dữ liệu và các ràng buộc. Relational, hierarchical, network và graph là các mô hình có thật; cần chọn mô hình khớp đặc điểm được hỏi, không nhầm mô hình với DBMS hay ngôn ngữ thao tác.

Đối chiếu phương án: B. Document-based data model (example: MongoDB, Cassandra, Couch DB); C. Key-value data model (example: Redis, DynamoDB, Riak); D. Hierarchical data model. Các lựa chọn này không đáp ứng định nghĩa/điều kiện được hỏi; khi ôn, hãy chỉ ra chính xác điểm lệch của từng lựa chọn thay vì học thuộc mỗi chữ cái.
\end{question}

\begin{question}{35}
% [ID gốc: 3.14]
% [Nguồn: DBI202 Order 3 Q14]
When the Relational Data Model first proposed?
\begin{choices}
\choice 1940
\choice 1950
\choice 1960
\choice 1970
\end{choices}
\textbf{Đáp án:} D
\textbf{Giải thích:} Đáp án theo dữ liệu đề: D — D. 1970.

Nền tảng kiến thức: Mô hình quan hệ biểu diễn dữ liệu bằng quan hệ (bảng), bộ (hàng), thuộc tính (cột) và miền giá trị. Khóa nhận diện bộ; khóa ngoại và các ràng buộc toàn vẹn ngăn dữ liệu tham chiếu sai hoặc vi phạm miền.

Phân tích: Data model là bộ khái niệm dùng để mô tả cấu trúc dữ liệu, quan hệ giữa dữ liệu và các ràng buộc. Relational, hierarchical, network và graph là các mô hình có thật; cần chọn mô hình khớp đặc điểm được hỏi, không nhầm mô hình với DBMS hay ngôn ngữ thao tác.

Đối chiếu phương án: A. 1940; B. 1950; C. 1960. Các lựa chọn này không đáp ứng định nghĩa/điều kiện được hỏi; khi ôn, hãy chỉ ra chính xác điểm lệch của từng lựa chọn thay vì học thuộc mỗi chữ cái.
\end{question}

\begin{question}{36}
% [ID gốc: 13.1]
% [Nguồn: DBI202 Order 13 Q1]
In query compiler, which unit transforms the initial query plan into the best available sequence of operations on the actual data?
\begin{choices}
\choice A query parser
\choice A query preprocessor
\choice A query optimizer
\choice A query processor
\end{choices}
\textbf{Đáp án:} C
\textbf{Giải thích:} Đáp án theo dữ liệu đề: C — C. A query optimizer.

Nền tảng kiến thức: Tối ưu hóa truy vấn biến biểu thức logic thành các kế hoạch tương đương rồi ước lượng chi phí để chọn kế hoạch có khả năng rẻ hơn. Đẩy selection/projection sớm và dùng thống kê có thể giảm dữ liệu trung gian, nhưng kế hoạch cụ thể phụ thuộc hệ quản trị và dữ liệu.

Phân tích: Đọc truy vấn theo các giai đoạn logic: FROM/JOIN tạo dữ liệu nguồn, WHERE lọc từng hàng, GROUP BY tạo nhóm, HAVING lọc nhóm, SELECT xác định biểu thức đầu ra và ORDER BY sắp kết quả. Với phép join, kiểm tra điều kiện ghép và các hàng không khớp trước khi tính kết quả.

Đối chiếu phương án: A. A query parser; B. A query preprocessor; D. A query processor. Các lựa chọn này không đáp ứng định nghĩa/điều kiện được hỏi; khi ôn, hãy chỉ ra chính xác điểm lệch của từng lựa chọn thay vì học thuộc mỗi chữ cái.
\end{question}

\begin{question}{37}
% [ID gốc: 13.2]
% [Nguồn: DBI202 Order 13 Q2]
What is query processing?
\begin{choices}
\choice The process of executing SQL statements.
\choice The process of defining database schemas.
\choice The process of designing user interfaces.
\choice The process of backing up data.
\end{choices}
\textbf{Đáp án:} A
\textbf{Giải thích:} Đáp án theo dữ liệu đề: A — A. The process of executing SQL statements..

Lời giải gốc: Xử lý truy vấn (query processing) là quá trình biên dịch và thực thi các câu lệnh SQL trong hệ thống cơ sở dữ liệu. ---

Nền tảng kiến thức: Tối ưu hóa truy vấn biến biểu thức logic thành các kế hoạch tương đương rồi ước lượng chi phí để chọn kế hoạch có khả năng rẻ hơn. Đẩy selection/projection sớm và dùng thống kê có thể giảm dữ liệu trung gian, nhưng kế hoạch cụ thể phụ thuộc hệ quản trị và dữ liệu.

Phân tích: Đọc truy vấn theo các giai đoạn logic: FROM/JOIN tạo dữ liệu nguồn, WHERE lọc từng hàng, GROUP BY tạo nhóm, HAVING lọc nhóm, SELECT xác định biểu thức đầu ra và ORDER BY sắp kết quả. Với phép join, kiểm tra điều kiện ghép và các hàng không khớp trước khi tính kết quả.

Đối chiếu phương án: B. The process of defining database schemas.; C. The process of designing user interfaces.; D. The process of backing up data.. Các lựa chọn này không đáp ứng định nghĩa/điều kiện được hỏi; khi ôn, hãy chỉ ra chính xác điểm lệch của từng lựa chọn thay vì học thuộc mỗi chữ cái.
\end{question}

\begin{question}{38}
% [ID gốc: 13.4]
% [Nguồn: DBI202 Order 13 Q4]
In a query compiler, which unit transforms the initial query plan into the best available sequence of operations on the actual data?
\begin{choices}
\choice A query parser
\choice A query preprocessor
\choice A query optimizer
\choice A query processor
\end{choices}
\textbf{Đáp án:} C
\textbf{Giải thích:} Đáp án theo dữ liệu đề: C — C. A query optimizer.

Lời giải gốc: Bộ tối ưu hóa truy vấn (Query optimizer) chịu trách nhiệm phân tích và chuyển kế hoạch truy vấn ban đầu thành chuỗi thao tác thực thi tối ưu nhất.

Nền tảng kiến thức: Tối ưu hóa truy vấn biến biểu thức logic thành các kế hoạch tương đương rồi ước lượng chi phí để chọn kế hoạch có khả năng rẻ hơn. Đẩy selection/projection sớm và dùng thống kê có thể giảm dữ liệu trung gian, nhưng kế hoạch cụ thể phụ thuộc hệ quản trị và dữ liệu.

Phân tích: Đọc truy vấn theo các giai đoạn logic: FROM/JOIN tạo dữ liệu nguồn, WHERE lọc từng hàng, GROUP BY tạo nhóm, HAVING lọc nhóm, SELECT xác định biểu thức đầu ra và ORDER BY sắp kết quả. Với phép join, kiểm tra điều kiện ghép và các hàng không khớp trước khi tính kết quả.

Đối chiếu phương án: A. A query parser; B. A query preprocessor; D. A query processor. Các lựa chọn này không đáp ứng định nghĩa/điều kiện được hỏi; khi ôn, hãy chỉ ra chính xác điểm lệch của từng lựa chọn thay vì học thuộc mỗi chữ cái.
\end{question}

\begin{question}{39}
% [ID gốc: 14.1]
% [Nguồn: DBI202 Order 14 Q1]
The query is parsed and optimized by a query compiler. The resulting query plan, or sequence of actions the DBMS will perform to answer the query, is passed to the execution engine.
QN=6 (6798) Today, in a modern DBMS, the queries are parsed and optimized by
\begin{choices}
\choice all of the others
\choice execution engine
\choice buffer manager
\choice query compiler
\end{choices}
\textbf{Đáp án:} D
\textbf{Giải thích:} Đáp án theo dữ liệu đề: D — D. query compiler.

Nền tảng kiến thức: Trong pipeline truy vấn, parser kiểm tra cú pháp và dựng cây; các bước phân tích ngữ nghĩa/tiền xử lý kiểm tra tên và kiểu; optimizer chọn kế hoạch; execution engine thực thi kế hoạch. Không nên nhầm vai trò biên dịch với vai trò chạy truy vấn.

Phân tích: Đọc truy vấn theo các giai đoạn logic: FROM/JOIN tạo dữ liệu nguồn, WHERE lọc từng hàng, GROUP BY tạo nhóm, HAVING lọc nhóm, SELECT xác định biểu thức đầu ra và ORDER BY sắp kết quả. Với phép join, kiểm tra điều kiện ghép và các hàng không khớp trước khi tính kết quả.

Đối chiếu phương án: A. all of the others; B. execution engine; C. buffer manager. Các lựa chọn này không đáp ứng định nghĩa/điều kiện được hỏi; khi ôn, hãy chỉ ra chính xác điểm lệch của từng lựa chọn thay vì học thuộc mỗi chữ cái.
\end{question}

\begin{question}{40}
% [ID gốc: 14.2]
% [Nguồn: DBI202 Order 14 Q2]
In query compiler, which unit builds a tree structure from the text form of the query?
\begin{choices}
\choice A query parser
\choice A query preprocessor
\choice A query optimizer
\choice A query processor
\end{choices}
\textbf{Đáp án:} A
\textbf{Giải thích:} Đáp án theo dữ liệu đề: A — A. A query parser.

Nền tảng kiến thức: Trong pipeline truy vấn, parser kiểm tra cú pháp và dựng cây; các bước phân tích ngữ nghĩa/tiền xử lý kiểm tra tên và kiểu; optimizer chọn kế hoạch; execution engine thực thi kế hoạch. Không nên nhầm vai trò biên dịch với vai trò chạy truy vấn.

Phân tích: Đọc truy vấn theo các giai đoạn logic: FROM/JOIN tạo dữ liệu nguồn, WHERE lọc từng hàng, GROUP BY tạo nhóm, HAVING lọc nhóm, SELECT xác định biểu thức đầu ra và ORDER BY sắp kết quả. Với phép join, kiểm tra điều kiện ghép và các hàng không khớp trước khi tính kết quả.

Đối chiếu phương án: B. A query preprocessor; C. A query optimizer; D. A query processor. Các lựa chọn này không đáp ứng định nghĩa/điều kiện được hỏi; khi ôn, hãy chỉ ra chính xác điểm lệch của từng lựa chọn thay vì học thuộc mỗi chữ cái.
\end{question}

\begin{question}{41}
% [ID gốc: 14.4]
% [Nguồn: DBI202 Order 14 Q4]
In a query compiler, which unit builds a tree structure from the text form of the query?
\begin{choices}
\choice A query parser
\choice A query preprocessor
\choice A query optimizer
\choice A query processor
\end{choices}
\textbf{Đáp án:} A
\textbf{Giải thích:} Đáp án theo dữ liệu đề: A — A. A query parser.

Lời giải gốc: Bộ phân tích cú pháp truy vấn (query parser) chịu trách nhiệm kiểm tra câu lệnh và xây dựng cấu trúc cây cú pháp từ văn bản truy vấn.

Nền tảng kiến thức: Trong pipeline truy vấn, parser kiểm tra cú pháp và dựng cây; các bước phân tích ngữ nghĩa/tiền xử lý kiểm tra tên và kiểu; optimizer chọn kế hoạch; execution engine thực thi kế hoạch. Không nên nhầm vai trò biên dịch với vai trò chạy truy vấn.

Phân tích: Đọc truy vấn theo các giai đoạn logic: FROM/JOIN tạo dữ liệu nguồn, WHERE lọc từng hàng, GROUP BY tạo nhóm, HAVING lọc nhóm, SELECT xác định biểu thức đầu ra và ORDER BY sắp kết quả. Với phép join, kiểm tra điều kiện ghép và các hàng không khớp trước khi tính kết quả.

Đối chiếu phương án: B. A query preprocessor; C. A query optimizer; D. A query processor. Các lựa chọn này không đáp ứng định nghĩa/điều kiện được hỏi; khi ôn, hãy chỉ ra chính xác điểm lệch của từng lựa chọn thay vì học thuộc mỗi chữ cái.
\end{question}

\begin{question}{42}
% [ID gốc: 25.1]
% [Nguồn: DBI202 Order 25 Q1]
The aim for the Object-relational database is to bridge the gap between conceptual data modelling techniques such as Entity-relationship diagram (ERD) and object-relational mapping (ORM), which often use classes and inheritance, and relational databases, which do not directly support them
\begin{choices}
\choice True
\choice False
\end{choices}
\textbf{Đáp án:} A
\textbf{Giải thích:} Đáp án theo dữ liệu đề: A — A. True.

Nền tảng kiến thức: CSDL đối tượng và đối tượng-quan hệ đưa khái niệm kiểu do người dùng định nghĩa, giá trị phức hợp hoặc đặc tính đối tượng vào hệ quản trị quan hệ. Mục tiêu là biểu diễn kiểu ứng dụng gần hơn mà vẫn tận dụng lưu trữ/truy vấn quan hệ.

Phân tích: Khi đọc mô hình ER, tách ba thứ: entity set là loại đối tượng, relationship set là liên kết giữa các loại đối tượng, còn cardinality/participation là ràng buộc số lượng hoặc mức tham gia. Đáp án đúng phải gọi đúng thành phần mà câu hỏi mô tả.

Đối chiếu phương án: B. False. Các lựa chọn này không đáp ứng định nghĩa/điều kiện được hỏi; khi ôn, hãy chỉ ra chính xác điểm lệch của từng lựa chọn thay vì học thuộc mỗi chữ cái.
\end{question}

\begin{question}{43}
% [ID gốc: 28.1]
% [Nguồn: DBI202 Order 28 Q1]
Which of the following is an example of a NoSQL database system that is not based on the relational model?
\begin{choices}
\choice MySQL
\choice PostgreSQL
\choice MongoDB
\choice Oracle Database
\end{choices}
\textbf{Đáp án:} C
\textbf{Giải thích:} Đáp án theo dữ liệu đề: C — C. MongoDB.

Lời giải gốc: MongoDB là hệ quản trị cơ sở dữ liệu phi quan hệ (NoSQL) lưu trữ dữ liệu theo tài liệu (document-oriented).

Nền tảng kiến thức: Chương bổ sung gom các chủ đề mở rộng như CSDL thời gian, di động, đa phương tiện, bộ nhớ chính và tích hợp nhiều nguồn. Trả lời cần dựa vào đặc điểm chuyên biệt được nêu trong câu, không áp máy móc giả định của bảng quan hệ truyền thống.

Phân tích: MongoDB là hệ quản trị NoSQL thiên về tài liệu JSON/BSON, nên không dựa vào quan hệ bảng hàng-cột như MySQL, PostgreSQL hay Oracle. “NoSQL” không có nghĩa là không thể truy vấn; nó chỉ chỉ một nhóm mô hình lưu trữ khác quan hệ truyền thống.

Đối chiếu phương án: A. MySQL; B. PostgreSQL; D. Oracle Database. Các lựa chọn này không đáp ứng định nghĩa/điều kiện được hỏi; khi ôn, hãy chỉ ra chính xác điểm lệch của từng lựa chọn thay vì học thuộc mỗi chữ cái.
\end{question}


\newpage
\section*{Chương 2: The Relational Data Model}

\begin{question}{44}
% [ID gốc: 1.8]
% [Nguồn: DBI202 Order 1 Q8]
Today, the two important data models are:
\begin{choices}
\choice The semi-structured data model (including XML and related standards) and the network data model
\choice The relational data model (including object-relational extensions) and the semi-structured data model (including XML and related standards)
\choice The network data model and the flat file data model
\choice The flat file data model and the semi-structured data model (including XML and related standards)
\end{choices}
\textbf{Đáp án:} B
\textbf{Giải thích:} Đáp án theo dữ liệu đề: B — B. The relational data model (including object-relational extensions) and the semi-structured data model (including XML and related standards).

Nền tảng kiến thức: Chương này phân biệt dữ liệu, cơ sở dữ liệu, hệ quản trị CSDL (DBMS), mô hình dữ liệu, lược đồ và trạng thái dữ liệu. Cơ sở dữ liệu lưu dữ liệu có liên hệ; DBMS là phần mềm định nghĩa, truy vấn, cập nhật và bảo vệ dữ liệu đó.

Phân tích: Mô hình đồ thị/mạng biểu diễn đối tượng thành nút và liên kết thành cạnh; mô hình phân cấp thường tổ chức quan hệ cha-con dạng cây. Hãy dựa vào cấu trúc liên kết được mô tả để phân biệt, không chỉ nhìn thấy từ “node”.

Đối chiếu phương án: A. The semi-structured data model (including XML and related standards) and the network data model; C. The network data model and the flat file data model; D. The flat file data model and the semi-structured data model (including XML and related standar…. Các lựa chọn này không đáp ứng định nghĩa/điều kiện được hỏi; khi ôn, hãy chỉ ra chính xác điểm lệch của từng lựa chọn thay vì học thuộc mỗi chữ cái.
\end{question}

\begin{question}{45}
% [ID gốc: 1.21]
% [Nguồn: DBI202 Order 1 Q21]
A data model is a notation for describing data or information. The description generally consists ?
\begin{choices}
\choice Structure of the data
\choice Operations on the data
\choice Constraints on the data.
\choice All of the others
\end{choices}
\textbf{Đáp án:} D
\textbf{Giải thích:} Đáp án theo dữ liệu đề: D — D. All of the others.

Nền tảng kiến thức: Chương này phân biệt dữ liệu, cơ sở dữ liệu, hệ quản trị CSDL (DBMS), mô hình dữ liệu, lược đồ và trạng thái dữ liệu. Cơ sở dữ liệu lưu dữ liệu có liên hệ; DBMS là phần mềm định nghĩa, truy vấn, cập nhật và bảo vệ dữ liệu đó.

Phân tích: Data model là bộ khái niệm dùng để mô tả cấu trúc dữ liệu, quan hệ giữa dữ liệu và các ràng buộc. Relational, hierarchical, network và graph là các mô hình có thật; cần chọn mô hình khớp đặc điểm được hỏi, không nhầm mô hình với DBMS hay ngôn ngữ thao tác.

Đối chiếu phương án: A. Structure of the data; B. Operations on the data; C. Constraints on the data.. Các lựa chọn này không đáp ứng định nghĩa/điều kiện được hỏi; khi ôn, hãy chỉ ra chính xác điểm lệch của từng lựa chọn thay vì học thuộc mỗi chữ cái.
\end{question}

\begin{question}{46}
% [ID gốc: 1.40]
% [Nguồn: DBI202 Order 1 Q40]
Which data model introduced the concept of tables with rows and columns?
\begin{choices}
\choice Hierarchical Model
\choice Network Model
\choice Object-Oriented Model
\choice Relational Model
\end{choices}
\textbf{Đáp án:} D
\textbf{Giải thích:} Đáp án theo dữ liệu đề: D — D. Relational Model.

Lời giải gốc: Mô hình quan hệ (Relational Model) giới thiệu cách tổ chức dữ liệu dưới dạng các bảng (tables) gồm các dòng (rows/tuples) và các cột (columns/attributes).

Nền tảng kiến thức: Chương này phân biệt dữ liệu, cơ sở dữ liệu, hệ quản trị CSDL (DBMS), mô hình dữ liệu, lược đồ và trạng thái dữ liệu. Cơ sở dữ liệu lưu dữ liệu có liên hệ; DBMS là phần mềm định nghĩa, truy vấn, cập nhật và bảo vệ dữ liệu đó.

Phân tích: Data model là bộ khái niệm dùng để mô tả cấu trúc dữ liệu, quan hệ giữa dữ liệu và các ràng buộc. Relational, hierarchical, network và graph là các mô hình có thật; cần chọn mô hình khớp đặc điểm được hỏi, không nhầm mô hình với DBMS hay ngôn ngữ thao tác.

Đối chiếu phương án: A. Hierarchical Model; B. Network Model; C. Object-Oriented Model. Các lựa chọn này không đáp ứng định nghĩa/điều kiện được hỏi; khi ôn, hãy chỉ ra chính xác điểm lệch của từng lựa chọn thay vì học thuộc mỗi chữ cái.
\end{question}

\begin{question}{47}
% [ID gốc: 1.42]
% [Nguồn: DBI202 Order 1 Q42]
In the relational model of data, a data model is a notation for describing data or information. What are the three descriptions of data and information? Choose three answers.
\begin{choices}
\choice Structure
\choice Relationship
\choice Constraints
\choice Operations
\end{choices}
\textbf{Đáp án:} ACD
\textbf{Giải thích:} Đáp án theo dữ liệu đề: A, C, D — A. Structure; C. Constraints; D. Operations.

Lời giải gốc: Một mô hình dữ liệu bao gồm ba khía cạnh mô tả chính: cấu trúc dữ liệu (Structure), các ràng buộc dữ liệu (Constraints) và các phép toán trên dữ liệu (Operations).

Nền tảng kiến thức: Chương này phân biệt dữ liệu, cơ sở dữ liệu, hệ quản trị CSDL (DBMS), mô hình dữ liệu, lược đồ và trạng thái dữ liệu. Cơ sở dữ liệu lưu dữ liệu có liên hệ; DBMS là phần mềm định nghĩa, truy vấn, cập nhật và bảo vệ dữ liệu đó.

Phân tích: Data model là bộ khái niệm dùng để mô tả cấu trúc dữ liệu, quan hệ giữa dữ liệu và các ràng buộc. Relational, hierarchical, network và graph là các mô hình có thật; cần chọn mô hình khớp đặc điểm được hỏi, không nhầm mô hình với DBMS hay ngôn ngữ thao tác.

Đối chiếu phương án: B. Relationship. Các lựa chọn này không đáp ứng định nghĩa/điều kiện được hỏi; khi ôn, hãy chỉ ra chính xác điểm lệch của từng lựa chọn thay vì học thuộc mỗi chữ cái.
\end{question}

\begin{question}{48}
% [ID gốc: 1.44]
% [Nguồn: DBI202 Order 1 Q44]
Database data models usually have a way to describe limitations on what the data can be, which is called _ _ _ _ _.
\begin{choices}
\choice Constraints on the data
\choice Operations on the data
\choice Structure of the data
\choice Definition of data
\end{choices}
\textbf{Đáp án:} A
\textbf{Giải thích:} Đáp án theo dữ liệu đề: A — A. Constraints on the data.

Lời giải gốc: Giới hạn hoặc quy tắc áp đặt lên dữ liệu để đảm bảo tính đúng đắn và nhất quán được gọi là các ràng buộc dữ liệu (Constraints on the data).

Nền tảng kiến thức: Chương này phân biệt dữ liệu, cơ sở dữ liệu, hệ quản trị CSDL (DBMS), mô hình dữ liệu, lược đồ và trạng thái dữ liệu. Cơ sở dữ liệu lưu dữ liệu có liên hệ; DBMS là phần mềm định nghĩa, truy vấn, cập nhật và bảo vệ dữ liệu đó.

Phân tích: Constraint là quy tắc giới hạn giá trị/hình dạng dữ liệu hợp lệ, như miền giá trị, khóa hoặc tham chiếu. Operations mô tả thao tác trên dữ liệu, còn structure mô tả thành phần dữ liệu; constraint trả lời trực tiếp câu hỏi “giá trị nào được phép”.

Đối chiếu phương án: B. Operations on the data; C. Structure of the data; D. Definition of data. Các lựa chọn này không đáp ứng định nghĩa/điều kiện được hỏi; khi ôn, hãy chỉ ra chính xác điểm lệch của từng lựa chọn thay vì học thuộc mỗi chữ cái.
\end{question}

\begin{question}{49}
% [ID gốc: 1.55]
% [Nguồn: DBI202 Order 1 Q55]
Generally, what parts does a data model include in the context of a database?
\begin{choices}
\choice Structure of the data, operations on the data, and constraints on the data
\choice Metadata of the data, action logs, and data files
\choice Keys on the table, records accepted in the table, and attributes of the table
\choice Data schema, domain of attributes, and keys
\end{choices}
\textbf{Đáp án:} A
\textbf{Giải thích:} Đáp án theo dữ liệu đề: A — A. Structure of the data, operations on the data, and constraints on the data.

Lời giải gốc: Một mô hình dữ liệu chuẩn bao gồm cấu trúc dữ liệu, các phép toán trên dữ liệu và các ràng buộc dữ liệu.

Nền tảng kiến thức: Chương này phân biệt dữ liệu, cơ sở dữ liệu, hệ quản trị CSDL (DBMS), mô hình dữ liệu, lược đồ và trạng thái dữ liệu. Cơ sở dữ liệu lưu dữ liệu có liên hệ; DBMS là phần mềm định nghĩa, truy vấn, cập nhật và bảo vệ dữ liệu đó.

Phân tích: Data model là bộ khái niệm dùng để mô tả cấu trúc dữ liệu, quan hệ giữa dữ liệu và các ràng buộc. Relational, hierarchical, network và graph là các mô hình có thật; cần chọn mô hình khớp đặc điểm được hỏi, không nhầm mô hình với DBMS hay ngôn ngữ thao tác.

Đối chiếu phương án: B. Metadata of the data, action logs, and data files; C. Keys on the table, records accepted in the table, and attributes of the table; D. Data schema, domain of attributes, and keys. Các lựa chọn này không đáp ứng định nghĩa/điều kiện được hỏi; khi ôn, hãy chỉ ra chính xác điểm lệch của từng lựa chọn thay vì học thuộc mỗi chữ cái.
\end{question}

\begin{question}{50}
% [ID gốc: 1.66]
% [Nguồn: DBI202 Order 1 Q66]
A data model is a notation for describing data or information. The description generally consists of:
\begin{choices}
\choice Structure of the data
\choice Operations on the data
\choice Constraints on the data
\choice All of the others
\end{choices}
\textbf{Đáp án:} D
\textbf{Giải thích:} Đáp án theo dữ liệu đề: D — D. All of the others.

Lời giải gốc: Một mô hình dữ liệu hoàn chỉnh bao gồm ba thành phần: cấu trúc dữ liệu, các phép toán trên dữ liệu và các ràng buộc dữ liệu.

Nền tảng kiến thức: Chương này phân biệt dữ liệu, cơ sở dữ liệu, hệ quản trị CSDL (DBMS), mô hình dữ liệu, lược đồ và trạng thái dữ liệu. Cơ sở dữ liệu lưu dữ liệu có liên hệ; DBMS là phần mềm định nghĩa, truy vấn, cập nhật và bảo vệ dữ liệu đó.

Phân tích: Data model là bộ khái niệm dùng để mô tả cấu trúc dữ liệu, quan hệ giữa dữ liệu và các ràng buộc. Relational, hierarchical, network và graph là các mô hình có thật; cần chọn mô hình khớp đặc điểm được hỏi, không nhầm mô hình với DBMS hay ngôn ngữ thao tác.

Đối chiếu phương án: A. Structure of the data; B. Operations on the data; C. Constraints on the data. Các lựa chọn này không đáp ứng định nghĩa/điều kiện được hỏi; khi ôn, hãy chỉ ra chính xác điểm lệch của từng lựa chọn thay vì học thuộc mỗi chữ cái.
\end{question}

\begin{question}{51}
% [ID gốc: 3.1]
% [Nguồn: DBI202 Order 3 Q1]
A ____ is a relation name, together with the attributes of that relation.
\begin{choices}
\choice schema
\choice database
\choice database instance
\choice schema instance
\end{choices}
\textbf{Đáp án:} A
\textbf{Giải thích:} Đáp án theo dữ liệu đề: A — A. schema.

Nền tảng kiến thức: Mô hình quan hệ biểu diễn dữ liệu bằng quan hệ (bảng), bộ (hàng), thuộc tính (cột) và miền giá trị. Khóa nhận diện bộ; khóa ngoại và các ràng buộc toàn vẹn ngăn dữ liệu tham chiếu sai hoặc vi phạm miền.

Phân tích: Schema mô tả cấu trúc tương đối ổn định (tên quan hệ, thuộc tính và kiểu/ràng buộc); instance là các bộ dữ liệu hiện có tại một thời điểm. Câu hỏi định nghĩa relation name cùng các thuộc tính đang hỏi schema, không phải nội dung hàng hiện tại.

Đối chiếu phương án: B. database; C. database instance; D. schema instance. Các lựa chọn này không đáp ứng định nghĩa/điều kiện được hỏi; khi ôn, hãy chỉ ra chính xác điểm lệch của từng lựa chọn thay vì học thuộc mỗi chữ cái.
\end{question}

\begin{question}{52}
% [ID gốc: 3.2]
% [Nguồn: DBI202 Order 3 Q2]
What is another term for a row in a relational table?
\begin{choices}
\choice Attribute
\choice Tuple
\choice Field
\choice Relation
\end{choices}
\textbf{Đáp án:} B
\textbf{Giải thích:} Đáp án theo dữ liệu đề: B — B. Tuple.

Nền tảng kiến thức: Mô hình quan hệ biểu diễn dữ liệu bằng quan hệ (bảng), bộ (hàng), thuộc tính (cột) và miền giá trị. Khóa nhận diện bộ; khóa ngoại và các ràng buộc toàn vẹn ngăn dữ liệu tham chiếu sai hoặc vi phạm miền.

Phân tích: Trong quan hệ, mỗi hàng dữ liệu là một tuple (bộ), mỗi cột là attribute (thuộc tính), còn relation là tập/bag các tuple theo một schema. Dùng đúng thuật ngữ này giúp tránh nhầm hàng với trường/cột hoặc toàn bộ bảng.

Đối chiếu phương án: A. Attribute; C. Field; D. Relation. Các lựa chọn này không đáp ứng định nghĩa/điều kiện được hỏi; khi ôn, hãy chỉ ra chính xác điểm lệch của từng lựa chọn thay vì học thuộc mỗi chữ cái.
\end{question}

\begin{question}{53}
% [ID gốc: 3.3]
% [Nguồn: DBI202 Order 3 Q3]
A superkey is a set of attributes that can uniquely identify a tuple in a relation. Since the only key of R is {A1, A2}, it means that this combination of attributes can uniquely identify a tuple in the relation. However, any superset of this key is also considered a superkey. So the possible superkeys for R are: {A1,A2}, {A1,A2,A3}, {A1,A2,A4} and {A1,A2,A3,A4}. Therefore, R has 4 superkeys.
QN=11 (7966) Suppose R is a relation with attributes A1, A2, A3, A4.
The only key of R is {A1, A2}. So, how many super-keys do R have?
\begin{choices}
\choice 4
\choice 8
\choice 12
\choice 16
\end{choices}
\textbf{Đáp án:} A
\textbf{Giải thích:} Đáp án theo dữ liệu đề: A — A. 4.

Nền tảng kiến thức: Mô hình quan hệ biểu diễn dữ liệu bằng quan hệ (bảng), bộ (hàng), thuộc tính (cột) và miền giá trị. Khóa nhận diện bộ; khóa ngoại và các ràng buộc toàn vẹn ngăn dữ liệu tham chiếu sai hoặc vi phạm miền.

Phân tích: Khóa là thuộc tính hoặc tập thuộc tính dùng để nhận diện duy nhất một bộ/hàng. Khóa có thể gồm một hay nhiều cột, nên không bắt buộc luôn là cột đơn hoặc luôn gồm nhiều cột; nó cũng không “nhận diện một cột”.

Đối chiếu phương án: B. 8; C. 12; D. 16. Các lựa chọn này không đáp ứng định nghĩa/điều kiện được hỏi; khi ôn, hãy chỉ ra chính xác điểm lệch của từng lựa chọn thay vì học thuộc mỗi chữ cái.
\end{question}

\begin{question}{54}
% [ID gốc: 3.11]
% [Nguồn: DBI202 Order 3 Q11]
Relation = bảng (table)
Tuple = hàng (row)
Attribute = tên cột
In relational data model, an attribute is a column
\begin{choices}
\choice True
\choice False
\end{choices}
\textbf{Đáp án:} A
\textbf{Giải thích:} Đáp án theo dữ liệu đề: B — B. False.

Nền tảng kiến thức: Mô hình quan hệ biểu diễn dữ liệu bằng quan hệ (bảng), bộ (hàng), thuộc tính (cột) và miền giá trị. Khóa nhận diện bộ; khóa ngoại và các ràng buộc toàn vẹn ngăn dữ liệu tham chiếu sai hoặc vi phạm miền.

Phân tích: Data model là bộ khái niệm dùng để mô tả cấu trúc dữ liệu, quan hệ giữa dữ liệu và các ràng buộc. Relational, hierarchical, network và graph là các mô hình có thật; cần chọn mô hình khớp đặc điểm được hỏi, không nhầm mô hình với DBMS hay ngôn ngữ thao tác.

Đối chiếu phương án: A. True. Các lựa chọn này không đáp ứng định nghĩa/điều kiện được hỏi; khi ôn, hãy chỉ ra chính xác điểm lệch của từng lựa chọn thay vì học thuộc mỗi chữ cái.
\end{question}

\begin{question}{55}
% [ID gốc: 3.12]
% [Nguồn: DBI202 Order 3 Q12]
Look at this line:
MOVIES (Id, Name, Description)
What the above represent?
\begin{choices}
\choice A relation
\choice A schema
\choice A schema instance
\choice A relation instance
\end{choices}
\textbf{Đáp án:} B
\textbf{Giải thích:} Đáp án theo dữ liệu đề: B — B. A schema.

Nền tảng kiến thức: Mô hình quan hệ biểu diễn dữ liệu bằng quan hệ (bảng), bộ (hàng), thuộc tính (cột) và miền giá trị. Khóa nhận diện bộ; khóa ngoại và các ràng buộc toàn vẹn ngăn dữ liệu tham chiếu sai hoặc vi phạm miền.

Phân tích: Khóa là thuộc tính hoặc tập thuộc tính dùng để nhận diện duy nhất một bộ/hàng. Khóa có thể gồm một hay nhiều cột, nên không bắt buộc luôn là cột đơn hoặc luôn gồm nhiều cột; nó cũng không “nhận diện một cột”.

Đối chiếu phương án: A. A relation; C. A schema instance; D. A relation instance b // Theo mình con này chọn BC The correct answers are [B] and [C]. R is a…. Các lựa chọn này không đáp ứng định nghĩa/điều kiện được hỏi; khi ôn, hãy chỉ ra chính xác điểm lệch của từng lựa chọn thay vì học thuộc mỗi chữ cái.
\end{question}

\begin{question}{56}
% [ID gốc: 3.13b]
% [Nguồn: DBI202 Order 3 Q13b]
In relational data model, a tuple is a record or a row.
\begin{choices}
\choice True
\choice False
\end{choices}
\textbf{Đáp án:} A
\textbf{Giải thích:} Trong mô hình dữ liệu quan hệ, một bộ (tuple) tương ứng chính xác với một bản ghi (record) hoặc một hàng (row) trong bảng quan hệ.
\end{question}

\begin{question}{57}
% [ID gốc: 3.13]
% [Nguồn: DBI202 Order 3 Q13]
Look at the following data for table R:

$$\begin{array}{|c|c|c|}\hline 1 & 1 & 2 \\\hline 1 & 1 & 2 \\\hline 1 & 1 & 3 \\\hline\end{array}$$

Choose all correct answers:
\begin{choices}
\choice R is not a relation
\choice R is a relation
\choice R has no key
\choice None of the others
\end{choices}
\textbf{Đáp án:} BC
\textbf{Giải thích:} Bảng R có các dòng trùng lặp hoàn toàn (1, 1, 2 xuất hiện 2 lần). Theo ngữ nghĩa quan hệ túi (bag semantics / multiset) của SQL, R vẫn là một quan hệ (relation), nhưng không có khóa hợp lệ.
\end{question}

\begin{question}{58}
% [ID gốc: 3.15]
% [Nguồn: DBI202 Order 3 Q15]
A relation is not a list of tuples, but rather a set of tuples with the same attributes. In a set, the order of elements does not matter and there are no duplicate elements. In contrast, a list is an ordered collection of elements where duplicates are allowed.
A relation is a list of tuples
\begin{choices}
\choice True
\choice False
\end{choices}
\textbf{Đáp án:} B
\textbf{Giải thích:} Đáp án theo dữ liệu đề: B — B. False.

Nền tảng kiến thức: Mô hình quan hệ biểu diễn dữ liệu bằng quan hệ (bảng), bộ (hàng), thuộc tính (cột) và miền giá trị. Khóa nhận diện bộ; khóa ngoại và các ràng buộc toàn vẹn ngăn dữ liệu tham chiếu sai hoặc vi phạm miền.

Phân tích: Các hàng dữ liệu của một relation được gọi là tuple; header chứa tên thuộc tính không phải tuple dữ liệu. Trong bảng, tuple tương ứng một bản ghi, còn attribute tương ứng cột.

Đối chiếu phương án: A. True. Các lựa chọn này không đáp ứng định nghĩa/điều kiện được hỏi; khi ôn, hãy chỉ ra chính xác điểm lệch của từng lựa chọn thay vì học thuộc mỗi chữ cái.
\end{question}

\begin{question}{59}
% [ID gốc: 3.16]
% [Nguồn: DBI202 Order 3 Q16]
Cần phân biệt được set và list: List thì quan trọng thứ tự, còn set thì không. Một relation gồm các hàng (tuples) thì thứ tự các hàng trong bảng thay đổi thì vẫn là bảng ban đầu
A relation is a set of tuples
\begin{choices}
\choice True
\choice False
\end{choices}
\textbf{Đáp án:} A
\textbf{Giải thích:} Đáp án theo dữ liệu đề: A — A. True.

Nền tảng kiến thức: Mô hình quan hệ biểu diễn dữ liệu bằng quan hệ (bảng), bộ (hàng), thuộc tính (cột) và miền giá trị. Khóa nhận diện bộ; khóa ngoại và các ràng buộc toàn vẹn ngăn dữ liệu tham chiếu sai hoặc vi phạm miền.

Phân tích: Quan hệ theo mô hình toán học là tập các tuple nên thứ tự hàng không mang ý nghĩa; đổi vị trí hàng không đổi quan hệ. List thì có thứ tự, vì vậy không đồng nhất hai cấu trúc; thứ tự chỉ được đảm bảo trong kết quả SQL khi có ORDER BY.

Đối chiếu phương án: B. False. Các lựa chọn này không đáp ứng định nghĩa/điều kiện được hỏi; khi ôn, hãy chỉ ra chính xác điểm lệch của từng lựa chọn thay vì học thuộc mỗi chữ cái.
\end{question}

\begin{question}{60}
% [ID gốc: 3.17]
% [Nguồn: DBI202 Order 3 Q17]
A key:
\begin{choices}
\choice must always be composed of two or more columns.
\choice can only be one column.
\choice identifies a row.
\choice identifies a column.
\end{choices}
\textbf{Đáp án:} C
\textbf{Giải thích:} Đáp án theo dữ liệu đề: C — C. identifies a row..

Nền tảng kiến thức: Mô hình quan hệ biểu diễn dữ liệu bằng quan hệ (bảng), bộ (hàng), thuộc tính (cột) và miền giá trị. Khóa nhận diện bộ; khóa ngoại và các ràng buộc toàn vẹn ngăn dữ liệu tham chiếu sai hoặc vi phạm miền.

Phân tích: Khóa là thuộc tính hoặc tập thuộc tính dùng để nhận diện duy nhất một bộ/hàng. Khóa có thể gồm một hay nhiều cột, nên không bắt buộc luôn là cột đơn hoặc luôn gồm nhiều cột; nó cũng không “nhận diện một cột”.

Đối chiếu phương án: A. must always be composed of two or more columns.; B. can only be one column.; D. identifies a column.. Các lựa chọn này không đáp ứng định nghĩa/điều kiện được hỏi; khi ôn, hãy chỉ ra chính xác điểm lệch của từng lựa chọn thay vì học thuộc mỗi chữ cái.
\end{question}

\begin{question}{61}
% [ID gốc: 3.22]
% [Nguồn: DBI202 Order 3 Q22]
Which of the following statements about relational databases is true?
\begin{choices}
\choice The database is built on a relational data model.
\choice The database is created from excel software.
\choice Relational databases are only relations with attributes that are not primary keys.
\choice Relational database is the relationship between rows and columns in a data table.
\end{choices}
\textbf{Đáp án:} A
\textbf{Giải thích:} Đáp án theo dữ liệu đề: A — A. The database is built on a relational data model..

Nền tảng kiến thức: Mô hình quan hệ biểu diễn dữ liệu bằng quan hệ (bảng), bộ (hàng), thuộc tính (cột) và miền giá trị. Khóa nhận diện bộ; khóa ngoại và các ràng buộc toàn vẹn ngăn dữ liệu tham chiếu sai hoặc vi phạm miền.

Phân tích: Data model là bộ khái niệm dùng để mô tả cấu trúc dữ liệu, quan hệ giữa dữ liệu và các ràng buộc. Relational, hierarchical, network và graph là các mô hình có thật; cần chọn mô hình khớp đặc điểm được hỏi, không nhầm mô hình với DBMS hay ngôn ngữ thao tác.

Đối chiếu phương án: B. The database is created from excel software.; C. Relational databases are only relations with attributes that are not primary keys.; D. Relational database is the relationship between rows and columns in a data table.. Các lựa chọn này không đáp ứng định nghĩa/điều kiện được hỏi; khi ôn, hãy chỉ ra chính xác điểm lệch của từng lựa chọn thay vì học thuộc mỗi chữ cái.
\end{question}

\begin{question}{62}
% [ID gốc: 3.25]
% [Nguồn: DBI202 Order 3 Q25]
What does the attribute domain define in the relational model?
\begin{choices}
\choice The primary key of the table
\choice The foreign key references
\choice The cardinality of relationships
\choice The data type and range of values that an attribute can hold
\end{choices}
\textbf{Đáp án:} D
\textbf{Giải thích:} Đáp án theo dữ liệu đề: D — D. The data type and range of values that an attribute can hold.

Nền tảng kiến thức: Mô hình quan hệ biểu diễn dữ liệu bằng quan hệ (bảng), bộ (hàng), thuộc tính (cột) và miền giá trị. Khóa nhận diện bộ; khóa ngoại và các ràng buộc toàn vẹn ngăn dữ liệu tham chiếu sai hoặc vi phạm miền.

Phân tích: Khi đọc mô hình ER, tách ba thứ: entity set là loại đối tượng, relationship set là liên kết giữa các loại đối tượng, còn cardinality/participation là ràng buộc số lượng hoặc mức tham gia. Đáp án đúng phải gọi đúng thành phần mà câu hỏi mô tả.

Đối chiếu phương án: A. The primary key of the table; B. The foreign key references; C. The cardinality of relationships. Các lựa chọn này không đáp ứng định nghĩa/điều kiện được hỏi; khi ôn, hãy chỉ ra chính xác điểm lệch của từng lựa chọn thay vì học thuộc mỗi chữ cái.
\end{question}

\begin{question}{63}
% [ID gốc: 3.26]
% [Nguồn: DBI202 Order 3 Q26]
What is a Relational Database?
\begin{choices}
\choice A collection of two or more tables
\choice A database that is able to process queries, forms, reports and macro
\choice The same as an Excel file database
\choice A collection of two or more tables that are joined using relationships
\end{choices}
\textbf{Đáp án:} D
\textbf{Giải thích:} Đáp án theo dữ liệu đề: D — D. A collection of two or more tables that are joined using relationships.

Nền tảng kiến thức: Mô hình quan hệ biểu diễn dữ liệu bằng quan hệ (bảng), bộ (hàng), thuộc tính (cột) và miền giá trị. Khóa nhận diện bộ; khóa ngoại và các ràng buộc toàn vẹn ngăn dữ liệu tham chiếu sai hoặc vi phạm miền.

Phân tích: OLTP ghi nhận giao dịch nghiệp vụ ngắn và cập nhật thường xuyên; decision support/OLAP tập trung truy vấn phân tích, tổng hợp và báo cáo. Hãy phân loại theo mục đích thao tác chứ không theo việc cả hai cùng đọc dữ liệu.

Đối chiếu phương án: A. A collection of two or more tables; B. A database that is able to process queries, forms, reports and macro; C. The same as an Excel file database. Các lựa chọn này không đáp ứng định nghĩa/điều kiện được hỏi; khi ôn, hãy chỉ ra chính xác điểm lệch của từng lựa chọn thay vì học thuộc mỗi chữ cái.
\end{question}

\begin{question}{64}
% [ID gốc: 3.31]
% [Nguồn: DBI202 Order 3 Q31]
_______is a set of one or more attributes taken collectively to uniquely identify a record.
\begin{choices}
\choice Primary Key
\choice Foreign key
\choice Super key
\choice Candidate key
\end{choices}
\textbf{Đáp án:} C
\textbf{Giải thích:} Đáp án theo dữ liệu đề: C — C. Super key.

Nền tảng kiến thức: Mô hình quan hệ biểu diễn dữ liệu bằng quan hệ (bảng), bộ (hàng), thuộc tính (cột) và miền giá trị. Khóa nhận diện bộ; khóa ngoại và các ràng buộc toàn vẹn ngăn dữ liệu tham chiếu sai hoặc vi phạm miền.

Phân tích: Khóa ngoại buộc giá trị tham chiếu phải khớp một khóa được tham chiếu ở quan hệ đích (hoặc NULL nếu ràng buộc cho phép). Đây là toàn vẹn tham chiếu; nó khác tính duy nhất của primary/candidate key.

Đối chiếu phương án: A. Primary Key; B. Foreign key; D. Candidate key. Các lựa chọn này không đáp ứng định nghĩa/điều kiện được hỏi; khi ôn, hãy chỉ ra chính xác điểm lệch của từng lựa chọn thay vì học thuộc mỗi chữ cái.
\end{question}

\begin{question}{65}
% [ID gốc: 3.32]
% [Nguồn: DBI202 Order 3 Q32]
What does an RDBMS consist of?
\begin{choices}
\choice Collection of Tables
\choice Collection of Records
\choice Collection of Keys
\choice Collection of Field
\end{choices}
\textbf{Đáp án:} A
\textbf{Giải thích:} Đáp án theo dữ liệu đề: A — A. Collection of Tables.

Nền tảng kiến thức: Mô hình quan hệ biểu diễn dữ liệu bằng quan hệ (bảng), bộ (hàng), thuộc tính (cột) và miền giá trị. Khóa nhận diện bộ; khóa ngoại và các ràng buộc toàn vẹn ngăn dữ liệu tham chiếu sai hoặc vi phạm miền.

Phân tích: RDBMS tổ chức dữ liệu logic thành các bảng/quan hệ gồm hàng và cột, liên kết bằng khóa và được truy vấn bằng SQL. Record/tuple là hàng trong một bảng; một collection of tables mô tả cấu trúc tổng thể phù hợp hơn một collection of records.

Đối chiếu phương án: B. Collection of Records; C. Collection of Keys; D. Collection of Field. Các lựa chọn này không đáp ứng định nghĩa/điều kiện được hỏi; khi ôn, hãy chỉ ra chính xác điểm lệch của từng lựa chọn thay vì học thuộc mỗi chữ cái.
\end{question}

\begin{question}{66}
% [ID gốc: 3.34]
% [Nguồn: DBI202 Order 3 Q34]
Which of the following statements is correct when converting from the ERD model to the relational model?
\begin{choices}
\choice The order of the columns is not important
\choice The order of both rows and columns is important
\choice The order of the rows is important
\choice The order of the columns is important
\end{choices}
\textbf{Đáp án:} A
\textbf{Giải thích:} Đáp án theo dữ liệu đề: A — A. The order of the columns is not important.

Nền tảng kiến thức: Mô hình quan hệ biểu diễn dữ liệu bằng quan hệ (bảng), bộ (hàng), thuộc tính (cột) và miền giá trị. Khóa nhận diện bộ; khóa ngoại và các ràng buộc toàn vẹn ngăn dữ liệu tham chiếu sai hoặc vi phạm miền.

Phân tích: ER diagram mô tả các entity set, thuộc tính và relationship set, cùng các ràng buộc cần thiết. Nó tập trung vào đối tượng dữ liệu và liên kết giữa chúng, không mô tả luồng xử lý/chương trình như sơ đồ quy trình.

Đối chiếu phương án: B. The order of both rows and columns is important; C. The order of the rows is important; D. The order of the columns is important. Các lựa chọn này không đáp ứng định nghĩa/điều kiện được hỏi; khi ôn, hãy chỉ ra chính xác điểm lệch của từng lựa chọn thay vì học thuộc mỗi chữ cái.
\end{question}

\begin{question}{67}
% [ID gốc: 3.35]
% [Nguồn: DBI202 Order 3 Q35]
If there is more than one key for relation schema in DBMS then each key relation schema is classified as
\begin{choices}
\choice Prime key
\choice Super key
\choice Candidate key
\choice Primary key
\end{choices}
\textbf{Đáp án:} C
\textbf{Giải thích:} Đáp án theo dữ liệu đề: C — C. Candidate key.

Nền tảng kiến thức: Mô hình quan hệ biểu diễn dữ liệu bằng quan hệ (bảng), bộ (hàng), thuộc tính (cột) và miền giá trị. Khóa nhận diện bộ; khóa ngoại và các ràng buộc toàn vẹn ngăn dữ liệu tham chiếu sai hoặc vi phạm miền.

Phân tích: Lược đồ quan hệ ghi tên quan hệ kèm danh sách thuộc tính, thường theo dạng `Boat(BoatID, BoatName, Color, Status)`. Nếu đề hỏi schema thuần túy, không nhất thiết phải kèm kiểu dữ liệu hay các bộ dữ liệu đang lưu.

Đối chiếu phương án: A. Prime key; B. Super key; D. Primary key. Các lựa chọn này không đáp ứng định nghĩa/điều kiện được hỏi; khi ôn, hãy chỉ ra chính xác điểm lệch của từng lựa chọn thay vì học thuộc mỗi chữ cái.
\end{question}

\begin{question}{68}
% [ID gốc: 3.38]
% [Nguồn: DBI202 Order 3 Q38]
Consider a table named "Boat". Which option is the correct relation schema for the table?
\begin{choices}
\choice Boat: BoatID, BoatName, Color, Status
\choice Boat(BoatID, BoatName, Color, Status)
\choice Boat(BoatID: String, BoatName: String, Color: String, Status: String)
\choice Boat: {BoatID, BoatName, Color, Status}
\end{choices}
\textbf{Đáp án:} B
\textbf{Giải thích:} Đáp án theo dữ liệu đề: B — B. Boat(BoatID, BoatName, Color, Status).

Lời giải gốc: Ký hiệu tiêu chuẩn cho lược đồ quan hệ là tên quan hệ kèm danh sách thuộc tính đặt trong dấu ngoặc đơn, dạng RelationName(Attribute1, Attribute2, ...).

Nền tảng kiến thức: Mô hình quan hệ biểu diễn dữ liệu bằng quan hệ (bảng), bộ (hàng), thuộc tính (cột) và miền giá trị. Khóa nhận diện bộ; khóa ngoại và các ràng buộc toàn vẹn ngăn dữ liệu tham chiếu sai hoặc vi phạm miền.

Phân tích: Lược đồ quan hệ ghi tên quan hệ kèm danh sách thuộc tính, thường theo dạng `Boat(BoatID, BoatName, Color, Status)`. Nếu đề hỏi schema thuần túy, không nhất thiết phải kèm kiểu dữ liệu hay các bộ dữ liệu đang lưu.

Đối chiếu phương án: A. Boat: BoatID, BoatName, Color, Status; C. Boat(BoatID: String, BoatName: String, Color: String, Status: String); D. Boat: {BoatID, BoatName, Color, Status}. Các lựa chọn này không đáp ứng định nghĩa/điều kiện được hỏi; khi ôn, hãy chỉ ra chính xác điểm lệch của từng lựa chọn thay vì học thuộc mỗi chữ cái.
\end{question}

\begin{question}{69}
% [ID gốc: 3.40]
% [Nguồn: DBI202 Order 3 Q40]
The rows of a relation, other than the header row containing the attribute names, are called ______.
\begin{choices}
\choice Schemas
\choice Tuples
\choice Domains
\choice Attributes
\end{choices}
\textbf{Đáp án:} B
\textbf{Giải thích:} Đáp án theo dữ liệu đề: B — B. Tuples.

Lời giải gốc: Trong mô hình dữ liệu quan hệ, các hàng (rows) dữ liệu trong một bảng được gọi thuật ngữ học thuật là bộ dữ liệu (Tuples). Các cột gọi là Attributes và tiêu đề là Schema.

Nền tảng kiến thức: Mô hình quan hệ biểu diễn dữ liệu bằng quan hệ (bảng), bộ (hàng), thuộc tính (cột) và miền giá trị. Khóa nhận diện bộ; khóa ngoại và các ràng buộc toàn vẹn ngăn dữ liệu tham chiếu sai hoặc vi phạm miền.

Phân tích: Các hàng dữ liệu của một relation được gọi là tuple; header chứa tên thuộc tính không phải tuple dữ liệu. Trong bảng, tuple tương ứng một bản ghi, còn attribute tương ứng cột.

Đối chiếu phương án: A. Schemas; C. Domains; D. Attributes. Các lựa chọn này không đáp ứng định nghĩa/điều kiện được hỏi; khi ôn, hãy chỉ ra chính xác điểm lệch của từng lựa chọn thay vì học thuộc mỗi chữ cái.
\end{question}

\begin{question}{70}
% [ID gốc: 3.43]
% [Nguồn: DBI202 Order 3 Q43]
A(n) _ _ _ _ _ is a property of the entire relation, rather than of the individual tuples in which each tuple is unique.
\begin{choices}
\choice Row
\choice Key
\choice Attribute
\choice Field
\end{choices}
\textbf{Đáp án:} B
\textbf{Giải thích:} Đáp án theo dữ liệu đề: B — B. Key.

Lời giải gốc: Khóa (Key) là thuộc tính áp dụng cho toàn bộ quan hệ để đảm bảo tính duy nhất của từng dòng (tuple).

Nền tảng kiến thức: Mô hình quan hệ biểu diễn dữ liệu bằng quan hệ (bảng), bộ (hàng), thuộc tính (cột) và miền giá trị. Khóa nhận diện bộ; khóa ngoại và các ràng buộc toàn vẹn ngăn dữ liệu tham chiếu sai hoặc vi phạm miền.

Phân tích: Các hàng dữ liệu của một relation được gọi là tuple; header chứa tên thuộc tính không phải tuple dữ liệu. Trong bảng, tuple tương ứng một bản ghi, còn attribute tương ứng cột.

Đối chiếu phương án: A. Row; C. Attribute; D. Field. Các lựa chọn này không đáp ứng định nghĩa/điều kiện được hỏi; khi ôn, hãy chỉ ra chính xác điểm lệch của từng lựa chọn thay vì học thuộc mỗi chữ cái.
\end{question}

\begin{question}{71}
% [ID gốc: 3.52]
% [Nguồn: DBI202 Order 3 Q52]
The subset of a super key is a candidate key under what condition?
\begin{choices}
\choice No proper subset is a super key
\choice All subsets are super keys
\choice Subset is a super key
\choice Each subset is a super key
\end{choices}
\textbf{Đáp án:} A
\textbf{Giải thích:} Đáp án theo dữ liệu đề: A — A. No proper subset is a super key.

Lời giải gốc: Khóa dự tuyển (Candidate key) định nghĩa là một siêu khóa tối tiểu (Minimal superkey). Do đó, một tập con của siêu khóa trở thành khóa dự tuyển khi và chỉ khi không có tập con thực sự nào của nó còn là một siêu khóa nữa.

Nền tảng kiến thức: Mô hình quan hệ biểu diễn dữ liệu bằng quan hệ (bảng), bộ (hàng), thuộc tính (cột) và miền giá trị. Khóa nhận diện bộ; khóa ngoại và các ràng buộc toàn vẹn ngăn dữ liệu tham chiếu sai hoặc vi phạm miền.

Phân tích: Kiểm tra khóa bằng hai tiêu chí: tập thuộc tính phải xác định duy nhất mọi bộ (tính siêu khóa), và nếu là candidate key thì phải tối thiểu, tức bỏ bất kỳ thuộc tính nào cũng mất tính duy nhất. PRIMARY KEY chọn một candidate key và không nhận NULL.

Đối chiếu phương án: B. All subsets are super keys; C. Subset is a super key; D. Each subset is a super key. Các lựa chọn này không đáp ứng định nghĩa/điều kiện được hỏi; khi ôn, hãy chỉ ra chính xác điểm lệch của từng lựa chọn thay vì học thuộc mỗi chữ cái.
\end{question}

\begin{question}{72}
% [ID gốc: 3.57]
% [Nguồn: DBI202 Order 3 Q57]
What is a superkey?
\begin{choices}
\choice A minimal set of attributes that uniquely identifies a row.
\choice Any set of attributes that contains a key.
\choice An attribute that is not part of any key.
\choice A foreign key in a table.
\end{choices}
\textbf{Đáp án:} B
\textbf{Giải thích:} Đáp án theo dữ liệu đề: B — B. Any set of attributes that contains a key..

Lời giải gốc: Siêu khóa (Superkey) là bất kỳ tập hợp thuộc tính nào có chứa một khóa (Candidate key) và có khả năng định danh duy nhất một dòng trong bảng.

Nền tảng kiến thức: Mô hình quan hệ biểu diễn dữ liệu bằng quan hệ (bảng), bộ (hàng), thuộc tính (cột) và miền giá trị. Khóa nhận diện bộ; khóa ngoại và các ràng buộc toàn vẹn ngăn dữ liệu tham chiếu sai hoặc vi phạm miền.

Phân tích: Kiểm tra khóa bằng hai tiêu chí: tập thuộc tính phải xác định duy nhất mọi bộ (tính siêu khóa), và nếu là candidate key thì phải tối thiểu, tức bỏ bất kỳ thuộc tính nào cũng mất tính duy nhất. PRIMARY KEY chọn một candidate key và không nhận NULL.

Đối chiếu phương án: A. A minimal set of attributes that uniquely identifies a row.; C. An attribute that is not part of any key.; D. A foreign key in a table.. Các lựa chọn này không đáp ứng định nghĩa/điều kiện được hỏi; khi ôn, hãy chỉ ra chính xác điểm lệch của từng lựa chọn thay vì học thuộc mỗi chữ cái.
\end{question}

\begin{question}{73}
% [ID gốc: 3.60]
% [Nguồn: DBI202 Order 3 Q60]
In a relation, rows are called:
\begin{choices}
\choice Attributes
\choice Fields
\choice Tuples
\choice Keys
\end{choices}
\textbf{Đáp án:} C
\textbf{Giải thích:} Đáp án theo dữ liệu đề: C — C. Tuples.

Lời giải gốc: Trong mô hình quan hệ, một bảng (Table) được gọi là Quan hệ (Relation), cột (Column) là Thuộc tính (Attribute), và hàng (Row) được gọi là Bộ (Tuple).

Nền tảng kiến thức: Mô hình quan hệ biểu diễn dữ liệu bằng quan hệ (bảng), bộ (hàng), thuộc tính (cột) và miền giá trị. Khóa nhận diện bộ; khóa ngoại và các ràng buộc toàn vẹn ngăn dữ liệu tham chiếu sai hoặc vi phạm miền.

Phân tích: Các hàng dữ liệu của một relation được gọi là tuple; header chứa tên thuộc tính không phải tuple dữ liệu. Trong bảng, tuple tương ứng một bản ghi, còn attribute tương ứng cột.

Đối chiếu phương án: A. Attributes; B. Fields; D. Keys. Các lựa chọn này không đáp ứng định nghĩa/điều kiện được hỏi; khi ôn, hãy chỉ ra chính xác điểm lệch của từng lựa chọn thay vì học thuộc mỗi chữ cái.
\end{question}

\begin{question}{74}
% [ID gốc: 3.69]
% [Nguồn: DBI202 Order 3 Q69]
The name of a relation and the set of attributes for a relation is called:
\begin{choices}
\choice Schema
\choice Tuple
\choice Domain
\choice Attribute
\end{choices}
\textbf{Đáp án:} A
\textbf{Giải thích:} Đáp án theo dữ liệu đề: A — A. Schema.

Lời giải gốc: Tên của quan hệ đi kèm tập hợp các thuộc tính của nó tạo thành lược đồ quan hệ (Schema).

Nền tảng kiến thức: Mô hình quan hệ biểu diễn dữ liệu bằng quan hệ (bảng), bộ (hàng), thuộc tính (cột) và miền giá trị. Khóa nhận diện bộ; khóa ngoại và các ràng buộc toàn vẹn ngăn dữ liệu tham chiếu sai hoặc vi phạm miền.

Phân tích: Trong quan hệ, mỗi hàng dữ liệu là một tuple (bộ), mỗi cột là attribute (thuộc tính), còn relation là tập/bag các tuple theo một schema. Dùng đúng thuật ngữ này giúp tránh nhầm hàng với trường/cột hoặc toàn bộ bảng.

Đối chiếu phương án: B. Tuple; C. Domain; D. Attribute. Các lựa chọn này không đáp ứng định nghĩa/điều kiện được hỏi; khi ôn, hãy chỉ ra chính xác điểm lệch của từng lựa chọn thay vì học thuộc mỗi chữ cái.
\end{question}

\begin{question}{75}
% [ID gốc: 3.71]
% [Nguồn: DBI202 Order 3 Q71]
What is a relation in the relational model?
\begin{choices}
\choice A collection of files.
\choice A table with rows and columns.
\choice A hierarchical structure.
\choice A network of data elements.
\end{choices}
\textbf{Đáp án:} B
\textbf{Giải thích:} Đáp án theo dữ liệu đề: B — B. A table with rows and columns..

Lời giải gốc: Trong mô hình quan hệ, một relation (quan hệ) chính là một bảng dữ liệu bao gồm các hàng và cột. ---

Nền tảng kiến thức: Mô hình quan hệ biểu diễn dữ liệu bằng quan hệ (bảng), bộ (hàng), thuộc tính (cột) và miền giá trị. Khóa nhận diện bộ; khóa ngoại và các ràng buộc toàn vẹn ngăn dữ liệu tham chiếu sai hoặc vi phạm miền.

Phân tích: Relation là một bảng logic gồm các tuple cùng schema; mỗi tuple có giá trị cho các thuộc tính đã định nghĩa. Theo lý thuyết tập hợp, thứ tự hàng không xác định quan hệ, dù DBMS có thể trả kết quả theo một thứ tự tình cờ nếu không dùng ORDER BY.

Đối chiếu phương án: A. A collection of files.; C. A hierarchical structure.; D. A network of data elements.. Các lựa chọn này không đáp ứng định nghĩa/điều kiện được hỏi; khi ôn, hãy chỉ ra chính xác điểm lệch của từng lựa chọn thay vì học thuộc mỗi chữ cái.
\end{question}

\begin{question}{76}
% [ID gốc: 4.2]
% [Nguồn: DBI202 Order 4 Q2]
Suppose relation R(A,B).
Choose the SQL statement that is equivalent to the following Relational Algebra expression:

$$\pi_{A+B \rightarrow C}(\sigma_{A+B > 10}(R))$$
\begin{choices}
\choice SELECT *
FROM R
\choice SELECT (A+B) AS C
FROM R
WHERE (A+B) > 10
\choice SELECT (A+B) AS C
FROM R
\choice There is not any choice
\end{choices}
\textbf{Đáp án:} B
\textbf{Giải thích:} Đáp án theo dữ liệu đề: B — B. SELECT (A+B) AS C FROM R WHERE (A+B) > 10.

Nền tảng kiến thức: Đại số quan hệ thao tác trên quan hệ: selection lọc hàng, projection chọn cột, join ghép bộ theo điều kiện, còn union/difference/intersection là các phép tập hợp. Cần phân biệt ngữ nghĩa tập hợp với bag/multiset vì bag cho phép bản sao và có quy tắc bội số riêng.

Phân tích: Đọc truy vấn theo các giai đoạn logic: FROM/JOIN tạo dữ liệu nguồn, WHERE lọc từng hàng, GROUP BY tạo nhóm, HAVING lọc nhóm, SELECT xác định biểu thức đầu ra và ORDER BY sắp kết quả. Với phép join, kiểm tra điều kiện ghép và các hàng không khớp trước khi tính kết quả.

Đối chiếu phương án: A. SELECT * FROM R; C. SELECT (A+B) AS C FROM R; D. There is not any choice. Các lựa chọn này không đáp ứng định nghĩa/điều kiện được hỏi; khi ôn, hãy chỉ ra chính xác điểm lệch của từng lựa chọn thay vì học thuộc mỗi chữ cái.
\end{question}

\begin{question}{77}
% [ID gốc: 4.4]
% [Nguồn: DBI202 Order 4 Q4]
Given relations R(A,B) and S(B,C,D). The result of natural join of the relations R and S has
\begin{choices}
\choice Only attribute B
\choice Only two attributes R.B and S.B
\choice Attributes A, B, C, D
\choice None of the others
\end{choices}
\textbf{Đáp án:} C
\textbf{Giải thích:} Đáp án theo dữ liệu đề: C — C. Attributes A, B, C, D.

Nền tảng kiến thức: Đại số quan hệ thao tác trên quan hệ: selection lọc hàng, projection chọn cột, join ghép bộ theo điều kiện, còn union/difference/intersection là các phép tập hợp. Cần phân biệt ngữ nghĩa tập hợp với bag/multiset vì bag cho phép bản sao và có quy tắc bội số riêng.

Phân tích: Join ghép các bộ thỏa điều kiện. Inner join loại bộ không ghép được; LEFT/RIGHT OUTER JOIN giữ mọi bộ từ phía được chỉ định và điền NULL phía thiếu. Natural join tự ghép các thuộc tính trùng tên và chỉ xuất một cột cho mỗi thuộc tính ghép.

Đối chiếu phương án: A. Only attribute B; B. Only two attributes R.B and S.B; D. None of the others. Các lựa chọn này không đáp ứng định nghĩa/điều kiện được hỏi; khi ôn, hãy chỉ ra chính xác điểm lệch của từng lựa chọn thay vì học thuộc mỗi chữ cái.
\end{question}

\begin{question}{78}
% [ID gốc: 4.7]
% [Nguồn: DBI202 Order 4 Q7]
Given table A with n rows and table B with m rows. How many rows of table result by Cartesian product on AxB?
\begin{choices}
\choice m*n rows
\choice m*n +1 rows
\choice m*n-1 rows
\choice m+n rows
\end{choices}
\textbf{Đáp án:} A
\textbf{Giải thích:} Đáp án theo dữ liệu đề: A — A. m*n rows.

Nền tảng kiến thức: Đại số quan hệ thao tác trên quan hệ: selection lọc hàng, projection chọn cột, join ghép bộ theo điều kiện, còn union/difference/intersection là các phép tập hợp. Cần phân biệt ngữ nghĩa tập hợp với bag/multiset vì bag cho phép bản sao và có quy tắc bội số riêng.

Phân tích: Tích Descartes tạo một bộ cho mỗi cặp (một hàng từ quan hệ thứ nhất, một hàng từ quan hệ thứ hai). Vì thế nếu hai đầu vào có m và n hàng thì kết quả có m×n hàng; điều kiện join thường lọc bớt các cặp này.

Đối chiếu phương án: B. m*n +1 rows; C. m*n-1 rows; D. m+n rows. Các lựa chọn này không đáp ứng định nghĩa/điều kiện được hỏi; khi ôn, hãy chỉ ra chính xác điểm lệch của từng lựa chọn thay vì học thuộc mỗi chữ cái.
\end{question}

\begin{question}{79}
% [ID gốc: 4.8]
% [Nguồn: DBI202 Order 4 Q8]
Which of the following is NOT the binary operation in the relational Algebra?
\begin{choices}
\choice Project
\choice Cartesian product
\choice Union
\choice Set difference
\end{choices}
\textbf{Đáp án:} A
\textbf{Giải thích:} Đáp án theo dữ liệu đề: A — A. Project.

Nền tảng kiến thức: Đại số quan hệ thao tác trên quan hệ: selection lọc hàng, projection chọn cột, join ghép bộ theo điều kiện, còn union/difference/intersection là các phép tập hợp. Cần phân biệt ngữ nghĩa tập hợp với bag/multiset vì bag cho phép bản sao và có quy tắc bội số riêng.

Phân tích: Phép tập hợp trên quan hệ cần so sánh các tuple theo thuộc tính tương ứng; union lấy hợp, intersection lấy phần chung và difference lấy phần có ở vế trái mà không có ở vế phải. Trước tiên kiểm tra tính tương thích schema, sau đó áp phép theo đúng hướng/toán tử.

Đối chiếu phương án: B. Cartesian product; C. Union; D. Set difference. Các lựa chọn này không đáp ứng định nghĩa/điều kiện được hỏi; khi ôn, hãy chỉ ra chính xác điểm lệch của từng lựa chọn thay vì học thuộc mỗi chữ cái.
\end{question}

\begin{question}{80}
% [ID gốc: 4.9]
% [Nguồn: DBI202 Order 4 Q9]
_____ returns a relation instance containing all tuples that occur in both R and S
\begin{choices}
\choice Intersection
\choice Set-difference
\choice Union
\choice Cross-product
\end{choices}
\textbf{Đáp án:} A
\textbf{Giải thích:} Đáp án theo dữ liệu đề: A — A. Intersection.

Nền tảng kiến thức: Đại số quan hệ thao tác trên quan hệ: selection lọc hàng, projection chọn cột, join ghép bộ theo điều kiện, còn union/difference/intersection là các phép tập hợp. Cần phân biệt ngữ nghĩa tập hợp với bag/multiset vì bag cho phép bản sao và có quy tắc bội số riêng.

Phân tích: Các hàng dữ liệu của một relation được gọi là tuple; header chứa tên thuộc tính không phải tuple dữ liệu. Trong bảng, tuple tương ứng một bản ghi, còn attribute tương ứng cột.

Đối chiếu phương án: B. Set-difference; C. Union; D. Cross-product. Các lựa chọn này không đáp ứng định nghĩa/điều kiện được hỏi; khi ôn, hãy chỉ ra chính xác điểm lệch của từng lựa chọn thay vì học thuộc mỗi chữ cái.
\end{question}

\begin{question}{81}
% [ID gốc: 4.10]
% [Nguồn: DBI202 Order 4 Q10]
The result which operation contains all pairs of tuples from the two relations, regardless of whether their attribute values match.
\begin{choices}
\choice Join
\choice Cartesian product
\choice Intersection
\choice Set difference
\end{choices}
\textbf{Đáp án:} B
\textbf{Giải thích:} Đáp án theo dữ liệu đề: B — B. Cartesian product.

Nền tảng kiến thức: Đại số quan hệ thao tác trên quan hệ: selection lọc hàng, projection chọn cột, join ghép bộ theo điều kiện, còn union/difference/intersection là các phép tập hợp. Cần phân biệt ngữ nghĩa tập hợp với bag/multiset vì bag cho phép bản sao và có quy tắc bội số riêng.

Phân tích: Các hàng dữ liệu của một relation được gọi là tuple; header chứa tên thuộc tính không phải tuple dữ liệu. Trong bảng, tuple tương ứng một bản ghi, còn attribute tương ứng cột.

Đối chiếu phương án: A. Join; C. Intersection; D. Set difference. Các lựa chọn này không đáp ứng định nghĩa/điều kiện được hỏi; khi ôn, hãy chỉ ra chính xác điểm lệch của từng lựa chọn thay vì học thuộc mỗi chữ cái.
\end{question}

\begin{question}{82}
% [ID gốc: 4.11]
% [Nguồn: DBI202 Order 4 Q11]
What is relational algebra?
\begin{choices}
\choice Relational algebra is a set of operations on relations.
\choice Relational algebra is the decomposition of relations.
\choice Relational algebra is eliminated from relations.
\choice Relational algebra is gone from relations.
\end{choices}
\textbf{Đáp án:} A
\textbf{Giải thích:} Đáp án theo dữ liệu đề: A — A. Relational algebra is a set of operations on relations..

Nền tảng kiến thức: Đại số quan hệ thao tác trên quan hệ: selection lọc hàng, projection chọn cột, join ghép bộ theo điều kiện, còn union/difference/intersection là các phép tập hợp. Cần phân biệt ngữ nghĩa tập hợp với bag/multiset vì bag cho phép bản sao và có quy tắc bội số riêng.

Phân tích: Phép tập hợp trên quan hệ cần so sánh các tuple theo thuộc tính tương ứng; union lấy hợp, intersection lấy phần chung và difference lấy phần có ở vế trái mà không có ở vế phải. Trước tiên kiểm tra tính tương thích schema, sau đó áp phép theo đúng hướng/toán tử.

Đối chiếu phương án: B. Relational algebra is the decomposition of relations.; C. Relational algebra is eliminated from relations.; D. Relational algebra is gone from relations.. Các lựa chọn này không đáp ứng định nghĩa/điều kiện được hỏi; khi ôn, hãy chỉ ra chính xác điểm lệch của từng lựa chọn thay vì học thuộc mỗi chữ cái.
\end{question}

\begin{question}{83}
% [ID gốc: 4.12]
% [Nguồn: DBI202 Order 4 Q12]
For the following relations and a Set operation:

$$\textbf{Employee} \qquad \textbf{Customer}$$
$$\begin{array}{|c|c|c|}
\hline
\text{Name} & \text{City} & \text{Gender} \\
\hline
\text{Nguyen Trung Truc} & \text{HCM} & \text{Male} \\
\hline
\text{Le Hong Son} & \text{HCM} & \text{Male} \\
\hline
\end{array}
\qquad
\begin{array}{|c|c|c|}
\hline
\text{Name} & \text{City} & \text{Gender} \\
\hline
\text{Nguyen Trung Truc} & \text{HCM} & \text{Male} \\
\hline
\text{Tran Tuyet Nhi} & \text{HN} & \text{Female} \\
\hline
\end{array}$$

What is the result of the Union operation (Employee \cup Customer)?
\begin{choices}
\choice $$\begin{array}{|c|c|c|}\hline \text{Name} & \text{City} & \text{Gender} \\\hline \text{Nguyen Trung Truc} & \text{HCM} & \text{Male} \\\hline \text{Le Hong Son} & \text{HCM} & \text{Male} \\\hline\end{array}$$
\choice $$\begin{array}{|c|c|c|}\hline \text{Name} & \text{City} & \text{Gender} \\\hline \text{Nguyen Trung Truc} & \text{HCM} & \text{Male} \\\hline \text{Tran Tuyet Nhi} & \text{HN} & \text{Female} \\\hline\end{array}$$
\choice $$\begin{array}{|c|c|c|}\hline \text{Name} & \text{City} & \text{Gender} \\\hline \text{Nguyen Trung Truc} & \text{HCM} & \text{Male} \\\hline \text{Le Hong Son} & \text{HCM} & \text{Male} \\\hline \text{Nguyen Trung Truc} & \text{HCM} & \text{Male} \\\hline \text{Tran Tuyet Nhi} & \text{HN} & \text{Female} \\\hline\end{array}$$
\choice $$\begin{array}{|c|c|c|}\hline \text{Name} & \text{City} & \text{Gender} \\\hline \text{Nguyen Trung Truc} & \text{HCM} & \text{Male} \\\hline \text{Le Hong Son} & \text{HCM} & \text{Male} \\\hline \text{Tran Tuyet Nhi} & \text{HN} & \text{Female} \\\hline\end{array}$$
\end{choices}
\textbf{Đáp án:} D
\textbf{Giải thích:} Đáp án theo dữ liệu đề: D — D. Name | City | Gender Nguyen Trung Truc | HCM | Male Le Hong Son | HCM | Male Tran Tuyet Nhi | HN | Female.

Nền tảng kiến thức: Đại số quan hệ thao tác trên quan hệ: selection lọc hàng, projection chọn cột, join ghép bộ theo điều kiện, còn union/difference/intersection là các phép tập hợp. Cần phân biệt ngữ nghĩa tập hợp với bag/multiset vì bag cho phép bản sao và có quy tắc bội số riêng.

Phân tích: Phép tập hợp trên quan hệ cần so sánh các tuple theo thuộc tính tương ứng; union lấy hợp, intersection lấy phần chung và difference lấy phần có ở vế trái mà không có ở vế phải. Trước tiên kiểm tra tính tương thích schema, sau đó áp phép theo đúng hướng/toán tử.

Đối chiếu phương án: A. Name | City | Gender Nguyen Trung Truc | HCM | Male Le Hong Son | HCM | Male; B. Name | City | Gender Nguyen Trung Truc | HCM | Male Tran Tuyet Nhi | HN | Female; C. Name | City | Gender Nguyen Trung Truc | HCM | Male Le Hong Son | HCM | Male Nguyen Trung Truc…. Các lựa chọn này không đáp ứng định nghĩa/điều kiện được hỏi; khi ôn, hãy chỉ ra chính xác điểm lệch của từng lựa chọn thay vì học thuộc mỗi chữ cái.
\end{question}

\begin{question}{84}
% [ID gốc: 4.13]
% [Nguồn: DBI202 Order 4 Q13]
For the following relation:

$$\textbf{Student}$$
$$\begin{array}{|c|c|c|c|}
\hline
\text{ID} & \text{Name} & \text{City} & \text{Gender} \\
\hline
\text{SE1234} & \text{Le Hong Phong} & \text{HCM} & \text{Male} \\
\hline
\text{SE1235} & \text{Le Hong Hanh} & \text{HCM} & \text{Female} \\
\hline
\text{SE1236} & \text{Tran Binh Trong} & \text{HN} & \text{Male} \\
\hline
\end{array}$$

Use Algebraic query language to list all name and city of students in HCM and Male?
\begin{choices}
\choice \(\pi_{\text{Name}, \text{City}}(\sigma_{\text{City}=\'HCM\' \land \text{Gender}=\'Male\'}(\text{Student}))\)
\choice \(\pi_{\text{Name}, \text{City}}[\sigma_{\text{City}=\'HCM\' \land \text{Gender}=\'Male\'}(\text{Student})]\)
\choice \(\sigma_{\text{Name}, \text{City}}(\pi_{\text{City}=\'HCM\' \land \text{Gender}=\'Male\'}(\text{Student}))\)
\choice None of the others
\end{choices}
\textbf{Đáp án:} A
\textbf{Giải thích:} Đáp án theo dữ liệu đề: A — A. Name, City (City='HCM' AND Gender='Male' (Student)).

Nền tảng kiến thức: Đại số quan hệ thao tác trên quan hệ: selection lọc hàng, projection chọn cột, join ghép bộ theo điều kiện, còn union/difference/intersection là các phép tập hợp. Cần phân biệt ngữ nghĩa tập hợp với bag/multiset vì bag cho phép bản sao và có quy tắc bội số riêng.

Phân tích: Đọc truy vấn theo các giai đoạn logic: FROM/JOIN tạo dữ liệu nguồn, WHERE lọc từng hàng, GROUP BY tạo nhóm, HAVING lọc nhóm, SELECT xác định biểu thức đầu ra và ORDER BY sắp kết quả. Với phép join, kiểm tra điều kiện ghép và các hàng không khớp trước khi tính kết quả.

Đối chiếu phương án: B. Name, City([City='HCM' AND Gender='Male' (Student)); C. Name, City (City='HCM' && Gender='Male' (Student)); D. Name, City (City='HCM' && Gender='Male' (Student). Các lựa chọn này không đáp ứng định nghĩa/điều kiện được hỏi; khi ôn, hãy chỉ ra chính xác điểm lệch của từng lựa chọn thay vì học thuộc mỗi chữ cái.
\end{question}

\begin{question}{85}
% [ID gốc: 4.14]
% [Nguồn: DBI202 Order 4 Q14]
The ___ operation allows the combining of two relations by merging pairs of tuples, one from each relation, into a single tuple.
\begin{choices}
\choice Select
\choice Join
\choice Union
\choice Intersection
\end{choices}
\textbf{Đáp án:} B
\textbf{Giải thích:} Đáp án theo dữ liệu đề: B — B. Join.

Nền tảng kiến thức: Đại số quan hệ thao tác trên quan hệ: selection lọc hàng, projection chọn cột, join ghép bộ theo điều kiện, còn union/difference/intersection là các phép tập hợp. Cần phân biệt ngữ nghĩa tập hợp với bag/multiset vì bag cho phép bản sao và có quy tắc bội số riêng.

Phân tích: Các hàng dữ liệu của một relation được gọi là tuple; header chứa tên thuộc tính không phải tuple dữ liệu. Trong bảng, tuple tương ứng một bản ghi, còn attribute tương ứng cột.

Đối chiếu phương án: A. Select; C. Union; D. Intersection. Các lựa chọn này không đáp ứng định nghĩa/điều kiện được hỏi; khi ôn, hãy chỉ ra chính xác điểm lệch của từng lựa chọn thay vì học thuộc mỗi chữ cái.
\end{question}

\begin{question}{86}
% [ID gốc: 4.16]
% [Nguồn: DBI202 Order 4 Q16]
For the following relations:

$$\textbf{Employee} \qquad \textbf{Customer}$$
$$\begin{array}{|c|c|c|}
\hline
\text{Name} & \text{City} & \text{Gender} \\
\hline
\text{Nguyen Trung Truc} & \text{HCM} & \text{Male} \\
\hline
\text{Le Hong Son} & \text{HCM} & \text{Male} \\
\hline
\end{array}
\qquad
\begin{array}{|c|c|c|}
\hline
\text{Name} & \text{City} & \text{Gender} \\
\hline
\text{Nguyen Trung Truc} & \text{HCM} & \text{Male} \\
\hline
\text{Tran Tuyet Nhi} & \text{HN} & \text{Female} \\
\hline
\end{array}$$

What is the result of the Intersection operation (Employee \cap Customer)?
\begin{choices}
\choice $$\begin{array}{|c|c|c|}\hline \text{Name} & \text{City} & \text{Gender} \\\hline \text{Nguyen Trung Truc} & \text{HCM} & \text{Male} \\\hline\end{array}$$
\choice $$\begin{array}{|c|c|c|}\hline \text{Name} & \text{City} & \text{Gender} \\\hline \text{Le Hong Son} & \text{HCM} & \text{Male} \\\hline\end{array}$$
\choice $$\begin{array}{|c|c|c|}\hline \text{Name} & \text{City} & \text{Gender} \\\hline \text{Tran Tuyet Nhi} & \text{HN} & \text{Female} \\\hline\end{array}$$
\choice None of the others
\end{choices}
\textbf{Đáp án:} A
\textbf{Giải thích:} Đáp án theo dữ liệu đề: C — C. maker (type='PC' (Product)) \ maker (type='Laptop' (Product)).

Nền tảng kiến thức: Đại số quan hệ thao tác trên quan hệ: selection lọc hàng, projection chọn cột, join ghép bộ theo điều kiện, còn union/difference/intersection là các phép tập hợp. Cần phân biệt ngữ nghĩa tập hợp với bag/multiset vì bag cho phép bản sao và có quy tắc bội số riêng.

Phân tích: Đọc truy vấn theo các giai đoạn logic: FROM/JOIN tạo dữ liệu nguồn, WHERE lọc từng hàng, GROUP BY tạo nhóm, HAVING lọc nhóm, SELECT xác định biểu thức đầu ra và ORDER BY sắp kết quả. Với phép join, kiểm tra điều kiện ghép và các hàng không khớp trước khi tính kết quả.

Đối chiếu phương án: A. maker (type='Laptop' (Product)) \[maker (type='Laptop' (Product)); B. maker (type='Laptop' (Product)) \maker (type='PC' (Product)); D. Omaker (type='PC' (Product)) \ [maker ([type='PC' (Product)). Các lựa chọn này không đáp ứng định nghĩa/điều kiện được hỏi; khi ôn, hãy chỉ ra chính xác điểm lệch của từng lựa chọn thay vì học thuộc mỗi chữ cái.
\end{question}

\begin{question}{87}
% [ID gốc: 4.17]
% [Nguồn: DBI202 Order 4 Q17]
In the relational model of data, what is the result of an algebraic query language?
\begin{choices}
\choice Data series
\choice Data file
\choice Relation
\choice database
\end{choices}
\textbf{Đáp án:} C
\textbf{Giải thích:} Đáp án theo dữ liệu đề: C — C. Relation.

Nền tảng kiến thức: Đại số quan hệ thao tác trên quan hệ: selection lọc hàng, projection chọn cột, join ghép bộ theo điều kiện, còn union/difference/intersection là các phép tập hợp. Cần phân biệt ngữ nghĩa tập hợp với bag/multiset vì bag cho phép bản sao và có quy tắc bội số riêng.

Phân tích: Đọc truy vấn theo các giai đoạn logic: FROM/JOIN tạo dữ liệu nguồn, WHERE lọc từng hàng, GROUP BY tạo nhóm, HAVING lọc nhóm, SELECT xác định biểu thức đầu ra và ORDER BY sắp kết quả. Với phép join, kiểm tra điều kiện ghép và các hàng không khớp trước khi tính kết quả.

Đối chiếu phương án: A. Data series; B. Data file; D. database. Các lựa chọn này không đáp ứng định nghĩa/điều kiện được hỏi; khi ôn, hãy chỉ ra chính xác điểm lệch của từng lựa chọn thay vì học thuộc mỗi chữ cái.
\end{question}

\begin{question}{88}
% [ID gốc: 4.20]
% [Nguồn: DBI202 Order 4 Q20]
The result of which operation contains all pairs of tuples from the two relations, regardless of whether their attribute values match?
\begin{choices}
\choice Join
\choice Cartesian product
\choice Intersection
\choice Set difference
\end{choices}
\textbf{Đáp án:} B
\textbf{Giải thích:} Đáp án theo dữ liệu đề: B — B. Cartesian product.

Lời giải gốc: Phép tích Descartes (Cartesian product) kết hợp mọi bộ dữ liệu của quan hệ này với mọi bộ dữ liệu của quan hệ kia mà không cần điều kiện khớp thuộc tính.

Nền tảng kiến thức: Đại số quan hệ thao tác trên quan hệ: selection lọc hàng, projection chọn cột, join ghép bộ theo điều kiện, còn union/difference/intersection là các phép tập hợp. Cần phân biệt ngữ nghĩa tập hợp với bag/multiset vì bag cho phép bản sao và có quy tắc bội số riêng.

Phân tích: Các hàng dữ liệu của một relation được gọi là tuple; header chứa tên thuộc tính không phải tuple dữ liệu. Trong bảng, tuple tương ứng một bản ghi, còn attribute tương ứng cột.

Đối chiếu phương án: A. Join; C. Intersection; D. Set difference. Các lựa chọn này không đáp ứng định nghĩa/điều kiện được hỏi; khi ôn, hãy chỉ ra chính xác điểm lệch của từng lựa chọn thay vì học thuộc mỗi chữ cái.
\end{question}

\begin{question}{89}
% [ID gốc: 4.21]
% [Nguồn: DBI202 Order 4 Q21]
Which syntax is commonly used for the rename operation in Relational Algebra?
\begin{choices}
\choice RENAME relation TO new_relation
\choice RENAME attribute IN relation TO new_attribute
\choice ALTER TABLE relation RENAME TO new_relation
\choice SELECT * FROM relation AS new_relation
\end{choices}
\textbf{Đáp án:} A
\textbf{Giải thích:} Đáp án theo dữ liệu đề: D — D. SELECT * FROM relation AS new_relation.

Lời giải gốc: Trong SQL (ngôn ngữ dựa trên đại số quan hệ), từ khóa AS được dùng làm bí danh (alias) để đổi tên tạm thời cho một bảng hoặc một tập kết quả của câu lệnh truy vấn.

Nền tảng kiến thức: Đại số quan hệ thao tác trên quan hệ: selection lọc hàng, projection chọn cột, join ghép bộ theo điều kiện, còn union/difference/intersection là các phép tập hợp. Cần phân biệt ngữ nghĩa tập hợp với bag/multiset vì bag cho phép bản sao và có quy tắc bội số riêng.

Phân tích: Phép tập hợp trên quan hệ cần so sánh các tuple theo thuộc tính tương ứng; union lấy hợp, intersection lấy phần chung và difference lấy phần có ở vế trái mà không có ở vế phải. Trước tiên kiểm tra tính tương thích schema, sau đó áp phép theo đúng hướng/toán tử.

Đối chiếu phương án: A. RENAME relation TO new_relation; B. RENAME attribute IN relation TO new_attribute; C. ALTER TABLE relation RENAME TO new_relation. Các lựa chọn này không đáp ứng định nghĩa/điều kiện được hỏi; khi ôn, hãy chỉ ra chính xác điểm lệch của từng lựa chọn thay vì học thuộc mỗi chữ cái.
\end{question}

\begin{question}{90}
% [ID gốc: 4.24]
% [Nguồn: DBI202 Order 4 Q24]
Which of the following operations cannot be used in relational algebra? (Choose two answers)
\begin{choices}
\choice Union
\choice Divide
\choice Select
\choice Loop
\end{choices}
\textbf{Đáp án:} BD
\textbf{Giải thích:} Đáp án theo dữ liệu đề: B, D — B. Divide; D. Loop.

Lời giải gốc: Trong đại số quan hệ tiêu chuẩn, phép toán Chia (Divide/Division) tuy có tồn tại dưới dạng phép toán mở rộng nâng cao, nhưng xét theo tính chất cốt lõi của câu hỏi (hoặc hệ thống chấm điểm của bài này) và phép Vòng lặp (Loop) hoàn toàn không thuộc về đại số quan hệ thuần túy (đại số quan hệ không chứa các cấu trúc điều khiển lập trình như Loop).

Nền tảng kiến thức: Đại số quan hệ thao tác trên quan hệ: selection lọc hàng, projection chọn cột, join ghép bộ theo điều kiện, còn union/difference/intersection là các phép tập hợp. Cần phân biệt ngữ nghĩa tập hợp với bag/multiset vì bag cho phép bản sao và có quy tắc bội số riêng.

Phân tích: Phép tập hợp trên quan hệ cần so sánh các tuple theo thuộc tính tương ứng; union lấy hợp, intersection lấy phần chung và difference lấy phần có ở vế trái mà không có ở vế phải. Trước tiên kiểm tra tính tương thích schema, sau đó áp phép theo đúng hướng/toán tử.

Đối chiếu phương án: A. Union; C. Select. Các lựa chọn này không đáp ứng định nghĩa/điều kiện được hỏi; khi ôn, hãy chỉ ra chính xác điểm lệch của từng lựa chọn thay vì học thuộc mỗi chữ cái.
\end{question}

\begin{question}{91}
% [ID gốc: 4.27]
% [Nguồn: DBI202 Order 4 Q27]
Which relational algebra operator removes attributes from a relation?
\begin{choices}
\choice Selection (σ)
\choice Projection (π)
\choice Join (⨝)
\choice Union (∪)
\end{choices}
\textbf{Đáp án:} B
\textbf{Giải thích:} Đáp án theo dữ liệu đề: B — B. Projection (π).

Lời giải gốc: Phép chiếu (Projection - π) được dùng để chọn ra một số thuộc tính (cột) nhất định và loại bỏ đi những thuộc tính không cần thiết khác khỏi quan hệ ban đầu.

Nền tảng kiến thức: Đại số quan hệ thao tác trên quan hệ: selection lọc hàng, projection chọn cột, join ghép bộ theo điều kiện, còn union/difference/intersection là các phép tập hợp. Cần phân biệt ngữ nghĩa tập hợp với bag/multiset vì bag cho phép bản sao và có quy tắc bội số riêng.

Phân tích: Phép tập hợp trên quan hệ cần so sánh các tuple theo thuộc tính tương ứng; union lấy hợp, intersection lấy phần chung và difference lấy phần có ở vế trái mà không có ở vế phải. Trước tiên kiểm tra tính tương thích schema, sau đó áp phép theo đúng hướng/toán tử.

Đối chiếu phương án: A. Selection (σ); C. Join (⨝); D. Union (∪). Các lựa chọn này không đáp ứng định nghĩa/điều kiện được hỏi; khi ôn, hãy chỉ ra chính xác điểm lệch của từng lựa chọn thay vì học thuộc mỗi chữ cái.
\end{question}

\begin{question}{92}
% [ID gốc: 4.28]
% [Nguồn: DBI202 Order 4 Q28]
What is the purpose of the RENAME operation in relational algebra? (Choose three answers)
\begin{choices}
\choice To avoid ambiguity, especially when joining relations with the same attribute names.
\choice To make query expressions clearer and easier to read.
\choice To resolve naming conflicts in operations like UNION
\choice To search, index, and join tables quickly.
\end{choices}
\textbf{Đáp án:} ABC
\textbf{Giải thích:} Đáp án theo dữ liệu đề: A, B, C — A. To avoid ambiguity, especially when joining relations with the same attribute names.; B. To make query expressions clearer and easier to read.; C. To resolve naming conflicts in operations like UNION.

Lời giải gốc: Thao tác đổi tên (RENAME) giúp tránh nhập nhằng dữ liệu khi kết nối các bảng có trùng tên cột (A), làm biểu thức truy vấn rõ ràng dễ đọc hơn (B) và giải quyết xung đột tên trong các phép toán tập hợp như UNION (C). Nó không có chức năng tăng tốc độ tìm kiếm hay tạo chỉ mục (D).

Nền tảng kiến thức: Đại số quan hệ thao tác trên quan hệ: selection lọc hàng, projection chọn cột, join ghép bộ theo điều kiện, còn union/difference/intersection là các phép tập hợp. Cần phân biệt ngữ nghĩa tập hợp với bag/multiset vì bag cho phép bản sao và có quy tắc bội số riêng.

Phân tích: Phép tập hợp trên quan hệ cần so sánh các tuple theo thuộc tính tương ứng; union lấy hợp, intersection lấy phần chung và difference lấy phần có ở vế trái mà không có ở vế phải. Trước tiên kiểm tra tính tương thích schema, sau đó áp phép theo đúng hướng/toán tử.

Đối chiếu phương án: D. To search, index, and join tables quickly.. Các lựa chọn này không đáp ứng định nghĩa/điều kiện được hỏi; khi ôn, hãy chỉ ra chính xác điểm lệch của từng lựa chọn thay vì học thuộc mỗi chữ cái.
\end{question}

\begin{question}{93}
% [ID gốc: 4.33]
% [Nguồn: DBI202 Order 4 Q33]
A natural join joins two relations by:
\begin{choices}
\choice All attributes
\choice Only key attributes
\choice Matching common attributes
\choice Cartesian product
\end{choices}
\textbf{Đáp án:} C
\textbf{Giải thích:} Đáp án theo dữ liệu đề: C — C. Matching common attributes.

Lời giải gốc: Phép kết nối tự nhiên (Natural join) thực hiện nối hai quan hệ dựa trên các thuộc tính chung có cùng tên và kiểu dữ liệu. ---

Nền tảng kiến thức: Đại số quan hệ thao tác trên quan hệ: selection lọc hàng, projection chọn cột, join ghép bộ theo điều kiện, còn union/difference/intersection là các phép tập hợp. Cần phân biệt ngữ nghĩa tập hợp với bag/multiset vì bag cho phép bản sao và có quy tắc bội số riêng.

Phân tích: Join ghép các bộ thỏa điều kiện. Inner join loại bộ không ghép được; LEFT/RIGHT OUTER JOIN giữ mọi bộ từ phía được chỉ định và điền NULL phía thiếu. Natural join tự ghép các thuộc tính trùng tên và chỉ xuất một cột cho mỗi thuộc tính ghép.

Đối chiếu phương án: A. All attributes; B. Only key attributes; D. Cartesian product. Các lựa chọn này không đáp ứng định nghĩa/điều kiện được hỏi; khi ôn, hãy chỉ ra chính xác điểm lệch của từng lựa chọn thay vì học thuộc mỗi chữ cái.
\end{question}

\begin{question}{94}
% [ID gốc: 4.35]
% [Nguồn: DBI202 Order 4 Q35]
The operation that combines two relations by merging pairs of tuples, one from each relation, into a single tuple is:
\begin{choices}
\choice Select
\choice Join
\choice Union
\choice Intersection
\end{choices}
\textbf{Đáp án:} B
\textbf{Giải thích:} Đáp án theo dữ liệu đề: B — B. Join.

Lời giải gốc: Phép kết nối (Join) cho phép kết hợp hai quan hệ bằng cách ghép nối các bộ dữ liệu từ hai bảng thỏa mãn điều kiện kết nối.

Nền tảng kiến thức: Đại số quan hệ thao tác trên quan hệ: selection lọc hàng, projection chọn cột, join ghép bộ theo điều kiện, còn union/difference/intersection là các phép tập hợp. Cần phân biệt ngữ nghĩa tập hợp với bag/multiset vì bag cho phép bản sao và có quy tắc bội số riêng.

Phân tích: Các hàng dữ liệu của một relation được gọi là tuple; header chứa tên thuộc tính không phải tuple dữ liệu. Trong bảng, tuple tương ứng một bản ghi, còn attribute tương ứng cột.

Đối chiếu phương án: A. Select; C. Union; D. Intersection. Các lựa chọn này không đáp ứng định nghĩa/điều kiện được hỏi; khi ôn, hãy chỉ ra chính xác điểm lệch của từng lựa chọn thay vì học thuộc mỗi chữ cái.
\end{question}

\begin{question}{95}
% [ID gốc: 4.36]
% [Nguồn: DBI202 Order 4 Q36]
In relational algebra, the natural join operation can be viewed as:
\begin{choices}
\choice A projection followed by a union
\choice A selection followed by a projection
\choice A Cartesian product followed by a selection and projection
\choice A product followed by a union
\end{choices}
\textbf{Đáp án:} C
\textbf{Giải thích:} Đáp án theo dữ liệu đề: C — C. A Cartesian product followed by a selection and projection.

Lời giải gốc: Phép kết tự nhiên (Natural Join) tương đương với việc thực hiện tích Descartes của hai bảng, sau đó lọc các hàng có giá trị thuộc tính chung bằng nhau (Selection) và loại bỏ các cột trùng lặp (Projection).

Nền tảng kiến thức: Đại số quan hệ thao tác trên quan hệ: selection lọc hàng, projection chọn cột, join ghép bộ theo điều kiện, còn union/difference/intersection là các phép tập hợp. Cần phân biệt ngữ nghĩa tập hợp với bag/multiset vì bag cho phép bản sao và có quy tắc bội số riêng.

Phân tích: Phép tập hợp trên quan hệ cần so sánh các tuple theo thuộc tính tương ứng; union lấy hợp, intersection lấy phần chung và difference lấy phần có ở vế trái mà không có ở vế phải. Trước tiên kiểm tra tính tương thích schema, sau đó áp phép theo đúng hướng/toán tử.

Đối chiếu phương án: A. A projection followed by a union; B. A selection followed by a projection; D. A product followed by a union. Các lựa chọn này không đáp ứng định nghĩa/điều kiện được hỏi; khi ôn, hãy chỉ ra chính xác điểm lệch của từng lựa chọn thay vì học thuộc mỗi chữ cái.
\end{question}

\begin{question}{96}
% [ID gốc: 4.37]
% [Nguồn: DBI202 Order 4 Q37]
Which operator is used to produce a new relation from a relation R that contains only some of R's columns?
\begin{choices}
\choice Selection
\choice Projection
\choice Intersection
\choice All of the others
\end{choices}
\textbf{Đáp án:} B
\textbf{Giải thích:} Đáp án theo dữ liệu đề: B — B. Projection.

Lời giải gốc: Phép chiếu (Projection) trong đại số quan hệ lọc ra các cột được chỉ định từ quan hệ ban đầu.

Nền tảng kiến thức: Đại số quan hệ thao tác trên quan hệ: selection lọc hàng, projection chọn cột, join ghép bộ theo điều kiện, còn union/difference/intersection là các phép tập hợp. Cần phân biệt ngữ nghĩa tập hợp với bag/multiset vì bag cho phép bản sao và có quy tắc bội số riêng.

Phân tích: Với lựa chọn tổng hợp, không chọn theo cảm tính: kiểm tra từng phát biểu độc lập. “All of the others/above” chỉ đúng khi mọi phát biểu còn lại đúng; “None” chỉ đúng khi không phát biểu nào trong nhóm thỏa yêu cầu. Đáp án lưu cần được đọc cùng các mệnh đề thành phần.

Đối chiếu phương án: A. Selection; C. Intersection; D. All of the others. Các lựa chọn này không đáp ứng định nghĩa/điều kiện được hỏi; khi ôn, hãy chỉ ra chính xác điểm lệch của từng lựa chọn thay vì học thuộc mỗi chữ cái.
\end{question}

\begin{question}{97}
% [ID gốc: 4.38]
% [Nguồn: DBI202 Order 4 Q38]
The intersection operator is used to get the ____ tuples.
\begin{choices}
\choice Different
\choice Common
\choice All
\choice Repeating
\end{choices}
\textbf{Đáp án:} B
\textbf{Giải thích:} Đáp án theo dữ liệu đề: B — B. Common.

Lời giải gốc: Phép giao (Intersection) được dùng để lấy các bộ dữ liệu chung xuất hiện ở cả hai quan hệ.

Nền tảng kiến thức: Đại số quan hệ thao tác trên quan hệ: selection lọc hàng, projection chọn cột, join ghép bộ theo điều kiện, còn union/difference/intersection là các phép tập hợp. Cần phân biệt ngữ nghĩa tập hợp với bag/multiset vì bag cho phép bản sao và có quy tắc bội số riêng.

Phân tích: Các hàng dữ liệu của một relation được gọi là tuple; header chứa tên thuộc tính không phải tuple dữ liệu. Trong bảng, tuple tương ứng một bản ghi, còn attribute tương ứng cột.

Đối chiếu phương án: A. Different; C. All; D. Repeating. Các lựa chọn này không đáp ứng định nghĩa/điều kiện được hỏi; khi ôn, hãy chỉ ra chính xác điểm lệch của từng lựa chọn thay vì học thuộc mỗi chữ cái.
\end{question}

\begin{question}{98}
% [ID gốc: 4.41]
% [Nguồn: DBI202 Order 4 Q41]
Which operation is commonly used to display selected columns (attributes) from the original relation?
\begin{choices}
\choice Union
\choice Intersection
\choice Closure
\choice Projection
\end{choices}
\textbf{Đáp án:} D
\textbf{Giải thích:} Đáp án theo dữ liệu đề: D — D. Projection.

Lời giải gốc: Phép chiếu (Projection) chọn ra các cột (thuộc tính) cụ thể từ quan hệ ban đầu.

Nền tảng kiến thức: Đại số quan hệ thao tác trên quan hệ: selection lọc hàng, projection chọn cột, join ghép bộ theo điều kiện, còn union/difference/intersection là các phép tập hợp. Cần phân biệt ngữ nghĩa tập hợp với bag/multiset vì bag cho phép bản sao và có quy tắc bội số riêng.

Phân tích: Phép tập hợp trên quan hệ cần so sánh các tuple theo thuộc tính tương ứng; union lấy hợp, intersection lấy phần chung và difference lấy phần có ở vế trái mà không có ở vế phải. Trước tiên kiểm tra tính tương thích schema, sau đó áp phép theo đúng hướng/toán tử.

Đối chiếu phương án: A. Union; B. Intersection; C. Closure. Các lựa chọn này không đáp ứng định nghĩa/điều kiện được hỏi; khi ôn, hãy chỉ ra chính xác điểm lệch của từng lựa chọn thay vì học thuộc mỗi chữ cái.
\end{question}

\begin{question}{99}
% [ID gốc: 4.42]
% [Nguồn: DBI202 Order 4 Q42]
Which of the following is a relational algebra operation?
\begin{choices}
\choice Sorting
\choice Indexing
\choice Selection
\choice Comping
\end{choices}
\textbf{Đáp án:} C
\textbf{Giải thích:} Đáp án theo dữ liệu đề: C — C. Selection.

Lời giải gốc: Phép chọn (Selection) là một trong các phép toán cơ bản của đại số quan hệ, dùng để lọc các bộ thỏa mãn điều kiện.

Nền tảng kiến thức: Đại số quan hệ thao tác trên quan hệ: selection lọc hàng, projection chọn cột, join ghép bộ theo điều kiện, còn union/difference/intersection là các phép tập hợp. Cần phân biệt ngữ nghĩa tập hợp với bag/multiset vì bag cho phép bản sao và có quy tắc bội số riêng.

Phân tích: Phép tập hợp trên quan hệ cần so sánh các tuple theo thuộc tính tương ứng; union lấy hợp, intersection lấy phần chung và difference lấy phần có ở vế trái mà không có ở vế phải. Trước tiên kiểm tra tính tương thích schema, sau đó áp phép theo đúng hướng/toán tử.

Đối chiếu phương án: A. Sorting; B. Indexing; D. Comping. Các lựa chọn này không đáp ứng định nghĩa/điều kiện được hỏi; khi ôn, hãy chỉ ra chính xác điểm lệch của từng lựa chọn thay vì học thuộc mỗi chữ cái.
\end{question}

\begin{question}{100}
% [ID gốc: 4.43]
% [Nguồn: DBI202 Order 4 Q43]
What does the relational algebra operation "Projection" do?
\begin{choices}
\choice Selects rows based on a condition.
\choice Combines rows from two tables.
\choice Selects columns from a table.
\choice Renames a table.
\end{choices}
\textbf{Đáp án:} C
\textbf{Giải thích:} Đáp án theo dữ liệu đề: C — C. Selects columns from a table..

Lời giải gốc: Phép chiếu (Projection) trong đại số quan hệ thực hiện việc trích xuất các cột (thuộc tính) chỉ định từ một quan hệ.

Nền tảng kiến thức: Đại số quan hệ thao tác trên quan hệ: selection lọc hàng, projection chọn cột, join ghép bộ theo điều kiện, còn union/difference/intersection là các phép tập hợp. Cần phân biệt ngữ nghĩa tập hợp với bag/multiset vì bag cho phép bản sao và có quy tắc bội số riêng.

Phân tích: Phép tập hợp trên quan hệ cần so sánh các tuple theo thuộc tính tương ứng; union lấy hợp, intersection lấy phần chung và difference lấy phần có ở vế trái mà không có ở vế phải. Trước tiên kiểm tra tính tương thích schema, sau đó áp phép theo đúng hướng/toán tử.

Đối chiếu phương án: A. Selects rows based on a condition.; B. Combines rows from two tables.; D. Renames a table.. Các lựa chọn này không đáp ứng định nghĩa/điều kiện được hỏi; khi ôn, hãy chỉ ra chính xác điểm lệch của từng lựa chọn thay vì học thuộc mỗi chữ cái.
\end{question}

\begin{question}{101}
% [ID gốc: 4.44]
% [Nguồn: DBI202 Order 4 Q44]
For two relations Employee and Customer sharing the columns (Name, City, Gender), what is the result of the Union operation between Employee and Customer?
\begin{choices}
\choice Only the rows common to both relations
\choice All rows from Employee only
\choice All rows from Customer only
\choice All distinct rows from both relations, with duplicates removed
\end{choices}
\textbf{Đáp án:} D
\textbf{Giải thích:} Đáp án theo dữ liệu đề: D — D. All distinct rows from both relations, with duplicates removed.

Lời giải gốc: Phép hợp (UNION) gộp các bản ghi từ hai quan hệ và tự động loại bỏ các dòng bị trùng lặp hoàn toàn, giữ lại danh sách đầy đủ các bản ghi độc nhất.

Nền tảng kiến thức: Đại số quan hệ thao tác trên quan hệ: selection lọc hàng, projection chọn cột, join ghép bộ theo điều kiện, còn union/difference/intersection là các phép tập hợp. Cần phân biệt ngữ nghĩa tập hợp với bag/multiset vì bag cho phép bản sao và có quy tắc bội số riêng.

Phân tích: Phép tập hợp trên quan hệ cần so sánh các tuple theo thuộc tính tương ứng; union lấy hợp, intersection lấy phần chung và difference lấy phần có ở vế trái mà không có ở vế phải. Trước tiên kiểm tra tính tương thích schema, sau đó áp phép theo đúng hướng/toán tử.

Đối chiếu phương án: A. Only the rows common to both relations; B. All rows from Employee only; C. All rows from Customer only. Các lựa chọn này không đáp ứng định nghĩa/điều kiện được hỏi; khi ôn, hãy chỉ ra chính xác điểm lệch của từng lựa chọn thay vì học thuộc mỗi chữ cái.
\end{question}

\begin{question}{102}
% [ID gốc: 5.11]
% [Nguồn: DBI202 Order 5 Q11]
When we apply set operators (UNION, INTERSECT, EXCEPT) to two relations R and S, which conditions on R and S must be satisfied?
\begin{choices}
\choice R and S must have schemas with identical sets of attributes and the types (domains) for each attributes must be the same in R and S
\choice Before we compute the set-theoretic union, intersection, or difference of sets of tuples, the columns of R and S must be ordered so that the order of attributes is the same for both relations
\choice All of the others
\end{choices}
\textbf{Đáp án:} C
\textbf{Giải thích:} Đáp án theo dữ liệu đề: C — C. All of the others.

Nền tảng kiến thức: SQL tách truy vấn dữ liệu (DML) khỏi định nghĩa cấu trúc (DDL) và quyền truy cập (DCL). Với SELECT, cần xác định bảng nguồn, điều kiện lọc, phép kết hợp, nhóm và điều kiện nhóm; các ràng buộc, view, thủ tục và trigger đảm nhiệm các vai trò khác nhau.

Phân tích: Phép tập hợp trên quan hệ cần so sánh các tuple theo thuộc tính tương ứng; union lấy hợp, intersection lấy phần chung và difference lấy phần có ở vế trái mà không có ở vế phải. Trước tiên kiểm tra tính tương thích schema, sau đó áp phép theo đúng hướng/toán tử.

Đối chiếu phương án: A. R and S must have schemas with identical sets of attributes and the types (domains) for each a…; B. Before we compute the set-theoretic union, intersection, or difference of sets of tuples, the…. Các lựa chọn này không đáp ứng định nghĩa/điều kiện được hỏi; khi ôn, hãy chỉ ra chính xác điểm lệch của từng lựa chọn thay vì học thuộc mỗi chữ cái.
\end{question}

\begin{question}{103}
% [ID gốc: 12.1]
% [Nguồn: DBI202 Order 12 Q1]
Consider the join of a relation R with a relation S. If R has m tuples and S has n tuples, then the maximum size of the join is:
\begin{choices}
\choice m * n
\choice m + n
\choice (m + n) / 2
\choice 2(m + n)
\end{choices}
\textbf{Đáp án:} A
\textbf{Giải thích:} Đáp án theo dữ liệu đề: A — A. m * n.

Lời giải gốc: Trong trường hợp xấu nhất (mọi bộ trong R đều khớp với mọi bộ trong S, tương tự phép tích Descartes), số bộ tối đa trong kết quả phép join là tích số bộ của hai quan hệ, tức m * n.

Nền tảng kiến thức: Đánh giá toán tử quan hệ chọn thuật toán theo kích thước quan hệ, số trang, mức chọn lọc, thứ tự dữ liệu và chỉ mục sẵn có. Với join, nested-loop, sort-merge và hash join có chi phí/điều kiện phù hợp khác nhau; kích thước kết quả không đồng nhất với chi phí thực thi.

Phân tích: Các hàng dữ liệu của một relation được gọi là tuple; header chứa tên thuộc tính không phải tuple dữ liệu. Trong bảng, tuple tương ứng một bản ghi, còn attribute tương ứng cột.

Đối chiếu phương án: B. m + n; C. (m + n) / 2; D. 2(m + n). Các lựa chọn này không đáp ứng định nghĩa/điều kiện được hỏi; khi ôn, hãy chỉ ra chính xác điểm lệch của từng lựa chọn thay vì học thuộc mỗi chữ cái.
\end{question}

\begin{question}{104}
% [ID gốc: 22.1]
% [Nguồn: DBI202 Order 22 Q1]
______ is a set of markup declarations that define a document type for SGML-family markup languages (SGML, XML, HTML)
\begin{choices}
\choice Document Type Definition (DTD)
\choice Document Definition Language
\choice Document Manipulation Language
\choice Document Control Language
\end{choices}
\textbf{Đáp án:} A
\textbf{Giải thích:} Đáp án theo dữ liệu đề: A — A. Document Type Definition (DTD).

Nền tảng kiến thức: CSDL Internet kết hợp ứng dụng web với lưu trữ/truy vấn dữ liệu; XML là dạng dữ liệu có cấu trúc phân cấp nhưng có thể linh hoạt hơn bảng cố định. DTD mô tả cấu trúc và các phần tử/thuộc tính hợp lệ; XML có thể được coi là dữ liệu bán cấu trúc.

Phân tích: XML biểu diễn dữ liệu phân cấp bằng phần tử/thuộc tính. Tài liệu well-formed tuân cú pháp XML; DTD mô tả cấu trúc/khai báo hợp lệ. Ký hiệu xuất hiện 0 hoặc 1 lần thường là `?` trong content model DTD, khác `*` (0 trở lên) và `+` (1 trở lên).

Đối chiếu phương án: B. Document Definition Language; C. Document Manipulation Language; D. Document Control Language. Các lựa chọn này không đáp ứng định nghĩa/điều kiện được hỏi; khi ôn, hãy chỉ ra chính xác điểm lệch của từng lựa chọn thay vì học thuộc mỗi chữ cái.
\end{question}

\begin{question}{105}
% [ID gốc: 22.2]
% [Nguồn: DBI202 Order 22 Q2]
A ____ document defines the XML document structure with a list of legal elements and attributes
\begin{choices}
\choice XML
\choice DTD
\choice XSLT
\choice SGML
\end{choices}
\textbf{Đáp án:} B
\textbf{Giải thích:} Đáp án theo dữ liệu đề: B — B. DTD.

Nền tảng kiến thức: CSDL Internet kết hợp ứng dụng web với lưu trữ/truy vấn dữ liệu; XML là dạng dữ liệu có cấu trúc phân cấp nhưng có thể linh hoạt hơn bảng cố định. DTD mô tả cấu trúc và các phần tử/thuộc tính hợp lệ; XML có thể được coi là dữ liệu bán cấu trúc.

Phân tích: XML biểu diễn dữ liệu phân cấp bằng phần tử/thuộc tính. Tài liệu well-formed tuân cú pháp XML; DTD mô tả cấu trúc/khai báo hợp lệ. Ký hiệu xuất hiện 0 hoặc 1 lần thường là `?` trong content model DTD, khác `*` (0 trở lên) và `+` (1 trở lên).

Đối chiếu phương án: A. XML; C. XSLT; D. SGML. Các lựa chọn này không đáp ứng định nghĩa/điều kiện được hỏi; khi ôn, hãy chỉ ra chính xác điểm lệch của từng lựa chọn thay vì học thuộc mỗi chữ cái.
\end{question}

\begin{question}{106}
% [ID gốc: 22.3]
% [Nguồn: DBI202 Order 22 Q3]
Choose the right symbol to declare zero or one occurrence of an element in XML DTD:
\begin{choices}
\choice ?
\choice *
\choice +
\choice None of the others
\end{choices}
\textbf{Đáp án:} A
\textbf{Giải thích:} Đáp án theo dữ liệu đề: A — A. .

Nền tảng kiến thức: CSDL Internet kết hợp ứng dụng web với lưu trữ/truy vấn dữ liệu; XML là dạng dữ liệu có cấu trúc phân cấp nhưng có thể linh hoạt hơn bảng cố định. DTD mô tả cấu trúc và các phần tử/thuộc tính hợp lệ; XML có thể được coi là dữ liệu bán cấu trúc.

Phân tích: XML biểu diễn dữ liệu phân cấp bằng phần tử/thuộc tính. Tài liệu well-formed tuân cú pháp XML; DTD mô tả cấu trúc/khai báo hợp lệ. Ký hiệu xuất hiện 0 hoặc 1 lần thường là `?` trong content model DTD, khác `*` (0 trở lên) và `+` (1 trở lên).

Đối chiếu phương án: B. ; C. [D]. Các lựa chọn này không đáp ứng định nghĩa/điều kiện được hỏi; khi ôn, hãy chỉ ra chính xác điểm lệch của từng lựa chọn thay vì học thuộc mỗi chữ cái.
\end{question}

\begin{question}{107}
% [ID gốc: 22.4]
% [Nguồn: DBI202 Order 22 Q4]
Which of the following statements is TRUE about a well-formed XML document?
\begin{choices}
\choice XML tags are case-insensitive
\choice All XML elements must have a closing tag and be properly nested
\choice XML attribute values do not need to be enclosed in quotation marks
\choice An XML document can have more than one root element
\end{choices}
\textbf{Đáp án:} B
\textbf{Giải thích:} Đáp án đúng: B — All XML elements must have a closing tag and be properly nested.

Nền tảng kiến thức: Một tài liệu XML được coi là chuẩn cú pháp (Well-formed XML) khi tuân thủ nghiêm ngặt các quy tắc:
1. Có duy nhất một phần tử gốc (single root element) bao bọc toàn bộ nội dung tài liệu.
2. Mọi thẻ mở đều phải có thẻ đóng tương ứng và phân biệt chữ hoa/thường (Case-sensitive).
3. Các thẻ lồng nhau phải đóng/mở theo đúng thứ tự (Properly nested), không được chồng lấn nhau.
4. Các giá trị thuộc tính phải luôn được đặt trong dấu ngoặc kép hoặc ngoặc đơn (ví dụ: title="value").

Phân tích các phương án:
- Phương án B đúng vì đây là quy tắc cốt lõi của XML well-formed.
- Phương án A sai vì XML luôn phân biệt chữ hoa và chữ thường (Case-sensitive).
- Phương án C sai vì thuộc tính trong XML bắt buộc phải nằm trong dấu ngoặc kép.
- Phương án D sai vì tài liệu XML chỉ được phép có duy nhất 1 phần tử gốc bao bọc bên ngoài.
\end{question}

\begin{question}{108}
% [ID gốc: 22.5]
% [Nguồn: DBI202 Order 22 Q5]
XML files are semi-structured data
\begin{choices}
\choice True
\choice False
\end{choices}
\textbf{Đáp án:} A
\textbf{Giải thích:} Đáp án theo dữ liệu đề: A — A. True.

Nền tảng kiến thức: CSDL Internet kết hợp ứng dụng web với lưu trữ/truy vấn dữ liệu; XML là dạng dữ liệu có cấu trúc phân cấp nhưng có thể linh hoạt hơn bảng cố định. DTD mô tả cấu trúc và các phần tử/thuộc tính hợp lệ; XML có thể được coi là dữ liệu bán cấu trúc.

Phân tích: XML biểu diễn dữ liệu phân cấp bằng phần tử/thuộc tính. Tài liệu well-formed tuân cú pháp XML; DTD mô tả cấu trúc/khai báo hợp lệ. Ký hiệu xuất hiện 0 hoặc 1 lần thường là `?` trong content model DTD, khác `*` (0 trở lên) và `+` (1 trở lên).

Đối chiếu phương án: B. False. Các lựa chọn này không đáp ứng định nghĩa/điều kiện được hỏi; khi ôn, hãy chỉ ra chính xác điểm lệch của từng lựa chọn thay vì học thuộc mỗi chữ cái.
\end{question}

\begin{question}{109}
% [ID gốc: 27.1]
% [Nguồn: DBI202 Order 27 Q1]
The Datalog rule H(x,y)  2 AND y  2 and y  2 AND y < 6. Relation S(x, y) has 3 tuples (2, 3), (3, 5), and (4, 6). What is about H?
\begin{choices}
\choice H has a tuple (3, 5)
\choice H has a tuple (2, 3)
\choice H has 2 tuples (2, 3) and (3, 5)
\choice H has 3 tuples (2, 3) and (3, 5) and (4, 6)
\end{choices}
\textbf{Đáp án:} A
\textbf{Giải thích:} Đáp án theo dữ liệu đề: A — A. H has a tuple (3, 5).

Nền tảng kiến thức: CSDL suy diễn dùng facts và rules để suy ra facts mới; Datalog thường biểu diễn head :- body, trong đó body là các điều kiện cần đồng thời đúng. Truy vấn đệ quy được tính lặp tới điểm cố định, nên cần kiểm tra biến, phép nối và điều kiện lọc trên từng bộ.

Phân tích: Các hàng dữ liệu của một relation được gọi là tuple; header chứa tên thuộc tính không phải tuple dữ liệu. Trong bảng, tuple tương ứng một bản ghi, còn attribute tương ứng cột.

Đối chiếu phương án: B. H has a tuple (2, 3); C. H has 2 tuples (2, 3) and (3, 5); D. H has 3 tuples (2, 3) and (3, 5) and (4, 6). Các lựa chọn này không đáp ứng định nghĩa/điều kiện được hỏi; khi ôn, hãy chỉ ra chính xác điểm lệch của từng lựa chọn thay vì học thuộc mỗi chữ cái.
\end{question}

\begin{question}{110}
% [ID gốc: 27.2]
% [Nguồn: DBI202 Order 27 Q2]
In Datalog rules with equality and inequality subgoals, which of the following conditions ensures that variable w is safe if z already appears in a positive relational subgoal?
\begin{choices}
\choice w = z
\choice w < z
\choice w <= max(x,y)
\choice w != z
\end{choices}
\textbf{Đáp án:} A
\textbf{Giải thích:} Trong Datalog, một biến w chưa xuất hiện trong quan hệ gốc vẫn an toàn (safe) nếu nó được gán bằng một biến z đã an toàn thông qua điều kiện đẳng thức w = z.
\end{question}

\begin{question}{111}
% [ID gốc: 27.3]
% [Nguồn: DBI202 Order 27 Q3]
Need to be deleted
Consider the Datalog rule H(xy)  2 AND y < 6. Relation S(x y) has 3 tuples (2.3). (3.5). and (4.6). What is about H?
\begin{choices}
\choice H has a tuple (2.3)
\choice H has 3 tuples (2.3) and (3.5) and (4.6)
\choice H has 2 tuples (2.3) and (3.5)
\choice H has a tuple (3.5)
\end{choices}
\textbf{Đáp án:} D
\textbf{Giải thích:} Đáp án theo dữ liệu đề: D — D. H has a tuple (3.5).

Nền tảng kiến thức: CSDL suy diễn dùng facts và rules để suy ra facts mới; Datalog thường biểu diễn head :- body, trong đó body là các điều kiện cần đồng thời đúng. Truy vấn đệ quy được tính lặp tới điểm cố định, nên cần kiểm tra biến, phép nối và điều kiện lọc trên từng bộ.

Phân tích: Các hàng dữ liệu của một relation được gọi là tuple; header chứa tên thuộc tính không phải tuple dữ liệu. Trong bảng, tuple tương ứng một bản ghi, còn attribute tương ứng cột.

Đối chiếu phương án: A. H has a tuple (2.3); B. H has 3 tuples (2.3) and (3.5) and (4.6); C. H has 2 tuples (2.3) and (3.5). Các lựa chọn này không đáp ứng định nghĩa/điều kiện được hỏi; khi ôn, hãy chỉ ra chính xác điểm lệch của từng lựa chọn thay vì học thuộc mỗi chữ cái.
\end{question}


\newpage
\section*{Chương 3: Design Theory for Relational Databases}

\begin{question}{112}
% [ID gốc: 15.1]
% [Nguồn: DBI202 Order 15 Q1]
A constraint of any kind on a relation is said to be trivial if it holds for every instance of the relation, regardless of what other constraints are assumed.
QN=9 (7961) Given a relation R(A,B,C,D). Which of the followings is trivial?
\begin{choices}
\choice A->AB
\choice A->->AB
\choice A->BCD
\choice A->->BCD
\end{choices}
\textbf{Đáp án:} D
\textbf{Giải thích:} Đáp án theo dữ liệu đề: D — D. A->->BCD.

Nền tảng kiến thức: Chuẩn hóa dùng phụ thuộc hàm và khóa để giảm dư thừa cũng như bất thường cập nhật. Tính bao đóng thuộc tính để tìm khóa; 2NF loại phụ thuộc bộ phận vào một phần khóa ghép, 3NF xử lý phụ thuộc không khóa kiểu bắc cầu, còn BCNF yêu cầu mọi vế trái không tầm thường là siêu khóa.

Phân tích: Để xác định một ràng buộc/phụ thuộc tầm thường, kiểm tra xem nó có đúng chỉ từ định nghĩa tập thuộc tính của quan hệ hay không, không cần giả định thêm về dữ liệu. Với MVD X ↠ Y, điều kiện tầm thường là Y ⊆ X hoặc X ∪ Y bằng toàn bộ thuộc tính; cần phân biệt ký hiệu ↠ với phụ thuộc hàm →.

Đối chiếu phương án: A. A->AB; B. A->->AB; C. A->BCD. Các lựa chọn này không đáp ứng định nghĩa/điều kiện được hỏi; khi ôn, hãy chỉ ra chính xác điểm lệch của từng lựa chọn thay vì học thuộc mỗi chữ cái.
\end{question}

\begin{question}{113}
% [ID gốc: 15.2]
% [Nguồn: DBI202 Order 15 Q2]
Từ A -> B và D, từ C -> B. Suy ra bộ AC là key
QN=10 (7972) Let R(ABCD) be a relation with functional dependencies
{A -> B,
C -> B,
B -> D}
What is the key for R (choose one)
\begin{choices}
\choice AB
\choice AC
\choice AD
\choice BD
\end{choices}
\textbf{Đáp án:} B
\textbf{Giải thích:} Đáp án theo dữ liệu đề: B — B. AC.

Nền tảng kiến thức: Chuẩn hóa dùng phụ thuộc hàm và khóa để giảm dư thừa cũng như bất thường cập nhật. Tính bao đóng thuộc tính để tìm khóa; 2NF loại phụ thuộc bộ phận vào một phần khóa ghép, 3NF xử lý phụ thuộc không khóa kiểu bắc cầu, còn BCNF yêu cầu mọi vế trái không tầm thường là siêu khóa.

Phân tích: Kiểm tra khóa bằng hai tiêu chí: tập thuộc tính phải xác định duy nhất mọi bộ (tính siêu khóa), và nếu là candidate key thì phải tối thiểu, tức bỏ bất kỳ thuộc tính nào cũng mất tính duy nhất. PRIMARY KEY chọn một candidate key và không nhận NULL.

Đối chiếu phương án: A. AB; C. AD; D. BD. Các lựa chọn này không đáp ứng định nghĩa/điều kiện được hỏi; khi ôn, hãy chỉ ra chính xác điểm lệch của từng lựa chọn thay vì học thuộc mỗi chữ cái.
\end{question}

\begin{question}{114}
% [ID gốc: 15.3]
% [Nguồn: DBI202 Order 15 Q3]
Consider the following functional dependencies
a,b -> c,d e,g,h -> f,j
a,c -> b,d p,q -> r,s
e,f,g -> h,i s -> t
f,g -> j q -> u
g,h -> i
Which of the following best describes the relation R(e,f,g,h,i,j)?
\begin{choices}
\choice R is in First Normal Form
\choice R is in Second Normal Form
\choice R is in Third Normal Form
\choice R is in Boyce Codd Normal Form
\end{choices}
\textbf{Đáp án:} A
\textbf{Giải thích:} Đáp án theo dữ liệu đề: A — A. R is in First Normal Form.

Nền tảng kiến thức: Chuẩn hóa dùng phụ thuộc hàm và khóa để giảm dư thừa cũng như bất thường cập nhật. Tính bao đóng thuộc tính để tìm khóa; 2NF loại phụ thuộc bộ phận vào một phần khóa ghép, 3NF xử lý phụ thuộc không khóa kiểu bắc cầu, còn BCNF yêu cầu mọi vế trái không tầm thường là siêu khóa.

Phân tích: Với phụ thuộc X → Y, mọi hai bộ có cùng giá trị X phải có cùng giá trị Y. Để kiểm tra khóa, tính X⁺ bằng cách áp dụng lặp các phụ thuộc; nếu bao đóng chứa toàn bộ thuộc tính thì X là siêu khóa, sau đó kiểm tra tính tối thiểu.

Đối chiếu phương án: B. R is in Second Normal Form; C. R is in Third Normal Form; D. R is in Boyce Codd Normal Form. Các lựa chọn này không đáp ứng định nghĩa/điều kiện được hỏi; khi ôn, hãy chỉ ra chính xác điểm lệch của từng lựa chọn thay vì học thuộc mỗi chữ cái.
\end{question}

\begin{question}{115}
% [ID gốc: 15.4]
% [Nguồn: DBI202 Order 15 Q4]
Based on the given functional dependencies {A -> B; B -> A; A -> D; D -> B}, we can derive that {A} and {B} are both candidate keys for relation R(ABCD). Since all attributes in R are either part of a candidate key or fully functionally dependent on it, we can conclude that relation R is in Third Normal Form (3NF)
QN=13 (7978) The relation R(ABCD) has following FDs:
{ A -> B ; B -> A ; A -> D ; D -> B }
\begin{choices}
\choice R is in 3NF
\choice R is not in 3NF
\choice R is not in 2NF
\choice None of the others
\end{choices}
\textbf{Đáp án:} C
\textbf{Giải thích:} Đáp án theo dữ liệu đề: A — A. R is in 3NF.

Nền tảng kiến thức: Chuẩn hóa dùng phụ thuộc hàm và khóa để giảm dư thừa cũng như bất thường cập nhật. Tính bao đóng thuộc tính để tìm khóa; 2NF loại phụ thuộc bộ phận vào một phần khóa ghép, 3NF xử lý phụ thuộc không khóa kiểu bắc cầu, còn BCNF yêu cầu mọi vế trái không tầm thường là siêu khóa.

Phân tích: Đánh giá dạng chuẩn theo thứ tự: 1NF yêu cầu giá trị nguyên tố/không lặp theo mô hình quan hệ; 2NF loại phụ thuộc bộ phận của thuộc tính không khóa vào một phần khóa ghép; 3NF kiểm soát phụ thuộc bắc cầu theo điều kiện prime attribute; BCNF yêu cầu mọi determinant không tầm thường là siêu khóa.

Đối chiếu phương án: B. R is not in 3NF; C. R is not in 2NF; D. None of the others. Các lựa chọn này không đáp ứng định nghĩa/điều kiện được hỏi; khi ôn, hãy chỉ ra chính xác điểm lệch của từng lựa chọn thay vì học thuộc mỗi chữ cái.
\end{question}

\begin{question}{116}
% [ID gốc: 15.5]
% [Nguồn: DBI202 Order 15 Q5]
Given a R(A,B,C,D) with the following FDs:
AB->D
BC->A
AD->B
CD->B
AD->C
Choose a correct statement about R:
\begin{choices}
\choice The FD set of R is not canonical and R is in BCNF
\choice The FD set of R is canonical and R is not in BCNF
\choice The FD set of R is not canonical and R is not in BCNF
\end{choices}
\textbf{Đáp án:} C
\textbf{Giải thích:} Đáp án theo dữ liệu đề: A — A. The FD set of R is not canonical and R is in BCNF.

Nền tảng kiến thức: Chuẩn hóa dùng phụ thuộc hàm và khóa để giảm dư thừa cũng như bất thường cập nhật. Tính bao đóng thuộc tính để tìm khóa; 2NF loại phụ thuộc bộ phận vào một phần khóa ghép, 3NF xử lý phụ thuộc không khóa kiểu bắc cầu, còn BCNF yêu cầu mọi vế trái không tầm thường là siêu khóa.

Phân tích: Với phụ thuộc X → Y, mọi hai bộ có cùng giá trị X phải có cùng giá trị Y. Để kiểm tra khóa, tính X⁺ bằng cách áp dụng lặp các phụ thuộc; nếu bao đóng chứa toàn bộ thuộc tính thì X là siêu khóa, sau đó kiểm tra tính tối thiểu.

Đối chiếu phương án: B. The FD set of R is canonical and R is not in BCNF; C. The FD set of R is not canonical and R is not in BCNF. Các lựa chọn này không đáp ứng định nghĩa/điều kiện được hỏi; khi ôn, hãy chỉ ra chính xác điểm lệch của từng lựa chọn thay vì học thuộc mỗi chữ cái.
\end{question}

\begin{question}{117}
% [ID gốc: 15.6]
% [Nguồn: DBI202 Order 15 Q6]
What is the difference between the 2NF and the 3NF?
\begin{choices}
\choice 2NF deals with partial functional dependency, while 3NF deals with transitive functional dependency
\choice 2NF deals with transitive functional dependency, while 3NF deals with partial functional dependency
\end{choices}
\textbf{Đáp án:} A
\textbf{Giải thích:} Đáp án theo dữ liệu đề: A — A. 2NF deals with partial functional dependency, while 3NF deals with transitive functional dependency.

Nền tảng kiến thức: Chuẩn hóa dùng phụ thuộc hàm và khóa để giảm dư thừa cũng như bất thường cập nhật. Tính bao đóng thuộc tính để tìm khóa; 2NF loại phụ thuộc bộ phận vào một phần khóa ghép, 3NF xử lý phụ thuộc không khóa kiểu bắc cầu, còn BCNF yêu cầu mọi vế trái không tầm thường là siêu khóa.

Phân tích: Phép tập hợp trên quan hệ cần so sánh các tuple theo thuộc tính tương ứng; union lấy hợp, intersection lấy phần chung và difference lấy phần có ở vế trái mà không có ở vế phải. Trước tiên kiểm tra tính tương thích schema, sau đó áp phép theo đúng hướng/toán tử.

Đối chiếu phương án: B. 2NF deals with transitive functional dependency, while 3NF deals with partial functional depen…. Các lựa chọn này không đáp ứng định nghĩa/điều kiện được hỏi; khi ôn, hãy chỉ ra chính xác điểm lệch của từng lựa chọn thay vì học thuộc mỗi chữ cái.
\end{question}

\begin{question}{118}
% [ID gốc: 15.7]
% [Nguồn: DBI202 Order 15 Q7]
Which of the following relations is in Boyce-Codd Normal Form (BCNF)?
\begin{choices}
\choice R(ABCD) with FD's: BD -> C ; AB -> D ; AC -> B ; BD -> A
\choice R(ABCD) with FD's: BC -> A ; AD -> C ; CD -> B ; BD -> C
\choice R(ABCD) with FD's: A -> C ; B -> A ; A -> D ; AD -> C
\choice R(ABCD) with FD's: A -> D ; C -> A ; D -> B ; AC -> B
\end{choices}
\textbf{Đáp án:} A
\textbf{Giải thích:} Đáp án theo dữ liệu đề: A — A. R(ABCD) with FD's: BD -> C ; AB -> D ; AC -> B ; BD -> A.

Nền tảng kiến thức: Chuẩn hóa dùng phụ thuộc hàm và khóa để giảm dư thừa cũng như bất thường cập nhật. Tính bao đóng thuộc tính để tìm khóa; 2NF loại phụ thuộc bộ phận vào một phần khóa ghép, 3NF xử lý phụ thuộc không khóa kiểu bắc cầu, còn BCNF yêu cầu mọi vế trái không tầm thường là siêu khóa.

Phân tích: Đánh giá dạng chuẩn theo thứ tự: 1NF yêu cầu giá trị nguyên tố/không lặp theo mô hình quan hệ; 2NF loại phụ thuộc bộ phận của thuộc tính không khóa vào một phần khóa ghép; 3NF kiểm soát phụ thuộc bắc cầu theo điều kiện prime attribute; BCNF yêu cầu mọi determinant không tầm thường là siêu khóa.

Đối chiếu phương án: B. R(ABCD) with FD's: BC -> A ; AD -> C ; CD -> B ; BD -> C; C. R(ABCD) with FD's: A -> C ; B -> A ; A -> D ; AD -> C; D. R(ABCD) with FD's: A -> D ; C -> A ; D -> B ; AC -> B. Các lựa chọn này không đáp ứng định nghĩa/điều kiện được hỏi; khi ôn, hãy chỉ ra chính xác điểm lệch của từng lựa chọn thay vì học thuộc mỗi chữ cái.
\end{question}

\begin{question}{119}
% [ID gốc: 15.8]
% [Nguồn: DBI202 Order 15 Q8]
3NFconcept is related to (choose 1 answer only):
\begin{choices}
\choice Atomic definition
\choice Full dependency definition
\choice Transitive dependency definition
\choice Super Key definition
\choice All of the others
\end{choices}
\textbf{Đáp án:} C
\textbf{Giải thích:} Đáp án theo dữ liệu đề: E — E. All of the others.

Nền tảng kiến thức: Chuẩn hóa dùng phụ thuộc hàm và khóa để giảm dư thừa cũng như bất thường cập nhật. Tính bao đóng thuộc tính để tìm khóa; 2NF loại phụ thuộc bộ phận vào một phần khóa ghép, 3NF xử lý phụ thuộc không khóa kiểu bắc cầu, còn BCNF yêu cầu mọi vế trái không tầm thường là siêu khóa.

Phân tích: 3NF tập trung loại phụ thuộc bắc cầu của thuộc tính không khóa vào khóa, đồng thời có ngoại lệ cho phụ thuộc khi vế phải là thuộc tính prime. Phân biệt với 2NF, vốn nhắm tới phụ thuộc bộ phận vào một phần của khóa ghép.

Đối chiếu phương án: A. Atomic definition; B. Full dependency definition; C. Transitive dependency definition; D. Super Key definition. Các lựa chọn này không đáp ứng định nghĩa/điều kiện được hỏi; khi ôn, hãy chỉ ra chính xác điểm lệch của từng lựa chọn thay vì học thuộc mỗi chữ cái.
\end{question}

\begin{question}{120}
% [ID gốc: 15.9]
% [Nguồn: DBI202 Order 15 Q9]
The relation R(ABCD) has following FDs:
{ACD -> B ;
AC -> D ;
D -> C ;
AC -> B}
Choose the correct statement about R:
\begin{choices}
\choice R is in 3NF
\choice R is in 2NF only, not higher
\choice R is in 1NF only, not higher
\choice None of the others
\end{choices}
\textbf{Đáp án:} A
\textbf{Giải thích:} Đáp án theo dữ liệu đề: A — A. R is in 3NF.

Nền tảng kiến thức: Chuẩn hóa dùng phụ thuộc hàm và khóa để giảm dư thừa cũng như bất thường cập nhật. Tính bao đóng thuộc tính để tìm khóa; 2NF loại phụ thuộc bộ phận vào một phần khóa ghép, 3NF xử lý phụ thuộc không khóa kiểu bắc cầu, còn BCNF yêu cầu mọi vế trái không tầm thường là siêu khóa.

Phân tích: Với phụ thuộc X → Y, mọi hai bộ có cùng giá trị X phải có cùng giá trị Y. Để kiểm tra khóa, tính X⁺ bằng cách áp dụng lặp các phụ thuộc; nếu bao đóng chứa toàn bộ thuộc tính thì X là siêu khóa, sau đó kiểm tra tính tối thiểu.

Đối chiếu phương án: B. R is in 2NF only, not higher; C. R is in 1NF only, not higher; D. None of the others. Các lựa chọn này không đáp ứng định nghĩa/điều kiện được hỏi; khi ôn, hãy chỉ ra chính xác điểm lệch của từng lựa chọn thay vì học thuộc mỗi chữ cái.
\end{question}

\begin{question}{121}
% [ID gốc: 15.10]
% [Nguồn: DBI202 Order 15 Q10]
Let R(A,B,C,D) with the following FDs: {AB->C, AC->B, AD->C}
Choose a correct statement about R:
\begin{choices}
\choice R is in BCNF
\choice R is in 3NF
\choice R is in 2NF
\choice None of the others
\end{choices}
\textbf{Đáp án:} C
\textbf{Giải thích:} Đáp án theo dữ liệu đề: C — C. R is in 2NF.

Nền tảng kiến thức: Chuẩn hóa dùng phụ thuộc hàm và khóa để giảm dư thừa cũng như bất thường cập nhật. Tính bao đóng thuộc tính để tìm khóa; 2NF loại phụ thuộc bộ phận vào một phần khóa ghép, 3NF xử lý phụ thuộc không khóa kiểu bắc cầu, còn BCNF yêu cầu mọi vế trái không tầm thường là siêu khóa.

Phân tích: Với phụ thuộc X → Y, mọi hai bộ có cùng giá trị X phải có cùng giá trị Y. Để kiểm tra khóa, tính X⁺ bằng cách áp dụng lặp các phụ thuộc; nếu bao đóng chứa toàn bộ thuộc tính thì X là siêu khóa, sau đó kiểm tra tính tối thiểu.

Đối chiếu phương án: A. R is in BCNF; B. R is in 3NF; D. None of the others. Các lựa chọn này không đáp ứng định nghĩa/điều kiện được hỏi; khi ôn, hãy chỉ ra chính xác điểm lệch của từng lựa chọn thay vì học thuộc mỗi chữ cái.
\end{question}

\begin{question}{122}
% [ID gốc: 15.11]
% [Nguồn: DBI202 Order 15 Q11]
Which of the following statements is correct?
\begin{choices}
\choice All relations in 3NF are also in BCNF
\choice For any relation schema, there is a dependency-preserving decomposition into 3NF
\choice All relations in 2NF are also in BCNF
\end{choices}
\textbf{Đáp án:} B
\textbf{Giải thích:} Đáp án theo dữ liệu đề: B — B. For any relation schema, there is a dependency-preserving decomposition into 3NF.

Nền tảng kiến thức: Chuẩn hóa dùng phụ thuộc hàm và khóa để giảm dư thừa cũng như bất thường cập nhật. Tính bao đóng thuộc tính để tìm khóa; 2NF loại phụ thuộc bộ phận vào một phần khóa ghép, 3NF xử lý phụ thuộc không khóa kiểu bắc cầu, còn BCNF yêu cầu mọi vế trái không tầm thường là siêu khóa.

Phân tích: Lược đồ quan hệ ghi tên quan hệ kèm danh sách thuộc tính, thường theo dạng `Boat(BoatID, BoatName, Color, Status)`. Nếu đề hỏi schema thuần túy, không nhất thiết phải kèm kiểu dữ liệu hay các bộ dữ liệu đang lưu.

Đối chiếu phương án: A. All relations in 3NF are also in BCNF; C. All relations in 2NF are also in BCNF. Các lựa chọn này không đáp ứng định nghĩa/điều kiện được hỏi; khi ôn, hãy chỉ ra chính xác điểm lệch của từng lựa chọn thay vì học thuộc mỗi chữ cái.
\end{question}

\begin{question}{123}
% [ID gốc: 15.12]
% [Nguồn: DBI202 Order 15 Q12]
Which of the following statements is true?
\begin{choices}
\choice BCNF condition guarantees the non existence of the anomalies.
\choice In BCNF condition, the left side of every non trivial FD must be a super key.
\choice Any two-attribute relation is in BCNF.
\choice All of the others.
\end{choices}
\textbf{Đáp án:} D
\textbf{Giải thích:} Đáp án theo dữ liệu đề: D — D. All of the others..

Nền tảng kiến thức: Chuẩn hóa dùng phụ thuộc hàm và khóa để giảm dư thừa cũng như bất thường cập nhật. Tính bao đóng thuộc tính để tìm khóa; 2NF loại phụ thuộc bộ phận vào một phần khóa ghép, 3NF xử lý phụ thuộc không khóa kiểu bắc cầu, còn BCNF yêu cầu mọi vế trái không tầm thường là siêu khóa.

Phân tích: Để xác định một ràng buộc/phụ thuộc tầm thường, kiểm tra xem nó có đúng chỉ từ định nghĩa tập thuộc tính của quan hệ hay không, không cần giả định thêm về dữ liệu. Với MVD X ↠ Y, điều kiện tầm thường là Y ⊆ X hoặc X ∪ Y bằng toàn bộ thuộc tính; cần phân biệt ký hiệu ↠ với phụ thuộc hàm →.

Đối chiếu phương án: A. BCNF condition guarantees the non existence of the anomalies.; B. In BCNF condition, the left side of every non trivial FD must be a super key.; C. Any two-attribute relation is in BCNF.. Các lựa chọn này không đáp ứng định nghĩa/điều kiện được hỏi; khi ôn, hãy chỉ ra chính xác điểm lệch của từng lựa chọn thay vì học thuộc mỗi chữ cái.
\end{question}

\begin{question}{124}
% [ID gốc: 15.13]
% [Nguồn: DBI202 Order 15 Q13]
How to eliminate anomalies when we design a database?
\begin{choices}
\choice We should decompose relation to eliminate anomalies
\choice We should join relations to eliminate anomalies
\choice We should union relations to eliminate anomalies
\choice None of the others
\end{choices}
\textbf{Đáp án:} A
\textbf{Giải thích:} Đáp án theo dữ liệu đề: A — A. We should decompose relation to eliminate anomalies.

Nền tảng kiến thức: Chuẩn hóa dùng phụ thuộc hàm và khóa để giảm dư thừa cũng như bất thường cập nhật. Tính bao đóng thuộc tính để tìm khóa; 2NF loại phụ thuộc bộ phận vào một phần khóa ghép, 3NF xử lý phụ thuộc không khóa kiểu bắc cầu, còn BCNF yêu cầu mọi vế trái không tầm thường là siêu khóa.

Phân tích: Phép tập hợp trên quan hệ cần so sánh các tuple theo thuộc tính tương ứng; union lấy hợp, intersection lấy phần chung và difference lấy phần có ở vế trái mà không có ở vế phải. Trước tiên kiểm tra tính tương thích schema, sau đó áp phép theo đúng hướng/toán tử.

Đối chiếu phương án: B. We should join relations to eliminate anomalies; C. We should union relations to eliminate anomalies; D. None of the others. Các lựa chọn này không đáp ứng định nghĩa/điều kiện được hỏi; khi ôn, hãy chỉ ra chính xác điểm lệch của từng lựa chọn thay vì học thuộc mỗi chữ cái.
\end{question}

\begin{question}{125}
% [ID gốc: 15.14]
% [Nguồn: DBI202 Order 15 Q14]
What is the goal of decomposition when designing a database?
\begin{choices}
\choice The goal of decomposition is to replace a relation by several relations that do not exhibit anomalies
\choice The goal of decomposition is to increase the speed of the database
\choice The goal of decomposition is to increase the security of the database
\choice None of the others
\end{choices}
\textbf{Đáp án:} A
\textbf{Giải thích:} Đáp án theo dữ liệu đề: A — A. The goal of decomposition is to replace a relation by several relations that do not exhibit anomalies.

Nền tảng kiến thức: Chuẩn hóa dùng phụ thuộc hàm và khóa để giảm dư thừa cũng như bất thường cập nhật. Tính bao đóng thuộc tính để tìm khóa; 2NF loại phụ thuộc bộ phận vào một phần khóa ghép, 3NF xử lý phụ thuộc không khóa kiểu bắc cầu, còn BCNF yêu cầu mọi vế trái không tầm thường là siêu khóa.

Phân tích: Decomposition tách một quan hệ thành các quan hệ nhỏ hơn để giảm dư thừa và các bất thường chèn/sửa/xóa. Một phép tách tốt còn cần lossless join (ghép lại không sinh tuple giả); tùy bài toán còn xét bảo toàn phụ thuộc.

Đối chiếu phương án: B. The goal of decomposition is to increase the speed of the database; C. The goal of decomposition is to increase the security of the database; D. None of the others. Các lựa chọn này không đáp ứng định nghĩa/điều kiện được hỏi; khi ôn, hãy chỉ ra chính xác điểm lệch của từng lựa chọn thay vì học thuộc mỗi chữ cái.
\end{question}

\begin{question}{126}
% [ID gốc: 15.15]
% [Nguồn: DBI202 Order 15 Q15]
Given a relation R(A.B.C.D). Which of the followings is trivial?
\begin{choices}
\choice A->BCD
\choice A->->BCD
\choice A->AB
\choice A->->AB
\end{choices}
\textbf{Đáp án:} B
\textbf{Giải thích:} Đáp án theo dữ liệu đề: B — B. A->->BCD.

Nền tảng kiến thức: Chuẩn hóa dùng phụ thuộc hàm và khóa để giảm dư thừa cũng như bất thường cập nhật. Tính bao đóng thuộc tính để tìm khóa; 2NF loại phụ thuộc bộ phận vào một phần khóa ghép, 3NF xử lý phụ thuộc không khóa kiểu bắc cầu, còn BCNF yêu cầu mọi vế trái không tầm thường là siêu khóa.

Phân tích: Để xác định một ràng buộc/phụ thuộc tầm thường, kiểm tra xem nó có đúng chỉ từ định nghĩa tập thuộc tính của quan hệ hay không, không cần giả định thêm về dữ liệu. Với MVD X ↠ Y, điều kiện tầm thường là Y ⊆ X hoặc X ∪ Y bằng toàn bộ thuộc tính; cần phân biệt ký hiệu ↠ với phụ thuộc hàm →.

Đối chiếu phương án: A. A->BCD; C. A->AB; D. A->->AB. Các lựa chọn này không đáp ứng định nghĩa/điều kiện được hỏi; khi ôn, hãy chỉ ra chính xác điểm lệch của từng lựa chọn thay vì học thuộc mỗi chữ cái.
\end{question}

\begin{question}{127}
% [ID gốc: 15.16]
% [Nguồn: DBI202 Order 15 Q16]
Give the relation R(XYZT) with the following FD's:
XY -> Z ; XYT -> Z ; XYZ -> T ; XZ -> T
Determine the primary key for R and then choose the correct statement
\begin{choices}
\choice R is in BCNF
\choice R is not in BCNF
\end{choices}
\textbf{Đáp án:} A
\textbf{Giải thích:} Đáp án theo dữ liệu đề: B — B. R is not in BCNF.

Nền tảng kiến thức: Chuẩn hóa dùng phụ thuộc hàm và khóa để giảm dư thừa cũng như bất thường cập nhật. Tính bao đóng thuộc tính để tìm khóa; 2NF loại phụ thuộc bộ phận vào một phần khóa ghép, 3NF xử lý phụ thuộc không khóa kiểu bắc cầu, còn BCNF yêu cầu mọi vế trái không tầm thường là siêu khóa.

Phân tích: Kiểm tra khóa bằng hai tiêu chí: tập thuộc tính phải xác định duy nhất mọi bộ (tính siêu khóa), và nếu là candidate key thì phải tối thiểu, tức bỏ bất kỳ thuộc tính nào cũng mất tính duy nhất. PRIMARY KEY chọn một candidate key và không nhận NULL.

Đối chiếu phương án: A. R is in BCNF. Các lựa chọn này không đáp ứng định nghĩa/điều kiện được hỏi; khi ôn, hãy chỉ ra chính xác điểm lệch của từng lựa chọn thay vì học thuộc mỗi chữ cái.
\end{question}

\begin{question}{128}
% [ID gốc: 15.17]
% [Nguồn: DBI202 Order 15 Q17]
Given the relation schema R(XYZT) and functional dependencies
F = {X->Z, T->Y}.
Which functional dependency causes a violation of second normal form (2NF) ?
\begin{choices}
\choice X->Z
\choice T->Y
\choice both X->Z and T->Y
\choice none of the above
\end{choices}
\textbf{Đáp án:} C
\textbf{Giải thích:} Đáp án theo dữ liệu đề: C — C. both X->Z and T->Y.

Nền tảng kiến thức: Chuẩn hóa dùng phụ thuộc hàm và khóa để giảm dư thừa cũng như bất thường cập nhật. Tính bao đóng thuộc tính để tìm khóa; 2NF loại phụ thuộc bộ phận vào một phần khóa ghép, 3NF xử lý phụ thuộc không khóa kiểu bắc cầu, còn BCNF yêu cầu mọi vế trái không tầm thường là siêu khóa.

Phân tích: Lược đồ quan hệ ghi tên quan hệ kèm danh sách thuộc tính, thường theo dạng `Boat(BoatID, BoatName, Color, Status)`. Nếu đề hỏi schema thuần túy, không nhất thiết phải kèm kiểu dữ liệu hay các bộ dữ liệu đang lưu.

Đối chiếu phương án: A. X->Z; B. T->Y; D. none of the above. Các lựa chọn này không đáp ứng định nghĩa/điều kiện được hỏi; khi ôn, hãy chỉ ra chính xác điểm lệch của từng lựa chọn thay vì học thuộc mỗi chữ cái.
\end{question}

\begin{question}{129}
% [ID gốc: 15.18]
% [Nguồn: DBI202 Order 15 Q18]
What are some benefits of normalization when converting an ERD to a relational model?
\begin{choices}
\choice Increased storage space efficiency
\choice Reduced update anomalies
\choice Simplified database queries
\choice Enhanced data security
\end{choices}
\textbf{Đáp án:} B
\textbf{Giải thích:} Đáp án theo dữ liệu đề: B — B. Reduced update anomalies.

Nền tảng kiến thức: Chuẩn hóa dùng phụ thuộc hàm và khóa để giảm dư thừa cũng như bất thường cập nhật. Tính bao đóng thuộc tính để tìm khóa; 2NF loại phụ thuộc bộ phận vào một phần khóa ghép, 3NF xử lý phụ thuộc không khóa kiểu bắc cầu, còn BCNF yêu cầu mọi vế trái không tầm thường là siêu khóa.

Phân tích: Khi đọc mô hình ER, tách ba thứ: entity set là loại đối tượng, relationship set là liên kết giữa các loại đối tượng, còn cardinality/participation là ràng buộc số lượng hoặc mức tham gia. Đáp án đúng phải gọi đúng thành phần mà câu hỏi mô tả.

Đối chiếu phương án: A. Increased storage space efficiency; C. Simplified database queries; D. Enhanced data security. Các lựa chọn này không đáp ứng định nghĩa/điều kiện được hỏi; khi ôn, hãy chỉ ra chính xác điểm lệch của từng lựa chọn thay vì học thuộc mỗi chữ cái.
\end{question}

\begin{question}{130}
% [ID gốc: 15.19]
% [Nguồn: DBI202 Order 15 Q19]
What is the purpose of data normalization in the context of the database design?
\begin{choices}
\choice Ensure data security
\choice Avoid information anomalies
\choice Make sure data is inherited
\choice Ensuring better data storage
\end{choices}
\textbf{Đáp án:} B
\textbf{Giải thích:} Đáp án theo dữ liệu đề: B — B. Avoid information anomalies.

Nền tảng kiến thức: Chuẩn hóa dùng phụ thuộc hàm và khóa để giảm dư thừa cũng như bất thường cập nhật. Tính bao đóng thuộc tính để tìm khóa; 2NF loại phụ thuộc bộ phận vào một phần khóa ghép, 3NF xử lý phụ thuộc không khóa kiểu bắc cầu, còn BCNF yêu cầu mọi vế trái không tầm thường là siêu khóa.

Phân tích: Phân biệt tệp dữ liệu chứa trang/bảng với tệp transaction log ghi lại thay đổi phục vụ commit và recovery. Một trang/khối là đơn vị I/O cơ bản; lựa chọn tổ chức vật lý ảnh hưởng số lần đọc ghi và cách tìm bản ghi.

Đối chiếu phương án: A. Ensure data security; C. Make sure data is inherited; D. Ensuring better data storage. Các lựa chọn này không đáp ứng định nghĩa/điều kiện được hỏi; khi ôn, hãy chỉ ra chính xác điểm lệch của từng lựa chọn thay vì học thuộc mỗi chữ cái.
\end{question}

\begin{question}{131}
% [ID gốc: 15.20]
% [Nguồn: DBI202 Order 15 Q20]
Given following schema R(A, B, C, D) and S={A→B, BC, C→D}Compute {A}+?
\begin{choices}
\choice AB
\choice ABC
\choice ABCD
\choice BCD
\end{choices}
\textbf{Đáp án:} C
\textbf{Giải thích:} Đáp án theo dữ liệu đề: C — C. ABCD.

Nền tảng kiến thức: Chuẩn hóa dùng phụ thuộc hàm và khóa để giảm dư thừa cũng như bất thường cập nhật. Tính bao đóng thuộc tính để tìm khóa; 2NF loại phụ thuộc bộ phận vào một phần khóa ghép, 3NF xử lý phụ thuộc không khóa kiểu bắc cầu, còn BCNF yêu cầu mọi vế trái không tầm thường là siêu khóa.

Phân tích: Tính bao đóng bằng cách bắt đầu A⁺ = {A}, rồi áp dụng lần lượt mọi phụ thuộc có toàn bộ vế trái đã nằm trong bao đóng. Ký hiệu phụ thuộc ở nguồn có vẻ bị thiếu mũi tên sau `BC`; nếu ý nguồn là A → B, B → C, C → D thì suy ra A⁺ = {A,B,C,D} theo chuỗi bắc cầu. Cần đối chiếu bản đề gốc để xác nhận phần ký hiệu bị thiếu trước khi chốt đáp án.

Đối chiếu phương án: A. AB; B. ABC; D. BCD. Các lựa chọn này không đáp ứng định nghĩa/điều kiện được hỏi; khi ôn, hãy chỉ ra chính xác điểm lệch của từng lựa chọn thay vì học thuộc mỗi chữ cái.
\end{question}

\begin{question}{132}
% [ID gốc: 15.21]
% [Nguồn: DBI202 Order 15 Q21]
Suppose that we have a relation that possesses data about an individual entity. In Addition, it should not have any multi-valued dependencies.
Which of the following is the highest normal form for this relation?
\begin{choices}
\choice 1NF
\choice 2NF
\choice 3NF
\choice 4NF
\end{choices}
\textbf{Đáp án:} D
\textbf{Giải thích:} Đáp án theo dữ liệu đề: D — D. 4NF.

Nền tảng kiến thức: Chuẩn hóa dùng phụ thuộc hàm và khóa để giảm dư thừa cũng như bất thường cập nhật. Tính bao đóng thuộc tính để tìm khóa; 2NF loại phụ thuộc bộ phận vào một phần khóa ghép, 3NF xử lý phụ thuộc không khóa kiểu bắc cầu, còn BCNF yêu cầu mọi vế trái không tầm thường là siêu khóa.

Phân tích: Đánh giá dạng chuẩn theo thứ tự: 1NF yêu cầu giá trị nguyên tố/không lặp theo mô hình quan hệ; 2NF loại phụ thuộc bộ phận của thuộc tính không khóa vào một phần khóa ghép; 3NF kiểm soát phụ thuộc bắc cầu theo điều kiện prime attribute; BCNF yêu cầu mọi determinant không tầm thường là siêu khóa.

Đối chiếu phương án: A. 1NF; B. 2NF; C. 3NF. Các lựa chọn này không đáp ứng định nghĩa/điều kiện được hỏi; khi ôn, hãy chỉ ra chính xác điểm lệch của từng lựa chọn thay vì học thuộc mỗi chữ cái.
\end{question}

\begin{question}{133}
% [ID gốc: 15.22]
% [Nguồn: DBI202 Order 15 Q22]
Which of the following statements is true about data normalization:
\begin{choices}
\choice 1NF: Is a relationship with a primary key and no repeating groups
\choice 2NF: Is 1NF without transitive functional dependencies
\choice 3NF is: 1 NF without partial functional dependencies
\choice All of the above answers are wrong.
\end{choices}
\textbf{Đáp án:} A
\textbf{Giải thích:} Đáp án theo dữ liệu đề: A — A. 1NF: Is a relationship with a primary key and no repeating groups.

Nền tảng kiến thức: Chuẩn hóa dùng phụ thuộc hàm và khóa để giảm dư thừa cũng như bất thường cập nhật. Tính bao đóng thuộc tính để tìm khóa; 2NF loại phụ thuộc bộ phận vào một phần khóa ghép, 3NF xử lý phụ thuộc không khóa kiểu bắc cầu, còn BCNF yêu cầu mọi vế trái không tầm thường là siêu khóa.

Phân tích: Với lựa chọn tổng hợp, không chọn theo cảm tính: kiểm tra từng phát biểu độc lập. “All of the others/above” chỉ đúng khi mọi phát biểu còn lại đúng; “None” chỉ đúng khi không phát biểu nào trong nhóm thỏa yêu cầu. Đáp án lưu cần được đọc cùng các mệnh đề thành phần.

Đối chiếu phương án: B. 2NF: Is 1NF without transitive functional dependencies; C. 3NF is: 1 NF without partial functional dependencies; D. All of the above answers are wrong.. Các lựa chọn này không đáp ứng định nghĩa/điều kiện được hỏi; khi ôn, hãy chỉ ra chính xác điểm lệch của từng lựa chọn thay vì học thuộc mỗi chữ cái.
\end{question}

\begin{question}{134}
% [ID gốc: 15.23]
% [Nguồn: DBI202 Order 15 Q23]
Which normalization level ensures that each attribute is dependent only on the primary key?
\begin{choices}
\choice First Normal Form (1NF)
\choice Second Normal Form (2NF)
\choice Third Normal Form (3NF)
\choice Boyce-Codd Normal Form (BCNF)
\end{choices}
\textbf{Đáp án:} C
\textbf{Giải thích:} Đáp án theo dữ liệu đề: C — C. Third Normal Form (3NF).

Nền tảng kiến thức: Chuẩn hóa dùng phụ thuộc hàm và khóa để giảm dư thừa cũng như bất thường cập nhật. Tính bao đóng thuộc tính để tìm khóa; 2NF loại phụ thuộc bộ phận vào một phần khóa ghép, 3NF xử lý phụ thuộc không khóa kiểu bắc cầu, còn BCNF yêu cầu mọi vế trái không tầm thường là siêu khóa.

Phân tích: Kiểm tra khóa bằng hai tiêu chí: tập thuộc tính phải xác định duy nhất mọi bộ (tính siêu khóa), và nếu là candidate key thì phải tối thiểu, tức bỏ bất kỳ thuộc tính nào cũng mất tính duy nhất. PRIMARY KEY chọn một candidate key và không nhận NULL.

Đối chiếu phương án: A. First Normal Form (1NF); B. Second Normal Form (2NF); D. Boyce-Codd Normal Form (BCNF). Các lựa chọn này không đáp ứng định nghĩa/điều kiện được hỏi; khi ôn, hãy chỉ ra chính xác điểm lệch của từng lựa chọn thay vì học thuộc mỗi chữ cái.
\end{question}

\begin{question}{135}
% [ID gốc: 15.24]
% [Nguồn: DBI202 Order 15 Q24]
Suppose that there is a relation R (A, B, C, D, E, G)
And functional dependencies S = {AB→C, DEG, CA, BEC, BC→D, CG → BD, ACD → B, CE → AG}
Compute {AB}+
\begin{choices}
\choice ABCD
\choice ABCDE
\choice ABCDG
\choice ABCDEG
\end{choices}
\textbf{Đáp án:} D
\textbf{Giải thích:} Đáp án theo dữ liệu đề: D — D. ABCDEG.

Nền tảng kiến thức: Chuẩn hóa dùng phụ thuộc hàm và khóa để giảm dư thừa cũng như bất thường cập nhật. Tính bao đóng thuộc tính để tìm khóa; 2NF loại phụ thuộc bộ phận vào một phần khóa ghép, 3NF xử lý phụ thuộc không khóa kiểu bắc cầu, còn BCNF yêu cầu mọi vế trái không tầm thường là siêu khóa.

Phân tích: Với phụ thuộc X → Y, mọi hai bộ có cùng giá trị X phải có cùng giá trị Y. Để kiểm tra khóa, tính X⁺ bằng cách áp dụng lặp các phụ thuộc; nếu bao đóng chứa toàn bộ thuộc tính thì X là siêu khóa, sau đó kiểm tra tính tối thiểu.

Đối chiếu phương án: A. ABCD; B. ABCDE; C. ABCDG. Các lựa chọn này không đáp ứng định nghĩa/điều kiện được hỏi; khi ôn, hãy chỉ ra chính xác điểm lệch của từng lựa chọn thay vì học thuộc mỗi chữ cái.
\end{question}

\begin{question}{136}
% [ID gốc: 15.25]
% [Nguồn: DBI202 Order 15 Q25]
Which of the following represents Armstrong's Axiom of Transitivity?
\begin{choices}
\choice If AB, then AC → BC
\choice If A B and B → C, then A → C
\choice If A B, then B → A
\choice If A B, then A → BC
\end{choices}
\textbf{Đáp án:} B
\textbf{Giải thích:} Đáp án theo dữ liệu đề: B — B. If A B and B → C, then A → C.

Nền tảng kiến thức: Chuẩn hóa dùng phụ thuộc hàm và khóa để giảm dư thừa cũng như bất thường cập nhật. Tính bao đóng thuộc tính để tìm khóa; 2NF loại phụ thuộc bộ phận vào một phần khóa ghép, 3NF xử lý phụ thuộc không khóa kiểu bắc cầu, còn BCNF yêu cầu mọi vế trái không tầm thường là siêu khóa.

Phân tích: Tiên đề bắc cầu của Armstrong: nếu X → Y và Y → Z thì suy ra X → Z. Phải kiểm tra hai phụ thuộc nối tiếp cùng vế giữa Y; các tiên đề phản xạ và tăng trưởng mô tả quy tắc khác.

Đối chiếu phương án: A. If AB, then AC → BC; C. If A B, then B → A; D. If A B, then A → BC. Các lựa chọn này không đáp ứng định nghĩa/điều kiện được hỏi; khi ôn, hãy chỉ ra chính xác điểm lệch của từng lựa chọn thay vì học thuộc mỗi chữ cái.
\end{question}

\begin{question}{137}
% [ID gốc: 15.26]
% [Nguồn: DBI202 Order 15 Q26]
Which statement is the best represents Armstrong's Axiom of Decomposition?
\begin{choices}
\choice If A B and A → C, then A → BC
\choice If A B and B → C, then A → C
\choice If A B, then AC → BC
\choice If A B, then B → C
\end{choices}
\textbf{Đáp án:} D
\textbf{Giải thích:} Đáp án theo dữ liệu đề: C — C. If A B, then AC → BC.

Nền tảng kiến thức: Chuẩn hóa dùng phụ thuộc hàm và khóa để giảm dư thừa cũng như bất thường cập nhật. Tính bao đóng thuộc tính để tìm khóa; 2NF loại phụ thuộc bộ phận vào một phần khóa ghép, 3NF xử lý phụ thuộc không khóa kiểu bắc cầu, còn BCNF yêu cầu mọi vế trái không tầm thường là siêu khóa.

Phân tích: Tiên đề bắc cầu của Armstrong: nếu X → Y và Y → Z thì suy ra X → Z. Phải kiểm tra hai phụ thuộc nối tiếp cùng vế giữa Y; các tiên đề phản xạ và tăng trưởng mô tả quy tắc khác.

Đối chiếu phương án: A. If A B and A → C, then A → BC; B. If A B and B → C, then A → C; D. If A B, then B → C. Các lựa chọn này không đáp ứng định nghĩa/điều kiện được hỏi; khi ôn, hãy chỉ ra chính xác điểm lệch của từng lựa chọn thay vì học thuộc mỗi chữ cái.
\end{question}

\begin{question}{138}
% [ID gốc: 15.27]
% [Nguồn: DBI202 Order 15 Q27]
What are logical constraints?
\begin{choices}
\choice The relationships between attributes are represented by relational operations.
\choice The relationships between attributes are represented by comparisons operator.
\choice The relationships between attributes are represented by mathematical expressions.
\choice The relationships between attributes are represented by functional dependencies.
\end{choices}
\textbf{Đáp án:} D
\textbf{Giải thích:} Đáp án theo dữ liệu đề: D — D. The relationships between attributes are represented by functional dependencies..

Nền tảng kiến thức: Chuẩn hóa dùng phụ thuộc hàm và khóa để giảm dư thừa cũng như bất thường cập nhật. Tính bao đóng thuộc tính để tìm khóa; 2NF loại phụ thuộc bộ phận vào một phần khóa ghép, 3NF xử lý phụ thuộc không khóa kiểu bắc cầu, còn BCNF yêu cầu mọi vế trái không tầm thường là siêu khóa.

Phân tích: Với phụ thuộc X → Y, mọi hai bộ có cùng giá trị X phải có cùng giá trị Y. Để kiểm tra khóa, tính X⁺ bằng cách áp dụng lặp các phụ thuộc; nếu bao đóng chứa toàn bộ thuộc tính thì X là siêu khóa, sau đó kiểm tra tính tối thiểu.

Đối chiếu phương án: A. The relationships between attributes are represented by relational operations.; B. The relationships between attributes are represented by comparisons operator.; C. The relationships between attributes are represented by mathematical expressions.. Các lựa chọn này không đáp ứng định nghĩa/điều kiện được hỏi; khi ôn, hãy chỉ ra chính xác điểm lệch của từng lựa chọn thay vì học thuộc mỗi chữ cái.
\end{question}

\begin{question}{139}
% [ID gốc: 15.28]
% [Nguồn: DBI202 Order 15 Q28]
If attribute A determines both attributes B and C, then it is also true that:
\begin{choices}
\choice (B,C) -> A.
\choice C -> A.
\choice B -> A.
\choice A -> B
\end{choices}
\textbf{Đáp án:} D
\textbf{Giải thích:} Đáp án theo dữ liệu đề: D — D. A -> B.

Nền tảng kiến thức: Chuẩn hóa dùng phụ thuộc hàm và khóa để giảm dư thừa cũng như bất thường cập nhật. Tính bao đóng thuộc tính để tìm khóa; 2NF loại phụ thuộc bộ phận vào một phần khóa ghép, 3NF xử lý phụ thuộc không khóa kiểu bắc cầu, còn BCNF yêu cầu mọi vế trái không tầm thường là siêu khóa.

Phân tích: Nếu A xác định đồng thời B và C thì theo quy tắc phân rã của phụ thuộc hàm, A → B và A → C đều suy ra; theo quy tắc hợp, có thể viết A → BC. Đây là hệ quả trực tiếp của ý nghĩa “cùng A thì cùng giá trị của mọi thuộc tính ở vế phải”.

Đối chiếu phương án: A. (B,C) -> A.; B. C -> A.; C. B -> A.. Các lựa chọn này không đáp ứng định nghĩa/điều kiện được hỏi; khi ôn, hãy chỉ ra chính xác điểm lệch của từng lựa chọn thay vì học thuộc mỗi chữ cái.
\end{question}

\begin{question}{140}
% [ID gốc: 15.29]
% [Nguồn: DBI202 Order 15 Q29]
Which operation is commonly used to display all the properties derived from the original property?
\begin{choices}
\choice Union
\choice Intersection
\choice Closure
\choice Projection
\end{choices}
\textbf{Đáp án:} C
\textbf{Giải thích:} Đáp án theo dữ liệu đề: C — C. Closure.

Nền tảng kiến thức: Chuẩn hóa dùng phụ thuộc hàm và khóa để giảm dư thừa cũng như bất thường cập nhật. Tính bao đóng thuộc tính để tìm khóa; 2NF loại phụ thuộc bộ phận vào một phần khóa ghép, 3NF xử lý phụ thuộc không khóa kiểu bắc cầu, còn BCNF yêu cầu mọi vế trái không tầm thường là siêu khóa.

Phân tích: Phép tập hợp trên quan hệ cần so sánh các tuple theo thuộc tính tương ứng; union lấy hợp, intersection lấy phần chung và difference lấy phần có ở vế trái mà không có ở vế phải. Trước tiên kiểm tra tính tương thích schema, sau đó áp phép theo đúng hướng/toán tử.

Đối chiếu phương án: A. Union; B. Intersection; D. Projection. Các lựa chọn này không đáp ứng định nghĩa/điều kiện được hỏi; khi ôn, hãy chỉ ra chính xác điểm lệch của từng lựa chọn thay vì học thuộc mỗi chữ cái.
\end{question}

\begin{question}{141}
% [ID gốc: 15.30]
% [Nguồn: DBI202 Order 15 Q30]
If a relation is in BCNF, then it is also in
\begin{choices}
\choice 1 NF
\choice 2 NF
\choice 3 NF
\choice All of the mentioned
\end{choices}
\textbf{Đáp án:} D
\textbf{Giải thích:} Đáp án theo dữ liệu đề: D — D. All of the mentioned.

Nền tảng kiến thức: Chuẩn hóa dùng phụ thuộc hàm và khóa để giảm dư thừa cũng như bất thường cập nhật. Tính bao đóng thuộc tính để tìm khóa; 2NF loại phụ thuộc bộ phận vào một phần khóa ghép, 3NF xử lý phụ thuộc không khóa kiểu bắc cầu, còn BCNF yêu cầu mọi vế trái không tầm thường là siêu khóa.

Phân tích: Đánh giá dạng chuẩn theo thứ tự: 1NF yêu cầu giá trị nguyên tố/không lặp theo mô hình quan hệ; 2NF loại phụ thuộc bộ phận của thuộc tính không khóa vào một phần khóa ghép; 3NF kiểm soát phụ thuộc bắc cầu theo điều kiện prime attribute; BCNF yêu cầu mọi determinant không tầm thường là siêu khóa.

Đối chiếu phương án: A. 1 NF; B. 2 NF; C. 3 NF. Các lựa chọn này không đáp ứng định nghĩa/điều kiện được hỏi; khi ôn, hãy chỉ ra chính xác điểm lệch của từng lựa chọn thay vì học thuộc mỗi chữ cái.
\end{question}

\begin{question}{142}
% [ID gốc: 15.31]
% [Nguồn: DBI202 Order 15 Q31]
_______has a right side that is a subset of its left side
\begin{choices}
\choice A trivial FD
\choice A non_trivial FD
\choice A key of relation
\choice A super key of relation
\end{choices}
\textbf{Đáp án:} A
\textbf{Giải thích:} Đáp án theo dữ liệu đề: A — A. A trivial FD.

Nền tảng kiến thức: Chuẩn hóa dùng phụ thuộc hàm và khóa để giảm dư thừa cũng như bất thường cập nhật. Tính bao đóng thuộc tính để tìm khóa; 2NF loại phụ thuộc bộ phận vào một phần khóa ghép, 3NF xử lý phụ thuộc không khóa kiểu bắc cầu, còn BCNF yêu cầu mọi vế trái không tầm thường là siêu khóa.

Phân tích: Khóa là thuộc tính hoặc tập thuộc tính dùng để nhận diện duy nhất một bộ/hàng. Khóa có thể gồm một hay nhiều cột, nên không bắt buộc luôn là cột đơn hoặc luôn gồm nhiều cột; nó cũng không “nhận diện một cột”.

Đối chiếu phương án: B. A non_trivial FD; C. A key of relation; D. A super key of relation. Các lựa chọn này không đáp ứng định nghĩa/điều kiện được hỏi; khi ôn, hãy chỉ ra chính xác điểm lệch của từng lựa chọn thay vì học thuộc mỗi chữ cái.
\end{question}

\begin{question}{143}
% [ID gốc: 15.32]
% [Nguồn: DBI202 Order 15 Q32]
Given R as a relation, which of the following statements is correct about normalizing data in 3NF form:
\begin{choices}
\choice If R: Has a primary key, no repeating groups, no multivalued attributes/composite attributes
\choice If R: Has a primary key, no repeating groups, no multivalued attributes/complex attributes, no partial functional dependencies.
\choice If R: Has a primary key, no repeating groups, no multivalued/composite attributes, no transitive dependencies.
\choice If R: Has a primary key, no repeating groups, no multivalued attributes/complex attributes, no transitive functional dependencies, no partial functional dependencies.
\choice If R: has complete functional dependency.
\end{choices}
\textbf{Đáp án:} D
\textbf{Giải thích:} Đáp án theo dữ liệu đề: D — D. If R: Has a primary key, no repeating groups, no multivalued attributes/complex attributes, no transitive functional dependencies, no partial functional dependencies..

Nền tảng kiến thức: Chuẩn hóa dùng phụ thuộc hàm và khóa để giảm dư thừa cũng như bất thường cập nhật. Tính bao đóng thuộc tính để tìm khóa; 2NF loại phụ thuộc bộ phận vào một phần khóa ghép, 3NF xử lý phụ thuộc không khóa kiểu bắc cầu, còn BCNF yêu cầu mọi vế trái không tầm thường là siêu khóa.

Phân tích: Đánh giá dạng chuẩn theo thứ tự: 1NF yêu cầu giá trị nguyên tố/không lặp theo mô hình quan hệ; 2NF loại phụ thuộc bộ phận của thuộc tính không khóa vào một phần khóa ghép; 3NF kiểm soát phụ thuộc bắc cầu theo điều kiện prime attribute; BCNF yêu cầu mọi determinant không tầm thường là siêu khóa.

Đối chiếu phương án: A. If R: Has a primary key, no repeating groups, no multivalued attributes/composite attributes; B. If R: Has a primary key, no repeating groups, no multivalued attributes/complex attributes, no…; C. If R: Has a primary key, no repeating groups, no multivalued/composite attributes, no transiti…; E. If R: has complete functional dependency.. Các lựa chọn này không đáp ứng định nghĩa/điều kiện được hỏi; khi ôn, hãy chỉ ra chính xác điểm lệch của từng lựa chọn thay vì học thuộc mỗi chữ cái.
\end{question}

\begin{question}{144}
% [ID gốc: 15.33]
% [Nguồn: DBI202 Order 15 Q33]
Which of the following statements is true regarding Third Normal Form (3NF)?
\begin{choices}
\choice It eliminates transitive dependencies
\choice It allows for redundant data
\choice It only considers partial dependencies
\choice It doesn't address any anomalies in a database
\end{choices}
\textbf{Đáp án:} A
\textbf{Giải thích:} Đáp án theo dữ liệu đề: A — A. It eliminates transitive dependencies.

Nền tảng kiến thức: Chuẩn hóa dùng phụ thuộc hàm và khóa để giảm dư thừa cũng như bất thường cập nhật. Tính bao đóng thuộc tính để tìm khóa; 2NF loại phụ thuộc bộ phận vào một phần khóa ghép, 3NF xử lý phụ thuộc không khóa kiểu bắc cầu, còn BCNF yêu cầu mọi vế trái không tầm thường là siêu khóa.

Phân tích: Đánh giá dạng chuẩn theo thứ tự: 1NF yêu cầu giá trị nguyên tố/không lặp theo mô hình quan hệ; 2NF loại phụ thuộc bộ phận của thuộc tính không khóa vào một phần khóa ghép; 3NF kiểm soát phụ thuộc bắc cầu theo điều kiện prime attribute; BCNF yêu cầu mọi determinant không tầm thường là siêu khóa.

Đối chiếu phương án: B. It allows for redundant data; C. It only considers partial dependencies; D. It doesn't address any anomalies in a database. Các lựa chọn này không đáp ứng định nghĩa/điều kiện được hỏi; khi ôn, hãy chỉ ra chính xác điểm lệch của từng lựa chọn thay vì học thuộc mỗi chữ cái.
\end{question}

\begin{question}{145}
% [ID gốc: 15.35]
% [Nguồn: DBI202 Order 15 Q35]
What is Database Normalization?
\begin{choices}
\choice A process where by a limit is put on the number of fields a record can contain
\choice The process of ensuring that each table has a key
\choice The process of ensuring that a relational database has at least two tables in it
\choice A process whereby the design of a table (relation) is decomposed into more tables that more precisely fit the relational mode
\end{choices}
\textbf{Đáp án:} D
\textbf{Giải thích:} Đáp án theo dữ liệu đề: D — D. A process whereby the design of a table (relation) is decomposed into more tables that more precisely fit the relational mode.

Nền tảng kiến thức: Chuẩn hóa dùng phụ thuộc hàm và khóa để giảm dư thừa cũng như bất thường cập nhật. Tính bao đóng thuộc tính để tìm khóa; 2NF loại phụ thuộc bộ phận vào một phần khóa ghép, 3NF xử lý phụ thuộc không khóa kiểu bắc cầu, còn BCNF yêu cầu mọi vế trái không tầm thường là siêu khóa.

Phân tích: Normalization tổ chức/ phân rã quan hệ dựa trên phụ thuộc và khóa để giảm dư thừa, tránh bất thường khi thêm/sửa/xóa. Không phải mục tiêu đơn thuần là tăng tốc hay tăng bảo mật; sau khi phân rã còn cần kiểm tra khả năng ghép nối không mất mát.

Đối chiếu phương án: A. A process where by a limit is put on the number of fields a record can contain; B. The process of ensuring that each table has a key; C. The process of ensuring that a relational database has at least two tables in it. Các lựa chọn này không đáp ứng định nghĩa/điều kiện được hỏi; khi ôn, hãy chỉ ra chính xác điểm lệch của từng lựa chọn thay vì học thuộc mỗi chữ cái.
\end{question}

\begin{question}{146}
% [ID gốc: 15.36]
% [Nguồn: DBI202 Order 15 Q36]
Consider a relation R(A,B,C,D,E) with functional dependencies:
AB->C, B->C, E->D and C->E.What is/are the key(s) for R
\begin{choices}
\choice AC
\choice AD
\choice AB
\choice AE
\end{choices}
\textbf{Đáp án:} C
\textbf{Giải thích:} Đáp án theo dữ liệu đề: C — C. AB.

Nền tảng kiến thức: Chuẩn hóa dùng phụ thuộc hàm và khóa để giảm dư thừa cũng như bất thường cập nhật. Tính bao đóng thuộc tính để tìm khóa; 2NF loại phụ thuộc bộ phận vào một phần khóa ghép, 3NF xử lý phụ thuộc không khóa kiểu bắc cầu, còn BCNF yêu cầu mọi vế trái không tầm thường là siêu khóa.

Phân tích: Với phụ thuộc X → Y, mọi hai bộ có cùng giá trị X phải có cùng giá trị Y. Để kiểm tra khóa, tính X⁺ bằng cách áp dụng lặp các phụ thuộc; nếu bao đóng chứa toàn bộ thuộc tính thì X là siêu khóa, sau đó kiểm tra tính tối thiểu.

Đối chiếu phương án: A. AC; B. AD; D. AE. Các lựa chọn này không đáp ứng định nghĩa/điều kiện được hỏi; khi ôn, hãy chỉ ra chính xác điểm lệch của từng lựa chọn thay vì học thuộc mỗi chữ cái.
\end{question}

\begin{question}{147}
% [ID gốc: 15.37]
% [Nguồn: DBI202 Order 15 Q37]
Suppose that we have a relation R (B, C,
\begin{choices}
\choice that satisfies FD's:D->B and D->
\choice Which of the following FD is equivalent to the given dependency functions?
\choice D->BC
\choice D->DB
\choice DB->B
\choice DC->DB
\end{choices}
\textbf{Đáp án:} C
\textbf{Giải thích:} Đáp án theo dữ liệu đề: A — A. D->BC.

Nền tảng kiến thức: Chuẩn hóa dùng phụ thuộc hàm và khóa để giảm dư thừa cũng như bất thường cập nhật. Tính bao đóng thuộc tính để tìm khóa; 2NF loại phụ thuộc bộ phận vào một phần khóa ghép, 3NF xử lý phụ thuộc không khóa kiểu bắc cầu, còn BCNF yêu cầu mọi vế trái không tầm thường là siêu khóa.

Phân tích: Với phụ thuộc X → Y, mọi hai bộ có cùng giá trị X phải có cùng giá trị Y. Để kiểm tra khóa, tính X⁺ bằng cách áp dụng lặp các phụ thuộc; nếu bao đóng chứa toàn bộ thuộc tính thì X là siêu khóa, sau đó kiểm tra tính tối thiểu.

Đối chiếu phương án: D. that satisfies FD's:D->B and D->; C. Which of the following FD is equivalent to the given dependency functions?; B. D->DB; C. DB->B; D. DC->DB. Các lựa chọn này không đáp ứng định nghĩa/điều kiện được hỏi; khi ôn, hãy chỉ ra chính xác điểm lệch của từng lựa chọn thay vì học thuộc mỗi chữ cái.
\end{question}

\begin{question}{148}
% [ID gốc: 15.38]
% [Nguồn: DBI202 Order 15 Q38]
Which normalization level ensures that there are no repeating groups in tables?
\begin{choices}
\choice First Normal Form (1NF)
\choice Second Normal Form (2NF)
\choice Third Normal Form (3NF)
\choice Boyce-Codd Normal Form (BCNF)
\end{choices}
\textbf{Đáp án:} A
\textbf{Giải thích:} Đáp án theo dữ liệu đề: A — A. First Normal Form (1NF).

Nền tảng kiến thức: Chuẩn hóa dùng phụ thuộc hàm và khóa để giảm dư thừa cũng như bất thường cập nhật. Tính bao đóng thuộc tính để tìm khóa; 2NF loại phụ thuộc bộ phận vào một phần khóa ghép, 3NF xử lý phụ thuộc không khóa kiểu bắc cầu, còn BCNF yêu cầu mọi vế trái không tầm thường là siêu khóa.

Phân tích: Đánh giá dạng chuẩn theo thứ tự: 1NF yêu cầu giá trị nguyên tố/không lặp theo mô hình quan hệ; 2NF loại phụ thuộc bộ phận của thuộc tính không khóa vào một phần khóa ghép; 3NF kiểm soát phụ thuộc bắc cầu theo điều kiện prime attribute; BCNF yêu cầu mọi determinant không tầm thường là siêu khóa.

Đối chiếu phương án: B. Second Normal Form (2NF); C. Third Normal Form (3NF); D. Boyce-Codd Normal Form (BCNF). Các lựa chọn này không đáp ứng định nghĩa/điều kiện được hỏi; khi ôn, hãy chỉ ra chính xác điểm lệch của từng lựa chọn thay vì học thuộc mỗi chữ cái.
\end{question}

\begin{question}{149}
% [ID gốc: 15.39]
% [Nguồn: DBI202 Order 15 Q39]
What is the primary goal of normalization in database design?
\begin{choices}
\choice To increase data redundancy
\choice To simplify data retrieval
\choice To introduce duplicate data
\choice To violate functional dependencies
\end{choices}
\textbf{Đáp án:} B
\textbf{Giải thích:} Đáp án theo dữ liệu đề: B — B. To simplify data retrieval.

Nền tảng kiến thức: Chuẩn hóa dùng phụ thuộc hàm và khóa để giảm dư thừa cũng như bất thường cập nhật. Tính bao đóng thuộc tính để tìm khóa; 2NF loại phụ thuộc bộ phận vào một phần khóa ghép, 3NF xử lý phụ thuộc không khóa kiểu bắc cầu, còn BCNF yêu cầu mọi vế trái không tầm thường là siêu khóa.

Phân tích: Với phụ thuộc X → Y, mọi hai bộ có cùng giá trị X phải có cùng giá trị Y. Để kiểm tra khóa, tính X⁺ bằng cách áp dụng lặp các phụ thuộc; nếu bao đóng chứa toàn bộ thuộc tính thì X là siêu khóa, sau đó kiểm tra tính tối thiểu.

Đối chiếu phương án: A. To increase data redundancy; C. To introduce duplicate data; D. To violate functional dependencies. Các lựa chọn này không đáp ứng định nghĩa/điều kiện được hỏi; khi ôn, hãy chỉ ra chính xác điểm lệch của từng lựa chọn thay vì học thuộc mỗi chữ cái.
\end{question}

\begin{question}{150}
% [ID gốc: 15.40]
% [Nguồn: DBI202 Order 15 Q40]
Consider the relation:
OrderDetail (orderNumber, productCode, productName, quantity, dateOrdered, salesDiscount)
And FDs{ orderNumber, productCode->quantity,
orderNumber->dateOrdered,
productCode->productName, salesDiscount}
And the key of OrderDetail is {orderNumber, productCode}
Which of the following is full dependency on the key?
\begin{choices}
\choice orderNumber, productCode->quantity
\choice orderNumber->dateOrdered
\choice productCode->productName
\choice productCode->productName, salesDiscount
\end{choices}
\textbf{Đáp án:} A
\textbf{Giải thích:} Đáp án theo dữ liệu đề: A — A. orderNumber, productCode->quantity.

Nền tảng kiến thức: Chuẩn hóa dùng phụ thuộc hàm và khóa để giảm dư thừa cũng như bất thường cập nhật. Tính bao đóng thuộc tính để tìm khóa; 2NF loại phụ thuộc bộ phận vào một phần khóa ghép, 3NF xử lý phụ thuộc không khóa kiểu bắc cầu, còn BCNF yêu cầu mọi vế trái không tầm thường là siêu khóa.

Phân tích: Với phụ thuộc X → Y, mọi hai bộ có cùng giá trị X phải có cùng giá trị Y. Để kiểm tra khóa, tính X⁺ bằng cách áp dụng lặp các phụ thuộc; nếu bao đóng chứa toàn bộ thuộc tính thì X là siêu khóa, sau đó kiểm tra tính tối thiểu.

Đối chiếu phương án: B. orderNumber->dateOrdered; C. productCode->productName; D. productCode->productName, salesDiscount. Các lựa chọn này không đáp ứng định nghĩa/điều kiện được hỏi; khi ôn, hãy chỉ ra chính xác điểm lệch của từng lựa chọn thay vì học thuộc mỗi chữ cái.
\end{question}

\begin{question}{151}
% [ID gốc: 15.41]
% [Nguồn: DBI202 Order 15 Q41]
Problems such as redundancy that occur when we try to cram too much into a single relation are called _____
\begin{choices}
\choice Normal Form
\choice Anomalies
\choice Decomposing Relations
\choice None of the others
\end{choices}
\textbf{Đáp án:} B
\textbf{Giải thích:} Đáp án theo dữ liệu đề: B — B. Anomalies.

Nền tảng kiến thức: Chuẩn hóa dùng phụ thuộc hàm và khóa để giảm dư thừa cũng như bất thường cập nhật. Tính bao đóng thuộc tính để tìm khóa; 2NF loại phụ thuộc bộ phận vào một phần khóa ghép, 3NF xử lý phụ thuộc không khóa kiểu bắc cầu, còn BCNF yêu cầu mọi vế trái không tầm thường là siêu khóa.

Phân tích: Với lựa chọn tổng hợp, không chọn theo cảm tính: kiểm tra từng phát biểu độc lập. “All of the others/above” chỉ đúng khi mọi phát biểu còn lại đúng; “None” chỉ đúng khi không phát biểu nào trong nhóm thỏa yêu cầu. Đáp án lưu cần được đọc cùng các mệnh đề thành phần.

Đối chiếu phương án: A. Normal Form; C. Decomposing Relations; D. None of the others. Các lựa chọn này không đáp ứng định nghĩa/điều kiện được hỏi; khi ôn, hãy chỉ ra chính xác điểm lệch của từng lựa chọn thay vì học thuộc mỗi chữ cái.
\end{question}

\begin{question}{152}
% [ID gốc: 15.42]
% [Nguồn: DBI202 Order 15 Q42]
The highest normal form for relation schema R(ABCD) with functional dependencies: F = {AB-> C; D->B, C->ABD} is:
\begin{choices}
\choice 2NF
\choice 1NF
\choice 3NF
\choice BCNF
\end{choices}
\textbf{Đáp án:} C
\textbf{Giải thích:} Đáp án theo dữ liệu đề: C — C. 3NF.

Lời giải gốc: Quan hệ đạt 3NF nhưng không đạt BCNF.

Nền tảng kiến thức: Chuẩn hóa dùng phụ thuộc hàm và khóa để giảm dư thừa cũng như bất thường cập nhật. Tính bao đóng thuộc tính để tìm khóa; 2NF loại phụ thuộc bộ phận vào một phần khóa ghép, 3NF xử lý phụ thuộc không khóa kiểu bắc cầu, còn BCNF yêu cầu mọi vế trái không tầm thường là siêu khóa.

Phân tích: Lược đồ quan hệ ghi tên quan hệ kèm danh sách thuộc tính, thường theo dạng `Boat(BoatID, BoatName, Color, Status)`. Nếu đề hỏi schema thuần túy, không nhất thiết phải kèm kiểu dữ liệu hay các bộ dữ liệu đang lưu.

Đối chiếu phương án: A. 2NF; B. 1NF; D. BCNF. Các lựa chọn này không đáp ứng định nghĩa/điều kiện được hỏi; khi ôn, hãy chỉ ra chính xác điểm lệch của từng lựa chọn thay vì học thuộc mỗi chữ cái.
\end{question}

\begin{question}{153}
% [ID gốc: 15.43]
% [Nguồn: DBI202 Order 15 Q43]
What is Boyce-Codd Normal Form (BCNF)?
\begin{choices}
\choice The weakest normal form.
\choice A normal form where every determinant is a candidate key.
\choice A normal form that allows redundancy.
\choice The same as first normal form.
\end{choices}
\textbf{Đáp án:} B
\textbf{Giải thích:} Đáp án theo dữ liệu đề: B — B. A normal form where every determinant is a candidate key..

Lời giải gốc: Dạng chuẩn Boyce-Codd (BCNF) là phiên bản nghiêm ngặt hơn của 3NF, trong đó với mọi phụ thuộc hàm không tầm thường, vế trái (determinant) bắt buộc phải là một siêu khóa hoặc khóa ứng viên.

Nền tảng kiến thức: Chuẩn hóa dùng phụ thuộc hàm và khóa để giảm dư thừa cũng như bất thường cập nhật. Tính bao đóng thuộc tính để tìm khóa; 2NF loại phụ thuộc bộ phận vào một phần khóa ghép, 3NF xử lý phụ thuộc không khóa kiểu bắc cầu, còn BCNF yêu cầu mọi vế trái không tầm thường là siêu khóa.

Phân tích: Đánh giá dạng chuẩn theo thứ tự: 1NF yêu cầu giá trị nguyên tố/không lặp theo mô hình quan hệ; 2NF loại phụ thuộc bộ phận của thuộc tính không khóa vào một phần khóa ghép; 3NF kiểm soát phụ thuộc bắc cầu theo điều kiện prime attribute; BCNF yêu cầu mọi determinant không tầm thường là siêu khóa.

Đối chiếu phương án: A. The weakest normal form.; C. A normal form that allows redundancy.; D. The same as first normal form.. Các lựa chọn này không đáp ứng định nghĩa/điều kiện được hỏi; khi ôn, hãy chỉ ra chính xác điểm lệch của từng lựa chọn thay vì học thuộc mỗi chữ cái.
\end{question}

\begin{question}{154}
% [ID gốc: 15.44]
% [Nguồn: DBI202 Order 15 Q44]
Suppose we have a relation R(ID, full_name, address) with functional dependencies {ID → full_name, ID → address}. What is the highest normal form of R in relational database design?
\begin{choices}
\choice 2NF
\choice 1NF
\choice 3NF
\choice BCNF
\end{choices}
\textbf{Đáp án:} D
\textbf{Giải thích:} Đáp án theo dữ liệu đề: D — D. BCNF.

Lời giải gốc: Với khóa chính là ID, mọi phụ thuộc hàm đều có vế trái là siêu khóa (ID → full_name, ID → address), do đó quan hệ đạt dạng chuẩn cao nhất là Boyce-Codd Normal Form (BCNF).

Nền tảng kiến thức: Chuẩn hóa dùng phụ thuộc hàm và khóa để giảm dư thừa cũng như bất thường cập nhật. Tính bao đóng thuộc tính để tìm khóa; 2NF loại phụ thuộc bộ phận vào một phần khóa ghép, 3NF xử lý phụ thuộc không khóa kiểu bắc cầu, còn BCNF yêu cầu mọi vế trái không tầm thường là siêu khóa.

Phân tích: Đánh giá dạng chuẩn theo thứ tự: 1NF yêu cầu giá trị nguyên tố/không lặp theo mô hình quan hệ; 2NF loại phụ thuộc bộ phận của thuộc tính không khóa vào một phần khóa ghép; 3NF kiểm soát phụ thuộc bắc cầu theo điều kiện prime attribute; BCNF yêu cầu mọi determinant không tầm thường là siêu khóa.

Đối chiếu phương án: A. 2NF; B. 1NF; C. 3NF. Các lựa chọn này không đáp ứng định nghĩa/điều kiện được hỏi; khi ôn, hãy chỉ ra chính xác điểm lệch của từng lựa chọn thay vì học thuộc mỗi chữ cái.
\end{question}

\begin{question}{155}
% [ID gốc: 15.45]
% [Nguồn: DBI202 Order 15 Q45]
What is the purpose of data normalization in the context of database design?
\begin{choices}
\choice Ensure data security
\choice Avoid information anomalies
\choice Make sure data is inherited
\choice Ensuring better data storage
\end{choices}
\textbf{Đáp án:} B
\textbf{Giải thích:} Đáp án theo dữ liệu đề: B — B. Avoid information anomalies.

Lời giải gốc: Chuẩn hóa dữ liệu (Data Normalization) nhằm mục đích giảm thiểu sự trùng lặp dữ liệu và tránh các hiện tượng bất thường (anomalies) như Insertion Anomaly, Deletion Anomaly, và Update Anomaly.

Nền tảng kiến thức: Chuẩn hóa dùng phụ thuộc hàm và khóa để giảm dư thừa cũng như bất thường cập nhật. Tính bao đóng thuộc tính để tìm khóa; 2NF loại phụ thuộc bộ phận vào một phần khóa ghép, 3NF xử lý phụ thuộc không khóa kiểu bắc cầu, còn BCNF yêu cầu mọi vế trái không tầm thường là siêu khóa.

Phân tích: Phân biệt tệp dữ liệu chứa trang/bảng với tệp transaction log ghi lại thay đổi phục vụ commit và recovery. Một trang/khối là đơn vị I/O cơ bản; lựa chọn tổ chức vật lý ảnh hưởng số lần đọc ghi và cách tìm bản ghi.

Đối chiếu phương án: A. Ensure data security; C. Make sure data is inherited; D. Ensuring better data storage. Các lựa chọn này không đáp ứng định nghĩa/điều kiện được hỏi; khi ôn, hãy chỉ ra chính xác điểm lệch của từng lựa chọn thay vì học thuộc mỗi chữ cái.
\end{question}

\begin{question}{156}
% [ID gốc: 15.46]
% [Nguồn: DBI202 Order 15 Q46]
Consider a relation R(A, B, C, D) and the following rules. Rule 1: If A → C then C → A. Rule 2: If B → C then BC → C. Rule 3: If AB → C then A → C and B → C. Which of the following rules is NOT valid?
\begin{choices}
\choice Only Rule 1
\choice None of the given rules are valid
\choice Rule 2, Rule 3
\choice Rule 1, Rule 3
\end{choices}
\textbf{Đáp án:} D
\textbf{Giải thích:} Đáp án theo dữ liệu đề: D — D. Rule 1, Rule 3.

Lời giải gốc: Rule 1 (đối xứng không phải là luật hợp lệ của phụ thuộc hàm) và Rule 3 (tách vế trái không đảm bảo tính đúng đắn) đều không hợp lệ. Rule 2 hợp lệ theo luật phản xạ/gia tăng.

Nền tảng kiến thức: Chuẩn hóa dùng phụ thuộc hàm và khóa để giảm dư thừa cũng như bất thường cập nhật. Tính bao đóng thuộc tính để tìm khóa; 2NF loại phụ thuộc bộ phận vào một phần khóa ghép, 3NF xử lý phụ thuộc không khóa kiểu bắc cầu, còn BCNF yêu cầu mọi vế trái không tầm thường là siêu khóa.

Phân tích: Rule 1 không hợp lệ vì phụ thuộc A → C không cho phép đảo chiều thành C → A. Rule 2 hợp lệ do C nằm trong BC nên BC → C đúng theo phản xạ; Rule 3 không hợp lệ vì không thể tự ý bỏ A hoặc B khỏi vế trái AB → C. Do đó các rule sai là 1 và 3.

Đối chiếu phương án: A. Only Rule 1; B. None of the given rules are valid; C. Rule 2, Rule 3. Các lựa chọn này không đáp ứng định nghĩa/điều kiện được hỏi; khi ôn, hãy chỉ ra chính xác điểm lệch của từng lựa chọn thay vì học thuộc mỗi chữ cái.
\end{question}

\begin{question}{157}
% [ID gốc: 15.47]
% [Nguồn: DBI202 Order 15 Q47]
Suppose we have a relation R(A, B, C, D, E, F, N) with functional dependencies: F → A, A → B, B → C, BC → A, A → D, E → A, D → E. Which of the following is NOT a key?
\begin{choices}
\choice AB
\choice FN
\choice E
\choice D
\end{choices}
\textbf{Đáp án:} A
\textbf{Giải thích:} Đáp án theo dữ liệu đề: A — A. AB.

Lời giải gốc: Tập thuộc tính {A, B} không chứa N và không thể suy ra N từ tập phụ thuộc hàm, do đó không thể là khóa của quan hệ.

Nền tảng kiến thức: Chuẩn hóa dùng phụ thuộc hàm và khóa để giảm dư thừa cũng như bất thường cập nhật. Tính bao đóng thuộc tính để tìm khóa; 2NF loại phụ thuộc bộ phận vào một phần khóa ghép, 3NF xử lý phụ thuộc không khóa kiểu bắc cầu, còn BCNF yêu cầu mọi vế trái không tầm thường là siêu khóa.

Phân tích: Khóa là thuộc tính hoặc tập thuộc tính dùng để nhận diện duy nhất một bộ/hàng. Khóa có thể gồm một hay nhiều cột, nên không bắt buộc luôn là cột đơn hoặc luôn gồm nhiều cột; nó cũng không “nhận diện một cột”.

Đối chiếu phương án: B. FN; C. E; D. D. Các lựa chọn này không đáp ứng định nghĩa/điều kiện được hỏi; khi ôn, hãy chỉ ra chính xác điểm lệch của từng lựa chọn thay vì học thuộc mỗi chữ cái.
\end{question}

\begin{question}{158}
% [ID gốc: 15.50]
% [Nguồn: DBI202 Order 15 Q50]
What is the main requirement of the First Normal Form (1NF) in database design?
\begin{choices}
\choice Each table contains only one foreign key.
\choice Each attribute contains only atomic values, not repeating groups or arrays of values.
\choice Each table must have at least one primary key.
\choice Non-key attributes must be completely dependent on the primary key.
\end{choices}
\textbf{Đáp án:} B
\textbf{Giải thích:} Đáp án theo dữ liệu đề: B — B. Each attribute contains only atomic values, not repeating groups or arrays of values..

Lời giải gốc: Dạng chuẩn 1 (1NF) yêu cầu mỗi thuộc tính trong bảng phải chứa các giá trị nguyên tử (atomic), không chứa các nhóm lặp hoặc mảng giá trị. ---

Nền tảng kiến thức: Chuẩn hóa dùng phụ thuộc hàm và khóa để giảm dư thừa cũng như bất thường cập nhật. Tính bao đóng thuộc tính để tìm khóa; 2NF loại phụ thuộc bộ phận vào một phần khóa ghép, 3NF xử lý phụ thuộc không khóa kiểu bắc cầu, còn BCNF yêu cầu mọi vế trái không tầm thường là siêu khóa.

Phân tích: Đánh giá dạng chuẩn theo thứ tự: 1NF yêu cầu giá trị nguyên tố/không lặp theo mô hình quan hệ; 2NF loại phụ thuộc bộ phận của thuộc tính không khóa vào một phần khóa ghép; 3NF kiểm soát phụ thuộc bắc cầu theo điều kiện prime attribute; BCNF yêu cầu mọi determinant không tầm thường là siêu khóa.

Đối chiếu phương án: A. Each table contains only one foreign key.; C. Each table must have at least one primary key.; D. Non-key attributes must be completely dependent on the primary key.. Các lựa chọn này không đáp ứng định nghĩa/điều kiện được hỏi; khi ôn, hãy chỉ ra chính xác điểm lệch của từng lựa chọn thay vì học thuộc mỗi chữ cái.
\end{question}

\begin{question}{159}
% [ID gốc: 15.51]
% [Nguồn: DBI202 Order 15 Q51]
Consider the relation R(A,B,C,D) and functional dependencies: A→B, C→D. Find key of the above relation.
\begin{choices}
\choice The key is the A column
\choice The key includes two columns (A,C)
\choice The key is the A column or C column
\choice The key is the C column
\end{choices}
\textbf{Đáp án:} B
\textbf{Giải thích:} Đáp án theo dữ liệu đề: B — B. The key includes two columns (A,C).

Lời giải gốc: Không có thuộc tính nào xác định được A và C. Do đó, A và C bắt buộc phải là một phần của khóa. Bao đóng của (A,C)+ = {A, B, C, D}, vừa đủ xác định toàn bộ bảng R. Khóa chính là (A,C).

Nền tảng kiến thức: Chuẩn hóa dùng phụ thuộc hàm và khóa để giảm dư thừa cũng như bất thường cập nhật. Tính bao đóng thuộc tính để tìm khóa; 2NF loại phụ thuộc bộ phận vào một phần khóa ghép, 3NF xử lý phụ thuộc không khóa kiểu bắc cầu, còn BCNF yêu cầu mọi vế trái không tầm thường là siêu khóa.

Phân tích: Với phụ thuộc X → Y, mọi hai bộ có cùng giá trị X phải có cùng giá trị Y. Để kiểm tra khóa, tính X⁺ bằng cách áp dụng lặp các phụ thuộc; nếu bao đóng chứa toàn bộ thuộc tính thì X là siêu khóa, sau đó kiểm tra tính tối thiểu.

Đối chiếu phương án: A. The key is the A column; C. The key is the A column or C column; D. The key is the C column. Các lựa chọn này không đáp ứng định nghĩa/điều kiện được hỏi; khi ôn, hãy chỉ ra chính xác điểm lệch của từng lựa chọn thay vì học thuộc mỗi chữ cái.
\end{question}

\begin{question}{160}
% [ID gốc: 15.54]
% [Nguồn: DBI202 Order 15 Q54]
Consider a user table with columns ID, NAME, Password, and email. Which functional dependency exists in this table?
\begin{choices}
\choice ID → NAME
\choice Password → email
\choice Password → NAME
\choice NAME → email
\end{choices}
\textbf{Đáp án:} A
\textbf{Giải thích:} Đáp án theo dữ liệu đề: A — A. ID → NAME.

Lời giải gốc: Trong cấu trúc quản lý dữ liệu thông thường, mã định danh ID xác định duy nhất tên người dùng NAME.

Nền tảng kiến thức: Chuẩn hóa dùng phụ thuộc hàm và khóa để giảm dư thừa cũng như bất thường cập nhật. Tính bao đóng thuộc tính để tìm khóa; 2NF loại phụ thuộc bộ phận vào một phần khóa ghép, 3NF xử lý phụ thuộc không khóa kiểu bắc cầu, còn BCNF yêu cầu mọi vế trái không tầm thường là siêu khóa.

Phân tích: Với phụ thuộc X → Y, mọi hai bộ có cùng giá trị X phải có cùng giá trị Y. Để kiểm tra khóa, tính X⁺ bằng cách áp dụng lặp các phụ thuộc; nếu bao đóng chứa toàn bộ thuộc tính thì X là siêu khóa, sau đó kiểm tra tính tối thiểu.

Đối chiếu phương án: B. Password → email; C. Password → NAME; D. NAME → email. Các lựa chọn này không đáp ứng định nghĩa/điều kiện được hỏi; khi ôn, hãy chỉ ra chính xác điểm lệch của từng lựa chọn thay vì học thuộc mỗi chữ cái.
\end{question}

\begin{question}{161}
% [ID gốc: 15.55]
% [Nguồn: DBI202 Order 15 Q55]
Consider a relation R(A,B,C,D) with functional dependencies: A → C, AB → A, C → CB, AD → ABC Which of the following is trivial functional dependency of the relation?
\begin{choices}
\choice AB → A
\choice AD → ABC
\choice C → CB
\choice AB → A, C→ CB
\end{choices}
\textbf{Đáp án:} A
\textbf{Giải thích:} Đáp án theo dữ liệu đề: A — A. AB → A.

Lời giải gốc: Phụ thuộc hàm $X \rightarrow Y$ được gọi là hiển nhiên (Trivial) nếu tập thuộc tính $Y$ là tập con của tập thuộc tính $X$ ($Y \subseteq X$). Ở đây $A \subseteq AB$, do đó AB $\rightarrow$ A là phụ thuộc hàm hiển nhiên.

Nền tảng kiến thức: Chuẩn hóa dùng phụ thuộc hàm và khóa để giảm dư thừa cũng như bất thường cập nhật. Tính bao đóng thuộc tính để tìm khóa; 2NF loại phụ thuộc bộ phận vào một phần khóa ghép, 3NF xử lý phụ thuộc không khóa kiểu bắc cầu, còn BCNF yêu cầu mọi vế trái không tầm thường là siêu khóa.

Phân tích: Để xác định một ràng buộc/phụ thuộc tầm thường, kiểm tra xem nó có đúng chỉ từ định nghĩa tập thuộc tính của quan hệ hay không, không cần giả định thêm về dữ liệu. Với MVD X ↠ Y, điều kiện tầm thường là Y ⊆ X hoặc X ∪ Y bằng toàn bộ thuộc tính; cần phân biệt ký hiệu ↠ với phụ thuộc hàm →.

Đối chiếu phương án: B. AD → ABC; C. C → CB; D. AB → A, C→ CB. Các lựa chọn này không đáp ứng định nghĩa/điều kiện được hỏi; khi ôn, hãy chỉ ra chính xác điểm lệch của từng lựa chọn thay vì học thuộc mỗi chữ cái.
\end{question}

\begin{question}{162}
% [ID gốc: 15.56]
% [Nguồn: DBI202 Order 15 Q56]
Which of the following best describes an anomaly that normalization helps to eliminate in relational database design?
\begin{choices}
\choice Inconsistent data types within a table
\choice Redundant data that leads to storage inefficiency
\choice Insertion, update, or deletion anomalies
\choice Lack of referential integrity between tables
\end{choices}
\textbf{Đáp án:} C
\textbf{Giải thích:} Đáp án theo dữ liệu đề: C — C. Insertion, update, or deletion anomalies.

Lời giải gốc: Chuẩn hóa dữ liệu giúp loại bỏ trực tiếp các bất thường (anomalies) xảy ra khi thêm, sửa hoặc xóa dữ liệu trong bảng.

Nền tảng kiến thức: Chuẩn hóa dùng phụ thuộc hàm và khóa để giảm dư thừa cũng như bất thường cập nhật. Tính bao đóng thuộc tính để tìm khóa; 2NF loại phụ thuộc bộ phận vào một phần khóa ghép, 3NF xử lý phụ thuộc không khóa kiểu bắc cầu, còn BCNF yêu cầu mọi vế trái không tầm thường là siêu khóa.

Phân tích: Khóa ngoại buộc giá trị tham chiếu phải khớp một khóa được tham chiếu ở quan hệ đích (hoặc NULL nếu ràng buộc cho phép). Đây là toàn vẹn tham chiếu; nó khác tính duy nhất của primary/candidate key.

Đối chiếu phương án: A. Inconsistent data types within a table; B. Redundant data that leads to storage inefficiency; D. Lack of referential integrity between tables. Các lựa chọn này không đáp ứng định nghĩa/điều kiện được hỏi; khi ôn, hãy chỉ ra chính xác điểm lệch của từng lựa chọn thay vì học thuộc mỗi chữ cái.
\end{question}

\begin{question}{163}
% [ID gốc: 15.57]
% [Nguồn: DBI202 Order 15 Q57]
Which of the following is the main requirement of 1NF (First Normal Form)?
\begin{choices}
\choice No column in the table should contain repeating values or arrays.
\choice Each table contains only one foreign key.
\choice Each table has at least one primary key.
\choice All non-key attributes must be fully dependent on the primary key.
\end{choices}
\textbf{Đáp án:} A
\textbf{Giải thích:} Đáp án theo dữ liệu đề: A — A. No column in the table should contain repeating values or arrays..

Lời giải gốc: Điều kiện cốt lõi của Dạng chuẩn 1 (1NF) là mọi thuộc tính phải chứa các giá trị nguyên tố (nguyên tử - atomic), nghĩa là một ô dữ liệu không được chứa danh sách, mảng hoặc các nhóm lặp lại của nhiều giá trị.

Nền tảng kiến thức: Chuẩn hóa dùng phụ thuộc hàm và khóa để giảm dư thừa cũng như bất thường cập nhật. Tính bao đóng thuộc tính để tìm khóa; 2NF loại phụ thuộc bộ phận vào một phần khóa ghép, 3NF xử lý phụ thuộc không khóa kiểu bắc cầu, còn BCNF yêu cầu mọi vế trái không tầm thường là siêu khóa.

Phân tích: Đánh giá dạng chuẩn theo thứ tự: 1NF yêu cầu giá trị nguyên tố/không lặp theo mô hình quan hệ; 2NF loại phụ thuộc bộ phận của thuộc tính không khóa vào một phần khóa ghép; 3NF kiểm soát phụ thuộc bắc cầu theo điều kiện prime attribute; BCNF yêu cầu mọi determinant không tầm thường là siêu khóa.

Đối chiếu phương án: B. Each table contains only one foreign key.; C. Each table has at least one primary key.; D. All non-key attributes must be fully dependent on the primary key.. Các lựa chọn này không đáp ứng định nghĩa/điều kiện được hỏi; khi ôn, hãy chỉ ra chính xác điểm lệch của từng lựa chọn thay vì học thuộc mỗi chữ cái.
\end{question}

\begin{question}{164}
% [ID gốc: 15.59]
% [Nguồn: DBI202 Order 15 Q59]
Consider the table: Which functional dependency exists in the above table?
\begin{choices}
\choice ID NAME
\choice Password → email
\choice Password →NAME
\choice NAME email
\end{choices}
\textbf{Đáp án:} A
\textbf{Giải thích:} Đáp án theo dữ liệu đề: A — A. ID NAME.

Lời giải gốc: Theo quy tắc thiết kế cơ sở dữ liệu tiêu chuẩn, khóa chính (ID) sẽ xác định duy nhất các thuộc tính còn lại. Do đó, phụ thuộc hàm ID -> NAME là hoàn toàn hợp lý.

Nền tảng kiến thức: Chuẩn hóa dùng phụ thuộc hàm và khóa để giảm dư thừa cũng như bất thường cập nhật. Tính bao đóng thuộc tính để tìm khóa; 2NF loại phụ thuộc bộ phận vào một phần khóa ghép, 3NF xử lý phụ thuộc không khóa kiểu bắc cầu, còn BCNF yêu cầu mọi vế trái không tầm thường là siêu khóa.

Phân tích: Với phụ thuộc X → Y, mọi hai bộ có cùng giá trị X phải có cùng giá trị Y. Để kiểm tra khóa, tính X⁺ bằng cách áp dụng lặp các phụ thuộc; nếu bao đóng chứa toàn bộ thuộc tính thì X là siêu khóa, sau đó kiểm tra tính tối thiểu.

Đối chiếu phương án: B. Password → email; C. Password →NAME; D. NAME email. Các lựa chọn này không đáp ứng định nghĩa/điều kiện được hỏi; khi ôn, hãy chỉ ra chính xác điểm lệch của từng lựa chọn thay vì học thuộc mỗi chữ cái.
\end{question}

\begin{question}{165}
% [ID gốc: 15.62]
% [Nguồn: DBI202 Order 15 Q62]
Consider a relation R(A,B,C,D,E) with functional dependencies: AB-&gt;C, B-&gt;C, E-&gt;D and C-&gt;E. What is/are the key(s) for R
\begin{choices}
\choice AC
\choice AD
\choice AB
\choice AE
\end{choices}
\textbf{Đáp án:} C
\textbf{Giải thích:} Đáp án theo dữ liệu đề: C — C. AB.

Lời giải gốc: Nhìn vào các phụ thuộc hàm, ta thấy thuộc tính A và B không bao giờ xuất hiện ở vế phải (không bị suy ra bởi ai). Vậy bất kỳ khóa nào cũng phải chứa ít nhất A và B. Ta tính bao đóng của tập (A,B): {A,B}+ = {A, B, C, E, D}. Bao đóng này bao phủ toàn bộ các thuộc tính, nên AB chính là khóa.

Nền tảng kiến thức: Chuẩn hóa dùng phụ thuộc hàm và khóa để giảm dư thừa cũng như bất thường cập nhật. Tính bao đóng thuộc tính để tìm khóa; 2NF loại phụ thuộc bộ phận vào một phần khóa ghép, 3NF xử lý phụ thuộc không khóa kiểu bắc cầu, còn BCNF yêu cầu mọi vế trái không tầm thường là siêu khóa.

Phân tích: Với phụ thuộc X → Y, mọi hai bộ có cùng giá trị X phải có cùng giá trị Y. Để kiểm tra khóa, tính X⁺ bằng cách áp dụng lặp các phụ thuộc; nếu bao đóng chứa toàn bộ thuộc tính thì X là siêu khóa, sau đó kiểm tra tính tối thiểu.

Đối chiếu phương án: A. AC; B. AD; D. AE. Các lựa chọn này không đáp ứng định nghĩa/điều kiện được hỏi; khi ôn, hãy chỉ ra chính xác điểm lệch của từng lựa chọn thay vì học thuộc mỗi chữ cái.
\end{question}

\begin{question}{166}
% [ID gốc: 15.64]
% [Nguồn: DBI202 Order 15 Q64]
Consider a relation with schema R(A, B, C, D) and FD's (A → B, B → C, B → D, D → A). Which of the following is the {A}+ ?
\begin{choices}
\choice {A}
\choice {A, B, C, D}
\choice {A, B, C}
\choice {A, B}
\end{choices}
\textbf{Đáp án:} B
\textbf{Giải thích:} Đáp án theo dữ liệu đề: B — B. {A, B, C, D}.

Lời giải gốc: Tìm bao đóng của A: Đầu tiên {A}+ = {A}. Từ A → B, thêm B thành {A, B}. Từ B → C và B → D, thêm C, D thành {A, B, C, D}. Vậy {A}+ = {A, B, C, D}.

Nền tảng kiến thức: Chuẩn hóa dùng phụ thuộc hàm và khóa để giảm dư thừa cũng như bất thường cập nhật. Tính bao đóng thuộc tính để tìm khóa; 2NF loại phụ thuộc bộ phận vào một phần khóa ghép, 3NF xử lý phụ thuộc không khóa kiểu bắc cầu, còn BCNF yêu cầu mọi vế trái không tầm thường là siêu khóa.

Phân tích: Với phụ thuộc X → Y, mọi hai bộ có cùng giá trị X phải có cùng giá trị Y. Để kiểm tra khóa, tính X⁺ bằng cách áp dụng lặp các phụ thuộc; nếu bao đóng chứa toàn bộ thuộc tính thì X là siêu khóa, sau đó kiểm tra tính tối thiểu.

Đối chiếu phương án: A. {A}; C. {A, B, C}; D. {A, B}. Các lựa chọn này không đáp ứng định nghĩa/điều kiện được hỏi; khi ôn, hãy chỉ ra chính xác điểm lệch của từng lựa chọn thay vì học thuộc mỗi chữ cái.
\end{question}

\begin{question}{167}
% [ID gốc: 15.65]
% [Nguồn: DBI202 Order 15 Q65]
In the relational model, what does "functional dependency" imply?
\begin{choices}
\choice One attribute uniquely determines another attribute
\choice Each relation depends on another relation
\choice Relations can be modified without affecting each other
\choice All attributes must be independent of each other
\end{choices}
\textbf{Đáp án:} A
\textbf{Giải thích:} Đáp án theo dữ liệu đề: A — A. One attribute uniquely determines another attribute.

Lời giải gốc: Phụ thuộc hàm biểu thị mối quan hệ mà một tập thuộc tính này xác định duy nhất giá trị của một tập thuộc tính khác.

Nền tảng kiến thức: Chuẩn hóa dùng phụ thuộc hàm và khóa để giảm dư thừa cũng như bất thường cập nhật. Tính bao đóng thuộc tính để tìm khóa; 2NF loại phụ thuộc bộ phận vào một phần khóa ghép, 3NF xử lý phụ thuộc không khóa kiểu bắc cầu, còn BCNF yêu cầu mọi vế trái không tầm thường là siêu khóa.

Phân tích: Với phụ thuộc X → Y, mọi hai bộ có cùng giá trị X phải có cùng giá trị Y. Để kiểm tra khóa, tính X⁺ bằng cách áp dụng lặp các phụ thuộc; nếu bao đóng chứa toàn bộ thuộc tính thì X là siêu khóa, sau đó kiểm tra tính tối thiểu.

Đối chiếu phương án: B. Each relation depends on another relation; C. Relations can be modified without affecting each other; D. All attributes must be independent of each other. Các lựa chọn này không đáp ứng định nghĩa/điều kiện được hỏi; khi ôn, hãy chỉ ra chính xác điểm lệch của từng lựa chọn thay vì học thuộc mỗi chữ cái.
\end{question}

\begin{question}{168}
% [ID gốc: 15.67]
% [Nguồn: DBI202 Order 15 Q67]
A relation R(X, Y, Z) has functional dependencies XY → Z and Z → Y. Assuming X, Y, Z are single attributes and {X, Y} is the only candidate key, what is the highest normal form that R satisfies?
\begin{choices}
\choice First Normal Form (1NF) only
\choice Second Normal Form (2NF) but not Third Normal Form (3NF)
\choice Third Normal Form (3NF) but not Boyce-Codd Normal Form (BCNF)
\choice Boyce-Codd Normal Form (BCNF)
\end{choices}
\textbf{Đáp án:} C
\textbf{Giải thích:} Đáp án theo dữ liệu đề: C — C. Third Normal Form (3NF) but not Boyce-Codd Normal Form (BCNF).

Lời giải gốc: Phụ thuộc hàm Z → Y có vế trái Z không phải là siêu khóa, vi phạm điều kiện BCNF. Tuy nhiên Y là thuộc tính khóa (nằm trong khóa {X, Y}) nên vẫn thỏa 3NF. Do đó dạng chuẩn cao nhất là 3NF.

Nền tảng kiến thức: Chuẩn hóa dùng phụ thuộc hàm và khóa để giảm dư thừa cũng như bất thường cập nhật. Tính bao đóng thuộc tính để tìm khóa; 2NF loại phụ thuộc bộ phận vào một phần khóa ghép, 3NF xử lý phụ thuộc không khóa kiểu bắc cầu, còn BCNF yêu cầu mọi vế trái không tầm thường là siêu khóa.

Phân tích: Đánh giá dạng chuẩn theo thứ tự: 1NF yêu cầu giá trị nguyên tố/không lặp theo mô hình quan hệ; 2NF loại phụ thuộc bộ phận của thuộc tính không khóa vào một phần khóa ghép; 3NF kiểm soát phụ thuộc bắc cầu theo điều kiện prime attribute; BCNF yêu cầu mọi determinant không tầm thường là siêu khóa.

Đối chiếu phương án: A. First Normal Form (1NF) only; B. Second Normal Form (2NF) but not Third Normal Form (3NF); D. Boyce-Codd Normal Form (BCNF). Các lựa chọn này không đáp ứng định nghĩa/điều kiện được hỏi; khi ôn, hãy chỉ ra chính xác điểm lệch của từng lựa chọn thay vì học thuộc mỗi chữ cái.
\end{question}

\begin{question}{169}
% [ID gốc: 15.69]
% [Nguồn: DBI202 Order 15 Q69]
In a student registration system, which normal form is violated if one field contains multiple course IDs for the same student?
\begin{choices}
\choice 1NF
\choice 2NF
\choice 3NF
\choice 4NF
\end{choices}
\textbf{Đáp án:} A
\textbf{Giải thích:} Đáp án theo dữ liệu đề: A — A. 1NF.

Lời giải gốc: Như giải thích ở Câu 12, việc một trường dữ liệu (field) chứa nhiều giá trị đồng thời (nhiều mã môn học - course IDs) đã vi phạm tính nguyên tố của dữ liệu, trực tiếp vi phạm Dạng chuẩn 1 (1NF).

Nền tảng kiến thức: Chuẩn hóa dùng phụ thuộc hàm và khóa để giảm dư thừa cũng như bất thường cập nhật. Tính bao đóng thuộc tính để tìm khóa; 2NF loại phụ thuộc bộ phận vào một phần khóa ghép, 3NF xử lý phụ thuộc không khóa kiểu bắc cầu, còn BCNF yêu cầu mọi vế trái không tầm thường là siêu khóa.

Phân tích: Đánh giá dạng chuẩn theo thứ tự: 1NF yêu cầu giá trị nguyên tố/không lặp theo mô hình quan hệ; 2NF loại phụ thuộc bộ phận của thuộc tính không khóa vào một phần khóa ghép; 3NF kiểm soát phụ thuộc bắc cầu theo điều kiện prime attribute; BCNF yêu cầu mọi determinant không tầm thường là siêu khóa.

Đối chiếu phương án: B. 2NF; C. 3NF; D. 4NF. Các lựa chọn này không đáp ứng định nghĩa/điều kiện được hỏi; khi ôn, hãy chỉ ra chính xác điểm lệch của từng lựa chọn thay vì học thuộc mỗi chữ cái.
\end{question}

\begin{question}{170}
% [ID gốc: 15.70]
% [Nguồn: DBI202 Order 15 Q70]
In relational database design, which normal form eliminates multi-valued attributes and ensures that data is organized at a fundamental level?
\begin{choices}
\choice 1NF
\choice 2NF
\choice 3NF
\choice 4NF
\end{choices}
\textbf{Đáp án:} A
\textbf{Giải thích:} Đáp án theo dữ liệu đề: A — A. 1NF.

Lời giải gốc: Dạng chuẩn 1 (1NF) yêu cầu các thuộc tính phải là nguyên tố, qua đó loại bỏ các thuộc tính đa trị hoặc nhóm lặp lại.

Nền tảng kiến thức: Chuẩn hóa dùng phụ thuộc hàm và khóa để giảm dư thừa cũng như bất thường cập nhật. Tính bao đóng thuộc tính để tìm khóa; 2NF loại phụ thuộc bộ phận vào một phần khóa ghép, 3NF xử lý phụ thuộc không khóa kiểu bắc cầu, còn BCNF yêu cầu mọi vế trái không tầm thường là siêu khóa.

Phân tích: Đánh giá dạng chuẩn theo thứ tự: 1NF yêu cầu giá trị nguyên tố/không lặp theo mô hình quan hệ; 2NF loại phụ thuộc bộ phận của thuộc tính không khóa vào một phần khóa ghép; 3NF kiểm soát phụ thuộc bắc cầu theo điều kiện prime attribute; BCNF yêu cầu mọi determinant không tầm thường là siêu khóa.

Đối chiếu phương án: B. 2NF; C. 3NF; D. 4NF. Các lựa chọn này không đáp ứng định nghĩa/điều kiện được hỏi; khi ôn, hãy chỉ ra chính xác điểm lệch của từng lựa chọn thay vì học thuộc mỗi chữ cái.
\end{question}

\begin{question}{171}
% [ID gốc: 15.72]
% [Nguồn: DBI202 Order 15 Q72]
Which normalization form eliminates transitive dependency?
\begin{choices}
\choice 1NF
\choice 2NF
\choice 3NF
\choice BCNF
\end{choices}
\textbf{Đáp án:} C
\textbf{Giải thích:} Đáp án theo dữ liệu đề: C — C. 3NF.

Lời giải gốc: Dạng chuẩn 3 (3NF) yêu cầu bảng phải đạt chuẩn 2NF và loại bỏ hoàn toàn các phụ thuộc hàm bắc cầu (Transitive dependencies) của các thuộc tính không phải khóa vào khóa chính.

Nền tảng kiến thức: Chuẩn hóa dùng phụ thuộc hàm và khóa để giảm dư thừa cũng như bất thường cập nhật. Tính bao đóng thuộc tính để tìm khóa; 2NF loại phụ thuộc bộ phận vào một phần khóa ghép, 3NF xử lý phụ thuộc không khóa kiểu bắc cầu, còn BCNF yêu cầu mọi vế trái không tầm thường là siêu khóa.

Phân tích: Đánh giá dạng chuẩn theo thứ tự: 1NF yêu cầu giá trị nguyên tố/không lặp theo mô hình quan hệ; 2NF loại phụ thuộc bộ phận của thuộc tính không khóa vào một phần khóa ghép; 3NF kiểm soát phụ thuộc bắc cầu theo điều kiện prime attribute; BCNF yêu cầu mọi determinant không tầm thường là siêu khóa.

Đối chiếu phương án: A. 1NF; B. 2NF; D. BCNF. Các lựa chọn này không đáp ứng định nghĩa/điều kiện được hỏi; khi ôn, hãy chỉ ra chính xác điểm lệch của từng lựa chọn thay vì học thuộc mỗi chữ cái.
\end{question}

\begin{question}{172}
% [ID gốc: 15.73]
% [Nguồn: DBI202 Order 15 Q73]
Suppose there is a relation R(A, B, C, D, E, G) and functional dependencies F = {AB → C, D → EG, C → A, BE → C, BC → D, CG → BD, ACD → B, CE → AG}. Compute the closure {AB}+.
\begin{choices}
\choice ABCD
\choice ABCDE
\choice ABCDG
\choice ABCDEG
\end{choices}
\textbf{Đáp án:} D
\textbf{Giải thích:} Đáp án theo dữ liệu đề: D — D. ABCDEG.

Lời giải gốc: Dựa vào tập phụ thuộc hàm, bao đóng của AB tính được chứa toàn bộ các thuộc tính A, B, C, D, E, G của quan hệ. Bắt đầu từ AB: AB → C, có C → A và BC → D nên thêm D, D → EG nên thêm E, G. Kết quả là ABCDEG.

Nền tảng kiến thức: Chuẩn hóa dùng phụ thuộc hàm và khóa để giảm dư thừa cũng như bất thường cập nhật. Tính bao đóng thuộc tính để tìm khóa; 2NF loại phụ thuộc bộ phận vào một phần khóa ghép, 3NF xử lý phụ thuộc không khóa kiểu bắc cầu, còn BCNF yêu cầu mọi vế trái không tầm thường là siêu khóa.

Phân tích: Với phụ thuộc X → Y, mọi hai bộ có cùng giá trị X phải có cùng giá trị Y. Để kiểm tra khóa, tính X⁺ bằng cách áp dụng lặp các phụ thuộc; nếu bao đóng chứa toàn bộ thuộc tính thì X là siêu khóa, sau đó kiểm tra tính tối thiểu.

Đối chiếu phương án: A. ABCD; B. ABCDE; C. ABCDG. Các lựa chọn này không đáp ứng định nghĩa/điều kiện được hỏi; khi ôn, hãy chỉ ra chính xác điểm lệch của từng lựa chọn thay vì học thuộc mỗi chữ cái.
\end{question}

\begin{question}{173}
% [ID gốc: 15.74]
% [Nguồn: DBI202 Order 15 Q74]
Consider a relation R(A, B, C, D) with functional dependencies {AB → C, C → D, D → A}. What are all the keys of R?
\begin{choices}
\choice AB, BC, BD
\choice AB, BC, AD
\choice BC, BD
\choice AB, AD
\end{choices}
\textbf{Đáp án:} A
\textbf{Giải thích:} Đáp án theo dữ liệu đề: A — A. AB, BC, BD.

Lời giải gốc: Dựa vào các phụ thuộc hàm, các khóa tối thiểu của quan hệ gồm các tập {A, B}, {B, C} và {B, D}. Thuộc tính B luôn phải có mặt vì không FD nào sinh ra B.

Nền tảng kiến thức: Chuẩn hóa dùng phụ thuộc hàm và khóa để giảm dư thừa cũng như bất thường cập nhật. Tính bao đóng thuộc tính để tìm khóa; 2NF loại phụ thuộc bộ phận vào một phần khóa ghép, 3NF xử lý phụ thuộc không khóa kiểu bắc cầu, còn BCNF yêu cầu mọi vế trái không tầm thường là siêu khóa.

Phân tích: Với phụ thuộc X → Y, mọi hai bộ có cùng giá trị X phải có cùng giá trị Y. Để kiểm tra khóa, tính X⁺ bằng cách áp dụng lặp các phụ thuộc; nếu bao đóng chứa toàn bộ thuộc tính thì X là siêu khóa, sau đó kiểm tra tính tối thiểu.

Đối chiếu phương án: B. AB, BC, AD; C. BC, BD; D. AB, AD. Các lựa chọn này không đáp ứng định nghĩa/điều kiện được hỏi; khi ôn, hãy chỉ ra chính xác điểm lệch của từng lựa chọn thay vì học thuộc mỗi chữ cái.
\end{question}

\begin{question}{174}
% [ID gốc: 15.75]
% [Nguồn: DBI202 Order 15 Q75]
Suppose that we have a relation R (B, C, D) that satisfies FD's: D-&gt;B and D-&gt;C. Which of the following FD is equivalent to the given dependency functions?
\begin{choices}
\choice D-&gt;BC
\choice D-&gt;DB
\choice DB-&gt;B
\choice DC-&gt;DB
\end{choices}
\textbf{Đáp án:} A
\textbf{Giải thích:} Đáp án theo dữ liệu đề: A — A. D-&gt;BC.

Lời giải gốc: Theo quy tắc hợp (Union rule) trong hệ tiên đề Armstrong, nếu $D \to B$ và $D \to C$ thì tương đương với $D \to BC$. ---

Nền tảng kiến thức: Chuẩn hóa dùng phụ thuộc hàm và khóa để giảm dư thừa cũng như bất thường cập nhật. Tính bao đóng thuộc tính để tìm khóa; 2NF loại phụ thuộc bộ phận vào một phần khóa ghép, 3NF xử lý phụ thuộc không khóa kiểu bắc cầu, còn BCNF yêu cầu mọi vế trái không tầm thường là siêu khóa.

Phân tích: Với phụ thuộc X → Y, mọi hai bộ có cùng giá trị X phải có cùng giá trị Y. Để kiểm tra khóa, tính X⁺ bằng cách áp dụng lặp các phụ thuộc; nếu bao đóng chứa toàn bộ thuộc tính thì X là siêu khóa, sau đó kiểm tra tính tối thiểu.

Đối chiếu phương án: B. D-&gt;DB; C. DB-&gt;B; D. DC-&gt;DB. Các lựa chọn này không đáp ứng định nghĩa/điều kiện được hỏi; khi ôn, hãy chỉ ra chính xác điểm lệch của từng lựa chọn thay vì học thuộc mỗi chữ cái.
\end{question}

\begin{question}{175}
% [ID gốc: 15.76]
% [Nguồn: DBI202 Order 15 Q76]
Consider a relation R(A, B, C, D, E, F) with the following set of functional dependencies:
$$F = \{ AB \rightarrow C, \quad C \rightarrow D, \quad D \rightarrow E, \quad E \rightarrow F \}$$
Consider the following candidate functional dependencies:
- \(f_1: A \rightarrow D\)
- \(f_2: AB \rightarrow D\)
- \(f_3: AB \rightarrow E\)
- \(f_4: AB \rightarrow F\)

Based on Armstrong's inference rules, which of the functional dependencies can be derived from F?
\begin{choices}
\choice f1, f3, f4
\choice f1, f2, f3, f4
\choice f2, f3, f4
\choice f1, f2, f4
\end{choices}
\textbf{Đáp án:} C
\textbf{Giải thích:} Áp dụng luật bắc cầu (Transitivity) của hệ tiên đề Armstrong: từ AB -> C và C -> D suy ra AB -> D (f2); từ AB -> D và D -> E suy ra AB -> E (f3); từ AB -> E và E -> F suy ra AB -> F (f4). Phụ thuộc hàm f1: A -> D không thể suy diễn được vì thiếu thuộc tính B để suy ra C.
\end{question}

\begin{question}{176}
% [ID gốc: 15.77]
% [Nguồn: DBI202 Order 15 Q77]
If A → B and B → C, which rule allows us to infer A → C?
\begin{choices}
\choice Reflexivity
\choice Augmentation
\choice Transitivity
\choice Projection
\end{choices}
\textbf{Đáp án:} C
\textbf{Giải thích:} Đáp án theo dữ liệu đề: C — C. Transitivity.

Lời giải gốc: Đây là luật bắc cầu (Transitivity rule) trong hệ tiên đề Armstrong.

Nền tảng kiến thức: Chuẩn hóa dùng phụ thuộc hàm và khóa để giảm dư thừa cũng như bất thường cập nhật. Tính bao đóng thuộc tính để tìm khóa; 2NF loại phụ thuộc bộ phận vào một phần khóa ghép, 3NF xử lý phụ thuộc không khóa kiểu bắc cầu, còn BCNF yêu cầu mọi vế trái không tầm thường là siêu khóa.

Phân tích: Tiên đề bắc cầu của Armstrong: nếu X → Y và Y → Z thì suy ra X → Z. Phải kiểm tra hai phụ thuộc nối tiếp cùng vế giữa Y; các tiên đề phản xạ và tăng trưởng mô tả quy tắc khác.

Đối chiếu phương án: A. Reflexivity; B. Augmentation; D. Projection. Các lựa chọn này không đáp ứng định nghĩa/điều kiện được hỏi; khi ôn, hãy chỉ ra chính xác điểm lệch của từng lựa chọn thay vì học thuộc mỗi chữ cái.
\end{question}

\begin{question}{177}
% [ID gốc: 15.80]
% [Nguồn: DBI202 Order 15 Q80]
Given the set of functional dependencies F={AB→C,C→D}. Determine the highest normal form of the relation R(A,B,C,D).
\begin{choices}
\choice 1NF
\choice 2NF
\choice 3NF
\choice BCNF
\end{choices}
\textbf{Đáp án:} B
\textbf{Giải thích:} Đáp án theo dữ liệu đề: B — B. 2NF.

Lời giải gốc: Khóa chính của quan hệ là $AB$. Do có phụ thuộc hàm bắc cầu $AB \to C \to D$ (trong đó $C$ không phải là siêu khóa), bảng không đạt chuẩn 3NF mà chỉ đạt chuẩn 2NF. ---

Nền tảng kiến thức: Chuẩn hóa dùng phụ thuộc hàm và khóa để giảm dư thừa cũng như bất thường cập nhật. Tính bao đóng thuộc tính để tìm khóa; 2NF loại phụ thuộc bộ phận vào một phần khóa ghép, 3NF xử lý phụ thuộc không khóa kiểu bắc cầu, còn BCNF yêu cầu mọi vế trái không tầm thường là siêu khóa.

Phân tích: Đánh giá dạng chuẩn theo thứ tự: 1NF yêu cầu giá trị nguyên tố/không lặp theo mô hình quan hệ; 2NF loại phụ thuộc bộ phận của thuộc tính không khóa vào một phần khóa ghép; 3NF kiểm soát phụ thuộc bắc cầu theo điều kiện prime attribute; BCNF yêu cầu mọi determinant không tầm thường là siêu khóa.

Đối chiếu phương án: A. 1NF; C. 3NF; D. BCNF. Các lựa chọn này không đáp ứng định nghĩa/điều kiện được hỏi; khi ôn, hãy chỉ ra chính xác điểm lệch của từng lựa chọn thay vì học thuộc mỗi chữ cái.
\end{question}

\begin{question}{178}
% [ID gốc: 15.82]
% [Nguồn: DBI202 Order 15 Q82]
What is a benefit of normalization when converting an ERD to a relational model?
\begin{choices}
\choice Increased storage space efficiency
\choice Reduced update anomalies
\choice Simplified database queries
\choice Enhanced data security
\end{choices}
\textbf{Đáp án:} B
\textbf{Giải thích:} Đáp án theo dữ liệu đề: B — B. Reduced update anomalies.

Lời giải gốc: Chuẩn hóa giúp giảm thiểu tối đa các bất thường khi cập nhật dữ liệu (update anomalies).

Nền tảng kiến thức: Chuẩn hóa dùng phụ thuộc hàm và khóa để giảm dư thừa cũng như bất thường cập nhật. Tính bao đóng thuộc tính để tìm khóa; 2NF loại phụ thuộc bộ phận vào một phần khóa ghép, 3NF xử lý phụ thuộc không khóa kiểu bắc cầu, còn BCNF yêu cầu mọi vế trái không tầm thường là siêu khóa.

Phân tích: Khi đọc mô hình ER, tách ba thứ: entity set là loại đối tượng, relationship set là liên kết giữa các loại đối tượng, còn cardinality/participation là ràng buộc số lượng hoặc mức tham gia. Đáp án đúng phải gọi đúng thành phần mà câu hỏi mô tả.

Đối chiếu phương án: A. Increased storage space efficiency; C. Simplified database queries; D. Enhanced data security. Các lựa chọn này không đáp ứng định nghĩa/điều kiện được hỏi; khi ôn, hãy chỉ ra chính xác điểm lệch của từng lựa chọn thay vì học thuộc mỗi chữ cái.
\end{question}

\begin{question}{179}
% [ID gốc: 15.83]
% [Nguồn: DBI202 Order 15 Q83]
The highest normal form for relation schema R(ABCD) with functional dependencies: F = {AB-&gt; C; D-&gt;B, C-&gt;ABD} is:
\begin{choices}
\choice 2NF
\choice 1NF
\choice 3NF
\choice BCNF
\end{choices}
\textbf{Đáp án:} C
\textbf{Giải thích:} Đáp án theo dữ liệu đề: C — C. 3NF.

Lời giải gốc: Khóa của R gồm: {C}, {A,B}, {A,D}. Nhờ đó, cả 4 thuộc tính đều là thuộc tính khóa (prime attributes). Bất kỳ lược đồ nào mà tất cả các thuộc tính đều là thuộc tính khóa thì chắc chắn đạt 3NF (không có sự phụ thuộc bắc cầu của thuộc tính không khóa). Nó vi phạm BCNF vì phụ thuộc D->B có vế trái (D) không phải là siêu khóa. Vậy dạng chuẩn cao nhất đạt được là 3NF.

Nền tảng kiến thức: Chuẩn hóa dùng phụ thuộc hàm và khóa để giảm dư thừa cũng như bất thường cập nhật. Tính bao đóng thuộc tính để tìm khóa; 2NF loại phụ thuộc bộ phận vào một phần khóa ghép, 3NF xử lý phụ thuộc không khóa kiểu bắc cầu, còn BCNF yêu cầu mọi vế trái không tầm thường là siêu khóa.

Phân tích: Lược đồ quan hệ ghi tên quan hệ kèm danh sách thuộc tính, thường theo dạng `Boat(BoatID, BoatName, Color, Status)`. Nếu đề hỏi schema thuần túy, không nhất thiết phải kèm kiểu dữ liệu hay các bộ dữ liệu đang lưu.

Đối chiếu phương án: A. 2NF; B. 1NF; D. BCNF. Các lựa chọn này không đáp ứng định nghĩa/điều kiện được hỏi; khi ôn, hãy chỉ ra chính xác điểm lệch của từng lựa chọn thay vì học thuộc mỗi chữ cái.
\end{question}

\begin{question}{180}
% [ID gốc: 15.84]
% [Nguồn: DBI202 Order 15 Q84]
Which condition best describes a table that is in Second Normal Form (2NF)?
\begin{choices}
\choice It has no primary key.
\choice It has atomic attributes.
\choice It has no partial dependency on a composite key.
\choice It allows repeating groups and no transitive dependency exists.
\end{choices}
\textbf{Đáp án:} C
\textbf{Giải thích:} Đáp án theo dữ liệu đề: C — C. It has no partial dependency on a composite key..

Lời giải gốc: Một bảng đạt dạng chuẩn 2 (2NF) khi nó ở dạng chuẩn 1 và không có phụ thuộc một phần (partial dependency) của các thuộc tính không khóa vào khóa chính phức hợp. ---

Nền tảng kiến thức: Chuẩn hóa dùng phụ thuộc hàm và khóa để giảm dư thừa cũng như bất thường cập nhật. Tính bao đóng thuộc tính để tìm khóa; 2NF loại phụ thuộc bộ phận vào một phần khóa ghép, 3NF xử lý phụ thuộc không khóa kiểu bắc cầu, còn BCNF yêu cầu mọi vế trái không tầm thường là siêu khóa.

Phân tích: Đánh giá dạng chuẩn theo thứ tự: 1NF yêu cầu giá trị nguyên tố/không lặp theo mô hình quan hệ; 2NF loại phụ thuộc bộ phận của thuộc tính không khóa vào một phần khóa ghép; 3NF kiểm soát phụ thuộc bắc cầu theo điều kiện prime attribute; BCNF yêu cầu mọi determinant không tầm thường là siêu khóa.

Đối chiếu phương án: A. It has no primary key.; B. It has atomic attributes.; D. It allows repeating groups and no transitive dependency exists.. Các lựa chọn này không đáp ứng định nghĩa/điều kiện được hỏi; khi ôn, hãy chỉ ra chính xác điểm lệch của từng lựa chọn thay vì học thuộc mỗi chữ cái.
\end{question}

\begin{question}{181}
% [ID gốc: 15.85]
% [Nguồn: DBI202 Order 15 Q85]
Given relation R(X, Y, Z, W) with the set of functional dependencies: F = {X→Y, Y→Z, Z→W} Which of the following is a key for R?
\begin{choices}
\choice X
\choice Y
\choice XZ
\choice XY
\end{choices}
\textbf{Đáp án:} A
\textbf{Giải thích:} Đáp án theo dữ liệu đề: A — A. X.

Lời giải gốc: Xét bao đóng của X: $X^+ = \{X, Y, Z, W\}$. Vì $X^+$ chứa tất cả các thuộc tính của R nên X là khóa (Key) của R.

Nền tảng kiến thức: Chuẩn hóa dùng phụ thuộc hàm và khóa để giảm dư thừa cũng như bất thường cập nhật. Tính bao đóng thuộc tính để tìm khóa; 2NF loại phụ thuộc bộ phận vào một phần khóa ghép, 3NF xử lý phụ thuộc không khóa kiểu bắc cầu, còn BCNF yêu cầu mọi vế trái không tầm thường là siêu khóa.

Phân tích: Khóa là thuộc tính hoặc tập thuộc tính dùng để nhận diện duy nhất một bộ/hàng. Khóa có thể gồm một hay nhiều cột, nên không bắt buộc luôn là cột đơn hoặc luôn gồm nhiều cột; nó cũng không “nhận diện một cột”.

Đối chiếu phương án: B. Y; C. XZ; D. XY. Các lựa chọn này không đáp ứng định nghĩa/điều kiện được hỏi; khi ôn, hãy chỉ ra chính xác điểm lệch của từng lựa chọn thay vì học thuộc mỗi chữ cái.
\end{question}

\begin{question}{182}
% [ID gốc: 15.89]
% [Nguồn: DBI202 Order 15 Q89]
In a functional dependency A → B, what can be said about B?
\begin{choices}
\choice B is also a superkey
\choice B is a determinant
\choice B is functionally dependent on A
\choice B is a foreign key
\end{choices}
\textbf{Đáp án:} C
\textbf{Giải thích:} Đáp án theo dữ liệu đề: C — C. B is functionally dependent on A.

Lời giải gốc: Trong phụ thuộc hàm A → B, thuộc tính B phụ thuộc hàm vào tập thuộc tính A.

Nền tảng kiến thức: Chuẩn hóa dùng phụ thuộc hàm và khóa để giảm dư thừa cũng như bất thường cập nhật. Tính bao đóng thuộc tính để tìm khóa; 2NF loại phụ thuộc bộ phận vào một phần khóa ghép, 3NF xử lý phụ thuộc không khóa kiểu bắc cầu, còn BCNF yêu cầu mọi vế trái không tầm thường là siêu khóa.

Phân tích: Với phụ thuộc X → Y, mọi hai bộ có cùng giá trị X phải có cùng giá trị Y. Để kiểm tra khóa, tính X⁺ bằng cách áp dụng lặp các phụ thuộc; nếu bao đóng chứa toàn bộ thuộc tính thì X là siêu khóa, sau đó kiểm tra tính tối thiểu.

Đối chiếu phương án: A. B is also a superkey; B. B is a determinant; D. B is a foreign key. Các lựa chọn này không đáp ứng định nghĩa/điều kiện được hỏi; khi ôn, hãy chỉ ra chính xác điểm lệch của từng lựa chọn thay vì học thuộc mỗi chữ cái.
\end{question}

\begin{question}{183}
% [ID gốc: 15.90]
% [Nguồn: DBI202 Order 15 Q90]
DB1202_SP26_RE | Multiple Choice Question 29 (Choose 2 answers) During database design, why is it necessary to identify candidate keys? Choose two answers
\begin{choices}
\choice To reduce data type complexity.
\choice To normalize tables faster
\choice To select the most appropriate primary key
\choice To remove NULL values
\end{choices}
\textbf{Đáp án:} BC
\textbf{Giải thích:} Đáp án theo dữ liệu đề: B, C — B. To normalize tables faster; C. To select the most appropriate primary key.

Lời giải gốc: Việc xác định các khóa ứng viên (Candidate keys) trong quá trình thiết kế là bắt buộc để: Giúp chúng ta phân tích các phụ thuộc hàm nhằm tiến hành chuẩn hóa các bảng dữ liệu nhanh chóng và chính xác (B đúng). Từ danh sách các khóa ứng viên đó, người thiết kế sẽ chọn ra một khóa phù hợp nhất để làm khóa chính (Primary key) cho bảng (C đúng).

Nền tảng kiến thức: Chuẩn hóa dùng phụ thuộc hàm và khóa để giảm dư thừa cũng như bất thường cập nhật. Tính bao đóng thuộc tính để tìm khóa; 2NF loại phụ thuộc bộ phận vào một phần khóa ghép, 3NF xử lý phụ thuộc không khóa kiểu bắc cầu, còn BCNF yêu cầu mọi vế trái không tầm thường là siêu khóa.

Phân tích: NULL biểu thị thiếu/không biết/không áp dụng, không phải số 0 hay chuỗi rỗng. So sánh với NULL bằng `=` không cho TRUE/FALSE mà cho UNKNOWN; dùng `IS NULL`/`IS NOT NULL`. Trong logic ba giá trị, AND/OR/NOT phải được tính theo cả trạng thái UNKNOWN.

Đối chiếu phương án: A. To reduce data type complexity.; D. To remove NULL values. Các lựa chọn này không đáp ứng định nghĩa/điều kiện được hỏi; khi ôn, hãy chỉ ra chính xác điểm lệch của từng lựa chọn thay vì học thuộc mỗi chữ cái.
\end{question}

\begin{question}{184}
% [ID gốc: 15.92]
% [Nguồn: DBI202 Order 15 Q92]
Suppose we a relation R(A,B,C,D,E,F,N) and functional dependencies: F-&gt;A, A-&gt;B, B-&gt;C, BC-&gt;A, A-&gt;D, E-&gt;A, D-&gt;E Which of the following is not a key? Choose correct answers
\begin{choices}
\choice AB
\choice FN
\choice E
\choice D
\end{choices}
\textbf{Đáp án:} ACD
\textbf{Giải thích:} Đáp án theo dữ liệu đề: A, C, D — A. AB; C. E; D. D.

Lời giải gốc: Để là khóa, bao đóng của tập thuộc tính phải chứa tất cả thuộc tính {A,B,C,D,E,F,N}. Ta thấy thuộc tính $N$ không xuất hiện ở vế phải của bất kỳ phụ thuộc hàm nào, nên bất kỳ khóa nào cũng bắt buộc phải chứa $N$. Các tập $AB$, $E$, $D$ không chứa $N$ nên chắc chắn không phải là khóa. Chỉ có $FN$ có $(FN)^+ = {F,N,A,B,C,D,E}$ là khóa. Do đó các phương án không phải khóa là A, C, D.

Nền tảng kiến thức: Chuẩn hóa dùng phụ thuộc hàm và khóa để giảm dư thừa cũng như bất thường cập nhật. Tính bao đóng thuộc tính để tìm khóa; 2NF loại phụ thuộc bộ phận vào một phần khóa ghép, 3NF xử lý phụ thuộc không khóa kiểu bắc cầu, còn BCNF yêu cầu mọi vế trái không tầm thường là siêu khóa.

Phân tích: Khóa là thuộc tính hoặc tập thuộc tính dùng để nhận diện duy nhất một bộ/hàng. Khóa có thể gồm một hay nhiều cột, nên không bắt buộc luôn là cột đơn hoặc luôn gồm nhiều cột; nó cũng không “nhận diện một cột”.

Đối chiếu phương án: B. FN. Các lựa chọn này không đáp ứng định nghĩa/điều kiện được hỏi; khi ôn, hãy chỉ ra chính xác điểm lệch của từng lựa chọn thay vì học thuộc mỗi chữ cái.
\end{question}

\begin{question}{185}
% [ID gốc: 15.94]
% [Nguồn: DBI202 Order 15 Q94]
Normalization is:
\begin{choices}
\choice The process of ensuring that each table has a key
\choice The process whereby a limit is put on the number of fields a record can contain
\choice The process of organizing data efficiently and minimizing repetition, helping maintain data consistency and reducing the risk of errors
\choice The process of ensuring that a relational database has at least two tables
\end{choices}
\textbf{Đáp án:} C
\textbf{Giải thích:} Đáp án theo dữ liệu đề: C — C. The process of organizing data efficiently and minimizing repetition, helping maintain data consistency and reducing the risk of errors.

Lời giải gốc: Chuẩn hóa là quá trình tổ chức dữ liệu hiệu quả nhằm giảm thiểu sự lặp lại, duy trì tính nhất quán và giảm thiểu lỗi dữ liệu.

Nền tảng kiến thức: Chuẩn hóa dùng phụ thuộc hàm và khóa để giảm dư thừa cũng như bất thường cập nhật. Tính bao đóng thuộc tính để tìm khóa; 2NF loại phụ thuộc bộ phận vào một phần khóa ghép, 3NF xử lý phụ thuộc không khóa kiểu bắc cầu, còn BCNF yêu cầu mọi vế trái không tầm thường là siêu khóa.

Phân tích: Khóa là thuộc tính hoặc tập thuộc tính dùng để nhận diện duy nhất một bộ/hàng. Khóa có thể gồm một hay nhiều cột, nên không bắt buộc luôn là cột đơn hoặc luôn gồm nhiều cột; nó cũng không “nhận diện một cột”.

Đối chiếu phương án: A. The process of ensuring that each table has a key; B. The process whereby a limit is put on the number of fields a record can contain; D. The process of ensuring that a relational database has at least two tables. Các lựa chọn này không đáp ứng định nghĩa/điều kiện được hỏi; khi ôn, hãy chỉ ra chính xác điểm lệch của từng lựa chọn thay vì học thuộc mỗi chữ cái.
\end{question}

\begin{question}{186}
% [ID gốc: 15.95]
% [Nguồn: DBI202 Order 15 Q95]
Suppose we have a relation that stores data about a single entity and has no multi-valued dependencies. Which of the following is the highest normal form for this relation?
\begin{choices}
\choice 1NF
\choice 2NF
\choice 3NF
\choice 4NF
\end{choices}
\textbf{Đáp án:} D
\textbf{Giải thích:} Đáp án theo dữ liệu đề: D — D. 4NF.

Lời giải gốc: Dạng chuẩn 4 (4NF) loại bỏ các phụ thuộc đa trị, đảm bảo mô hình đạt mức tối ưu cao hơn 3NF.

Nền tảng kiến thức: Chuẩn hóa dùng phụ thuộc hàm và khóa để giảm dư thừa cũng như bất thường cập nhật. Tính bao đóng thuộc tính để tìm khóa; 2NF loại phụ thuộc bộ phận vào một phần khóa ghép, 3NF xử lý phụ thuộc không khóa kiểu bắc cầu, còn BCNF yêu cầu mọi vế trái không tầm thường là siêu khóa.

Phân tích: Đánh giá dạng chuẩn theo thứ tự: 1NF yêu cầu giá trị nguyên tố/không lặp theo mô hình quan hệ; 2NF loại phụ thuộc bộ phận của thuộc tính không khóa vào một phần khóa ghép; 3NF kiểm soát phụ thuộc bắc cầu theo điều kiện prime attribute; BCNF yêu cầu mọi determinant không tầm thường là siêu khóa.

Đối chiếu phương án: A. 1NF; B. 2NF; C. 3NF. Các lựa chọn này không đáp ứng định nghĩa/điều kiện được hỏi; khi ôn, hãy chỉ ra chính xác điểm lệch của từng lựa chọn thay vì học thuộc mỗi chữ cái.
\end{question}

\begin{question}{187}
% [ID gốc: 15.96]
% [Nguồn: DBI202 Order 15 Q96]
Which statement best represents Armstrong's Axiom of Decomposition?
\begin{choices}
\choice If A → BC, then A → B and A → C
\choice If A → B and B → C, then A → C
\choice If A → B, then AC → BC
\choice If A → B, then B → C
\end{choices}
\textbf{Đáp án:} A
\textbf{Giải thích:} Đáp án theo dữ liệu đề: A — A. If A → BC, then A → B and A → C.

Lời giải gốc: Luật phân rã (Decomposition rule) trong hệ tiên đề Armstrong phát biểu rằng nếu A → BC thì A → B và A → C.

Nền tảng kiến thức: Chuẩn hóa dùng phụ thuộc hàm và khóa để giảm dư thừa cũng như bất thường cập nhật. Tính bao đóng thuộc tính để tìm khóa; 2NF loại phụ thuộc bộ phận vào một phần khóa ghép, 3NF xử lý phụ thuộc không khóa kiểu bắc cầu, còn BCNF yêu cầu mọi vế trái không tầm thường là siêu khóa.

Phân tích: Tiên đề bắc cầu của Armstrong: nếu X → Y và Y → Z thì suy ra X → Z. Phải kiểm tra hai phụ thuộc nối tiếp cùng vế giữa Y; các tiên đề phản xạ và tăng trưởng mô tả quy tắc khác.

Đối chiếu phương án: B. If A → B and B → C, then A → C; C. If A → B, then AC → BC; D. If A → B, then B → C. Các lựa chọn này không đáp ứng định nghĩa/điều kiện được hỏi; khi ôn, hãy chỉ ra chính xác điểm lệch của từng lựa chọn thay vì học thuộc mỗi chữ cái.
\end{question}

\begin{question}{188}
% [ID gốc: 15.97]
% [Nguồn: DBI202 Order 15 Q97]
Which of the following is an example of a violation of 2NF?
\begin{choices}
\choice A non-prime attribute is functionally dependent on part of a primary key.
\choice A non-prime attribute is functionally dependent on a non-prime attribute.
\choice A non-prime attribute is functionally dependent on the entire candidate key.
\choice There is a transitive dependency between candidate keys.
\end{choices}
\textbf{Đáp án:} A
\textbf{Giải thích:} Đáp án theo dữ liệu đề: A — A. A non-prime attribute is functionally dependent on part of a primary key..

Lời giải gốc: Vi phạm chuẩn 2 xảy ra khi thuộc tính không khóa phụ thuộc hàm vào một phần của khóa chính phức hợp. ---

Nền tảng kiến thức: Chuẩn hóa dùng phụ thuộc hàm và khóa để giảm dư thừa cũng như bất thường cập nhật. Tính bao đóng thuộc tính để tìm khóa; 2NF loại phụ thuộc bộ phận vào một phần khóa ghép, 3NF xử lý phụ thuộc không khóa kiểu bắc cầu, còn BCNF yêu cầu mọi vế trái không tầm thường là siêu khóa.

Phân tích: Đánh giá dạng chuẩn theo thứ tự: 1NF yêu cầu giá trị nguyên tố/không lặp theo mô hình quan hệ; 2NF loại phụ thuộc bộ phận của thuộc tính không khóa vào một phần khóa ghép; 3NF kiểm soát phụ thuộc bắc cầu theo điều kiện prime attribute; BCNF yêu cầu mọi determinant không tầm thường là siêu khóa.

Đối chiếu phương án: B. A non-prime attribute is functionally dependent on a non-prime attribute.; C. A non-prime attribute is functionally dependent on the entire candidate key.; D. There is a transitive dependency between candidate keys.. Các lựa chọn này không đáp ứng định nghĩa/điều kiện được hỏi; khi ôn, hãy chỉ ra chính xác điểm lệch của từng lựa chọn thay vì học thuộc mỗi chữ cái.
\end{question}

\begin{question}{189}
% [ID gốc: 15.98]
% [Nguồn: DBI202 Order 15 Q98]
Consider the relation R(A, B, C, D) with functional dependencies A → B and C → D. Find the key of this relation.
\begin{choices}
\choice The key is column A
\choice The key includes two columns (A, C)
\choice The key is column A or column C
\choice The key is column C
\end{choices}
\textbf{Đáp án:} B
\textbf{Giải thích:} Đáp án theo dữ liệu đề: B — B. The key includes two columns (A, C).

Lời giải gốc: Từ tập phụ thuộc hàm, bao đóng của {A, C} sinh ra toàn bộ các thuộc tính {A, B, C, D}, do đó {A, C} là khóa của quan hệ.

Nền tảng kiến thức: Chuẩn hóa dùng phụ thuộc hàm và khóa để giảm dư thừa cũng như bất thường cập nhật. Tính bao đóng thuộc tính để tìm khóa; 2NF loại phụ thuộc bộ phận vào một phần khóa ghép, 3NF xử lý phụ thuộc không khóa kiểu bắc cầu, còn BCNF yêu cầu mọi vế trái không tầm thường là siêu khóa.

Phân tích: Với phụ thuộc X → Y, mọi hai bộ có cùng giá trị X phải có cùng giá trị Y. Để kiểm tra khóa, tính X⁺ bằng cách áp dụng lặp các phụ thuộc; nếu bao đóng chứa toàn bộ thuộc tính thì X là siêu khóa, sau đó kiểm tra tính tối thiểu.

Đối chiếu phương án: A. The key is column A; C. The key is column A or column C; D. The key is column C. Các lựa chọn này không đáp ứng định nghĩa/điều kiện được hỏi; khi ôn, hãy chỉ ra chính xác điểm lệch của từng lựa chọn thay vì học thuộc mỗi chữ cái.
\end{question}

\begin{question}{190}
% [ID gốc: 15.99]
% [Nguồn: DBI202 Order 15 Q99]
A table can be in 3NF but not in 1NF.
\begin{choices}
\choice True
\choice False
\end{choices}
\textbf{Đáp án:} B
\textbf{Giải thích:} Đáp án theo dữ liệu đề: B — B. False.

Lời giải gốc: Phát biểu này sai vì để đạt các dạng chuẩn cao hơn như 3NF, một bảng bắt buộc phải thỏa mãn các dạng chuẩn thấp hơn là 1NF và 2NF theo thứ tự phân cấp.

Nền tảng kiến thức: Chuẩn hóa dùng phụ thuộc hàm và khóa để giảm dư thừa cũng như bất thường cập nhật. Tính bao đóng thuộc tính để tìm khóa; 2NF loại phụ thuộc bộ phận vào một phần khóa ghép, 3NF xử lý phụ thuộc không khóa kiểu bắc cầu, còn BCNF yêu cầu mọi vế trái không tầm thường là siêu khóa.

Phân tích: Đánh giá dạng chuẩn theo thứ tự: 1NF yêu cầu giá trị nguyên tố/không lặp theo mô hình quan hệ; 2NF loại phụ thuộc bộ phận của thuộc tính không khóa vào một phần khóa ghép; 3NF kiểm soát phụ thuộc bắc cầu theo điều kiện prime attribute; BCNF yêu cầu mọi determinant không tầm thường là siêu khóa.

Đối chiếu phương án: A. True. Các lựa chọn này không đáp ứng định nghĩa/điều kiện được hỏi; khi ôn, hãy chỉ ra chính xác điểm lệch của từng lựa chọn thay vì học thuộc mỗi chữ cái.
\end{question}

\begin{question}{191}
% [ID gốc: 15.100]
% [Nguồn: DBI202 Order 15 Q100]
Which of the following violates 1NF (First Normal Form)?
\begin{choices}
\choice Having a NULL in a column
\choice Repeating groups of values in a column
\choice Having a primary key and all single attributes
\choice Not exist multi-value attributes
\end{choices}
\textbf{Đáp án:} B
\textbf{Giải thích:} Đáp án theo dữ liệu đề: B — B. Repeating groups of values in a column.

Lời giải gốc: Chuẩn 1NF yêu cầu mọi thuộc tính phải chứa các giá trị nguyên tố (Atomic values). Việc chứa nhóm giá trị lặp lại (Repeating groups) hoặc đa trị trong cùng một cột sẽ vi phạm 1NF.

Nền tảng kiến thức: Chuẩn hóa dùng phụ thuộc hàm và khóa để giảm dư thừa cũng như bất thường cập nhật. Tính bao đóng thuộc tính để tìm khóa; 2NF loại phụ thuộc bộ phận vào một phần khóa ghép, 3NF xử lý phụ thuộc không khóa kiểu bắc cầu, còn BCNF yêu cầu mọi vế trái không tầm thường là siêu khóa.

Phân tích: Đánh giá dạng chuẩn theo thứ tự: 1NF yêu cầu giá trị nguyên tố/không lặp theo mô hình quan hệ; 2NF loại phụ thuộc bộ phận của thuộc tính không khóa vào một phần khóa ghép; 3NF kiểm soát phụ thuộc bắc cầu theo điều kiện prime attribute; BCNF yêu cầu mọi determinant không tầm thường là siêu khóa.

Đối chiếu phương án: A. Having a NULL in a column; C. Having a primary key and all single attributes; D. Not exist multi-value attributes. Các lựa chọn này không đáp ứng định nghĩa/điều kiện được hỏi; khi ôn, hãy chỉ ra chính xác điểm lệch của từng lựa chọn thay vì học thuộc mỗi chữ cái.
\end{question}

\begin{question}{192}
% [ID gốc: 15.101]
% [Nguồn: DBI202 Order 15 Q101]
Given the relational schema R(A, B, C, D, E) and the set of functional dependencies: F = {AB -&gt; C, B -&gt; D, C -&gt; E, D -&gt; B}. What is the closure of the attribute set {A, B}?
\begin{choices}
\choice {AB, CD}
\choice {A, B, C, D, E}
\choice {AB, CE}
\choice {A, B, D}
\end{choices}
\textbf{Đáp án:} B
\textbf{Giải thích:} Đáp án theo dữ liệu đề: B — B. {A, B, C, D, E}.

Lời giải gốc: Tính bao đóng $\{A, B\}^+$: Ban đầu: $\{A, B\}$ Dùng $AB \rightarrow C \Rightarrow \{A, B, C\}$ Dùng $B \rightarrow D \Rightarrow \{A, B, C, D\}$ Dùng $C \rightarrow E \Rightarrow \{A, B, C, D, E\}$ Vậy bao đóng của $\{A, B\}$ là $\{A, B, C, D, E\}$.

Nền tảng kiến thức: Chuẩn hóa dùng phụ thuộc hàm và khóa để giảm dư thừa cũng như bất thường cập nhật. Tính bao đóng thuộc tính để tìm khóa; 2NF loại phụ thuộc bộ phận vào một phần khóa ghép, 3NF xử lý phụ thuộc không khóa kiểu bắc cầu, còn BCNF yêu cầu mọi vế trái không tầm thường là siêu khóa.

Phân tích: Với phụ thuộc X → Y, mọi hai bộ có cùng giá trị X phải có cùng giá trị Y. Để kiểm tra khóa, tính X⁺ bằng cách áp dụng lặp các phụ thuộc; nếu bao đóng chứa toàn bộ thuộc tính thì X là siêu khóa, sau đó kiểm tra tính tối thiểu.

Đối chiếu phương án: A. {AB, CD}; C. {AB, CE}; D. {A, B, D}. Các lựa chọn này không đáp ứng định nghĩa/điều kiện được hỏi; khi ôn, hãy chỉ ra chính xác điểm lệch của từng lựa chọn thay vì học thuộc mỗi chữ cái.
\end{question}

\begin{question}{193}
% [ID gốc: 15.102]
% [Nguồn: DBI202 Order 15 Q102]
In SQL, which of the following statements is correct?
\begin{choices}
\choice All relations in 2NF are also in BCNF
\choice All relations in 1NF are also in BCNF
\choice All relations in 3NF are also in BCNF
\choice For any relation schema, there is a dependency-preserving decomposition into 3NF
\end{choices}
\textbf{Đáp án:} D
\textbf{Giải thích:} Đáp án theo dữ liệu đề: D — D. For any relation schema, there is a dependency-preserving decomposition into 3NF.

Lời giải gốc: Đối với bất kỳ lược đồ quan hệ nào, luôn tồn tại phép phân rã bảo toàn phụ thuộc hàm đạt dạng chuẩn 3NF.

Nền tảng kiến thức: Chuẩn hóa dùng phụ thuộc hàm và khóa để giảm dư thừa cũng như bất thường cập nhật. Tính bao đóng thuộc tính để tìm khóa; 2NF loại phụ thuộc bộ phận vào một phần khóa ghép, 3NF xử lý phụ thuộc không khóa kiểu bắc cầu, còn BCNF yêu cầu mọi vế trái không tầm thường là siêu khóa.

Phân tích: Lược đồ quan hệ ghi tên quan hệ kèm danh sách thuộc tính, thường theo dạng `Boat(BoatID, BoatName, Color, Status)`. Nếu đề hỏi schema thuần túy, không nhất thiết phải kèm kiểu dữ liệu hay các bộ dữ liệu đang lưu.

Đối chiếu phương án: A. All relations in 2NF are also in BCNF; B. All relations in 1NF are also in BCNF; C. All relations in 3NF are also in BCNF. Các lựa chọn này không đáp ứng định nghĩa/điều kiện được hỏi; khi ôn, hãy chỉ ra chính xác điểm lệch của từng lựa chọn thay vì học thuộc mỗi chữ cái.
\end{question}

\begin{question}{194}
% [ID gốc: 15.104]
% [Nguồn: DBI202 Order 15 Q104]
Consider a relation R(A, B, C, D, E, F) with functional dependencies {AB → C, C → D, D → AE}. Which of the following is the closure of {B, D}?
\begin{choices}
\choice {ABCDE}
\choice {BCDE}
\choice {ABDEF}
\choice {ABCEF}
\choice {BCEF}
\end{choices}
\textbf{Đáp án:} A
\textbf{Giải thích:} Đáp án theo dữ liệu đề: A — A. {ABCDE}.

Lời giải gốc: Bao đóng của {B, D}: từ D → AE thêm A và E, nay có {A, B, D, E}; từ AB → C thêm C, được {A, B, C, D, E}. Kết quả là {ABCDE}.

Nền tảng kiến thức: Chuẩn hóa dùng phụ thuộc hàm và khóa để giảm dư thừa cũng như bất thường cập nhật. Tính bao đóng thuộc tính để tìm khóa; 2NF loại phụ thuộc bộ phận vào một phần khóa ghép, 3NF xử lý phụ thuộc không khóa kiểu bắc cầu, còn BCNF yêu cầu mọi vế trái không tầm thường là siêu khóa.

Phân tích: Với phụ thuộc X → Y, mọi hai bộ có cùng giá trị X phải có cùng giá trị Y. Để kiểm tra khóa, tính X⁺ bằng cách áp dụng lặp các phụ thuộc; nếu bao đóng chứa toàn bộ thuộc tính thì X là siêu khóa, sau đó kiểm tra tính tối thiểu.

Đối chiếu phương án: B. {BCDE}; C. {ABDEF}; D. {ABCEF}; E. {BCEF}. Các lựa chọn này không đáp ứng định nghĩa/điều kiện được hỏi; khi ôn, hãy chỉ ra chính xác điểm lệch của từng lựa chọn thay vì học thuộc mỗi chữ cái.
\end{question}

\begin{question}{195}
% [ID gốc: 15.106]
% [Nguồn: DBI202 Order 15 Q106]
A relation is in Third Normal Form (3NF) if:
\begin{choices}
\choice It is in 2NF and has only one attribute per tuple.
\choice It is in 2NF and has no transitive dependencies.
\choice It is in BCNF.
\choice All attributes are keys.
\end{choices}
\textbf{Đáp án:} B
\textbf{Giải thích:} Đáp án theo dữ liệu đề: B — B. It is in 2NF and has no transitive dependencies..

Lời giải gốc: Một quan hệ đạt chuẩn 3 (3NF) khi nó ở dạng chuẩn 2 và không chứa các phụ thuộc bắc cầu (transitive dependencies). ---

Nền tảng kiến thức: Chuẩn hóa dùng phụ thuộc hàm và khóa để giảm dư thừa cũng như bất thường cập nhật. Tính bao đóng thuộc tính để tìm khóa; 2NF loại phụ thuộc bộ phận vào một phần khóa ghép, 3NF xử lý phụ thuộc không khóa kiểu bắc cầu, còn BCNF yêu cầu mọi vế trái không tầm thường là siêu khóa.

Phân tích: Đánh giá dạng chuẩn theo thứ tự: 1NF yêu cầu giá trị nguyên tố/không lặp theo mô hình quan hệ; 2NF loại phụ thuộc bộ phận của thuộc tính không khóa vào một phần khóa ghép; 3NF kiểm soát phụ thuộc bắc cầu theo điều kiện prime attribute; BCNF yêu cầu mọi determinant không tầm thường là siêu khóa.

Đối chiếu phương án: A. It is in 2NF and has only one attribute per tuple.; C. It is in BCNF.; D. All attributes are keys.. Các lựa chọn này không đáp ứng định nghĩa/điều kiện được hỏi; khi ôn, hãy chỉ ra chính xác điểm lệch của từng lựa chọn thay vì học thuộc mỗi chữ cái.
\end{question}

\begin{question}{196}
% [ID gốc: 15.107]
% [Nguồn: DBI202 Order 15 Q107]
Given R as a relation, which of the following statements is correct about normalizing data to 3NF?
\begin{choices}
\choice R has a primary key, no repeating groups, no multivalued/composite attributes
\choice R has a primary key, no repeating groups, no multivalued/complex attributes, no partial functional dependencies
\choice R has a primary key, no repeating groups, no multivalued/composite attributes, no transitive dependencies
\choice R has a primary key, no repeating groups, no multivalued/complex attributes, no transitive functional dependencies, no partial functional dependencies
\choice R has complete functional dependency
\end{choices}
\textbf{Đáp án:} D
\textbf{Giải thích:} Đáp án theo dữ liệu đề: D — D. R has a primary key, no repeating groups, no multivalued/complex attributes, no transitive functional dependencies, no partial functional dependencies.

Lời giải gốc: Dạng chuẩn 3 (3NF) đòi hỏi quan hệ đạt chuẩn 2 (không có phụ thuộc hàm bộ phận) và đồng thời không chứa phụ thuộc hàm bắc cầu (transitive functional dependencies).

Nền tảng kiến thức: Chuẩn hóa dùng phụ thuộc hàm và khóa để giảm dư thừa cũng như bất thường cập nhật. Tính bao đóng thuộc tính để tìm khóa; 2NF loại phụ thuộc bộ phận vào một phần khóa ghép, 3NF xử lý phụ thuộc không khóa kiểu bắc cầu, còn BCNF yêu cầu mọi vế trái không tầm thường là siêu khóa.

Phân tích: Đánh giá dạng chuẩn theo thứ tự: 1NF yêu cầu giá trị nguyên tố/không lặp theo mô hình quan hệ; 2NF loại phụ thuộc bộ phận của thuộc tính không khóa vào một phần khóa ghép; 3NF kiểm soát phụ thuộc bắc cầu theo điều kiện prime attribute; BCNF yêu cầu mọi determinant không tầm thường là siêu khóa.

Đối chiếu phương án: A. R has a primary key, no repeating groups, no multivalued/composite attributes; B. R has a primary key, no repeating groups, no multivalued/complex attributes, no partial functi…; C. R has a primary key, no repeating groups, no multivalued/composite attributes, no transitive d…; E. R has complete functional dependency. Các lựa chọn này không đáp ứng định nghĩa/điều kiện được hỏi; khi ôn, hãy chỉ ra chính xác điểm lệch của từng lựa chọn thay vì học thuộc mỗi chữ cái.
\end{question}

\begin{question}{197}
% [ID gốc: 15.113]
% [Nguồn: DBI202 Order 15 Q113]
Consider the relation R(A, B, C, D) and functional dependencies: A -&gt; B, C -&gt; D. Find key of the above relation.
\begin{choices}
\choice The key is the A column
\choice The key includes two columns (A, C)
\choice The key is the A column or C column
\choice The key is the C column
\end{choices}
\textbf{Đáp án:} B
\textbf{Giải thích:} Đáp án theo dữ liệu đề: B — B. The key includes two columns (A, C).

Lời giải gốc: Tính bao đóng của $(AC)^+$: $AC \rightarrow AC$, từ $A \rightarrow B$ có $ABC$, từ $C \rightarrow D$ có $ABCD$. Do đó $\{A, C\}$ là khóa của quan hệ R.

Nền tảng kiến thức: Chuẩn hóa dùng phụ thuộc hàm và khóa để giảm dư thừa cũng như bất thường cập nhật. Tính bao đóng thuộc tính để tìm khóa; 2NF loại phụ thuộc bộ phận vào một phần khóa ghép, 3NF xử lý phụ thuộc không khóa kiểu bắc cầu, còn BCNF yêu cầu mọi vế trái không tầm thường là siêu khóa.

Phân tích: Với phụ thuộc X → Y, mọi hai bộ có cùng giá trị X phải có cùng giá trị Y. Để kiểm tra khóa, tính X⁺ bằng cách áp dụng lặp các phụ thuộc; nếu bao đóng chứa toàn bộ thuộc tính thì X là siêu khóa, sau đó kiểm tra tính tối thiểu.

Đối chiếu phương án: A. The key is the A column; C. The key is the A column or C column; D. The key is the C column. Các lựa chọn này không đáp ứng định nghĩa/điều kiện được hỏi; khi ôn, hãy chỉ ra chính xác điểm lệch của từng lựa chọn thay vì học thuộc mỗi chữ cái.
\end{question}

\begin{question}{198}
% [ID gốc: 15.114]
% [Nguồn: DBI202 Order 15 Q114]
Suppose that we have a relation: R(A,B,C,D) and FD : D-&gt;BC Which of the following is equivalent to the given dependency function?
\begin{choices}
\choice D-&gt;B and D-&gt;C
\choice C-&gt;B and C-&gt;B
\choice D-&gt;B and B-&gt;C
\choice C-&gt;D and C-&gt;B
\end{choices}
\textbf{Đáp án:} A
\textbf{Giải thích:} Đáp án theo dữ liệu đề: A — A. D-&gt;B and D-&gt;C.

Lời giải gốc: Theo quy tắc phân tách (Decomposition rule), phụ thuộc hàm $D \to BC$ tương đương với hai phụ thuộc hàm riêng biệt là $D \to B$ và $D \to C$. ---

Nền tảng kiến thức: Chuẩn hóa dùng phụ thuộc hàm và khóa để giảm dư thừa cũng như bất thường cập nhật. Tính bao đóng thuộc tính để tìm khóa; 2NF loại phụ thuộc bộ phận vào một phần khóa ghép, 3NF xử lý phụ thuộc không khóa kiểu bắc cầu, còn BCNF yêu cầu mọi vế trái không tầm thường là siêu khóa.

Phân tích: Với phụ thuộc X → Y, mọi hai bộ có cùng giá trị X phải có cùng giá trị Y. Để kiểm tra khóa, tính X⁺ bằng cách áp dụng lặp các phụ thuộc; nếu bao đóng chứa toàn bộ thuộc tính thì X là siêu khóa, sau đó kiểm tra tính tối thiểu.

Đối chiếu phương án: B. C-&gt;B and C-&gt;B; C. D-&gt;B and B-&gt;C; D. C-&gt;D and C-&gt;B. Các lựa chọn này không đáp ứng định nghĩa/điều kiện được hỏi; khi ôn, hãy chỉ ra chính xác điểm lệch của từng lựa chọn thay vì học thuộc mỗi chữ cái.
\end{question}

\begin{question}{199}
% [ID gốc: 15.115]
% [Nguồn: DBI202 Order 15 Q115]
To convert a table to First Normal Form (1NF), what step must be taken?
\begin{choices}
\choice Remove redundant data
\choice Ensure each attribute has a unique value
\choice Split the table into smaller tables
\choice Ensure each attribute contains atomic values
\end{choices}
\textbf{Đáp án:} D
\textbf{Giải thích:} Đáp án theo dữ liệu đề: D — D. Ensure each attribute contains atomic values.

Lời giải gốc: Để đạt dạng chuẩn 1 (1NF), mỗi thuộc tính trong bảng phải chứa các giá trị nguyên tố (atomic values), không chứa nhóm lặp hay thuộc tính đa trị.

Nền tảng kiến thức: Chuẩn hóa dùng phụ thuộc hàm và khóa để giảm dư thừa cũng như bất thường cập nhật. Tính bao đóng thuộc tính để tìm khóa; 2NF loại phụ thuộc bộ phận vào một phần khóa ghép, 3NF xử lý phụ thuộc không khóa kiểu bắc cầu, còn BCNF yêu cầu mọi vế trái không tầm thường là siêu khóa.

Phân tích: Đánh giá dạng chuẩn theo thứ tự: 1NF yêu cầu giá trị nguyên tố/không lặp theo mô hình quan hệ; 2NF loại phụ thuộc bộ phận của thuộc tính không khóa vào một phần khóa ghép; 3NF kiểm soát phụ thuộc bắc cầu theo điều kiện prime attribute; BCNF yêu cầu mọi determinant không tầm thường là siêu khóa.

Đối chiếu phương án: A. Remove redundant data; B. Ensure each attribute has a unique value; C. Split the table into smaller tables. Các lựa chọn này không đáp ứng định nghĩa/điều kiện được hỏi; khi ôn, hãy chỉ ra chính xác điểm lệch của từng lựa chọn thay vì học thuộc mỗi chữ cái.
\end{question}

\begin{question}{200}
% [ID gốc: 15.116]
% [Nguồn: DBI202 Order 15 Q116]
In SQL normalization theory, which of the following statements is correct?
\begin{choices}
\choice All relations in 2NF are also in BCNF.
\choice All relations in 1NF are also in BCNF.
\choice All relations in 3NF are also in BCNF.
\choice For any relation schema, there is a dependency-preserving decomposition into 3NF.
\end{choices}
\textbf{Đáp án:} D
\textbf{Giải thích:} Đáp án theo dữ liệu đề: D — D. For any relation schema, there is a dependency-preserving decomposition into 3NF..

Lời giải gốc: Mọi lược đồ quan hệ đều luôn có thể phân rã thành dạng chuẩn 3 (3NF) mà vẫn bảo toàn được các phụ thuộc hàm. Điều này không luôn đúng với BCNF.

Nền tảng kiến thức: Chuẩn hóa dùng phụ thuộc hàm và khóa để giảm dư thừa cũng như bất thường cập nhật. Tính bao đóng thuộc tính để tìm khóa; 2NF loại phụ thuộc bộ phận vào một phần khóa ghép, 3NF xử lý phụ thuộc không khóa kiểu bắc cầu, còn BCNF yêu cầu mọi vế trái không tầm thường là siêu khóa.

Phân tích: Lược đồ quan hệ ghi tên quan hệ kèm danh sách thuộc tính, thường theo dạng `Boat(BoatID, BoatName, Color, Status)`. Nếu đề hỏi schema thuần túy, không nhất thiết phải kèm kiểu dữ liệu hay các bộ dữ liệu đang lưu.

Đối chiếu phương án: A. All relations in 2NF are also in BCNF.; B. All relations in 1NF are also in BCNF.; C. All relations in 3NF are also in BCNF.. Các lựa chọn này không đáp ứng định nghĩa/điều kiện được hỏi; khi ôn, hãy chỉ ra chính xác điểm lệch của từng lựa chọn thay vì học thuộc mỗi chữ cái.
\end{question}

\begin{question}{201}
% [ID gốc: 15.117]
% [Nguồn: DBI202 Order 15 Q117]
Consider the relation R(A, B, C, D, E, F, G, K) with functional dependencies AB → C, B → C, C → AD, D → E, CF → B, G → K. What is the closure of {AB}?
\begin{choices}
\choice {ABCDE}
\choice {ABCDEF}
\choice {CDE}
\choice {ABCD}
\choice {ABCFG}
\end{choices}
\textbf{Đáp án:} A
\textbf{Giải thích:} Đáp án theo dữ liệu đề: A — A. {ABCDE}.

Lời giải gốc: Bao đóng của {AB}: AB → C thêm C; C → AD thêm D; D → E thêm E. Kết quả là {A, B, C, D, E}, tức {ABCDE}.

Nền tảng kiến thức: Chuẩn hóa dùng phụ thuộc hàm và khóa để giảm dư thừa cũng như bất thường cập nhật. Tính bao đóng thuộc tính để tìm khóa; 2NF loại phụ thuộc bộ phận vào một phần khóa ghép, 3NF xử lý phụ thuộc không khóa kiểu bắc cầu, còn BCNF yêu cầu mọi vế trái không tầm thường là siêu khóa.

Phân tích: Với phụ thuộc X → Y, mọi hai bộ có cùng giá trị X phải có cùng giá trị Y. Để kiểm tra khóa, tính X⁺ bằng cách áp dụng lặp các phụ thuộc; nếu bao đóng chứa toàn bộ thuộc tính thì X là siêu khóa, sau đó kiểm tra tính tối thiểu.

Đối chiếu phương án: B. {ABCDEF}; C. {CDE}; D. {ABCD}; E. {ABCFG}. Các lựa chọn này không đáp ứng định nghĩa/điều kiện được hỏi; khi ôn, hãy chỉ ra chính xác điểm lệch của từng lựa chọn thay vì học thuộc mỗi chữ cái.
\end{question}

\begin{question}{202}
% [ID gốc: 16.1]
% [Nguồn: DBI202 Order 16 Q1]
What is "de-normalization"?
\begin{choices}
\choice De-normalization means allowing redundancy in a table
\choice De-normalization means allowing duplicate columns appeared in a table
\choice De-normalization means allowing a table to have more than 3 columns
\choice De-normalization means allowing a table to have more than 1024 columns
\end{choices}
\textbf{Đáp án:} A
\textbf{Giải thích:} Đáp án theo dữ liệu đề: A — A. De-normalization means allowing redundancy in a table.

Nền tảng kiến thức: Thiết kế vật lý chọn chỉ mục, cách tổ chức tệp và các tối ưu theo workload thực tế. Denormalization hoặc thêm chỉ mục có thể giảm chi phí đọc nhưng đổi lại tăng lưu trữ, chi phí ghi và nguy cơ phải đồng bộ dữ liệu dư thừa.

Phân tích: Denormalization chủ động chấp nhận một phần dư thừa hoặc lưu trước kết quả để giảm join/tính toán ở workload đọc. Nó có thể tăng tốc truy xuất, nhưng làm tăng dung lượng và buộc hệ thống cập nhật đồng bộ các bản sao; đây là quyết định theo workload, không phải quy tắc luôn có lợi.

Đối chiếu phương án: B. De-normalization means allowing duplicate columns appeared in a table; C. De-normalization means allowing a table to have more than 3 columns; D. De-normalization means allowing a table to have more than 1024 columns. Các lựa chọn này không đáp ứng định nghĩa/điều kiện được hỏi; khi ôn, hãy chỉ ra chính xác điểm lệch của từng lựa chọn thay vì học thuộc mỗi chữ cái.
\end{question}

\begin{question}{203}
% [ID gốc: 16.2]
% [Nguồn: DBI202 Order 16 Q2]
What is the benefit of "de-normalization"?
\begin{choices}
\choice "de-normalization" has no benefit
\choice The main benefit of de-normalization is improved performance with simplified data retrieval (this is done by reduction in the number of joins needed for data processing)
\choice The main benefit of de-normalization is eliminating redundant information from a table and organizing the data so that future changes to the table are easier
\end{choices}
\textbf{Đáp án:} B
\textbf{Giải thích:} Đáp án theo dữ liệu đề: B — B. The main benefit of de-normalization is improved performance with simplified data retrieval (this is done by reduction in the number of joins needed for data processing).

Nền tảng kiến thức: Thiết kế vật lý chọn chỉ mục, cách tổ chức tệp và các tối ưu theo workload thực tế. Denormalization hoặc thêm chỉ mục có thể giảm chi phí đọc nhưng đổi lại tăng lưu trữ, chi phí ghi và nguy cơ phải đồng bộ dữ liệu dư thừa.

Phân tích: Denormalization chủ động chấp nhận một phần dư thừa hoặc lưu trước kết quả để giảm join/tính toán ở workload đọc. Nó có thể tăng tốc truy xuất, nhưng làm tăng dung lượng và buộc hệ thống cập nhật đồng bộ các bản sao; đây là quyết định theo workload, không phải quy tắc luôn có lợi.

Đối chiếu phương án: A. "de-normalization" has no benefit; C. The main benefit of de-normalization is eliminating redundant information from a table and org…. Các lựa chọn này không đáp ứng định nghĩa/điều kiện được hỏi; khi ôn, hãy chỉ ra chính xác điểm lệch của từng lựa chọn thay vì học thuộc mỗi chữ cái.
\end{question}


\newpage
\section*{Chương 4: High-Level Database Models}

\begin{question}{204}
% [ID gốc: 1.19]
% [Nguồn: DBI202 Order 1 Q19]
Which of the following is not a data model in the database?
\begin{choices}
\choice Conceptual data model
\choice Object-Oriented data model
\choice Network data model
\choice Hierarchical data model
\end{choices}
\textbf{Đáp án:} A
\textbf{Giải thích:} Đáp án theo dữ liệu đề: A — A. Conceptual data model.

Nền tảng kiến thức: Chương này phân biệt dữ liệu, cơ sở dữ liệu, hệ quản trị CSDL (DBMS), mô hình dữ liệu, lược đồ và trạng thái dữ liệu. Cơ sở dữ liệu lưu dữ liệu có liên hệ; DBMS là phần mềm định nghĩa, truy vấn, cập nhật và bảo vệ dữ liệu đó.

Phân tích: Data model là bộ khái niệm dùng để mô tả cấu trúc dữ liệu, quan hệ giữa dữ liệu và các ràng buộc. Relational, hierarchical, network và graph là các mô hình có thật; cần chọn mô hình khớp đặc điểm được hỏi, không nhầm mô hình với DBMS hay ngôn ngữ thao tác.

Đối chiếu phương án: B. Object-Oriented data model; C. Network data model; D. Hierarchical data model. Các lựa chọn này không đáp ứng định nghĩa/điều kiện được hỏi; khi ôn, hãy chỉ ra chính xác điểm lệch của từng lựa chọn thay vì học thuộc mỗi chữ cái.
\end{question}

\begin{question}{205}
% [ID gốc: 1.20]
% [Nguồn: DBI202 Order 1 Q20]
The conceptual model is:
\begin{choices}
\choice Independent of both hardware and software
\choice Dependent on both hardware and software.
\choice Dependent on software.
\choice Dependent on hardware.
\end{choices}
\textbf{Đáp án:} A
\textbf{Giải thích:} Đáp án theo dữ liệu đề: A — A. Independent of both hardware and software.

Nền tảng kiến thức: Chương này phân biệt dữ liệu, cơ sở dữ liệu, hệ quản trị CSDL (DBMS), mô hình dữ liệu, lược đồ và trạng thái dữ liệu. Cơ sở dữ liệu lưu dữ liệu có liên hệ; DBMS là phần mềm định nghĩa, truy vấn, cập nhật và bảo vệ dữ liệu đó.

Phân tích: Mô hình khái niệm mô tả dữ liệu và quy tắc nghiệp vụ ở mức độc lập với chi tiết cài đặt cụ thể. Vì vậy, nó được dùng để trao đổi yêu cầu trước khi quyết định DBMS, cấu trúc tệp hay cấu hình phần cứng.

Đối chiếu phương án: B. Dependent on both hardware and software.; C. Dependent on software.; D. Dependent on hardware.. Các lựa chọn này không đáp ứng định nghĩa/điều kiện được hỏi; khi ôn, hãy chỉ ra chính xác điểm lệch của từng lựa chọn thay vì học thuộc mỗi chữ cái.
\end{question}

\begin{question}{206}
% [ID gốc: 1.51]
% [Nguồn: DBI202 Order 1 Q51]
Which of the following is expressed by an E-R diagram?
\begin{choices}
\choice Relation between process and relationship
\choice Relation between entity and process
\choice Relation between processes
\choice Relation between entities
\end{choices}
\textbf{Đáp án:} D
\textbf{Giải thích:} Đáp án theo dữ liệu đề: D — D. Relation between entities.

Lời giải gốc: Biểu đồ ER tập trung biểu diễn các dữ liệu tĩnh mang tính cấu trúc, mô tả mối quan hệ giữa các thực thể dữ liệu với nhau (Relation between entities), không mô tả các quy trình xử lý động (processes).

Nền tảng kiến thức: Chương này phân biệt dữ liệu, cơ sở dữ liệu, hệ quản trị CSDL (DBMS), mô hình dữ liệu, lược đồ và trạng thái dữ liệu. Cơ sở dữ liệu lưu dữ liệu có liên hệ; DBMS là phần mềm định nghĩa, truy vấn, cập nhật và bảo vệ dữ liệu đó.

Phân tích: ER diagram mô tả các entity set, thuộc tính và relationship set, cùng các ràng buộc cần thiết. Nó tập trung vào đối tượng dữ liệu và liên kết giữa chúng, không mô tả luồng xử lý/chương trình như sơ đồ quy trình.

Đối chiếu phương án: A. Relation between process and relationship; B. Relation between entity and process; C. Relation between processes. Các lựa chọn này không đáp ứng định nghĩa/điều kiện được hỏi; khi ôn, hãy chỉ ra chính xác điểm lệch của từng lựa chọn thay vì học thuộc mỗi chữ cái.
\end{question}

\begin{question}{207}
% [ID gốc: 2.1]
% [Nguồn: DBI202 Order 2 Q1]
A weak entity set is an entity set that does not have a sufficient attribute to form a primary key. The key for a weak entity set E is formed by the combination of zero or more attributes of E and key attributes from supporting entity sets. This means that the key for a weak entity set is partially derived from the key of another entity set that it is related to through an identifying relationship.
QN=18 (8022) The key for a weak entity set E is .........
\begin{choices}
\choice Zero or more attributes of E
\choice The set of attributes of supporting relationships for E
\choice The set of attributes of supporting entity sets
\choice Zero or more attributes of E and key attributes from supporting entity sets
\end{choices}
\textbf{Đáp án:} D
\textbf{Giải thích:} Đáp án theo dữ liệu đề: D — D. Zero or more attributes of E and key attributes from supporting entity sets.

Nền tảng kiến thức: Mô hình ER mô tả yêu cầu ở mức khái niệm: thực thể là đối tượng cần lưu, thuộc tính mô tả thực thể hoặc mối quan hệ, còn lực lượng (cardinality) và participation diễn tả các ràng buộc kết hợp. Khóa của thực thể yếu thường ghép khóa bộ phận với khóa của thực thể chủ.

Phân tích: Ba thành phần nền tảng của ER là entity set, attribute và relationship set. Entity mô tả loại đối tượng, attribute mô tả dữ liệu của đối tượng/liên kết, relationship biểu diễn kết nối giữa entity set.

Đối chiếu phương án: A. Zero or more attributes of E; B. The set of attributes of supporting relationships for E; C. The set of attributes of supporting entity sets. Các lựa chọn này không đáp ứng định nghĩa/điều kiện được hỏi; khi ôn, hãy chỉ ra chính xác điểm lệch của từng lựa chọn thay vì học thuộc mỗi chữ cái.
\end{question}

\begin{question}{208}
% [ID gốc: 2.2]
% [Nguồn: DBI202 Order 2 Q2]
A class in UML is similar to...........
\begin{choices}
\choice An entity set in E/R model
\choice An attribute in E/R model
\choice A Relationship in E/R model
\choice None of the others
\end{choices}
\textbf{Đáp án:} A
\textbf{Giải thích:} Đáp án theo dữ liệu đề: A — A. An entity set in E/R model.

Nền tảng kiến thức: Mô hình ER mô tả yêu cầu ở mức khái niệm: thực thể là đối tượng cần lưu, thuộc tính mô tả thực thể hoặc mối quan hệ, còn lực lượng (cardinality) và participation diễn tả các ràng buộc kết hợp. Khóa của thực thể yếu thường ghép khóa bộ phận với khóa của thực thể chủ.

Phân tích: ER diagram mô tả các entity set, thuộc tính và relationship set, cùng các ràng buộc cần thiết. Nó tập trung vào đối tượng dữ liệu và liên kết giữa chúng, không mô tả luồng xử lý/chương trình như sơ đồ quy trình.

Đối chiếu phương án: B. An attribute in E/R model; C. A Relationship in E/R model; D. None of the others. Các lựa chọn này không đáp ứng định nghĩa/điều kiện được hỏi; khi ôn, hãy chỉ ra chính xác điểm lệch của từng lựa chọn thay vì học thuộc mỗi chữ cái.
\end{question}

\begin{question}{209}
% [ID gốc: 2.3]
% [Nguồn: DBI202 Order 2 Q3]
A(n) _____ provides a means to analyze business requirements so as to standardize organizational vocabulary, enforce business rules, and ensure adequate data quality
\begin{choices}
\choice Entity Relationship Diagram
\choice Relation Schema Design
\choice All of the others
\choice None of the others
\end{choices}
\textbf{Đáp án:} A
\textbf{Giải thích:} Đáp án theo dữ liệu đề: A — A. Entity Relationship Diagram.

Nền tảng kiến thức: Mô hình ER mô tả yêu cầu ở mức khái niệm: thực thể là đối tượng cần lưu, thuộc tính mô tả thực thể hoặc mối quan hệ, còn lực lượng (cardinality) và participation diễn tả các ràng buộc kết hợp. Khóa của thực thể yếu thường ghép khóa bộ phận với khóa của thực thể chủ.

Phân tích: Với lựa chọn tổng hợp, không chọn theo cảm tính: kiểm tra từng phát biểu độc lập. “All of the others/above” chỉ đúng khi mọi phát biểu còn lại đúng; “None” chỉ đúng khi không phát biểu nào trong nhóm thỏa yêu cầu. Đáp án lưu cần được đọc cùng các mệnh đề thành phần.

Đối chiếu phương án: B. Relation Schema Design; C. All of the others; D. None of the others. Các lựa chọn này không đáp ứng định nghĩa/điều kiện được hỏi; khi ôn, hãy chỉ ra chính xác điểm lệch của từng lựa chọn thay vì học thuộc mỗi chữ cái.
\end{question}

\begin{question}{210}
% [ID gốc: 2.4]
% [Nguồn: DBI202 Order 2 Q4]
An association class in UML is similar to ______ in the ER model
\begin{choices}
\choice attributes on a relationship
\choice Attributes
\choice Entities
\choice entity sets
\end{choices}
\textbf{Đáp án:} A
\textbf{Giải thích:} Đáp án theo dữ liệu đề: A — A. attributes on a relationship.

Nền tảng kiến thức: Mô hình ER mô tả yêu cầu ở mức khái niệm: thực thể là đối tượng cần lưu, thuộc tính mô tả thực thể hoặc mối quan hệ, còn lực lượng (cardinality) và participation diễn tả các ràng buộc kết hợp. Khóa của thực thể yếu thường ghép khóa bộ phận với khóa của thực thể chủ.

Phân tích: ER diagram mô tả các entity set, thuộc tính và relationship set, cùng các ràng buộc cần thiết. Nó tập trung vào đối tượng dữ liệu và liên kết giữa chúng, không mô tả luồng xử lý/chương trình như sơ đồ quy trình.

Đối chiếu phương án: B. Attributes; C. Entities; D. entity sets. Các lựa chọn này không đáp ứng định nghĩa/điều kiện được hỏi; khi ôn, hãy chỉ ra chính xác điểm lệch của từng lựa chọn thay vì học thuộc mỗi chữ cái.
\end{question}

\begin{question}{211}
% [ID gốc: 2.5]
% [Nguồn: DBI202 Order 2 Q5]
In UML, what is the difference between an aggregation and a composition?
\begin{choices}
\choice In composition, when the owning object is destroyed, so are the contained objects. In aggregation, this is not necessarily true.
\choice In aggregation, when the owning object is destroyed, so are the contained objects. In composition, this is not necessarily true.
\choice All of the others
\choice There is no difference between an aggregation and a composition
\end{choices}
\textbf{Đáp án:} A
\textbf{Giải thích:} Đáp án theo dữ liệu đề: A — A. In composition, when the owning object is destroyed, so are the contained objects. In aggregation, this is not necessarily true..

Nền tảng kiến thức: Mô hình ER mô tả yêu cầu ở mức khái niệm: thực thể là đối tượng cần lưu, thuộc tính mô tả thực thể hoặc mối quan hệ, còn lực lượng (cardinality) và participation diễn tả các ràng buộc kết hợp. Khóa của thực thể yếu thường ghép khóa bộ phận với khóa của thực thể chủ.

Phân tích: Phép tập hợp trên quan hệ cần so sánh các tuple theo thuộc tính tương ứng; union lấy hợp, intersection lấy phần chung và difference lấy phần có ở vế trái mà không có ở vế phải. Trước tiên kiểm tra tính tương thích schema, sau đó áp phép theo đúng hướng/toán tử.

Đối chiếu phương án: B. In aggregation, when the owning object is destroyed, so are the contained objects. In composit…; C. All of the others; D. There is no difference between an aggregation and a composition. Các lựa chọn này không đáp ứng định nghĩa/điều kiện được hỏi; khi ôn, hãy chỉ ra chính xác điểm lệch của từng lựa chọn thay vì học thuộc mỗi chữ cái.
\end{question}

\begin{question}{212}
% [ID gốc: 2.6]
% [Nguồn: DBI202 Order 2 Q6]
Consider the following statement:
"Entity Sets and their attributes should reflect the reality"
That above statement describes which principle?
\begin{choices}
\choice Avoid Redundancy Principle
\choice Simplicity Count Principle
\choice Limit the use of weak entity sets principle
\choice Faithfulness principle
\end{choices}
\textbf{Đáp án:} D
\textbf{Giải thích:} Đáp án theo dữ liệu đề: D — D. Faithfulness principle.

Nền tảng kiến thức: Mô hình ER mô tả yêu cầu ở mức khái niệm: thực thể là đối tượng cần lưu, thuộc tính mô tả thực thể hoặc mối quan hệ, còn lực lượng (cardinality) và participation diễn tả các ràng buộc kết hợp. Khóa của thực thể yếu thường ghép khóa bộ phận với khóa của thực thể chủ.

Phân tích: Thực thể yếu chưa có khóa độc lập; phải ghép khóa bộ phận của nó với khóa của thực thể chủ qua quan hệ định danh. Vì vậy, chỉ lấy thuộc tính của chủ hoặc chỉ lấy khóa bộ phận đều chưa nhận diện duy nhất được nó.

Đối chiếu phương án: A. Avoid Redundancy Principle; B. Simplicity Count Principle; C. Limit the use of weak entity sets principle. Các lựa chọn này không đáp ứng định nghĩa/điều kiện được hỏi; khi ôn, hãy chỉ ra chính xác điểm lệch của từng lựa chọn thay vì học thuộc mỗi chữ cái.
\end{question}

\begin{question}{213}
% [ID gốc: 2.7]
% [Nguồn: DBI202 Order 2 Q7]
Statement (a) is correct. It's important to pick the right kind of elements when designing a database to ensure that it accurately represents the data and supports efficient querying and updating.
Statement (b) is also correct. When an entity has no non-primary key attributes (i.e., it has only one attribute which is also its primary key), it can be converted into an attribute of another entity or relationship. This can simplify the schema and reduce redundancy.
Statement © is incorrect. Weak entities should be used when an entity cannot be identified by its own attributes alone and depends on another entity for its existence. However, using weak entities whenever possible is not a good design principle.
Statement (d) is incorrect. Redundancies in a database can lead to inconsistencies and should generally be avoided.
QN=25 (8010) Look at the following statements:
(a)We should pick the right kind of elements
(b) When an entity has no non-primary key attribute, we should convert that entity to an attribute
(c) Use weak entities when-ever possible
(d) Always make redundancies
\begin{choices}
\choice (a) and (b) are correct
\choice (a) and (c) are correct
\choice (a) and (d) are correct
\choice (b) and (c) are correct
\end{choices}
\textbf{Đáp án:} A
\textbf{Giải thích:} Đáp án theo dữ liệu đề: A, B — A. (a) and (b) are correct; B. (a) and (c) are correct.

Nền tảng kiến thức: Mô hình ER mô tả yêu cầu ở mức khái niệm: thực thể là đối tượng cần lưu, thuộc tính mô tả thực thể hoặc mối quan hệ, còn lực lượng (cardinality) và participation diễn tả các ràng buộc kết hợp. Khóa của thực thể yếu thường ghép khóa bộ phận với khóa của thực thể chủ.

Phân tích: Thực thể yếu chưa có khóa độc lập; phải ghép khóa bộ phận của nó với khóa của thực thể chủ qua quan hệ định danh. Vì vậy, chỉ lấy thuộc tính của chủ hoặc chỉ lấy khóa bộ phận đều chưa nhận diện duy nhất được nó.

Đối chiếu phương án: C. (a) and (d) are correct; D. (b) and (c) are correct. Các lựa chọn này không đáp ứng định nghĩa/điều kiện được hỏi; khi ôn, hãy chỉ ra chính xác điểm lệch của từng lựa chọn thay vì học thuộc mỗi chữ cái.
\end{question}

\begin{question}{214}
% [ID gốc: 2.8]
% [Nguồn: DBI202 Order 2 Q8]
The ER Diagram uses three principle element types:
\begin{choices}
\choice Entity sets, Attributes and Constraints
\choice Entity sets, Attributes, and Relationships
\choice Attributes, Constraints, and Relationships
\choice Entity sets, Constraints, and Relationships
\end{choices}
\textbf{Đáp án:} B
\textbf{Giải thích:} Đáp án theo dữ liệu đề: B — B. Entity sets, Attributes, and Relationships.

Nền tảng kiến thức: Mô hình ER mô tả yêu cầu ở mức khái niệm: thực thể là đối tượng cần lưu, thuộc tính mô tả thực thể hoặc mối quan hệ, còn lực lượng (cardinality) và participation diễn tả các ràng buộc kết hợp. Khóa của thực thể yếu thường ghép khóa bộ phận với khóa của thực thể chủ.

Phân tích: ER diagram mô tả các entity set, thuộc tính và relationship set, cùng các ràng buộc cần thiết. Nó tập trung vào đối tượng dữ liệu và liên kết giữa chúng, không mô tả luồng xử lý/chương trình như sơ đồ quy trình.

Đối chiếu phương án: A. Entity sets, Attributes and Constraints; C. Attributes, Constraints, and Relationships; D. Entity sets, Constraints, and Relationships. Các lựa chọn này không đáp ứng định nghĩa/điều kiện được hỏi; khi ôn, hãy chỉ ra chính xác điểm lệch của từng lựa chọn thay vì học thuộc mỗi chữ cái.
\end{question}

\begin{question}{215}
% [ID gốc: 2.9]
% [Nguồn: DBI202 Order 2 Q9]
Class = Entity set
Association = Binary relationship
Association class = Attributes on a relationship
Subclass = Isa hierarchy
Aggregation = Many-one relationship
Composition = Many-one relationship with referential integrity
QN=27 (8030) The binary relationship between classes in UML is called ...
\begin{choices}
\choice Aggregation
\choice Composition
\choice Association
\choice Relation
\end{choices}
\textbf{Đáp án:} C
\textbf{Giải thích:} Đáp án theo dữ liệu đề: C — C. Association.

Nền tảng kiến thức: Mô hình ER mô tả yêu cầu ở mức khái niệm: thực thể là đối tượng cần lưu, thuộc tính mô tả thực thể hoặc mối quan hệ, còn lực lượng (cardinality) và participation diễn tả các ràng buộc kết hợp. Khóa của thực thể yếu thường ghép khóa bộ phận với khóa của thực thể chủ.

Phân tích: Binary relationship/association nối hai kiểu thực thể/lớp; n-ary relationship có từ ba kiểu tham gia trở lên. Phân loại theo số entity set/class trong chính liên kết, không theo số thuộc tính của chúng.

Đối chiếu phương án: A. Aggregation; B. Composition; D. Relation. Các lựa chọn này không đáp ứng định nghĩa/điều kiện được hỏi; khi ôn, hãy chỉ ra chính xác điểm lệch của từng lựa chọn thay vì học thuộc mỗi chữ cái.
\end{question}

\begin{question}{216}
% [ID gốc: 2.10]
% [Nguồn: DBI202 Order 2 Q10]
In UML, a _____ between classes is called an association
\begin{choices}
\choice binary relationship
\choice multi-way relationship
\choice none of the others
\end{choices}
\textbf{Đáp án:} A
\textbf{Giải thích:} A binary relationship between classes in UML is called an association.
\end{question}

\begin{question}{217}
% [ID gốc: 2.11]
% [Nguồn: DBI202 Order 2 Q11]
When an entity is called "weak"?
(a) When it is used to represent weak objects, such as patients, elders
(b) When all of its attributes can not identify itself
(c) When It always needs support from other entity
(d) When It is not really necessary in the ERD
\begin{choices}
\choice (a) and (b) are correct
\choice (b) and (c) are correct
\choice (c) and (d) are correct
\choice (d) and (a) are correct
\end{choices}
\textbf{Đáp án:} B
\textbf{Giải thích:} A weak entity is one that cannot be identified by its own attributes alone and always needs support from an identifying owner entity (b and c).
\end{question}

\begin{question}{218}
% [ID gốc: 2.12]
% [Nguồn: DBI202 Order 2 Q12]
Consider the following statement:
"When drawing ERD, if an entity has no non-key attribute and it is the 'one' in many-one relationship, then we should set it to the attribute of other entities"
That above statement describes which principle?
\begin{choices}
\choice Avoid Redundancy Principle
\choice Picking the right kind of element principle
\choice Limit the use of weak entity sets principle
\choice Faithfulness principle
\end{choices}
\textbf{Đáp án:} B
\textbf{Giải thích:} This describes Picking the right kind of element principle.
\end{question}

\begin{question}{219}
% [ID gốc: 2.13]
% [Nguồn: DBI202 Order 2 Q13]
When drawing an ERD:
(a) Entity sets and their attributes should reflect reality
(b) Redundancy is not important, so we can ignore it
(c) Avoid introducing more elements into the design than is absolutely necessary
(d) Choosing the right relationships
\begin{choices}
\choice (a) and (b) and (c) are correct
\choice (a) and (c) and (d) are correct
\choice (a) and (b) and (d) are correct
\choice (b) and (c) and (d) are correct
\end{choices}
\textbf{Đáp án:} B
\textbf{Giải thích:} Redundancy is critical and must be avoided (b is false). Principles a, c, d are correct.
\end{question}

\begin{question}{220}
% [ID gốc: 2.14]
% [Nguồn: DBI202 Order 2 Q14]
A weak entity
\begin{choices}
\choice does not depend on other entities
\choice does not have enough key attribute(s)
\choice is an entity that has not any key-attribute
\choice None of the others
\end{choices}
\textbf{Đáp án:} B
\textbf{Giải thích:} A weak entity does not possess sufficient key attributes to form a primary key on its own.
\end{question}

\begin{question}{221}
% [ID gốc: 2.15]
% [Nguồn: DBI202 Order 2 Q15]
Which of the followings is true?
\begin{choices}
\choice The Entity Relationship (ER) model represents the structure of data graphically
\choice The ER model is a low level database design
IC]
The ER model represents the operation on data
\choice All of the others
\end{choices}
\textbf{Đáp án:} A
\textbf{Giải thích:} Đáp án theo dữ liệu đề: A — A. The Entity Relationship (ER) model represents the structure of data graphically.

Nền tảng kiến thức: Mô hình ER mô tả yêu cầu ở mức khái niệm: thực thể là đối tượng cần lưu, thuộc tính mô tả thực thể hoặc mối quan hệ, còn lực lượng (cardinality) và participation diễn tả các ràng buộc kết hợp. Khóa của thực thể yếu thường ghép khóa bộ phận với khóa của thực thể chủ.

Phân tích: ER diagram mô tả các entity set, thuộc tính và relationship set, cùng các ràng buộc cần thiết. Nó tập trung vào đối tượng dữ liệu và liên kết giữa chúng, không mô tả luồng xử lý/chương trình như sơ đồ quy trình.

Đối chiếu phương án: B. The ER model is a low level database design IC] The ER model represents the operation on data; D. All of the others. Các lựa chọn này không đáp ứng định nghĩa/điều kiện được hỏi; khi ôn, hãy chỉ ra chính xác điểm lệch của từng lựa chọn thay vì học thuộc mỗi chữ cái.
\end{question}

\begin{question}{222}
% [ID gốc: 2.17]
% [Nguồn: DBI202 Order 2 Q17]
5. The ER model is meant to
\begin{choices}
\choice replace relational design
\choice enable detailed descriptions of data query processing
\choice enable low level descriptions of data
\choice be close to a users perception of the data
\end{choices}
\textbf{Đáp án:} D
\textbf{Giải thích:} Đáp án theo dữ liệu đề: D — D. be close to a users perception of the data.

Nền tảng kiến thức: Mô hình ER mô tả yêu cầu ở mức khái niệm: thực thể là đối tượng cần lưu, thuộc tính mô tả thực thể hoặc mối quan hệ, còn lực lượng (cardinality) và participation diễn tả các ràng buộc kết hợp. Khóa của thực thể yếu thường ghép khóa bộ phận với khóa của thực thể chủ.

Phân tích: ER diagram mô tả các entity set, thuộc tính và relationship set, cùng các ràng buộc cần thiết. Nó tập trung vào đối tượng dữ liệu và liên kết giữa chúng, không mô tả luồng xử lý/chương trình như sơ đồ quy trình.

Đối chiếu phương án: A. replace relational design; B. enable detailed descriptions of data query processing; C. enable low level descriptions of data. Các lựa chọn này không đáp ứng định nghĩa/điều kiện được hỏi; khi ôn, hãy chỉ ra chính xác điểm lệch của từng lựa chọn thay vì học thuộc mỗi chữ cái.
\end{question}

\begin{question}{223}
% [ID gốc: 2.18]
% [Nguồn: DBI202 Order 2 Q18]
In databse design precesses, data requirements are expressed through......
(a) Schema Design
(b)ERD
(c)UML
(d) Table Design
\begin{choices}
\choice (c) and (d)
\choice (b) and (c)
\choice (a) and (b)
\choice (d) and (a)
\end{choices}
\textbf{Đáp án:} B
\textbf{Giải thích:} Đáp án theo dữ liệu đề: B, C — B. (b) and (c); C. (a) and (b).

Nền tảng kiến thức: Mô hình ER mô tả yêu cầu ở mức khái niệm: thực thể là đối tượng cần lưu, thuộc tính mô tả thực thể hoặc mối quan hệ, còn lực lượng (cardinality) và participation diễn tả các ràng buộc kết hợp. Khóa của thực thể yếu thường ghép khóa bộ phận với khóa của thực thể chủ.

Phân tích: Khi đọc mô hình ER, tách ba thứ: entity set là loại đối tượng, relationship set là liên kết giữa các loại đối tượng, còn cardinality/participation là ràng buộc số lượng hoặc mức tham gia. Đáp án đúng phải gọi đúng thành phần mà câu hỏi mô tả.

Đối chiếu phương án: A. (c) and (d); D. (d) and (a). Các lựa chọn này không đáp ứng định nghĩa/điều kiện được hỏi; khi ôn, hãy chỉ ra chính xác điểm lệch của từng lựa chọn thay vì học thuộc mỗi chữ cái.
\end{question}

\begin{question}{224}
% [ID gốc: 2.19]
% [Nguồn: DBI202 Order 2 Q19]
What is a weak entity in an ERD?
\begin{choices}
\choice An entity with no attributes
\choice An entity that participates in a one-to-many relationship
\choice An entity that cannot exist without being related to another entity
\choice An entity with a composite key
\end{choices}
\textbf{Đáp án:} C
\textbf{Giải thích:} Đáp án theo dữ liệu đề: C — C. An entity that cannot exist without being related to another entity.

Nền tảng kiến thức: Mô hình ER mô tả yêu cầu ở mức khái niệm: thực thể là đối tượng cần lưu, thuộc tính mô tả thực thể hoặc mối quan hệ, còn lực lượng (cardinality) và participation diễn tả các ràng buộc kết hợp. Khóa của thực thể yếu thường ghép khóa bộ phận với khóa của thực thể chủ.

Phân tích: Thực thể yếu chưa có khóa độc lập; phải ghép khóa bộ phận của nó với khóa của thực thể chủ qua quan hệ định danh. Vì vậy, chỉ lấy thuộc tính của chủ hoặc chỉ lấy khóa bộ phận đều chưa nhận diện duy nhất được nó.

Đối chiếu phương án: A. An entity with no attributes; B. An entity that participates in a one-to-many relationship; D. An entity with a composite key. Các lựa chọn này không đáp ứng định nghĩa/điều kiện được hỏi; khi ôn, hãy chỉ ra chính xác điểm lệch của từng lựa chọn thay vì học thuộc mỗi chữ cái.
\end{question}

\begin{question}{225}
% [ID gốc: 2.20]
% [Nguồn: DBI202 Order 2 Q20]
The key for a weak entity set is_____?
\begin{choices}
\choice Zero or more attributes
\choice The set of attributes of supporting relationships
\choice It is possible for an entity set's key to be composed of attributes, some or all of which belong to another entity set
\choice The set of attributes of supporting entity sets
\end{choices}
\textbf{Đáp án:} C
\textbf{Giải thích:} Đáp án theo dữ liệu đề: C — C. It is possible for an entity set's key to be composed of attributes, some or all of which belong to another entity set.

Nền tảng kiến thức: Mô hình ER mô tả yêu cầu ở mức khái niệm: thực thể là đối tượng cần lưu, thuộc tính mô tả thực thể hoặc mối quan hệ, còn lực lượng (cardinality) và participation diễn tả các ràng buộc kết hợp. Khóa của thực thể yếu thường ghép khóa bộ phận với khóa của thực thể chủ.

Phân tích: Thực thể yếu chưa có khóa độc lập; phải ghép khóa bộ phận của nó với khóa của thực thể chủ qua quan hệ định danh. Vì vậy, chỉ lấy thuộc tính của chủ hoặc chỉ lấy khóa bộ phận đều chưa nhận diện duy nhất được nó.

Đối chiếu phương án: A. Zero or more attributes; B. The set of attributes of supporting relationships; D. The set of attributes of supporting entity sets. Các lựa chọn này không đáp ứng định nghĩa/điều kiện được hỏi; khi ôn, hãy chỉ ra chính xác điểm lệch của từng lựa chọn thay vì học thuộc mỗi chữ cái.
\end{question}

\begin{question}{226}
% [ID gốc: 2.22]
% [Nguồn: DBI202 Order 2 Q22]
What happens to derived attributes when converting an ERD to a relational model?
\begin{choices}
\choice They become primary keys in the related tables
\choice They are calculated and stored as regular attributes in the related tables
\choice They become foreign keys in the related tables
\choice They are ignored during the conversion process
\end{choices}
\textbf{Đáp án:} D
\textbf{Giải thích:} Đáp án theo dữ liệu đề: D — D. They are ignored during the conversion process.

Nền tảng kiến thức: Mô hình ER mô tả yêu cầu ở mức khái niệm: thực thể là đối tượng cần lưu, thuộc tính mô tả thực thể hoặc mối quan hệ, còn lực lượng (cardinality) và participation diễn tả các ràng buộc kết hợp. Khóa của thực thể yếu thường ghép khóa bộ phận với khóa của thực thể chủ.

Phân tích: Trong ký hiệu ER phổ biến, thuộc tính dẫn xuất (derived) được vẽ bằng oval nét đứt vì giá trị có thể tính từ thuộc tính khác, ví dụ tuổi tính từ ngày sinh. Nó khác thuộc tính đa trị (oval kép) hoặc khóa (thường gạch dưới).

Đối chiếu phương án: A. They become primary keys in the related tables; B. They are calculated and stored as regular attributes in the related tables; C. They become foreign keys in the related tables. Các lựa chọn này không đáp ứng định nghĩa/điều kiện được hỏi; khi ôn, hãy chỉ ra chính xác điểm lệch của từng lựa chọn thay vì học thuộc mỗi chữ cái.
\end{question}

\begin{question}{227}
% [ID gốc: 2.23]
% [Nguồn: DBI202 Order 2 Q23]
To covert an entity- relationship diagram to a logical diagram, which of the following will create a new relation that will typically include primary keys from both entities and any additional attributes related to the relationship?
\begin{choices}
\choice A many-to-many relationship
\choice A recursive relationship
\choice A one-to-many relationship
\choice A one-to-one relationship
\end{choices}
\textbf{Đáp án:} A
\textbf{Giải thích:} Đáp án theo dữ liệu đề: A — A. A many-to-many relationship.

Nền tảng kiến thức: Mô hình ER mô tả yêu cầu ở mức khái niệm: thực thể là đối tượng cần lưu, thuộc tính mô tả thực thể hoặc mối quan hệ, còn lực lượng (cardinality) và participation diễn tả các ràng buộc kết hợp. Khóa của thực thể yếu thường ghép khóa bộ phận với khóa của thực thể chủ.

Phân tích: Khi chuyển ER sang quan hệ, quan hệ nhiều-nhiều thường cần một bảng kết nối riêng; bảng này mang khóa của các thực thể tham gia làm khóa ngoại và thường ghép chúng thành khóa chính. Quan hệ một-nhiều thường có thể biểu diễn bằng khóa ngoại ở phía nhiều.

Đối chiếu phương án: B. A recursive relationship; C. A one-to-many relationship; D. A one-to-one relationship. Các lựa chọn này không đáp ứng định nghĩa/điều kiện được hỏi; khi ôn, hãy chỉ ra chính xác điểm lệch của từng lựa chọn thay vì học thuộc mỗi chữ cái.
\end{question}

\begin{question}{228}
% [ID gốc: 2.24]
% [Nguồn: DBI202 Order 2 Q24]
What is a derived attribute in an ERD?
\begin{choices}
\choice An attribute that is calculated from other attributes
\choice An attribute that cannot be calculated
\choice An attribute that is the primary key of an entity
\choice An attribute that represents a relationship
\end{choices}
\textbf{Đáp án:} A
\textbf{Giải thích:} Đáp án theo dữ liệu đề: A — A. An attribute that is calculated from other attributes.

Nền tảng kiến thức: Mô hình ER mô tả yêu cầu ở mức khái niệm: thực thể là đối tượng cần lưu, thuộc tính mô tả thực thể hoặc mối quan hệ, còn lực lượng (cardinality) và participation diễn tả các ràng buộc kết hợp. Khóa của thực thể yếu thường ghép khóa bộ phận với khóa của thực thể chủ.

Phân tích: Trong ký hiệu ER phổ biến, thuộc tính dẫn xuất (derived) được vẽ bằng oval nét đứt vì giá trị có thể tính từ thuộc tính khác, ví dụ tuổi tính từ ngày sinh. Nó khác thuộc tính đa trị (oval kép) hoặc khóa (thường gạch dưới).

Đối chiếu phương án: B. An attribute that cannot be calculated; C. An attribute that is the primary key of an entity; D. An attribute that represents a relationship. Các lựa chọn này không đáp ứng định nghĩa/điều kiện được hỏi; khi ôn, hãy chỉ ra chính xác điểm lệch của từng lựa chọn thay vì học thuộc mỗi chữ cái.
\end{question}

\begin{question}{229}
% [ID gốc: 2.25]
% [Nguồn: DBI202 Order 2 Q25]
In relational database design, which of the following is known as a set of entities of the same type that share the same properties or attributes?
\begin{choices}
\choice Entity Relation model
\choice Entity set
\choice Field set
\choice Record set
\end{choices}
\textbf{Đáp án:} B
\textbf{Giải thích:} Đáp án theo dữ liệu đề: B — B. Entity set.

Nền tảng kiến thức: Mô hình ER mô tả yêu cầu ở mức khái niệm: thực thể là đối tượng cần lưu, thuộc tính mô tả thực thể hoặc mối quan hệ, còn lực lượng (cardinality) và participation diễn tả các ràng buộc kết hợp. Khóa của thực thể yếu thường ghép khóa bộ phận với khóa của thực thể chủ.

Phân tích: Khi đọc mô hình ER, tách ba thứ: entity set là loại đối tượng, relationship set là liên kết giữa các loại đối tượng, còn cardinality/participation là ràng buộc số lượng hoặc mức tham gia. Đáp án đúng phải gọi đúng thành phần mà câu hỏi mô tả.

Đối chiếu phương án: A. Entity Relation model; C. Field set; D. Record set. Các lựa chọn này không đáp ứng định nghĩa/điều kiện được hỏi; khi ôn, hãy chỉ ra chính xác điểm lệch của từng lựa chọn thay vì học thuộc mỗi chữ cái.
\end{question}

\begin{question}{230}
% [ID gốc: 2.26]
% [Nguồn: DBI202 Order 2 Q26]
An E/R diagram is a graphical representation of sets, __________?
\begin{choices}
\choice Entity
\choice Attributes
\choice Relationships
\choice All of the others
\end{choices}
\textbf{Đáp án:} D
\textbf{Giải thích:} Đáp án theo dữ liệu đề: D — D. All of the others.

Nền tảng kiến thức: Mô hình ER mô tả yêu cầu ở mức khái niệm: thực thể là đối tượng cần lưu, thuộc tính mô tả thực thể hoặc mối quan hệ, còn lực lượng (cardinality) và participation diễn tả các ràng buộc kết hợp. Khóa của thực thể yếu thường ghép khóa bộ phận với khóa của thực thể chủ.

Phân tích: ER diagram mô tả các entity set, thuộc tính và relationship set, cùng các ràng buộc cần thiết. Nó tập trung vào đối tượng dữ liệu và liên kết giữa chúng, không mô tả luồng xử lý/chương trình như sơ đồ quy trình.

Đối chiếu phương án: A. Entity; B. Attributes; C. Relationships. Các lựa chọn này không đáp ứng định nghĩa/điều kiện được hỏi; khi ôn, hãy chỉ ra chính xác điểm lệch của từng lựa chọn thay vì học thuộc mỗi chữ cái.
\end{question}

\begin{question}{231}
% [ID gốc: 2.27]
% [Nguồn: DBI202 Order 2 Q27]
To convert from ER Diagram to Relational Model, 1-M relationship:
\begin{choices}
\choice Put key attribute of one-side to M-side
\choice Put key attribute of M-side to one-side
\choice Generate 1 relation, Primary key of this relation combined from two relations
\choice Build a table with two columns, one column for each participating entity set's primary key
\end{choices}
\textbf{Đáp án:} A
\textbf{Giải thích:} Đáp án theo dữ liệu đề: A — A. Put key attribute of one-side to M-side.

Nền tảng kiến thức: Mô hình ER mô tả yêu cầu ở mức khái niệm: thực thể là đối tượng cần lưu, thuộc tính mô tả thực thể hoặc mối quan hệ, còn lực lượng (cardinality) và participation diễn tả các ràng buộc kết hợp. Khóa của thực thể yếu thường ghép khóa bộ phận với khóa của thực thể chủ.

Phân tích: ER diagram mô tả các entity set, thuộc tính và relationship set, cùng các ràng buộc cần thiết. Nó tập trung vào đối tượng dữ liệu và liên kết giữa chúng, không mô tả luồng xử lý/chương trình như sơ đồ quy trình.

Đối chiếu phương án: B. Put key attribute of M-side to one-side; C. Generate 1 relation, Primary key of this relation combined from two relations; D. Build a table with two columns, one column for each participating entity set's primary key. Các lựa chọn này không đáp ứng định nghĩa/điều kiện được hỏi; khi ôn, hãy chỉ ra chính xác điểm lệch của từng lựa chọn thay vì học thuộc mỗi chữ cái.
\end{question}

\begin{question}{232}
% [ID gốc: 2.28]
% [Nguồn: DBI202 Order 2 Q28]
The descriptive property possessed by each entity set is ______
\begin{choices}
\choice Entity
\choice Attribute
\choice Relation
\choice Model
\end{choices}
\textbf{Đáp án:} B
\textbf{Giải thích:} Đáp án theo dữ liệu đề: B — B. Attribute.

Nền tảng kiến thức: Mô hình ER mô tả yêu cầu ở mức khái niệm: thực thể là đối tượng cần lưu, thuộc tính mô tả thực thể hoặc mối quan hệ, còn lực lượng (cardinality) và participation diễn tả các ràng buộc kết hợp. Khóa của thực thể yếu thường ghép khóa bộ phận với khóa của thực thể chủ.

Phân tích: Khi đọc mô hình ER, tách ba thứ: entity set là loại đối tượng, relationship set là liên kết giữa các loại đối tượng, còn cardinality/participation là ràng buộc số lượng hoặc mức tham gia. Đáp án đúng phải gọi đúng thành phần mà câu hỏi mô tả.

Đối chiếu phương án: A. Entity; C. Relation; D. Model. Các lựa chọn này không đáp ứng định nghĩa/điều kiện được hỏi; khi ôn, hãy chỉ ra chính xác điểm lệch của từng lựa chọn thay vì học thuộc mỗi chữ cái.
\end{question}

\begin{question}{233}
% [ID gốc: 2.29]
% [Nguồn: DBI202 Order 2 Q29]
In SQL, ERD use three types of principle elements to form relationships:
\begin{choices}
\choice Entity sets, Constraints and Relationships
\choice Attributes, Constraints, and Relationships
\choice Entity sets, Attributes, and Relationships
\choice Entity sets, Attributes and Constraints
\end{choices}
\textbf{Đáp án:} C
\textbf{Giải thích:} Đáp án theo dữ liệu đề: C — C. Entity sets, Attributes, and Relationships.

Nền tảng kiến thức: Mô hình ER mô tả yêu cầu ở mức khái niệm: thực thể là đối tượng cần lưu, thuộc tính mô tả thực thể hoặc mối quan hệ, còn lực lượng (cardinality) và participation diễn tả các ràng buộc kết hợp. Khóa của thực thể yếu thường ghép khóa bộ phận với khóa của thực thể chủ.

Phân tích: Khi đọc mô hình ER, tách ba thứ: entity set là loại đối tượng, relationship set là liên kết giữa các loại đối tượng, còn cardinality/participation là ràng buộc số lượng hoặc mức tham gia. Đáp án đúng phải gọi đúng thành phần mà câu hỏi mô tả.

Đối chiếu phương án: A. Entity sets, Constraints and Relationships; B. Attributes, Constraints, and Relationships; D. Entity sets, Attributes and Constraints. Các lựa chọn này không đáp ứng định nghĩa/điều kiện được hỏi; khi ôn, hãy chỉ ra chính xác điểm lệch của từng lựa chọn thay vì học thuộc mỗi chữ cái.
\end{question}

\begin{question}{234}
% [ID gốc: 2.30]
% [Nguồn: DBI202 Order 2 Q30]
Which of the following is known as a set of entities of the same type that share same properties, or attributes?
\begin{choices}
\choice Relation set
\choice Tuples
\choice Entity set
\choice Entity Relation mode
\end{choices}
\textbf{Đáp án:} C
\textbf{Giải thích:} Đáp án theo dữ liệu đề: C — C. Entity set.

Nền tảng kiến thức: Mô hình ER mô tả yêu cầu ở mức khái niệm: thực thể là đối tượng cần lưu, thuộc tính mô tả thực thể hoặc mối quan hệ, còn lực lượng (cardinality) và participation diễn tả các ràng buộc kết hợp. Khóa của thực thể yếu thường ghép khóa bộ phận với khóa của thực thể chủ.

Phân tích: Các hàng dữ liệu của một relation được gọi là tuple; header chứa tên thuộc tính không phải tuple dữ liệu. Trong bảng, tuple tương ứng một bản ghi, còn attribute tương ứng cột.

Đối chiếu phương án: A. Relation set; B. Tuples; D. Entity Relation mode. Các lựa chọn này không đáp ứng định nghĩa/điều kiện được hỏi; khi ôn, hãy chỉ ra chính xác điểm lệch của từng lựa chọn thay vì học thuộc mỗi chữ cái.
\end{question}

\begin{question}{235}
% [ID gốc: 2.31]
% [Nguồn: DBI202 Order 2 Q31]
How are weak entity sets represented in an Entity Relationship Diagram (ERD)?
\begin{choices}
\choice With rectangle
\choice With diamond shape
\choice With a double-line rectangle
\choice With a dashed oval
\end{choices}
\textbf{Đáp án:} C
\textbf{Giải thích:} Đáp án theo dữ liệu đề: C — C. With a double-line rectangle.

Nền tảng kiến thức: Mô hình ER mô tả yêu cầu ở mức khái niệm: thực thể là đối tượng cần lưu, thuộc tính mô tả thực thể hoặc mối quan hệ, còn lực lượng (cardinality) và participation diễn tả các ràng buộc kết hợp. Khóa của thực thể yếu thường ghép khóa bộ phận với khóa của thực thể chủ.

Phân tích: Thực thể yếu chưa có khóa độc lập; phải ghép khóa bộ phận của nó với khóa của thực thể chủ qua quan hệ định danh. Vì vậy, chỉ lấy thuộc tính của chủ hoặc chỉ lấy khóa bộ phận đều chưa nhận diện duy nhất được nó.

Đối chiếu phương án: A. With rectangle; B. With diamond shape; D. With a dashed oval. Các lựa chọn này không đáp ứng định nghĩa/điều kiện được hỏi; khi ôn, hãy chỉ ra chính xác điểm lệch của từng lựa chọn thay vì học thuộc mỗi chữ cái.
\end{question}

\begin{question}{236}
% [ID gốc: 2.32]
% [Nguồn: DBI202 Order 2 Q32]
How are composite attributes represented in the relational model?
\begin{choices}
\choice As individual attributes in the related tables
\choice By creating a new table for each composite attribute
\choice By using a single attribute in the related tables
\choice By using a composite key in the related table
\end{choices}
\textbf{Đáp án:} A
\textbf{Giải thích:} Đáp án theo dữ liệu đề: A — A. As individual attributes in the related tables.

Nền tảng kiến thức: Mô hình ER mô tả yêu cầu ở mức khái niệm: thực thể là đối tượng cần lưu, thuộc tính mô tả thực thể hoặc mối quan hệ, còn lực lượng (cardinality) và participation diễn tả các ràng buộc kết hợp. Khóa của thực thể yếu thường ghép khóa bộ phận với khóa của thực thể chủ.

Phân tích: Composite attribute gồm các thành phần có ý nghĩa riêng, chẳng hạn FullName tách thành FirstName và LastName. Khi ánh xạ sang quan hệ, thường lưu các thành phần nguyên tố thành các cột riêng để truy vấn/sắp lọc từng phần.

Đối chiếu phương án: B. By creating a new table for each composite attribute; C. By using a single attribute in the related tables; D. By using a composite key in the related table. Các lựa chọn này không đáp ứng định nghĩa/điều kiện được hỏi; khi ôn, hãy chỉ ra chính xác điểm lệch của từng lựa chọn thay vì học thuộc mỗi chữ cái.
\end{question}

\begin{question}{237}
% [ID gốc: 2.33]
% [Nguồn: DBI202 Order 2 Q33]
In which of the following is a single-entity instance of one type related to many entity instances of another type?
\begin{choices}
\choice One-to-One Relationship
\choice One-to-Many Relationship
\choice Many-to-Many Relationship
\choice Composite Relationship
\end{choices}
\textbf{Đáp án:} B
\textbf{Giải thích:} Đáp án theo dữ liệu đề: B — B. One-to-Many Relationship.

Nền tảng kiến thức: Mô hình ER mô tả yêu cầu ở mức khái niệm: thực thể là đối tượng cần lưu, thuộc tính mô tả thực thể hoặc mối quan hệ, còn lực lượng (cardinality) và participation diễn tả các ràng buộc kết hợp. Khóa của thực thể yếu thường ghép khóa bộ phận với khóa của thực thể chủ.

Phân tích: Khi chuyển ER sang quan hệ, quan hệ nhiều-nhiều thường cần một bảng kết nối riêng; bảng này mang khóa của các thực thể tham gia làm khóa ngoại và thường ghép chúng thành khóa chính. Quan hệ một-nhiều thường có thể biểu diễn bằng khóa ngoại ở phía nhiều.

Đối chiếu phương án: A. One-to-One Relationship; C. Many-to-Many Relationship; D. Composite Relationship. Các lựa chọn này không đáp ứng định nghĩa/điều kiện được hỏi; khi ôn, hãy chỉ ra chính xác điểm lệch của từng lựa chọn thay vì học thuộc mỗi chữ cái.
\end{question}

\begin{question}{238}
% [ID gốc: 2.34]
% [Nguồn: DBI202 Order 2 Q34]
In an ERD, what does a line connecting two entities represent?
\begin{choices}
\choice An attribute
\choice A relationship
\choice A primary key
\choice A foreign key
\end{choices}
\textbf{Đáp án:} B
\textbf{Giải thích:} Đáp án theo dữ liệu đề: B — B. A relationship.

Nền tảng kiến thức: Mô hình ER mô tả yêu cầu ở mức khái niệm: thực thể là đối tượng cần lưu, thuộc tính mô tả thực thể hoặc mối quan hệ, còn lực lượng (cardinality) và participation diễn tả các ràng buộc kết hợp. Khóa của thực thể yếu thường ghép khóa bộ phận với khóa của thực thể chủ.

Phân tích: Khi đọc mô hình ER, tách ba thứ: entity set là loại đối tượng, relationship set là liên kết giữa các loại đối tượng, còn cardinality/participation là ràng buộc số lượng hoặc mức tham gia. Đáp án đúng phải gọi đúng thành phần mà câu hỏi mô tả.

Đối chiếu phương án: A. An attribute; C. A primary key; D. A foreign key. Các lựa chọn này không đáp ứng định nghĩa/điều kiện được hỏi; khi ôn, hãy chỉ ra chính xác điểm lệch của từng lựa chọn thay vì học thuộc mỗi chữ cái.
\end{question}

\begin{question}{239}
% [ID gốc: 2.35]
% [Nguồn: DBI202 Order 2 Q35]
Given that Cartoons and Murder Mysteries are special Movie genres, what is the relationship between entities Cartoons, Murder Mysteries and Movie?
\begin{choices}
\choice isa relationship
\choice ais relationship
\choice supporting relationship
\choice referential integrity relationship
\end{choices}
\textbf{Đáp án:} A
\textbf{Giải thích:} Đáp án theo dữ liệu đề: A — A. isa relationship.

Nền tảng kiến thức: Mô hình ER mô tả yêu cầu ở mức khái niệm: thực thể là đối tượng cần lưu, thuộc tính mô tả thực thể hoặc mối quan hệ, còn lực lượng (cardinality) và participation diễn tả các ràng buộc kết hợp. Khóa của thực thể yếu thường ghép khóa bộ phận với khóa của thực thể chủ.

Phân tích: Khóa ngoại buộc giá trị tham chiếu phải khớp một khóa được tham chiếu ở quan hệ đích (hoặc NULL nếu ràng buộc cho phép). Đây là toàn vẹn tham chiếu; nó khác tính duy nhất của primary/candidate key.

Đối chiếu phương án: B. ais relationship; C. supporting relationship; D. referential integrity relationship. Các lựa chọn này không đáp ứng định nghĩa/điều kiện được hỏi; khi ôn, hãy chỉ ra chính xác điểm lệch của từng lựa chọn thay vì học thuộc mỗi chữ cái.
\end{question}

\begin{question}{240}
% [ID gốc: 2.36]
% [Nguồn: DBI202 Order 2 Q36]
Let E1 and E2 be two entities in an E/R diagram with simple single-valued attributes. R1 and R2 are two relationships between E1 and E2, where R1 is one-to-many and R2 is many-to-many. R1 and R2 do not have any attributes of their own. What is the minimum
\begin{choices}
\choice 2
\choice 3
\choice 4
\choice 5
\end{choices}
\textbf{Đáp án:} B
\textbf{Giải thích:} Đáp án theo dữ liệu đề: B — B. 3.

Nền tảng kiến thức: Mô hình ER mô tả yêu cầu ở mức khái niệm: thực thể là đối tượng cần lưu, thuộc tính mô tả thực thể hoặc mối quan hệ, còn lực lượng (cardinality) và participation diễn tả các ràng buộc kết hợp. Khóa của thực thể yếu thường ghép khóa bộ phận với khóa của thực thể chủ.

Phân tích: ER diagram mô tả các entity set, thuộc tính và relationship set, cùng các ràng buộc cần thiết. Nó tập trung vào đối tượng dữ liệu và liên kết giữa chúng, không mô tả luồng xử lý/chương trình như sơ đồ quy trình.

Đối chiếu phương án: A. 2; C. 4; D. 5. Các lựa chọn này không đáp ứng định nghĩa/điều kiện được hỏi; khi ôn, hãy chỉ ra chính xác điểm lệch của từng lựa chọn thay vì học thuộc mỗi chữ cái.
\end{question}

\begin{question}{241}
% [ID gốc: 2.39]
% [Nguồn: DBI202 Order 2 Q39]
What is the purpose of a foreign key in an ERD?
\begin{choices}
\choice It defines the data type of attributes
\choice It determines the cardinality of relationships
\choice It serves as the primary identifier for a table
\choice It establishes a link between two tables
\end{choices}
\textbf{Đáp án:} D
\textbf{Giải thích:} Đáp án theo dữ liệu đề: D — D. It establishes a link between two tables.

Lời giải gốc: Khóa ngoại được dùng để thiết lập mối liên kết và đảm bảo tính toàn vẹn tham chiếu giữa hai bảng trong cơ sở dữ liệu.

Nền tảng kiến thức: Mô hình ER mô tả yêu cầu ở mức khái niệm: thực thể là đối tượng cần lưu, thuộc tính mô tả thực thể hoặc mối quan hệ, còn lực lượng (cardinality) và participation diễn tả các ràng buộc kết hợp. Khóa của thực thể yếu thường ghép khóa bộ phận với khóa của thực thể chủ.

Phân tích: Khi đọc mô hình ER, tách ba thứ: entity set là loại đối tượng, relationship set là liên kết giữa các loại đối tượng, còn cardinality/participation là ràng buộc số lượng hoặc mức tham gia. Đáp án đúng phải gọi đúng thành phần mà câu hỏi mô tả.

Đối chiếu phương án: A. It defines the data type of attributes; B. It determines the cardinality of relationships; C. It serves as the primary identifier for a table. Các lựa chọn này không đáp ứng định nghĩa/điều kiện được hỏi; khi ôn, hãy chỉ ra chính xác điểm lệch của từng lựa chọn thay vì học thuộc mỗi chữ cái.
\end{question}

\begin{question}{242}
% [ID gốc: 2.40]
% [Nguồn: DBI202 Order 2 Q40]
In a binary relationship with cardinality (1, N), what does it signify?
\begin{choices}
\choice One entity can be related to multiple instances of another entity, but not vice versa
\choice Multiple instances of one entity can relate to multiple instances of another entity
\choice One entity can relate to only one instance of another entity
\choice Each instance of one entity must relate to a unique instance of another entity
\end{choices}
\textbf{Đáp án:} A
\textbf{Giải thích:} Đáp án theo dữ liệu đề: A — A. One entity can be related to multiple instances of another entity, but not vice versa.

Lời giải gốc: Quan hệ 1-N (One-to-Many) nghĩa là một thực thể thuộc tập thực thể này có thể liên kết với nhiều thực thể thuộc tập thực thể kia, nhưng chiều ngược lại mỗi thực thể chỉ liên kết với duy nhất một thực thể.

Nền tảng kiến thức: Mô hình ER mô tả yêu cầu ở mức khái niệm: thực thể là đối tượng cần lưu, thuộc tính mô tả thực thể hoặc mối quan hệ, còn lực lượng (cardinality) và participation diễn tả các ràng buộc kết hợp. Khóa của thực thể yếu thường ghép khóa bộ phận với khóa của thực thể chủ.

Phân tích: Binary relationship/association nối hai kiểu thực thể/lớp; n-ary relationship có từ ba kiểu tham gia trở lên. Phân loại theo số entity set/class trong chính liên kết, không theo số thuộc tính của chúng.

Đối chiếu phương án: B. Multiple instances of one entity can relate to multiple instances of another entity; C. One entity can relate to only one instance of another entity; D. Each instance of one entity must relate to a unique instance of another entity. Các lựa chọn này không đáp ứng định nghĩa/điều kiện được hỏi; khi ôn, hãy chỉ ra chính xác điểm lệch của từng lựa chọn thay vì học thuộc mỗi chữ cái.
\end{question}

\begin{question}{243}
% [ID gốc: 2.41]
% [Nguồn: DBI202 Order 2 Q41]
An E/R diagram is a graphical representation of sets of what?
\begin{choices}
\choice Entities
\choice Attributes
\choice Relationships
\choice All of the others
\end{choices}
\textbf{Đáp án:} D
\textbf{Giải thích:} Đáp án theo dữ liệu đề: D — D. All of the others.

Lời giải gốc: Sơ đồ thực thể kết hợp (E/R diagram) là mô hình đồ họa biểu diễn tập thực thể, tập thuộc tính và tập mối quan hệ giữa chúng.

Nền tảng kiến thức: Mô hình ER mô tả yêu cầu ở mức khái niệm: thực thể là đối tượng cần lưu, thuộc tính mô tả thực thể hoặc mối quan hệ, còn lực lượng (cardinality) và participation diễn tả các ràng buộc kết hợp. Khóa của thực thể yếu thường ghép khóa bộ phận với khóa của thực thể chủ.

Phân tích: ER diagram mô tả các entity set, thuộc tính và relationship set, cùng các ràng buộc cần thiết. Nó tập trung vào đối tượng dữ liệu và liên kết giữa chúng, không mô tả luồng xử lý/chương trình như sơ đồ quy trình.

Đối chiếu phương án: A. Entities; B. Attributes; C. Relationships. Các lựa chọn này không đáp ứng định nghĩa/điều kiện được hỏi; khi ôn, hãy chỉ ra chính xác điểm lệch của từng lựa chọn thay vì học thuộc mỗi chữ cái.
\end{question}

\begin{question}{244}
% [ID gốc: 2.42]
% [Nguồn: DBI202 Order 2 Q42]
Which of the following best describes an entity in an ERD (Entity-Relationship Diagram)?
\begin{choices}
\choice A table with rows and columns
\choice A real-world object or concept with attributes
\choice A set of rows, columns, and constraints
\choice A table with rows, columns, and domain constraints
\end{choices}
\textbf{Đáp án:} B
\textbf{Giải thích:} Đáp án theo dữ liệu đề: B — B. A real-world object or concept with attributes.

Lời giải gốc: Trong sơ đồ ERD, thực thể (Entity) đại diện cho một đối tượng hoặc khái niệm trong thế giới thực có các thuộc tính mô tả.

Nền tảng kiến thức: Mô hình ER mô tả yêu cầu ở mức khái niệm: thực thể là đối tượng cần lưu, thuộc tính mô tả thực thể hoặc mối quan hệ, còn lực lượng (cardinality) và participation diễn tả các ràng buộc kết hợp. Khóa của thực thể yếu thường ghép khóa bộ phận với khóa của thực thể chủ.

Phân tích: Khi đọc mô hình ER, tách ba thứ: entity set là loại đối tượng, relationship set là liên kết giữa các loại đối tượng, còn cardinality/participation là ràng buộc số lượng hoặc mức tham gia. Đáp án đúng phải gọi đúng thành phần mà câu hỏi mô tả.

Đối chiếu phương án: A. A table with rows and columns; C. A set of rows, columns, and constraints; D. A table with rows, columns, and domain constraints. Các lựa chọn này không đáp ứng định nghĩa/điều kiện được hỏi; khi ôn, hãy chỉ ra chính xác điểm lệch của từng lựa chọn thay vì học thuộc mỗi chữ cái.
\end{question}

\begin{question}{245}
% [ID gốc: 2.44]
% [Nguồn: DBI202 Order 2 Q44]
Which of the following statements is correct when converting from the ERD model to the relational model?
\begin{choices}
\choice The order of the columns is not important
\choice The order of both rows and columns is important
\choice The order of the rows is important
\choice The order of the columns is important
\end{choices}
\textbf{Đáp án:} A
\textbf{Giải thích:} Đáp án theo dữ liệu đề: A — A. The order of the columns is not important.

Lời giải gốc: Trong mô hình quan hệ, thứ tự sắp xếp vật lý của các hàng (tuples) và cột (attributes) về mặt lý thuyết là không quan trọng.

Nền tảng kiến thức: Mô hình ER mô tả yêu cầu ở mức khái niệm: thực thể là đối tượng cần lưu, thuộc tính mô tả thực thể hoặc mối quan hệ, còn lực lượng (cardinality) và participation diễn tả các ràng buộc kết hợp. Khóa của thực thể yếu thường ghép khóa bộ phận với khóa của thực thể chủ.

Phân tích: ER diagram mô tả các entity set, thuộc tính và relationship set, cùng các ràng buộc cần thiết. Nó tập trung vào đối tượng dữ liệu và liên kết giữa chúng, không mô tả luồng xử lý/chương trình như sơ đồ quy trình.

Đối chiếu phương án: B. The order of both rows and columns is important; C. The order of the rows is important; D. The order of the columns is important. Các lựa chọn này không đáp ứng định nghĩa/điều kiện được hỏi; khi ôn, hãy chỉ ra chính xác điểm lệch của từng lựa chọn thay vì học thuộc mỗi chữ cái.
\end{question}

\begin{question}{246}
% [ID gốc: 2.45]
% [Nguồn: DBI202 Order 2 Q45]
In the entity-relationship model, what is the entity?
\begin{choices}
\choice Relationships
\choice Relationships between objects
\choice Objects and relationships between objects
\choice Data objects
\end{choices}
\textbf{Đáp án:} D
\textbf{Giải thích:} Đáp án theo dữ liệu đề: D — D. Data objects.

Lời giải gốc: Thực thể (Entity) trong mô hình ER đại diện cho các đối tượng dữ liệu (Data objects) hoặc sự vật trong thế giới thực.

Nền tảng kiến thức: Mô hình ER mô tả yêu cầu ở mức khái niệm: thực thể là đối tượng cần lưu, thuộc tính mô tả thực thể hoặc mối quan hệ, còn lực lượng (cardinality) và participation diễn tả các ràng buộc kết hợp. Khóa của thực thể yếu thường ghép khóa bộ phận với khóa của thực thể chủ.

Phân tích: Entity là một đối tượng/sự vật có thể phân biệt trong miền ứng dụng và cần lưu thông tin; entity set là tập các entity cùng loại. Ví dụ một sinh viên cụ thể là entity, còn Student là entity set/type.

Đối chiếu phương án: A. Relationships; B. Relationships between objects; C. Objects and relationships between objects. Các lựa chọn này không đáp ứng định nghĩa/điều kiện được hỏi; khi ôn, hãy chỉ ra chính xác điểm lệch của từng lựa chọn thay vì học thuộc mỗi chữ cái.
\end{question}

\begin{question}{247}
% [ID gốc: 2.47]
% [Nguồn: DBI202 Order 2 Q47]
How are many-to-many relationships typically represented in the relational model?
\begin{choices}
\choice By adding attributes to one of the related tables
\choice By creating a new table with foreign keys from both related tables
\choice By denormalizing the database
\choice By creating a new table for each related entity
\end{choices}
\textbf{Đáp án:} B
\textbf{Giải thích:} Đáp án theo dữ liệu đề: B — B. By creating a new table with foreign keys from both related tables.

Lời giải gốc: Mối quan hệ nhiều-nhiều được biểu diễn bằng cách tạo một bảng liên kết (junction table) chứa các khóa ngoại trỏ đến khóa chính của hai bảng liên quan.

Nền tảng kiến thức: Mô hình ER mô tả yêu cầu ở mức khái niệm: thực thể là đối tượng cần lưu, thuộc tính mô tả thực thể hoặc mối quan hệ, còn lực lượng (cardinality) và participation diễn tả các ràng buộc kết hợp. Khóa của thực thể yếu thường ghép khóa bộ phận với khóa của thực thể chủ.

Phân tích: Khi chuyển ER sang quan hệ, quan hệ nhiều-nhiều thường cần một bảng kết nối riêng; bảng này mang khóa của các thực thể tham gia làm khóa ngoại và thường ghép chúng thành khóa chính. Quan hệ một-nhiều thường có thể biểu diễn bằng khóa ngoại ở phía nhiều.

Đối chiếu phương án: A. By adding attributes to one of the related tables; C. By denormalizing the database; D. By creating a new table for each related entity. Các lựa chọn này không đáp ứng định nghĩa/điều kiện được hỏi; khi ôn, hãy chỉ ra chính xác điểm lệch của từng lựa chọn thay vì học thuộc mỗi chữ cái.
\end{question}

\begin{question}{248}
% [ID gốc: 2.51]
% [Nguồn: DBI202 Order 2 Q51]
Which of the following best describes a "many-to-many" relationship?
\begin{choices}
\choice One entity is related to one other entity
\choice One entity is related to many other entities
\choice Many entities are related to many other entities
\choice None of the others
\end{choices}
\textbf{Đáp án:} C
\textbf{Giải thích:} Đáp án theo dữ liệu đề: C — C. Many entities are related to many other entities.

Lời giải gốc: Mối quan hệ nhiều-nhiều (Many-to-many) là mối quan hệ trong đó nhiều thể hiện của thực thể này liên kết với nhiều thể hiện của thực thể kia.

Nền tảng kiến thức: Mô hình ER mô tả yêu cầu ở mức khái niệm: thực thể là đối tượng cần lưu, thuộc tính mô tả thực thể hoặc mối quan hệ, còn lực lượng (cardinality) và participation diễn tả các ràng buộc kết hợp. Khóa của thực thể yếu thường ghép khóa bộ phận với khóa của thực thể chủ.

Phân tích: Khi chuyển ER sang quan hệ, quan hệ nhiều-nhiều thường cần một bảng kết nối riêng; bảng này mang khóa của các thực thể tham gia làm khóa ngoại và thường ghép chúng thành khóa chính. Quan hệ một-nhiều thường có thể biểu diễn bằng khóa ngoại ở phía nhiều.

Đối chiếu phương án: A. One entity is related to one other entity; B. One entity is related to many other entities; D. None of the others. Các lựa chọn này không đáp ứng định nghĩa/điều kiện được hỏi; khi ôn, hãy chỉ ra chính xác điểm lệch của từng lựa chọn thay vì học thuộc mỗi chữ cái.
\end{question}

\begin{question}{249}
% [ID gốc: 2.52]
% [Nguồn: DBI202 Order 2 Q52]
To convert from an ER Diagram to the Relational Model, a 1-M relationship is represented by:
\begin{choices}
\choice Putting the key attribute of the one-side into the M-side
\choice Putting the key attribute of the M-side into the one-side
\choice Generating 1 relation whose primary key is combined from two relations
\choice Building a table with two columns, one for each participating entity set's primary key
\end{choices}
\textbf{Đáp án:} A
\textbf{Giải thích:} Đáp án theo dữ liệu đề: A — A. Putting the key attribute of the one-side into the M-side.

Lời giải gốc: Khi chuyển đổi mối quan hệ một-nhiều (1-M) sang mô hình quan hệ, ta đưa khóa chính của phía "1" sang làm khóa ngoại ở phía "M".

Nền tảng kiến thức: Mô hình ER mô tả yêu cầu ở mức khái niệm: thực thể là đối tượng cần lưu, thuộc tính mô tả thực thể hoặc mối quan hệ, còn lực lượng (cardinality) và participation diễn tả các ràng buộc kết hợp. Khóa của thực thể yếu thường ghép khóa bộ phận với khóa của thực thể chủ.

Phân tích: ER diagram mô tả các entity set, thuộc tính và relationship set, cùng các ràng buộc cần thiết. Nó tập trung vào đối tượng dữ liệu và liên kết giữa chúng, không mô tả luồng xử lý/chương trình như sơ đồ quy trình.

Đối chiếu phương án: B. Putting the key attribute of the M-side into the one-side; C. Generating 1 relation whose primary key is combined from two relations; D. Building a table with two columns, one for each participating entity set's primary key. Các lựa chọn này không đáp ứng định nghĩa/điều kiện được hỏi; khi ôn, hãy chỉ ra chính xác điểm lệch của từng lựa chọn thay vì học thuộc mỗi chữ cái.
\end{question}

\begin{question}{250}
% [ID gốc: 2.54]
% [Nguồn: DBI202 Order 2 Q54]
Which of the following is an example of a composite attribute?
\begin{choices}
\choice StudentID with unique value
\choice FullName (split into FirstName and LastName)
\choice EmailAddress with multi-values
\choice DateOfBirth value is accepted with format "dd/mm/yyyy"
\end{choices}
\textbf{Đáp án:} B
\textbf{Giải thích:} Đáp án theo dữ liệu đề: B — B. FullName (split into FirstName and LastName).

Lời giải gốc: Thuộc tính phức hợp (Composite attribute) là thuộc tính có thể chia nhỏ thành các thuộc tính đơn giản hơn, ví dụ FullName gồm FirstName và LastName.

Nền tảng kiến thức: Mô hình ER mô tả yêu cầu ở mức khái niệm: thực thể là đối tượng cần lưu, thuộc tính mô tả thực thể hoặc mối quan hệ, còn lực lượng (cardinality) và participation diễn tả các ràng buộc kết hợp. Khóa của thực thể yếu thường ghép khóa bộ phận với khóa của thực thể chủ.

Phân tích: Composite attribute gồm các thành phần có ý nghĩa riêng, chẳng hạn FullName tách thành FirstName và LastName. Khi ánh xạ sang quan hệ, thường lưu các thành phần nguyên tố thành các cột riêng để truy vấn/sắp lọc từng phần.

Đối chiếu phương án: A. StudentID with unique value; C. EmailAddress with multi-values; D. DateOfBirth value is accepted with format "dd/mm/yyyy". Các lựa chọn này không đáp ứng định nghĩa/điều kiện được hỏi; khi ôn, hãy chỉ ra chính xác điểm lệch của từng lựa chọn thay vì học thuộc mỗi chữ cái.
\end{question}

\begin{question}{251}
% [ID gốc: 2.55]
% [Nguồn: DBI202 Order 2 Q55]
Given that Cartoons and Murder Mysteries are special Movie genres, what is the relationship between the entities Cartoons, Murder Mysteries, and Movie?
\begin{choices}
\choice Is-a relationship
\choice A-is relationship
\choice Supporting relationship
\choice Referential integrity relationship
\end{choices}
\textbf{Đáp án:} A
\textbf{Giải thích:} Đáp án theo dữ liệu đề: A — A. Is-a relationship.

Lời giải gốc: Đây là mối quan hệ tổng quát hóa/chuyên biệt hóa (Generalization/Specialization), thể hiện qua mối quan hệ "is-a" (Cartoons là một dạng của Movie).

Nền tảng kiến thức: Mô hình ER mô tả yêu cầu ở mức khái niệm: thực thể là đối tượng cần lưu, thuộc tính mô tả thực thể hoặc mối quan hệ, còn lực lượng (cardinality) và participation diễn tả các ràng buộc kết hợp. Khóa của thực thể yếu thường ghép khóa bộ phận với khóa của thực thể chủ.

Phân tích: Khóa ngoại buộc giá trị tham chiếu phải khớp một khóa được tham chiếu ở quan hệ đích (hoặc NULL nếu ràng buộc cho phép). Đây là toàn vẹn tham chiếu; nó khác tính duy nhất của primary/candidate key.

Đối chiếu phương án: B. A-is relationship; C. Supporting relationship; D. Referential integrity relationship. Các lựa chọn này không đáp ứng định nghĩa/điều kiện được hỏi; khi ôn, hãy chỉ ra chính xác điểm lệch của từng lựa chọn thay vì học thuộc mỗi chữ cái.
\end{question}

\begin{question}{252}
% [ID gốc: 2.58]
% [Nguồn: DBI202 Order 2 Q58]
In an ERD, which of the following best describes a one-to-many (1:M) relationship?
\begin{choices}
\choice Each instance of Entity A is related to exactly one instance of Entity B
\choice Each instance of Entity A is related to many instances of Entity B
\choice Each instance of both entities is related to exactly one other instance
\choice Entities are not related at all
\end{choices}
\textbf{Đáp án:} B
\textbf{Giải thích:} Đáp án theo dữ liệu đề: B — B. Each instance of Entity A is related to many instances of Entity B.

Lời giải gốc: Mối quan hệ 1-N được định nghĩa là: Mỗi một bản ghi ở thực thể A có thể liên kết (tham chiếu) tới nhiều bản ghi ở thực thể B.

Nền tảng kiến thức: Mô hình ER mô tả yêu cầu ở mức khái niệm: thực thể là đối tượng cần lưu, thuộc tính mô tả thực thể hoặc mối quan hệ, còn lực lượng (cardinality) và participation diễn tả các ràng buộc kết hợp. Khóa của thực thể yếu thường ghép khóa bộ phận với khóa của thực thể chủ.

Phân tích: Khi đọc mô hình ER, tách ba thứ: entity set là loại đối tượng, relationship set là liên kết giữa các loại đối tượng, còn cardinality/participation là ràng buộc số lượng hoặc mức tham gia. Đáp án đúng phải gọi đúng thành phần mà câu hỏi mô tả.

Đối chiếu phương án: A. Each instance of Entity A is related to exactly one instance of Entity B; C. Each instance of both entities is related to exactly one other instance; D. Entities are not related at all. Các lựa chọn này không đáp ứng định nghĩa/điều kiện được hỏi; khi ôn, hãy chỉ ra chính xác điểm lệch của từng lựa chọn thay vì học thuộc mỗi chữ cái.
\end{question}

\begin{question}{253}
% [ID gốc: 2.59]
% [Nguồn: DBI202 Order 2 Q59]
In SQL, ERD use three types of principal elements to form relationships:
\begin{choices}
\choice Entity sets, Constraints and Relationships
\choice Attributes, Constraints, and Relationships
\choice Entity sets, Attributes, and Relationships
\choice Entity sets, Attributes and Constraints
\end{choices}
\textbf{Đáp án:} C
\textbf{Giải thích:} Đáp án theo dữ liệu đề: C — C. Entity sets, Attributes, and Relationships.

Lời giải gốc: Ba thành phần chính cơ bản của Sơ đồ thực thể mối quan hệ (ERD) bao gồm: Tập thực thể (Entity sets), Thuộc tính (Attributes) và Mối quan hệ (Relationships).

Nền tảng kiến thức: Mô hình ER mô tả yêu cầu ở mức khái niệm: thực thể là đối tượng cần lưu, thuộc tính mô tả thực thể hoặc mối quan hệ, còn lực lượng (cardinality) và participation diễn tả các ràng buộc kết hợp. Khóa của thực thể yếu thường ghép khóa bộ phận với khóa của thực thể chủ.

Phân tích: Khi đọc mô hình ER, tách ba thứ: entity set là loại đối tượng, relationship set là liên kết giữa các loại đối tượng, còn cardinality/participation là ràng buộc số lượng hoặc mức tham gia. Đáp án đúng phải gọi đúng thành phần mà câu hỏi mô tả.

Đối chiếu phương án: A. Entity sets, Constraints and Relationships; B. Attributes, Constraints, and Relationships; D. Entity sets, Attributes and Constraints. Các lựa chọn này không đáp ứng định nghĩa/điều kiện được hỏi; khi ôn, hãy chỉ ra chính xác điểm lệch của từng lựa chọn thay vì học thuộc mỗi chữ cái.
\end{question}

\begin{question}{254}
% [ID gốc: 2.60]
% [Nguồn: DBI202 Order 2 Q60]
What is the relationship between a weak entity set and its owner entity in an E/R diagram?
\begin{choices}
\choice One-to-one relationship.
\choice Many-to-many relationship.
\choice Many-to-one relationship.
\choice One-to-many relationship.
\end{choices}
\textbf{Đáp án:} C
\textbf{Giải thích:} Đáp án theo dữ liệu đề: D — D. One-to-many relationship..

Lời giải gốc: Một thực thể cha (owner) có thể sở hữu/định danh cho nhiều thực thể con thuộc tập thực thể yếu (weak entity set). Ngược lại, một thực thể yếu chỉ thuộc về duy nhất một thực thể cha. Do đó mối quan hệ từ phía Owner sang Weak Entity là mối quan hệ 1 - Nhiều (One-to-many).

Nền tảng kiến thức: Mô hình ER mô tả yêu cầu ở mức khái niệm: thực thể là đối tượng cần lưu, thuộc tính mô tả thực thể hoặc mối quan hệ, còn lực lượng (cardinality) và participation diễn tả các ràng buộc kết hợp. Khóa của thực thể yếu thường ghép khóa bộ phận với khóa của thực thể chủ.

Phân tích: ER diagram mô tả các entity set, thuộc tính và relationship set, cùng các ràng buộc cần thiết. Nó tập trung vào đối tượng dữ liệu và liên kết giữa chúng, không mô tả luồng xử lý/chương trình như sơ đồ quy trình.

Đối chiếu phương án: A. One-to-one relationship.; B. Many-to-many relationship.; C. Many-to-one relationship.. Các lựa chọn này không đáp ứng định nghĩa/điều kiện được hỏi; khi ôn, hãy chỉ ra chính xác điểm lệch của từng lựa chọn thay vì học thuộc mỗi chữ cái.
\end{question}

\begin{question}{255}
% [ID gốc: 2.61]
% [Nguồn: DBI202 Order 2 Q61]
Which of the following is a set of entities of the same type that share the same properties or attributes?
\begin{choices}
\choice Relation set
\choice Tuples
\choice Entity set
\choice Entity Relationship model
\end{choices}
\textbf{Đáp án:} C
\textbf{Giải thích:} Đáp án theo dữ liệu đề: C — C. Entity set.

Lời giải gốc: Tập thực thể (Entity set) là tập hợp tất cả các thực thể có cùng kiểu và chia sẻ các thuộc tính giống nhau.

Nền tảng kiến thức: Mô hình ER mô tả yêu cầu ở mức khái niệm: thực thể là đối tượng cần lưu, thuộc tính mô tả thực thể hoặc mối quan hệ, còn lực lượng (cardinality) và participation diễn tả các ràng buộc kết hợp. Khóa của thực thể yếu thường ghép khóa bộ phận với khóa của thực thể chủ.

Phân tích: Các hàng dữ liệu của một relation được gọi là tuple; header chứa tên thuộc tính không phải tuple dữ liệu. Trong bảng, tuple tương ứng một bản ghi, còn attribute tương ứng cột.

Đối chiếu phương án: A. Relation set; B. Tuples; D. Entity Relationship model. Các lựa chọn này không đáp ứng định nghĩa/điều kiện được hỏi; khi ôn, hãy chỉ ra chính xác điểm lệch của từng lựa chọn thay vì học thuộc mỗi chữ cái.
\end{question}

\begin{question}{256}
% [ID gốc: 2.62]
% [Nguồn: DBI202 Order 2 Q62]
What type of attribute in an ER diagram is represented with a dashed oval?
\begin{choices}
\choice Multivalued attribute
\choice Derived attribute
\choice Composite attribute
\choice Primary attribute
\end{choices}
\textbf{Đáp án:} B
\textbf{Giải thích:} Đáp án theo dữ liệu đề: B — B. Derived attribute.

Lời giải gốc: Trong biểu đồ ER (Chen's notation), thuộc tính dẫn xuất (Derived attribute) - thuộc tính có thể được tính toán ra từ thuộc tính khác (ví dụ: Tuổi suy ra từ Ngày sinh) - được biểu diễn bằng hình oval nét đứt.

Nền tảng kiến thức: Mô hình ER mô tả yêu cầu ở mức khái niệm: thực thể là đối tượng cần lưu, thuộc tính mô tả thực thể hoặc mối quan hệ, còn lực lượng (cardinality) và participation diễn tả các ràng buộc kết hợp. Khóa của thực thể yếu thường ghép khóa bộ phận với khóa của thực thể chủ.

Phân tích: ER diagram mô tả các entity set, thuộc tính và relationship set, cùng các ràng buộc cần thiết. Nó tập trung vào đối tượng dữ liệu và liên kết giữa chúng, không mô tả luồng xử lý/chương trình như sơ đồ quy trình.

Đối chiếu phương án: A. Multivalued attribute; C. Composite attribute; D. Primary attribute. Các lựa chọn này không đáp ứng định nghĩa/điều kiện được hỏi; khi ôn, hãy chỉ ra chính xác điểm lệch của từng lựa chọn thay vì học thuộc mỗi chữ cái.
\end{question}

\begin{question}{257}
% [ID gốc: 2.65]
% [Nguồn: DBI202 Order 2 Q65]
What is the main purpose of a conceptual data model in database design?
\begin{choices}
\choice To write SQL queries for the database
\choice To define the physical structure of tables and indexes
\choice To represent high-level business entities and relationships
\choice To optimize database performance by normalization and decomposition
\end{choices}
\textbf{Đáp án:} C
\textbf{Giải thích:} Đáp án theo dữ liệu đề: C — C. To represent high-level business entities and relationships.

Lời giải gốc: Mục đích chính của mô hình dữ liệu khái niệm là mô tả các thực thể nghiệp vụ cấp cao và mối quan hệ giữa chúng, độc lập với hệ quản trị cơ sở dữ liệu vật lý.

Nền tảng kiến thức: Mô hình ER mô tả yêu cầu ở mức khái niệm: thực thể là đối tượng cần lưu, thuộc tính mô tả thực thể hoặc mối quan hệ, còn lực lượng (cardinality) và participation diễn tả các ràng buộc kết hợp. Khóa của thực thể yếu thường ghép khóa bộ phận với khóa của thực thể chủ.

Phân tích: Data model là bộ khái niệm dùng để mô tả cấu trúc dữ liệu, quan hệ giữa dữ liệu và các ràng buộc. Relational, hierarchical, network và graph là các mô hình có thật; cần chọn mô hình khớp đặc điểm được hỏi, không nhầm mô hình với DBMS hay ngôn ngữ thao tác.

Đối chiếu phương án: A. To write SQL queries for the database; B. To define the physical structure of tables and indexes; D. To optimize database performance by normalization and decomposition. Các lựa chọn này không đáp ứng định nghĩa/điều kiện được hỏi; khi ôn, hãy chỉ ra chính xác điểm lệch của từng lựa chọn thay vì học thuộc mỗi chữ cái.
\end{question}

\begin{question}{258}
% [ID gốc: 2.69]
% [Nguồn: DBI202 Order 2 Q69]
______ Real-world thing, distinguishable from other objects, in ERD.
\begin{choices}
\choice Entity
\choice Diagram
\choice Relationship
\choice Primary entity
\end{choices}
\textbf{Đáp án:} A
\textbf{Giải thích:} Đáp án theo dữ liệu đề: A — A. Entity.

Lời giải gốc: Thực thể (Entity) là một đối tượng tồn tại trong thế giới thực và có thể phân biệt rõ ràng được với các đối tượng khác xung quanh.

Nền tảng kiến thức: Mô hình ER mô tả yêu cầu ở mức khái niệm: thực thể là đối tượng cần lưu, thuộc tính mô tả thực thể hoặc mối quan hệ, còn lực lượng (cardinality) và participation diễn tả các ràng buộc kết hợp. Khóa của thực thể yếu thường ghép khóa bộ phận với khóa của thực thể chủ.

Phân tích: Khi đọc mô hình ER, tách ba thứ: entity set là loại đối tượng, relationship set là liên kết giữa các loại đối tượng, còn cardinality/participation là ràng buộc số lượng hoặc mức tham gia. Đáp án đúng phải gọi đúng thành phần mà câu hỏi mô tả.

Đối chiếu phương án: B. Diagram; C. Relationship; D. Primary entity. Các lựa chọn này không đáp ứng định nghĩa/điều kiện được hỏi; khi ôn, hãy chỉ ra chính xác điểm lệch của từng lựa chọn thay vì học thuộc mỗi chữ cái.
\end{question}

\begin{question}{259}
% [ID gốc: 2.71]
% [Nguồn: DBI202 Order 2 Q71]
If two entities in an ER model share the same attribute name but represent different meanings, what should be done during schema transformation?
\begin{choices}
\choice Always add entity name to the attribute name to make the meaning clear.
\choice Rename one of the attributes to avoid ambiguity
\choice Merge both entities
\choice Ignore the attribute names
\end{choices}
\textbf{Đáp án:} B
\textbf{Giải thích:} Đáp án theo dữ liệu đề: B — B. Rename one of the attributes to avoid ambiguity.

Lời giải gốc: Khi hai thuộc tính trùng tên nhưng mang ý nghĩa khác nhau ở các thực thể khác nhau, việc đổi tên (Rename) thuộc tính là bắt buộc để tránh sự mơ hồ khi chuyển sang lược đồ quan hệ.

Nền tảng kiến thức: Mô hình ER mô tả yêu cầu ở mức khái niệm: thực thể là đối tượng cần lưu, thuộc tính mô tả thực thể hoặc mối quan hệ, còn lực lượng (cardinality) và participation diễn tả các ràng buộc kết hợp. Khóa của thực thể yếu thường ghép khóa bộ phận với khóa của thực thể chủ.

Phân tích: ER diagram mô tả các entity set, thuộc tính và relationship set, cùng các ràng buộc cần thiết. Nó tập trung vào đối tượng dữ liệu và liên kết giữa chúng, không mô tả luồng xử lý/chương trình như sơ đồ quy trình.

Đối chiếu phương án: A. Always add entity name to the attribute name to make the meaning clear.; C. Merge both entities; D. Ignore the attribute names. Các lựa chọn này không đáp ứng định nghĩa/điều kiện được hỏi; khi ôn, hãy chỉ ra chính xác điểm lệch của từng lựa chọn thay vì học thuộc mỗi chữ cái.
\end{question}

\begin{question}{260}
% [ID gốc: 2.80]
% [Nguồn: DBI202 Order 2 Q80]
When converting an ER diagram to relational schema, how is a one-to-many relationship between entities implemented?
\begin{choices}
\choice Create two tables with separate keys.
\choice Merge both entities into one. We accept the duplicated values in all columns of the “one” side.
\choice Add the primary key of the "one" side as a foreign key in the "many" side.
\choice Create a view combining both tables.
\end{choices}
\textbf{Đáp án:} C
\textbf{Giải thích:} Đáp án theo dữ liệu đề: C — C. Add the primary key of the "one" side as a foreign key in the "many" side..

Lời giải gốc: Khi chuyển đổi quan hệ một-nhiều, ta đưa khóa chính của phía "một" vào phía "nhiều" để làm khóa ngoại. ---

Nền tảng kiến thức: Mô hình ER mô tả yêu cầu ở mức khái niệm: thực thể là đối tượng cần lưu, thuộc tính mô tả thực thể hoặc mối quan hệ, còn lực lượng (cardinality) và participation diễn tả các ràng buộc kết hợp. Khóa của thực thể yếu thường ghép khóa bộ phận với khóa của thực thể chủ.

Phân tích: ER diagram mô tả các entity set, thuộc tính và relationship set, cùng các ràng buộc cần thiết. Nó tập trung vào đối tượng dữ liệu và liên kết giữa chúng, không mô tả luồng xử lý/chương trình như sơ đồ quy trình.

Đối chiếu phương án: A. Create two tables with separate keys.; B. Merge both entities into one. We accept the duplicated values in all columns of the “one” side.; D. Create a view combining both tables.. Các lựa chọn này không đáp ứng định nghĩa/điều kiện được hỏi; khi ôn, hãy chỉ ra chính xác điểm lệch của từng lựa chọn thay vì học thuộc mỗi chữ cái.
\end{question}

\begin{question}{261}
% [ID gốc: 2.81]
% [Nguồn: DBI202 Order 2 Q81]
A real-world thing, distinguishable from other objects, in an ERD is called:
\begin{choices}
\choice Entity
\choice Diagram
\choice Relationship
\choice Primary entity
\end{choices}
\textbf{Đáp án:} A
\textbf{Giải thích:} Đáp án theo dữ liệu đề: A — A. Entity.

Lời giải gốc: Thực thể (Entity) là đối tượng trong thế giới thực có thể phân biệt rõ ràng với các đối tượng khác.

Nền tảng kiến thức: Mô hình ER mô tả yêu cầu ở mức khái niệm: thực thể là đối tượng cần lưu, thuộc tính mô tả thực thể hoặc mối quan hệ, còn lực lượng (cardinality) và participation diễn tả các ràng buộc kết hợp. Khóa của thực thể yếu thường ghép khóa bộ phận với khóa của thực thể chủ.

Phân tích: Khi đọc mô hình ER, tách ba thứ: entity set là loại đối tượng, relationship set là liên kết giữa các loại đối tượng, còn cardinality/participation là ràng buộc số lượng hoặc mức tham gia. Đáp án đúng phải gọi đúng thành phần mà câu hỏi mô tả.

Đối chiếu phương án: B. Diagram; C. Relationship; D. Primary entity. Các lựa chọn này không đáp ứng định nghĩa/điều kiện được hỏi; khi ôn, hãy chỉ ra chính xác điểm lệch của từng lựa chọn thay vì học thuộc mỗi chữ cái.
\end{question}

\begin{question}{262}
% [ID gốc: 2.82]
% [Nguồn: DBI202 Order 2 Q82]
In an ERD, what role does a primary key play?
\begin{choices}
\choice It uniquely identifies each record in a table
\choice It establishes relationships between tables
\choice It determines the data type of attributes
\choice It defines the cardinality of relationships
\end{choices}
\textbf{Đáp án:} A
\textbf{Giải thích:} Đáp án theo dữ liệu đề: A — A. It uniquely identifies each record in a table.

Lời giải gốc: Khóa chính (Primary key) định danh duy nhất từng bản ghi trong bảng dữ liệu.

Nền tảng kiến thức: Mô hình ER mô tả yêu cầu ở mức khái niệm: thực thể là đối tượng cần lưu, thuộc tính mô tả thực thể hoặc mối quan hệ, còn lực lượng (cardinality) và participation diễn tả các ràng buộc kết hợp. Khóa của thực thể yếu thường ghép khóa bộ phận với khóa của thực thể chủ.

Phân tích: Khi đọc mô hình ER, tách ba thứ: entity set là loại đối tượng, relationship set là liên kết giữa các loại đối tượng, còn cardinality/participation là ràng buộc số lượng hoặc mức tham gia. Đáp án đúng phải gọi đúng thành phần mà câu hỏi mô tả.

Đối chiếu phương án: B. It establishes relationships between tables; C. It determines the data type of attributes; D. It defines the cardinality of relationships. Các lựa chọn này không đáp ứng định nghĩa/điều kiện được hỏi; khi ôn, hãy chỉ ra chính xác điểm lệch của từng lựa chọn thay vì học thuộc mỗi chữ cái.
\end{question}

\begin{question}{263}
% [ID gốc: 2.83]
% [Nguồn: DBI202 Order 2 Q83]
An ERD uses three types of principal elements to form relationships:
\begin{choices}
\choice Entity sets, Constraints, and Relationships
\choice Attributes, Constraints, and Relationships
\choice Entity sets, Attributes, and Relationships
\choice Entity sets, Attributes, and Constraints
\end{choices}
\textbf{Đáp án:} C
\textbf{Giải thích:} Đáp án theo dữ liệu đề: C — C. Entity sets, Attributes, and Relationships.

Lời giải gốc: Ba thành phần chính cấu thành sơ đồ thực thể kết hợp (ERD) gồm các tập thực thể (Entity sets), các thuộc tính (Attributes) và các mối quan hệ (Relationships).

Nền tảng kiến thức: Mô hình ER mô tả yêu cầu ở mức khái niệm: thực thể là đối tượng cần lưu, thuộc tính mô tả thực thể hoặc mối quan hệ, còn lực lượng (cardinality) và participation diễn tả các ràng buộc kết hợp. Khóa của thực thể yếu thường ghép khóa bộ phận với khóa của thực thể chủ.

Phân tích: Khi đọc mô hình ER, tách ba thứ: entity set là loại đối tượng, relationship set là liên kết giữa các loại đối tượng, còn cardinality/participation là ràng buộc số lượng hoặc mức tham gia. Đáp án đúng phải gọi đúng thành phần mà câu hỏi mô tả.

Đối chiếu phương án: A. Entity sets, Constraints, and Relationships; B. Attributes, Constraints, and Relationships; D. Entity sets, Attributes, and Constraints. Các lựa chọn này không đáp ứng định nghĩa/điều kiện được hỏi; khi ôn, hãy chỉ ra chính xác điểm lệch của từng lựa chọn thay vì học thuộc mỗi chữ cái.
\end{question}

\begin{question}{264}
% [ID gốc: 2.84]
% [Nguồn: DBI202 Order 2 Q84]
If R is neither many-one from E to F nor from F to E, then we say R is:
\begin{choices}
\choice Many-one
\choice One-one
\choice Many-many
\choice All of the others
\end{choices}
\textbf{Đáp án:} C
\textbf{Giải thích:} Đáp án theo dữ liệu đề: C — C. Many-many.

Lời giải gốc: Nếu quan hệ không phải là nhiều-một theo cả hai chiều thì đó là quan hệ nhiều-nhiều (Many-many).

Nền tảng kiến thức: Mô hình ER mô tả yêu cầu ở mức khái niệm: thực thể là đối tượng cần lưu, thuộc tính mô tả thực thể hoặc mối quan hệ, còn lực lượng (cardinality) và participation diễn tả các ràng buộc kết hợp. Khóa của thực thể yếu thường ghép khóa bộ phận với khóa của thực thể chủ.

Phân tích: Với lựa chọn tổng hợp, không chọn theo cảm tính: kiểm tra từng phát biểu độc lập. “All of the others/above” chỉ đúng khi mọi phát biểu còn lại đúng; “None” chỉ đúng khi không phát biểu nào trong nhóm thỏa yêu cầu. Đáp án lưu cần được đọc cùng các mệnh đề thành phần.

Đối chiếu phương án: A. Many-one; B. One-one; D. All of the others. Các lựa chọn này không đáp ứng định nghĩa/điều kiện được hỏi; khi ôn, hãy chỉ ra chính xác điểm lệch của từng lựa chọn thay vì học thuộc mỗi chữ cái.
\end{question}

\begin{question}{265}
% [ID gốc: 2.94]
% [Nguồn: DBI202 Order 2 Q94]
What is a composite key in an ERD?
\begin{choices}
\choice A key that represents a many-to-many relationship
\choice A key that uniquely identifies each record in a table
\choice A key used to link two entities
\choice A key made up of multiple attributes
\end{choices}
\textbf{Đáp án:} D
\textbf{Giải thích:} Đáp án theo dữ liệu đề: D — D. A key made up of multiple attributes.

Lời giải gốc: Khóa phức hợp (composite key) là khóa được tạo thành từ sự kết hợp của hai hay nhiều thuộc tính.

Nền tảng kiến thức: Mô hình ER mô tả yêu cầu ở mức khái niệm: thực thể là đối tượng cần lưu, thuộc tính mô tả thực thể hoặc mối quan hệ, còn lực lượng (cardinality) và participation diễn tả các ràng buộc kết hợp. Khóa của thực thể yếu thường ghép khóa bộ phận với khóa của thực thể chủ.

Phân tích: Khi đọc mô hình ER, tách ba thứ: entity set là loại đối tượng, relationship set là liên kết giữa các loại đối tượng, còn cardinality/participation là ràng buộc số lượng hoặc mức tham gia. Đáp án đúng phải gọi đúng thành phần mà câu hỏi mô tả.

Đối chiếu phương án: A. A key that represents a many-to-many relationship; B. A key that uniquely identifies each record in a table; C. A key used to link two entities. Các lựa chọn này không đáp ứng định nghĩa/điều kiện được hỏi; khi ôn, hãy chỉ ra chính xác điểm lệch của từng lựa chọn thay vì học thuộc mỗi chữ cái.
\end{question}

\begin{question}{266}
% [ID gốc: 2.96]
% [Nguồn: DBI202 Order 2 Q96]
In relational database design, which of the following is a set of entities of the same type that share the same properties or attributes?
\begin{choices}
\choice Entity Relationship model
\choice Entity set
\choice Field set
\choice Record set
\end{choices}
\textbf{Đáp án:} B
\textbf{Giải thích:} Đáp án theo dữ liệu đề: B — B. Entity set.

Lời giải gốc: Tập thực thể (Entity set) là tập hợp tất cả các thực thể cùng loại có chung các thuộc tính.

Nền tảng kiến thức: Mô hình ER mô tả yêu cầu ở mức khái niệm: thực thể là đối tượng cần lưu, thuộc tính mô tả thực thể hoặc mối quan hệ, còn lực lượng (cardinality) và participation diễn tả các ràng buộc kết hợp. Khóa của thực thể yếu thường ghép khóa bộ phận với khóa của thực thể chủ.

Phân tích: Khi đọc mô hình ER, tách ba thứ: entity set là loại đối tượng, relationship set là liên kết giữa các loại đối tượng, còn cardinality/participation là ràng buộc số lượng hoặc mức tham gia. Đáp án đúng phải gọi đúng thành phần mà câu hỏi mô tả.

Đối chiếu phương án: A. Entity Relationship model; C. Field set; D. Record set. Các lựa chọn này không đáp ứng định nghĩa/điều kiện được hỏi; khi ôn, hãy chỉ ra chính xác điểm lệch của từng lựa chọn thay vì học thuộc mỗi chữ cái.
\end{question}

\begin{question}{267}
% [ID gốc: 2.97]
% [Nguồn: DBI202 Order 2 Q97]
In the entity-relationship model, what is an entity?
\begin{choices}
\choice Relationships
\choice Relationships between objects
\choice Objects and relationships between objects
\choice Data objects
\end{choices}
\textbf{Đáp án:} D
\textbf{Giải thích:} Đáp án theo dữ liệu đề: D — D. Data objects.

Lời giải gốc: Trong mô hình thực thể kết hợp, thực thể đại diện cho một đối tượng có thực trong thế giới thực có thể phân biệt được (data object).

Nền tảng kiến thức: Mô hình ER mô tả yêu cầu ở mức khái niệm: thực thể là đối tượng cần lưu, thuộc tính mô tả thực thể hoặc mối quan hệ, còn lực lượng (cardinality) và participation diễn tả các ràng buộc kết hợp. Khóa của thực thể yếu thường ghép khóa bộ phận với khóa của thực thể chủ.

Phân tích: Entity là một đối tượng/sự vật có thể phân biệt trong miền ứng dụng và cần lưu thông tin; entity set là tập các entity cùng loại. Ví dụ một sinh viên cụ thể là entity, còn Student là entity set/type.

Đối chiếu phương án: A. Relationships; B. Relationships between objects; C. Objects and relationships between objects. Các lựa chọn này không đáp ứng định nghĩa/điều kiện được hỏi; khi ôn, hãy chỉ ra chính xác điểm lệch của từng lựa chọn thay vì học thuộc mỗi chữ cái.
\end{question}

\begin{question}{268}
% [ID gốc: 2.102]
% [Nguồn: DBI202 Order 2 Q102]
Converting an E/R design to a relational database schema includes:
\begin{choices}
\choice Turning each entity set into a relation with the same set of attributes
\choice Replacing a relationship by a relation whose attributes are the keys for the connected entity sets
\choice All of the others
\choice None of the others
\end{choices}
\textbf{Đáp án:} C
\textbf{Giải thích:} Đáp án theo dữ liệu đề: C — C. All of the others.

Lời giải gốc: Quy trình chuyển đổi mô hình E/R sang lược đồ quan hệ bao gồm việc chuyển tập thực thể thành quan hệ và xử lý các mối quan hệ thành bảng trung gian hoặc khóa ngoại.

Nền tảng kiến thức: Mô hình ER mô tả yêu cầu ở mức khái niệm: thực thể là đối tượng cần lưu, thuộc tính mô tả thực thể hoặc mối quan hệ, còn lực lượng (cardinality) và participation diễn tả các ràng buộc kết hợp. Khóa của thực thể yếu thường ghép khóa bộ phận với khóa của thực thể chủ.

Phân tích: Ba thành phần nền tảng của ER là entity set, attribute và relationship set. Entity mô tả loại đối tượng, attribute mô tả dữ liệu của đối tượng/liên kết, relationship biểu diễn kết nối giữa entity set.

Đối chiếu phương án: A. Turning each entity set into a relation with the same set of attributes; B. Replacing a relationship by a relation whose attributes are the keys for the connected entity…; D. None of the others. Các lựa chọn này không đáp ứng định nghĩa/điều kiện được hỏi; khi ôn, hãy chỉ ra chính xác điểm lệch của từng lựa chọn thay vì học thuộc mỗi chữ cái.
\end{question}

\begin{question}{269}
% [ID gốc: 2.104]
% [Nguồn: DBI202 Order 2 Q104]
Which diagram is used to model entities and relationships?
\begin{choices}
\choice UML class diagram
\choice Flowchart
\choice E/R diagram
\choice Dependency graph
\end{choices}
\textbf{Đáp án:} C
\textbf{Giải thích:} Đáp án theo dữ liệu đề: C — C. E/R diagram.

Lời giải gốc: Sơ đồ ER (E/R diagram - Entity-Relationship diagram) là công cụ chuyên dụng dùng để mô hình hóa các thực thể và mối quan hệ giữa chúng.

Nền tảng kiến thức: Mô hình ER mô tả yêu cầu ở mức khái niệm: thực thể là đối tượng cần lưu, thuộc tính mô tả thực thể hoặc mối quan hệ, còn lực lượng (cardinality) và participation diễn tả các ràng buộc kết hợp. Khóa của thực thể yếu thường ghép khóa bộ phận với khóa của thực thể chủ.

Phân tích: Sơ đồ ER/E-R dùng để mô hình hóa entity và relationship ở mức khái niệm. Khi một câu hỏi hỏi quan hệ giữa các entity, đó là liên kết dữ liệu; sơ đồ không nhằm biểu diễn quan hệ giữa các tiến trình thực thi.

Đối chiếu phương án: A. UML class diagram; B. Flowchart; D. Dependency graph. Các lựa chọn này không đáp ứng định nghĩa/điều kiện được hỏi; khi ôn, hãy chỉ ra chính xác điểm lệch của từng lựa chọn thay vì học thuộc mỗi chữ cái.
\end{question}

\begin{question}{270}
% [ID gốc: 2.105]
% [Nguồn: DBI202 Order 2 Q105]
In ER modeling, what does a weak entity require in order to be uniquely identified?
\begin{choices}
\choice A surrogate key
\choice A composite attribute with an identity constraint
\choice A relationship with an identifying entity
\choice A multi-valued attribute
\end{choices}
\textbf{Đáp án:} C
\textbf{Giải thích:} Đáp án theo dữ liệu đề: C — C. A relationship with an identifying entity.

Lời giải gốc: Thực thể yếu (Weak Entity) không có đủ thuộc tính để tự làm khóa chính, nó cần một mối quan hệ định danh (Identifying relationship) kết nối với một thực thể chủ/thực thể mạnh (Identifying/Owner entity).

Nền tảng kiến thức: Mô hình ER mô tả yêu cầu ở mức khái niệm: thực thể là đối tượng cần lưu, thuộc tính mô tả thực thể hoặc mối quan hệ, còn lực lượng (cardinality) và participation diễn tả các ràng buộc kết hợp. Khóa của thực thể yếu thường ghép khóa bộ phận với khóa của thực thể chủ.

Phân tích: ER diagram mô tả các entity set, thuộc tính và relationship set, cùng các ràng buộc cần thiết. Nó tập trung vào đối tượng dữ liệu và liên kết giữa chúng, không mô tả luồng xử lý/chương trình như sơ đồ quy trình.

Đối chiếu phương án: A. A surrogate key; B. A composite attribute with an identity constraint; D. A multi-valued attribute. Các lựa chọn này không đáp ứng định nghĩa/điều kiện được hỏi; khi ôn, hãy chỉ ra chính xác điểm lệch của từng lựa chọn thay vì học thuộc mỗi chữ cái.
\end{question}

\begin{question}{271}
% [ID gốc: 3.18]
% [Nguồn: DBI202 Order 3 Q18]
In the relational model, how are relationships represented between tables?
\begin{choices}
\choice By adding attributes to each table
\choice By creating a new table for each relationship
\choice By adding foreign keys to related tables
\choice By creating a primary key in each table
\end{choices}
\textbf{Đáp án:} C
\textbf{Giải thích:} Đáp án theo dữ liệu đề: C — C. By adding foreign keys to related tables.

Nền tảng kiến thức: Mô hình quan hệ biểu diễn dữ liệu bằng quan hệ (bảng), bộ (hàng), thuộc tính (cột) và miền giá trị. Khóa nhận diện bộ; khóa ngoại và các ràng buộc toàn vẹn ngăn dữ liệu tham chiếu sai hoặc vi phạm miền.

Phân tích: Kiểm tra khóa bằng hai tiêu chí: tập thuộc tính phải xác định duy nhất mọi bộ (tính siêu khóa), và nếu là candidate key thì phải tối thiểu, tức bỏ bất kỳ thuộc tính nào cũng mất tính duy nhất. PRIMARY KEY chọn một candidate key và không nhận NULL.

Đối chiếu phương án: A. By adding attributes to each table; B. By creating a new table for each relationship; D. By creating a primary key in each table. Các lựa chọn này không đáp ứng định nghĩa/điều kiện được hỏi; khi ôn, hãy chỉ ra chính xác điểm lệch của từng lựa chọn thay vì học thuộc mỗi chữ cái.
\end{question}

\begin{question}{272}
% [ID gốc: 15.103]
% [Nguồn: DBI202 Order 15 Q103]
Assume there is a table "Employee" with functional dependencies. You are required to specify the primary key and unique key for the table. Choose the most appropriate option.
\begin{choices}
\choice The primary key is the ID column, the unique key is the TAX_NUMBER column
\choice The primary key is the TAX_NUMBER column, the unique key is the ID column
\choice The primary key includes two columns (ID, TAX_NUMBER) and there should not be a unique key
\choice The primary key is the TAX_NUMBER column, and there should not be a unique key
\end{choices}
\textbf{Đáp án:} A
\textbf{Giải thích:} Đáp án theo dữ liệu đề: A — A. The primary key is the ID column, the unique key is the TAX_NUMBER column.

Lời giải gốc: Cột ID thường làm khóa chính chuẩn xác cho định danh bản ghi, trong khi TAX_NUMBER (mã số thuế) đảm bảo tính duy nhất và làm khóa ứng viên/unique key.

Nền tảng kiến thức: Chuẩn hóa dùng phụ thuộc hàm và khóa để giảm dư thừa cũng như bất thường cập nhật. Tính bao đóng thuộc tính để tìm khóa; 2NF loại phụ thuộc bộ phận vào một phần khóa ghép, 3NF xử lý phụ thuộc không khóa kiểu bắc cầu, còn BCNF yêu cầu mọi vế trái không tầm thường là siêu khóa.

Phân tích: Kiểm tra khóa bằng hai tiêu chí: tập thuộc tính phải xác định duy nhất mọi bộ (tính siêu khóa), và nếu là candidate key thì phải tối thiểu, tức bỏ bất kỳ thuộc tính nào cũng mất tính duy nhất. PRIMARY KEY chọn một candidate key và không nhận NULL.

Đối chiếu phương án: B. The primary key is the TAX_NUMBER column, the unique key is the ID column; C. The primary key includes two columns (ID, TAX_NUMBER) and there should not be a unique key; D. The primary key is the TAX_NUMBER column, and there should not be a unique key. Các lựa chọn này không đáp ứng định nghĩa/điều kiện được hỏi; khi ôn, hãy chỉ ra chính xác điểm lệch của từng lựa chọn thay vì học thuộc mỗi chữ cái.
\end{question}


\newpage
\section*{Chương 5: Algebraic and Logical Query Languages}

\begin{question}{273}
% [ID gốc: 4.1]
% [Nguồn: DBI202 Order 4 Q1]
Consider the law, that holds for set relations:
T intersect(R union S) = (T intersect R) union (T intersect S)
The above law still hold for bag relations?
\begin{choices}
\choice True
\choice False
\end{choices}
\textbf{Đáp án:} B
\textbf{Giải thích:} Đáp án theo dữ liệu đề: B — B. False.

Nền tảng kiến thức: Đại số quan hệ thao tác trên quan hệ: selection lọc hàng, projection chọn cột, join ghép bộ theo điều kiện, còn union/difference/intersection là các phép tập hợp. Cần phân biệt ngữ nghĩa tập hợp với bag/multiset vì bag cho phép bản sao và có quy tắc bội số riêng.

Phân tích: Phép tập hợp trên quan hệ cần so sánh các tuple theo thuộc tính tương ứng; union lấy hợp, intersection lấy phần chung và difference lấy phần có ở vế trái mà không có ở vế phải. Trước tiên kiểm tra tính tương thích schema, sau đó áp phép theo đúng hướng/toán tử.

Đối chiếu phương án: A. True. Các lựa chọn này không đáp ứng định nghĩa/điều kiện được hỏi; khi ôn, hãy chỉ ra chính xác điểm lệch của từng lựa chọn thay vì học thuộc mỗi chữ cái.
\end{question}

\begin{question}{274}
% [ID gốc: 4.3]
% [Nguồn: DBI202 Order 4 Q3]
A tuple which fails to participate in a join is called:
\begin{choices}
\choice dangling tuple
\choice child tuple
\choice parent tuple
\choice neighbor tuple
\end{choices}
\textbf{Đáp án:} A
\textbf{Giải thích:} Đáp án theo dữ liệu đề: A — A. dangling tuple.

Nền tảng kiến thức: Đại số quan hệ thao tác trên quan hệ: selection lọc hàng, projection chọn cột, join ghép bộ theo điều kiện, còn union/difference/intersection là các phép tập hợp. Cần phân biệt ngữ nghĩa tập hợp với bag/multiset vì bag cho phép bản sao và có quy tắc bội số riêng.

Phân tích: Trong quan hệ, mỗi hàng dữ liệu là một tuple (bộ), mỗi cột là attribute (thuộc tính), còn relation là tập/bag các tuple theo một schema. Dùng đúng thuật ngữ này giúp tránh nhầm hàng với trường/cột hoặc toàn bộ bảng.

Đối chiếu phương án: B. child tuple; C. parent tuple; D. neighbor tuple. Các lựa chọn này không đáp ứng định nghĩa/điều kiện được hỏi; khi ôn, hãy chỉ ra chính xác điểm lệch của từng lựa chọn thay vì học thuộc mỗi chữ cái.
\end{question}

\begin{question}{275}
% [ID gốc: 4.5]
% [Nguồn: DBI202 Order 4 Q5]
In the context of databases, a bag (or multiset) is an unordered collection of elements where duplicates are allowed. This is in contrast to a set, which is an unordered collection of elements without duplicates
--------------------
To understand the distinction between sets, bags, and lists, remember that a set has unordered elements, and only one occurrence of each element. A bag allows more than one occurrence of an element, but the elements and their occurrences are unordered. A list allows more than one occurrence of an element, but the occurrences are ordered. Thus, {1, 2, 1} and {2, 1, 1} are the same bag, but (1, 2, 1) and (2, 1, 1) are not the same list.
To apply a selection to bag:
\begin{choices}
\choice We apply the selection condition to each tuple independently.
\choice None of others.
\choice We apply the selection condition to each attribute independently.
\choice All of the other.
\end{choices}
\textbf{Đáp án:} A
\textbf{Giải thích:} Đáp án theo dữ liệu đề: A — A. We apply the selection condition to each tuple independently..

Nền tảng kiến thức: Đại số quan hệ thao tác trên quan hệ: selection lọc hàng, projection chọn cột, join ghép bộ theo điều kiện, còn union/difference/intersection là các phép tập hợp. Cần phân biệt ngữ nghĩa tập hợp với bag/multiset vì bag cho phép bản sao và có quy tắc bội số riêng.

Phân tích: Đọc truy vấn theo các giai đoạn logic: FROM/JOIN tạo dữ liệu nguồn, WHERE lọc từng hàng, GROUP BY tạo nhóm, HAVING lọc nhóm, SELECT xác định biểu thức đầu ra và ORDER BY sắp kết quả. Với phép join, kiểm tra điều kiện ghép và các hàng không khớp trước khi tính kết quả.

Đối chiếu phương án: B. None of others.; C. We apply the selection condition to each attribute independently.; D. All of the other.. Các lựa chọn này không đáp ứng định nghĩa/điều kiện được hỏi; khi ôn, hãy chỉ ra chính xác điểm lệch của từng lựa chọn thay vì học thuộc mỗi chữ cái.
\end{question}

\begin{question}{276}
% [ID gốc: 4.6]
% [Nguồn: DBI202 Order 4 Q6]
Consider the distrubutive law, that holds for set relation:
(R U S)- T =(R-T)U(S-T)
The above distributive still hold for bag relations?
\begin{choices}
\choice TRUE
\choice FALSE
a
\choice (2,3) appears twice in the result.
\choice (2,3) appears three times in the result.
\choice (3,3) appears once in the result.
\choice (2,2) appears once in the result.
\end{choices}
\textbf{Đáp án:} B
\textbf{Giải thích:} Đáp án theo dữ liệu đề: B — B. (2,3) appears three times in the result..

Nền tảng kiến thức: Đại số quan hệ thao tác trên quan hệ: selection lọc hàng, projection chọn cột, join ghép bộ theo điều kiện, còn union/difference/intersection là các phép tập hợp. Cần phân biệt ngữ nghĩa tập hợp với bag/multiset vì bag cho phép bản sao và có quy tắc bội số riêng.

Phân tích: Phép tập hợp trên set loại bản sao và thường đòi hai quan hệ union-compatible. Với bag, cần theo dõi số lần xuất hiện của mỗi bộ; vì vậy không thể tự động áp dụng mọi đẳng thức tập hợp mà bỏ qua multiplicity.

Đối chiếu phương án: A. TRUE; A. (2,3) appears twice in the result.; C. (3,3) appears once in the result.; D. (2,2) appears once in the result.. Các lựa chọn này không đáp ứng định nghĩa/điều kiện được hỏi; khi ôn, hãy chỉ ra chính xác điểm lệch của từng lựa chọn thay vì học thuộc mỗi chữ cái.
\end{question}

\begin{question}{277}
% [ID gốc: 4.18]
% [Nguồn: DBI202 Order 4 Q18]
Choose the correct statement.
\begin{choices}
\choice Outer join behave like join operations except on tuples that do not occur in the join results
\choice Projection behave like product operations except on tuples that do not occur in the join results
\choice Intersection behave like project operations except on tuples that occur in the join results
\choice Union behave like join operations except on tuples that occur in the join results
\end{choices}
\textbf{Đáp án:} A
\textbf{Giải thích:} Đáp án theo dữ liệu đề: A — A. Outer join behave like join operations except on tuples that do not occur in the join results.

Nền tảng kiến thức: Đại số quan hệ thao tác trên quan hệ: selection lọc hàng, projection chọn cột, join ghép bộ theo điều kiện, còn union/difference/intersection là các phép tập hợp. Cần phân biệt ngữ nghĩa tập hợp với bag/multiset vì bag cho phép bản sao và có quy tắc bội số riêng.

Phân tích: Các hàng dữ liệu của một relation được gọi là tuple; header chứa tên thuộc tính không phải tuple dữ liệu. Trong bảng, tuple tương ứng một bản ghi, còn attribute tương ứng cột.

Đối chiếu phương án: B. Projection behave like product operations except on tuples that do not occur in the join resul…; C. Intersection behave like project operations except on tuples that occur in the join results; D. Union behave like join operations except on tuples that occur in the join results. Các lựa chọn này không đáp ứng định nghĩa/điều kiện được hỏi; khi ôn, hãy chỉ ra chính xác điểm lệch của từng lựa chọn thay vì học thuộc mỗi chữ cái.
\end{question}

\begin{question}{278}
% [ID gốc: 4.25]
% [Nguồn: DBI202 Order 4 Q25]
Which operation is not part of basic relational algebra?
\begin{choices}
\choice Projection
\choice Join
\choice Sorting
\choice Selection
\end{choices}
\textbf{Đáp án:} C
\textbf{Giải thích:} Đáp án theo dữ liệu đề: C — C. Sorting.

Lời giải gốc: Sắp xếp (Sorting) là một thao tác hiển thị dữ liệu chứ không thuộc nhóm các phép toán cơ bản cấu thành nên đại số quan hệ (như Chọn σ, Chiếu π, Kết ⨝).

Nền tảng kiến thức: Đại số quan hệ thao tác trên quan hệ: selection lọc hàng, projection chọn cột, join ghép bộ theo điều kiện, còn union/difference/intersection là các phép tập hợp. Cần phân biệt ngữ nghĩa tập hợp với bag/multiset vì bag cho phép bản sao và có quy tắc bội số riêng.

Phân tích: Phép tập hợp trên quan hệ cần so sánh các tuple theo thuộc tính tương ứng; union lấy hợp, intersection lấy phần chung và difference lấy phần có ở vế trái mà không có ở vế phải. Trước tiên kiểm tra tính tương thích schema, sau đó áp phép theo đúng hướng/toán tử.

Đối chiếu phương án: A. Projection; B. Join; D. Selection. Các lựa chọn này không đáp ứng định nghĩa/điều kiện được hỏi; khi ôn, hãy chỉ ra chính xác điểm lệch của từng lựa chọn thay vì học thuộc mỗi chữ cái.
\end{question}

\begin{question}{279}
% [ID gốc: 5.4]
% [Nguồn: DBI202 Order 5 Q4]
Cho hai quan hệ \textbf{R(A, B)} và \textbf{S(B, C, D)}:

$$\textbf{R} \qquad\qquad \textbf{S}$$
$$\begin{array}{|c|c|}
\hline
\textbf{A} & \textbf{B} \\
\hline
10 & \text{x} \\
\hline
30 & \text{y} \\
\hline
50 & \text{z} \\
\hline
70 & \text{t} \\
\hline
\end{array}
\qquad\qquad
\begin{array}{|c|c|c|}
\hline
\textbf{B} & \textbf{C} & \textbf{D} \\
\hline
\text{x} & 40 & 60 \\
\hline
\text{z} & 60 & 80 \\
\hline
\text{z} & 50 & 50 \\
\hline
\text{z} & 70 & 90 \\
\hline
\end{array}$$

Compute the right outer join of R and S, where the condition is: R.B = S.B.
Then identify 2 tuples of R that do not appear in the computed result:
\begin{choices}
\choice (10,x)
\choice (30,y)
\choice (50,z)
\choice (70,t)
\end{choices}
\textbf{Đáp án:} B
\textbf{Giải thích:} Trong phép Right Outer Join (R RIGHT JOIN S ON R.B = S.B), mọi hàng của S được giữ lại. Bảng S chỉ chứa các giá trị B in {x, z}. Do đó, các bộ của R có B = 'y' là (30, y) và B = 't' là (70, t) sẽ không xuất hiện trong kết quả.
\end{question}


\newpage
\section*{Chương 6: The Database Language SQL}

\begin{question}{280}
% [ID gốc: 2.91]
% [Nguồn: DBI202 Order 2 Q91]
A primary key is combined with a foreign key creates:
\begin{choices}
\choice Parent-Child relationship between the tables that connect them.
\choice Many to many relationship between the tables that connect them.
\choice Network model between the tables that connect them.
\choice None of the mentioned.
\end{choices}
\textbf{Đáp án:} A
\textbf{Giải thích:} Đáp án theo dữ liệu đề: A — A. Parent-Child relationship between the tables that connect them..

Lời giải gốc: Sự kết hợp giữa khóa chính (từ bảng cha) và khóa ngoại (ở bảng con) tạo ra mối quan hệ Cha - Con (Parent-Child relationship) giữa các bảng. ---

Nền tảng kiến thức: Mô hình ER mô tả yêu cầu ở mức khái niệm: thực thể là đối tượng cần lưu, thuộc tính mô tả thực thể hoặc mối quan hệ, còn lực lượng (cardinality) và participation diễn tả các ràng buộc kết hợp. Khóa của thực thể yếu thường ghép khóa bộ phận với khóa của thực thể chủ.

Phân tích: Khóa ngoại buộc giá trị tham chiếu phải khớp một khóa được tham chiếu ở quan hệ đích (hoặc NULL nếu ràng buộc cho phép). Đây là toàn vẹn tham chiếu; nó khác tính duy nhất của primary/candidate key.

Đối chiếu phương án: B. Many to many relationship between the tables that connect them.; C. Network model between the tables that connect them.; D. None of the mentioned.. Các lựa chọn này không đáp ứng định nghĩa/điều kiện được hỏi; khi ôn, hãy chỉ ra chính xác điểm lệch của từng lựa chọn thay vì học thuộc mỗi chữ cái.
\end{question}

\begin{question}{281}
% [ID gốc: 3.4]
% [Nguồn: DBI202 Order 3 Q4]
Three basic types of database integrity constraints are:
(a)Entity integrity
(b) Domain integrity
(c) Referential integrity
(d) Primary key integrity
\begin{choices}
\choice (a) and (b) and (c) are true
\choice (a) and (b) and (d) are true
\choice (a) and (c) and (d) are true
\choice (b) and (c) and (d) are true
\end{choices}
\textbf{Đáp án:} A
\textbf{Giải thích:} Đáp án theo dữ liệu đề: A, B, C — A. (a) and (b) and (c) are true; B. (a) and (b) and (d) are true; C. (a) and (c) and (d) are true.

Nền tảng kiến thức: Mô hình quan hệ biểu diễn dữ liệu bằng quan hệ (bảng), bộ (hàng), thuộc tính (cột) và miền giá trị. Khóa nhận diện bộ; khóa ngoại và các ràng buộc toàn vẹn ngăn dữ liệu tham chiếu sai hoặc vi phạm miền.

Phân tích: Toàn vẹn dữ liệu nghĩa là các giá trị vẫn chính xác, hợp lệ và nhất quán với ràng buộc. Khôi phục sau lỗi, xử lý nhiều truy vấn đồng thời và mã hóa bảo mật là các khả năng quan trọng khác của DBMS nhưng không phải định nghĩa trực tiếp của integrity.

Đối chiếu phương án: D. (b) and (c) and (d) are true. Các lựa chọn này không đáp ứng định nghĩa/điều kiện được hỏi; khi ôn, hãy chỉ ra chính xác điểm lệch của từng lựa chọn thay vì học thuộc mỗi chữ cái.
\end{question}

\begin{question}{282}
% [ID gốc: 3.5]
% [Nguồn: DBI202 Order 3 Q5]
A(an) _____ asserts that a value appearing in one relation must also appear in the primary-key component(s) of another relation
\begin{choices}
\choice Unique key constraint
\choice Primary key constraint
\choice Foreign key constraint
\choice Candidate key constraint
\end{choices}
\textbf{Đáp án:} C
\textbf{Giải thích:} Đáp án theo dữ liệu đề: C — C. Foreign key constraint.

Nền tảng kiến thức: Mô hình quan hệ biểu diễn dữ liệu bằng quan hệ (bảng), bộ (hàng), thuộc tính (cột) và miền giá trị. Khóa nhận diện bộ; khóa ngoại và các ràng buộc toàn vẹn ngăn dữ liệu tham chiếu sai hoặc vi phạm miền.

Phân tích: Khóa ngoại buộc giá trị tham chiếu phải khớp một khóa được tham chiếu ở quan hệ đích (hoặc NULL nếu ràng buộc cho phép). Đây là toàn vẹn tham chiếu; nó khác tính duy nhất của primary/candidate key.

Đối chiếu phương án: A. Unique key constraint; B. Primary key constraint; D. Candidate key constraint. Các lựa chọn này không đáp ứng định nghĩa/điều kiện được hỏi; khi ôn, hãy chỉ ra chính xác điểm lệch của từng lựa chọn thay vì học thuộc mỗi chữ cái.
\end{question}

\begin{question}{283}
% [ID gốc: 3.6]
% [Nguồn: DBI202 Order 3 Q6]
Choose the most correct statement about PRIMARY KEY:
\begin{choices}
\choice The PRIMARY KEY constraint uniquely identifies each record in a database table
\choice Primary keys must contain unique values
\choice A primary key column cannot contain NULL values
\choice Each table should have a primary key, and each table can have only ONE primary key
\choice All of the others
\end{choices}
\textbf{Đáp án:} E
\textbf{Giải thích:} Đáp án theo dữ liệu đề: E — E. All of the others.

Nền tảng kiến thức: Mô hình quan hệ biểu diễn dữ liệu bằng quan hệ (bảng), bộ (hàng), thuộc tính (cột) và miền giá trị. Khóa nhận diện bộ; khóa ngoại và các ràng buộc toàn vẹn ngăn dữ liệu tham chiếu sai hoặc vi phạm miền.

Phân tích: Kiểm tra khóa bằng hai tiêu chí: tập thuộc tính phải xác định duy nhất mọi bộ (tính siêu khóa), và nếu là candidate key thì phải tối thiểu, tức bỏ bất kỳ thuộc tính nào cũng mất tính duy nhất. PRIMARY KEY chọn một candidate key và không nhận NULL.

Đối chiếu phương án: A. The PRIMARY KEY constraint uniquely identifies each record in a database table; B. Primary keys must contain unique values; C. A primary key column cannot contain NULL values; D. Each table should have a primary key, and each table can have only ONE primary key. Các lựa chọn này không đáp ứng định nghĩa/điều kiện được hỏi; khi ôn, hãy chỉ ra chính xác điểm lệch của từng lựa chọn thay vì học thuộc mỗi chữ cái.
\end{question}

\begin{question}{284}
% [ID gốc: 3.7]
% [Nguồn: DBI202 Order 3 Q7]
Referential integrity Constraints control relationships between ....
\begin{choices}
\choice attributes in a table
\choice operations of an object
\choice database instances
\choice tables in a database
\end{choices}
\textbf{Đáp án:} D
\textbf{Giải thích:} Đáp án theo dữ liệu đề: D — D. tables in a database.

Nền tảng kiến thức: Mô hình quan hệ biểu diễn dữ liệu bằng quan hệ (bảng), bộ (hàng), thuộc tính (cột) và miền giá trị. Khóa nhận diện bộ; khóa ngoại và các ràng buộc toàn vẹn ngăn dữ liệu tham chiếu sai hoặc vi phạm miền.

Phân tích: Toàn vẹn dữ liệu nghĩa là các giá trị vẫn chính xác, hợp lệ và nhất quán với ràng buộc. Khôi phục sau lỗi, xử lý nhiều truy vấn đồng thời và mã hóa bảo mật là các khả năng quan trọng khác của DBMS nhưng không phải định nghĩa trực tiếp của integrity.

Đối chiếu phương án: A. attributes in a table; B. operations of an object; C. database instances. Các lựa chọn này không đáp ứng định nghĩa/điều kiện được hỏi; khi ôn, hãy chỉ ra chính xác điểm lệch của từng lựa chọn thay vì học thuộc mỗi chữ cái.
\end{question}

\begin{question}{285}
% [ID gốc: 3.8]
% [Nguồn: DBI202 Order 3 Q8]
Database integrity ensures that:
\begin{choices}
\choice data entered into the database is accurate, valid, and consistent
\choice data entered into the database is accurate, and consistent
\choice data entered into the database is valid, and consistent
\choice data entered into the database is accurate, and valid
\end{choices}
\textbf{Đáp án:} A
\textbf{Giải thích:} Đáp án theo dữ liệu đề: A — A. data entered into the database is accurate, valid, and consistent.

Nền tảng kiến thức: Mô hình quan hệ biểu diễn dữ liệu bằng quan hệ (bảng), bộ (hàng), thuộc tính (cột) và miền giá trị. Khóa nhận diện bộ; khóa ngoại và các ràng buộc toàn vẹn ngăn dữ liệu tham chiếu sai hoặc vi phạm miền.

Phân tích: Toàn vẹn dữ liệu nghĩa là các giá trị vẫn chính xác, hợp lệ và nhất quán với ràng buộc. Khôi phục sau lỗi, xử lý nhiều truy vấn đồng thời và mã hóa bảo mật là các khả năng quan trọng khác của DBMS nhưng không phải định nghĩa trực tiếp của integrity.

Đối chiếu phương án: B. data entered into the database is accurate, and consistent; C. data entered into the database is valid, and consistent; D. data entered into the database is accurate, and valid. Các lựa chọn này không đáp ứng định nghĩa/điều kiện được hỏi; khi ôn, hãy chỉ ra chính xác điểm lệch của từng lựa chọn thay vì học thuộc mỗi chữ cái.
\end{question}

\begin{question}{286}
% [ID gốc: 3.9]
% [Nguồn: DBI202 Order 3 Q9]
Cho hai bảng dữ liệu \textbf{ARTISTS} và \textbf{ALBUMS}:

$$\textbf{ARTISTS} \qquad\qquad \textbf{ALBUMS}$$
$$\begin{array}{|c|c|}
\hline
\textbf{Artist\_id} & \textbf{Artist\_Name} \\
\hline
1 & \text{Bono} \\
\hline
2 & \text{Cher} \\
\hline
\end{array}
\qquad\qquad
\begin{array}{|c|c|}
\hline
\textbf{artist\_id} & \textbf{Album\_name} \\
\hline
3 & \text{Eat the rich} \\
\hline
\end{array}$$

Look at the tables above and then choose the correct statement:
\begin{choices}
\choice The above picture is an example of a database that has not enforced referential integrity
\choice The above picture is an example of a database that has not enforced entity integrity
\choice The above picture is an example of a database that has not enforced domain integrity
\choice All of the others
\end{choices}
\textbf{Đáp án:} A
\textbf{Giải thích:} Bảng ALBUMS có artist_id = 3 tham chiếu đến Artist_id trong bảng ARTISTS, nhưng giá trị 3 không tồn tại trong bảng ARTISTS. Đây là một ví dụ về việc cơ sở dữ liệu vi phạm ràng buộc toàn vẹn tham chiếu (referential integrity).
\end{question}

\begin{question}{287}
% [ID gốc: 3.10]
% [Nguồn: DBI202 Order 3 Q10]
When declaring foreign key constraint for relation A thai references relation B. the referenced attribute(s) of the relation B must be declare as_______
\begin{choices}
\choice UNIQUE or PRIMARY KEY
\choice FOREIGN KEY
\choice INDEX KEY
\choice All of the others
\end{choices}
\textbf{Đáp án:} A
\textbf{Giải thích:} Đáp án theo dữ liệu đề: A — A. UNIQUE or PRIMARY KEY.

Nền tảng kiến thức: Mô hình quan hệ biểu diễn dữ liệu bằng quan hệ (bảng), bộ (hàng), thuộc tính (cột) và miền giá trị. Khóa nhận diện bộ; khóa ngoại và các ràng buộc toàn vẹn ngăn dữ liệu tham chiếu sai hoặc vi phạm miền.

Phân tích: Khóa ngoại buộc giá trị tham chiếu phải khớp một khóa được tham chiếu ở quan hệ đích (hoặc NULL nếu ràng buộc cho phép). Đây là toàn vẹn tham chiếu; nó khác tính duy nhất của primary/candidate key.

Đối chiếu phương án: B. FOREIGN KEY; C. INDEX KEY; D. All of the others. Các lựa chọn này không đáp ứng định nghĩa/điều kiện được hỏi; khi ôn, hãy chỉ ra chính xác điểm lệch của từng lựa chọn thay vì học thuộc mỗi chữ cái.
\end{question}

\begin{question}{288}
% [ID gốc: 3.20]
% [Nguồn: DBI202 Order 3 Q20]
Domain constraints, functional dependency and referential integrity are special forms of _____
\begin{choices}
\choice Foreign key
\choice Primary key
\choice Assertion
\choice Referential constraint
\end{choices}
\textbf{Đáp án:} C
\textbf{Giải thích:} Đáp án theo dữ liệu đề: C — C. Assertion.

Nền tảng kiến thức: Mô hình quan hệ biểu diễn dữ liệu bằng quan hệ (bảng), bộ (hàng), thuộc tính (cột) và miền giá trị. Khóa nhận diện bộ; khóa ngoại và các ràng buộc toàn vẹn ngăn dữ liệu tham chiếu sai hoặc vi phạm miền.

Phân tích: Khóa ngoại buộc giá trị tham chiếu phải khớp một khóa được tham chiếu ở quan hệ đích (hoặc NULL nếu ràng buộc cho phép). Đây là toàn vẹn tham chiếu; nó khác tính duy nhất của primary/candidate key.

Đối chiếu phương án: A. Foreign key; B. Primary key; D. Referential constraint. Các lựa chọn này không đáp ứng định nghĩa/điều kiện được hỏi; khi ôn, hãy chỉ ra chính xác điểm lệch của từng lựa chọn thay vì học thuộc mỗi chữ cái.
\end{question}

\begin{question}{289}
% [ID gốc: 3.21]
% [Nguồn: DBI202 Order 3 Q21]
Which of the following is false about a foreign key?
\begin{choices}
\choice Whereas only one foreign key is allowed in a table.
\choice It refers to the field in a table which is the primary key of another table.
\choice It can contain duplicate values and a table in a relational database.
\choice It can also contain NULL values.
\end{choices}
\textbf{Đáp án:} A
\textbf{Giải thích:} Đáp án theo dữ liệu đề: A — A. Whereas only one foreign key is allowed in a table..

Nền tảng kiến thức: Mô hình quan hệ biểu diễn dữ liệu bằng quan hệ (bảng), bộ (hàng), thuộc tính (cột) và miền giá trị. Khóa nhận diện bộ; khóa ngoại và các ràng buộc toàn vẹn ngăn dữ liệu tham chiếu sai hoặc vi phạm miền.

Phân tích: Khóa ngoại buộc giá trị tham chiếu phải khớp một khóa được tham chiếu ở quan hệ đích (hoặc NULL nếu ràng buộc cho phép). Đây là toàn vẹn tham chiếu; nó khác tính duy nhất của primary/candidate key.

Đối chiếu phương án: B. It refers to the field in a table which is the primary key of another table.; C. It can contain duplicate values and a table in a relational database.; D. It can also contain NULL values.. Các lựa chọn này không đáp ứng định nghĩa/điều kiện được hỏi; khi ôn, hãy chỉ ra chính xác điểm lệch của từng lựa chọn thay vì học thuộc mỗi chữ cái.
\end{question}

\begin{question}{290}
% [ID gốc: 3.23]
% [Nguồn: DBI202 Order 3 Q23]
What is the foreign key?
\begin{choices}
\choice It is a relational constraint that enforces referential integrity between two tables
\choice It is a relational constraint that that ensures the values in a specific column or set of columns are unique across all the rows in a table
\choice The key that need to make a data in another table
\choice A collection of data that allows to access and use of that data
\end{choices}
\textbf{Đáp án:} A
\textbf{Giải thích:} Đáp án theo dữ liệu đề: A — A. It is a relational constraint that enforces referential integrity between two tables.

Nền tảng kiến thức: Mô hình quan hệ biểu diễn dữ liệu bằng quan hệ (bảng), bộ (hàng), thuộc tính (cột) và miền giá trị. Khóa nhận diện bộ; khóa ngoại và các ràng buộc toàn vẹn ngăn dữ liệu tham chiếu sai hoặc vi phạm miền.

Phân tích: Khóa ngoại buộc giá trị tham chiếu phải khớp một khóa được tham chiếu ở quan hệ đích (hoặc NULL nếu ràng buộc cho phép). Đây là toàn vẹn tham chiếu; nó khác tính duy nhất của primary/candidate key.

Đối chiếu phương án: B. It is a relational constraint that that ensures the values in a specific column or set of colu…; C. The key that need to make a data in another table; D. A collection of data that allows to access and use of that data. Các lựa chọn này không đáp ứng định nghĩa/điều kiện được hỏi; khi ôn, hãy chỉ ra chính xác điểm lệch của từng lựa chọn thay vì học thuộc mỗi chữ cái.
\end{question}

\begin{question}{291}
% [ID gốc: 3.24]
% [Nguồn: DBI202 Order 3 Q24]
Foreign key is the one in which the
of one relation is referenced in another relation.
\begin{choices}
\choice Foreign key
\choice Primary key
\choice References
\choice Check constraint
\end{choices}
\textbf{Đáp án:} B
\textbf{Giải thích:} Đáp án theo dữ liệu đề: B — B. Primary key.

Nền tảng kiến thức: Mô hình quan hệ biểu diễn dữ liệu bằng quan hệ (bảng), bộ (hàng), thuộc tính (cột) và miền giá trị. Khóa nhận diện bộ; khóa ngoại và các ràng buộc toàn vẹn ngăn dữ liệu tham chiếu sai hoặc vi phạm miền.

Phân tích: Khóa ngoại buộc giá trị tham chiếu phải khớp một khóa được tham chiếu ở quan hệ đích (hoặc NULL nếu ràng buộc cho phép). Đây là toàn vẹn tham chiếu; nó khác tính duy nhất của primary/candidate key.

Đối chiếu phương án: A. Foreign key; C. References; D. Check constraint. Các lựa chọn này không đáp ứng định nghĩa/điều kiện được hỏi; khi ôn, hãy chỉ ra chính xác điểm lệch của từng lựa chọn thay vì học thuộc mỗi chữ cái.
\end{question}

\begin{question}{292}
% [ID gốc: 3.28]
% [Nguồn: DBI202 Order 3 Q28]
The values appearing in given attributes of any tuple in the referencing relation must likewise occur in specified attributes of at least one tuple in the referenced relation, according to ________ constraint integrity
\begin{choices}
\choice Referential
\choice Primary
\choice Referencing
\choice Specific
\end{choices}
\textbf{Đáp án:} A
\textbf{Giải thích:} Đáp án theo dữ liệu đề: A — A. Referential.

Nền tảng kiến thức: Mô hình quan hệ biểu diễn dữ liệu bằng quan hệ (bảng), bộ (hàng), thuộc tính (cột) và miền giá trị. Khóa nhận diện bộ; khóa ngoại và các ràng buộc toàn vẹn ngăn dữ liệu tham chiếu sai hoặc vi phạm miền.

Phân tích: Trong quan hệ, mỗi hàng dữ liệu là một tuple (bộ), mỗi cột là attribute (thuộc tính), còn relation là tập/bag các tuple theo một schema. Dùng đúng thuật ngữ này giúp tránh nhầm hàng với trường/cột hoặc toàn bộ bảng.

Đối chiếu phương án: B. Primary; C. Referencing; D. Specific. Các lựa chọn này không đáp ứng định nghĩa/điều kiện được hỏi; khi ôn, hãy chỉ ra chính xác điểm lệch của từng lựa chọn thay vì học thuộc mỗi chữ cái.
\end{question}

\begin{question}{293}
% [ID gốc: 3.29]
% [Nguồn: DBI202 Order 3 Q29]
Which key does not accept the null value?
\begin{choices}
\choice Unique Key
\choice Primary Key
\choice Foreign Key
\choice Candidate key
\end{choices}
\textbf{Đáp án:} B
\textbf{Giải thích:} Đáp án theo dữ liệu đề: B — B. Primary Key.

Nền tảng kiến thức: Mô hình quan hệ biểu diễn dữ liệu bằng quan hệ (bảng), bộ (hàng), thuộc tính (cột) và miền giá trị. Khóa nhận diện bộ; khóa ngoại và các ràng buộc toàn vẹn ngăn dữ liệu tham chiếu sai hoặc vi phạm miền.

Phân tích: NULL biểu thị thiếu/không biết/không áp dụng, không phải số 0 hay chuỗi rỗng. So sánh với NULL bằng `=` không cho TRUE/FALSE mà cho UNKNOWN; dùng `IS NULL`/`IS NOT NULL`. Trong logic ba giá trị, AND/OR/NOT phải được tính theo cả trạng thái UNKNOWN.

Đối chiếu phương án: A. Unique Key; C. Foreign Key; D. Candidate key. Các lựa chọn này không đáp ứng định nghĩa/điều kiện được hỏi; khi ôn, hãy chỉ ra chính xác điểm lệch của từng lựa chọn thay vì học thuộc mỗi chữ cái.
\end{question}

\begin{question}{294}
% [ID gốc: 3.30]
% [Nguồn: DBI202 Order 3 Q30]
Regarding difference between a primary key vs. unique key, which of the following is false?
\begin{choices}
\choice There can be only one unique key for a table
\choice A primary key is a key that uniquely identifies each record in a table but cannot store NULL values
\choice A unique key prevents duplicate values in a column and can store NULL values.
\choice Only one primary key is allowed to use in a table.
\end{choices}
\textbf{Đáp án:} A
\textbf{Giải thích:} Đáp án theo dữ liệu đề: A — A. There can be only one unique key for a table.

Nền tảng kiến thức: Mô hình quan hệ biểu diễn dữ liệu bằng quan hệ (bảng), bộ (hàng), thuộc tính (cột) và miền giá trị. Khóa nhận diện bộ; khóa ngoại và các ràng buộc toàn vẹn ngăn dữ liệu tham chiếu sai hoặc vi phạm miền.

Phân tích: Phép tập hợp trên quan hệ cần so sánh các tuple theo thuộc tính tương ứng; union lấy hợp, intersection lấy phần chung và difference lấy phần có ở vế trái mà không có ở vế phải. Trước tiên kiểm tra tính tương thích schema, sau đó áp phép theo đúng hướng/toán tử.

Đối chiếu phương án: B. A primary key is a key that uniquely identifies each record in a table but cannot store NULL v…; C. A unique key prevents duplicate values in a column and can store NULL values.; D. Only one primary key is allowed to use in a table.. Các lựa chọn này không đáp ứng định nghĩa/điều kiện được hỏi; khi ôn, hãy chỉ ra chính xác điểm lệch của từng lựa chọn thay vì học thuộc mỗi chữ cái.
\end{question}

\begin{question}{295}
% [ID gốc: 3.33]
% [Nguồn: DBI202 Order 3 Q33]
What is the purpose of a foreign key in an ERD?
\begin{choices}
\choice It defines the data type of attributes
\choice It determines the cardinality of relationships
\choice It serves as the primary identifier for a table
\choice It establishes a link between two tables
\end{choices}
\textbf{Đáp án:} D
\textbf{Giải thích:} Đáp án theo dữ liệu đề: D — D. It establishes a link between two tables.

Nền tảng kiến thức: Mô hình quan hệ biểu diễn dữ liệu bằng quan hệ (bảng), bộ (hàng), thuộc tính (cột) và miền giá trị. Khóa nhận diện bộ; khóa ngoại và các ràng buộc toàn vẹn ngăn dữ liệu tham chiếu sai hoặc vi phạm miền.

Phân tích: Khi đọc mô hình ER, tách ba thứ: entity set là loại đối tượng, relationship set là liên kết giữa các loại đối tượng, còn cardinality/participation là ràng buộc số lượng hoặc mức tham gia. Đáp án đúng phải gọi đúng thành phần mà câu hỏi mô tả.

Đối chiếu phương án: A. It defines the data type of attributes; B. It determines the cardinality of relationships; C. It serves as the primary identifier for a table. Các lựa chọn này không đáp ứng định nghĩa/điều kiện được hỏi; khi ôn, hãy chỉ ra chính xác điểm lệch của từng lựa chọn thay vì học thuộc mỗi chữ cái.
\end{question}

\begin{question}{296}
% [ID gốc: 3.36]
% [Nguồn: DBI202 Order 3 Q36]
In the context of a DBMS, what is meant by "data integrity"?
\begin{choices}
\choice The ability to perform multiple queries simultaneously
\choice Ensuring data is consistent, accurate, and reliable
\choice The ability to recover data after a failure
\choice Encrypting data to protect it from unauthorized access
\end{choices}
\textbf{Đáp án:} B
\textbf{Giải thích:} Đáp án theo dữ liệu đề: B — B. Ensuring data is consistent, accurate, and reliable.

Lời giải gốc: Toàn vẹn dữ liệu (Data integrity) trong hệ quản trị cơ sở dữ liệu (DBMS) là việc đảm bảo dữ liệu luôn chính xác, nhất quán và đáng tin cậy trong suốt vòng đời của nó.

Nền tảng kiến thức: Mô hình quan hệ biểu diễn dữ liệu bằng quan hệ (bảng), bộ (hàng), thuộc tính (cột) và miền giá trị. Khóa nhận diện bộ; khóa ngoại và các ràng buộc toàn vẹn ngăn dữ liệu tham chiếu sai hoặc vi phạm miền.

Phân tích: Toàn vẹn dữ liệu nghĩa là các giá trị vẫn chính xác, hợp lệ và nhất quán với ràng buộc. Khôi phục sau lỗi, xử lý nhiều truy vấn đồng thời và mã hóa bảo mật là các khả năng quan trọng khác của DBMS nhưng không phải định nghĩa trực tiếp của integrity.

Đối chiếu phương án: A. The ability to perform multiple queries simultaneously; C. The ability to recover data after a failure; D. Encrypting data to protect it from unauthorized access. Các lựa chọn này không đáp ứng định nghĩa/điều kiện được hỏi; khi ôn, hãy chỉ ra chính xác điểm lệch của từng lựa chọn thay vì học thuộc mỗi chữ cái.
\end{question}

\begin{question}{297}
% [ID gốc: 3.39]
% [Nguồn: DBI202 Order 3 Q39]
In the relational model, what is the main purpose of using a primary key in a table?
\begin{choices}
\choice To allow duplicate records
\choice To identify each row uniquely
\choice To store data temporarily
\choice To improve database aesthetics
\end{choices}
\textbf{Đáp án:} B
\textbf{Giải thích:} Đáp án theo dữ liệu đề: B — B. To identify each row uniquely.

Lời giải gốc: Mục đích chính của Khóa chính (Primary Key) là để định danh duy nhất mỗi hàng (bản ghi) trong một bảng cơ sở dữ liệu.

Nền tảng kiến thức: Mô hình quan hệ biểu diễn dữ liệu bằng quan hệ (bảng), bộ (hàng), thuộc tính (cột) và miền giá trị. Khóa nhận diện bộ; khóa ngoại và các ràng buộc toàn vẹn ngăn dữ liệu tham chiếu sai hoặc vi phạm miền.

Phân tích: Kiểm tra khóa bằng hai tiêu chí: tập thuộc tính phải xác định duy nhất mọi bộ (tính siêu khóa), và nếu là candidate key thì phải tối thiểu, tức bỏ bất kỳ thuộc tính nào cũng mất tính duy nhất. PRIMARY KEY chọn một candidate key và không nhận NULL.

Đối chiếu phương án: A. To allow duplicate records; C. To store data temporarily; D. To improve database aesthetics. Các lựa chọn này không đáp ứng định nghĩa/điều kiện được hỏi; khi ôn, hãy chỉ ra chính xác điểm lệch của từng lựa chọn thay vì học thuộc mỗi chữ cái.
\end{question}

\begin{question}{298}
% [ID gốc: 3.42]
% [Nguồn: DBI202 Order 3 Q42]
Choose the correct statement about the impact of declaring a key for a relation R using PRIMARY KEY or UNIQUE.
\begin{choices}
\choice PRIMARY KEY: does not allow NULL values. UNIQUE: allows NULL values.
\choice PRIMARY KEY: allows NULL values. UNIQUE: does not allow NULL values.
\choice PRIMARY KEY: does not allow NULL values. UNIQUE: does not allow NULL values.
\choice PRIMARY KEY: allows NULL values. UNIQUE: allows NULL values.
\end{choices}
\textbf{Đáp án:} A
\textbf{Giải thích:} Đáp án theo dữ liệu đề: A — A. PRIMARY KEY: does not allow NULL values. UNIQUE: allows NULL values..

Lời giải gốc: Khóa chính (PRIMARY KEY) bắt buộc có giá trị duy nhất và không cho phép giá trị NULL, trong khi khóa độc nhất (UNIQUE) đảm bảo tính duy nhất nhưng vẫn cho phép chứa giá trị NULL (một giá trị NULL trong SQL Server).

Nền tảng kiến thức: Mô hình quan hệ biểu diễn dữ liệu bằng quan hệ (bảng), bộ (hàng), thuộc tính (cột) và miền giá trị. Khóa nhận diện bộ; khóa ngoại và các ràng buộc toàn vẹn ngăn dữ liệu tham chiếu sai hoặc vi phạm miền.

Phân tích: Khóa là thuộc tính hoặc tập thuộc tính dùng để nhận diện duy nhất một bộ/hàng. Khóa có thể gồm một hay nhiều cột, nên không bắt buộc luôn là cột đơn hoặc luôn gồm nhiều cột; nó cũng không “nhận diện một cột”.

Đối chiếu phương án: B. PRIMARY KEY: allows NULL values. UNIQUE: does not allow NULL values.; C. PRIMARY KEY: does not allow NULL values. UNIQUE: does not allow NULL values.; D. PRIMARY KEY: allows NULL values. UNIQUE: allows NULL values.. Các lựa chọn này không đáp ứng định nghĩa/điều kiện được hỏi; khi ôn, hãy chỉ ra chính xác điểm lệch của từng lựa chọn thay vì học thuộc mỗi chữ cái.
\end{question}

\begin{question}{299}
% [ID gốc: 3.45]
% [Nguồn: DBI202 Order 3 Q45]
Which two statements about foreign key constraints are true? (Choose two correct answers)
\begin{choices}
\choice They enforce referential integrity
\choice They allow orphaned rows in child tables
\choice They can be defined using ALTER TABLE
\choice They replace the need for primary keys
\end{choices}
\textbf{Đáp án:} AC
\textbf{Giải thích:} Đáp án theo dữ liệu đề: A — A. They enforce referential integrity.

Lời giải gốc: Khóa ngoại (Foreign Key) là công cụ chính để duy trì tính toàn vẹn tham chiếu (referential integrity) giữa các bảng. Nó có thể được tạo ra lúc định nghĩa bảng (CREATE TABLE) hoặc bổ sung sau này thông qua lệnh ALTER TABLE.

Nền tảng kiến thức: Mô hình quan hệ biểu diễn dữ liệu bằng quan hệ (bảng), bộ (hàng), thuộc tính (cột) và miền giá trị. Khóa nhận diện bộ; khóa ngoại và các ràng buộc toàn vẹn ngăn dữ liệu tham chiếu sai hoặc vi phạm miền.

Phân tích: Khóa ngoại buộc giá trị tham chiếu phải khớp một khóa được tham chiếu ở quan hệ đích (hoặc NULL nếu ràng buộc cho phép). Đây là toàn vẹn tham chiếu; nó khác tính duy nhất của primary/candidate key.

Đối chiếu phương án: B. They allow orphaned rows in child tables; C. They can be defined using ALTER TABLE; D. They replace the need for primary keys. Các lựa chọn này không đáp ứng định nghĩa/điều kiện được hỏi; khi ôn, hãy chỉ ra chính xác điểm lệch của từng lựa chọn thay vì học thuộc mỗi chữ cái.
\end{question}

\begin{question}{300}
% [ID gốc: 3.46]
% [Nguồn: DBI202 Order 3 Q46]
Which statement is used to disable a foreign key constraint named fk_name during INSERT and UPDATE transactions in table B?
\begin{choices}
\choice ALTER TABLE B DROP CONSTRAINT fk_name
\choice ALTER TABLE B DISABLE CONSTRAINT fk_name
\choice ALTER TABLE B DISABLE fk_name
\choice ALTER TABLE B NOCHECK CONSTRAINT fk_name
\end{choices}
\textbf{Đáp án:} D
\textbf{Giải thích:} Đáp án theo dữ liệu đề: D — D. ALTER TABLE B NOCHECK CONSTRAINT fk_name.

Lời giải gốc: Trong SQL Server, cú pháp để vô hiệu hóa tạm thời ràng buộc kiểm tra/khóa ngoại mà không xóa nó là ALTER TABLE table_name CHECK CONSTRAINT constraint_name hoặc NOCHECK CONSTRAINT constraint_name.

Nền tảng kiến thức: Mô hình quan hệ biểu diễn dữ liệu bằng quan hệ (bảng), bộ (hàng), thuộc tính (cột) và miền giá trị. Khóa nhận diện bộ; khóa ngoại và các ràng buộc toàn vẹn ngăn dữ liệu tham chiếu sai hoặc vi phạm miền.

Phân tích: Khóa ngoại buộc giá trị tham chiếu phải khớp một khóa được tham chiếu ở quan hệ đích (hoặc NULL nếu ràng buộc cho phép). Đây là toàn vẹn tham chiếu; nó khác tính duy nhất của primary/candidate key.

Đối chiếu phương án: A. ALTER TABLE B DROP CONSTRAINT fk_name; B. ALTER TABLE B DISABLE CONSTRAINT fk_name; C. ALTER TABLE B DISABLE fk_name. Các lựa chọn này không đáp ứng định nghĩa/điều kiện được hỏi; khi ôn, hãy chỉ ra chính xác điểm lệch của từng lựa chọn thay vì học thuộc mỗi chữ cái.
\end{question}

\begin{question}{301}
% [ID gốc: 3.47]
% [Nguồn: DBI202 Order 3 Q47]
Which constraint ensures that the values of certain attributes in a relation uniquely identify tuples in that relation?
\begin{choices}
\choice Foreign key constraint
\choice Primary key constraint
\choice Check constraint
\choice Unique constraint
\end{choices}
\textbf{Đáp án:} B
\textbf{Giải thích:} Đáp án theo dữ liệu đề: B — B. Primary key constraint.

Lời giải gốc: Ràng buộc khóa chính (Primary key constraint) đảm bảo các giá trị trong thuộc tính xác định duy nhất từng bộ dữ liệu trong quan hệ và không chứa giá trị NULL.

Nền tảng kiến thức: Mô hình quan hệ biểu diễn dữ liệu bằng quan hệ (bảng), bộ (hàng), thuộc tính (cột) và miền giá trị. Khóa nhận diện bộ; khóa ngoại và các ràng buộc toàn vẹn ngăn dữ liệu tham chiếu sai hoặc vi phạm miền.

Phân tích: Khóa là thuộc tính hoặc tập thuộc tính dùng để nhận diện duy nhất một bộ/hàng. Khóa có thể gồm một hay nhiều cột, nên không bắt buộc luôn là cột đơn hoặc luôn gồm nhiều cột; nó cũng không “nhận diện một cột”.

Đối chiếu phương án: A. Foreign key constraint; C. Check constraint; D. Unique constraint. Các lựa chọn này không đáp ứng định nghĩa/điều kiện được hỏi; khi ôn, hãy chỉ ra chính xác điểm lệch của từng lựa chọn thay vì học thuộc mỗi chữ cái.
\end{question}

\begin{question}{302}
% [ID gốc: 3.50]
% [Nguồn: DBI202 Order 3 Q50]
In the relational model, how are relationships represented between tables?
\begin{choices}
\choice By adding attributes to each table
\choice By creating a new table for each relationship
\choice By adding foreign keys to related tables
\choice By creating a primary key in each table
\end{choices}
\textbf{Đáp án:} C
\textbf{Giải thích:} Đáp án theo dữ liệu đề: C — C. By adding foreign keys to related tables.

Lời giải gốc: Mối quan hệ giữa các bảng trong mô hình quan hệ được thiết lập thông qua việc sử dụng các khóa ngoại (Foreign keys).

Nền tảng kiến thức: Mô hình quan hệ biểu diễn dữ liệu bằng quan hệ (bảng), bộ (hàng), thuộc tính (cột) và miền giá trị. Khóa nhận diện bộ; khóa ngoại và các ràng buộc toàn vẹn ngăn dữ liệu tham chiếu sai hoặc vi phạm miền.

Phân tích: Kiểm tra khóa bằng hai tiêu chí: tập thuộc tính phải xác định duy nhất mọi bộ (tính siêu khóa), và nếu là candidate key thì phải tối thiểu, tức bỏ bất kỳ thuộc tính nào cũng mất tính duy nhất. PRIMARY KEY chọn một candidate key và không nhận NULL.

Đối chiếu phương án: A. By adding attributes to each table; B. By creating a new table for each relationship; D. By creating a primary key in each table. Các lựa chọn này không đáp ứng định nghĩa/điều kiện được hỏi; khi ôn, hãy chỉ ra chính xác điểm lệch của từng lựa chọn thay vì học thuộc mỗi chữ cái.
\end{question}

\begin{question}{303}
% [ID gốc: 3.55]
% [Nguồn: DBI202 Order 3 Q55]
A primary key combined with a foreign key creates:
\begin{choices}
\choice A parent-child relationship between the tables that connect them
\choice A many-to-many relationship between the tables that connect them
\choice A network model between the tables that connect them
\choice None of the mentioned
\end{choices}
\textbf{Đáp án:} A
\textbf{Giải thích:} Đáp án theo dữ liệu đề: A — A. A parent-child relationship between the tables that connect them.

Lời giải gốc: Việc kết hợp khóa chính (ở bảng cha) và khóa ngoại (ở bảng con) tạo ra mối quan hệ cha-con (parent-child relationship) giữa hai bảng.

Nền tảng kiến thức: Mô hình quan hệ biểu diễn dữ liệu bằng quan hệ (bảng), bộ (hàng), thuộc tính (cột) và miền giá trị. Khóa nhận diện bộ; khóa ngoại và các ràng buộc toàn vẹn ngăn dữ liệu tham chiếu sai hoặc vi phạm miền.

Phân tích: Khóa ngoại buộc giá trị tham chiếu phải khớp một khóa được tham chiếu ở quan hệ đích (hoặc NULL nếu ràng buộc cho phép). Đây là toàn vẹn tham chiếu; nó khác tính duy nhất của primary/candidate key.

Đối chiếu phương án: B. A many-to-many relationship between the tables that connect them; C. A network model between the tables that connect them; D. None of the mentioned. Các lựa chọn này không đáp ứng định nghĩa/điều kiện được hỏi; khi ôn, hãy chỉ ra chính xác điểm lệch của từng lựa chọn thay vì học thuộc mỗi chữ cái.
\end{question}

\begin{question}{304}
% [ID gốc: 3.59]
% [Nguồn: DBI202 Order 3 Q59]
Which of the following statements is/are TRUE? (1) A column with a FOREIGN KEY constraint does not allow duplicate values. (2) A FOREIGN KEY constraint can be defined to refer to columns with the UNIQUE constraint in another table.
\begin{choices}
\choice Both statements are True
\choice Both statements are False
\choice Only statement 1 is True
\choice Only statement 2 is True
\end{choices}
\textbf{Đáp án:} D
\textbf{Giải thích:} Đáp án theo dữ liệu đề: D — D. Only statement 2 is True.

Lời giải gốc: Khóa ngoại (FOREIGN KEY) hoàn toàn cho phép các giá trị trùng lặp và có thể tham chiếu đến cột có khóa chính hoặc cột có ràng buộc UNIQUE ở bảng khác. Do đó phát biểu 1 sai, phát biểu 2 đúng.

Nền tảng kiến thức: Mô hình quan hệ biểu diễn dữ liệu bằng quan hệ (bảng), bộ (hàng), thuộc tính (cột) và miền giá trị. Khóa nhận diện bộ; khóa ngoại và các ràng buộc toàn vẹn ngăn dữ liệu tham chiếu sai hoặc vi phạm miền.

Phân tích: Khóa ngoại buộc giá trị tham chiếu phải khớp một khóa được tham chiếu ở quan hệ đích (hoặc NULL nếu ràng buộc cho phép). Đây là toàn vẹn tham chiếu; nó khác tính duy nhất của primary/candidate key.

Đối chiếu phương án: A. Both statements are True; B. Both statements are False; C. Only statement 1 is True. Các lựa chọn này không đáp ứng định nghĩa/điều kiện được hỏi; khi ôn, hãy chỉ ra chính xác điểm lệch của từng lựa chọn thay vì học thuộc mỗi chữ cái.
\end{question}

\begin{question}{305}
% [ID gốc: 3.61]
% [Nguồn: DBI202 Order 3 Q61]
There is a foreign-key F in Relation1 that references set of attributes G of Relation2. Suppose a tuple t of Relation1 has non-NULL values in all the attributes of F, call the list of t's values in these attributes t[F], then in Relation2 there must be some tuple that ___ with t[F] on the attributes G.
\begin{choices}
\choice agrees
\choice disagrees
\choice is not relative
\choice differs with
\end{choices}
\textbf{Đáp án:} A
\textbf{Giải thích:} Đáp án theo dữ liệu đề: A — A. agrees.

Lời giải gốc: Khái niệm ràng buộc toàn vẹn tham chiếu (Referential integrity): Giá trị của một khóa ngoại F (nếu không null) bắt buộc phải trùng khớp/khớp giá trị (agrees) với giá trị của thuộc tính G tương ứng đang làm khóa chính ở bảng được tham chiếu (Relation2).

Nền tảng kiến thức: Mô hình quan hệ biểu diễn dữ liệu bằng quan hệ (bảng), bộ (hàng), thuộc tính (cột) và miền giá trị. Khóa nhận diện bộ; khóa ngoại và các ràng buộc toàn vẹn ngăn dữ liệu tham chiếu sai hoặc vi phạm miền.

Phân tích: NULL biểu thị thiếu/không biết/không áp dụng, không phải số 0 hay chuỗi rỗng. So sánh với NULL bằng `=` không cho TRUE/FALSE mà cho UNKNOWN; dùng `IS NULL`/`IS NOT NULL`. Trong logic ba giá trị, AND/OR/NOT phải được tính theo cả trạng thái UNKNOWN.

Đối chiếu phương án: B. disagrees; C. is not relative; D. differs with. Các lựa chọn này không đáp ứng định nghĩa/điều kiện được hỏi; khi ôn, hãy chỉ ra chính xác điểm lệch của từng lựa chọn thay vì học thuộc mỗi chữ cái.
\end{question}

\begin{question}{306}
% [ID gốc: 3.64]
% [Nguồn: DBI202 Order 3 Q64]
To ensure that each row in a table refers to a valid entry in another table, which constraint should be used?
\begin{choices}
\choice CHECK
\choice UNIQUE
\choice PRIMARY KEY
\choice FOREIGN KEY
\end{choices}
\textbf{Đáp án:} D
\textbf{Giải thích:} Đáp án theo dữ liệu đề: D — D. FOREIGN KEY.

Lời giải gốc: Ràng buộc khóa ngoại (FOREIGN KEY) đảm bảo tính toàn vẹn tham chiếu giữa các bảng.

Nền tảng kiến thức: Mô hình quan hệ biểu diễn dữ liệu bằng quan hệ (bảng), bộ (hàng), thuộc tính (cột) và miền giá trị. Khóa nhận diện bộ; khóa ngoại và các ràng buộc toàn vẹn ngăn dữ liệu tham chiếu sai hoặc vi phạm miền.

Phân tích: Khóa ngoại buộc giá trị tham chiếu phải khớp một khóa được tham chiếu ở quan hệ đích (hoặc NULL nếu ràng buộc cho phép). Đây là toàn vẹn tham chiếu; nó khác tính duy nhất của primary/candidate key.

Đối chiếu phương án: A. CHECK; B. UNIQUE; C. PRIMARY KEY. Các lựa chọn này không đáp ứng định nghĩa/điều kiện được hỏi; khi ôn, hãy chỉ ra chính xác điểm lệch của từng lựa chọn thay vì học thuộc mỗi chữ cái.
\end{question}

\begin{question}{307}
% [ID gốc: 3.66]
% [Nguồn: DBI202 Order 3 Q66]
In a relational DBMS, what is the difference between PRIMARY KEY and UNIQUE KEY?
\begin{choices}
\choice UNIQUE KEY doesn't accept the NULL value but PRIMARY KEY does
\choice UNIQUE KEY and PRIMARY KEY are the same
\choice A table can have more than one PRIMARY KEY constraint but only one UNIQUE KEY
\choice A table can have more than one UNIQUE KEY constraint but only one PRIMARY KEY
\end{choices}
\textbf{Đáp án:} D
\textbf{Giải thích:} Đáp án theo dữ liệu đề: D — D. A table can have more than one UNIQUE KEY constraint but only one PRIMARY KEY.

Lời giải gốc: Một bảng chỉ có duy nhất 1 Primary Key nhưng có thể có nhiều ràng buộc Unique Key. Ngoài ra, Primary Key không chấp nhận giá trị NULL còn Unique Key chấp nhận giá trị NULL (tùy thuộc vào RDBMS).

Nền tảng kiến thức: Mô hình quan hệ biểu diễn dữ liệu bằng quan hệ (bảng), bộ (hàng), thuộc tính (cột) và miền giá trị. Khóa nhận diện bộ; khóa ngoại và các ràng buộc toàn vẹn ngăn dữ liệu tham chiếu sai hoặc vi phạm miền.

Phân tích: Phép tập hợp trên quan hệ cần so sánh các tuple theo thuộc tính tương ứng; union lấy hợp, intersection lấy phần chung và difference lấy phần có ở vế trái mà không có ở vế phải. Trước tiên kiểm tra tính tương thích schema, sau đó áp phép theo đúng hướng/toán tử.

Đối chiếu phương án: A. UNIQUE KEY doesn't accept the NULL value but PRIMARY KEY does; B. UNIQUE KEY and PRIMARY KEY are the same; C. A table can have more than one PRIMARY KEY constraint but only one UNIQUE KEY. Các lựa chọn này không đáp ứng định nghĩa/điều kiện được hỏi; khi ôn, hãy chỉ ra chính xác điểm lệch của từng lựa chọn thay vì học thuộc mỗi chữ cái.
\end{question}

\begin{question}{308}
% [ID gốc: 3.67]
% [Nguồn: DBI202 Order 3 Q67]
Which key does not accept a null value?
\begin{choices}
\choice Unique Key
\choice Primary Key
\choice Foreign Key
\choice Candidate Key
\end{choices}
\textbf{Đáp án:} B
\textbf{Giải thích:} Đáp án theo dữ liệu đề: B — B. Primary Key.

Lời giải gốc: Khóa chính (Primary Key) bắt buộc phải có giá trị duy nhất và không được phép chứa giá trị NULL.

Nền tảng kiến thức: Mô hình quan hệ biểu diễn dữ liệu bằng quan hệ (bảng), bộ (hàng), thuộc tính (cột) và miền giá trị. Khóa nhận diện bộ; khóa ngoại và các ràng buộc toàn vẹn ngăn dữ liệu tham chiếu sai hoặc vi phạm miền.

Phân tích: NULL biểu thị thiếu/không biết/không áp dụng, không phải số 0 hay chuỗi rỗng. So sánh với NULL bằng `=` không cho TRUE/FALSE mà cho UNKNOWN; dùng `IS NULL`/`IS NOT NULL`. Trong logic ba giá trị, AND/OR/NOT phải được tính theo cả trạng thái UNKNOWN.

Đối chiếu phương án: A. Unique Key; C. Foreign Key; D. Candidate Key. Các lựa chọn này không đáp ứng định nghĩa/điều kiện được hỏi; khi ôn, hãy chỉ ra chính xác điểm lệch của từng lựa chọn thay vì học thuộc mỗi chữ cái.
\end{question}

\begin{question}{309}
% [ID gốc: 4.22]
% [Nguồn: DBI202 Order 4 Q22]
In SQL, the statement SELECT * FROM R, S is equivalent to which of the following?
\begin{choices}
\choice SELECT * FROM R INNER JOIN S
\choice SELECT * FROM R NATURAL JOIN S
\choice SELECT * FROM R CROSS JOIN S
\choice SELECT * FROM R UNION JOIN S
\end{choices}
\textbf{Đáp án:} C
\textbf{Giải thích:} Đáp án theo dữ liệu đề: C — C. SELECT * FROM R CROSS JOIN S.

Lời giải gốc: Cú pháp liệt kê các bảng cách nhau bằng dấu phẩy trong mệnh đề FROM tương đương với phép kết chéo (CROSS JOIN) giữa hai bảng.

Nền tảng kiến thức: Đại số quan hệ thao tác trên quan hệ: selection lọc hàng, projection chọn cột, join ghép bộ theo điều kiện, còn union/difference/intersection là các phép tập hợp. Cần phân biệt ngữ nghĩa tập hợp với bag/multiset vì bag cho phép bản sao và có quy tắc bội số riêng.

Phân tích: Đọc truy vấn theo các giai đoạn logic: FROM/JOIN tạo dữ liệu nguồn, WHERE lọc từng hàng, GROUP BY tạo nhóm, HAVING lọc nhóm, SELECT xác định biểu thức đầu ra và ORDER BY sắp kết quả. Với phép join, kiểm tra điều kiện ghép và các hàng không khớp trước khi tính kết quả.

Đối chiếu phương án: A. SELECT * FROM R INNER JOIN S; B. SELECT * FROM R NATURAL JOIN S; D. SELECT * FROM R UNION JOIN S. Các lựa chọn này không đáp ứng định nghĩa/điều kiện được hỏi; khi ôn, hãy chỉ ra chính xác điểm lệch của từng lựa chọn thay vì học thuộc mỗi chữ cái.
\end{question}

\begin{question}{310}
% [ID gốc: 4.26]
% [Nguồn: DBI202 Order 4 Q26]
Consider the two relations from our running movie database: Movies (title, year, length, genre, studioName, producerC#) MovieExec(name, address, cert#, netWorth) The producerC# component of each Movies tuple must also appear in the cert# component of some MovieExec tuple. How can we express this integrity constraint in relational algebra?
\begin{choices}
\choice \(\pi_{\text{producerC#}}(\text{Movies}) = \pi_{\text{cert#}}(\text{MovieExec})\)
\choice \(\pi_{\text{producerC#}}(\text{Movies}) \subseteq \pi_{\text{cert#}}(\text{MovieExec})\)
\choice \(\pi_{\text{cert#}}(\text{MovieExec}) \subseteq \pi_{\text{producerC#}}(\text{Movies})\)
\choice \(\pi_{\text{producerC#}}(\text{Movies}) \neq \pi_{\text{cert#}}(\text{MovieExec})\)
\end{choices}
\textbf{Đáp án:} B
\textbf{Giải thích:} Đáp án theo dữ liệu đề: B — B. п_(producerC#) (Movies) n_(cert#) (MovieExec).

Lời giải gốc: Ràng buộc khóa ngoại (Referential Integrity Constraint) phát biểu rằng tập hợp tất cả các mã nhà sản xuất trong bảng Movies phải là tập con ($\subseteq$) của tập hợp các chứng chỉ/mã nhà quản lý trong bảng MovieExec.

Nền tảng kiến thức: Đại số quan hệ thao tác trên quan hệ: selection lọc hàng, projection chọn cột, join ghép bộ theo điều kiện, còn union/difference/intersection là các phép tập hợp. Cần phân biệt ngữ nghĩa tập hợp với bag/multiset vì bag cho phép bản sao và có quy tắc bội số riêng.

Phân tích: Phép tập hợp trên quan hệ cần so sánh các tuple theo thuộc tính tương ứng; union lấy hợp, intersection lấy phần chung và difference lấy phần có ở vế trái mà không có ở vế phải. Trước tiên kiểm tra tính tương thích schema, sau đó áp phép theo đúng hướng/toán tử.

Đối chiếu phương án: A. п_(producerC#) (Movies)=n_(cert#) (MovieExec); C. п_(cert#) (MovieExec)n_(producerC#) (Movies); D. n_(producerC#) (Movies)≠n_(cert#) (MovieExec). Các lựa chọn này không đáp ứng định nghĩa/điều kiện được hỏi; khi ôn, hãy chỉ ra chính xác điểm lệch của từng lựa chọn thay vì học thuộc mỗi chữ cái.
\end{question}

\begin{question}{311}
% [ID gốc: 4.32]
% [Nguồn: DBI202 Order 4 Q32]
Which operation is NOT typically associated with relational data manipulation?
\begin{choices}
\choice Deleting attributes
\choice Inserting tuples into a relation
\choice Modifying tuples in a relation
\choice Deleting tuples from a relation
\end{choices}
\textbf{Đáp án:} A
\textbf{Giải thích:} Đáp án theo dữ liệu đề: A — A. Deleting attributes.

Lời giải gốc: Thao tác dữ liệu (DML) tác động lên các bộ (tuples/dòng) của quan hệ như thêm, sửa, xóa dữ liệu; việc xóa cấu trúc cột (attributes) thuộc về nhóm DDL.

Nền tảng kiến thức: Đại số quan hệ thao tác trên quan hệ: selection lọc hàng, projection chọn cột, join ghép bộ theo điều kiện, còn union/difference/intersection là các phép tập hợp. Cần phân biệt ngữ nghĩa tập hợp với bag/multiset vì bag cho phép bản sao và có quy tắc bội số riêng.

Phân tích: DML thao tác trên các tuple: INSERT thêm hàng, UPDATE sửa giá trị hàng và DELETE xóa hàng. Thêm/xóa thuộc tính làm thay đổi schema nên thuộc DDL/ALTER TABLE, không phải xóa dữ liệu hàng.

Đối chiếu phương án: B. Inserting tuples into a relation; C. Modifying tuples in a relation; D. Deleting tuples from a relation. Các lựa chọn này không đáp ứng định nghĩa/điều kiện được hỏi; khi ôn, hãy chỉ ra chính xác điểm lệch của từng lựa chọn thay vì học thuộc mỗi chữ cái.
\end{question}

\begin{question}{312}
% [ID gốc: 5.2]
% [Nguồn: DBI202 Order 5 Q2]
A _____ is a language for defining data structures
\begin{choices}
\choice DDL
\choice DML
\choice DCL
\choice None of the others
\end{choices}
\textbf{Đáp án:} A
\textbf{Giải thích:} Đáp án theo dữ liệu đề: A — A. DDL.

Nền tảng kiến thức: SQL tách truy vấn dữ liệu (DML) khỏi định nghĩa cấu trúc (DDL) và quyền truy cập (DCL). Với SELECT, cần xác định bảng nguồn, điều kiện lọc, phép kết hợp, nhóm và điều kiện nhóm; các ràng buộc, view, thủ tục và trigger đảm nhiệm các vai trò khác nhau.

Phân tích: DDL (Data Definition Language) khai báo/thay đổi cấu trúc như bảng và schema; DML truy vấn/thêm/sửa/xóa dữ liệu; DCL quản lý quyền. Chọn theo đối tượng bị tác động: cấu trúc, các hàng hay quyền truy cập.

Đối chiếu phương án: B. DML; C. DCL; D. None of the others. Các lựa chọn này không đáp ứng định nghĩa/điều kiện được hỏi; khi ôn, hãy chỉ ra chính xác điểm lệch của từng lựa chọn thay vì học thuộc mỗi chữ cái.
\end{question}

\begin{question}{313}
% [ID gốc: 5.3]
% [Nguồn: DBI202 Order 5 Q3]
Which statement is used to remove a relation named R?
\begin{choices}
\choice DROP TABLE R;
\choice REMOVE TABLE R;
\choice DELETE TABLE R;
\choice TRUNCATE TABLE R;
\end{choices}
\textbf{Đáp án:} A
\textbf{Giải thích:} Đáp án theo dữ liệu đề: A — A. DROP TABLE R;.

Nền tảng kiến thức: SQL tách truy vấn dữ liệu (DML) khỏi định nghĩa cấu trúc (DDL) và quyền truy cập (DCL). Với SELECT, cần xác định bảng nguồn, điều kiện lọc, phép kết hợp, nhóm và điều kiện nhóm; các ràng buộc, view, thủ tục và trigger đảm nhiệm các vai trò khác nhau.

Phân tích: DROP TABLE loại bỏ định nghĩa bảng cùng dữ liệu phụ thuộc theo quy tắc DBMS; DELETE xóa các hàng nhưng giữ bảng, còn TRUNCATE xóa nhanh toàn bộ hàng trong khi vẫn giữ cấu trúc. Cần chọn lệnh đúng cấp độ tác động.

Đối chiếu phương án: B. REMOVE TABLE R;; C. DELETE TABLE R;; D. TRUNCATE TABLE R;. Các lựa chọn này không đáp ứng định nghĩa/điều kiện được hỏi; khi ôn, hãy chỉ ra chính xác điểm lệch của từng lựa chọn thay vì học thuộc mỗi chữ cái.
\end{question}

\begin{question}{314}
% [ID gốc: 5.5]
% [Nguồn: DBI202 Order 5 Q5]
Which of the following is NOT a standard aggregation operator?
\begin{choices}
\choice GROUP
\choice SUM
\choice COUNT
\choice AVG
\end{choices}
\textbf{Đáp án:} A
\textbf{Giải thích:} Đáp án theo dữ liệu đề: A — A. GROUP.

Nền tảng kiến thức: SQL tách truy vấn dữ liệu (DML) khỏi định nghĩa cấu trúc (DDL) và quyền truy cập (DCL). Với SELECT, cần xác định bảng nguồn, điều kiện lọc, phép kết hợp, nhóm và điều kiện nhóm; các ràng buộc, view, thủ tục và trigger đảm nhiệm các vai trò khác nhau.

Phân tích: Trong UML, composition biểu thị sở hữu mạnh: vòng đời phần phụ thuộc chủ thể sở hữu; aggregation biểu thị quan hệ whole–part yếu hơn và phần có thể tồn tại độc lập. Hãy đối chiếu phương án với điều kiện vòng đời được nêu.

Đối chiếu phương án: B. SUM; C. COUNT; D. AVG. Các lựa chọn này không đáp ứng định nghĩa/điều kiện được hỏi; khi ôn, hãy chỉ ra chính xác điểm lệch của từng lựa chọn thay vì học thuộc mỗi chữ cái.
\end{question}

\begin{question}{315}
% [ID gốc: 5.6]
% [Nguồn: DBI202 Order 5 Q6]
Choose the correct statement:
\begin{choices}
\choice Aggregate functions perform a calculation on a set of values and return a single value
\choice Aggregate functions perform a calculation on a single value and return a single value
\choice Aggregate functions perform a calculation on a set of values and return a relation
\choice Aggregate functions perform a calculation on a set of values and return a row
\end{choices}
\textbf{Đáp án:} A
\textbf{Giải thích:} Đáp án theo dữ liệu đề: A — A. Aggregate functions perform a calculation on a set of values and return a single value.

Nền tảng kiến thức: SQL tách truy vấn dữ liệu (DML) khỏi định nghĩa cấu trúc (DDL) và quyền truy cập (DCL). Với SELECT, cần xác định bảng nguồn, điều kiện lọc, phép kết hợp, nhóm và điều kiện nhóm; các ràng buộc, view, thủ tục và trigger đảm nhiệm các vai trò khác nhau.

Phân tích: Aggregate function tính một kết quả từ một tập nhiều hàng/giá trị, ví dụ COUNT, SUM, AVG, MIN và MAX. Trong truy vấn nhóm, GROUP BY xác định từng tập đầu vào; HAVING lọc các nhóm sau khi tính tổng hợp.

Đối chiếu phương án: B. Aggregate functions perform a calculation on a single value and return a single value; C. Aggregate functions perform a calculation on a set of values and return a relation; D. Aggregate functions perform a calculation on a set of values and return a row. Các lựa chọn này không đáp ứng định nghĩa/điều kiện được hỏi; khi ôn, hãy chỉ ra chính xác điểm lệch của từng lựa chọn thay vì học thuộc mỗi chữ cái.
\end{question}

\begin{question}{316}
% [ID gốc: 5.7]
% [Nguồn: DBI202 Order 5 Q7]
Cho quan hệ \textbf{S(B, C, D)} với dữ liệu như sau:

$$\begin{array}{|c|c|c|}
\hline
\textbf{B} & \textbf{C} & \textbf{D} \\
\hline
\text{a} & 4 & 6 \\
\hline
\text{c} & 6 & 8 \\
\hline
\text{c} & 5 & 5 \\
\hline
\text{e} & 9 & 9 \\
\hline
\text{f} & 10 & 10 \\
\hline
\text{c} & 7 & 9 \\
\hline
\end{array}$$

Compute the result of the following query:
\begin{codebox}
SELECT SUM(D)
FROM S
WHERE D < 8 AND D > 5;
\end{codebox}
Choose the row that appears in the computed result:
\begin{choices}
\choice 6
\choice 7
\choice 8
\choice 9
\end{choices}
\textbf{Đáp án:} A
\textbf{Giải thích:} Điều kiện WHERE D < 8 AND D > 5 (5 < D < 8) chỉ khớp với hàng duy nhất (a, 4, 6) có D = 6. Do đó SUM(D) = 6.
\end{question}

\begin{question}{317}
% [ID gốc: 5.9]
% [Nguồn: DBI202 Order 5 Q9]
Column A of a relation has the following list of values in the six rows of the table:
NULL, NULL, 10, 10, 20, 30
In SQL Server, which of the following is the correct value of
AVG(DISTINCT A)?
\begin{choices}
\choice (3,1,2)
\choice (1,3,4)
\choice (2,4,6)
\choice (3,2,6)
\end{choices}
\textbf{Đáp án:} C
\textbf{Giải thích:} Đáp án theo dữ liệu đề: C — C. (2,4,6).

Nền tảng kiến thức: SQL tách truy vấn dữ liệu (DML) khỏi định nghĩa cấu trúc (DDL) và quyền truy cập (DCL). Với SELECT, cần xác định bảng nguồn, điều kiện lọc, phép kết hợp, nhóm và điều kiện nhóm; các ràng buộc, view, thủ tục và trigger đảm nhiệm các vai trò khác nhau.

Phân tích: NULL biểu thị thiếu/không biết/không áp dụng, không phải số 0 hay chuỗi rỗng. So sánh với NULL bằng `=` không cho TRUE/FALSE mà cho UNKNOWN; dùng `IS NULL`/`IS NOT NULL`. Trong logic ba giá trị, AND/OR/NOT phải được tính theo cả trạng thái UNKNOWN.

Đối chiếu phương án: A. 15; B. NULL; D. 10 a [Dap an sai, chua fix] Dap an la (3, 1, 6) cơ QN=29 (8513) The table Arc(x,y) currently h…; A. (3,1,2); B. (1,3,4). Các lựa chọn này không đáp ứng định nghĩa/điều kiện được hỏi; khi ôn, hãy chỉ ra chính xác điểm lệch của từng lựa chọn thay vì học thuộc mỗi chữ cái.
\end{question}

\begin{question}{318}
% [ID gốc: 5.10]
% [Nguồn: DBI202 Order 5 Q10]
How we can understand about this statement:
"All aggregate functions are deterministic"?
\begin{choices}
\choice This means aggregate functions return the same value any time they are called by using a specific set of input values
\choice This means aggregate functions return the different values each time they are called by using a specific set of input values
\choice This means aggregate functions return the same value any time they are called by using any set of input values
\choice None of the others
\end{choices}
\textbf{Đáp án:} A
\textbf{Giải thích:} Đáp án theo dữ liệu đề: A — A. This means aggregate functions return the same value any time they are called by using a specific set of input values.

Nền tảng kiến thức: SQL tách truy vấn dữ liệu (DML) khỏi định nghĩa cấu trúc (DDL) và quyền truy cập (DCL). Với SELECT, cần xác định bảng nguồn, điều kiện lọc, phép kết hợp, nhóm và điều kiện nhóm; các ràng buộc, view, thủ tục và trigger đảm nhiệm các vai trò khác nhau.

Phân tích: Aggregate function tính một kết quả từ một tập nhiều hàng/giá trị, ví dụ COUNT, SUM, AVG, MIN và MAX. Trong truy vấn nhóm, GROUP BY xác định từng tập đầu vào; HAVING lọc các nhóm sau khi tính tổng hợp.

Đối chiếu phương án: B. This means aggregate functions return the different values each time they are called by using…; C. This means aggregate functions return the same value any time they are called by using any set…; D. None of the others. Các lựa chọn này không đáp ứng định nghĩa/điều kiện được hỏi; khi ôn, hãy chỉ ra chính xác điểm lệch của từng lựa chọn thay vì học thuộc mỗi chữ cái.
\end{question}

\begin{question}{319}
% [ID gốc: 5.12]
% [Nguồn: DBI202 Order 5 Q12]
How many JOIN conditions are necessary to combine five tables based on common attributes?
\begin{choices}
\choice 1
\choice 2
\choice 3
\choice 4
\end{choices}
\textbf{Đáp án:} D
\textbf{Giải thích:} Đáp án theo dữ liệu đề: D — D. 4.

Nền tảng kiến thức: SQL tách truy vấn dữ liệu (DML) khỏi định nghĩa cấu trúc (DDL) và quyền truy cập (DCL). Với SELECT, cần xác định bảng nguồn, điều kiện lọc, phép kết hợp, nhóm và điều kiện nhóm; các ràng buộc, view, thủ tục và trigger đảm nhiệm các vai trò khác nhau.

Phân tích: Join ghép các bộ thỏa điều kiện. Inner join loại bộ không ghép được; LEFT/RIGHT OUTER JOIN giữ mọi bộ từ phía được chỉ định và điền NULL phía thiếu. Natural join tự ghép các thuộc tính trùng tên và chỉ xuất một cột cho mỗi thuộc tính ghép.

Đối chiếu phương án: A. 1; B. 2; C. 3. Các lựa chọn này không đáp ứng định nghĩa/điều kiện được hỏi; khi ôn, hãy chỉ ra chính xác điểm lệch của từng lựa chọn thay vì học thuộc mỗi chữ cái.
\end{question}

\begin{question}{320}
% [ID gốc: 5.13]
% [Nguồn: DBI202 Order 5 Q13]
The given Boolean expression x = 2 AND z = 10) will have the truth value FALSE when both x = 2 AND z = 10) are FALSE. This means that x > 4 and (y >= 2 AND z = 10) must both be TRUE.
From the given options, only a. x = 5, y = 3, z = 10 satisfies these conditions. In this case, x > 4 is TRUE because x = 5 and (y >= 2 AND z = 10) is also TRUE because y =3 and z=10. So the correct answer is a. x =5 , y=3 ,z=10.
QN=34 (8539) For what values of x, y, and z, including NULL, does the Boolean expression
x = 2 AND z = 10)
have the truth value FALSE? Identify one of those values from the list below
\begin{choices}
\choice x = 5, y = 3, z = 10
\choice x = 3, y = 1, z = 10.
\choice x = NULL, y = 2, z = 10.
\choice x = NULL, y = 3, z = 10.
\end{choices}
\textbf{Đáp án:} A
\textbf{Giải thích:} Đáp án theo dữ liệu đề: A — A. x = 5, y = 3, z = 10.

Nền tảng kiến thức: SQL tách truy vấn dữ liệu (DML) khỏi định nghĩa cấu trúc (DDL) và quyền truy cập (DCL). Với SELECT, cần xác định bảng nguồn, điều kiện lọc, phép kết hợp, nhóm và điều kiện nhóm; các ràng buộc, view, thủ tục và trigger đảm nhiệm các vai trò khác nhau.

Phân tích: NULL biểu thị thiếu/không biết/không áp dụng, không phải số 0 hay chuỗi rỗng. So sánh với NULL bằng `=` không cho TRUE/FALSE mà cho UNKNOWN; dùng `IS NULL`/`IS NOT NULL`. Trong logic ba giá trị, AND/OR/NOT phải được tính theo cả trạng thái UNKNOWN.

Đối chiếu phương án: B. x = 3, y = 1, z = 10.; C. x = NULL, y = 2, z = 10.; D. x = NULL, y = 3, z = 10.. Các lựa chọn này không đáp ứng định nghĩa/điều kiện được hỏi; khi ôn, hãy chỉ ra chính xác điểm lệch của từng lựa chọn thay vì học thuộc mỗi chữ cái.
\end{question}

\begin{question}{321}
% [ID gốc: 5.14]
% [Nguồn: DBI202 Order 5 Q14]
What operator tests column for the absence of data?
\begin{choices}
\choice IS BLANK operator
\choice IS ZERO operator
\choice IS NULL operator
\choice ISNULL operator
\end{choices}
\textbf{Đáp án:} C
\textbf{Giải thích:} Đáp án theo dữ liệu đề: C — C. IS NULL operator.

Nền tảng kiến thức: SQL tách truy vấn dữ liệu (DML) khỏi định nghĩa cấu trúc (DDL) và quyền truy cập (DCL). Với SELECT, cần xác định bảng nguồn, điều kiện lọc, phép kết hợp, nhóm và điều kiện nhóm; các ràng buộc, view, thủ tục và trigger đảm nhiệm các vai trò khác nhau.

Phân tích: NULL biểu thị thiếu/không biết/không áp dụng, không phải số 0 hay chuỗi rỗng. So sánh với NULL bằng `=` không cho TRUE/FALSE mà cho UNKNOWN; dùng `IS NULL`/`IS NOT NULL`. Trong logic ba giá trị, AND/OR/NOT phải được tính theo cả trạng thái UNKNOWN.

Đối chiếu phương án: A. IS BLANK operator; B. IS ZERO operator; D. ISNULL operator. Các lựa chọn này không đáp ứng định nghĩa/điều kiện được hỏi; khi ôn, hãy chỉ ra chính xác điểm lệch của từng lựa chọn thay vì học thuộc mỗi chữ cái.
\end{question}

\begin{question}{322}
% [ID gốc: 5.15]
% [Nguồn: DBI202 Order 5 Q15]
What does NULL mean?
\begin{choices}
\choice The value NULL means UNKNOWN
\choice The value NULL means '' (empty string)
\choice The value NULL means ' ' (space character)
\choice The value NULL means ZERO
\end{choices}
\textbf{Đáp án:} A
\textbf{Giải thích:} Đáp án theo dữ liệu đề: A — A. The value NULL means UNKNOWN.

Nền tảng kiến thức: SQL tách truy vấn dữ liệu (DML) khỏi định nghĩa cấu trúc (DDL) và quyền truy cập (DCL). Với SELECT, cần xác định bảng nguồn, điều kiện lọc, phép kết hợp, nhóm và điều kiện nhóm; các ràng buộc, view, thủ tục và trigger đảm nhiệm các vai trò khác nhau.

Phân tích: NULL biểu thị thiếu/không biết/không áp dụng, không phải số 0 hay chuỗi rỗng. So sánh với NULL bằng `=` không cho TRUE/FALSE mà cho UNKNOWN; dùng `IS NULL`/`IS NOT NULL`. Trong logic ba giá trị, AND/OR/NOT phải được tính theo cả trạng thái UNKNOWN.

Đối chiếu phương án: B. The value NULL means '' (empty string); C. The value NULL means ' ' (space character); D. The value NULL means ZERO. Các lựa chọn này không đáp ứng định nghĩa/điều kiện được hỏi; khi ôn, hãy chỉ ra chính xác điểm lệch của từng lựa chọn thay vì học thuộc mỗi chữ cái.
\end{question}

\begin{question}{323}
% [ID gốc: 5.16]
% [Nguồn: DBI202 Order 5 Q16]
In Microsoft SQL Server, when you concatenate a NULL value with any non-null value using the + operator, the result is always NULL. This is because NULL represents missing or unknown data and any operation that involves NULL will also result in NULL.
In this case, the query SELECT 'Nancy' + NULL + 'Smith' concatenates the string 'Nancy' with a NULL value and then with the string 'Smith'. Since one of the values being concatenated is NULL, the result of the entire expression will be NULL.
QN=37 (8550) What is the value of the query:
SELECT 'Nancy' + NULL + 'Smith'
when evaluated on Microsoft SQL Server?
\begin{choices}
\choice 'Nancy Smith'
\choice 'Nancy'
\choice 'NancySmith'
\choice NULL
\end{choices}
\textbf{Đáp án:} D
\textbf{Giải thích:} Đáp án theo dữ liệu đề: D — D. NULL.

Nền tảng kiến thức: SQL tách truy vấn dữ liệu (DML) khỏi định nghĩa cấu trúc (DDL) và quyền truy cập (DCL). Với SELECT, cần xác định bảng nguồn, điều kiện lọc, phép kết hợp, nhóm và điều kiện nhóm; các ràng buộc, view, thủ tục và trigger đảm nhiệm các vai trò khác nhau.

Phân tích: NULL biểu thị thiếu/không biết/không áp dụng, không phải số 0 hay chuỗi rỗng. So sánh với NULL bằng `=` không cho TRUE/FALSE mà cho UNKNOWN; dùng `IS NULL`/`IS NOT NULL`. Trong logic ba giá trị, AND/OR/NOT phải được tính theo cả trạng thái UNKNOWN.

Đối chiếu phương án: A. 'Nancy Smith'; B. 'Nancy'; C. 'NancySmith'. Các lựa chọn này không đáp ứng định nghĩa/điều kiện được hỏi; khi ôn, hãy chỉ ra chính xác điểm lệch của từng lựa chọn thay vì học thuộc mỗi chữ cái.
\end{question}

\begin{question}{324}
% [ID gốc: 5.17]
% [Nguồn: DBI202 Order 5 Q17]
Which SQL keyword is used to sort the result-set?
\begin{choices}
\choice ORDER BY
\choice SORT BY
\choice ORDER
\choice SORT
\end{choices}
\textbf{Đáp án:} A
\textbf{Giải thích:} Đáp án theo dữ liệu đề: A — A. ORDER BY.

Nền tảng kiến thức: SQL tách truy vấn dữ liệu (DML) khỏi định nghĩa cấu trúc (DDL) và quyền truy cập (DCL). Với SELECT, cần xác định bảng nguồn, điều kiện lọc, phép kết hợp, nhóm và điều kiện nhóm; các ràng buộc, view, thủ tục và trigger đảm nhiệm các vai trò khác nhau.

Phân tích: Đọc truy vấn theo các giai đoạn logic: FROM/JOIN tạo dữ liệu nguồn, WHERE lọc từng hàng, GROUP BY tạo nhóm, HAVING lọc nhóm, SELECT xác định biểu thức đầu ra và ORDER BY sắp kết quả. Với phép join, kiểm tra điều kiện ghép và các hàng không khớp trước khi tính kết quả.

Đối chiếu phương án: B. SORT BY; C. ORDER; D. SORT. Các lựa chọn này không đáp ứng định nghĩa/điều kiện được hỏi; khi ôn, hãy chỉ ra chính xác điểm lệch của từng lựa chọn thay vì học thuộc mỗi chữ cái.
\end{question}

\begin{question}{325}
% [ID gốc: 5.18]
% [Nguồn: DBI202 Order 5 Q18]
What is the difference between the WHERE and HAVING SQL clauses?
\begin{choices}
\choice The WHERE and the HAVING clauses are identical
\choice The HAVING SQL clause condition(s) is applied to all rows in the result set before the WHERE clause is applied (if present). The WHERE clause is used only with SELECT SQL statements and specifies a search condition for an aggregate or a group
\choice The WHERE SQL clause condition(s) is applied to all rows in the result set before the HAVING clause is applied (if present). The HAVING clause is used only with SELECT SQL statements and specifies a search condition for an aggregate or a group
\choice None of the others
\end{choices}
\textbf{Đáp án:} C
\textbf{Giải thích:} Đáp án theo dữ liệu đề: C — C. The WHERE SQL clause condition(s) is applied to all rows in the result set before the HAVING clause is applied (if present). The HAVING clause is used only with SELECT SQL statements and specifies a search condition for….

Nền tảng kiến thức: SQL tách truy vấn dữ liệu (DML) khỏi định nghĩa cấu trúc (DDL) và quyền truy cập (DCL). Với SELECT, cần xác định bảng nguồn, điều kiện lọc, phép kết hợp, nhóm và điều kiện nhóm; các ràng buộc, view, thủ tục và trigger đảm nhiệm các vai trò khác nhau.

Phân tích: Phép tập hợp trên quan hệ cần so sánh các tuple theo thuộc tính tương ứng; union lấy hợp, intersection lấy phần chung và difference lấy phần có ở vế trái mà không có ở vế phải. Trước tiên kiểm tra tính tương thích schema, sau đó áp phép theo đúng hướng/toán tử.

Đối chiếu phương án: A. The WHERE and the HAVING clauses are identical; B. The HAVING SQL clause condition(s) is applied to all rows in the result set before the WHERE c…; D. None of the others. Các lựa chọn này không đáp ứng định nghĩa/điều kiện được hỏi; khi ôn, hãy chỉ ra chính xác điểm lệch của từng lựa chọn thay vì học thuộc mỗi chữ cái.
\end{question}

\begin{question}{326}
% [ID gốc: 5.20]
% [Nguồn: DBI202 Order 5 Q20]
"ON DELETE CASCADE" is a referential action for a foreign key that allows you to automatically delete data from child tables when you delete data from the parent table. It goes on the foreign key of the child table so that when you delete a parent row, all children rows are delete. This ensures referential integrity by automatically removing any dependent rows in the child relation when a row in the parent relation is deleted.
QN=44 (8585) Suppose R and S are 2 relations. R is the parent of S.
And the relationship between R and S is set to
"ON DELETE CASCADE".
This means that:
\begin{choices}
\choice We can delete a row from R if that row has children in S (and in this case, the database server will raise up an error)
\choice We can delete a row from R although that row has children in S (and in this case, all the children will be deleted too)
\end{choices}
\textbf{Đáp án:} B
\textbf{Giải thích:} Đáp án theo dữ liệu đề: B — B. We can delete a row from R although that row has children in S (and in this case, all the children will be deleted too).

Nền tảng kiến thức: SQL tách truy vấn dữ liệu (DML) khỏi định nghĩa cấu trúc (DDL) và quyền truy cập (DCL). Với SELECT, cần xác định bảng nguồn, điều kiện lọc, phép kết hợp, nhóm và điều kiện nhóm; các ràng buộc, view, thủ tục và trigger đảm nhiệm các vai trò khác nhau.

Phân tích: Khóa ngoại buộc giá trị tham chiếu phải khớp một khóa được tham chiếu ở quan hệ đích (hoặc NULL nếu ràng buộc cho phép). Đây là toàn vẹn tham chiếu; nó khác tính duy nhất của primary/candidate key.

Đối chiếu phương án: A. We can delete a row from R if that row has children in S (and in this case, the database serve…. Các lựa chọn này không đáp ứng định nghĩa/điều kiện được hỏi; khi ôn, hãy chỉ ra chính xác điểm lệch của từng lựa chọn thay vì học thuộc mỗi chữ cái.
\end{question}

\begin{question}{327}
% [ID gốc: 5.21]
% [Nguồn: DBI202 Order 5 Q21]
\% for multi-character
_ for single-character
QN=47 (8578) What one is a wildcard used for pattern matching?
\begin{choices}
\choice "%" (for multi-character substitution)
\choice "?" (for single character substitution)
\choice "?" (for multi-character substitution)
\end{choices}
\textbf{Đáp án:} A
\textbf{Giải thích:} Đáp án theo dữ liệu đề: A — A. "%" (for multi-character substitution).

Nền tảng kiến thức: SQL tách truy vấn dữ liệu (DML) khỏi định nghĩa cấu trúc (DDL) và quyền truy cập (DCL). Với SELECT, cần xác định bảng nguồn, điều kiện lọc, phép kết hợp, nhóm và điều kiện nhóm; các ràng buộc, view, thủ tục và trigger đảm nhiệm các vai trò khác nhau.

Phân tích: Trong LIKE của SQL Server, `%` khớp 0 hoặc nhiều ký tự và `_` khớp đúng một ký tự. ESCAPE cho phép coi ký tự wildcard như ký tự thường bằng cách đặt trước chúng một ký tự thoát; nó không thay thế mẫu bằng chuỗi bất kỳ.

Đối chiếu phương án: B. "?" (for single character substitution); C. "?" (for multi-character substitution). Các lựa chọn này không đáp ứng định nghĩa/điều kiện được hỏi; khi ôn, hãy chỉ ra chính xác điểm lệch của từng lựa chọn thay vì học thuộc mỗi chữ cái.
\end{question}

\begin{question}{328}
% [ID gốc: 5.22]
% [Nguồn: DBI202 Order 5 Q22]
To create a DEFAULT constraint on the "City" column of the table PERSON which is already created, use the following SQL:
\begin{choices}
\choice ALTER TABLE Person
ALTER COLUMN City SET DEFAULT 'SANDNES'
\choice ALTER TABLE Person
EDIT COLUMN City SET DEFAULT 'SANDNES'
\choice ALTER TABLE Person
UPDATE COLUMN City SET DEFAULT 'SANDNES'
\choice ALTER TABLE Person
MODIFY COLUMN City SET DEFAULT 'SANDNES'
\end{choices}
\textbf{Đáp án:} A
\textbf{Giải thích:} Đáp án theo dữ liệu đề: A — A. ALTER TABLE Person ALTER COLUMN City SET DEFAULT 'SANDNES'.

Nền tảng kiến thức: SQL tách truy vấn dữ liệu (DML) khỏi định nghĩa cấu trúc (DDL) và quyền truy cập (DCL). Với SELECT, cần xác định bảng nguồn, điều kiện lọc, phép kết hợp, nhóm và điều kiện nhóm; các ràng buộc, view, thủ tục và trigger đảm nhiệm các vai trò khác nhau.

Phân tích: ALTER TABLE thay đổi định nghĩa bảng đã tồn tại, như thêm cột hoặc sửa cấu hình cột/ràng buộc theo cú pháp của DBMS. Nó khác UPDATE, vốn sửa các giá trị trong những hàng hiện hữu; cú pháp chính xác có thể khác giữa SQL Server và chuẩn SQL.

Đối chiếu phương án: B. ALTER TABLE Person EDIT COLUMN City SET DEFAULT 'SANDNES'; C. ALTER TABLE Person UPDATE COLUMN City SET DEFAULT 'SANDNES'; D. ALTER TABLE Person MODIFY COLUMN City SET DEFAULT 'SANDNES'. Các lựa chọn này không đáp ứng định nghĩa/điều kiện được hỏi; khi ôn, hãy chỉ ra chính xác điểm lệch của từng lựa chọn thay vì học thuộc mỗi chữ cái.
\end{question}

\begin{question}{329}
% [ID gốc: 5.23]
% [Nguồn: DBI202 Order 5 Q23]
Cho quan hệ \textbf{R(A, B, C)} với dữ liệu như sau:

$$\begin{array}{|c|c|c|}
\hline
\textbf{A} & \textbf{B} & \textbf{C} \\
\hline
1 & 2 & 3 \\
\hline
4 & 5 & 6 \\
\hline
7 & 8 & 9 \\
\hline
\end{array}$$

Which is the result of the following query:
\begin{codebox}
SELECT SUM(B)
FROM R
WHERE C > 10;
\end{codebox}
\begin{choices}
\choice 0
\choice NULL
\choice 3
\choice 10
\end{choices}
\textbf{Đáp án:} B
\textbf{Giải thích:} Điều kiện WHERE C > 10 không có hàng nào thỏa mãn (các giá trị của C là 3, 6, 9 đều <= 10). Trong chuẩn SQL, khi hàm tổng hợp SUM() tính trên một tập rỗng (empty set), kết quả trả về là NULL.
\end{question}

\begin{question}{330}
% [ID gốc: 5.24]
% [Nguồn: DBI202 Order 5 Q24]
The table Arc(x,y) currently has the following tuples (note there are duplicates):
(See picture)
Compute the result of the query:
SELECT x, COUNT(y)
FROM Arc
GROUP BY x;
Which of the following tuples is in the result?
[file:8522.jpg]
\begin{choices}
\choice (4,4)
\choice (4,5)
\choice (4,6)
\choice (4,7)
\choice (4,8)
\end{choices}
\textbf{Đáp án:} A
\textbf{Giải thích:} Đáp án theo dữ liệu đề: A — A. (4,4).

Nền tảng kiến thức: SQL tách truy vấn dữ liệu (DML) khỏi định nghĩa cấu trúc (DDL) và quyền truy cập (DCL). Với SELECT, cần xác định bảng nguồn, điều kiện lọc, phép kết hợp, nhóm và điều kiện nhóm; các ràng buộc, view, thủ tục và trigger đảm nhiệm các vai trò khác nhau.

Phân tích: Các hàng dữ liệu của một relation được gọi là tuple; header chứa tên thuộc tính không phải tuple dữ liệu. Trong bảng, tuple tương ứng một bản ghi, còn attribute tương ứng cột.

Đối chiếu phương án: B. (4,5); C. (4,6); D. (4,7); E. (4,8). Các lựa chọn này không đáp ứng định nghĩa/điều kiện được hỏi; khi ôn, hãy chỉ ra chính xác điểm lệch của từng lựa chọn thay vì học thuộc mỗi chữ cái.
\end{question}

\begin{question}{331}
% [ID gốc: 5.26]
% [Nguồn: DBI202 Order 5 Q26]
Given the relation Employee(SSN, FNAME, LNAME, SALARY, DepartmentNo). Select the right query below to count the number of employees in each department.
\begin{choices}
\choice SELECT COUNT(*) FROM Employee
\choice SELECT DepartmentNo, COUNT(*) FROM Employee GROUP BY DepartmentNo
\choice SELECT DepartmentNo, COUNT(*) FROM Employee
\choice None of the others
\end{choices}
\textbf{Đáp án:} B
\textbf{Giải thích:} Đáp án theo dữ liệu đề: B — B. SELECT DepartmentNo, COUNT(*) FROM Employee GROUP BY DepartmentNo.

Nền tảng kiến thức: SQL tách truy vấn dữ liệu (DML) khỏi định nghĩa cấu trúc (DDL) và quyền truy cập (DCL). Với SELECT, cần xác định bảng nguồn, điều kiện lọc, phép kết hợp, nhóm và điều kiện nhóm; các ràng buộc, view, thủ tục và trigger đảm nhiệm các vai trò khác nhau.

Phân tích: Đọc truy vấn theo các giai đoạn logic: FROM/JOIN tạo dữ liệu nguồn, WHERE lọc từng hàng, GROUP BY tạo nhóm, HAVING lọc nhóm, SELECT xác định biểu thức đầu ra và ORDER BY sắp kết quả. Với phép join, kiểm tra điều kiện ghép và các hàng không khớp trước khi tính kết quả.

Đối chiếu phương án: A. SELECT COUNT(*) FROM Employee; C. SELECT DepartmentNo, COUNT(*) FROM Employee; D. None of the others. Các lựa chọn này không đáp ứng định nghĩa/điều kiện được hỏi; khi ôn, hãy chỉ ra chính xác điểm lệch của từng lựa chọn thay vì học thuộc mỗi chữ cái.
\end{question}

\begin{question}{332}
% [ID gốc: 5.28]
% [Nguồn: DBI202 Order 5 Q28]
Cho hai quan hệ \textbf{R(A, B)} và \textbf{S(C, D)}:

$$\textbf{R} \qquad\qquad \textbf{S}$$
$$\begin{array}{|c|c|}
\hline
\textbf{A} & \textbf{B} \\
\hline
1 & 4 \\
\hline
2 & 5 \\
\hline
8 & 6 \\
\hline
4 & 7 \\
\hline
\end{array}
\qquad\qquad
\begin{array}{|c|c|}
\hline
\textbf{C} & \textbf{D} \\
\hline
0 & 1 \\
\hline
1 & 0 \\
\hline
2 & 1 \\
\hline
5 & 8 \\
\hline
\end{array}$$

Compute the result of the query:
\begin{codebox}
SELECT A, B, C, D
FROM R RIGHT OUTER JOIN S
ON R.A = S.C;
\end{codebox}
Identify, in the list below, 2 rows that appear in the result:
\begin{choices}
\choice (null, null, 0, 1)
\choice (null, null, 5, 8)
\choice (8, 6, null, null)
\choice (4, 7, null, null)
\end{choices}
\textbf{Đáp án:} AB
\textbf{Giải thích:} Đáp án theo dữ liệu đề: A, B — A. (null, null, 0, 1); B. (null, null, 5, 8).

Nền tảng kiến thức: SQL tách truy vấn dữ liệu (DML) khỏi định nghĩa cấu trúc (DDL) và quyền truy cập (DCL). Với SELECT, cần xác định bảng nguồn, điều kiện lọc, phép kết hợp, nhóm và điều kiện nhóm; các ràng buộc, view, thủ tục và trigger đảm nhiệm các vai trò khác nhau.

Phân tích: Join ghép các bộ thỏa điều kiện. Inner join loại bộ không ghép được; LEFT/RIGHT OUTER JOIN giữ mọi bộ từ phía được chỉ định và điền NULL phía thiếu. Natural join tự ghép các thuộc tính trùng tên và chỉ xuất một cột cho mỗi thuộc tính ghép.

Đối chiếu phương án: A. (null, null, 0, 1); B. (null, null, 5, 8). Các lựa chọn này không đáp ứng định nghĩa/điều kiện được hỏi; khi ôn, hãy chỉ ra chính xác điểm lệch của từng lựa chọn thay vì học thuộc mỗi chữ cái.
\end{question}

\begin{question}{333}
% [ID gốc: 5.30]
% [Nguồn: DBI202 Order 5 Q30]
(TRUE OR NULL) return:
\begin{choices}
\choice TRUE
\choice FALSE
\choice NULL
\choice None of the others
\end{choices}
\textbf{Đáp án:} A
\textbf{Giải thích:} Đáp án theo dữ liệu đề: A — A. TRUE.

Nền tảng kiến thức: SQL tách truy vấn dữ liệu (DML) khỏi định nghĩa cấu trúc (DDL) và quyền truy cập (DCL). Với SELECT, cần xác định bảng nguồn, điều kiện lọc, phép kết hợp, nhóm và điều kiện nhóm; các ràng buộc, view, thủ tục và trigger đảm nhiệm các vai trò khác nhau.

Phân tích: NULL biểu thị thiếu/không biết/không áp dụng, không phải số 0 hay chuỗi rỗng. So sánh với NULL bằng `=` không cho TRUE/FALSE mà cho UNKNOWN; dùng `IS NULL`/`IS NOT NULL`. Trong logic ba giá trị, AND/OR/NOT phải được tính theo cả trạng thái UNKNOWN.

Đối chiếu phương án: B. FALSE; C. NULL; D. None of the others. Các lựa chọn này không đáp ứng định nghĩa/điều kiện được hỏi; khi ôn, hãy chỉ ra chính xác điểm lệch của từng lựa chọn thay vì học thuộc mỗi chữ cái.
\end{question}

\begin{question}{334}
% [ID gốc: 5.31]
% [Nguồn: DBI202 Order 5 Q31]
Choose one correct statement:
\begin{choices}
\choice Two null values are equal
\choice Comparisons between two null values, or between a NULL and any other value, return UNKNOWN
\choice Comparisons between two null values, or between a NULL and any other value, return FALSE
\end{choices}
\textbf{Đáp án:} B
\textbf{Giải thích:} Đáp án theo dữ liệu đề: B — B. Comparisons between two null values, or between a NULL and any other value, return UNKNOWN.

Nền tảng kiến thức: SQL tách truy vấn dữ liệu (DML) khỏi định nghĩa cấu trúc (DDL) và quyền truy cập (DCL). Với SELECT, cần xác định bảng nguồn, điều kiện lọc, phép kết hợp, nhóm và điều kiện nhóm; các ràng buộc, view, thủ tục và trigger đảm nhiệm các vai trò khác nhau.

Phân tích: NULL biểu thị thiếu/không biết/không áp dụng, không phải số 0 hay chuỗi rỗng. So sánh với NULL bằng `=` không cho TRUE/FALSE mà cho UNKNOWN; dùng `IS NULL`/`IS NOT NULL`. Trong logic ba giá trị, AND/OR/NOT phải được tính theo cả trạng thái UNKNOWN.

Đối chiếu phương án: A. Two null values are equal; C. Comparisons between two null values, or between a NULL and any other value, return FALSE. Các lựa chọn này không đáp ứng định nghĩa/điều kiện được hỏi; khi ôn, hãy chỉ ra chính xác điểm lệch của từng lựa chọn thay vì học thuộc mỗi chữ cái.
\end{question}

\begin{question}{335}
% [ID gốc: 5.32]
% [Nguồn: DBI202 Order 5 Q32]
Muốn UNKNOWN thì buộc một cái phải NULL, chọn A luôn
QN=43 (8554) For what values of x, y, and z, including NULL, does the Boolean expression
x = 15 OR z = 35)
have the truth value UNKNOWN? Identify one of those values from the list below.
\begin{choices}
\choice x = NULL, y = 0, z = 4.
\choice x = -2, y = -2, z = -2.
\choice x = 3, y = 0, z = 7.
\choice x = 2, y = -1, z = 6.
\end{choices}
\textbf{Đáp án:} A
\textbf{Giải thích:} Đáp án theo dữ liệu đề: A — A. x = NULL, y = 0, z = 4..

Nền tảng kiến thức: SQL tách truy vấn dữ liệu (DML) khỏi định nghĩa cấu trúc (DDL) và quyền truy cập (DCL). Với SELECT, cần xác định bảng nguồn, điều kiện lọc, phép kết hợp, nhóm và điều kiện nhóm; các ràng buộc, view, thủ tục và trigger đảm nhiệm các vai trò khác nhau.

Phân tích: NULL biểu thị thiếu/không biết/không áp dụng, không phải số 0 hay chuỗi rỗng. So sánh với NULL bằng `=` không cho TRUE/FALSE mà cho UNKNOWN; dùng `IS NULL`/`IS NOT NULL`. Trong logic ba giá trị, AND/OR/NOT phải được tính theo cả trạng thái UNKNOWN.

Đối chiếu phương án: B. x = -2, y = -2, z = -2.; C. x = 3, y = 0, z = 7.; D. x = 2, y = -1, z = 6.. Các lựa chọn này không đáp ứng định nghĩa/điều kiện được hỏi; khi ôn, hãy chỉ ra chính xác điểm lệch của từng lựa chọn thay vì học thuộc mỗi chữ cái.
\end{question}

\begin{question}{336}
% [ID gốc: 5.34]
% [Nguồn: DBI202 Order 5 Q34]
Foreign key constraints are created by using "_____" keyword to refer to the primary key of another table
\begin{choices}
\choice REFERENCES
\choice POINT TO
\choice REFER
\choice None of the others
\end{choices}
\textbf{Đáp án:} A
\textbf{Giải thích:} Đáp án theo dữ liệu đề: A — A. REFERENCES.

Nền tảng kiến thức: SQL tách truy vấn dữ liệu (DML) khỏi định nghĩa cấu trúc (DDL) và quyền truy cập (DCL). Với SELECT, cần xác định bảng nguồn, điều kiện lọc, phép kết hợp, nhóm và điều kiện nhóm; các ràng buộc, view, thủ tục và trigger đảm nhiệm các vai trò khác nhau.

Phân tích: Khóa ngoại buộc giá trị tham chiếu phải khớp một khóa được tham chiếu ở quan hệ đích (hoặc NULL nếu ràng buộc cho phép). Đây là toàn vẹn tham chiếu; nó khác tính duy nhất của primary/candidate key.

Đối chiếu phương án: B. POINT TO; C. REFER; D. None of the others. Các lựa chọn này không đáp ứng định nghĩa/điều kiện được hỏi; khi ôn, hãy chỉ ra chính xác điểm lệch của từng lựa chọn thay vì học thuộc mỗi chữ cái.
\end{question}

\begin{question}{337}
% [ID gốc: 5.38]
% [Nguồn: DBI202 Order 5 Q38]
Pay attention into the following query:
SELECT * FROM R
WHERE A LIKE '\%a%';
So, in the above case, the wildcard \% represents what?
\begin{choices}
\choice \% (percent sign) represents zero, one, or more character
\choice \% (percent sign) represents exactly 1 character
\end{choices}
\textbf{Đáp án:} A
\textbf{Giải thích:} Đáp án theo dữ liệu đề: A — A. \% (percent sign) represents zero, one, or more character.

Nền tảng kiến thức: SQL tách truy vấn dữ liệu (DML) khỏi định nghĩa cấu trúc (DDL) và quyền truy cập (DCL). Với SELECT, cần xác định bảng nguồn, điều kiện lọc, phép kết hợp, nhóm và điều kiện nhóm; các ràng buộc, view, thủ tục và trigger đảm nhiệm các vai trò khác nhau.

Phân tích: Trong LIKE của SQL Server, `%` khớp 0 hoặc nhiều ký tự và `_` khớp đúng một ký tự. ESCAPE cho phép coi ký tự wildcard như ký tự thường bằng cách đặt trước chúng một ký tự thoát; nó không thay thế mẫu bằng chuỗi bất kỳ.

Đối chiếu phương án: B. \% (percent sign) represents exactly 1 character. Các lựa chọn này không đáp ứng định nghĩa/điều kiện được hỏi; khi ôn, hãy chỉ ra chính xác điểm lệch của từng lựa chọn thay vì học thuộc mỗi chữ cái.
\end{question}

\begin{question}{338}
% [ID gốc: 5.39]
% [Nguồn: DBI202 Order 5 Q39]
With SQL, how do you select all the records from a table named "Persons" where the value of the column "FirstName" starts with an "a"?
\begin{choices}
\choice SELECT * FROM Persons WHERE FirstName='a'
\choice SELECT * FROM Persons WHERE FirstName LIKE '\%a'
\choice SELECT * FROM Persons WHERE FirstName='\%a%'
\choice SELECT * FROM Persons WHERE FirstName LIKE 'a%'
\end{choices}
\textbf{Đáp án:} D
\textbf{Giải thích:} Đáp án theo dữ liệu đề: D — D. SELECT * FROM Persons WHERE FirstName LIKE 'a%'.

Nền tảng kiến thức: SQL tách truy vấn dữ liệu (DML) khỏi định nghĩa cấu trúc (DDL) và quyền truy cập (DCL). Với SELECT, cần xác định bảng nguồn, điều kiện lọc, phép kết hợp, nhóm và điều kiện nhóm; các ràng buộc, view, thủ tục và trigger đảm nhiệm các vai trò khác nhau.

Phân tích: Đọc truy vấn theo các giai đoạn logic: FROM/JOIN tạo dữ liệu nguồn, WHERE lọc từng hàng, GROUP BY tạo nhóm, HAVING lọc nhóm, SELECT xác định biểu thức đầu ra và ORDER BY sắp kết quả. Với phép join, kiểm tra điều kiện ghép và các hàng không khớp trước khi tính kết quả.

Đối chiếu phương án: A. SELECT * FROM Persons WHERE FirstName='a'; B. SELECT * FROM Persons WHERE FirstName LIKE '\%a'; C. SELECT * FROM Persons WHERE FirstName='\%a%'. Các lựa chọn này không đáp ứng định nghĩa/điều kiện được hỏi; khi ôn, hãy chỉ ra chính xác điểm lệch của từng lựa chọn thay vì học thuộc mỗi chữ cái.
\end{question}

\begin{question}{339}
% [ID gốc: 5.40]
% [Nguồn: DBI202 Order 5 Q40]
What does the keyword ESCAPE mean?
\begin{choices}
\choice The keyword ESCAPE is used to match any string that contains the characters "%" or "_"
\choice There is no keyword named ESCAPE
\choice The keyword ESCAPE is used to match any string that contains the characters "?" or "*"
\choice The keyword ESCAPE is used to match any string that contains the characters "@" or "$"
\end{choices}
\textbf{Đáp án:} A
\textbf{Giải thích:} Đáp án theo dữ liệu đề: A — A. The keyword ESCAPE is used to match any string that contains the characters "%" or "_".

Nền tảng kiến thức: SQL tách truy vấn dữ liệu (DML) khỏi định nghĩa cấu trúc (DDL) và quyền truy cập (DCL). Với SELECT, cần xác định bảng nguồn, điều kiện lọc, phép kết hợp, nhóm và điều kiện nhóm; các ràng buộc, view, thủ tục và trigger đảm nhiệm các vai trò khác nhau.

Phân tích: Trong LIKE của SQL Server, `%` khớp 0 hoặc nhiều ký tự và `_` khớp đúng một ký tự. ESCAPE cho phép coi ký tự wildcard như ký tự thường bằng cách đặt trước chúng một ký tự thoát; nó không thay thế mẫu bằng chuỗi bất kỳ.

Đối chiếu phương án: B. There is no keyword named ESCAPE; C. The keyword ESCAPE is used to match any string that contains the characters "?" or "*"; D. The keyword ESCAPE is used to match any string that contains the characters "@" or "$". Các lựa chọn này không đáp ứng định nghĩa/điều kiện được hỏi; khi ôn, hãy chỉ ra chính xác điểm lệch của từng lựa chọn thay vì học thuộc mỗi chữ cái.
\end{question}

\begin{question}{340}
% [ID gốc: 5.41]
% [Nguồn: DBI202 Order 5 Q41]
The SQL BETWEEN operator:
\begin{choices}
\choice Specifies a range to test.
\choice Specifies which tables we are selecting from.
\choice Specifies that a column is a primary key.
\choice None of the others
\end{choices}
\textbf{Đáp án:} A
\textbf{Giải thích:} Đáp án theo dữ liệu đề: A — A. Specifies a range to test..

Nền tảng kiến thức: SQL tách truy vấn dữ liệu (DML) khỏi định nghĩa cấu trúc (DDL) và quyền truy cập (DCL). Với SELECT, cần xác định bảng nguồn, điều kiện lọc, phép kết hợp, nhóm và điều kiện nhóm; các ràng buộc, view, thủ tục và trigger đảm nhiệm các vai trò khác nhau.

Phân tích: Với lựa chọn tổng hợp, không chọn theo cảm tính: kiểm tra từng phát biểu độc lập. “All of the others/above” chỉ đúng khi mọi phát biểu còn lại đúng; “None” chỉ đúng khi không phát biểu nào trong nhóm thỏa yêu cầu. Đáp án lưu cần được đọc cùng các mệnh đề thành phần.

Đối chiếu phương án: B. Specifies which tables we are selecting from.; C. Specifies that a column is a primary key.; D. None of the others. Các lựa chọn này không đáp ứng định nghĩa/điều kiện được hỏi; khi ôn, hãy chỉ ra chính xác điểm lệch của từng lựa chọn thay vì học thuộc mỗi chữ cái.
\end{question}

\begin{question}{341}
% [ID gốc: 5.46]
% [Nguồn: DBI202 Order 5 Q46]
What is the difference between Where and Having Clause?
\begin{choices}
\choice WHERE is for Rows and HAVING is for Groups
\choice WHERE is for Groups and HAVING is for Rows
\choice WHERE and HAVING are the same
\choice We can use aggregation functions in WHERE clause
\end{choices}
\textbf{Đáp án:} A
\textbf{Giải thích:} WHERE filters individual rows before grouping; HAVING filters groups after aggregation.
\end{question}

\begin{question}{342}
% [ID gốc: 5.47]
% [Nguồn: DBI202 Order 5 Q47]
What does the following SQL statement do:
SELECT Customer, COUNT(Order)
FROM Sales
GROUP BY Customer
HAVING COUNT(Order) > 5
\begin{choices}
\choice Selects the total number of orders from the Sales table, if this number is greater than 5
\choice Selects all Customers from the Sales table
\choice Selects all customers from table Sales that have made more than 5 orders.
\choice None of the other three
\end{choices}
\textbf{Đáp án:} C
\textbf{Giải thích:} HAVING COUNT(Order) > 5 filters customers who placed strictly more than 5 orders.
\end{question}

\begin{question}{343}
% [ID gốc: 5.48]
% [Nguồn: DBI202 Order 5 Q48]
The relation R(A,B) may have duplicate tuples. Choose the query that create non-duplicated results, regardless of what tuples R contains?
\begin{choices}
\choice SELECT A, B FROM R
\choice SELECT A, B FROM R GROUP BY A, B
\choice SELECT A FROM R WHERE A NOT IN (SELECT B FROM R)
\choice SELECT TOP 3 A, B FROM R
\end{choices}
\textbf{Đáp án:} B
\textbf{Giải thích:} GROUP BY A, B collapses all identical rows into distinct tuples.
\end{question}

\begin{question}{344}
% [ID gốc: 5.49]
% [Nguồn: DBI202 Order 5 Q49]
The table R(X,Y) currently has the following tuples (note there are duplicates).The relation S(A,B,C) has the following tuples:
(See picture)
Which is the result when we execute the following query:
SELECT * FROM R
UNION
SELECT * FROM S
[file:8524.jpg]
\begin{choices}
\choice There is an error a
\choice The output has 15 rows
\choice The output has 16 rows
\choice The output has 17 rows
\end{choices}
\textbf{Đáp án:} A
\textbf{Giải thích:} Đáp án theo dữ liệu đề: A — A. There is an error a.

Nền tảng kiến thức: SQL tách truy vấn dữ liệu (DML) khỏi định nghĩa cấu trúc (DDL) và quyền truy cập (DCL). Với SELECT, cần xác định bảng nguồn, điều kiện lọc, phép kết hợp, nhóm và điều kiện nhóm; các ràng buộc, view, thủ tục và trigger đảm nhiệm các vai trò khác nhau.

Phân tích: Các hàng dữ liệu của một relation được gọi là tuple; header chứa tên thuộc tính không phải tuple dữ liệu. Trong bảng, tuple tương ứng một bản ghi, còn attribute tương ứng cột.

Đối chiếu phương án: B. The output has 15 rows; C. The output has 16 rows; D. The output has 17 rows. Các lựa chọn này không đáp ứng định nghĩa/điều kiện được hỏi; khi ôn, hãy chỉ ra chính xác điểm lệch của từng lựa chọn thay vì học thuộc mỗi chữ cái.
\end{question}

\begin{question}{345}
% [ID gốc: 5.50]
% [Nguồn: DBI202 Order 5 Q50]
What is a "join"?
\begin{choices}
\choice 'join' used to connect two or more tables logically with or without common field(s)
\choice 'join' used to connect two or more tables logically with common field(s)
\choice 'join' used to connect two or more tables logically without common field(s)
\choice 'join' used to connect two or more tables logically with aliases
\end{choices}
\textbf{Đáp án:} A
\textbf{Giải thích:} Đáp án theo dữ liệu đề: A — A. 'join' used to connect two or more tables logically with or without common field(s).

Nền tảng kiến thức: SQL tách truy vấn dữ liệu (DML) khỏi định nghĩa cấu trúc (DDL) và quyền truy cập (DCL). Với SELECT, cần xác định bảng nguồn, điều kiện lọc, phép kết hợp, nhóm và điều kiện nhóm; các ràng buộc, view, thủ tục và trigger đảm nhiệm các vai trò khác nhau.

Phân tích: Join ghép các bộ thỏa điều kiện. Inner join loại bộ không ghép được; LEFT/RIGHT OUTER JOIN giữ mọi bộ từ phía được chỉ định và điền NULL phía thiếu. Natural join tự ghép các thuộc tính trùng tên và chỉ xuất một cột cho mỗi thuộc tính ghép.

Đối chiếu phương án: B. 'join' used to connect two or more tables logically with common field(s); C. 'join' used to connect two or more tables logically without common field(s); D. 'join' used to connect two or more tables logically with aliases. Các lựa chọn này không đáp ứng định nghĩa/điều kiện được hỏi; khi ôn, hãy chỉ ra chính xác điểm lệch của từng lựa chọn thay vì học thuộc mỗi chữ cái.
\end{question}

\begin{question}{346}
% [ID gốc: 5.51]
% [Nguồn: DBI202 Order 5 Q51]
When join R and S, we want to display all the records in R we must use:
\begin{choices}
\choice LEFT OUTER JOIN a
\choice RIGHT OUTER JOIN
\choice INNER JOIN
\choice None of the others
\end{choices}
\textbf{Đáp án:} A
\textbf{Giải thích:} Đáp án theo dữ liệu đề: A — A. LEFT OUTER JOIN a.

Nền tảng kiến thức: SQL tách truy vấn dữ liệu (DML) khỏi định nghĩa cấu trúc (DDL) và quyền truy cập (DCL). Với SELECT, cần xác định bảng nguồn, điều kiện lọc, phép kết hợp, nhóm và điều kiện nhóm; các ràng buộc, view, thủ tục và trigger đảm nhiệm các vai trò khác nhau.

Phân tích: Join ghép các bộ thỏa điều kiện. Inner join loại bộ không ghép được; LEFT/RIGHT OUTER JOIN giữ mọi bộ từ phía được chỉ định và điền NULL phía thiếu. Natural join tự ghép các thuộc tính trùng tên và chỉ xuất một cột cho mỗi thuộc tính ghép.

Đối chiếu phương án: B. RIGHT OUTER JOIN; C. INNER JOIN; D. None of the others. Các lựa chọn này không đáp ứng định nghĩa/điều kiện được hỏi; khi ôn, hãy chỉ ra chính xác điểm lệch của từng lựa chọn thay vì học thuộc mỗi chữ cái.
\end{question}

\begin{question}{347}
% [ID gốc: 5.52]
% [Nguồn: DBI202 Order 5 Q52]
Relation R(x,y) currently consists of only one tuple (NULL, NULL). Which of the following queries will produce a nonempty output? That is, at least one tuple will be produced, although the tuple(s) may have NULL's.
\begin{choices}
\choice SELECT * FROM R WHERE x IS NULL
\choice SELECT * FROM R WHERE x = NULL
\choice SELECT * FROM R WHERE x = y
\choice SELECT * FROM R WHERE x <> 10
\end{choices}
\textbf{Đáp án:} A
\textbf{Giải thích:} Đáp án theo dữ liệu đề: A — A. SELECT * FROM R WHERE x IS NULL.

Nền tảng kiến thức: SQL tách truy vấn dữ liệu (DML) khỏi định nghĩa cấu trúc (DDL) và quyền truy cập (DCL). Với SELECT, cần xác định bảng nguồn, điều kiện lọc, phép kết hợp, nhóm và điều kiện nhóm; các ràng buộc, view, thủ tục và trigger đảm nhiệm các vai trò khác nhau.

Phân tích: NULL biểu thị thiếu/không biết/không áp dụng, không phải số 0 hay chuỗi rỗng. So sánh với NULL bằng `=` không cho TRUE/FALSE mà cho UNKNOWN; dùng `IS NULL`/`IS NOT NULL`. Trong logic ba giá trị, AND/OR/NOT phải được tính theo cả trạng thái UNKNOWN.

Đối chiếu phương án: B. SELECT * FROM R WHERE x = NULL; C. SELECT * FROM R WHERE x = y; D. SELECT * FROM R WHERE x <> 10. Các lựa chọn này không đáp ứng định nghĩa/điều kiện được hỏi; khi ôn, hãy chỉ ra chính xác điểm lệch của từng lựa chọn thay vì học thuộc mỗi chữ cái.
\end{question}

\begin{question}{348}
% [ID gốc: 5.53]
% [Nguồn: DBI202 Order 5 Q53]
(TRUE AND NULL) returns: ???
\begin{choices}
\choice NULL
\choice TRUE
\choice FALSE
\choice None of the others
\end{choices}
\textbf{Đáp án:} A
\textbf{Giải thích:} Đáp án theo dữ liệu đề: A — A. NULL.

Nền tảng kiến thức: SQL tách truy vấn dữ liệu (DML) khỏi định nghĩa cấu trúc (DDL) và quyền truy cập (DCL). Với SELECT, cần xác định bảng nguồn, điều kiện lọc, phép kết hợp, nhóm và điều kiện nhóm; các ràng buộc, view, thủ tục và trigger đảm nhiệm các vai trò khác nhau.

Phân tích: NULL biểu thị thiếu/không biết/không áp dụng, không phải số 0 hay chuỗi rỗng. So sánh với NULL bằng `=` không cho TRUE/FALSE mà cho UNKNOWN; dùng `IS NULL`/`IS NOT NULL`. Trong logic ba giá trị, AND/OR/NOT phải được tính theo cả trạng thái UNKNOWN.

Đối chiếu phương án: B. TRUE; C. FALSE; D. None of the others. Các lựa chọn này không đáp ứng định nghĩa/điều kiện được hỏi; khi ôn, hãy chỉ ra chính xác điểm lệch của từng lựa chọn thay vì học thuộc mỗi chữ cái.
\end{question}

\begin{question}{349}
% [ID gốc: 5.54]
% [Nguồn: DBI202 Order 5 Q54]
In order to modify or delete an existing constraint, it is necessary that:
\begin{choices}
\choice the constraint has a name
\choice the constraint has an ID
\choice the constraint has some attributes
\end{choices}
\textbf{Đáp án:} A
\textbf{Giải thích:} Đáp án theo dữ liệu đề: A — A. the constraint has a name.

Nền tảng kiến thức: SQL tách truy vấn dữ liệu (DML) khỏi định nghĩa cấu trúc (DDL) và quyền truy cập (DCL). Với SELECT, cần xác định bảng nguồn, điều kiện lọc, phép kết hợp, nhóm và điều kiện nhóm; các ràng buộc, view, thủ tục và trigger đảm nhiệm các vai trò khác nhau.

Phân tích: Muốn sửa/xóa một constraint cụ thể, cần định danh constraint đó để phát lệnh ALTER TABLE DROP CONSTRAINT hoặc thao tác tương ứng. Đặt tên tường minh khi tạo giúp quản trị ổn định hơn so với tên tự sinh phụ thuộc môi trường.

Đối chiếu phương án: B. the constraint has an ID; C. the constraint has some attributes. Các lựa chọn này không đáp ứng định nghĩa/điều kiện được hỏi; khi ôn, hãy chỉ ra chính xác điểm lệch của từng lựa chọn thay vì học thuộc mỗi chữ cái.
\end{question}

\begin{question}{350}
% [ID gốc: 5.55]
% [Nguồn: DBI202 Order 5 Q55]
The IN SQL keyword
\begin{choices}
\choice Is used with the DISTINCT SQL keyword only
\choice Is used with the INSERT SQL keyword only
\choice Determines if a value
any of the values in a list or a sub-query
\choice Defines the tables we are selecting or deleting data from
\end{choices}
\textbf{Đáp án:} C
\textbf{Giải thích:} Đáp án theo dữ liệu đề: C — C. Determines if a value any of the values in a list or a sub-query.

Nền tảng kiến thức: SQL tách truy vấn dữ liệu (DML) khỏi định nghĩa cấu trúc (DDL) và quyền truy cập (DCL). Với SELECT, cần xác định bảng nguồn, điều kiện lọc, phép kết hợp, nhóm và điều kiện nhóm; các ràng buộc, view, thủ tục và trigger đảm nhiệm các vai trò khác nhau.

Phân tích: Trong T-SQL, biến cục bộ do người dùng khai báo bắt đầu bằng một dấu `@` (ví dụ `@count`); tên bắt đầu `@@` được dùng cho các hàm/biến hệ thống. Đừng nhầm tiền tố biến cục bộ với các biến hệ thống như `@@ROWCOUNT`.

Đối chiếu phương án: A. Is used with the DISTINCT SQL keyword only; B. Is used with the INSERT SQL keyword only; D. Defines the tables we are selecting or deleting data from. Các lựa chọn này không đáp ứng định nghĩa/điều kiện được hỏi; khi ôn, hãy chỉ ra chính xác điểm lệch của từng lựa chọn thay vì học thuộc mỗi chữ cái.
\end{question}

\begin{question}{351}
% [ID gốc: 5.57]
% [Nguồn: DBI202 Order 5 Q57]
Given the relation Employee(SSN. FNAME. LNAME. SALARY). Select the right query below to find the employee(s) who has the lowest salary in the company
\begin{choices}
\choice SELECT LNAME. FNAME. SALARY FROM Employee WHERE SALARY IN (SELECT MIN(SALARY) FROM Employee)
\choice SELECT LNAME. FNAME. SALARY FROM Employee WHERE SALARY >= ALL (SELECT SALARY FROM Employee)
\choice SELECT LNAME. FNAME. SALARY FROM Employee WHERE SALARY < MAX (SELECT SALARY FROM Employee)
\choice None of the others
\end{choices}
\textbf{Đáp án:} A
\textbf{Giải thích:} Đáp án theo dữ liệu đề: A — A. SELECT LNAME. FNAME. SALARY FROM Employee WHERE SALARY IN (SELECT MIN(SALARY) FROM Employee).

Nền tảng kiến thức: SQL tách truy vấn dữ liệu (DML) khỏi định nghĩa cấu trúc (DDL) và quyền truy cập (DCL). Với SELECT, cần xác định bảng nguồn, điều kiện lọc, phép kết hợp, nhóm và điều kiện nhóm; các ràng buộc, view, thủ tục và trigger đảm nhiệm các vai trò khác nhau.

Phân tích: Đọc truy vấn theo các giai đoạn logic: FROM/JOIN tạo dữ liệu nguồn, WHERE lọc từng hàng, GROUP BY tạo nhóm, HAVING lọc nhóm, SELECT xác định biểu thức đầu ra và ORDER BY sắp kết quả. Với phép join, kiểm tra điều kiện ghép và các hàng không khớp trước khi tính kết quả.

Đối chiếu phương án: B. SELECT LNAME. FNAME. SALARY FROM Employee WHERE SALARY >= ALL (SELECT SALARY FROM Employee); C. SELECT LNAME. FNAME. SALARY FROM Employee WHERE SALARY < MAX (SELECT SALARY FROM Employee); D. None of the others. Các lựa chọn này không đáp ứng định nghĩa/điều kiện được hỏi; khi ôn, hãy chỉ ra chính xác điểm lệch của từng lựa chọn thay vì học thuộc mỗi chữ cái.
\end{question}

\begin{question}{352}
% [ID gốc: 5.59]
% [Nguồn: DBI202 Order 5 Q59]
In SQL language, the command/statement that let you add an attribute to a relation is_
\begin{choices}
\choice Alter
\choice None of the others
\choice Insert
\choice Update
\end{choices}
\textbf{Đáp án:} A
\textbf{Giải thích:} Đáp án theo dữ liệu đề: A — A. Alter.

Nền tảng kiến thức: SQL tách truy vấn dữ liệu (DML) khỏi định nghĩa cấu trúc (DDL) và quyền truy cập (DCL). Với SELECT, cần xác định bảng nguồn, điều kiện lọc, phép kết hợp, nhóm và điều kiện nhóm; các ràng buộc, view, thủ tục và trigger đảm nhiệm các vai trò khác nhau.

Phân tích: Với lựa chọn tổng hợp, không chọn theo cảm tính: kiểm tra từng phát biểu độc lập. “All of the others/above” chỉ đúng khi mọi phát biểu còn lại đúng; “None” chỉ đúng khi không phát biểu nào trong nhóm thỏa yêu cầu. Đáp án lưu cần được đọc cùng các mệnh đề thành phần.

Đối chiếu phương án: B. None of the others; C. Insert; D. Update. Các lựa chọn này không đáp ứng định nghĩa/điều kiện được hỏi; khi ôn, hãy chỉ ra chính xác điểm lệch của từng lựa chọn thay vì học thuộc mỗi chữ cái.
\end{question}

\begin{question}{353}
% [ID gốc: 5.60]
% [Nguồn: DBI202 Order 5 Q60]
Select the valid query to declare the foreign key presC# of the relation Studiolname. address. presC#) that references the cert of the relation MovieExeclname. address. cert#. netWorth):
\begin{choices}
\choice All of the others.
\choice CREATE TABLE Studio (name CHAR(30) PRIMARY KEY. address VARCHAR(256). presC# INT FOREIGN KEY):
\choice CREATE TABLE Studio (name CHAR(30) PRIMARY KEY. address VARCHAR(256). presC# INT UNIQUE KEY REFERENCES MovieExec):
\choice CREATE TABLE Studio (name CHAR(30) PRIMARY KEY. address VARCHAR(256). presC# INT REFERENCES Movie Exec (cert#));
\end{choices}
\textbf{Đáp án:} D
\textbf{Giải thích:} Đáp án theo dữ liệu đề: D — D. CREATE TABLE Studio (name CHAR(30) PRIMARY KEY. address VARCHAR(256). presC# INT REFERENCES Movie Exec (cert#));.

Nền tảng kiến thức: SQL tách truy vấn dữ liệu (DML) khỏi định nghĩa cấu trúc (DDL) và quyền truy cập (DCL). Với SELECT, cần xác định bảng nguồn, điều kiện lọc, phép kết hợp, nhóm và điều kiện nhóm; các ràng buộc, view, thủ tục và trigger đảm nhiệm các vai trò khác nhau.

Phân tích: Khóa ngoại buộc giá trị tham chiếu phải khớp một khóa được tham chiếu ở quan hệ đích (hoặc NULL nếu ràng buộc cho phép). Đây là toàn vẹn tham chiếu; nó khác tính duy nhất của primary/candidate key.

Đối chiếu phương án: A. All of the others.; B. CREATE TABLE Studio (name CHAR(30) PRIMARY KEY. address VARCHAR(256). presC# INT FOREIGN KEY):; C. CREATE TABLE Studio (name CHAR(30) PRIMARY KEY. address VARCHAR(256). presC# INT UNIQUE KEY RE…. Các lựa chọn này không đáp ứng định nghĩa/điều kiện được hỏi; khi ôn, hãy chỉ ra chính xác điểm lệch của từng lựa chọn thay vì học thuộc mỗi chữ cái.
\end{question}

\begin{question}{354}
% [ID gốc: 5.61]
% [Nguồn: DBI202 Order 5 Q61]
Select the right statement
\begin{choices}
\choice All of the others
\choice Every constraint has a name. If we don't define constraint's name explicitly, then DBMS automatically generates a name for it
\choice We can create constraint on a tuple as a whole.
\choice We can create constraint on a single attribute
\end{choices}
\textbf{Đáp án:} A
\textbf{Giải thích:} Đáp án theo dữ liệu đề: A — A. All of the others.

Nền tảng kiến thức: SQL tách truy vấn dữ liệu (DML) khỏi định nghĩa cấu trúc (DDL) và quyền truy cập (DCL). Với SELECT, cần xác định bảng nguồn, điều kiện lọc, phép kết hợp, nhóm và điều kiện nhóm; các ràng buộc, view, thủ tục và trigger đảm nhiệm các vai trò khác nhau.

Phân tích: Đọc truy vấn theo các giai đoạn logic: FROM/JOIN tạo dữ liệu nguồn, WHERE lọc từng hàng, GROUP BY tạo nhóm, HAVING lọc nhóm, SELECT xác định biểu thức đầu ra và ORDER BY sắp kết quả. Với phép join, kiểm tra điều kiện ghép và các hàng không khớp trước khi tính kết quả.

Đối chiếu phương án: B. Every constraint has a name. If we don't define constraint's name explicitly, then DBMS automa…; C. We can create constraint on a tuple as a whole.; D. We can create constraint on a single attribute. Các lựa chọn này không đáp ứng định nghĩa/điều kiện được hỏi; khi ôn, hãy chỉ ra chính xác điểm lệch của từng lựa chọn thay vì học thuộc mỗi chữ cái.
\end{question}

\begin{question}{355}
% [ID gốc: 5.62]
% [Nguồn: DBI202 Order 5 Q62]
(FALSE AND NULL) return:
\begin{choices}
\choice NULL
\choice TRUE
\choice FALSE
\choice None of the others
a 
Give the relation Employee(SSN,FNAME,LNAME,SALARY).Select the right query below to find the employee(s) who has the lowest salary in the company
\choice SELECT LNAME,FNAME,SALARY FROM Employee WHERE SALARY IN ( SELECT MIN(SALARY) FROM Employee)
\choice SELECT LNAME,FNAME,SALARY FROM Employee WHERE SALARY >=ALL( SELECT MIN(SALARY) FROM Employee)
\choice SELECT LNAME,FNAME,SALARY FROM Employee WHERE SALARY < MAX ( SELECT MIN(SALARY) FROM Employee)
\choice None of the others
\end{choices}
\textbf{Đáp án:} C
\textbf{Giải thích:} Đáp án theo dữ liệu đề: A — A. SELECT LNAME,FNAME,SALARY FROM Employee WHERE SALARY IN ( SELECT MIN(SALARY) FROM Employee).

Nền tảng kiến thức: SQL tách truy vấn dữ liệu (DML) khỏi định nghĩa cấu trúc (DDL) và quyền truy cập (DCL). Với SELECT, cần xác định bảng nguồn, điều kiện lọc, phép kết hợp, nhóm và điều kiện nhóm; các ràng buộc, view, thủ tục và trigger đảm nhiệm các vai trò khác nhau.

Phân tích: NULL biểu thị thiếu/không biết/không áp dụng, không phải số 0 hay chuỗi rỗng. So sánh với NULL bằng `=` không cho TRUE/FALSE mà cho UNKNOWN; dùng `IS NULL`/`IS NOT NULL`. Trong logic ba giá trị, AND/OR/NOT phải được tính theo cả trạng thái UNKNOWN.

Đối chiếu phương án: B. TRUE; C. FALSE; D. None of the others a Give the relation Employee(SSN,FNAME,LNAME,SALARY).Select the right query…; B. SELECT LNAME,FNAME,SALARY FROM Employee WHERE SALARY >=ALL( SELECT MIN(SALARY) FROM Employee); C. SELECT LNAME,FNAME,SALARY FROM Employee WHERE SALARY < MAX ( SELECT MIN(SALARY) FROM Employee). Các lựa chọn này không đáp ứng định nghĩa/điều kiện được hỏi; khi ôn, hãy chỉ ra chính xác điểm lệch của từng lựa chọn thay vì học thuộc mỗi chữ cái.
\end{question}

\begin{question}{356}
% [ID gốc: 5.64]
% [Nguồn: DBI202 Order 5 Q64]
What SQL keyword can be use to return data from two non-related tables as a combined set of rows?
\begin{choices}
\choice DISTINCT ALL
\choice DISTINCT
\choice UNION ALL
\choice COMBINED
\end{choices}
\textbf{Đáp án:} C
\textbf{Giải thích:} Đáp án theo dữ liệu đề: C — C. UNION ALL.

Nền tảng kiến thức: SQL tách truy vấn dữ liệu (DML) khỏi định nghĩa cấu trúc (DDL) và quyền truy cập (DCL). Với SELECT, cần xác định bảng nguồn, điều kiện lọc, phép kết hợp, nhóm và điều kiện nhóm; các ràng buộc, view, thủ tục và trigger đảm nhiệm các vai trò khác nhau.

Phân tích: Trong T-SQL, biến cục bộ do người dùng khai báo bắt đầu bằng một dấu `@` (ví dụ `@count`); tên bắt đầu `@@` được dùng cho các hàm/biến hệ thống. Đừng nhầm tiền tố biến cục bộ với các biến hệ thống như `@@ROWCOUNT`.

Đối chiếu phương án: A. DISTINCT ALL; B. DISTINCT; D. COMBINED. Các lựa chọn này không đáp ứng định nghĩa/điều kiện được hỏi; khi ôn, hãy chỉ ra chính xác điểm lệch của từng lựa chọn thay vì học thuộc mỗi chữ cái.
\end{question}

\begin{question}{357}
% [ID gốc: 5.67]
% [Nguồn: DBI202 Order 5 Q67]
This simple SQL queries, uses the three keywords that characterize SQL is ____ ?
\begin{choices}
\choice Count, Sum and Avg
\choice Insert, Update and Delete
\choice Select, From, and Where
\choice And, Or, and Not
\end{choices}
\textbf{Đáp án:} C
\textbf{Giải thích:} Đáp án theo dữ liệu đề: C — C. Select, From, and Where.

Nền tảng kiến thức: SQL tách truy vấn dữ liệu (DML) khỏi định nghĩa cấu trúc (DDL) và quyền truy cập (DCL). Với SELECT, cần xác định bảng nguồn, điều kiện lọc, phép kết hợp, nhóm và điều kiện nhóm; các ràng buộc, view, thủ tục và trigger đảm nhiệm các vai trò khác nhau.

Phân tích: Đọc truy vấn theo các giai đoạn logic: FROM/JOIN tạo dữ liệu nguồn, WHERE lọc từng hàng, GROUP BY tạo nhóm, HAVING lọc nhóm, SELECT xác định biểu thức đầu ra và ORDER BY sắp kết quả. Với phép join, kiểm tra điều kiện ghép và các hàng không khớp trước khi tính kết quả.

Đối chiếu phương án: A. Count, Sum and Avg; B. Insert, Update and Delete; D. And, Or, and Not. Các lựa chọn này không đáp ứng định nghĩa/điều kiện được hỏi; khi ôn, hãy chỉ ra chính xác điểm lệch của từng lựa chọn thay vì học thuộc mỗi chữ cái.
\end{question}

\begin{question}{358}
% [ID gốc: 5.68]
% [Nguồn: DBI202 Order 5 Q68]
To modify the schema of an existing relation use ______
\begin{choices}
\choice Create table
\choice Modify table
\choice Alter table
\choice Drop table
\end{choices}
\textbf{Đáp án:} C
\textbf{Giải thích:} Đáp án theo dữ liệu đề: C — C. Alter table.

Nền tảng kiến thức: SQL tách truy vấn dữ liệu (DML) khỏi định nghĩa cấu trúc (DDL) và quyền truy cập (DCL). Với SELECT, cần xác định bảng nguồn, điều kiện lọc, phép kết hợp, nhóm và điều kiện nhóm; các ràng buộc, view, thủ tục và trigger đảm nhiệm các vai trò khác nhau.

Phân tích: ALTER TABLE thay đổi định nghĩa bảng đã tồn tại, như thêm cột hoặc sửa cấu hình cột/ràng buộc theo cú pháp của DBMS. Nó khác UPDATE, vốn sửa các giá trị trong những hàng hiện hữu; cú pháp chính xác có thể khác giữa SQL Server và chuẩn SQL.

Đối chiếu phương án: A. Create table; B. Modify table; D. Drop table. Các lựa chọn này không đáp ứng định nghĩa/điều kiện được hỏi; khi ôn, hãy chỉ ra chính xác điểm lệch của từng lựa chọn thay vì học thuộc mỗi chữ cái.
\end{question}

\begin{question}{359}
% [ID gốc: 5.69]
% [Nguồn: DBI202 Order 5 Q69]
Choose the incorrect statement about DELETE and TRUNCAT in SQL Server.
\begin{choices}
\choice DELETE is used for unconditional removal of data records from Tables; TRUNCATE is used for conditional removal of data records from Tables.
\choice DELETE is a DML command and the operations are logged, whereas TRUNCATE is a DDL command and the operations are not logged.
\choice DELETE removes records and records each deletion in the transaction log whereas TRUNCATE deallocates pages and records each deallocation in the transaction log.
\choice TRUNCATE is generally considered quicker as it makes less use of the transaction log
\end{choices}
\textbf{Đáp án:} A
\textbf{Giải thích:} Đáp án theo dữ liệu đề: A — A. DELETE is used for unconditional removal of data records from Tables; TRUNCATE is used for conditional removal of data records from Tables..

Nền tảng kiến thức: SQL tách truy vấn dữ liệu (DML) khỏi định nghĩa cấu trúc (DDL) và quyền truy cập (DCL). Với SELECT, cần xác định bảng nguồn, điều kiện lọc, phép kết hợp, nhóm và điều kiện nhóm; các ràng buộc, view, thủ tục và trigger đảm nhiệm các vai trò khác nhau.

Phân tích: DDL (Data Definition Language) khai báo/thay đổi cấu trúc như bảng và schema; DML truy vấn/thêm/sửa/xóa dữ liệu; DCL quản lý quyền. Chọn theo đối tượng bị tác động: cấu trúc, các hàng hay quyền truy cập.

Đối chiếu phương án: B. DELETE is a DML command and the operations are logged, whereas TRUNCATE is a DDL command and t…; C. DELETE removes records and records each deletion in the transaction log whereas TRUNCATE deall…; D. TRUNCATE is generally considered quicker as it makes less use of the transaction log. Các lựa chọn này không đáp ứng định nghĩa/điều kiện được hỏi; khi ôn, hãy chỉ ra chính xác điểm lệch của từng lựa chọn thay vì học thuộc mỗi chữ cái.
\end{question}

\begin{question}{360}
% [ID gốc: 5.74]
% [Nguồn: DBI202 Order 5 Q74]
Which command is used to remove a relation from an SQL?
\begin{choices}
\choice Drop table
\choice Delete
\choice Purge
\choice Remove
\end{choices}
\textbf{Đáp án:} A
\textbf{Giải thích:} Đáp án theo dữ liệu đề: A — A. Drop table.

Nền tảng kiến thức: SQL tách truy vấn dữ liệu (DML) khỏi định nghĩa cấu trúc (DDL) và quyền truy cập (DCL). Với SELECT, cần xác định bảng nguồn, điều kiện lọc, phép kết hợp, nhóm và điều kiện nhóm; các ràng buộc, view, thủ tục và trigger đảm nhiệm các vai trò khác nhau.

Phân tích: DROP TABLE loại bỏ định nghĩa bảng cùng dữ liệu phụ thuộc theo quy tắc DBMS; DELETE xóa các hàng nhưng giữ bảng, còn TRUNCATE xóa nhanh toàn bộ hàng trong khi vẫn giữ cấu trúc. Cần chọn lệnh đúng cấp độ tác động.

Đối chiếu phương án: B. Delete; C. Purge; D. Remove. Các lựa chọn này không đáp ứng định nghĩa/điều kiện được hỏi; khi ôn, hãy chỉ ra chính xác điểm lệch của từng lựa chọn thay vì học thuộc mỗi chữ cái.
\end{question}

\begin{question}{361}
% [ID gốc: 5.75]
% [Nguồn: DBI202 Order 5 Q75]
In an employee table to include the attributes whose value always have some value which of the following constraint must be used?
\begin{choices}
\choice Null
\choice Not null
\choice Unique
\choice Distinct
\end{choices}
\textbf{Đáp án:} B
\textbf{Giải thích:} Đáp án theo dữ liệu đề: B — B. Not null.

Nền tảng kiến thức: SQL tách truy vấn dữ liệu (DML) khỏi định nghĩa cấu trúc (DDL) và quyền truy cập (DCL). Với SELECT, cần xác định bảng nguồn, điều kiện lọc, phép kết hợp, nhóm và điều kiện nhóm; các ràng buộc, view, thủ tục và trigger đảm nhiệm các vai trò khác nhau.

Phân tích: NOT NULL là constraint cấm giá trị NULL ở cột, phù hợp khi mỗi hàng bắt buộc phải có giá trị. Nó không tự bảo đảm giá trị duy nhất (UNIQUE/PRIMARY KEY) và cũng không kiểm tra miền/range như CHECK.

Đối chiếu phương án: A. Null; C. Unique; D. Distinct. Các lựa chọn này không đáp ứng định nghĩa/điều kiện được hỏi; khi ôn, hãy chỉ ra chính xác điểm lệch của từng lựa chọn thay vì học thuộc mỗi chữ cái.
\end{question}

\begin{question}{362}
% [ID gốc: 5.76]
% [Nguồn: DBI202 Order 5 Q76]
Choose the correct statement that belongs to DML.
\begin{choices}
\choice SELECT, DELETE, INSERT
\choice CREATE, UPDATE, SELECT
\choice TRUNCAT, DELETE, UPDATE
\choice UPDATE, DELETE, ALTER
\end{choices}
\textbf{Đáp án:} A
\textbf{Giải thích:} Đáp án theo dữ liệu đề: A — A. SELECT, DELETE, INSERT.

Nền tảng kiến thức: SQL tách truy vấn dữ liệu (DML) khỏi định nghĩa cấu trúc (DDL) và quyền truy cập (DCL). Với SELECT, cần xác định bảng nguồn, điều kiện lọc, phép kết hợp, nhóm và điều kiện nhóm; các ràng buộc, view, thủ tục và trigger đảm nhiệm các vai trò khác nhau.

Phân tích: DDL (Data Definition Language) khai báo/thay đổi cấu trúc như bảng và schema; DML truy vấn/thêm/sửa/xóa dữ liệu; DCL quản lý quyền. Chọn theo đối tượng bị tác động: cấu trúc, các hàng hay quyền truy cập.

Đối chiếu phương án: B. CREATE, UPDATE, SELECT; C. TRUNCAT, DELETE, UPDATE; D. UPDATE, DELETE, ALTER. Các lựa chọn này không đáp ứng định nghĩa/điều kiện được hỏi; khi ôn, hãy chỉ ra chính xác điểm lệch của từng lựa chọn thay vì học thuộc mỗi chữ cái.
\end{question}

\begin{question}{363}
% [ID gốc: 5.77]
% [Nguồn: DBI202 Order 5 Q77]
A Delete command operates on _____ relation.
\begin{choices}
\choice One
\choice Two
\choice Several
\choice Null
\end{choices}
\textbf{Đáp án:} A
\textbf{Giải thích:} Đáp án theo dữ liệu đề: A — A. One.

Nền tảng kiến thức: SQL tách truy vấn dữ liệu (DML) khỏi định nghĩa cấu trúc (DDL) và quyền truy cập (DCL). Với SELECT, cần xác định bảng nguồn, điều kiện lọc, phép kết hợp, nhóm và điều kiện nhóm; các ràng buộc, view, thủ tục và trigger đảm nhiệm các vai trò khác nhau.

Phân tích: Một câu lệnh DELETE tác động lên một target table/relation trong mỗi câu lệnh, dù có thể xóa nhiều hàng của bảng đó theo WHERE. Các hàng tham chiếu ở bảng khác chỉ bị ảnh hưởng nếu có cascade/trigger được cấu hình.

Đối chiếu phương án: B. Two; C. Several; D. Null. Các lựa chọn này không đáp ứng định nghĩa/điều kiện được hỏi; khi ôn, hãy chỉ ra chính xác điểm lệch của từng lựa chọn thay vì học thuộc mỗi chữ cái.
\end{question}

\begin{question}{364}
% [ID gốc: 5.78]
% [Nguồn: DBI202 Order 5 Q78]
You have a table named Books that has columns named Book Title and Description. There is a full-text index on these columns. You need to return rows from the table in which the word 'computer' exists in either column. Which code segment should you use?
\begin{choices}
\choice SELECT * FROM Books WHERE FREETEXT(*, 'computer')
\choice SELECT * FROM Books WHERE Book Title LIKE '\%computer%'
\choice SELECT * FROM Books WHERE Book Title = '\%computer%' OR Description = '\%computer%'
\choice SELECT * FROM Books WHERE FREETEXT(Book Title, 'computer')
\end{choices}
\textbf{Đáp án:} A
\textbf{Giải thích:} Đáp án theo dữ liệu đề: A — A. SELECT * FROM Books WHERE FREETEXT(*, 'computer').

Nền tảng kiến thức: SQL tách truy vấn dữ liệu (DML) khỏi định nghĩa cấu trúc (DDL) và quyền truy cập (DCL). Với SELECT, cần xác định bảng nguồn, điều kiện lọc, phép kết hợp, nhóm và điều kiện nhóm; các ràng buộc, view, thủ tục và trigger đảm nhiệm các vai trò khác nhau.

Phân tích: Index lưu cấu trúc tra cứu phụ để tìm vị trí bản ghi nhanh hơn thay vì quét toàn bộ bảng. Bù lại, nó dùng thêm dung lượng và phải được duy trì khi INSERT/UPDATE/DELETE; vì thế “càng nhiều index càng tốt” không đúng cho mọi workload.

Đối chiếu phương án: B. SELECT * FROM Books WHERE Book Title LIKE '\%computer%'; C. SELECT * FROM Books WHERE Book Title = '\%computer%' OR Description = '\%computer%'; D. SELECT * FROM Books WHERE FREETEXT(Book Title, 'computer'). Các lựa chọn này không đáp ứng định nghĩa/điều kiện được hỏi; khi ôn, hãy chỉ ra chính xác điểm lệch của từng lựa chọn thay vì học thuộc mỗi chữ cái.
\end{question}

\begin{question}{365}
% [ID gốc: 5.79]
% [Nguồn: DBI202 Order 5 Q79]
In SQL, to update a relation's schema, which one of the following statements can be used?
\begin{choices}
\choice Insert
\choice Update
\choice Alter
\choice Alter table
\end{choices}
\textbf{Đáp án:} D
\textbf{Giải thích:} Đáp án theo dữ liệu đề: D — D. Alter table.

Nền tảng kiến thức: SQL tách truy vấn dữ liệu (DML) khỏi định nghĩa cấu trúc (DDL) và quyền truy cập (DCL). Với SELECT, cần xác định bảng nguồn, điều kiện lọc, phép kết hợp, nhóm và điều kiện nhóm; các ràng buộc, view, thủ tục và trigger đảm nhiệm các vai trò khác nhau.

Phân tích: ALTER TABLE thay đổi định nghĩa bảng đã tồn tại, như thêm cột hoặc sửa cấu hình cột/ràng buộc theo cú pháp của DBMS. Nó khác UPDATE, vốn sửa các giá trị trong những hàng hiện hữu; cú pháp chính xác có thể khác giữa SQL Server và chuẩn SQL.

Đối chiếu phương án: A. Insert; B. Update; C. Alter. Các lựa chọn này không đáp ứng định nghĩa/điều kiện được hỏi; khi ôn, hãy chỉ ra chính xác điểm lệch của từng lựa chọn thay vì học thuộc mỗi chữ cái.
\end{question}

\begin{question}{366}
% [ID gốc: 5.80]
% [Nguồn: DBI202 Order 5 Q80]
The DROP TABLE statement:
\begin{choices}
\choice deletes the table structure only.
\choice deletes the table structure along with the table data.
\choice works whether or not referential integrity constraints would be violated.
\choice is not an SQL statement.
\end{choices}
\textbf{Đáp án:} B
\textbf{Giải thích:} Đáp án theo dữ liệu đề: B — B. deletes the table structure along with the table data..

Nền tảng kiến thức: SQL tách truy vấn dữ liệu (DML) khỏi định nghĩa cấu trúc (DDL) và quyền truy cập (DCL). Với SELECT, cần xác định bảng nguồn, điều kiện lọc, phép kết hợp, nhóm và điều kiện nhóm; các ràng buộc, view, thủ tục và trigger đảm nhiệm các vai trò khác nhau.

Phân tích: DROP TABLE loại bỏ định nghĩa bảng cùng dữ liệu phụ thuộc theo quy tắc DBMS; DELETE xóa các hàng nhưng giữ bảng, còn TRUNCATE xóa nhanh toàn bộ hàng trong khi vẫn giữ cấu trúc. Cần chọn lệnh đúng cấp độ tác động.

Đối chiếu phương án: A. deletes the table structure only.; C. works whether or not referential integrity constraints would be violated.; D. is not an SQL statement.. Các lựa chọn này không đáp ứng định nghĩa/điều kiện được hỏi; khi ôn, hãy chỉ ra chính xác điểm lệch của từng lựa chọn thay vì học thuộc mỗi chữ cái.
\end{question}

\begin{question}{367}
% [ID gốc: 5.81]
% [Nguồn: DBI202 Order 5 Q81]
Which of the following aggregate functions does not ignore nulls in its results?
\begin{choices}
\choice MIN
\choice MAX
\choice COUNT (*)
\choice COUNT
\end{choices}
\textbf{Đáp án:} C
\textbf{Giải thích:} Đáp án theo dữ liệu đề: C — C. COUNT (*).

Nền tảng kiến thức: SQL tách truy vấn dữ liệu (DML) khỏi định nghĩa cấu trúc (DDL) và quyền truy cập (DCL). Với SELECT, cần xác định bảng nguồn, điều kiện lọc, phép kết hợp, nhóm và điều kiện nhóm; các ràng buộc, view, thủ tục và trigger đảm nhiệm các vai trò khác nhau.

Phân tích: Aggregate function tính một kết quả từ một tập nhiều hàng/giá trị, ví dụ COUNT, SUM, AVG, MIN và MAX. Trong truy vấn nhóm, GROUP BY xác định từng tập đầu vào; HAVING lọc các nhóm sau khi tính tổng hợp.

Đối chiếu phương án: A. MIN; B. MAX; D. COUNT. Các lựa chọn này không đáp ứng định nghĩa/điều kiện được hỏi; khi ôn, hãy chỉ ra chính xác điểm lệch của từng lựa chọn thay vì học thuộc mỗi chữ cái.
\end{question}

\begin{question}{368}
% [ID gốc: 5.82]
% [Nguồn: DBI202 Order 5 Q82]
Which of the following statements is false regarding SQL Correlated sub-queries?
\begin{choices}
\choice A correlated subquery is not evaluated once for each row processed by the parent statement.
\choice Correlated subqueries are used for row-by-row processing.
\choice Each subquery is executed once for every row of the outer query.
\choice Subqueries always process the innermost query first and the work outward.
\end{choices}
\textbf{Đáp án:} A
\textbf{Giải thích:} Đáp án theo dữ liệu đề: A — A. A correlated subquery is not evaluated once for each row processed by the parent statement..

Nền tảng kiến thức: SQL tách truy vấn dữ liệu (DML) khỏi định nghĩa cấu trúc (DDL) và quyền truy cập (DCL). Với SELECT, cần xác định bảng nguồn, điều kiện lọc, phép kết hợp, nhóm và điều kiện nhóm; các ràng buộc, view, thủ tục và trigger đảm nhiệm các vai trò khác nhau.

Phân tích: Đọc truy vấn theo các giai đoạn logic: FROM/JOIN tạo dữ liệu nguồn, WHERE lọc từng hàng, GROUP BY tạo nhóm, HAVING lọc nhóm, SELECT xác định biểu thức đầu ra và ORDER BY sắp kết quả. Với phép join, kiểm tra điều kiện ghép và các hàng không khớp trước khi tính kết quả.

Đối chiếu phương án: B. Correlated subqueries are used for row-by-row processing.; C. Each subquery is executed once for every row of the outer query.; D. Subqueries always process the innermost query first and the work outward.. Các lựa chọn này không đáp ứng định nghĩa/điều kiện được hỏi; khi ôn, hãy chỉ ra chính xác điểm lệch của từng lựa chọn thay vì học thuộc mỗi chữ cái.
\end{question}

\begin{question}{369}
% [ID gốc: 5.83]
% [Nguồn: DBI202 Order 5 Q83]
Suppose relation \textbf{R(A, B)} has the following tuples:

$$\begin{array}{|c|c|}
\hline
\textbf{A} & \textbf{B} \\
\hline
1 & \text{a} \\
\hline
2 & \text{b} \\
\hline
3 & \text{a} \\
\hline
4 & \text{d} \\
\hline
5 & \text{e} \\
\hline
6 & \text{a} \\
\hline
7 & \text{g} \\
\hline
\end{array}$$

Choose the SQL statement to count distinct rows from column B in relation R:
\begin{choices}
\choice SELECT COUNT DISTINCT B AS 'Total_rows' FROM R;
\choice SELECT (COUNT DISTINCT B) AS 'Total_rows' from R;
\choice SELECT DISTINCT B AS 'Total_rows' from R;
\choice SELECT COUNT (DISTINCT(B)) AS 'total_rows' FROM R;
\end{choices}
\textbf{Đáp án:} D
\textbf{Giải thích:} Trong SQL, cú pháp đếm số giá trị phân biệt (loại bỏ trùng lặp) của một cột là COUNT(DISTINCT cột). Do đó SELECT COUNT(DISTINCT(B)) AS 'total_rows' FROM R; là câu lệnh chính xác.
\end{question}

\begin{question}{370}
% [ID gốc: 5.84]
% [Nguồn: DBI202 Order 5 Q84]
Which statement is used to change the data type of a column named E in table A?
\begin{choices}
\choice ALTER TABLE A ALTER COLUMN E [DataType]
\choice ALTER TABLE A ALTER COLUMN E
\choice ALTER TABLE A ALTER E [DataType]
\choice ALTER TABLE A DROP COLUMN E [DataType]
\end{choices}
\textbf{Đáp án:} A
\textbf{Giải thích:} Đáp án theo dữ liệu đề: A — A. ALTER TABLE A ALTER COLUMN E [DataType].

Nền tảng kiến thức: SQL tách truy vấn dữ liệu (DML) khỏi định nghĩa cấu trúc (DDL) và quyền truy cập (DCL). Với SELECT, cần xác định bảng nguồn, điều kiện lọc, phép kết hợp, nhóm và điều kiện nhóm; các ràng buộc, view, thủ tục và trigger đảm nhiệm các vai trò khác nhau.

Phân tích: ALTER TABLE thay đổi định nghĩa bảng đã tồn tại, như thêm cột hoặc sửa cấu hình cột/ràng buộc theo cú pháp của DBMS. Nó khác UPDATE, vốn sửa các giá trị trong những hàng hiện hữu; cú pháp chính xác có thể khác giữa SQL Server và chuẩn SQL.

Đối chiếu phương án: B. ALTER TABLE A ALTER COLUMN E; C. ALTER TABLE A ALTER E [DataType]; D. ALTER TABLE A DROP COLUMN E [DataType]. Các lựa chọn này không đáp ứng định nghĩa/điều kiện được hỏi; khi ôn, hãy chỉ ra chính xác điểm lệch của từng lựa chọn thay vì học thuộc mỗi chữ cái.
\end{question}

\begin{question}{371}
% [ID gốc: 5.85]
% [Nguồn: DBI202 Order 5 Q85]
The ability to query data, as well as insert, delete, and alter tuples, is offered by ____
\begin{choices}
\choice TCL (Transaction Control Language)
\choice DCL (Data Control Language)
\choice DDL (Data Definition Language)
\choice DML (Data Manipulation Language)
\end{choices}
\textbf{Đáp án:} D
\textbf{Giải thích:} Đáp án theo dữ liệu đề: D — D. DML (Data Manipulation Language).

Nền tảng kiến thức: SQL tách truy vấn dữ liệu (DML) khỏi định nghĩa cấu trúc (DDL) và quyền truy cập (DCL). Với SELECT, cần xác định bảng nguồn, điều kiện lọc, phép kết hợp, nhóm và điều kiện nhóm; các ràng buộc, view, thủ tục và trigger đảm nhiệm các vai trò khác nhau.

Phân tích: Các hàng dữ liệu của một relation được gọi là tuple; header chứa tên thuộc tính không phải tuple dữ liệu. Trong bảng, tuple tương ứng một bản ghi, còn attribute tương ứng cột.

Đối chiếu phương án: A. TCL (Transaction Control Language); B. DCL (Data Control Language); C. DDL (Data Definition Language). Các lựa chọn này không đáp ứng định nghĩa/điều kiện được hỏi; khi ôn, hãy chỉ ra chính xác điểm lệch của từng lựa chọn thay vì học thuộc mỗi chữ cái.
\end{question}

\begin{question}{372}
% [ID gốc: 5.86]
% [Nguồn: DBI202 Order 5 Q86]
Which of the following is correct?
\begin{choices}
\choice SELECT datepart(DATE, '10-jan-24')
\choice SELECT datepart('10-jan-24', DATE)
\choice SELECT datepart('10-jan-24', DAY)
\choice SELECT datepart(DAY,'10-jan-24')
\end{choices}
\textbf{Đáp án:} D
\textbf{Giải thích:} Đáp án theo dữ liệu đề: D — D. SELECT datepart(DAY,'10-jan-24').

Nền tảng kiến thức: SQL tách truy vấn dữ liệu (DML) khỏi định nghĩa cấu trúc (DDL) và quyền truy cập (DCL). Với SELECT, cần xác định bảng nguồn, điều kiện lọc, phép kết hợp, nhóm và điều kiện nhóm; các ràng buộc, view, thủ tục và trigger đảm nhiệm các vai trò khác nhau.

Phân tích: Đọc truy vấn theo các giai đoạn logic: FROM/JOIN tạo dữ liệu nguồn, WHERE lọc từng hàng, GROUP BY tạo nhóm, HAVING lọc nhóm, SELECT xác định biểu thức đầu ra và ORDER BY sắp kết quả. Với phép join, kiểm tra điều kiện ghép và các hàng không khớp trước khi tính kết quả.

Đối chiếu phương án: A. SELECT datepart(DATE, '10-jan-24'); B. SELECT datepart('10-jan-24', DATE); C. SELECT datepart('10-jan-24', DAY). Các lựa chọn này không đáp ứng định nghĩa/điều kiện được hỏi; khi ôn, hãy chỉ ra chính xác điểm lệch của từng lựa chọn thay vì học thuộc mỗi chữ cái.
\end{question}

\begin{question}{373}
% [ID gốc: 5.88]
% [Nguồn: DBI202 Order 5 Q88]
You have a column named TelephoneNumber that stores numbers as varchar(20). You need to write a query that returns the first three characters of a telephone number. Which expression should you use?
\begin{choices}
\choice CHARINDEX('[0-9][0-9][0-9]', TelephoneNumber, 3)
\choice SUBSTRING (Telephone Number, 3, 1)
\choice SUBSTRING(TelephoneNumber, 3, 3)
\choice LEFT(TelephoneNumber, 3)
\end{choices}
\textbf{Đáp án:} D
\textbf{Giải thích:} Đáp án theo dữ liệu đề: D — D. LEFT(TelephoneNumber, 3).

Nền tảng kiến thức: SQL tách truy vấn dữ liệu (DML) khỏi định nghĩa cấu trúc (DDL) và quyền truy cập (DCL). Với SELECT, cần xác định bảng nguồn, điều kiện lọc, phép kết hợp, nhóm và điều kiện nhóm; các ràng buộc, view, thủ tục và trigger đảm nhiệm các vai trò khác nhau.

Phân tích: Đọc truy vấn theo các giai đoạn logic: FROM/JOIN tạo dữ liệu nguồn, WHERE lọc từng hàng, GROUP BY tạo nhóm, HAVING lọc nhóm, SELECT xác định biểu thức đầu ra và ORDER BY sắp kết quả. Với phép join, kiểm tra điều kiện ghép và các hàng không khớp trước khi tính kết quả.

Đối chiếu phương án: A. CHARINDEX('[0-9][0-9][0-9]', TelephoneNumber, 3); B. SUBSTRING (Telephone Number, 3, 1); C. SUBSTRING(TelephoneNumber, 3, 3). Các lựa chọn này không đáp ứng định nghĩa/điều kiện được hỏi; khi ôn, hãy chỉ ra chính xác điểm lệch của từng lựa chọn thay vì học thuộc mỗi chữ cái.
\end{question}

\begin{question}{374}
% [ID gốc: 5.89]
% [Nguồn: DBI202 Order 5 Q89]
Which SQL statement selects all rows from table called Contest, with column ContestDate having values greater or equal to May 25,2006?
\begin{choices}
\choice SELECT * FROM Contest WHERE ContestDate >= '05/25/2006'
\choice SELECT * FROM Contest GROUPBY ContestDate >= '05/25/2006'.
\choice SELECT * FROM Contest WHERE ContestDate = '05/25/2006'.
\end{choices}
\textbf{Đáp án:} A
\textbf{Giải thích:} Đáp án theo dữ liệu đề: A — A. SELECT * FROM Contest WHERE ContestDate >= '05/25/2006'.

Nền tảng kiến thức: SQL tách truy vấn dữ liệu (DML) khỏi định nghĩa cấu trúc (DDL) và quyền truy cập (DCL). Với SELECT, cần xác định bảng nguồn, điều kiện lọc, phép kết hợp, nhóm và điều kiện nhóm; các ràng buộc, view, thủ tục và trigger đảm nhiệm các vai trò khác nhau.

Phân tích: Đọc truy vấn theo các giai đoạn logic: FROM/JOIN tạo dữ liệu nguồn, WHERE lọc từng hàng, GROUP BY tạo nhóm, HAVING lọc nhóm, SELECT xác định biểu thức đầu ra và ORDER BY sắp kết quả. Với phép join, kiểm tra điều kiện ghép và các hàng không khớp trước khi tính kết quả.

Đối chiếu phương án: B. SELECT * FROM Contest GROUPBY ContestDate >= '05/25/2006'.; C. SELECT * FROM Contest WHERE ContestDate = '05/25/2006'.. Các lựa chọn này không đáp ứng định nghĩa/điều kiện được hỏi; khi ôn, hãy chỉ ra chính xác điểm lệch của từng lựa chọn thay vì học thuộc mỗi chữ cái.
\end{question}

\begin{question}{375}
% [ID gốc: 5.90]
% [Nguồn: DBI202 Order 5 Q90]
The latter clause consists of the keyword followed by a condition about the group.
\begin{choices}
\choice Having
\choice Group-by
\choice Order by
\choice Where
\end{choices}
\textbf{Đáp án:} A
\textbf{Giải thích:} Đáp án theo dữ liệu đề: A — A. Having.

Nền tảng kiến thức: SQL tách truy vấn dữ liệu (DML) khỏi định nghĩa cấu trúc (DDL) và quyền truy cập (DCL). Với SELECT, cần xác định bảng nguồn, điều kiện lọc, phép kết hợp, nhóm và điều kiện nhóm; các ràng buộc, view, thủ tục và trigger đảm nhiệm các vai trò khác nhau.

Phân tích: Đọc truy vấn theo các giai đoạn logic: FROM/JOIN tạo dữ liệu nguồn, WHERE lọc từng hàng, GROUP BY tạo nhóm, HAVING lọc nhóm, SELECT xác định biểu thức đầu ra và ORDER BY sắp kết quả. Với phép join, kiểm tra điều kiện ghép và các hàng không khớp trước khi tính kết quả.

Đối chiếu phương án: B. Group-by; C. Order by; D. Where. Các lựa chọn này không đáp ứng định nghĩa/điều kiện được hỏi; khi ôn, hãy chỉ ra chính xác điểm lệch của từng lựa chọn thay vì học thuộc mỗi chữ cái.
\end{question}

\begin{question}{376}
% [ID gốc: 5.91]
% [Nguồn: DBI202 Order 5 Q91]
In MS-SQL Server, which of the following is true about join operator?
SELECT * FROM tbICUSTOMER, tbIORDER
\begin{choices}
\choice Cartesian join
\choice Equi-join
\choice Natural join
\choice Outer join
\end{choices}
\textbf{Đáp án:} A
\textbf{Giải thích:} Đáp án theo dữ liệu đề: A — A. Cartesian join.

Nền tảng kiến thức: SQL tách truy vấn dữ liệu (DML) khỏi định nghĩa cấu trúc (DDL) và quyền truy cập (DCL). Với SELECT, cần xác định bảng nguồn, điều kiện lọc, phép kết hợp, nhóm và điều kiện nhóm; các ràng buộc, view, thủ tục và trigger đảm nhiệm các vai trò khác nhau.

Phân tích: Join ghép các bộ thỏa điều kiện. Inner join loại bộ không ghép được; LEFT/RIGHT OUTER JOIN giữ mọi bộ từ phía được chỉ định và điền NULL phía thiếu. Natural join tự ghép các thuộc tính trùng tên và chỉ xuất một cột cho mỗi thuộc tính ghép.

Đối chiếu phương án: B. Equi-join; C. Natural join; D. Outer join. Các lựa chọn này không đáp ứng định nghĩa/điều kiện được hỏi; khi ôn, hãy chỉ ra chính xác điểm lệch của từng lựa chọn thay vì học thuộc mỗi chữ cái.
\end{question}

\begin{question}{377}
% [ID gốc: 5.92]
% [Nguồn: DBI202 Order 5 Q92]
UPDATE 
____ salary= salary * 1.05;
Fill in with correct keyword to update the instructor relation.
\begin{choices}
\choice Where
\choice Set
\choice In
\choice Select
\end{choices}
\textbf{Đáp án:} B
\textbf{Giải thích:} Đáp án theo dữ liệu đề: B — B. Set.

Nền tảng kiến thức: SQL tách truy vấn dữ liệu (DML) khỏi định nghĩa cấu trúc (DDL) và quyền truy cập (DCL). Với SELECT, cần xác định bảng nguồn, điều kiện lọc, phép kết hợp, nhóm và điều kiện nhóm; các ràng buộc, view, thủ tục và trigger đảm nhiệm các vai trò khác nhau.

Phân tích: Đọc truy vấn theo các giai đoạn logic: FROM/JOIN tạo dữ liệu nguồn, WHERE lọc từng hàng, GROUP BY tạo nhóm, HAVING lọc nhóm, SELECT xác định biểu thức đầu ra và ORDER BY sắp kết quả. Với phép join, kiểm tra điều kiện ghép và các hàng không khớp trước khi tính kết quả.

Đối chiếu phương án: A. Where; C. In; D. Select. Các lựa chọn này không đáp ứng định nghĩa/điều kiện được hỏi; khi ôn, hãy chỉ ra chính xác điểm lệch của từng lựa chọn thay vì học thuộc mỗi chữ cái.
\end{question}

\begin{question}{378}
% [ID gốc: 5.93]
% [Nguồn: DBI202 Order 5 Q93]
Which one of the following SQL statements returns the name of the employee (employeename, dept_code,
salary) receiving the maximum salary in a particular department?
\begin{choices}
\choice Select employeename, dept_code, salary From employee Where employee.salary in (select (salary) from Employee group by dept_code having Max(salary))
\choice Select employeename, dept_code, salary From employee Where employee.salary in (select Max(salary)from Employee group by dept_code)
\choice Select employeename, dept_code, salary From employee Where employee.salary in (select Max(salary)from Employee group by dept_code having Max(salary))
\choice Select employeename, dept_code, salary From employee Where employee.salary =(select Max(salary)from Employee group by dept_code
\end{choices}
\textbf{Đáp án:} B
\textbf{Giải thích:} Đáp án theo dữ liệu đề: B — B. Select employeename, dept_code, salary From employee Where employee.salary in (select Max(salary)from Employee group by dept_code).

Nền tảng kiến thức: SQL tách truy vấn dữ liệu (DML) khỏi định nghĩa cấu trúc (DDL) và quyền truy cập (DCL). Với SELECT, cần xác định bảng nguồn, điều kiện lọc, phép kết hợp, nhóm và điều kiện nhóm; các ràng buộc, view, thủ tục và trigger đảm nhiệm các vai trò khác nhau.

Phân tích: Đọc truy vấn theo các giai đoạn logic: FROM/JOIN tạo dữ liệu nguồn, WHERE lọc từng hàng, GROUP BY tạo nhóm, HAVING lọc nhóm, SELECT xác định biểu thức đầu ra và ORDER BY sắp kết quả. Với phép join, kiểm tra điều kiện ghép và các hàng không khớp trước khi tính kết quả.

Đối chiếu phương án: A. Select employeename, dept_code, salary From employee Where employee.salary in (select (salary)…; C. Select employeename, dept_code, salary From employee Where employee.salary in (select Max(sala…; D. Select employeename, dept_code, salary From employee Where employee.salary =(select Max(salary…. Các lựa chọn này không đáp ứng định nghĩa/điều kiện được hỏi; khi ôn, hãy chỉ ra chính xác điểm lệch của từng lựa chọn thay vì học thuộc mỗi chữ cái.
\end{question}

\begin{question}{379}
% [ID gốc: 5.94]
% [Nguồn: DBI202 Order 5 Q94]
Insert into Students (32, N'lona Bush', N'Female');In the given query which of the keyword has to be inserted ?
\begin{choices}
\choice Into
\choice Set
\choice Where
\choice Values
\end{choices}
\textbf{Đáp án:} D
\textbf{Giải thích:} Đáp án theo dữ liệu đề: D — D. Values.

Nền tảng kiến thức: SQL tách truy vấn dữ liệu (DML) khỏi định nghĩa cấu trúc (DDL) và quyền truy cập (DCL). Với SELECT, cần xác định bảng nguồn, điều kiện lọc, phép kết hợp, nhóm và điều kiện nhóm; các ràng buộc, view, thủ tục và trigger đảm nhiệm các vai trò khác nhau.

Phân tích: Đọc truy vấn theo các giai đoạn logic: FROM/JOIN tạo dữ liệu nguồn, WHERE lọc từng hàng, GROUP BY tạo nhóm, HAVING lọc nhóm, SELECT xác định biểu thức đầu ra và ORDER BY sắp kết quả. Với phép join, kiểm tra điều kiện ghép và các hàng không khớp trước khi tính kết quả.

Đối chiếu phương án: A. Into; B. Set; C. Where. Các lựa chọn này không đáp ứng định nghĩa/điều kiện được hỏi; khi ôn, hãy chỉ ra chính xác điểm lệch của từng lựa chọn thay vì học thuộc mỗi chữ cái.
\end{question}

\begin{question}{380}
% [ID gốc: 5.95]
% [Nguồn: DBI202 Order 5 Q95]
Foreign key constraints are created by using the " "keyword to refer to the primary key of another table.
\begin{choices}
\choice Refer
\choice References
\choice Referential
\choice All of the others
\end{choices}
\textbf{Đáp án:} B
\textbf{Giải thích:} Đáp án theo dữ liệu đề: B — B. References.

Nền tảng kiến thức: SQL tách truy vấn dữ liệu (DML) khỏi định nghĩa cấu trúc (DDL) và quyền truy cập (DCL). Với SELECT, cần xác định bảng nguồn, điều kiện lọc, phép kết hợp, nhóm và điều kiện nhóm; các ràng buộc, view, thủ tục và trigger đảm nhiệm các vai trò khác nhau.

Phân tích: Khóa ngoại buộc giá trị tham chiếu phải khớp một khóa được tham chiếu ở quan hệ đích (hoặc NULL nếu ràng buộc cho phép). Đây là toàn vẹn tham chiếu; nó khác tính duy nhất của primary/candidate key.

Đối chiếu phương án: A. Refer; C. Referential; D. All of the others. Các lựa chọn này không đáp ứng định nghĩa/điều kiện được hỏi; khi ôn, hãy chỉ ra chính xác điểm lệch của từng lựa chọn thay vì học thuộc mỗi chữ cái.
\end{question}

\begin{question}{381}
% [ID gốc: 5.107]
% [Nguồn: DBI202 Order 5 Q107]
Which of the following belongs to DDL (data definition language)?
\begin{choices}
\choice create, alter, drop
\choice Insert, update, delete, select
\choice Truncate, select, invoke
\choice Deny, create, drop
\end{choices}
\textbf{Đáp án:} A
\textbf{Giải thích:} Đáp án theo dữ liệu đề: A — A. create, alter, drop.

Nền tảng kiến thức: SQL tách truy vấn dữ liệu (DML) khỏi định nghĩa cấu trúc (DDL) và quyền truy cập (DCL). Với SELECT, cần xác định bảng nguồn, điều kiện lọc, phép kết hợp, nhóm và điều kiện nhóm; các ràng buộc, view, thủ tục và trigger đảm nhiệm các vai trò khác nhau.

Phân tích: DDL (Data Definition Language) khai báo/thay đổi cấu trúc như bảng và schema; DML truy vấn/thêm/sửa/xóa dữ liệu; DCL quản lý quyền. Chọn theo đối tượng bị tác động: cấu trúc, các hàng hay quyền truy cập.

Đối chiếu phương án: B. Insert, update, delete, select; C. Truncate, select, invoke; D. Deny, create, drop. Các lựa chọn này không đáp ứng định nghĩa/điều kiện được hỏi; khi ôn, hãy chỉ ra chính xác điểm lệch của từng lựa chọn thay vì học thuộc mỗi chữ cái.
\end{question}

\begin{question}{382}
% [ID gốc: 5.108]
% [Nguồn: DBI202 Order 5 Q108]
"SELECT TOP 100  FROM  " is the statement and suppose if there are 10 rows in the table then that what will be the result?
\begin{choices}
\choice Return 10 rows with data and 90 blank rows without data
\choice Return all the rows and all columns in the table
\choice Return the field of all the 10 rows
\choice Return error
\end{choices}
\textbf{Đáp án:} B
\textbf{Giải thích:} Đáp án theo dữ liệu đề: C — C. Return the field of all the 10 rows.

Nền tảng kiến thức: SQL tách truy vấn dữ liệu (DML) khỏi định nghĩa cấu trúc (DDL) và quyền truy cập (DCL). Với SELECT, cần xác định bảng nguồn, điều kiện lọc, phép kết hợp, nhóm và điều kiện nhóm; các ràng buộc, view, thủ tục và trigger đảm nhiệm các vai trò khác nhau.

Phân tích: Đọc truy vấn theo các giai đoạn logic: FROM/JOIN tạo dữ liệu nguồn, WHERE lọc từng hàng, GROUP BY tạo nhóm, HAVING lọc nhóm, SELECT xác định biểu thức đầu ra và ORDER BY sắp kết quả. Với phép join, kiểm tra điều kiện ghép và các hàng không khớp trước khi tính kết quả.

Đối chiếu phương án: A. Return 10 rows with data and 90 blank rows without data; B. Return all the rows and all columns in the table; D. Return error. Các lựa chọn này không đáp ứng định nghĩa/điều kiện được hỏi; khi ôn, hãy chỉ ra chính xác điểm lệch của từng lựa chọn thay vì học thuộc mỗi chữ cái.
\end{question}

\begin{question}{383}
% [ID gốc: 5.109]
% [Nguồn: DBI202 Order 5 Q109]
The predicate in a where clause can involve Boolean operations such as and. The result of true and unknown is false and unknown is while unknown and unknown is
\begin{choices}
\choice Unknown, unknown, false
\choice True, false, unknown
\choice True, unknown, unknown
\choice Unknown, false, unknown
\end{choices}
\textbf{Đáp án:} D
\textbf{Giải thích:} Đáp án theo dữ liệu đề: D — D. Unknown, false, unknown.

Nền tảng kiến thức: SQL tách truy vấn dữ liệu (DML) khỏi định nghĩa cấu trúc (DDL) và quyền truy cập (DCL). Với SELECT, cần xác định bảng nguồn, điều kiện lọc, phép kết hợp, nhóm và điều kiện nhóm; các ràng buộc, view, thủ tục và trigger đảm nhiệm các vai trò khác nhau.

Phân tích: Đọc truy vấn theo các giai đoạn logic: FROM/JOIN tạo dữ liệu nguồn, WHERE lọc từng hàng, GROUP BY tạo nhóm, HAVING lọc nhóm, SELECT xác định biểu thức đầu ra và ORDER BY sắp kết quả. Với phép join, kiểm tra điều kiện ghép và các hàng không khớp trước khi tính kết quả.

Đối chiếu phương án: A. Unknown, unknown, false; B. True, false, unknown; C. True, unknown, unknown. Các lựa chọn này không đáp ứng định nghĩa/điều kiện được hỏi; khi ôn, hãy chỉ ra chính xác điểm lệch của từng lựa chọn thay vì học thuộc mỗi chữ cái.
\end{question}

\begin{question}{384}
% [ID gốc: 5.111]
% [Nguồn: DBI202 Order 5 Q111]
UPDATE instructor
salary= salary * 1.05;
Fill in with correct keyword to update the instructor relation.
\begin{choices}
\choice Where
\choice Set
\choice In
\choice Select
\end{choices}
\textbf{Đáp án:} B
\textbf{Giải thích:} Đáp án theo dữ liệu đề: B — B. Set.

Nền tảng kiến thức: SQL tách truy vấn dữ liệu (DML) khỏi định nghĩa cấu trúc (DDL) và quyền truy cập (DCL). Với SELECT, cần xác định bảng nguồn, điều kiện lọc, phép kết hợp, nhóm và điều kiện nhóm; các ràng buộc, view, thủ tục và trigger đảm nhiệm các vai trò khác nhau.

Phân tích: Đọc truy vấn theo các giai đoạn logic: FROM/JOIN tạo dữ liệu nguồn, WHERE lọc từng hàng, GROUP BY tạo nhóm, HAVING lọc nhóm, SELECT xác định biểu thức đầu ra và ORDER BY sắp kết quả. Với phép join, kiểm tra điều kiện ghép và các hàng không khớp trước khi tính kết quả.

Đối chiếu phương án: A. Where; C. In; D. Select. Các lựa chọn này không đáp ứng định nghĩa/điều kiện được hỏi; khi ôn, hãy chỉ ra chính xác điểm lệch của từng lựa chọn thay vì học thuộc mỗi chữ cái.
\end{question}

\begin{question}{385}
% [ID gốc: 5.113]
% [Nguồn: DBI202 Order 5 Q113]
In an UPDATE statement, what does the WHERE clause determine?
\begin{choices}
\choice The order of update
\choice The columns to be updated
\choice The condition for selecting rows to be updated
\choice The table to be updated
\end{choices}
\textbf{Đáp án:} C
\textbf{Giải thích:} Đáp án theo dữ liệu đề: C — C. The condition for selecting rows to be updated.

Nền tảng kiến thức: SQL tách truy vấn dữ liệu (DML) khỏi định nghĩa cấu trúc (DDL) và quyền truy cập (DCL). Với SELECT, cần xác định bảng nguồn, điều kiện lọc, phép kết hợp, nhóm và điều kiện nhóm; các ràng buộc, view, thủ tục và trigger đảm nhiệm các vai trò khác nhau.

Phân tích: Đọc truy vấn theo các giai đoạn logic: FROM/JOIN tạo dữ liệu nguồn, WHERE lọc từng hàng, GROUP BY tạo nhóm, HAVING lọc nhóm, SELECT xác định biểu thức đầu ra và ORDER BY sắp kết quả. Với phép join, kiểm tra điều kiện ghép và các hàng không khớp trước khi tính kết quả.

Đối chiếu phương án: A. The order of update; B. The columns to be updated; D. The table to be updated. Các lựa chọn này không đáp ứng định nghĩa/điều kiện được hỏi; khi ôn, hãy chỉ ra chính xác điểm lệch của từng lựa chọn thay vì học thuộc mỗi chữ cái.
\end{question}

\begin{question}{386}
% [ID gốc: 5.121]
% [Nguồn: DBI202 Order 5 Q121]
What data type is best suited for Vietnamese name? (in SQL Server)
\begin{choices}
\choice Varchar(30)
\choice Nvarchar(30)
\choice Text(30)
\choice Char(30)
\end{choices}
\textbf{Đáp án:} B
\textbf{Giải thích:} Đáp án theo dữ liệu đề: B — B. Nvarchar(30).

Nền tảng kiến thức: SQL tách truy vấn dữ liệu (DML) khỏi định nghĩa cấu trúc (DDL) và quyền truy cập (DCL). Với SELECT, cần xác định bảng nguồn, điều kiện lọc, phép kết hợp, nhóm và điều kiện nhóm; các ràng buộc, view, thủ tục và trigger đảm nhiệm các vai trò khác nhau.

Phân tích: NVARCHAR lưu văn bản Unicode nên phù hợp với tiếng Việt có dấu; VARCHAR không Unicode có thể phụ thuộc code page. Kích thước trong ngoặc là giới hạn số đơn vị ký tự theo SQL Server, cần chọn dựa trên độ dài dữ liệu thực tế.

Đối chiếu phương án: A. Varchar(30); C. Text(30); D. Char(30). Các lựa chọn này không đáp ứng định nghĩa/điều kiện được hỏi; khi ôn, hãy chỉ ra chính xác điểm lệch của từng lựa chọn thay vì học thuộc mỗi chữ cái.
\end{question}

\begin{question}{387}
% [ID gốc: 5.122]
% [Nguồn: DBI202 Order 5 Q122]
Assuming UserProfile is a table containing a column Proession which accepts NULL value. What is the result of the below query. SET ANSI_NULLS OFF SELECT Profession FROM UserProfile WHERE (Profession <>NULL)
\begin{choices}
\choice Gives an error.
\choice <> is not an operator in SQL Server
\choice Returns all Professions which has not null values
\choice None of the above.
\end{choices}
\textbf{Đáp án:} C
\textbf{Giải thích:} Đáp án theo dữ liệu đề: C — C. Returns all Professions which has not null values.

Nền tảng kiến thức: SQL tách truy vấn dữ liệu (DML) khỏi định nghĩa cấu trúc (DDL) và quyền truy cập (DCL). Với SELECT, cần xác định bảng nguồn, điều kiện lọc, phép kết hợp, nhóm và điều kiện nhóm; các ràng buộc, view, thủ tục và trigger đảm nhiệm các vai trò khác nhau.

Phân tích: NULL biểu thị thiếu/không biết/không áp dụng, không phải số 0 hay chuỗi rỗng. So sánh với NULL bằng `=` không cho TRUE/FALSE mà cho UNKNOWN; dùng `IS NULL`/`IS NOT NULL`. Trong logic ba giá trị, AND/OR/NOT phải được tính theo cả trạng thái UNKNOWN.

Đối chiếu phương án: A. Gives an error.; B. <> is not an operator in SQL Server; D. None of the above.. Các lựa chọn này không đáp ứng định nghĩa/điều kiện được hỏi; khi ôn, hãy chỉ ra chính xác điểm lệch của từng lựa chọn thay vì học thuộc mỗi chữ cái.
\end{question}

\begin{question}{388}
% [ID gốc: 5.123]
% [Nguồn: DBI202 Order 5 Q123]
In MS-SQL Server, which of the following is true about join operator?
SELECT *
FROM tbICUSTOMER, tbIORDER
WHERE tbICUSTOMER.CUSTOMER_ID = tblORDER. CUSTOMER_ID
\begin{choices}
\choice Full outer-join
\choice Natural join
\choice Outer join
\choice Cartesian join
\end{choices}
\textbf{Đáp án:} D
\textbf{Giải thích:} Đáp án theo dữ liệu đề: B — B. Natural join.

Nền tảng kiến thức: SQL tách truy vấn dữ liệu (DML) khỏi định nghĩa cấu trúc (DDL) và quyền truy cập (DCL). Với SELECT, cần xác định bảng nguồn, điều kiện lọc, phép kết hợp, nhóm và điều kiện nhóm; các ràng buộc, view, thủ tục và trigger đảm nhiệm các vai trò khác nhau.

Phân tích: Join ghép các bộ thỏa điều kiện. Inner join loại bộ không ghép được; LEFT/RIGHT OUTER JOIN giữ mọi bộ từ phía được chỉ định và điền NULL phía thiếu. Natural join tự ghép các thuộc tính trùng tên và chỉ xuất một cột cho mỗi thuộc tính ghép.

Đối chiếu phương án: A. Full outer-join; C. Outer join; D. Cartesian join. Các lựa chọn này không đáp ứng định nghĩa/điều kiện được hỏi; khi ôn, hãy chỉ ra chính xác điểm lệch của từng lựa chọn thay vì học thuộc mỗi chữ cái.
\end{question}

\begin{question}{389}
% [ID gốc: 5.124]
% [Nguồn: DBI202 Order 5 Q124]
In SQL query, choose the correct statement about the order of execution.
\begin{choices}
\choice SELECT-FROM-WHERE - GROUP BY - HAVING - ORDER BY
\choice SELECT - FROM - GROUP BY - WHERE - HAVING - ORDER BY
\choice SELECT-FROM-HAVING GROUP BY - WHERE-ORDER BY
\choice SELECT-FROM-OREDER BY WHERE HAVING - GROUP BY
\end{choices}
\textbf{Đáp án:} A
\textbf{Giải thích:} Đáp án theo dữ liệu đề: A — A. SELECT-FROM-WHERE - GROUP BY - HAVING - ORDER BY.

Nền tảng kiến thức: SQL tách truy vấn dữ liệu (DML) khỏi định nghĩa cấu trúc (DDL) và quyền truy cập (DCL). Với SELECT, cần xác định bảng nguồn, điều kiện lọc, phép kết hợp, nhóm và điều kiện nhóm; các ràng buộc, view, thủ tục và trigger đảm nhiệm các vai trò khác nhau.

Phân tích: Đọc truy vấn theo các giai đoạn logic: FROM/JOIN tạo dữ liệu nguồn, WHERE lọc từng hàng, GROUP BY tạo nhóm, HAVING lọc nhóm, SELECT xác định biểu thức đầu ra và ORDER BY sắp kết quả. Với phép join, kiểm tra điều kiện ghép và các hàng không khớp trước khi tính kết quả.

Đối chiếu phương án: B. SELECT - FROM - GROUP BY - WHERE - HAVING - ORDER BY; C. SELECT-FROM-HAVING GROUP BY - WHERE-ORDER BY; D. SELECT-FROM-OREDER BY WHERE HAVING - GROUP BY. Các lựa chọn này không đáp ứng định nghĩa/điều kiện được hỏi; khi ôn, hãy chỉ ra chính xác điểm lệch của từng lựa chọn thay vì học thuộc mỗi chữ cái.
\end{question}

\begin{question}{390}
% [ID gốc: 5.125]
% [Nguồn: DBI202 Order 5 Q125]
In MS-SQL Server, there are three primary SQL commands (INSERT, UPDATE, and DELETE) to modify data
in a table. Which of the following is false about function of these commands?
\begin{choices}
\choice These commands allow you to change the data type of a column, add or drop columns, and rename columns, among other operations.
\choice These commands enable the alteration of data within database tables.
\choice Each command has its specific syntax, usage, and allowing you to insert new rows, modifies existing rows, and remove rows from the table.
\choice These commands belong to the Data Manipulation Language (DML) subset of SQL.
\end{choices}
\textbf{Đáp án:} A
\textbf{Giải thích:} Đáp án theo dữ liệu đề: A — A. These commands allow you to change the data type of a column, add or drop columns, and rename columns, among other operations..

Nền tảng kiến thức: SQL tách truy vấn dữ liệu (DML) khỏi định nghĩa cấu trúc (DDL) và quyền truy cập (DCL). Với SELECT, cần xác định bảng nguồn, điều kiện lọc, phép kết hợp, nhóm và điều kiện nhóm; các ràng buộc, view, thủ tục và trigger đảm nhiệm các vai trò khác nhau.

Phân tích: DDL (Data Definition Language) khai báo/thay đổi cấu trúc như bảng và schema; DML truy vấn/thêm/sửa/xóa dữ liệu; DCL quản lý quyền. Chọn theo đối tượng bị tác động: cấu trúc, các hàng hay quyền truy cập.

Đối chiếu phương án: B. These commands enable the alteration of data within database tables.; C. Each command has its specific syntax, usage, and allowing you to insert new rows, modifies exi…; D. These commands belong to the Data Manipulation Language (DML) subset of SQL.. Các lựa chọn này không đáp ứng định nghĩa/điều kiện được hỏi; khi ôn, hãy chỉ ra chính xác điểm lệch của từng lựa chọn thay vì học thuộc mỗi chữ cái.
\end{question}

\begin{question}{391}
% [ID gốc: 5.131]
% [Nguồn: DBI202 Order 5 Q131]
What data type is best suited for a Vietnamese name in SQL Server?
\begin{choices}
\choice Varchar(30)
\choice Nvarchar(30)
\choice Text(30)
\choice Char(30)
\end{choices}
\textbf{Đáp án:} B
\textbf{Giải thích:} Đáp án theo dữ liệu đề: B — B. Nvarchar(30).

Lời giải gốc: Kiểu dữ liệu Nvarchar hỗ trợ mã Unicode, cho phép lưu trữ chính xác các ký tự tiếng Việt có dấu mà không bị lỗi font, phù hợp nhất để lưu tên tiếng Việt trong SQL Server.

Nền tảng kiến thức: SQL tách truy vấn dữ liệu (DML) khỏi định nghĩa cấu trúc (DDL) và quyền truy cập (DCL). Với SELECT, cần xác định bảng nguồn, điều kiện lọc, phép kết hợp, nhóm và điều kiện nhóm; các ràng buộc, view, thủ tục và trigger đảm nhiệm các vai trò khác nhau.

Phân tích: NVARCHAR lưu văn bản Unicode nên phù hợp với tiếng Việt có dấu; VARCHAR không Unicode có thể phụ thuộc code page. Kích thước trong ngoặc là giới hạn số đơn vị ký tự theo SQL Server, cần chọn dựa trên độ dài dữ liệu thực tế.

Đối chiếu phương án: A. Varchar(30); C. Text(30); D. Char(30). Các lựa chọn này không đáp ứng định nghĩa/điều kiện được hỏi; khi ôn, hãy chỉ ra chính xác điểm lệch của từng lựa chọn thay vì học thuộc mỗi chữ cái.
\end{question}

\begin{question}{392}
% [ID gốc: 5.132]
% [Nguồn: DBI202 Order 5 Q132]
(Choose 3 answers) Which one of the following is NOT a DDL command? Choose correct answers.
\begin{choices}
\choice DELETE
\choice CREATE DATABASE
\choice INSERT
\choice USE
\end{choices}
\textbf{Đáp án:} ACD
\textbf{Giải thích:} Đáp án theo dữ liệu đề: A, C, D — A. DELETE; C. INSERT; D. USE.

Lời giải gốc: CREATE DATABASE là lệnh DDL duy nhất trong danh sách (định nghĩa cấu trúc). DELETE và INSERT thuộc DML vì thao tác trên dữ liệu, còn USE chỉ là lệnh tiện ích chọn database đang làm việc cho phiên hiện tại, không định nghĩa hay thay đổi cấu trúc nào.

Nền tảng kiến thức: SQL tách truy vấn dữ liệu (DML) khỏi định nghĩa cấu trúc (DDL) và quyền truy cập (DCL). Với SELECT, cần xác định bảng nguồn, điều kiện lọc, phép kết hợp, nhóm và điều kiện nhóm; các ràng buộc, view, thủ tục và trigger đảm nhiệm các vai trò khác nhau.

Phân tích: DDL (Data Definition Language) khai báo/thay đổi cấu trúc như bảng và schema; DML truy vấn/thêm/sửa/xóa dữ liệu; DCL quản lý quyền. Chọn theo đối tượng bị tác động: cấu trúc, các hàng hay quyền truy cập.

Đối chiếu phương án: B. CREATE DATABASE. Các lựa chọn này không đáp ứng định nghĩa/điều kiện được hỏi; khi ôn, hãy chỉ ra chính xác điểm lệch của từng lựa chọn thay vì học thuộc mỗi chữ cái.
\end{question}

\begin{question}{393}
% [ID gốc: 5.136]
% [Nguồn: DBI202 Order 5 Q136]
In SQL, to update a relation schema, which one of the following statements can be used?
\begin{choices}
\choice Insert
\choice Update
\choice Alter
\choice Alter table
\end{choices}
\textbf{Đáp án:} D
\textbf{Giải thích:} Đáp án theo dữ liệu đề: D — D. Alter table.

Lời giải gốc: Lệnh ALTER TABLE được sử dụng trong SQL để cập nhật cấu trúc (schema) của một quan hệ hoặc bảng đã tồn tại trong cơ sở dữ liệu. ---

Nền tảng kiến thức: SQL tách truy vấn dữ liệu (DML) khỏi định nghĩa cấu trúc (DDL) và quyền truy cập (DCL). Với SELECT, cần xác định bảng nguồn, điều kiện lọc, phép kết hợp, nhóm và điều kiện nhóm; các ràng buộc, view, thủ tục và trigger đảm nhiệm các vai trò khác nhau.

Phân tích: Lược đồ quan hệ ghi tên quan hệ kèm danh sách thuộc tính, thường theo dạng `Boat(BoatID, BoatName, Color, Status)`. Nếu đề hỏi schema thuần túy, không nhất thiết phải kèm kiểu dữ liệu hay các bộ dữ liệu đang lưu.

Đối chiếu phương án: A. Insert; B. Update; C. Alter. Các lựa chọn này không đáp ứng định nghĩa/điều kiện được hỏi; khi ôn, hãy chỉ ra chính xác điểm lệch của từng lựa chọn thay vì học thuộc mỗi chữ cái.
\end{question}

\begin{question}{394}
% [ID gốc: 5.140]
% [Nguồn: DBI202 Order 5 Q140]
A DELETE command operates on how many relations?
\begin{choices}
\choice One
\choice Two
\choice Several
\choice Null
\end{choices}
\textbf{Đáp án:} A
\textbf{Giải thích:} Đáp án theo dữ liệu đề: A — A. One.

Lời giải gốc: Câu lệnh DELETE trong SQL tiêu chuẩn chỉ thực hiện thao tác xóa dữ liệu trên một bảng (relation) tại mỗi câu lệnh.

Nền tảng kiến thức: SQL tách truy vấn dữ liệu (DML) khỏi định nghĩa cấu trúc (DDL) và quyền truy cập (DCL). Với SELECT, cần xác định bảng nguồn, điều kiện lọc, phép kết hợp, nhóm và điều kiện nhóm; các ràng buộc, view, thủ tục và trigger đảm nhiệm các vai trò khác nhau.

Phân tích: Một câu lệnh DELETE tác động lên một target table/relation trong mỗi câu lệnh, dù có thể xóa nhiều hàng của bảng đó theo WHERE. Các hàng tham chiếu ở bảng khác chỉ bị ảnh hưởng nếu có cascade/trigger được cấu hình.

Đối chiếu phương án: B. Two; C. Several; D. Null. Các lựa chọn này không đáp ứng định nghĩa/điều kiện được hỏi; khi ôn, hãy chỉ ra chính xác điểm lệch của từng lựa chọn thay vì học thuộc mỗi chữ cái.
\end{question}

\begin{question}{395}
% [ID gốc: 5.146]
% [Nguồn: DBI202 Order 5 Q146]
Suppose relation R(A, B, C, D) has tuples. Which tuples are contained in the result of the query below?<pre>SELECT C, AVG(D) FROM R GROUP BY C HAVING AVG(D) > 1000 ORDER BY C;</pre>
\begin{choices}
\choice (101, 11500); (102, 2500)
\choice (101, 2300); (102, 1250)
\choice (101, 5); (102, 2)
\choice (101, 11500); (102, 1250)
\end{choices}
\textbf{Đáp án:} D
\textbf{Giải thích:} Filtering grouped tuples with AVG(D) > 1000 yields D.
\end{question}

\begin{question}{396}
% [ID gốc: 5.147]
% [Nguồn: DBI202 Order 5 Q147]
To completely remove a table and its data in SQL, we use:
\begin{choices}
\choice DELETE
\choice DROP
\choice REMOVE
\choice RASE
\end{choices}
\textbf{Đáp án:} B
\textbf{Giải thích:} Đáp án theo dữ liệu đề: B — B. DROP.

Lời giải gốc: Lệnh DROP được dùng để xóa hoàn toàn bảng cùng với toàn bộ dữ liệu và cấu trúc của nó khỏi cơ sở dữ liệu. ---

Nền tảng kiến thức: SQL tách truy vấn dữ liệu (DML) khỏi định nghĩa cấu trúc (DDL) và quyền truy cập (DCL). Với SELECT, cần xác định bảng nguồn, điều kiện lọc, phép kết hợp, nhóm và điều kiện nhóm; các ràng buộc, view, thủ tục và trigger đảm nhiệm các vai trò khác nhau.

Phân tích: DROP TABLE loại bỏ định nghĩa bảng cùng dữ liệu phụ thuộc theo quy tắc DBMS; DELETE xóa các hàng nhưng giữ bảng, còn TRUNCATE xóa nhanh toàn bộ hàng trong khi vẫn giữ cấu trúc. Cần chọn lệnh đúng cấp độ tác động.

Đối chiếu phương án: A. DELETE; C. REMOVE; D. RASE. Các lựa chọn này không đáp ứng định nghĩa/điều kiện được hỏi; khi ôn, hãy chỉ ra chính xác điểm lệch của từng lựa chọn thay vì học thuộc mỗi chữ cái.
\end{question}

\begin{question}{397}
% [ID gốc: 5.151]
% [Nguồn: DBI202 Order 5 Q151]
Which SQL clause is used to add a new column to an existing table?
\begin{choices}
\choice ADD COLUMN
\choice ALTER TABLE ADD COLUMN
\choice INSERT INTO
\choice CREATE COLUMN
\end{choices}
\textbf{Đáp án:} B
\textbf{Giải thích:} Đáp án theo dữ liệu đề: B — B. ALTER TABLE ADD COLUMN.

Lời giải gốc: Lệnh ALTER TABLE table_name ADD column_name datatype được dùng để thêm một cột mới vào bảng hiện có.

Nền tảng kiến thức: SQL tách truy vấn dữ liệu (DML) khỏi định nghĩa cấu trúc (DDL) và quyền truy cập (DCL). Với SELECT, cần xác định bảng nguồn, điều kiện lọc, phép kết hợp, nhóm và điều kiện nhóm; các ràng buộc, view, thủ tục và trigger đảm nhiệm các vai trò khác nhau.

Phân tích: ALTER TABLE thay đổi định nghĩa bảng đã tồn tại, như thêm cột hoặc sửa cấu hình cột/ràng buộc theo cú pháp của DBMS. Nó khác UPDATE, vốn sửa các giá trị trong những hàng hiện hữu; cú pháp chính xác có thể khác giữa SQL Server và chuẩn SQL.

Đối chiếu phương án: A. ADD COLUMN; C. INSERT INTO; D. CREATE COLUMN. Các lựa chọn này không đáp ứng định nghĩa/điều kiện được hỏi; khi ôn, hãy chỉ ra chính xác điểm lệch của từng lựa chọn thay vì học thuộc mỗi chữ cái.
\end{question}

\begin{question}{398}
% [ID gốc: 5.152]
% [Nguồn: DBI202 Order 5 Q152]
Which command is used to create a stored procedure in SQL Server?
\begin{choices}
\choice CREATE FUNCTION
\choice CREATE TABLE
\choice CREATE PROCEDURE
\choice CREATE VIEW
\end{choices}
\textbf{Đáp án:} C
\textbf{Giải thích:} Đáp án theo dữ liệu đề: C — C. CREATE PROCEDURE.

Lời giải gốc: Lệnh CREATE PROCEDURE được sử dụng để tạo một thủ tục lưu trữ (stored procedure) trong SQL Server. ---

Nền tảng kiến thức: SQL tách truy vấn dữ liệu (DML) khỏi định nghĩa cấu trúc (DDL) và quyền truy cập (DCL). Với SELECT, cần xác định bảng nguồn, điều kiện lọc, phép kết hợp, nhóm và điều kiện nhóm; các ràng buộc, view, thủ tục và trigger đảm nhiệm các vai trò khác nhau.

Phân tích: Stored procedure là tập lệnh SQL có tên lưu trong CSDL, được gọi để thực thi một thủ tục; nó có thể nhận tham số và chứa nhiều câu lệnh. Đừng nhầm với trigger tự chạy khi có sự kiện hoặc function có quy tắc trả về giá trị.

Đối chiếu phương án: A. CREATE FUNCTION; B. CREATE TABLE; D. CREATE VIEW. Các lựa chọn này không đáp ứng định nghĩa/điều kiện được hỏi; khi ôn, hãy chỉ ra chính xác điểm lệch của từng lựa chọn thay vì học thuộc mỗi chữ cái.
\end{question}

\begin{question}{399}
% [ID gốc: 5.153]
% [Nguồn: DBI202 Order 5 Q153]
Fill in the correct keyword to update the instructor relation.<pre>UPDATE instructor ____ salary = salary * 1.05;</pre>A. WHERE
\begin{choices}
\choice SET
\choice FROM
\choice SELECT
\end{choices}
\textbf{Đáp án:} A
\textbf{Giải thích:} Đáp án theo dữ liệu đề: B — B. SET.

Lời giải gốc: Trong cú pháp câu lệnh UPDATE của SQL, từ khóa SET được dùng để chỉ định cột cần cập nhật giá trị mới.

Nền tảng kiến thức: SQL tách truy vấn dữ liệu (DML) khỏi định nghĩa cấu trúc (DDL) và quyền truy cập (DCL). Với SELECT, cần xác định bảng nguồn, điều kiện lọc, phép kết hợp, nhóm và điều kiện nhóm; các ràng buộc, view, thủ tục và trigger đảm nhiệm các vai trò khác nhau.

Phân tích: Đọc truy vấn theo các giai đoạn logic: FROM/JOIN tạo dữ liệu nguồn, WHERE lọc từng hàng, GROUP BY tạo nhóm, HAVING lọc nhóm, SELECT xác định biểu thức đầu ra và ORDER BY sắp kết quả. Với phép join, kiểm tra điều kiện ghép và các hàng không khớp trước khi tính kết quả.

Đối chiếu phương án: C. FROM; D. SELECT. Các lựa chọn này không đáp ứng định nghĩa/điều kiện được hỏi; khi ôn, hãy chỉ ra chính xác điểm lệch của từng lựa chọn thay vì học thuộc mỗi chữ cái.
\end{question}

\begin{question}{400}
% [ID gốc: 5.155]
% [Nguồn: DBI202 Order 5 Q155]
Which SQL command is used to define a table schema?
\begin{choices}
\choice ALTER TABLE
\choice SELECT
\choice CREATE TABLE
\choice INSERT
\end{choices}
\textbf{Đáp án:} C
\textbf{Giải thích:} Đáp án theo dữ liệu đề: C — C. CREATE TABLE.

Lời giải gốc: Lệnh CREATE TABLE thuộc nhóm ngôn ngữ định nghĩa dữ liệu (DDL), được dùng để tạo mới và định nghĩa lược đồ cấu trúc cấu thành của một bảng.

Nền tảng kiến thức: SQL tách truy vấn dữ liệu (DML) khỏi định nghĩa cấu trúc (DDL) và quyền truy cập (DCL). Với SELECT, cần xác định bảng nguồn, điều kiện lọc, phép kết hợp, nhóm và điều kiện nhóm; các ràng buộc, view, thủ tục và trigger đảm nhiệm các vai trò khác nhau.

Phân tích: ALTER TABLE thay đổi định nghĩa bảng đã tồn tại, như thêm cột hoặc sửa cấu hình cột/ràng buộc theo cú pháp của DBMS. Nó khác UPDATE, vốn sửa các giá trị trong những hàng hiện hữu; cú pháp chính xác có thể khác giữa SQL Server và chuẩn SQL.

Đối chiếu phương án: A. ALTER TABLE; B. SELECT; D. INSERT. Các lựa chọn này không đáp ứng định nghĩa/điều kiện được hỏi; khi ôn, hãy chỉ ra chính xác điểm lệch của từng lựa chọn thay vì học thuộc mỗi chữ cái.
\end{question}

\begin{question}{401}
% [ID gốc: 5.156]
% [Nguồn: DBI202 Order 5 Q156]
When designing a database for a library system, which constraint ensures that every book has exactly one unique ISBN?
\begin{choices}
\choice Foreign key constraint
\choice NOT NULL constraint
\choice Unique constraint
\choice Default constraint
\end{choices}
\textbf{Đáp án:} C
\textbf{Giải thích:} Đáp án theo dữ liệu đề: C — C. Unique constraint.

Lời giải gốc: Ràng buộc UNIQUE đảm bảo rằng các giá trị trong cột ISBN là duy nhất cho mỗi cuốn sách. ---

Nền tảng kiến thức: SQL tách truy vấn dữ liệu (DML) khỏi định nghĩa cấu trúc (DDL) và quyền truy cập (DCL). Với SELECT, cần xác định bảng nguồn, điều kiện lọc, phép kết hợp, nhóm và điều kiện nhóm; các ràng buộc, view, thủ tục và trigger đảm nhiệm các vai trò khác nhau.

Phân tích: NULL biểu thị thiếu/không biết/không áp dụng, không phải số 0 hay chuỗi rỗng. So sánh với NULL bằng `=` không cho TRUE/FALSE mà cho UNKNOWN; dùng `IS NULL`/`IS NOT NULL`. Trong logic ba giá trị, AND/OR/NOT phải được tính theo cả trạng thái UNKNOWN.

Đối chiếu phương án: A. Foreign key constraint; B. NOT NULL constraint; D. Default constraint. Các lựa chọn này không đáp ứng định nghĩa/điều kiện được hỏi; khi ôn, hãy chỉ ra chính xác điểm lệch của từng lựa chọn thay vì học thuộc mỗi chữ cái.
\end{question}

\begin{question}{402}
% [ID gốc: 5.157]
% [Nguồn: DBI202 Order 5 Q157]
What is the difference between DELETE and TRUNCATE in SQL?
\begin{choices}
\choice DELETE removes rows based on a condition, and TRUNCATE removes all rows from the table
\choice DELETE is faster and does not generate logs
\choice TRUNCATE allows rollback, but DELETE does not
\choice TRUNCATE can delete specific rows, while DELETE removes all rows
\end{choices}
\textbf{Đáp án:} A
\textbf{Giải thích:} Đáp án theo dữ liệu đề: A — A. DELETE removes rows based on a condition, and TRUNCATE removes all rows from the table.

Lời giải gốc: Lệnh DELETE cho phép xóa các hàng có điều kiện (dùng WHERE) và ghi log từng dòng, trong khi TRUNCATE xóa toàn bộ dữ liệu nhanh hơn mà không cần điều kiện.

Nền tảng kiến thức: SQL tách truy vấn dữ liệu (DML) khỏi định nghĩa cấu trúc (DDL) và quyền truy cập (DCL). Với SELECT, cần xác định bảng nguồn, điều kiện lọc, phép kết hợp, nhóm và điều kiện nhóm; các ràng buộc, view, thủ tục và trigger đảm nhiệm các vai trò khác nhau.

Phân tích: Phép tập hợp trên quan hệ cần so sánh các tuple theo thuộc tính tương ứng; union lấy hợp, intersection lấy phần chung và difference lấy phần có ở vế trái mà không có ở vế phải. Trước tiên kiểm tra tính tương thích schema, sau đó áp phép theo đúng hướng/toán tử.

Đối chiếu phương án: B. DELETE is faster and does not generate logs; C. TRUNCATE allows rollback, but DELETE does not; D. TRUNCATE can delete specific rows, while DELETE removes all rows. Các lựa chọn này không đáp ứng định nghĩa/điều kiện được hỏi; khi ôn, hãy chỉ ra chính xác điểm lệch của từng lựa chọn thay vì học thuộc mỗi chữ cái.
\end{question}

\begin{question}{403}
% [ID gốc: 5.161]
% [Nguồn: DBI202 Order 5 Q161]
Given two relations: PC(model, speed, ram, hd, type) and Laptop(model, speed, ram, hd, screen, price). Which subquery gets the model of laptops with the same price as any PC?
\begin{choices}
\choice SELECT model FROM Laptop WHERE price IN (SELECT price FROM PC);
\choice SELECT model FROM Laptop WHERE price = ALL (SELECT price FROM PC);
\choice SELECT model FROM Laptop WHERE price &gt; ALL (SELECT price FROM PC);
\choice SELECT model FROM Laptop WHERE price &lt; ALL (SELECT price FROM PC);
\end{choices}
\textbf{Đáp án:} A
\textbf{Giải thích:} Đáp án theo dữ liệu đề: A — A. SELECT model FROM Laptop WHERE price IN (SELECT price FROM PC);.

Lời giải gốc: Để tìm các laptop có giá bằng với bất kỳ PC nào, ta sử dụng toán tử IN kết hợp với tập hợp giá của các PC. ---

Nền tảng kiến thức: SQL tách truy vấn dữ liệu (DML) khỏi định nghĩa cấu trúc (DDL) và quyền truy cập (DCL). Với SELECT, cần xác định bảng nguồn, điều kiện lọc, phép kết hợp, nhóm và điều kiện nhóm; các ràng buộc, view, thủ tục và trigger đảm nhiệm các vai trò khác nhau.

Phân tích: Đọc truy vấn theo các giai đoạn logic: FROM/JOIN tạo dữ liệu nguồn, WHERE lọc từng hàng, GROUP BY tạo nhóm, HAVING lọc nhóm, SELECT xác định biểu thức đầu ra và ORDER BY sắp kết quả. Với phép join, kiểm tra điều kiện ghép và các hàng không khớp trước khi tính kết quả.

Đối chiếu phương án: B. SELECT model FROM Laptop WHERE price = ALL (SELECT price FROM PC);; C. SELECT model FROM Laptop WHERE price &gt; ALL (SELECT price FROM PC);; D. SELECT model FROM Laptop WHERE price &lt; ALL (SELECT price FROM PC);. Các lựa chọn này không đáp ứng định nghĩa/điều kiện được hỏi; khi ôn, hãy chỉ ra chính xác điểm lệch của từng lựa chọn thay vì học thuộc mỗi chữ cái.
\end{question}

\begin{question}{404}
% [ID gốc: 5.166]
% [Nguồn: DBI202 Order 5 Q166]
Which function gets the day value (of the month) in SQL?
\begin{choices}
\choice DAYOFMONTH(date)
\choice DATE(date)
\choice DAY(date)
\choice MONTHDAY(date)
\choice None of the above answers
\end{choices}
\textbf{Đáp án:} C
\textbf{Giải thích:} Đáp án theo dữ liệu đề: C — C. DAY(date).

Lời giải gốc: Hàm DAY(date) trong SQL Server được dùng để trích xuất phần ngày trong tháng từ một giá trị ngày tháng.

Nền tảng kiến thức: SQL tách truy vấn dữ liệu (DML) khỏi định nghĩa cấu trúc (DDL) và quyền truy cập (DCL). Với SELECT, cần xác định bảng nguồn, điều kiện lọc, phép kết hợp, nhóm và điều kiện nhóm; các ràng buộc, view, thủ tục và trigger đảm nhiệm các vai trò khác nhau.

Phân tích: Với lựa chọn tổng hợp, không chọn theo cảm tính: kiểm tra từng phát biểu độc lập. “All of the others/above” chỉ đúng khi mọi phát biểu còn lại đúng; “None” chỉ đúng khi không phát biểu nào trong nhóm thỏa yêu cầu. Đáp án lưu cần được đọc cùng các mệnh đề thành phần.

Đối chiếu phương án: A. DAYOFMONTH(date); B. DATE(date); D. MONTHDAY(date); E. None of the above answers. Các lựa chọn này không đáp ứng định nghĩa/điều kiện được hỏi; khi ôn, hãy chỉ ra chính xác điểm lệch của từng lựa chọn thay vì học thuộc mỗi chữ cái.
\end{question}

\begin{question}{405}
% [ID gốc: 5.167]
% [Nguồn: DBI202 Order 5 Q167]
Assuming UserProfile is a table containing a column Profession which accepts NULL value. What is the result of the below query? SET ANSI_NULLS OFF SELECT Profession FROM UserProfile WHERE (Profession &lt;&gt; NULL)
\begin{choices}
\choice Gives an error.
\choice &lt;&gt; is not an operator in SQL Server
\choice Returns all Professions which has not null values.
\choice None of the above.
\end{choices}
\textbf{Đáp án:} C
\textbf{Giải thích:} Đáp án theo dữ liệu đề: C — C. Returns all Professions which has not null values..

Lời giải gốc: Khi thiết lập SET ANSI_NULLS OFF, toán tử <> có thể đánh giá giá trị NULL theo cú pháp cũ, trả về tất cả các giá trị nghề nghiệp khác NULL (không chứa giá trị rỗng). ---

Nền tảng kiến thức: SQL tách truy vấn dữ liệu (DML) khỏi định nghĩa cấu trúc (DDL) và quyền truy cập (DCL). Với SELECT, cần xác định bảng nguồn, điều kiện lọc, phép kết hợp, nhóm và điều kiện nhóm; các ràng buộc, view, thủ tục và trigger đảm nhiệm các vai trò khác nhau.

Phân tích: NULL biểu thị thiếu/không biết/không áp dụng, không phải số 0 hay chuỗi rỗng. So sánh với NULL bằng `=` không cho TRUE/FALSE mà cho UNKNOWN; dùng `IS NULL`/`IS NOT NULL`. Trong logic ba giá trị, AND/OR/NOT phải được tính theo cả trạng thái UNKNOWN.

Đối chiếu phương án: A. Gives an error.; B. &lt;&gt; is not an operator in SQL Server; D. None of the above.. Các lựa chọn này không đáp ứng định nghĩa/điều kiện được hỏi; khi ôn, hãy chỉ ra chính xác điểm lệch của từng lựa chọn thay vì học thuộc mỗi chữ cái.
\end{question}

\begin{question}{406}
% [ID gốc: 5.172]
% [Nguồn: DBI202 Order 5 Q172]
In SQL, to update a relation's schema, which statement can be used?
\begin{choices}
\choice INSERT
\choice UPDATE
\choice ALTER
\choice ALTER TABLE
\end{choices}
\textbf{Đáp án:} D
\textbf{Giải thích:} Đáp án theo dữ liệu đề: D — D. ALTER TABLE.

Lời giải gốc: Lệnh ALTER TABLE thuộc nhóm DDL, dùng để chỉnh sửa cấu trúc (schema) của một quan hệ/bảng.

Nền tảng kiến thức: SQL tách truy vấn dữ liệu (DML) khỏi định nghĩa cấu trúc (DDL) và quyền truy cập (DCL). Với SELECT, cần xác định bảng nguồn, điều kiện lọc, phép kết hợp, nhóm và điều kiện nhóm; các ràng buộc, view, thủ tục và trigger đảm nhiệm các vai trò khác nhau.

Phân tích: ALTER TABLE thay đổi định nghĩa bảng đã tồn tại, như thêm cột hoặc sửa cấu hình cột/ràng buộc theo cú pháp của DBMS. Nó khác UPDATE, vốn sửa các giá trị trong những hàng hiện hữu; cú pháp chính xác có thể khác giữa SQL Server và chuẩn SQL.

Đối chiếu phương án: A. INSERT; B. UPDATE; C. ALTER. Các lựa chọn này không đáp ứng định nghĩa/điều kiện được hỏi; khi ôn, hãy chỉ ra chính xác điểm lệch của từng lựa chọn thay vì học thuộc mỗi chữ cái.
\end{question}

\begin{question}{407}
% [ID gốc: 5.174]
% [Nguồn: DBI202 Order 5 Q174]
Which of the following is NOT the OPTIONAL clause in a SELECT statement?
\begin{choices}
\choice WHERE
\choice ORDER BY
\choice FROM
\choice HAVING
\end{choices}
\textbf{Đáp án:} C
\textbf{Giải thích:} Đáp án theo dữ liệu đề: C — C. FROM.

Lời giải gốc: Trong câu lệnh SELECT, mệnh đề FROM là bắt buộc phải có, trong khi WHERE, ORDER BY, và HAVING là các mệnh đề tùy chọn. ---

Nền tảng kiến thức: SQL tách truy vấn dữ liệu (DML) khỏi định nghĩa cấu trúc (DDL) và quyền truy cập (DCL). Với SELECT, cần xác định bảng nguồn, điều kiện lọc, phép kết hợp, nhóm và điều kiện nhóm; các ràng buộc, view, thủ tục và trigger đảm nhiệm các vai trò khác nhau.

Phân tích: Đọc truy vấn theo các giai đoạn logic: FROM/JOIN tạo dữ liệu nguồn, WHERE lọc từng hàng, GROUP BY tạo nhóm, HAVING lọc nhóm, SELECT xác định biểu thức đầu ra và ORDER BY sắp kết quả. Với phép join, kiểm tra điều kiện ghép và các hàng không khớp trước khi tính kết quả.

Đối chiếu phương án: A. WHERE; B. ORDER BY; D. HAVING. Các lựa chọn này không đáp ứng định nghĩa/điều kiện được hỏi; khi ôn, hãy chỉ ra chính xác điểm lệch của từng lựa chọn thay vì học thuộc mỗi chữ cái.
\end{question}

\begin{question}{408}
% [ID gốc: 5.175]
% [Nguồn: DBI202 Order 5 Q175]
Which command is used to remove a relation from SQL?
\begin{choices}
\choice Drop table
\choice Delete
\choice Purge
\choice Remove
\end{choices}
\textbf{Đáp án:} A
\textbf{Giải thích:} Đáp án theo dữ liệu đề: A — A. Drop table.

Lời giải gốc: Lệnh DROP TABLE được sử dụng trong SQL để xóa hoàn toàn cấu trúc bảng (relation) cùng toàn bộ dữ liệu của nó khỏi cơ sở dữ liệu.

Nền tảng kiến thức: SQL tách truy vấn dữ liệu (DML) khỏi định nghĩa cấu trúc (DDL) và quyền truy cập (DCL). Với SELECT, cần xác định bảng nguồn, điều kiện lọc, phép kết hợp, nhóm và điều kiện nhóm; các ràng buộc, view, thủ tục và trigger đảm nhiệm các vai trò khác nhau.

Phân tích: DROP TABLE loại bỏ định nghĩa bảng cùng dữ liệu phụ thuộc theo quy tắc DBMS; DELETE xóa các hàng nhưng giữ bảng, còn TRUNCATE xóa nhanh toàn bộ hàng trong khi vẫn giữ cấu trúc. Cần chọn lệnh đúng cấp độ tác động.

Đối chiếu phương án: B. Delete; C. Purge; D. Remove. Các lựa chọn này không đáp ứng định nghĩa/điều kiện được hỏi; khi ôn, hãy chỉ ra chính xác điểm lệch của từng lựa chọn thay vì học thuộc mỗi chữ cái.
\end{question}

\begin{question}{409}
% [ID gốc: 5.179]
% [Nguồn: DBI202 Order 5 Q179]
Given two relations: Printer(maker, model, type) PC(model, speed, ram,hd,, price) What does this query return? SELECT maker FROM Product WHERE model NOT IN ( SELECT model FROM PC );
\begin{choices}
\choice Makers of all devices
\choice Makers that also build PCs
\choice Makers of non-PC devices
\choice Makers of PCs only
\end{choices}
\textbf{Đáp án:} C
\textbf{Giải thích:} Đáp án theo dữ liệu đề: C — C. Makers of non-PC devices.

Lời giải gốc: Truy vấn con lấy danh sách tất cả các model của PC. Mệnh đề WHERE model NOT IN (...) lọc ra các sản phẩm không thuộc PC, do đó kết quả trả về nhà sản xuất của các thiết bị không phải PC.

Nền tảng kiến thức: SQL tách truy vấn dữ liệu (DML) khỏi định nghĩa cấu trúc (DDL) và quyền truy cập (DCL). Với SELECT, cần xác định bảng nguồn, điều kiện lọc, phép kết hợp, nhóm và điều kiện nhóm; các ràng buộc, view, thủ tục và trigger đảm nhiệm các vai trò khác nhau.

Phân tích: ER diagram mô tả các entity set, thuộc tính và relationship set, cùng các ràng buộc cần thiết. Nó tập trung vào đối tượng dữ liệu và liên kết giữa chúng, không mô tả luồng xử lý/chương trình như sơ đồ quy trình.

Đối chiếu phương án: A. Makers of all devices; B. Makers that also build PCs; D. Makers of PCs only. Các lựa chọn này không đáp ứng định nghĩa/điều kiện được hỏi; khi ôn, hãy chỉ ra chính xác điểm lệch của từng lựa chọn thay vì học thuộc mỗi chữ cái.
\end{question}

\begin{question}{410}
% [ID gốc: 5.180]
% [Nguồn: DBI202 Order 5 Q180]
What is the main difference between DELETE and TRUNCATE?
\begin{choices}
\choice DELETE cannot be undone (ROLLBACK), while TRUNCATE can.
\choice TRUNCATE deletes all data in a table faster and logs minimally, while DELETE can delete data based on the WHERE condition
\choice DELETE is for tables, while TRUNCATE is for Views.
\choice There is no difference.
\end{choices}
\textbf{Đáp án:} B
\textbf{Giải thích:} Đáp án theo dữ liệu đề: B — B. TRUNCATE deletes all data in a table faster and logs minimally, while DELETE can delete data based on the WHERE condition.

Lời giải gốc: TRUNCATE xóa toàn bộ dữ liệu trong bảng nhanh hơn do ghi log tối thiểu (minimal logging) và không áp dụng điều kiện WHERE, trong khi DELETE ghi log chi tiết từng dòng và có thể xóa theo điều kiện WHERE.

Nền tảng kiến thức: SQL tách truy vấn dữ liệu (DML) khỏi định nghĩa cấu trúc (DDL) và quyền truy cập (DCL). Với SELECT, cần xác định bảng nguồn, điều kiện lọc, phép kết hợp, nhóm và điều kiện nhóm; các ràng buộc, view, thủ tục và trigger đảm nhiệm các vai trò khác nhau.

Phân tích: Phép tập hợp trên quan hệ cần so sánh các tuple theo thuộc tính tương ứng; union lấy hợp, intersection lấy phần chung và difference lấy phần có ở vế trái mà không có ở vế phải. Trước tiên kiểm tra tính tương thích schema, sau đó áp phép theo đúng hướng/toán tử.

Đối chiếu phương án: A. DELETE cannot be undone (ROLLBACK), while TRUNCATE can.; C. DELETE is for tables, while TRUNCATE is for Views.; D. There is no difference.. Các lựa chọn này không đáp ứng định nghĩa/điều kiện được hỏi; khi ôn, hãy chỉ ra chính xác điểm lệch của từng lựa chọn thay vì học thuộc mỗi chữ cái.
\end{question}

\begin{question}{411}
% [ID gốc: 5.183]
% [Nguồn: DBI202 Order 5 Q183]
In SQL, after grouping data using the GROUP BY clause, which clause is used to filter the groups based on specific conditions?
\begin{choices}
\choice WHERE
\choice HAVING
\choice GROUP BY
\choice WITH
\end{choices}
\textbf{Đáp án:} B
\textbf{Giải thích:} Đáp án theo dữ liệu đề: B — B. HAVING.

Lời giải gốc: Mệnh đề HAVING được dùng để lọc các nhóm sau khi đã gom nhóm bằng GROUP BY, hỗ trợ cả các hàm tổng hợp.

Nền tảng kiến thức: SQL tách truy vấn dữ liệu (DML) khỏi định nghĩa cấu trúc (DDL) và quyền truy cập (DCL). Với SELECT, cần xác định bảng nguồn, điều kiện lọc, phép kết hợp, nhóm và điều kiện nhóm; các ràng buộc, view, thủ tục và trigger đảm nhiệm các vai trò khác nhau.

Phân tích: Đọc truy vấn theo các giai đoạn logic: FROM/JOIN tạo dữ liệu nguồn, WHERE lọc từng hàng, GROUP BY tạo nhóm, HAVING lọc nhóm, SELECT xác định biểu thức đầu ra và ORDER BY sắp kết quả. Với phép join, kiểm tra điều kiện ghép và các hàng không khớp trước khi tính kết quả.

Đối chiếu phương án: A. WHERE; C. GROUP BY; D. WITH. Các lựa chọn này không đáp ứng định nghĩa/điều kiện được hỏi; khi ôn, hãy chỉ ra chính xác điểm lệch của từng lựa chọn thay vì học thuộc mỗi chữ cái.
\end{question}

\begin{question}{412}
% [ID gốc: 5.184]
% [Nguồn: DBI202 Order 5 Q184]
What is a NOT NULL constraint?
\begin{choices}
\choice A constraint that allows only NULL values.
\choice A constraint that prevents NULL values in a column
\choice A constraint that sets default values to NULL.
\choice A constraint that converts NULL values to a specific data type.
\end{choices}
\textbf{Đáp án:} B
\textbf{Giải thích:} Đáp án theo dữ liệu đề: B — B. A constraint that prevents NULL values in a column.

Lời giải gốc: Ràng buộc NOT NULL đảm bảo cột không được phép chứa giá trị trống/rỗng (NULL).

Nền tảng kiến thức: SQL tách truy vấn dữ liệu (DML) khỏi định nghĩa cấu trúc (DDL) và quyền truy cập (DCL). Với SELECT, cần xác định bảng nguồn, điều kiện lọc, phép kết hợp, nhóm và điều kiện nhóm; các ràng buộc, view, thủ tục và trigger đảm nhiệm các vai trò khác nhau.

Phân tích: NULL biểu thị thiếu/không biết/không áp dụng, không phải số 0 hay chuỗi rỗng. So sánh với NULL bằng `=` không cho TRUE/FALSE mà cho UNKNOWN; dùng `IS NULL`/`IS NOT NULL`. Trong logic ba giá trị, AND/OR/NOT phải được tính theo cả trạng thái UNKNOWN.

Đối chiếu phương án: A. A constraint that allows only NULL values.; C. A constraint that sets default values to NULL.; D. A constraint that converts NULL values to a specific data type.. Các lựa chọn này không đáp ứng định nghĩa/điều kiện được hỏi; khi ôn, hãy chỉ ra chính xác điểm lệch của từng lựa chọn thay vì học thuộc mỗi chữ cái.
\end{question}

\begin{question}{413}
% [ID gốc: 5.185]
% [Nguồn: DBI202 Order 5 Q185]
The \% wildcard, when used with the LIKE keyword, represents a string of characters that:
\begin{choices}
\choice Has a minimum length of 0 characters
\choice Has a minimum length of 1 character, a maximum length of n characters
\choice Has a minimum length of 1 character, a maximum length of 255 characters
\choice Has a minimum length of 1 character, a maximum length of 256 characters
\choice None of the above answers
\end{choices}
\textbf{Đáp án:} A
\textbf{Giải thích:} Đáp án theo dữ liệu đề: A — A. Has a minimum length of 0 characters.

Lời giải gốc: Ký tự đại diện \% trong mệnh đề LIKE đại diện cho một chuỗi gồm 0 hoặc nhiều ký tự bất kỳ.

Nền tảng kiến thức: SQL tách truy vấn dữ liệu (DML) khỏi định nghĩa cấu trúc (DDL) và quyền truy cập (DCL). Với SELECT, cần xác định bảng nguồn, điều kiện lọc, phép kết hợp, nhóm và điều kiện nhóm; các ràng buộc, view, thủ tục và trigger đảm nhiệm các vai trò khác nhau.

Phân tích: Trong LIKE của SQL Server, `%` khớp 0 hoặc nhiều ký tự và `_` khớp đúng một ký tự. ESCAPE cho phép coi ký tự wildcard như ký tự thường bằng cách đặt trước chúng một ký tự thoát; nó không thay thế mẫu bằng chuỗi bất kỳ.

Đối chiếu phương án: B. Has a minimum length of 1 character, a maximum length of n characters; C. Has a minimum length of 1 character, a maximum length of 255 characters; D. Has a minimum length of 1 character, a maximum length of 256 characters; E. None of the above answers. Các lựa chọn này không đáp ứng định nghĩa/điều kiện được hỏi; khi ôn, hãy chỉ ra chính xác điểm lệch của từng lựa chọn thay vì học thuộc mỗi chữ cái.
\end{question}

\begin{question}{414}
% [ID gốc: 5.187]
% [Nguồn: DBI202 Order 5 Q187]
Consider two relations MovieStar(name, address, gender, birthdate) and MovieExec(name, address, cert, netWorth). You want the names and addresses of all female movie stars who are also movie executives with a net worth over $10,000,000. What is the correct statement?
\begin{choices}
\choice SELECT name, address FROM MovieStar WHERE gender = 'F' UNION SELECT name, address FROM MovieExec WHERE netWorth &gt; 10000000
\choice SELECT name, address FROM MovieStar WHERE gender = 'F' INTERSECT SELECT name, address FROM MovieExec WHERE netWorth &gt; 10000000
\choice SELECT name, address FROM MovieStar WHERE gender = 'F' EXCEPT SELECT name, address FROM MovieExec WHERE netWorth &gt; 10000000
\choice None of the others
\end{choices}
\textbf{Đáp án:} B
\textbf{Giải thích:} Đáp án theo dữ liệu đề: B — B. SELECT name, address FROM MovieStar WHERE gender = 'F' INTERSECT SELECT name, address FROM MovieExec WHERE netWorth &gt; 10000000.

Lời giải gốc: Để tìm những người vừa là nữ diễn viên vừa là nhà điều hành phim thỏa mãn điều kiện, ta dùng phép giao INTERSECT giữa hai tập kết quả.

Nền tảng kiến thức: SQL tách truy vấn dữ liệu (DML) khỏi định nghĩa cấu trúc (DDL) và quyền truy cập (DCL). Với SELECT, cần xác định bảng nguồn, điều kiện lọc, phép kết hợp, nhóm và điều kiện nhóm; các ràng buộc, view, thủ tục và trigger đảm nhiệm các vai trò khác nhau.

Phân tích: Phép tập hợp trên quan hệ cần so sánh các tuple theo thuộc tính tương ứng; union lấy hợp, intersection lấy phần chung và difference lấy phần có ở vế trái mà không có ở vế phải. Trước tiên kiểm tra tính tương thích schema, sau đó áp phép theo đúng hướng/toán tử.

Đối chiếu phương án: A. SELECT name, address FROM MovieStar WHERE gender = 'F' UNION SELECT name, address FROM MovieEx…; C. SELECT name, address FROM MovieStar WHERE gender = 'F' EXCEPT SELECT name, address FROM MovieE…; D. None of the others. Các lựa chọn này không đáp ứng định nghĩa/điều kiện được hỏi; khi ôn, hãy chỉ ra chính xác điểm lệch của từng lựa chọn thay vì học thuộc mỗi chữ cái.
\end{question}

\begin{question}{415}
% [ID gốc: 5.193]
% [Nguồn: DBI202 Order 5 Q193]
Given the table Orders(order_id, customer_id, amount), which SQL query returns the customer with the lowest total order amount?
\begin{choices}
\choice SELECT TOP 1 customer_id FROM Orders GROUP BY customer_id ORDER BY SUM(amount) DESC;
\choice SELECT TOP 1 customer_id FROM Orders GROUP BY customer_id ORDER BY SUM(amount);
\choice SELECT customer_id FROM Orders WHERE amount = MIN(amount);
\choice SELECT customer_id, MAX(amount) FROM Orders;
\end{choices}
\textbf{Đáp án:} B
\textbf{Giải thích:} Đáp án theo dữ liệu đề: B — B. SELECT TOP 1 customer_id FROM Orders GROUP BY customer_id ORDER BY SUM(amount);.

Lời giải gốc: Câu lệnh nhóm theo khách hàng, tính tổng tiền và sắp xếp tăng dần (ORDER BY SUM(amount)), sau đó lấy dòng đầu tiên bằng TOP 1 sẽ trả về khách hàng có tổng giá trị đơn hàng thấp nhất. ---

Nền tảng kiến thức: SQL tách truy vấn dữ liệu (DML) khỏi định nghĩa cấu trúc (DDL) và quyền truy cập (DCL). Với SELECT, cần xác định bảng nguồn, điều kiện lọc, phép kết hợp, nhóm và điều kiện nhóm; các ràng buộc, view, thủ tục và trigger đảm nhiệm các vai trò khác nhau.

Phân tích: Đọc truy vấn theo các giai đoạn logic: FROM/JOIN tạo dữ liệu nguồn, WHERE lọc từng hàng, GROUP BY tạo nhóm, HAVING lọc nhóm, SELECT xác định biểu thức đầu ra và ORDER BY sắp kết quả. Với phép join, kiểm tra điều kiện ghép và các hàng không khớp trước khi tính kết quả.

Đối chiếu phương án: A. SELECT TOP 1 customer_id FROM Orders GROUP BY customer_id ORDER BY SUM(amount) DESC;; C. SELECT customer_id FROM Orders WHERE amount = MIN(amount);; D. SELECT customer_id, MAX(amount) FROM Orders;. Các lựa chọn này không đáp ứng định nghĩa/điều kiện được hỏi; khi ôn, hãy chỉ ra chính xác điểm lệch của từng lựa chọn thay vì học thuộc mỗi chữ cái.
\end{question}

\begin{question}{416}
% [ID gốc: 5.195]
% [Nguồn: DBI202 Order 5 Q195]
Which one of the following is NOT a DDL command?
\begin{choices}
\choice DELETE
\choice CREATE DATABASE
\choice INSERT
\choice USE
\end{choices}
\textbf{Đáp án:} AC
\textbf{Giải thích:} Đáp án theo dữ liệu đề: A — A. DELETE.

Lời giải gốc: Lệnh DELETE thuộc nhóm ngôn ngữ thao tác dữ liệu (DML), trong khi CREATE thuộc nhóm ngôn ngữ định nghĩa dữ liệu (DDL). Lưu ý: INSERT cũng là DML, nên nếu đề yêu cầu chọn nhiều đáp án thì cả A và C đều không phải DDL.

Nền tảng kiến thức: SQL tách truy vấn dữ liệu (DML) khỏi định nghĩa cấu trúc (DDL) và quyền truy cập (DCL). Với SELECT, cần xác định bảng nguồn, điều kiện lọc, phép kết hợp, nhóm và điều kiện nhóm; các ràng buộc, view, thủ tục và trigger đảm nhiệm các vai trò khác nhau.

Phân tích: DDL (Data Definition Language) khai báo/thay đổi cấu trúc như bảng và schema; DML truy vấn/thêm/sửa/xóa dữ liệu; DCL quản lý quyền. Chọn theo đối tượng bị tác động: cấu trúc, các hàng hay quyền truy cập.

Đối chiếu phương án: B. CREATE DATABASE; C. INSERT; D. USE. Các lựa chọn này không đáp ứng định nghĩa/điều kiện được hỏi; khi ôn, hãy chỉ ra chính xác điểm lệch của từng lựa chọn thay vì học thuộc mỗi chữ cái.
\end{question}

\begin{question}{417}
% [ID gốc: 5.198]
% [Nguồn: DBI202 Order 5 Q198]
Regarding the difference between WHERE and HAVING in SQL, which of the following is true?
\begin{choices}
\choice WHERE filters data on aggregated values, HAVING filters on non-aggregated ones.
\choice WHERE is written after GROUP BY, while HAVING comes before GROUP BY.
\choice WHERE filters data on non-aggregated values, HAVING filters on aggregated ones.
\choice Both HAVING and WHERE are written before GROUP BY, and HAVING comes before WHERE.
\end{choices}
\textbf{Đáp án:} C
\textbf{Giải thích:} Đáp án theo dữ liệu đề: C — C. WHERE filters data on non-aggregated values, HAVING filters on aggregated ones..

Lời giải gốc: Mệnh đề WHERE lọc dữ liệu trên các giá trị thông thường trước khi nhóm, còn HAVING lọc dữ liệu trên các giá trị đã qua hàm tổng hợp (aggregate functions) sau khi nhóm.

Nền tảng kiến thức: SQL tách truy vấn dữ liệu (DML) khỏi định nghĩa cấu trúc (DDL) và quyền truy cập (DCL). Với SELECT, cần xác định bảng nguồn, điều kiện lọc, phép kết hợp, nhóm và điều kiện nhóm; các ràng buộc, view, thủ tục và trigger đảm nhiệm các vai trò khác nhau.

Phân tích: Phép tập hợp trên quan hệ cần so sánh các tuple theo thuộc tính tương ứng; union lấy hợp, intersection lấy phần chung và difference lấy phần có ở vế trái mà không có ở vế phải. Trước tiên kiểm tra tính tương thích schema, sau đó áp phép theo đúng hướng/toán tử.

Đối chiếu phương án: A. WHERE filters data on aggregated values, HAVING filters on non-aggregated ones.; B. WHERE is written after GROUP BY, while HAVING comes before GROUP BY.; D. Both HAVING and WHERE are written before GROUP BY, and HAVING comes before WHERE.. Các lựa chọn này không đáp ứng định nghĩa/điều kiện được hỏi; khi ôn, hãy chỉ ra chính xác điểm lệch của từng lựa chọn thay vì học thuộc mỗi chữ cái.
\end{question}

\begin{question}{418}
% [ID gốc: 5.199]
% [Nguồn: DBI202 Order 5 Q199]
With SQL, how do you select all records from a table named "Customers" where the value of the column "FirstName" ends with an "a"?
\begin{choices}
\choice SELECT * FROM Customers WHERE FirstName = 'a'
\choice SELECT * FROM Customers WHERE FirstName LIKE 'a%'
\choice SELECT * FROM Customers WHERE FirstName LIKE '\%a'
\choice SELECT * FROM Customers WHERE FirstName = '\%a%'
\end{choices}
\textbf{Đáp án:} C
\textbf{Giải thích:} Đáp án theo dữ liệu đề: C — C. SELECT * FROM Customers WHERE FirstName LIKE '\%a'.

Lời giải gốc: Toán tử LIKE '\%a' dùng để tìm các chuỗi ký tự kết thúc bằng chữ 'a'. Ngược lại, LIKE 'a%' tìm chuỗi bắt đầu bằng 'a'.

Nền tảng kiến thức: SQL tách truy vấn dữ liệu (DML) khỏi định nghĩa cấu trúc (DDL) và quyền truy cập (DCL). Với SELECT, cần xác định bảng nguồn, điều kiện lọc, phép kết hợp, nhóm và điều kiện nhóm; các ràng buộc, view, thủ tục và trigger đảm nhiệm các vai trò khác nhau.

Phân tích: Đọc truy vấn theo các giai đoạn logic: FROM/JOIN tạo dữ liệu nguồn, WHERE lọc từng hàng, GROUP BY tạo nhóm, HAVING lọc nhóm, SELECT xác định biểu thức đầu ra và ORDER BY sắp kết quả. Với phép join, kiểm tra điều kiện ghép và các hàng không khớp trước khi tính kết quả.

Đối chiếu phương án: A. SELECT * FROM Customers WHERE FirstName = 'a'; B. SELECT * FROM Customers WHERE FirstName LIKE 'a%'; D. SELECT * FROM Customers WHERE FirstName = '\%a%'. Các lựa chọn này không đáp ứng định nghĩa/điều kiện được hỏi; khi ôn, hãy chỉ ra chính xác điểm lệch của từng lựa chọn thay vì học thuộc mỗi chữ cái.
\end{question}

\begin{question}{419}
% [ID gốc: 5.200]
% [Nguồn: DBI202 Order 5 Q200]
What does the ON UPDATE CASCADE option do in SQL Server when defined for a foreign key constraint?
\begin{choices}
\choice Automatically updates the foreign key value when the referenced primary key is updated
\choice Automatically deletes the foreign key row when the referenced row is updated
\choice Prevents the update of foreign key values
\choice Sets the foreign key value to NULL when the referenced primary key is updated
\end{choices}
\textbf{Đáp án:} A
\textbf{Giải thích:} Đáp án theo dữ liệu đề: A — A. Automatically updates the foreign key value when the referenced primary key is updated.

Lời giải gốc: Tùy chọn ON UPDATE CASCADE tự động cập nhật giá trị khóa ngoại ở bảng con tương ứng khi giá trị khóa chính ở bảng cha được chỉnh sửa.

Nền tảng kiến thức: SQL tách truy vấn dữ liệu (DML) khỏi định nghĩa cấu trúc (DDL) và quyền truy cập (DCL). Với SELECT, cần xác định bảng nguồn, điều kiện lọc, phép kết hợp, nhóm và điều kiện nhóm; các ràng buộc, view, thủ tục và trigger đảm nhiệm các vai trò khác nhau.

Phân tích: Khóa ngoại buộc giá trị tham chiếu phải khớp một khóa được tham chiếu ở quan hệ đích (hoặc NULL nếu ràng buộc cho phép). Đây là toàn vẹn tham chiếu; nó khác tính duy nhất của primary/candidate key.

Đối chiếu phương án: B. Automatically deletes the foreign key row when the referenced row is updated; C. Prevents the update of foreign key values; D. Sets the foreign key value to NULL when the referenced primary key is updated. Các lựa chọn này không đáp ứng định nghĩa/điều kiện được hỏi; khi ôn, hãy chỉ ra chính xác điểm lệch của từng lựa chọn thay vì học thuộc mỗi chữ cái.
\end{question}

\begin{question}{420}
% [ID gốc: 5.204]
% [Nguồn: DBI202 Order 5 Q204]
In SQL, after grouping data using the GROUP BY clause, the ______ clause is used to filter the groups based on specific conditions.
\begin{choices}
\choice Where
\choice Having
\choice Group by
\choice With
\end{choices}
\textbf{Đáp án:} B
\textbf{Giải thích:} Đáp án theo dữ liệu đề: B — B. Having.

Lời giải gốc: Điều kiện lọc dùng cho các hàm tổng hợp (aggregate functions) sau khi đã gom nhóm dữ liệu bằng GROUP BY phải được đặt ở mệnh đề HAVING. Mệnh đề WHERE chỉ lọc các hàng dữ liệu thô trước khi gom nhóm.

Nền tảng kiến thức: SQL tách truy vấn dữ liệu (DML) khỏi định nghĩa cấu trúc (DDL) và quyền truy cập (DCL). Với SELECT, cần xác định bảng nguồn, điều kiện lọc, phép kết hợp, nhóm và điều kiện nhóm; các ràng buộc, view, thủ tục và trigger đảm nhiệm các vai trò khác nhau.

Phân tích: Đọc truy vấn theo các giai đoạn logic: FROM/JOIN tạo dữ liệu nguồn, WHERE lọc từng hàng, GROUP BY tạo nhóm, HAVING lọc nhóm, SELECT xác định biểu thức đầu ra và ORDER BY sắp kết quả. Với phép join, kiểm tra điều kiện ghép và các hàng không khớp trước khi tính kết quả.

Đối chiếu phương án: A. Where; C. Group by; D. With. Các lựa chọn này không đáp ứng định nghĩa/điều kiện được hỏi; khi ôn, hãy chỉ ra chính xác điểm lệch của từng lựa chọn thay vì học thuộc mỗi chữ cái.
\end{question}

\begin{question}{421}
% [ID gốc: 5.206]
% [Nguồn: DBI202 Order 5 Q206]
In SQL query, choose the correct statement about the logical order of execution.
\begin{choices}
\choice SELECT - FROM - WHERE - GROUP BY - HAVING - ORDER BY
\choice SELECT - FROM - GROUP BY - WHERE - HAVING - ORDER BY
\choice SELECT - FROM - HAVING - GROUP BY - WHERE - ORDER BY
\choice SELECT - FROM - ORDER BY - WHERE - HAVING - GROUP BY
\end{choices}
\textbf{Đáp án:} A
\textbf{Giải thích:} Đáp án theo dữ liệu đề: A — A. SELECT - FROM - WHERE - GROUP BY - HAVING - ORDER BY.

Lời giải gốc: Thứ tự xử lý logic của các mệnh đề SQL là: FROM -> WHERE -> GROUP BY -> HAVING -> SELECT -> ORDER BY. Phương án A trình bày đúng trình tự các mệnh đề viết theo cú pháp tiêu chuẩn.

Nền tảng kiến thức: SQL tách truy vấn dữ liệu (DML) khỏi định nghĩa cấu trúc (DDL) và quyền truy cập (DCL). Với SELECT, cần xác định bảng nguồn, điều kiện lọc, phép kết hợp, nhóm và điều kiện nhóm; các ràng buộc, view, thủ tục và trigger đảm nhiệm các vai trò khác nhau.

Phân tích: Đọc truy vấn theo các giai đoạn logic: FROM/JOIN tạo dữ liệu nguồn, WHERE lọc từng hàng, GROUP BY tạo nhóm, HAVING lọc nhóm, SELECT xác định biểu thức đầu ra và ORDER BY sắp kết quả. Với phép join, kiểm tra điều kiện ghép và các hàng không khớp trước khi tính kết quả.

Đối chiếu phương án: B. SELECT - FROM - GROUP BY - WHERE - HAVING - ORDER BY; C. SELECT - FROM - HAVING - GROUP BY - WHERE - ORDER BY; D. SELECT - FROM - ORDER BY - WHERE - HAVING - GROUP BY. Các lựa chọn này không đáp ứng định nghĩa/điều kiện được hỏi; khi ôn, hãy chỉ ra chính xác điểm lệch của từng lựa chọn thay vì học thuộc mỗi chữ cái.
\end{question}

\begin{question}{422}
% [ID gốc: 5.207]
% [Nguồn: DBI202 Order 5 Q207]
In a relational database, which SQL command is used to modify the structure of an existing table?
\begin{choices}
\choice UPDATE
\choice ALTER
\choice MODIFY
\choice CHANGE
\end{choices}
\textbf{Đáp án:} B
\textbf{Giải thích:} Đáp án theo dữ liệu đề: B — B. ALTER.

Lời giải gốc: Để chỉnh sửa cấu trúc của một bảng đã tồn tại (như thêm/xóa cột, đổi kiểu dữ liệu), người ta sử dụng câu lệnh ALTER TABLE.

Nền tảng kiến thức: SQL tách truy vấn dữ liệu (DML) khỏi định nghĩa cấu trúc (DDL) và quyền truy cập (DCL). Với SELECT, cần xác định bảng nguồn, điều kiện lọc, phép kết hợp, nhóm và điều kiện nhóm; các ràng buộc, view, thủ tục và trigger đảm nhiệm các vai trò khác nhau.

Phân tích: ALTER TABLE thay đổi định nghĩa bảng đã tồn tại, như thêm cột hoặc sửa cấu hình cột/ràng buộc theo cú pháp của DBMS. Nó khác UPDATE, vốn sửa các giá trị trong những hàng hiện hữu; cú pháp chính xác có thể khác giữa SQL Server và chuẩn SQL.

Đối chiếu phương án: A. UPDATE; C. MODIFY; D. CHANGE. Các lựa chọn này không đáp ứng định nghĩa/điều kiện được hỏi; khi ôn, hãy chỉ ra chính xác điểm lệch của từng lựa chọn thay vì học thuộc mỗi chữ cái.
\end{question}

\begin{question}{423}
% [ID gốc: 5.211]
% [Nguồn: DBI202 Order 5 Q211]
Which of the following SQL statements returns all rows from Cars where year is greater than 2019?
\begin{choices}
\choice SELECT * FROM Cars WHERE year &gt; 2019;
\choice SELECT * FROM Cars HAVING year &gt; 2019;
\choice SELECT * FROM Cars WHEN year &gt; 2019;
\choice SELECT * FROM Cars WHERE year =&gt; 2019;
\end{choices}
\textbf{Đáp án:} A
\textbf{Giải thích:} Đáp án theo dữ liệu đề: A — A. SELECT * FROM Cars WHERE year &gt; 2019;.

Lời giải gốc: Cú pháp chuẩn để lọc các dòng dựa trên điều kiện là sử dụng mệnh đề WHERE với toán tử so sánh >.

Nền tảng kiến thức: SQL tách truy vấn dữ liệu (DML) khỏi định nghĩa cấu trúc (DDL) và quyền truy cập (DCL). Với SELECT, cần xác định bảng nguồn, điều kiện lọc, phép kết hợp, nhóm và điều kiện nhóm; các ràng buộc, view, thủ tục và trigger đảm nhiệm các vai trò khác nhau.

Phân tích: Đọc truy vấn theo các giai đoạn logic: FROM/JOIN tạo dữ liệu nguồn, WHERE lọc từng hàng, GROUP BY tạo nhóm, HAVING lọc nhóm, SELECT xác định biểu thức đầu ra và ORDER BY sắp kết quả. Với phép join, kiểm tra điều kiện ghép và các hàng không khớp trước khi tính kết quả.

Đối chiếu phương án: B. SELECT * FROM Cars HAVING year &gt; 2019;; C. SELECT * FROM Cars WHEN year &gt; 2019;; D. SELECT * FROM Cars WHERE year =&gt; 2019;. Các lựa chọn này không đáp ứng định nghĩa/điều kiện được hỏi; khi ôn, hãy chỉ ra chính xác điểm lệch của từng lựa chọn thay vì học thuộc mỗi chữ cái.
\end{question}

\begin{question}{424}
% [ID gốc: 5.212]
% [Nguồn: DBI202 Order 5 Q212]
What is the value of the query: SELECT 'Agnes' + NULL + 'Allison'
\begin{choices}
\choice Agnes
\choice Allison
\choice NULL
\choice Agnes Allison
\end{choices}
\textbf{Đáp án:} C
\textbf{Giải thích:} Đáp án theo dữ liệu đề: C — C. NULL.

Lời giải gốc: Trong SQL Server, khi thực hiện phép nối chuỗi (+) với bất kỳ giá trị NULL nào, kết quả trả về sẽ luôn là NULL.

Nền tảng kiến thức: SQL tách truy vấn dữ liệu (DML) khỏi định nghĩa cấu trúc (DDL) và quyền truy cập (DCL). Với SELECT, cần xác định bảng nguồn, điều kiện lọc, phép kết hợp, nhóm và điều kiện nhóm; các ràng buộc, view, thủ tục và trigger đảm nhiệm các vai trò khác nhau.

Phân tích: NULL biểu thị thiếu/không biết/không áp dụng, không phải số 0 hay chuỗi rỗng. So sánh với NULL bằng `=` không cho TRUE/FALSE mà cho UNKNOWN; dùng `IS NULL`/`IS NOT NULL`. Trong logic ba giá trị, AND/OR/NOT phải được tính theo cả trạng thái UNKNOWN.

Đối chiếu phương án: A. Agnes; B. Allison; D. Agnes Allison. Các lựa chọn này không đáp ứng định nghĩa/điều kiện được hỏi; khi ôn, hãy chỉ ra chính xác điểm lệch của từng lựa chọn thay vì học thuộc mỗi chữ cái.
\end{question}

\begin{question}{425}
% [ID gốc: 5.214]
% [Nguồn: DBI202 Order 5 Q214]
The ORDER BY clause can only be used in cases ______
\begin{choices}
\choice GROUP BY queries
\choice SELECT queries
\choice HAVING queries
\choice None of the others
\end{choices}
\textbf{Đáp án:} B
\textbf{Giải thích:} Đáp án theo dữ liệu đề: B — B. SELECT queries.

Lời giải gốc: ORDER BY dùng để sắp xếp thứ tự hiển thị dữ liệu đầu ra và nó là một phần bổ trợ nằm trong cấu trúc của câu lệnh truy vấn dữ liệu SELECT.

Nền tảng kiến thức: SQL tách truy vấn dữ liệu (DML) khỏi định nghĩa cấu trúc (DDL) và quyền truy cập (DCL). Với SELECT, cần xác định bảng nguồn, điều kiện lọc, phép kết hợp, nhóm và điều kiện nhóm; các ràng buộc, view, thủ tục và trigger đảm nhiệm các vai trò khác nhau.

Phân tích: Đọc truy vấn theo các giai đoạn logic: FROM/JOIN tạo dữ liệu nguồn, WHERE lọc từng hàng, GROUP BY tạo nhóm, HAVING lọc nhóm, SELECT xác định biểu thức đầu ra và ORDER BY sắp kết quả. Với phép join, kiểm tra điều kiện ghép và các hàng không khớp trước khi tính kết quả.

Đối chiếu phương án: A. GROUP BY queries; C. HAVING queries; D. None of the others. Các lựa chọn này không đáp ứng định nghĩa/điều kiện được hỏi; khi ôn, hãy chỉ ra chính xác điểm lệch của từng lựa chọn thay vì học thuộc mỗi chữ cái.
\end{question}

\begin{question}{426}
% [ID gốc: 5.217]
% [Nguồn: DBI202 Order 5 Q217]
In an employee table, to ensure an attribute always has a value, which constraint must be used?
\begin{choices}
\choice NULL
\choice NOT NULL
\choice UNIQUE
\choice DISTINCT
\end{choices}
\textbf{Đáp án:} B
\textbf{Giải thích:} Đáp án theo dữ liệu đề: B — B. NOT NULL.

Lời giải gốc: Ràng buộc NOT NULL đảm bảo rằng cột dữ liệu không được phép để trống (luôn phải có giá trị).

Nền tảng kiến thức: SQL tách truy vấn dữ liệu (DML) khỏi định nghĩa cấu trúc (DDL) và quyền truy cập (DCL). Với SELECT, cần xác định bảng nguồn, điều kiện lọc, phép kết hợp, nhóm và điều kiện nhóm; các ràng buộc, view, thủ tục và trigger đảm nhiệm các vai trò khác nhau.

Phân tích: NOT NULL là constraint cấm giá trị NULL ở cột, phù hợp khi mỗi hàng bắt buộc phải có giá trị. Nó không tự bảo đảm giá trị duy nhất (UNIQUE/PRIMARY KEY) và cũng không kiểm tra miền/range như CHECK.

Đối chiếu phương án: A. NULL; C. UNIQUE; D. DISTINCT. Các lựa chọn này không đáp ứng định nghĩa/điều kiện được hỏi; khi ôn, hãy chỉ ra chính xác điểm lệch của từng lựa chọn thay vì học thuộc mỗi chữ cái.
\end{question}

\begin{question}{427}
% [ID gốc: 5.218]
% [Nguồn: DBI202 Order 5 Q218]
Suppose we need to sort all students in the STUDENTS(id, name, age) table in descending order of name and descending order of age. Which command should we use?
\begin{choices}
\choice SELECT * FROM Students ORDER BY name, age DESC
\choice SELECT * FROM Students ORDER BY name DESC, age DESC
\choice Both A and B are correct
\choice Both A and B are wrong
\choice This request cannot be fulfilled
\end{choices}
\textbf{Đáp án:} B
\textbf{Giải thích:} Đáp án theo dữ liệu đề: B — B. SELECT * FROM Students ORDER BY name DESC, age DESC.

Lời giải gốc: Để sắp xếp giảm dần trên cả hai cột name và age, từ khóa DESC phải được áp dụng tường minh cho từng cột. Ở phương án A, cột name không có DESC nên mặc định sắp xếp tăng dần.

Nền tảng kiến thức: SQL tách truy vấn dữ liệu (DML) khỏi định nghĩa cấu trúc (DDL) và quyền truy cập (DCL). Với SELECT, cần xác định bảng nguồn, điều kiện lọc, phép kết hợp, nhóm và điều kiện nhóm; các ràng buộc, view, thủ tục và trigger đảm nhiệm các vai trò khác nhau.

Phân tích: Đọc truy vấn theo các giai đoạn logic: FROM/JOIN tạo dữ liệu nguồn, WHERE lọc từng hàng, GROUP BY tạo nhóm, HAVING lọc nhóm, SELECT xác định biểu thức đầu ra và ORDER BY sắp kết quả. Với phép join, kiểm tra điều kiện ghép và các hàng không khớp trước khi tính kết quả.

Đối chiếu phương án: A. SELECT * FROM Students ORDER BY name, age DESC; C. Both A and B are correct; D. Both A and B are wrong; E. This request cannot be fulfilled. Các lựa chọn này không đáp ứng định nghĩa/điều kiện được hỏi; khi ôn, hãy chỉ ra chính xác điểm lệch của từng lựa chọn thay vì học thuộc mỗi chữ cái.
\end{question}

\begin{question}{428}
% [ID gốc: 5.219]
% [Nguồn: DBI202 Order 5 Q219]
While executing a join operation between two tables, what type of join returns all rows from the left table and matched rows from the right table?
\begin{choices}
\choice Inner Join
\choice Right Join
\choice Full Outer Join
\choice Left Join
\end{choices}
\textbf{Đáp án:} D
\textbf{Giải thích:} Đáp án theo dữ liệu đề: D — D. Left Join.

Lời giải gốc: Phép nối LEFT JOIN (hoặc Left Outer Join) sẽ trả về toàn bộ tất cả các dòng dữ liệu ở bảng bên trái, kết hợp với các dòng có dữ liệu trùng khớp ở bảng bên phải. Nếu bên phải không có dữ liệu khớp, các cột bên phải sẽ mang giá trị NULL.

Nền tảng kiến thức: SQL tách truy vấn dữ liệu (DML) khỏi định nghĩa cấu trúc (DDL) và quyền truy cập (DCL). Với SELECT, cần xác định bảng nguồn, điều kiện lọc, phép kết hợp, nhóm và điều kiện nhóm; các ràng buộc, view, thủ tục và trigger đảm nhiệm các vai trò khác nhau.

Phân tích: Join ghép các bộ thỏa điều kiện. Inner join loại bộ không ghép được; LEFT/RIGHT OUTER JOIN giữ mọi bộ từ phía được chỉ định và điền NULL phía thiếu. Natural join tự ghép các thuộc tính trùng tên và chỉ xuất một cột cho mỗi thuộc tính ghép.

Đối chiếu phương án: A. Inner Join; B. Right Join; C. Full Outer Join. Các lựa chọn này không đáp ứng định nghĩa/điều kiện được hỏi; khi ôn, hãy chỉ ra chính xác điểm lệch của từng lựa chọn thay vì học thuộc mỗi chữ cái.
\end{question}

\begin{question}{429}
% [ID gốc: 5.222]
% [Nguồn: DBI202 Order 5 Q222]
Suppose relation R1(A, B) and relation R2(C, A) each have tuples. Which tuples are contained in the result of the query below?<pre>SELECT R1.B, R2.C FROM R1 RIGHT JOIN R2 ON R1.A = R2.A;</pre>
\begin{choices}
\choice (a, 101); (a, 104); (a, 109); (b, 102); (c, 103); (d, NULL); (e, NULL)
\choice (2, 101); (b, 102); (c, 103); (a, 104); (a, 105)
\choice (a, 101); (b, 102); (e, 103)
\choice (a, 101); (a, 104); (a, 109); (b, 102); (c, 103); (e, 105); (f, 106); (g, 107); (h, 108)
\end{choices}
\textbf{Đáp án:} D
\textbf{Giải thích:} RIGHT JOIN preserves all tuples from R2.
\end{question}

\begin{question}{430}
% [ID gốc: 5.223]
% [Nguồn: DBI202 Order 5 Q223]
Which keyword is used in SQL to filter grouped results?
\begin{choices}
\choice WHERE
\choice TOP
\choice GROUP BY
\choice HAVING
\end{choices}
\textbf{Đáp án:} D
\textbf{Giải thích:} Đáp án theo dữ liệu đề: D — D. HAVING.

Lời giải gốc: Tương tự câu 30, từ khóa để lọc các kết quả sau khi đã áp dụng gom nhóm là HAVING.

Nền tảng kiến thức: SQL tách truy vấn dữ liệu (DML) khỏi định nghĩa cấu trúc (DDL) và quyền truy cập (DCL). Với SELECT, cần xác định bảng nguồn, điều kiện lọc, phép kết hợp, nhóm và điều kiện nhóm; các ràng buộc, view, thủ tục và trigger đảm nhiệm các vai trò khác nhau.

Phân tích: Đọc truy vấn theo các giai đoạn logic: FROM/JOIN tạo dữ liệu nguồn, WHERE lọc từng hàng, GROUP BY tạo nhóm, HAVING lọc nhóm, SELECT xác định biểu thức đầu ra và ORDER BY sắp kết quả. Với phép join, kiểm tra điều kiện ghép và các hàng không khớp trước khi tính kết quả.

Đối chiếu phương án: A. WHERE; B. TOP; C. GROUP BY. Các lựa chọn này không đáp ứng định nghĩa/điều kiện được hỏi; khi ôn, hãy chỉ ra chính xác điểm lệch của từng lựa chọn thay vì học thuộc mỗi chữ cái.
\end{question}

\begin{question}{431}
% [ID gốc: 5.224]
% [Nguồn: DBI202 Order 5 Q224]
In SQL, what happens when you use a RIGHT JOIN between two tables with unmatched records in the left table?
\begin{choices}
\choice All records from the left table are returned.
\choice Only matching records are returned.
\choice NULL values are filled for unmatched rows from the left.
\choice NULL values are filled for unmatched rows from the right.
\end{choices}
\textbf{Đáp án:} C
\textbf{Giải thích:} Đáp án theo dữ liệu đề: D — D. NULL values are filled for unmatched rows from the right..

Lời giải gốc: Phép RIGHT JOIN giữ lại toàn bộ các hàng từ bảng bên phải và điền giá trị NULL cho các hàng không có cặp tương ứng từ bảng bên trái. ---

Nền tảng kiến thức: SQL tách truy vấn dữ liệu (DML) khỏi định nghĩa cấu trúc (DDL) và quyền truy cập (DCL). Với SELECT, cần xác định bảng nguồn, điều kiện lọc, phép kết hợp, nhóm và điều kiện nhóm; các ràng buộc, view, thủ tục và trigger đảm nhiệm các vai trò khác nhau.

Phân tích: Join ghép các bộ thỏa điều kiện. Inner join loại bộ không ghép được; LEFT/RIGHT OUTER JOIN giữ mọi bộ từ phía được chỉ định và điền NULL phía thiếu. Natural join tự ghép các thuộc tính trùng tên và chỉ xuất một cột cho mỗi thuộc tính ghép.

Đối chiếu phương án: A. All records from the left table are returned.; B. Only matching records are returned.; C. NULL values are filled for unmatched rows from the left.. Các lựa chọn này không đáp ứng định nghĩa/điều kiện được hỏi; khi ôn, hãy chỉ ra chính xác điểm lệch của từng lựa chọn thay vì học thuộc mỗi chữ cái.
\end{question}

\begin{question}{432}
% [ID gốc: 5.225]
% [Nguồn: DBI202 Order 5 Q225]
Which of the following is used to delete rows where "model" is "ford" from the Cars table?
\begin{choices}
\choice DELETE FROM Cars WHERE model = 'ford';
\choice DELETE Cars WHERE model = 'ford';
\choice DELETE FROM Cars WITH TIE model IS 'ford';
\choice DELETE model = 'ford' WHERE Cars;
\end{choices}
\textbf{Đáp án:} A
\textbf{Giải thích:} Đáp án theo dữ liệu đề: A — A. DELETE FROM Cars WHERE model = 'ford';.

Lời giải gốc: Cú pháp chuẩn của lệnh xóa bản ghi trong SQL là DELETE FROM table_name WHERE condition;.

Nền tảng kiến thức: SQL tách truy vấn dữ liệu (DML) khỏi định nghĩa cấu trúc (DDL) và quyền truy cập (DCL). Với SELECT, cần xác định bảng nguồn, điều kiện lọc, phép kết hợp, nhóm và điều kiện nhóm; các ràng buộc, view, thủ tục và trigger đảm nhiệm các vai trò khác nhau.

Phân tích: Đọc truy vấn theo các giai đoạn logic: FROM/JOIN tạo dữ liệu nguồn, WHERE lọc từng hàng, GROUP BY tạo nhóm, HAVING lọc nhóm, SELECT xác định biểu thức đầu ra và ORDER BY sắp kết quả. Với phép join, kiểm tra điều kiện ghép và các hàng không khớp trước khi tính kết quả.

Đối chiếu phương án: B. DELETE Cars WHERE model = 'ford';; C. DELETE FROM Cars WITH TIE model IS 'ford';; D. DELETE model = 'ford' WHERE Cars;. Các lựa chọn này không đáp ứng định nghĩa/điều kiện được hỏi; khi ôn, hãy chỉ ra chính xác điểm lệch của từng lựa chọn thay vì học thuộc mỗi chữ cái.
\end{question}

\begin{question}{433}
% [ID gốc: 5.226]
% [Nguồn: DBI202 Order 5 Q226]
How many joins are required to join 6 tables in a straight chain?
\begin{choices}
\choice 4
\choice 5
\choice 6
\choice 12
\end{choices}
\textbf{Đáp án:} B
\textbf{Giải thích:} Đáp án theo dữ liệu đề: B — B. 5.

Lời giải gốc: Để kết nối thẳng n bảng với nhau, số lượng phép kết nối (joins) cần thiết luôn là n - 1. Với 6 bảng, cần 6 - 1 = 5 phép join.

Nền tảng kiến thức: SQL tách truy vấn dữ liệu (DML) khỏi định nghĩa cấu trúc (DDL) và quyền truy cập (DCL). Với SELECT, cần xác định bảng nguồn, điều kiện lọc, phép kết hợp, nhóm và điều kiện nhóm; các ràng buộc, view, thủ tục và trigger đảm nhiệm các vai trò khác nhau.

Phân tích: Join ghép các bộ thỏa điều kiện. Inner join loại bộ không ghép được; LEFT/RIGHT OUTER JOIN giữ mọi bộ từ phía được chỉ định và điền NULL phía thiếu. Natural join tự ghép các thuộc tính trùng tên và chỉ xuất một cột cho mỗi thuộc tính ghép.

Đối chiếu phương án: A. 4; C. 6; D. 12. Các lựa chọn này không đáp ứng định nghĩa/điều kiện được hỏi; khi ôn, hãy chỉ ra chính xác điểm lệch của từng lựa chọn thay vì học thuộc mỗi chữ cái.
\end{question}

\begin{question}{434}
% [ID gốc: 5.227]
% [Nguồn: DBI202 Order 5 Q227]
Which statement best describes the DELETE command?
\begin{choices}
\choice Deletes all rows and cannot be rolled back
\choice Deletes the table structure and all its data
\choice Deletes specific rows and can be rolled back
\end{choices}
\textbf{Đáp án:} C
\textbf{Giải thích:} Đáp án theo dữ liệu đề: C — C. Deletes specific rows and can be rolled back.

Lời giải gốc: Lệnh DELETE thuộc nhóm DML, dùng để xóa các hàng cụ thể (hoặc toàn bộ hàng nếu không có WHERE) dựa trên điều kiện và hành động này được ghi vào log nên hoàn toàn có thể khôi phục (rollback) lại được nếu nằm trong một transaction chưa commit. (Khác với TRUNCATE không thể rollback dễ dàng hoặc DROP làm mất luôn cấu trúc bảng).

Nền tảng kiến thức: SQL tách truy vấn dữ liệu (DML) khỏi định nghĩa cấu trúc (DDL) và quyền truy cập (DCL). Với SELECT, cần xác định bảng nguồn, điều kiện lọc, phép kết hợp, nhóm và điều kiện nhóm; các ràng buộc, view, thủ tục và trigger đảm nhiệm các vai trò khác nhau.

Phân tích: DELETE loại bỏ các hàng thỏa điều kiện WHERE; bỏ WHERE thì có thể xóa mọi hàng nhưng vẫn giữ cấu trúc bảng. Đây khác DROP TABLE (xóa đối tượng bảng) và UPDATE (sửa giá trị); có thể xóa một hoặc nhiều hàng trong cùng một bảng.

Đối chiếu phương án: A. Deletes all rows and cannot be rolled back; B. Deletes the table structure and all its data. Các lựa chọn này không đáp ứng định nghĩa/điều kiện được hỏi; khi ôn, hãy chỉ ra chính xác điểm lệch của từng lựa chọn thay vì học thuộc mỗi chữ cái.
\end{question}

\begin{question}{435}
% [ID gốc: 5.228]
% [Nguồn: DBI202 Order 5 Q228]
What is the result of a CROSS JOIN in SQL?
\begin{choices}
\choice Union of tables
\choice Intersection of tables
\choice Cartesian product
\choice Natural join
\end{choices}
\textbf{Đáp án:} C
\textbf{Giải thích:} Đáp án theo dữ liệu đề: C — C. Cartesian product.

Lời giải gốc: Truy vấn CROSS JOIN (hay phép tích chéo) sẽ ghép từng dòng của bảng A với toàn bộ các dòng của bảng B, tạo thành một Tích Đề-các (Cartesian product). Số lượng dòng kết quả bằng (số dòng bảng A) x (số dòng bảng B).

Nền tảng kiến thức: SQL tách truy vấn dữ liệu (DML) khỏi định nghĩa cấu trúc (DDL) và quyền truy cập (DCL). Với SELECT, cần xác định bảng nguồn, điều kiện lọc, phép kết hợp, nhóm và điều kiện nhóm; các ràng buộc, view, thủ tục và trigger đảm nhiệm các vai trò khác nhau.

Phân tích: Join ghép các bộ thỏa điều kiện. Inner join loại bộ không ghép được; LEFT/RIGHT OUTER JOIN giữ mọi bộ từ phía được chỉ định và điền NULL phía thiếu. Natural join tự ghép các thuộc tính trùng tên và chỉ xuất một cột cho mỗi thuộc tính ghép.

Đối chiếu phương án: A. Union of tables; B. Intersection of tables; D. Natural join. Các lựa chọn này không đáp ứng định nghĩa/điều kiện được hỏi; khi ôn, hãy chỉ ra chính xác điểm lệch của từng lựa chọn thay vì học thuộc mỗi chữ cái.
\end{question}

\begin{question}{436}
% [ID gốc: 5.229]
% [Nguồn: DBI202 Order 5 Q229]
The \% wildcard when used with the LIKE keyword is used to represent a string of characters:
\begin{choices}
\choice Has a minimum length of 0 characters.
\choice Has a minimum length of 1 character, a maximum length of n characters.
\choice Has a minimum length of 1 character, a maximum length of 255 characters.
\choice Has a minimum length of 1 character, a maximum length of 256 characters.
\choice None of the above answers.
\end{choices}
\textbf{Đáp án:} A
\textbf{Giải thích:} Đáp án theo dữ liệu đề: A — A. Has a minimum length of 0 characters..

Lời giải gốc: Ký tự đại diện \% trong lệnh LIKE đại diện cho một chuỗi gồm bất kỳ ký tự nào, có chiều dài tối thiểu là 0 ký tự. ---

Nền tảng kiến thức: SQL tách truy vấn dữ liệu (DML) khỏi định nghĩa cấu trúc (DDL) và quyền truy cập (DCL). Với SELECT, cần xác định bảng nguồn, điều kiện lọc, phép kết hợp, nhóm và điều kiện nhóm; các ràng buộc, view, thủ tục và trigger đảm nhiệm các vai trò khác nhau.

Phân tích: Trong LIKE của SQL Server, `%` khớp 0 hoặc nhiều ký tự và `_` khớp đúng một ký tự. ESCAPE cho phép coi ký tự wildcard như ký tự thường bằng cách đặt trước chúng một ký tự thoát; nó không thay thế mẫu bằng chuỗi bất kỳ.

Đối chiếu phương án: B. Has a minimum length of 1 character, a maximum length of n characters.; C. Has a minimum length of 1 character, a maximum length of 255 characters.; D. Has a minimum length of 1 character, a maximum length of 256 characters.; E. None of the above answers.. Các lựa chọn này không đáp ứng định nghĩa/điều kiện được hỏi; khi ôn, hãy chỉ ra chính xác điểm lệch của từng lựa chọn thay vì học thuộc mỗi chữ cái.
\end{question}

\begin{question}{437}
% [ID gốc: 5.232]
% [Nguồn: DBI202 Order 5 Q232]
Suppose relation R(A, B) has tuples. Which tuples are contained in the result of the query below?<pre>SELECT * FROM R WHERE A LIKE 'a%' UNION SELECT * FROM R WHERE B &gt; 7;</pre>A. (3, 5); (h, 8)
\begin{choices}
\choice (3, 5); (3, 9); (h, 8)
\choice (3, 9); (b, 7); (h, 8)
\choice (2, 9)
\end{choices}
\textbf{Đáp án:} C
\textbf{Giải thích:} Đáp án theo dữ liệu đề: C — C. (3, 9); (b, 7); (h, 8).

Lời giải gốc: Phép UNION kết hợp các bản ghi thỏa mãn điều kiện A bắt đầu bằng 'a' hoặc có giá trị B lớn hơn 7, đồng thời loại bỏ các bản ghi trùng lặp.

Nền tảng kiến thức: SQL tách truy vấn dữ liệu (DML) khỏi định nghĩa cấu trúc (DDL) và quyền truy cập (DCL). Với SELECT, cần xác định bảng nguồn, điều kiện lọc, phép kết hợp, nhóm và điều kiện nhóm; các ràng buộc, view, thủ tục và trigger đảm nhiệm các vai trò khác nhau.

Phân tích: Các hàng dữ liệu của một relation được gọi là tuple; header chứa tên thuộc tính không phải tuple dữ liệu. Trong bảng, tuple tương ứng một bản ghi, còn attribute tương ứng cột.

Đối chiếu phương án: B. (3, 5); (3, 9); (h, 8); D. (2, 9). Các lựa chọn này không đáp ứng định nghĩa/điều kiện được hỏi; khi ôn, hãy chỉ ra chính xác điểm lệch của từng lựa chọn thay vì học thuộc mỗi chữ cái.
\end{question}

\begin{question}{438}
% [ID gốc: 5.233]
% [Nguồn: DBI202 Order 5 Q233]
Which of the following SQL statements returns only matching records between Customers and SalesInvoice based on custID?
\begin{choices}
\choice SELECT * FROM Customers LEFT JOIN SalesInvoice ON Customers.custID = SalesInvoice.custID;
\choice SELECT * FROM Customers INNER JOIN SalesInvoice ON Customers.custID = SalesInvoice.custID;
\choice SELECT * FROM Customers RIGHT JOIN SalesInvoice ON Customers.custID = SalesInvoice.custID;
\choice SELECT * FROM Customers FULL JOIN SalesInvoice ON Customers.custID = SalesInvoice.custID;
\end{choices}
\textbf{Đáp án:} B
\textbf{Giải thích:} Đáp án theo dữ liệu đề: B — B. SELECT * FROM Customers INNER JOIN SalesInvoice ON Customers.custID = SalesInvoice.custID;.

Lời giải gốc: Phép kết nội (INNER JOIN) chỉ trả về các bản ghi khớp nhau ở cả hai bảng dựa trên điều kiện kết nối.

Nền tảng kiến thức: SQL tách truy vấn dữ liệu (DML) khỏi định nghĩa cấu trúc (DDL) và quyền truy cập (DCL). Với SELECT, cần xác định bảng nguồn, điều kiện lọc, phép kết hợp, nhóm và điều kiện nhóm; các ràng buộc, view, thủ tục và trigger đảm nhiệm các vai trò khác nhau.

Phân tích: Join ghép các bộ thỏa điều kiện. Inner join loại bộ không ghép được; LEFT/RIGHT OUTER JOIN giữ mọi bộ từ phía được chỉ định và điền NULL phía thiếu. Natural join tự ghép các thuộc tính trùng tên và chỉ xuất một cột cho mỗi thuộc tính ghép.

Đối chiếu phương án: A. SELECT * FROM Customers LEFT JOIN SalesInvoice ON Customers.custID = SalesInvoice.custID;; C. SELECT * FROM Customers RIGHT JOIN SalesInvoice ON Customers.custID = SalesInvoice.custID;; D. SELECT * FROM Customers FULL JOIN SalesInvoice ON Customers.custID = SalesInvoice.custID;. Các lựa chọn này không đáp ứng định nghĩa/điều kiện được hỏi; khi ôn, hãy chỉ ra chính xác điểm lệch của từng lựa chọn thay vì học thuộc mỗi chữ cái.
\end{question}

\begin{question}{439}
% [ID gốc: 5.234]
% [Nguồn: DBI202 Order 5 Q234]
What does the following SQL query return? SELECT name FROM Products WHERE price &gt; 100;
\begin{choices}
\choice All product names with price less than 100
\choice All product names regardless of price
\choice Product names with price greater than 100
\choice Syntax error
\end{choices}
\textbf{Đáp án:} C
\textbf{Giải thích:} Đáp án theo dữ liệu đề: C — C. Product names with price greater than 100.

Lời giải gốc: Câu truy vấn chọn cột name từ bảng Products thỏa mãn điều kiện giá price lớn hơn 100.

Nền tảng kiến thức: SQL tách truy vấn dữ liệu (DML) khỏi định nghĩa cấu trúc (DDL) và quyền truy cập (DCL). Với SELECT, cần xác định bảng nguồn, điều kiện lọc, phép kết hợp, nhóm và điều kiện nhóm; các ràng buộc, view, thủ tục và trigger đảm nhiệm các vai trò khác nhau.

Phân tích: Đọc truy vấn theo các giai đoạn logic: FROM/JOIN tạo dữ liệu nguồn, WHERE lọc từng hàng, GROUP BY tạo nhóm, HAVING lọc nhóm, SELECT xác định biểu thức đầu ra và ORDER BY sắp kết quả. Với phép join, kiểm tra điều kiện ghép và các hàng không khớp trước khi tính kết quả.

Đối chiếu phương án: A. All product names with price less than 100; B. All product names regardless of price; D. Syntax error. Các lựa chọn này không đáp ứng định nghĩa/điều kiện được hỏi; khi ôn, hãy chỉ ra chính xác điểm lệch của từng lựa chọn thay vì học thuộc mỗi chữ cái.
\end{question}

\begin{question}{440}
% [ID gốc: 5.235]
% [Nguồn: DBI202 Order 5 Q235]
The Data Definition Language (DDL) in SQL is used to:
\begin{choices}
\choice Insert, update, and delete data from a database
\choice Define and modify the structure of database objects such as tables, views, and indexes
\choice Query and retrieve data from a database
\choice Manage user permissions and access rights
\end{choices}
\textbf{Đáp án:} B
\textbf{Giải thích:} Đáp án theo dữ liệu đề: B — B. Define and modify the structure of database objects such as tables, views, and indexes.

Lời giải gốc: Ngôn ngữ định nghĩa dữ liệu (DDL) được dùng để tạo, sửa đổi và xóa bỏ cấu trúc các đối tượng cơ sở dữ liệu.

Nền tảng kiến thức: SQL tách truy vấn dữ liệu (DML) khỏi định nghĩa cấu trúc (DDL) và quyền truy cập (DCL). Với SELECT, cần xác định bảng nguồn, điều kiện lọc, phép kết hợp, nhóm và điều kiện nhóm; các ràng buộc, view, thủ tục và trigger đảm nhiệm các vai trò khác nhau.

Phân tích: DDL (Data Definition Language) khai báo/thay đổi cấu trúc như bảng và schema; DML truy vấn/thêm/sửa/xóa dữ liệu; DCL quản lý quyền. Chọn theo đối tượng bị tác động: cấu trúc, các hàng hay quyền truy cập.

Đối chiếu phương án: A. Insert, update, and delete data from a database; C. Query and retrieve data from a database; D. Manage user permissions and access rights. Các lựa chọn này không đáp ứng định nghĩa/điều kiện được hỏi; khi ôn, hãy chỉ ra chính xác điểm lệch của từng lựa chọn thay vì học thuộc mỗi chữ cái.
\end{question}

\begin{question}{441}
% [ID gốc: 5.236]
% [Nguồn: DBI202 Order 5 Q236]
Given two relations: R(A, B, C, D) with the primary key is A, S(D, E, F) with the unique key is D. How to make a R.D foreign key referencing S.D?
\begin{choices}
\choice ADD CONSTRAINT fk_R_D FOREIGN KEY R(D) REFERENCES S(D)
\choice ALTER TABLE R ALTER CONSTRAINT fk_R_D FOREIGN KEY (D) REFERENCES S(D)
\choice Can not add a foreign key in the R table because the S's D column is not the primary key.
\choice ALTER TABLE R ADD CONSTRAINT fk_R_D FOREIGN KEY (D) REFERENCES S(D)
\end{choices}
\textbf{Đáp án:} D
\textbf{Giải thích:} Đáp án theo dữ liệu đề: D — D. ALTER TABLE R ADD CONSTRAINT fk_R_D FOREIGN KEY (D) REFERENCES S(D).

Lời giải gốc: Cú pháp chuẩn để thêm khóa ngoại tham chiếu đến bảng khác trong SQL là ALTER TABLE R ADD CONSTRAINT fk_R_D FOREIGN KEY (D) REFERENCES S(D). ---

Nền tảng kiến thức: SQL tách truy vấn dữ liệu (DML) khỏi định nghĩa cấu trúc (DDL) và quyền truy cập (DCL). Với SELECT, cần xác định bảng nguồn, điều kiện lọc, phép kết hợp, nhóm và điều kiện nhóm; các ràng buộc, view, thủ tục và trigger đảm nhiệm các vai trò khác nhau.

Phân tích: Khóa ngoại buộc giá trị tham chiếu phải khớp một khóa được tham chiếu ở quan hệ đích (hoặc NULL nếu ràng buộc cho phép). Đây là toàn vẹn tham chiếu; nó khác tính duy nhất của primary/candidate key.

Đối chiếu phương án: A. ADD CONSTRAINT fk_R_D FOREIGN KEY R(D) REFERENCES S(D); B. ALTER TABLE R ALTER CONSTRAINT fk_R_D FOREIGN KEY (D) REFERENCES S(D); C. Can not add a foreign key in the R table because the S's D column is not the primary key.. Các lựa chọn này không đáp ứng định nghĩa/điều kiện được hỏi; khi ôn, hãy chỉ ra chính xác điểm lệch của từng lựa chọn thay vì học thuộc mỗi chữ cái.
\end{question}

\begin{question}{442}
% [ID gốc: 5.239]
% [Nguồn: DBI202 Order 5 Q239]
In the query below, which keyword must be inserted?<pre>INSERT INTO Students ____ (32, N'Iona Bush', N'Female');</pre>
\begin{choices}
\choice INTO
\choice SET
\choice WHERE
\choice VALUES
\end{choices}
\textbf{Đáp án:} D
\textbf{Giải thích:} Standard SQL INSERT syntax requires the VALUES keyword before column data.
\end{question}

\begin{question}{443}
% [ID gốc: 5.250]
% [Nguồn: DBI202 Order 5 Q250]
Which SQL statements will correctly update the salary of all employees by 10\% where job title is 'Manager'?
\begin{choices}
\choice UPDATE Employees SET salary = salary * 1.1 WHERE job = 'Manager';
\choice UPDATE Employees WHERE job = 'Manager' SET salary = salary * 1.1;
\choice UPDATE Employees SET salary *= 1.1 WHERE job = 'Manager';
\choice MODIFY Employees SET salary = salary + salary*0.1 WHERE job = 'Manager';
\end{choices}
\textbf{Đáp án:} A
\textbf{Giải thích:} Đáp án theo dữ liệu đề: A — A. UPDATE Employees SET salary = salary * 1.1 WHERE job = 'Manager';.

Lời giải gốc: Cú pháp chuẩn của câu lệnh UPDATE trong SQL: UPDATE table_name SET column = value WHERE condition;.

Nền tảng kiến thức: SQL tách truy vấn dữ liệu (DML) khỏi định nghĩa cấu trúc (DDL) và quyền truy cập (DCL). Với SELECT, cần xác định bảng nguồn, điều kiện lọc, phép kết hợp, nhóm và điều kiện nhóm; các ràng buộc, view, thủ tục và trigger đảm nhiệm các vai trò khác nhau.

Phân tích: Đọc truy vấn theo các giai đoạn logic: FROM/JOIN tạo dữ liệu nguồn, WHERE lọc từng hàng, GROUP BY tạo nhóm, HAVING lọc nhóm, SELECT xác định biểu thức đầu ra và ORDER BY sắp kết quả. Với phép join, kiểm tra điều kiện ghép và các hàng không khớp trước khi tính kết quả.

Đối chiếu phương án: B. UPDATE Employees WHERE job = 'Manager' SET salary = salary * 1.1;; C. UPDATE Employees SET salary *= 1.1 WHERE job = 'Manager';; D. MODIFY Employees SET salary = salary + salary*0.1 WHERE job = 'Manager';. Các lựa chọn này không đáp ứng định nghĩa/điều kiện được hỏi; khi ôn, hãy chỉ ra chính xác điểm lệch của từng lựa chọn thay vì học thuộc mỗi chữ cái.
\end{question}

\begin{question}{444}
% [ID gốc: 5.252]
% [Nguồn: DBI202 Order 5 Q252]
Given a relation: R(A,B,C,D) with the key is A, the foreign key is D And the query: "select A, count(D) from R group by D" What is the output of the query?
\begin{choices}
\choice It will return all rows from the R table
\choice It returns only the A column where the value in D column is grouped
\choice This query will cause an error in standard SQL
\choice There is no GROUP BY with a HAVING clause.
\end{choices}
\textbf{Đáp án:} C
\textbf{Giải thích:} Đáp án theo dữ liệu đề: C — C. This query will cause an error in standard SQL.

Lời giải gốc: Trong chuẩn SQL, các cột xuất hiện ở mệnh đề SELECT không nằm trong hàm gom nhóm (Aggregate Function) bắt buộc phải xuất hiện trong mệnh đề GROUP BY. Do cột A không nằm trong GROUP BY D nên truy vấn bị lỗi.

Nền tảng kiến thức: SQL tách truy vấn dữ liệu (DML) khỏi định nghĩa cấu trúc (DDL) và quyền truy cập (DCL). Với SELECT, cần xác định bảng nguồn, điều kiện lọc, phép kết hợp, nhóm và điều kiện nhóm; các ràng buộc, view, thủ tục và trigger đảm nhiệm các vai trò khác nhau.

Phân tích: Khóa ngoại buộc giá trị tham chiếu phải khớp một khóa được tham chiếu ở quan hệ đích (hoặc NULL nếu ràng buộc cho phép). Đây là toàn vẹn tham chiếu; nó khác tính duy nhất của primary/candidate key.

Đối chiếu phương án: A. It will return all rows from the R table; B. It returns only the A column where the value in D column is grouped; D. There is no GROUP BY with a HAVING clause.. Các lựa chọn này không đáp ứng định nghĩa/điều kiện được hỏi; khi ôn, hãy chỉ ra chính xác điểm lệch của từng lựa chọn thay vì học thuộc mỗi chữ cái.
\end{question}

\begin{question}{445}
% [ID gốc: 5.257]
% [Nguồn: DBI202 Order 5 Q257]
Which of the following is NOT true about an aggregate function?
\begin{choices}
\choice All aggregate functions are nondeterministic.
\choice It performs a calculation on a set of values and returns a single value.
\choice Aggregate functions such as COUNT(*) ignore null values.
\choice Aggregate functions are often used with the GROUP BY clause of the SELECT statement.
\end{choices}
\textbf{Đáp án:} A
\textbf{Giải thích:} Đáp án theo dữ liệu đề: A — A. All aggregate functions are nondeterministic..

Lời giải gốc: Phát biểu này không đúng vì các hàm tổng hợp chuẩn (như SUM, AVG, MIN, MAX) là hàm xác định (deterministic) - cùng đầu vào luôn cho cùng kết quả.

Nền tảng kiến thức: SQL tách truy vấn dữ liệu (DML) khỏi định nghĩa cấu trúc (DDL) và quyền truy cập (DCL). Với SELECT, cần xác định bảng nguồn, điều kiện lọc, phép kết hợp, nhóm và điều kiện nhóm; các ràng buộc, view, thủ tục và trigger đảm nhiệm các vai trò khác nhau.

Phân tích: Aggregate function tính một kết quả từ một tập nhiều hàng/giá trị, ví dụ COUNT, SUM, AVG, MIN và MAX. Trong truy vấn nhóm, GROUP BY xác định từng tập đầu vào; HAVING lọc các nhóm sau khi tính tổng hợp.

Đối chiếu phương án: B. It performs a calculation on a set of values and returns a single value.; C. Aggregate functions such as COUNT(*) ignore null values.; D. Aggregate functions are often used with the GROUP BY clause of the SELECT statement.. Các lựa chọn này không đáp ứng định nghĩa/điều kiện được hỏi; khi ôn, hãy chỉ ra chính xác điểm lệch của từng lựa chọn thay vì học thuộc mỗi chữ cái.
\end{question}

\begin{question}{446}
% [ID gốc: 5.260]
% [Nguồn: DBI202 Order 5 Q260]
Which of the following queries deletes employees who are not assigned to any department?
\begin{choices}
\choice DELETE FROM Employees WHERE DepartmentID IS NULL;
\choice DELETE FROM Employees WHERE DepartmentID = '';
\choice DELETE FROM Employees WHERE DepartmentID = 0;
\choice DELETE FROM Employees WHERE DepartmentID IS NOT NULL;
\end{choices}
\textbf{Đáp án:} A
\textbf{Giải thích:} Đáp án theo dữ liệu đề: A — A. DELETE FROM Employees WHERE DepartmentID IS NULL;.

Lời giải gốc: Để tìm hoặc xóa các bản ghi mang giá trị trống/không xác định, ta dùng toán tử chuẩn IS NULL. Không thể so sánh NULL bằng dấu = như trong phương án B, C.

Nền tảng kiến thức: SQL tách truy vấn dữ liệu (DML) khỏi định nghĩa cấu trúc (DDL) và quyền truy cập (DCL). Với SELECT, cần xác định bảng nguồn, điều kiện lọc, phép kết hợp, nhóm và điều kiện nhóm; các ràng buộc, view, thủ tục và trigger đảm nhiệm các vai trò khác nhau.

Phân tích: NULL biểu thị thiếu/không biết/không áp dụng, không phải số 0 hay chuỗi rỗng. So sánh với NULL bằng `=` không cho TRUE/FALSE mà cho UNKNOWN; dùng `IS NULL`/`IS NOT NULL`. Trong logic ba giá trị, AND/OR/NOT phải được tính theo cả trạng thái UNKNOWN.

Đối chiếu phương án: B. DELETE FROM Employees WHERE DepartmentID = '';; C. DELETE FROM Employees WHERE DepartmentID = 0;; D. DELETE FROM Employees WHERE DepartmentID IS NOT NULL;. Các lựa chọn này không đáp ứng định nghĩa/điều kiện được hỏi; khi ôn, hãy chỉ ra chính xác điểm lệch của từng lựa chọn thay vì học thuộc mỗi chữ cái.
\end{question}

\begin{question}{447}
% [ID gốc: 5.267]
% [Nguồn: DBI202 Order 5 Q267]
Which of the following wildcard characters matches any sequence of zero or more characters in a string comparison in SQL?
\begin{choices}
\choice &amp;
\choice $
\choice _
\choice %
\end{choices}
\textbf{Đáp án:} D
\textbf{Giải thích:} Đáp án theo dữ liệu đề: D — D. %.

Lời giải gốc: Trong SQL, ký tự đại diện \% khớp với một chuỗi gồm 0 hoặc nhiều ký tự bất kỳ (ký tự _ chỉ khớp với đúng 1 ký tự).

Nền tảng kiến thức: SQL tách truy vấn dữ liệu (DML) khỏi định nghĩa cấu trúc (DDL) và quyền truy cập (DCL). Với SELECT, cần xác định bảng nguồn, điều kiện lọc, phép kết hợp, nhóm và điều kiện nhóm; các ràng buộc, view, thủ tục và trigger đảm nhiệm các vai trò khác nhau.

Phân tích: Trong LIKE của SQL Server, `%` khớp 0 hoặc nhiều ký tự và `_` khớp đúng một ký tự. ESCAPE cho phép coi ký tự wildcard như ký tự thường bằng cách đặt trước chúng một ký tự thoát; nó không thay thế mẫu bằng chuỗi bất kỳ.

Đối chiếu phương án: A. &amp;; B. $; C. _. Các lựa chọn này không đáp ứng định nghĩa/điều kiện được hỏi; khi ôn, hãy chỉ ra chính xác điểm lệch của từng lựa chọn thay vì học thuộc mỗi chữ cái.
\end{question}

\begin{question}{448}
% [ID gốc: 5.269]
% [Nguồn: DBI202 Order 5 Q269]
Which of the following is NOT a type of join?
\begin{choices}
\choice Inner join
\choice Outer join
\choice Cross join
\choice Super join
\end{choices}
\textbf{Đáp án:} D
\textbf{Giải thích:} Đáp án theo dữ liệu đề: D — D. Super join.

Lời giải gốc: Super join không phải là một loại phép kết nối trong SQL, trong khi Inner join, Outer join và Cross join là các phép join chuẩn.

Nền tảng kiến thức: SQL tách truy vấn dữ liệu (DML) khỏi định nghĩa cấu trúc (DDL) và quyền truy cập (DCL). Với SELECT, cần xác định bảng nguồn, điều kiện lọc, phép kết hợp, nhóm và điều kiện nhóm; các ràng buộc, view, thủ tục và trigger đảm nhiệm các vai trò khác nhau.

Phân tích: Join ghép các bộ thỏa điều kiện. Inner join loại bộ không ghép được; LEFT/RIGHT OUTER JOIN giữ mọi bộ từ phía được chỉ định và điền NULL phía thiếu. Natural join tự ghép các thuộc tính trùng tên và chỉ xuất một cột cho mỗi thuộc tính ghép.

Đối chiếu phương án: A. Inner join; B. Outer join; C. Cross join. Các lựa chọn này không đáp ứng định nghĩa/điều kiện được hỏi; khi ôn, hãy chỉ ra chính xác điểm lệch của từng lựa chọn thay vì học thuộc mỗi chữ cái.
\end{question}

\begin{question}{449}
% [ID gốc: 5.271]
% [Nguồn: DBI202 Order 5 Q271]
Suppose you have a table Employees (id, name, salary). To increase the salary of employees with a salary below 5000 by 10%, which of the following statements is correc
\begin{choices}
\choice UPDATE Employees SET salary = salary + 0.1 WHERE salary &lt; 5000
\choice UPDATE Employees SET salary = salary * 1.1 WHERE salary &lt; 5000
\choice UPDATE Employees
SET salary = salary * 0.1 WHERE salary &lt; 5000
\choice UPDATE Employees
SET salary = 1.1 WHERE salary &lt; 5000
\end{choices}
\textbf{Đáp án:} B
\textbf{Giải thích:} Đáp án theo dữ liệu đề: B — B. UPDATE Employees SET salary = salary * 1.1 WHERE salary &lt; 5000.

Lời giải gốc: Tăng lương thêm 10\% đồng nghĩa với việc nhân mức lương hiện tại với hệ số $1 + 0.1 = 1.1$. ---

Nền tảng kiến thức: SQL tách truy vấn dữ liệu (DML) khỏi định nghĩa cấu trúc (DDL) và quyền truy cập (DCL). Với SELECT, cần xác định bảng nguồn, điều kiện lọc, phép kết hợp, nhóm và điều kiện nhóm; các ràng buộc, view, thủ tục và trigger đảm nhiệm các vai trò khác nhau.

Phân tích: Đọc truy vấn theo các giai đoạn logic: FROM/JOIN tạo dữ liệu nguồn, WHERE lọc từng hàng, GROUP BY tạo nhóm, HAVING lọc nhóm, SELECT xác định biểu thức đầu ra và ORDER BY sắp kết quả. Với phép join, kiểm tra điều kiện ghép và các hàng không khớp trước khi tính kết quả.

Đối chiếu phương án: A. UPDATE Employees SET salary = salary + 0.1 WHERE salary &lt; 5000; C. UPDATE Employees SET salary = salary * 0.1 WHERE salary &lt; 5000; D. UPDATE Employees SET salary = 1.1 WHERE salary &lt; 5000. Các lựa chọn này không đáp ứng định nghĩa/điều kiện được hỏi; khi ôn, hãy chỉ ra chính xác điểm lệch của từng lựa chọn thay vì học thuộc mỗi chữ cái.
\end{question}

\begin{question}{450}
% [ID gốc: 5.273]
% [Nguồn: DBI202 Order 5 Q273]
Which of the following is NOT a type of SQL constraint?
\begin{choices}
\choice PRIMARY KEY
\choice FOREIGN KEY
\choice INDEX KEY
\choice CHECK
\end{choices}
\textbf{Đáp án:} C
\textbf{Giải thích:} Đáp án theo dữ liệu đề: C — C. INDEX KEY.

Lời giải gốc: INDEX KEY không phải là một loại ràng buộc (constraint) chuẩn trong SQL, trong khi PRIMARY KEY, FOREIGN KEY, và CHECK đều là các ràng buộc dữ liệu.

Nền tảng kiến thức: SQL tách truy vấn dữ liệu (DML) khỏi định nghĩa cấu trúc (DDL) và quyền truy cập (DCL). Với SELECT, cần xác định bảng nguồn, điều kiện lọc, phép kết hợp, nhóm và điều kiện nhóm; các ràng buộc, view, thủ tục và trigger đảm nhiệm các vai trò khác nhau.

Phân tích: Khóa ngoại buộc giá trị tham chiếu phải khớp một khóa được tham chiếu ở quan hệ đích (hoặc NULL nếu ràng buộc cho phép). Đây là toàn vẹn tham chiếu; nó khác tính duy nhất của primary/candidate key.

Đối chiếu phương án: A. PRIMARY KEY; B. FOREIGN KEY; D. CHECK. Các lựa chọn này không đáp ứng định nghĩa/điều kiện được hỏi; khi ôn, hãy chỉ ra chính xác điểm lệch của từng lựa chọn thay vì học thuộc mỗi chữ cái.
\end{question}

\begin{question}{451}
% [ID gốc: 15.49]
% [Nguồn: DBI202 Order 15 Q49]
DB1202_SP26_RE | Multiple Choice Question 15 (Choose 1 answer) Domain constraints, functional dependency and referential integrity are special forms of
\begin{choices}
\choice Foreign key
\choice Primary key
\choice Assertion
\choice Referential constraint
\end{choices}
\textbf{Đáp án:} C
\textbf{Giải thích:} Đáp án theo dữ liệu đề: C — C. Assertion.

Lời giải gốc: Trong lý thuyết CSDL, "Assertion" (Khẳng định/Lời xác nhận) là một thuật ngữ chung nhất để định nghĩa các quy tắc ràng buộc trạng thái mà cơ sở dữ liệu phải luôn tuân thủ. Các ràng buộc toàn vẹn, miền giá trị hay khóa ngoại đều là dạng đặc biệt của assertion.

Nền tảng kiến thức: Chuẩn hóa dùng phụ thuộc hàm và khóa để giảm dư thừa cũng như bất thường cập nhật. Tính bao đóng thuộc tính để tìm khóa; 2NF loại phụ thuộc bộ phận vào một phần khóa ghép, 3NF xử lý phụ thuộc không khóa kiểu bắc cầu, còn BCNF yêu cầu mọi vế trái không tầm thường là siêu khóa.

Phân tích: Khóa ngoại buộc giá trị tham chiếu phải khớp một khóa được tham chiếu ở quan hệ đích (hoặc NULL nếu ràng buộc cho phép). Đây là toàn vẹn tham chiếu; nó khác tính duy nhất của primary/candidate key.

Đối chiếu phương án: A. Foreign key; B. Primary key; D. Referential constraint. Các lựa chọn này không đáp ứng định nghĩa/điều kiện được hỏi; khi ôn, hãy chỉ ra chính xác điểm lệch của từng lựa chọn thay vì học thuộc mỗi chữ cái.
\end{question}

\begin{question}{452}
% [ID gốc: 17.1]
% [Nguồn: DBI202 Order 17 Q1]
Choose a right answer
\begin{choices}
\choice When a privilege is granted, it cannot be revoked
\choice Privileges cannot be granted on a view
\choice An authorization ID may be granted privileges from others or may give its privileges to others
\choice All of the others
\end{choices}
\textbf{Đáp án:} C
\textbf{Giải thích:} Đáp án theo dữ liệu đề: C — C. An authorization ID may be granted privileges from others or may give its privileges to others.

Nền tảng kiến thức: Bảo mật CSDL dựa trên xác thực và phân quyền theo nguyên tắc quyền tối thiểu. GRANT cấp quyền; REVOKE thu hồi; WITH GRANT OPTION cho phép người nhận tiếp tục cấp quyền trong phạm vi được cho.

Phân tích: Với lựa chọn tổng hợp, không chọn theo cảm tính: kiểm tra từng phát biểu độc lập. “All of the others/above” chỉ đúng khi mọi phát biểu còn lại đúng; “None” chỉ đúng khi không phát biểu nào trong nhóm thỏa yêu cầu. Đáp án lưu cần được đọc cùng các mệnh đề thành phần.

Đối chiếu phương án: A. When a privilege is granted, it cannot be revoked; B. Privileges cannot be granted on a view; D. All of the others. Các lựa chọn này không đáp ứng định nghĩa/điều kiện được hỏi; khi ôn, hãy chỉ ra chính xác điểm lệch của từng lựa chọn thay vì học thuộc mỗi chữ cái.
\end{question}

\begin{question}{453}
% [ID gốc: 17.2]
% [Nguồn: DBI202 Order 17 Q2]
Choose the right statement to grant the INSERT and SELECT privileges on table Movies to users torn' and jerry'
\begin{choices}
\choice GRANT SELECT, INSERT on Movies
\choice GRANT SELECT, INSERT on Movies TO torn, jerry CASCADE
\choice GRANT SELECT, INSERT on Movies TO torn, jerry WITH GRANT OPTION
\choice GRANT SELECT, INSERT on Movies TO ALL
\end{choices}
\textbf{Đáp án:} C
\textbf{Giải thích:} Đáp án theo dữ liệu đề: C — C. GRANT SELECT, INSERT on Movies TO torn, jerry WITH GRANT OPTION.

Nền tảng kiến thức: Bảo mật CSDL dựa trên xác thực và phân quyền theo nguyên tắc quyền tối thiểu. GRANT cấp quyền; REVOKE thu hồi; WITH GRANT OPTION cho phép người nhận tiếp tục cấp quyền trong phạm vi được cho.

Phân tích: GRANT cấp một quyền xác định trên đối tượng cho principal/role. WITH GRANT OPTION mới cho phép người nhận chuyển tiếp quyền; quyền cấp không đồng nghĩa với quyền sở hữu bảng và có thể bị thu hồi theo cơ chế REVOKE.

Đối chiếu phương án: A. GRANT SELECT, INSERT on Movies; B. GRANT SELECT, INSERT on Movies TO torn, jerry CASCADE; D. GRANT SELECT, INSERT on Movies TO ALL. Các lựa chọn này không đáp ứng định nghĩa/điều kiện được hỏi; khi ôn, hãy chỉ ra chính xác điểm lệch của từng lựa chọn thay vì học thuộc mỗi chữ cái.
\end{question}


\newpage
\section*{Chương 7: Constraints and Triggers}

\begin{question}{454}
% [ID gốc: 1.25]
% [Nguồn: DBI202 Order 1 Q25]
Which of the following is the appropriate characteristic of a database?
\begin{choices}
\choice Because a database is created to suit the format of the data, it cannot respond flexibly to data format changes.
\choice The procedure for making backups is complicated.
\choice It is difficult to share data between operations due to an exclusive control function.
\choice It can be accessed by multiple users at the same time due to an exclusive control function.
\end{choices}
\textbf{Đáp án:} D
\textbf{Giải thích:} Đáp án theo dữ liệu đề: D — D. It can be accessed by multiple users at the same time due to an exclusive control function..

Nền tảng kiến thức: Chương này phân biệt dữ liệu, cơ sở dữ liệu, hệ quản trị CSDL (DBMS), mô hình dữ liệu, lược đồ và trạng thái dữ liệu. Cơ sở dữ liệu lưu dữ liệu có liên hệ; DBMS là phần mềm định nghĩa, truy vấn, cập nhật và bảo vệ dữ liệu đó.

Phân tích: Stored procedure là tập lệnh SQL có tên lưu trong CSDL, được gọi để thực thi một thủ tục; nó có thể nhận tham số và chứa nhiều câu lệnh. Đừng nhầm với trigger tự chạy khi có sự kiện hoặc function có quy tắc trả về giá trị.

Đối chiếu phương án: A. Because a database is created to suit the format of the data, it cannot respond flexibly to da…; B. The procedure for making backups is complicated.; C. It is difficult to share data between operations due to an exclusive control function.. Các lựa chọn này không đáp ứng định nghĩa/điều kiện được hỏi; khi ôn, hãy chỉ ra chính xác điểm lệch của từng lựa chọn thay vì học thuộc mỗi chữ cái.
\end{question}

\begin{question}{455}
% [ID gốc: 3.54]
% [Nguồn: DBI202 Order 3 Q54]
You need to develop an application to manage university courses, which requires structured data (tables, rows, columns), strict consistency, complex queries, and enforcement of ACID properties. Which data model is most appropriate?
\begin{choices}
\choice Relational data model (e.g. Microsoft SQL Server, MySQL)
\choice Document-based data model (e.g. MongoDB, Cassandra, CouchDB)
\choice Key-value data model (e.g. Redis, DynamoDB, Riak)
\choice Hierarchical data model
\end{choices}
\textbf{Đáp án:} A
\textbf{Giải thích:} Đáp án theo dữ liệu đề: A — A. Relational data model (e.g. Microsoft SQL Server, MySQL).

Lời giải gốc: Mô hình dữ liệu quan hệ phù hợp nhất cho các ứng dụng yêu cầu tính cấu trúc bảng, tính nhất quán nghiêm ngặt và hỗ trợ giao dịch ACID.

Nền tảng kiến thức: Mô hình quan hệ biểu diễn dữ liệu bằng quan hệ (bảng), bộ (hàng), thuộc tính (cột) và miền giá trị. Khóa nhận diện bộ; khóa ngoại và các ràng buộc toàn vẹn ngăn dữ liệu tham chiếu sai hoặc vi phạm miền.

Phân tích: Data model là bộ khái niệm dùng để mô tả cấu trúc dữ liệu, quan hệ giữa dữ liệu và các ràng buộc. Relational, hierarchical, network và graph là các mô hình có thật; cần chọn mô hình khớp đặc điểm được hỏi, không nhầm mô hình với DBMS hay ngôn ngữ thao tác.

Đối chiếu phương án: B. Document-based data model (e.g. MongoDB, Cassandra, CouchDB); C. Key-value data model (e.g. Redis, DynamoDB, Riak); D. Hierarchical data model. Các lựa chọn này không đáp ứng định nghĩa/điều kiện được hỏi; khi ôn, hãy chỉ ra chính xác điểm lệch của từng lựa chọn thay vì học thuộc mỗi chữ cái.
\end{question}

\begin{question}{456}
% [ID gốc: 5.1]
% [Nguồn: DBI202 Order 5 Q1]
What is an SQL virtual table that is constructed from other tables?
\begin{choices}
\choice Just another table
\choice A view
\choice A relation
\choice Query results
\end{choices}
\textbf{Đáp án:} B
\textbf{Giải thích:} Đáp án theo dữ liệu đề: B — B. A view.

Nền tảng kiến thức: SQL tách truy vấn dữ liệu (DML) khỏi định nghĩa cấu trúc (DDL) và quyền truy cập (DCL). Với SELECT, cần xác định bảng nguồn, điều kiện lọc, phép kết hợp, nhóm và điều kiện nhóm; các ràng buộc, view, thủ tục và trigger đảm nhiệm các vai trò khác nhau.

Phân tích: View là quan hệ ảo được định nghĩa bằng truy vấn trên một hay nhiều bảng/view; thông thường nó lưu định nghĩa truy vấn chứ không lưu bản sao kết quả như materialized view. Do đó đáp án phải chỉ đối tượng được tạo từ truy vấn, không nhầm với bảng vật lý.

Đối chiếu phương án: A. Just another table; C. A relation; D. Query results. Các lựa chọn này không đáp ứng định nghĩa/điều kiện được hỏi; khi ôn, hãy chỉ ra chính xác điểm lệch của từng lựa chọn thay vì học thuộc mỗi chữ cái.
\end{question}

\begin{question}{457}
% [ID gốc: 5.35]
% [Nguồn: DBI202 Order 5 Q35]
The SQL operations that are performed while a connection is active form a ____
\begin{choices}
\choice Session
\choice Catalog
\choice Module
\choice None of the others
\end{choices}
\textbf{Đáp án:} A
\textbf{Giải thích:} Đáp án theo dữ liệu đề: A — A. Session.

Nền tảng kiến thức: SQL tách truy vấn dữ liệu (DML) khỏi định nghĩa cấu trúc (DDL) và quyền truy cập (DCL). Với SELECT, cần xác định bảng nguồn, điều kiện lọc, phép kết hợp, nhóm và điều kiện nhóm; các ràng buộc, view, thủ tục và trigger đảm nhiệm các vai trò khác nhau.

Phân tích: Với lựa chọn tổng hợp, không chọn theo cảm tính: kiểm tra từng phát biểu độc lập. “All of the others/above” chỉ đúng khi mọi phát biểu còn lại đúng; “None” chỉ đúng khi không phát biểu nào trong nhóm thỏa yêu cầu. Đáp án lưu cần được đọc cùng các mệnh đề thành phần.

Đối chiếu phương án: B. Catalog; C. Module; D. None of the others. Các lựa chọn này không đáp ứng định nghĩa/điều kiện được hỏi; khi ôn, hãy chỉ ra chính xác điểm lệch của từng lựa chọn thay vì học thuộc mỗi chữ cái.
\end{question}

\begin{question}{458}
% [ID gốc: 5.37]
% [Nguồn: DBI202 Order 5 Q37]
A ____ is the framework under which data may exist and SQL operations on data may be executed
\begin{choices}
\choice SQL Environment
\choice SQL Statement
\choice SQL Parser
\choice SQL Optimizer
\end{choices}
\textbf{Đáp án:} A
\textbf{Giải thích:} Đáp án theo dữ liệu đề: A — A. SQL Environment.

Nền tảng kiến thức: SQL tách truy vấn dữ liệu (DML) khỏi định nghĩa cấu trúc (DDL) và quyền truy cập (DCL). Với SELECT, cần xác định bảng nguồn, điều kiện lọc, phép kết hợp, nhóm và điều kiện nhóm; các ràng buộc, view, thủ tục và trigger đảm nhiệm các vai trò khác nhau.

Phân tích: Trong T-SQL, biến cục bộ do người dùng khai báo bắt đầu bằng một dấu `@` (ví dụ `@count`); tên bắt đầu `@@` được dùng cho các hàm/biến hệ thống. Đừng nhầm tiền tố biến cục bộ với các biến hệ thống như `@@ROWCOUNT`.

Đối chiếu phương án: B. SQL Statement; C. SQL Parser; D. SQL Optimizer. Các lựa chọn này không đáp ứng định nghĩa/điều kiện được hỏi; khi ôn, hãy chỉ ra chính xác điểm lệch của từng lựa chọn thay vì học thuộc mỗi chữ cái.
\end{question}

\begin{question}{459}
% [ID gốc: 5.42]
% [Nguồn: DBI202 Order 5 Q42]
What is a view?
\begin{choices}
\choice A view is a database diagram
\choice A view is a special stored procedure executed when certain event occurs
\choice A view is a virtual table which results of executing a pre-compiled query. A view is not a part of the physical database schema, while the regular tables are.
\choice None of the others
\end{choices}
\textbf{Đáp án:} C
\textbf{Giải thích:} Đáp án theo dữ liệu đề: C — C. A view is a virtual table which results of executing a pre-compiled query. A view is not a part of the physical database schema, while the regular tables are..

Nền tảng kiến thức: SQL tách truy vấn dữ liệu (DML) khỏi định nghĩa cấu trúc (DDL) và quyền truy cập (DCL). Với SELECT, cần xác định bảng nguồn, điều kiện lọc, phép kết hợp, nhóm và điều kiện nhóm; các ràng buộc, view, thủ tục và trigger đảm nhiệm các vai trò khác nhau.

Phân tích: View là quan hệ ảo được định nghĩa bằng truy vấn trên một hay nhiều bảng/view; thông thường nó lưu định nghĩa truy vấn chứ không lưu bản sao kết quả như materialized view. Do đó đáp án phải chỉ đối tượng được tạo từ truy vấn, không nhầm với bảng vật lý.

Đối chiếu phương án: A. A view is a database diagram; B. A view is a special stored procedure executed when certain event occurs; D. None of the others. Các lựa chọn này không đáp ứng định nghĩa/điều kiện được hỏi; khi ôn, hãy chỉ ra chính xác điểm lệch của từng lựa chọn thay vì học thuộc mỗi chữ cái.
\end{question}

\begin{question}{460}
% [ID gốc: 5.43]
% [Nguồn: DBI202 Order 5 Q43]
An updatable view is a view that allows data modification through INSERT, UPDATE, and DELETE statements. When an update is made to an updatable view, the changes are translated into the underlying base table.
Suppose an updatable view ParamountMovies is associated with Movies relation. Choose a correct answer.
\begin{choices}
\choice Drop Movies relation also delete the view ParamountMovies
\choice Drop ParamountMovies also delete Movies
\choice An update on ParamountMovies is translated into Movies
\choice None of the others
\end{choices}
\textbf{Đáp án:} C
\textbf{Giải thích:} Đáp án theo dữ liệu đề: C — C. An update on ParamountMovies is translated into Movies.

Nền tảng kiến thức: SQL tách truy vấn dữ liệu (DML) khỏi định nghĩa cấu trúc (DDL) và quyền truy cập (DCL). Với SELECT, cần xác định bảng nguồn, điều kiện lọc, phép kết hợp, nhóm và điều kiện nhóm; các ràng buộc, view, thủ tục và trigger đảm nhiệm các vai trò khác nhau.

Phân tích: View là quan hệ ảo được định nghĩa bằng truy vấn trên một hay nhiều bảng/view; thông thường nó lưu định nghĩa truy vấn chứ không lưu bản sao kết quả như materialized view. Do đó đáp án phải chỉ đối tượng được tạo từ truy vấn, không nhầm với bảng vật lý.

Đối chiếu phương án: A. Drop Movies relation also delete the view ParamountMovies; B. Drop ParamountMovies also delete Movies; D. None of the others. Các lựa chọn này không đáp ứng định nghĩa/điều kiện được hỏi; khi ôn, hãy chỉ ra chính xác điểm lệch của từng lựa chọn thay vì học thuộc mỗi chữ cái.
\end{question}

\begin{question}{461}
% [ID gốc: 5.56]
% [Nguồn: DBI202 Order 5 Q56]
Choose a right answer.
\begin{choices}
\choice A SOL environment is the framework under which data may exist and SQL operations on data may be executed
\choice All of the others
\choice Within a SQL environment are two special kinds of processes: SQL clients and SQL servers.
\choice A SQL environment is a DBMS running at some installation.
\end{choices}
\textbf{Đáp án:} B
\textbf{Giải thích:} Đáp án theo dữ liệu đề: B — B. All of the others.

Nền tảng kiến thức: SQL tách truy vấn dữ liệu (DML) khỏi định nghĩa cấu trúc (DDL) và quyền truy cập (DCL). Với SELECT, cần xác định bảng nguồn, điều kiện lọc, phép kết hợp, nhóm và điều kiện nhóm; các ràng buộc, view, thủ tục và trigger đảm nhiệm các vai trò khác nhau.

Phân tích: Với lựa chọn tổng hợp, không chọn theo cảm tính: kiểm tra từng phát biểu độc lập. “All of the others/above” chỉ đúng khi mọi phát biểu còn lại đúng; “None” chỉ đúng khi không phát biểu nào trong nhóm thỏa yêu cầu. Đáp án lưu cần được đọc cùng các mệnh đề thành phần.

Đối chiếu phương án: A. A SOL environment is the framework under which data may exist and SQL operations on data may b…; C. Within a SQL environment are two special kinds of processes: SQL clients and SQL servers.; D. A SQL environment is a DBMS running at some installation.. Các lựa chọn này không đáp ứng định nghĩa/điều kiện được hỏi; khi ôn, hãy chỉ ra chính xác điểm lệch của từng lựa chọn thay vì học thuộc mỗi chữ cái.
\end{question}

\begin{question}{462}
% [ID gốc: 5.96]
% [Nguồn: DBI202 Order 5 Q96]
Which one of the following is not true for a view?
\begin{choices}
\choice View never contains derived columns.
\choice A view definition is permanently stored as part of the database.
\choice View is a virtual table.
\choice View is derived from other tables
\end{choices}
\textbf{Đáp án:} A
\textbf{Giải thích:} Đáp án theo dữ liệu đề: A — A. View never contains derived columns..

Nền tảng kiến thức: SQL tách truy vấn dữ liệu (DML) khỏi định nghĩa cấu trúc (DDL) và quyền truy cập (DCL). Với SELECT, cần xác định bảng nguồn, điều kiện lọc, phép kết hợp, nhóm và điều kiện nhóm; các ràng buộc, view, thủ tục và trigger đảm nhiệm các vai trò khác nhau.

Phân tích: View là quan hệ ảo được định nghĩa bằng truy vấn trên một hay nhiều bảng/view; thông thường nó lưu định nghĩa truy vấn chứ không lưu bản sao kết quả như materialized view. Do đó đáp án phải chỉ đối tượng được tạo từ truy vấn, không nhầm với bảng vật lý.

Đối chiếu phương án: B. A view definition is permanently stored as part of the database.; C. View is a virtual table.; D. View is derived from other tables. Các lựa chọn này không đáp ứng định nghĩa/điều kiện được hỏi; khi ôn, hãy chỉ ra chính xác điểm lệch của từng lựa chọn thay vì học thuộc mỗi chữ cái.
\end{question}

\begin{question}{463}
% [ID gốc: 5.97]
% [Nguồn: DBI202 Order 5 Q97]
Which is the SQL statement that creates an updatable view? In this question, all the tables in the SQL
statements are updatable.
\begin{choices}
\choice CREATE VIEW Ordered_Product(Product_No) AS SELECT DISTINCT Product_No FROM Order
\choice CREATE VIEW Order_List(Order_No, Product_Name, Order_Quantity) AS SELECT Order_No,Product_Name, Order_Quantity FROM Order, Product WHERE Order Product_No=Product.Product_No
\choice CREATE VIEW Product_Order(Product_No, Order_Quantity) AS SELECT Product_No,SUM(Order_Quantity) FROM Order GROUP BY Product_No
\choice CREATE VIEW Expensive_Product(Product_No, Product_Name) AS SELECT Product_No,Product_Name FROM Product WHERE Product_Unit_Price>1000
\end{choices}
\textbf{Đáp án:} D
\textbf{Giải thích:} Đáp án theo dữ liệu đề: D — D. CREATE VIEW Expensive_Product(Product_No, Product_Name) AS SELECT Product_No,Product_Name FROM Product WHERE Product_Unit_Price>1000.

Nền tảng kiến thức: SQL tách truy vấn dữ liệu (DML) khỏi định nghĩa cấu trúc (DDL) và quyền truy cập (DCL). Với SELECT, cần xác định bảng nguồn, điều kiện lọc, phép kết hợp, nhóm và điều kiện nhóm; các ràng buộc, view, thủ tục và trigger đảm nhiệm các vai trò khác nhau.

Phân tích: View là quan hệ ảo được định nghĩa bằng truy vấn trên một hay nhiều bảng/view; thông thường nó lưu định nghĩa truy vấn chứ không lưu bản sao kết quả như materialized view. Do đó đáp án phải chỉ đối tượng được tạo từ truy vấn, không nhầm với bảng vật lý.

Đối chiếu phương án: A. CREATE VIEW Ordered_Product(Product_No) AS SELECT DISTINCT Product_No FROM Order; B. CREATE VIEW Order_List(Order_No, Product_Name, Order_Quantity) AS SELECT Order_No,Product_Name…; C. CREATE VIEW Product_Order(Product_No, Order_Quantity) AS SELECT Product_No,SUM(Order_Quantity)…. Các lựa chọn này không đáp ứng định nghĩa/điều kiện được hỏi; khi ôn, hãy chỉ ra chính xác điểm lệch của từng lựa chọn thay vì học thuộc mỗi chữ cái.
\end{question}

\begin{question}{464}
% [ID gốc: 5.112]
% [Nguồn: DBI202 Order 5 Q112]
Your database contains a table named Purchases. The table includes a DATETIME column named PurchaseTime that stores the date and time each purchase is made. There is a non-clustered index on the Purchase Time column. The business team wants a report that displays the total number of purchases made on the current day. You need to write a query that will return the correct results in the most efficient manner.Which Transact-SQL query should you use?
\begin{choices}
\choice SELECT COUNT(*) FROM Purchases WHERE PurchaseTime = CONVERT(DATE, GETDATE())
\choice SELECT COUNT(*) FROM Purchases WHERE PurchaseTime = GETDATE()
\choice SELECT COUNT(*) FROM Purchases WHERE CONVERT(VARCHAR, PurchaseTime, 112) = CONVERT(VARCHAR, GETDATE(), 112)
\choice SELECT COUNT(*) FROM Purchases WHERE PurchaseTime >= CONVERT(DATE, GETDATE()) AND PurchaseTime < DATEADD(DAY, 1, CONVERT(DATE, GETDATE())
\end{choices}
\textbf{Đáp án:} D
\textbf{Giải thích:} Đáp án theo dữ liệu đề: D — D. SELECT COUNT(*) FROM Purchases WHERE PurchaseTime >= CONVERT(DATE, GETDATE()) AND PurchaseTime < DATEADD(DAY, 1, CONVERT(DATE, GETDATE()).

Nền tảng kiến thức: SQL tách truy vấn dữ liệu (DML) khỏi định nghĩa cấu trúc (DDL) và quyền truy cập (DCL). Với SELECT, cần xác định bảng nguồn, điều kiện lọc, phép kết hợp, nhóm và điều kiện nhóm; các ràng buộc, view, thủ tục và trigger đảm nhiệm các vai trò khác nhau.

Phân tích: Trong T-SQL, biến cục bộ do người dùng khai báo bắt đầu bằng một dấu `@` (ví dụ `@count`); tên bắt đầu `@@` được dùng cho các hàm/biến hệ thống. Đừng nhầm tiền tố biến cục bộ với các biến hệ thống như `@@ROWCOUNT`.

Đối chiếu phương án: A. SELECT COUNT(*) FROM Purchases WHERE PurchaseTime = CONVERT(DATE, GETDATE()); B. SELECT COUNT(*) FROM Purchases WHERE PurchaseTime = GETDATE(); C. SELECT COUNT(*) FROM Purchases WHERE CONVERT(VARCHAR, PurchaseTime, 112) = CONVERT(VARCHAR, GE…. Các lựa chọn này không đáp ứng định nghĩa/điều kiện được hỏi; khi ôn, hãy chỉ ra chính xác điểm lệch của từng lựa chọn thay vì học thuộc mỗi chữ cái.
\end{question}

\begin{question}{465}
% [ID gốc: 5.114]
% [Nguồn: DBI202 Order 5 Q114]
Through VIEW, is it possible to update data in the original table (that VIEW sees)?
\begin{choices}
\choice Always updateable
\choice Can never update
\choice Some cases can be updated, some cases cannot
\choice If the View sees a table, it can be updated
\choice All of the mentioned answers are wrong
\end{choices}
\textbf{Đáp án:} C
\textbf{Giải thích:} Đáp án theo dữ liệu đề: C — C. Some cases can be updated, some cases cannot.

Nền tảng kiến thức: SQL tách truy vấn dữ liệu (DML) khỏi định nghĩa cấu trúc (DDL) và quyền truy cập (DCL). Với SELECT, cần xác định bảng nguồn, điều kiện lọc, phép kết hợp, nhóm và điều kiện nhóm; các ràng buộc, view, thủ tục và trigger đảm nhiệm các vai trò khác nhau.

Phân tích: View là quan hệ ảo được định nghĩa bằng truy vấn trên một hay nhiều bảng/view; thông thường nó lưu định nghĩa truy vấn chứ không lưu bản sao kết quả như materialized view. Do đó đáp án phải chỉ đối tượng được tạo từ truy vấn, không nhầm với bảng vật lý.

Đối chiếu phương án: A. Always updateable; B. Can never update; D. If the View sees a table, it can be updated; E. All of the mentioned answers are wrong. Các lựa chọn này không đáp ứng định nghĩa/điều kiện được hỏi; khi ôn, hãy chỉ ra chính xác điểm lệch của từng lựa chọn thay vì học thuộc mỗi chữ cái.
\end{question}

\begin{question}{466}
% [ID gốc: 5.115]
% [Nguồn: DBI202 Order 5 Q115]
What is the difference between a view and a table in SQL?
\begin{choices}
\choice Tables can be updated, while the view is displaying data from the query
\choice Views can store data permanently, while tables are temporary
\choice Tables are created with the CREATE TABLE statement, while views use CREATE VIEW
\choice Views and tables are functionally equivalent in SQL
\end{choices}
\textbf{Đáp án:} C
\textbf{Giải thích:} Đáp án theo dữ liệu đề: C — C. Tables are created with the CREATE TABLE statement, while views use CREATE VIEW.

Nền tảng kiến thức: SQL tách truy vấn dữ liệu (DML) khỏi định nghĩa cấu trúc (DDL) và quyền truy cập (DCL). Với SELECT, cần xác định bảng nguồn, điều kiện lọc, phép kết hợp, nhóm và điều kiện nhóm; các ràng buộc, view, thủ tục và trigger đảm nhiệm các vai trò khác nhau.

Phân tích: Phép tập hợp trên quan hệ cần so sánh các tuple theo thuộc tính tương ứng; union lấy hợp, intersection lấy phần chung và difference lấy phần có ở vế trái mà không có ở vế phải. Trước tiên kiểm tra tính tương thích schema, sau đó áp phép theo đúng hướng/toán tử.

Đối chiếu phương án: A. Tables can be updated, while the view is displaying data from the query; B. Views can store data permanently, while tables are temporary; D. Views and tables are functionally equivalent in SQL. Các lựa chọn này không đáp ứng định nghĩa/điều kiện được hỏi; khi ôn, hãy chỉ ra chính xác điểm lệch của từng lựa chọn thay vì học thuộc mỗi chữ cái.
\end{question}

\begin{question}{467}
% [ID gốc: 5.126]
% [Nguồn: DBI202 Order 5 Q126]
In SQL, what is a view?
\begin{choices}
\choice A virtual table based on the result of a SELECT query
\choice A physical table containing aggregated data
\choice A temporary table used for transactions
\choice A stored procedure with read-only access
\end{choices}
\textbf{Đáp án:} A
\textbf{Giải thích:} Đáp án theo dữ liệu đề: A — A. A virtual table based on the result of a SELECT query.

Nền tảng kiến thức: SQL tách truy vấn dữ liệu (DML) khỏi định nghĩa cấu trúc (DDL) và quyền truy cập (DCL). Với SELECT, cần xác định bảng nguồn, điều kiện lọc, phép kết hợp, nhóm và điều kiện nhóm; các ràng buộc, view, thủ tục và trigger đảm nhiệm các vai trò khác nhau.

Phân tích: View là quan hệ ảo được định nghĩa bằng truy vấn trên một hay nhiều bảng/view; thông thường nó lưu định nghĩa truy vấn chứ không lưu bản sao kết quả như materialized view. Do đó đáp án phải chỉ đối tượng được tạo từ truy vấn, không nhầm với bảng vật lý.

Đối chiếu phương án: B. A physical table containing aggregated data; C. A temporary table used for transactions; D. A stored procedure with read-only access. Các lựa chọn này không đáp ứng định nghĩa/điều kiện được hỏi; khi ôn, hãy chỉ ra chính xác điểm lệch của từng lựa chọn thay vì học thuộc mỗi chữ cái.
\end{question}

\begin{question}{468}
% [ID gốc: 5.137]
% [Nguồn: DBI202 Order 5 Q137]
A view's definition can be modified using which statement, while a view's data can be modified using which statement?
\begin{choices}
\choice ALTER VIEW, ALTER
\choice ALTER VIEW, EDIT
\choice ALTER VIEW, UPDATE
\choice UPDATE, ALTER VIEW
\end{choices}
\textbf{Đáp án:} C
\textbf{Giải thích:} Đáp án theo dữ liệu đề: C — C. ALTER VIEW, UPDATE.

Lời giải gốc: Cấu trúc/định nghĩa của một view được chỉnh sửa bằng lệnh ALTER VIEW, trong khi dữ liệu ở bảng cơ sở thông qua view có thể được thao tác bằng câu lệnh UPDATE.

Nền tảng kiến thức: SQL tách truy vấn dữ liệu (DML) khỏi định nghĩa cấu trúc (DDL) và quyền truy cập (DCL). Với SELECT, cần xác định bảng nguồn, điều kiện lọc, phép kết hợp, nhóm và điều kiện nhóm; các ràng buộc, view, thủ tục và trigger đảm nhiệm các vai trò khác nhau.

Phân tích: View là quan hệ ảo được định nghĩa bằng truy vấn trên một hay nhiều bảng/view; thông thường nó lưu định nghĩa truy vấn chứ không lưu bản sao kết quả như materialized view. Do đó đáp án phải chỉ đối tượng được tạo từ truy vấn, không nhầm với bảng vật lý.

Đối chiếu phương án: A. ALTER VIEW, ALTER; B. ALTER VIEW, EDIT; D. UPDATE, ALTER VIEW. Các lựa chọn này không đáp ứng định nghĩa/điều kiện được hỏi; khi ôn, hãy chỉ ra chính xác điểm lệch của từng lựa chọn thay vì học thuộc mỗi chữ cái.
\end{question}

\begin{question}{469}
% [ID gốc: 5.170]
% [Nguồn: DBI202 Order 5 Q170]
Which SQL statement creates an updatable view?
\begin{choices}
\choice CREATE VIEW Ordered_Product(Product_No) AS SELECT DISTINCT Product_No FROM "Order"
\choice CREATE VIEW Order_List(Order_No, Product_Name, Order_Quantity) AS SELECT Order_No, Product_Name, Order_Quantity FROM "Order", Product WHERE "Order".Product_No = Product.Product_No
\choice CREATE VIEW Product_Order(Product_No, Order_Quantity) AS SELECT Product_No, SUM(Order_Quantity) FROM "Order" GROUP BY Product_No
\choice CREATE VIEW Expensive_Product(Product_No, Product_Name) AS SELECT Product_No, Product_Name FROM Product WHERE Product_Unit_Price &gt; 1000
\end{choices}
\textbf{Đáp án:} D
\textbf{Giải thích:} Đáp án theo dữ liệu đề: D — D. CREATE VIEW Expensive_Product(Product_No, Product_Name) AS SELECT Product_No, Product_Name FROM Product WHERE Product_Unit_Price &gt; 1000.

Lời giải gốc: View có thể cập nhật được thường là view tham chiếu trực tiếp từ một bảng cơ sở duy nhất, không chứa DISTINCT, hàm gộp nhóm hay GROUP BY. Chỉ có phương án D thỏa mãn.

Nền tảng kiến thức: SQL tách truy vấn dữ liệu (DML) khỏi định nghĩa cấu trúc (DDL) và quyền truy cập (DCL). Với SELECT, cần xác định bảng nguồn, điều kiện lọc, phép kết hợp, nhóm và điều kiện nhóm; các ràng buộc, view, thủ tục và trigger đảm nhiệm các vai trò khác nhau.

Phân tích: View là quan hệ ảo được định nghĩa bằng truy vấn trên một hay nhiều bảng/view; thông thường nó lưu định nghĩa truy vấn chứ không lưu bản sao kết quả như materialized view. Do đó đáp án phải chỉ đối tượng được tạo từ truy vấn, không nhầm với bảng vật lý.

Đối chiếu phương án: A. CREATE VIEW Ordered_Product(Product_No) AS SELECT DISTINCT Product_No FROM "Order"; B. CREATE VIEW Order_List(Order_No, Product_Name, Order_Quantity) AS SELECT Order_No, Product_Nam…; C. CREATE VIEW Product_Order(Product_No, Order_Quantity) AS SELECT Product_No, SUM(Order_Quantity…. Các lựa chọn này không đáp ứng định nghĩa/điều kiện được hỏi; khi ôn, hãy chỉ ra chính xác điểm lệch của từng lựa chọn thay vì học thuộc mỗi chữ cái.
\end{question}

\begin{question}{470}
% [ID gốc: 5.171]
% [Nguồn: DBI202 Order 5 Q171]
In SQL, a view is a virtual table based on the result-set of an SQL statement. Which statement is NOT correct about a view?
\begin{choices}
\choice The UPDATE command is applicable to all views.
\choice The fields in a view are fields from one or more real tables in the database.
\choice The database engine recreates the view every time a user queries it.
\choice A view always shows up-to-date data.
\end{choices}
\textbf{Đáp án:} A
\textbf{Giải thích:} Đáp án theo dữ liệu đề: A — A. The UPDATE command is applicable to all views..

Lời giải gốc: Không phải mọi view đều có thể cập nhật (UPDATE). Các view chứa phép kết phức tạp, hàm gom nhóm hoặc DISTINCT thường không cho phép cập nhật trực tiếp.

Nền tảng kiến thức: SQL tách truy vấn dữ liệu (DML) khỏi định nghĩa cấu trúc (DDL) và quyền truy cập (DCL). Với SELECT, cần xác định bảng nguồn, điều kiện lọc, phép kết hợp, nhóm và điều kiện nhóm; các ràng buộc, view, thủ tục và trigger đảm nhiệm các vai trò khác nhau.

Phân tích: View là quan hệ ảo được định nghĩa bằng truy vấn trên một hay nhiều bảng/view; thông thường nó lưu định nghĩa truy vấn chứ không lưu bản sao kết quả như materialized view. Do đó đáp án phải chỉ đối tượng được tạo từ truy vấn, không nhầm với bảng vật lý.

Đối chiếu phương án: B. The fields in a view are fields from one or more real tables in the database.; C. The database engine recreates the view every time a user queries it.; D. A view always shows up-to-date data.. Các lựa chọn này không đáp ứng định nghĩa/điều kiện được hỏi; khi ôn, hãy chỉ ra chính xác điểm lệch của từng lựa chọn thay vì học thuộc mỗi chữ cái.
\end{question}

\begin{question}{471}
% [ID gốc: 5.246]
% [Nguồn: DBI202 Order 5 Q246]
Suppose we have defined VIEW ParamountMovies from the Movies table. When we execute the command: DROP TABLE Movies Which of the following statements is correct?
\begin{choices}
\choice The Movies table goes away, but VIEW ParamountMovies can still be used.
\choice The Movies table goes away and VIEW ParamountMovies cannot be used.
\choice The Movies table cannot be dropped because VIEW ParamountMovies still references it.
\choice The Movies table cannot be dropped, but VIEW ParamountMovies is dropped.
\end{choices}
\textbf{Đáp án:} B
\textbf{Giải thích:} Đáp án theo dữ liệu đề: B — B. The Movies table goes away and VIEW ParamountMovies cannot be used..

Lời giải gốc: Bảng Movies sẽ bị xóa. Do View chỉ là bảng ảo tham chiếu tới dữ liệu bảng gốc, nên khi bảng gốc bị xóa, ParamountMovies sẽ bị lỗi không sử dụng được nữa (Invalid/Broken View).

Nền tảng kiến thức: SQL tách truy vấn dữ liệu (DML) khỏi định nghĩa cấu trúc (DDL) và quyền truy cập (DCL). Với SELECT, cần xác định bảng nguồn, điều kiện lọc, phép kết hợp, nhóm và điều kiện nhóm; các ràng buộc, view, thủ tục và trigger đảm nhiệm các vai trò khác nhau.

Phân tích: View là quan hệ ảo được định nghĩa bằng truy vấn trên một hay nhiều bảng/view; thông thường nó lưu định nghĩa truy vấn chứ không lưu bản sao kết quả như materialized view. Do đó đáp án phải chỉ đối tượng được tạo từ truy vấn, không nhầm với bảng vật lý.

Đối chiếu phương án: A. The Movies table goes away, but VIEW ParamountMovies can still be used.; C. The Movies table cannot be dropped because VIEW ParamountMovies still references it.; D. The Movies table cannot be dropped, but VIEW ParamountMovies is dropped.. Các lựa chọn này không đáp ứng định nghĩa/điều kiện được hỏi; khi ôn, hãy chỉ ra chính xác điểm lệch của từng lựa chọn thay vì học thuộc mỗi chữ cái.
\end{question}

\begin{question}{472}
% [ID gốc: 5.264]
% [Nguồn: DBI202 Order 5 Q264]
Through a VIEW, is it possible to update data in the original table (that the VIEW sees)?
\begin{choices}
\choice Always updateable
\choice Can never update
\choice Some cases can be updated, some cases cannot
\choice If the VIEW sees a table, it can be updated
\choice All of the mentioned answers are wrong
\end{choices}
\textbf{Đáp án:} C
\textbf{Giải thích:} Đáp án theo dữ liệu đề: C — C. Some cases can be updated, some cases cannot.

Lời giải gốc: Dữ liệu qua VIEW có thể cập nhật được trong một số trường hợp đơn giản, nhưng không thể cập nhật nếu VIEW chứa các phép tính toán hoặc gom nhóm phức tạp.

Nền tảng kiến thức: SQL tách truy vấn dữ liệu (DML) khỏi định nghĩa cấu trúc (DDL) và quyền truy cập (DCL). Với SELECT, cần xác định bảng nguồn, điều kiện lọc, phép kết hợp, nhóm và điều kiện nhóm; các ràng buộc, view, thủ tục và trigger đảm nhiệm các vai trò khác nhau.

Phân tích: View là quan hệ ảo được định nghĩa bằng truy vấn trên một hay nhiều bảng/view; thông thường nó lưu định nghĩa truy vấn chứ không lưu bản sao kết quả như materialized view. Do đó đáp án phải chỉ đối tượng được tạo từ truy vấn, không nhầm với bảng vật lý.

Đối chiếu phương án: A. Always updateable; B. Can never update; D. If the VIEW sees a table, it can be updated; E. All of the mentioned answers are wrong. Các lựa chọn này không đáp ứng định nghĩa/điều kiện được hỏi; khi ôn, hãy chỉ ra chính xác điểm lệch của từng lựa chọn thay vì học thuộc mỗi chữ cái.
\end{question}

\begin{question}{473}
% [ID gốc: 7.1]
% [Nguồn: DBI202 Order 7 Q1]
Which files are created automatically by SQL Server when you create a new database?
\begin{choices}
\choice Primary data file and secondary data file
\choice Primary data file and log file
\choice Index file and log file
\choice Secondary data file and log file
\choice Backup file and metadata file
\end{choices}
\textbf{Đáp án:} B
\textbf{Giải thích:} Đáp án theo dữ liệu đề: B — B. Primary data file and log file.

Lời giải gốc: Khi tạo một cơ sở dữ liệu mới trong SQL Server, hệ thống luôn tự động tạo ít nhất một tệp dữ liệu chính (Primary data file - .mdf) và một tệp nhật ký giao dịch (Log file - .ldf).

Nền tảng kiến thức: DBMS lưu dữ liệu lâu dài trên thiết bị lưu trữ theo trang/khối; thao tác đọc ghi đĩa thường đắt hơn xử lý trong bộ nhớ. Khi phân tích câu hỏi, phân biệt tệp dữ liệu với log và tập trung vào số lần I/O.

Phân tích: Index lưu cấu trúc tra cứu phụ để tìm vị trí bản ghi nhanh hơn thay vì quét toàn bộ bảng. Bù lại, nó dùng thêm dung lượng và phải được duy trì khi INSERT/UPDATE/DELETE; vì thế “càng nhiều index càng tốt” không đúng cho mọi workload.

Đối chiếu phương án: A. Primary data file and secondary data file; C. Index file and log file; D. Secondary data file and log file; E. Backup file and metadata file. Các lựa chọn này không đáp ứng định nghĩa/điều kiện được hỏi; khi ôn, hãy chỉ ra chính xác điểm lệch của từng lựa chọn thay vì học thuộc mỗi chữ cái.
\end{question}

\begin{question}{474}
% [ID gốc: 8.1]
% [Nguồn: DBI202 Order 8 Q1]
The most useful index on a relation is an index on its key. This is because:
\begin{choices}
\choice The index on non key attribute(s) runs slower
\choice The search operation based on the primary key is commonly used.
\choice The index on non key attribute(s) makes update operations to the relation more complex and time-consuming
\choice All of the others
\end{choices}
\textbf{Đáp án:} B
\textbf{Giải thích:} Đáp án theo dữ liệu đề: B — B. The search operation based on the primary key is commonly used..

Nền tảng kiến thức: Tổ chức tệp quyết định cách bản ghi nằm trên trang; chỉ mục ánh xạ giá trị khóa tìm kiếm tới vị trí bản ghi để giảm số trang phải đọc. Chỉ mục tăng tốc một số truy vấn nhưng tiêu tốn dung lượng và cần cập nhật khi dữ liệu thay đổi.

Phân tích: Index lưu cấu trúc tra cứu phụ để tìm vị trí bản ghi nhanh hơn thay vì quét toàn bộ bảng. Bù lại, nó dùng thêm dung lượng và phải được duy trì khi INSERT/UPDATE/DELETE; vì thế “càng nhiều index càng tốt” không đúng cho mọi workload.

Đối chiếu phương án: A. The index on non key attribute(s) runs slower; C. The index on non key attribute(s) makes update operations to the relation more complex and tim…; D. All of the others. Các lựa chọn này không đáp ứng định nghĩa/điều kiện được hỏi; khi ôn, hãy chỉ ra chính xác điểm lệch của từng lựa chọn thay vì học thuộc mỗi chữ cái.
\end{question}

\begin{question}{475}
% [ID gốc: 8.2]
% [Nguồn: DBI202 Order 8 Q2]
What structure can you implement for the database to speed up table reads?
\begin{choices}
\choice indexes
\choice triggers
\choice store-procedures
\choice synonyms
\end{choices}
\textbf{Đáp án:} A
\textbf{Giải thích:} Đáp án theo dữ liệu đề: A — A. indexes.

Nền tảng kiến thức: Tổ chức tệp quyết định cách bản ghi nằm trên trang; chỉ mục ánh xạ giá trị khóa tìm kiếm tới vị trí bản ghi để giảm số trang phải đọc. Chỉ mục tăng tốc một số truy vấn nhưng tiêu tốn dung lượng và cần cập nhật khi dữ liệu thay đổi.

Phân tích: Index lưu cấu trúc tra cứu phụ để tìm vị trí bản ghi nhanh hơn thay vì quét toàn bộ bảng. Bù lại, nó dùng thêm dung lượng và phải được duy trì khi INSERT/UPDATE/DELETE; vì thế “càng nhiều index càng tốt” không đúng cho mọi workload.

Đối chiếu phương án: B. triggers; C. store-procedures; D. synonyms. Các lựa chọn này không đáp ứng định nghĩa/điều kiện được hỏi; khi ôn, hãy chỉ ra chính xác điểm lệch của từng lựa chọn thay vì học thuộc mỗi chữ cái.
\end{question}

\begin{question}{476}
% [ID gốc: 8.3]
% [Nguồn: DBI202 Order 8 Q3]
Often, a DBMS allows us to build a single index on multiple attributes. This type of index takes values for several attributes and efficiently finds the tuples with the given values for these attributes. (P.343)
Often, a DBMS allows us to build a single index on ____ attribute(s)
(choose the most correct answer)
\begin{choices}
\choice only one
\choice two
\choice three
\choice multiple
\end{choices}
\textbf{Đáp án:} D
\textbf{Giải thích:} Đáp án theo dữ liệu đề: D — D. multiple.

Nền tảng kiến thức: Tổ chức tệp quyết định cách bản ghi nằm trên trang; chỉ mục ánh xạ giá trị khóa tìm kiếm tới vị trí bản ghi để giảm số trang phải đọc. Chỉ mục tăng tốc một số truy vấn nhưng tiêu tốn dung lượng và cần cập nhật khi dữ liệu thay đổi.

Phân tích: Các hàng dữ liệu của một relation được gọi là tuple; header chứa tên thuộc tính không phải tuple dữ liệu. Trong bảng, tuple tương ứng một bản ghi, còn attribute tương ứng cột.

Đối chiếu phương án: A. only one; B. two; C. three. Các lựa chọn này không đáp ứng định nghĩa/điều kiện được hỏi; khi ôn, hãy chỉ ra chính xác điểm lệch của từng lựa chọn thay vì học thuộc mỗi chữ cái.
\end{question}

\begin{question}{477}
% [ID gốc: 8.4]
% [Nguồn: DBI202 Order 8 Q4]
Which SQL statement is used to remove an existing index from a table?
\begin{choices}
\choice Atomicity, Consistency, Isolation, Durability
\choice REMOVE INDEX
\choice DELETE INDEX
\choice UNINDEX TABLE
\end{choices}
\textbf{Đáp án:} A
\textbf{Giải thích:} Đáp án theo dữ liệu đề: A — A. Atomicity, Consistency, Isolation, Durability.

Nền tảng kiến thức: Tổ chức tệp quyết định cách bản ghi nằm trên trang; chỉ mục ánh xạ giá trị khóa tìm kiếm tới vị trí bản ghi để giảm số trang phải đọc. Chỉ mục tăng tốc một số truy vấn nhưng tiêu tốn dung lượng và cần cập nhật khi dữ liệu thay đổi.

Phân tích: DROP TABLE loại bỏ định nghĩa bảng cùng dữ liệu phụ thuộc theo quy tắc DBMS; DELETE xóa các hàng nhưng giữ bảng, còn TRUNCATE xóa nhanh toàn bộ hàng trong khi vẫn giữ cấu trúc. Cần chọn lệnh đúng cấp độ tác động.

Đối chiếu phương án: B. REMOVE INDEX; C. DELETE INDEX; D. UNINDEX TABLE. Các lựa chọn này không đáp ứng định nghĩa/điều kiện được hỏi; khi ôn, hãy chỉ ra chính xác điểm lệch của từng lựa chọn thay vì học thuộc mỗi chữ cái.
\end{question}

\begin{question}{478}
% [ID gốc: 8.5]
% [Nguồn: DBI202 Order 8 Q5]
What is not the purpose of the index in MS-SQL Server?
\begin{choices}
\choice To use less hardware resources
\choice To enhance the query performance
\choice To provide an index to a record
\choice To perform fast in searching
\end{choices}
\textbf{Đáp án:} A
\textbf{Giải thích:} Đáp án theo dữ liệu đề: A — A. To use less hardware resources.

Nền tảng kiến thức: Tổ chức tệp quyết định cách bản ghi nằm trên trang; chỉ mục ánh xạ giá trị khóa tìm kiếm tới vị trí bản ghi để giảm số trang phải đọc. Chỉ mục tăng tốc một số truy vấn nhưng tiêu tốn dung lượng và cần cập nhật khi dữ liệu thay đổi.

Phân tích: Index lưu cấu trúc tra cứu phụ để tìm vị trí bản ghi nhanh hơn thay vì quét toàn bộ bảng. Bù lại, nó dùng thêm dung lượng và phải được duy trì khi INSERT/UPDATE/DELETE; vì thế “càng nhiều index càng tốt” không đúng cho mọi workload.

Đối chiếu phương án: B. To enhance the query performance; C. To provide an index to a record; D. To perform fast in searching. Các lựa chọn này không đáp ứng định nghĩa/điều kiện được hỏi; khi ôn, hãy chỉ ra chính xác điểm lệch của từng lựa chọn thay vì học thuộc mỗi chữ cái.
\end{question}

\begin{question}{479}
% [ID gốc: 8.6]
% [Nguồn: DBI202 Order 8 Q6]
What is NOT a purpose of the index in MS SQL Server?
\begin{choices}
\choice To use fewer hardware resources
\choice To enhance query performance
\choice To provide an index to a record
\choice To perform fast searching
\end{choices}
\textbf{Đáp án:} A
\textbf{Giải thích:} Đáp án theo dữ liệu đề: A — A. To use fewer hardware resources.

Lời giải gốc: Chỉ mục (index) giúp tăng tốc độ tìm kiếm và tối ưu hiệu suất truy vấn, nhưng đổi lại nó tiêu tốn thêm tài nguyên phần cứng (dung lượng lưu trữ và bộ nhớ) để duy trì cấu trúc index.

Nền tảng kiến thức: Tổ chức tệp quyết định cách bản ghi nằm trên trang; chỉ mục ánh xạ giá trị khóa tìm kiếm tới vị trí bản ghi để giảm số trang phải đọc. Chỉ mục tăng tốc một số truy vấn nhưng tiêu tốn dung lượng và cần cập nhật khi dữ liệu thay đổi.

Phân tích: Index lưu cấu trúc tra cứu phụ để tìm vị trí bản ghi nhanh hơn thay vì quét toàn bộ bảng. Bù lại, nó dùng thêm dung lượng và phải được duy trì khi INSERT/UPDATE/DELETE; vì thế “càng nhiều index càng tốt” không đúng cho mọi workload.

Đối chiếu phương án: B. To enhance query performance; C. To provide an index to a record; D. To perform fast searching. Các lựa chọn này không đáp ứng định nghĩa/điều kiện được hỏi; khi ôn, hãy chỉ ra chính xác điểm lệch của từng lựa chọn thay vì học thuộc mỗi chữ cái.
\end{question}

\begin{question}{480}
% [ID gốc: 8.7]
% [Nguồn: DBI202 Order 8 Q7]
Select the most correct answer:
\begin{choices}
\choice An index is not a data structure used to speed access to tuples of a relation, given values of one or more attributes.
\choice The key to an index cannot be any attribute or set of attributes, and need not be the key to the relation.
\choice We can think of the index as a binary search tree of (key, locations) pairs in which a key is not associated with a set of locations of the tuples.
\choice An index is a data structure used to speed access to tuples of a relation, given values of one or more attributes.
\end{choices}
\textbf{Đáp án:} D
\textbf{Giải thích:} Đáp án theo dữ liệu đề: D — D. An index is a data structure used to speed access to tuples of a relation, given values of one or more attributes..

Lời giải gốc: Chỉ mục (Index) là một cấu trúc dữ liệu được dùng để tăng tốc độ truy cập vào các bộ của một quan hệ dựa trên giá trị của một hoặc nhiều thuộc tính.

Nền tảng kiến thức: Tổ chức tệp quyết định cách bản ghi nằm trên trang; chỉ mục ánh xạ giá trị khóa tìm kiếm tới vị trí bản ghi để giảm số trang phải đọc. Chỉ mục tăng tốc một số truy vấn nhưng tiêu tốn dung lượng và cần cập nhật khi dữ liệu thay đổi.

Phân tích: Đọc truy vấn theo các giai đoạn logic: FROM/JOIN tạo dữ liệu nguồn, WHERE lọc từng hàng, GROUP BY tạo nhóm, HAVING lọc nhóm, SELECT xác định biểu thức đầu ra và ORDER BY sắp kết quả. Với phép join, kiểm tra điều kiện ghép và các hàng không khớp trước khi tính kết quả.

Đối chiếu phương án: A. An index is not a data structure used to speed access to tuples of a relation, given values of…; B. The key to an index cannot be any attribute or set of attributes, and need not be the key to t…; C. We can think of the index as a binary search tree of (key, locations) pairs in which a key is…. Các lựa chọn này không đáp ứng định nghĩa/điều kiện được hỏi; khi ôn, hãy chỉ ra chính xác điểm lệch của từng lựa chọn thay vì học thuộc mỗi chữ cái.
\end{question}

\begin{question}{481}
% [ID gốc: 8.10]
% [Nguồn: DBI202 Order 8 Q10]
What is an index?
\begin{choices}
\choice An index is a database table attribute that speeds up data search within a table.
\choice An index is the same as an alias.
\choice An index is a special way to join 2 or more tables.
\choice An index is a database row attribute that speeds up data search within a table.
\end{choices}
\textbf{Đáp án:} A
\textbf{Giải thích:} Đáp án theo dữ liệu đề: A — A. An index is a database table attribute that speeds up data search within a table..

Lời giải gốc: Chỉ mục (Index) là một cấu trúc dữ liệu liên kết với các thuộc tính của bảng giúp tăng tốc độ tìm kiếm và truy vấn dữ liệu.

Nền tảng kiến thức: Tổ chức tệp quyết định cách bản ghi nằm trên trang; chỉ mục ánh xạ giá trị khóa tìm kiếm tới vị trí bản ghi để giảm số trang phải đọc. Chỉ mục tăng tốc một số truy vấn nhưng tiêu tốn dung lượng và cần cập nhật khi dữ liệu thay đổi.

Phân tích: Index lưu cấu trúc tra cứu phụ để tìm vị trí bản ghi nhanh hơn thay vì quét toàn bộ bảng. Bù lại, nó dùng thêm dung lượng và phải được duy trì khi INSERT/UPDATE/DELETE; vì thế “càng nhiều index càng tốt” không đúng cho mọi workload.

Đối chiếu phương án: B. An index is the same as an alias.; C. An index is a special way to join 2 or more tables.; D. An index is a database row attribute that speeds up data search within a table.. Các lựa chọn này không đáp ứng định nghĩa/điều kiện được hỏi; khi ôn, hãy chỉ ra chính xác điểm lệch của từng lựa chọn thay vì học thuộc mỗi chữ cái.
\end{question}

\begin{question}{482}
% [ID gốc: 8.11]
% [Nguồn: DBI202 Order 8 Q11]
Which SQL command is used to create a non-clustered index on the EmployeeName column in the Employees table?
\begin{choices}
\choice CREATE INDEX EmployeeName_Index ON Employees(EmployeeName);
\choice CREATE NONCLUSTERED INDEX EmployeeName_Index ON Employees(EmployeeName);
\choice CREATE PRIMARY KEY INDEX EmployeeName_Index ON Employees(EmployeeName);
\choice CREATE CLUSTERED INDEX EmployeeName_Index ON Employees(EmployeeName);
\end{choices}
\textbf{Đáp án:} B
\textbf{Giải thích:} Đáp án theo dữ liệu đề: B — B. CREATE NONCLUSTERED INDEX EmployeeName_Index ON Employees(EmployeeName);.

Lời giải gốc: Cú pháp tạo Non-clustered index rõ ràng nhất trong SQL Server là dùng từ khóa CREATE NONCLUSTERED INDEX.

Nền tảng kiến thức: Tổ chức tệp quyết định cách bản ghi nằm trên trang; chỉ mục ánh xạ giá trị khóa tìm kiếm tới vị trí bản ghi để giảm số trang phải đọc. Chỉ mục tăng tốc một số truy vấn nhưng tiêu tốn dung lượng và cần cập nhật khi dữ liệu thay đổi.

Phân tích: Index lưu cấu trúc tra cứu phụ để tìm vị trí bản ghi nhanh hơn thay vì quét toàn bộ bảng. Bù lại, nó dùng thêm dung lượng và phải được duy trì khi INSERT/UPDATE/DELETE; vì thế “càng nhiều index càng tốt” không đúng cho mọi workload.

Đối chiếu phương án: A. CREATE INDEX EmployeeName_Index ON Employees(EmployeeName);; C. CREATE PRIMARY KEY INDEX EmployeeName_Index ON Employees(EmployeeName);; D. CREATE CLUSTERED INDEX EmployeeName_Index ON Employees(EmployeeName);. Các lựa chọn này không đáp ứng định nghĩa/điều kiện được hỏi; khi ôn, hãy chỉ ra chính xác điểm lệch của từng lựa chọn thay vì học thuộc mỗi chữ cái.
\end{question}

\begin{question}{483}
% [ID gốc: 8.12]
% [Nguồn: DBI202 Order 8 Q12]
Which indexes can be automatically created by defining constraints in a table? (Choose two correct answers)
\begin{choices}
\choice PRIMARY KEY
\choice CHECK
\choice UNIQUE
\choice FOREIGN KEY
\end{choices}
\textbf{Đáp án:} AC
\textbf{Giải thích:} Đáp án theo dữ liệu đề: A, C — A. PRIMARY KEY; C. UNIQUE.

Lời giải gốc: Hệ quản trị cơ sở dữ liệu (như SQL Server) sẽ tự động tạo một Chỉ mục cụm (Clustered Index) khi người dùng định nghĩa ràng buộc PRIMARY KEY, và tự động tạo Chỉ mục không cụm (Non-clustered Index) khi định nghĩa ràng buộc duy nhất UNIQUE.

Nền tảng kiến thức: Tổ chức tệp quyết định cách bản ghi nằm trên trang; chỉ mục ánh xạ giá trị khóa tìm kiếm tới vị trí bản ghi để giảm số trang phải đọc. Chỉ mục tăng tốc một số truy vấn nhưng tiêu tốn dung lượng và cần cập nhật khi dữ liệu thay đổi.

Phân tích: Index lưu cấu trúc tra cứu phụ để tìm vị trí bản ghi nhanh hơn thay vì quét toàn bộ bảng. Bù lại, nó dùng thêm dung lượng và phải được duy trì khi INSERT/UPDATE/DELETE; vì thế “càng nhiều index càng tốt” không đúng cho mọi workload.

Đối chiếu phương án: B. CHECK; D. FOREIGN KEY. Các lựa chọn này không đáp ứng định nghĩa/điều kiện được hỏi; khi ôn, hãy chỉ ra chính xác điểm lệch của từng lựa chọn thay vì học thuộc mỗi chữ cái.
\end{question}

\begin{question}{484}
% [ID gốc: 16.3]
% [Nguồn: DBI202 Order 16 Q3]
Which of the following statements is the most correct?
\begin{choices}
\choice Poorly designed indexes and a lack of indexes are primary sources of database application bottlenecks.
\choice Designing efficient indexes is paramount to achieve good database and application performance
\choice The selection of the right indexes for a database and its workload is a complex balancing act between query speed and update cost
\choice All of the others
\end{choices}
\textbf{Đáp án:} D
\textbf{Giải thích:} Đáp án theo dữ liệu đề: D — D. All of the others.

Nền tảng kiến thức: Thiết kế vật lý chọn chỉ mục, cách tổ chức tệp và các tối ưu theo workload thực tế. Denormalization hoặc thêm chỉ mục có thể giảm chi phí đọc nhưng đổi lại tăng lưu trữ, chi phí ghi và nguy cơ phải đồng bộ dữ liệu dư thừa.

Phân tích: Với lựa chọn tổng hợp, không chọn theo cảm tính: kiểm tra từng phát biểu độc lập. “All of the others/above” chỉ đúng khi mọi phát biểu còn lại đúng; “None” chỉ đúng khi không phát biểu nào trong nhóm thỏa yêu cầu. Đáp án lưu cần được đọc cùng các mệnh đề thành phần.

Đối chiếu phương án: A. Poorly designed indexes and a lack of indexes are primary sources of database application bott…; B. Designing efficient indexes is paramount to achieve good database and application performance; C. The selection of the right indexes for a database and its workload is a complex balancing act…. Các lựa chọn này không đáp ứng định nghĩa/điều kiện được hỏi; khi ôn, hãy chỉ ra chính xác điểm lệch của từng lựa chọn thay vì học thuộc mỗi chữ cái.
\end{question}

\begin{question}{485}
% [ID gốc: 16.4]
% [Nguồn: DBI202 Order 16 Q4]
A materialized view is a database object that contains the results of a query. It is a precomputed data set that can be used to replicate data to other locations, eliminating the need to run expensive queries every time you retrieve data from a database. Since materialized views store their output data, querying them is faster than running queries directly on base relations (tables).
Materialized views can be set to refresh on a schedule so that the updated information won't fall through the cracks. When used in conjunction with other development tools, a materialized view can create a more secure environment for your database.
Given relations Movies(title, year, length, genre. studioName. producer#). MovieExec(name. address. cert#). and Studio(name. address. presC#). Suppose we have the materialized view that finds the name - of the producer of a given movie as follows:
CREATE MATERIALIZED VIEW MovieProd AS
SELECT title, year, name
FROM Movies. MovieExec
WHERE producer# - cert#
Which of the following modification will affect the MovieProd materialized view?
\begin{choices}
\choice Insert a new tuple into Studio
\choice Delete a tuple from Studio
\choice Insert a new tuple into Movies or delete a tuple from Movies
\choice None of the others
\end{choices}
\textbf{Đáp án:} C
\textbf{Giải thích:} Đáp án theo dữ liệu đề: C — C. Insert a new tuple into Movies or delete a tuple from Movies.

Nền tảng kiến thức: Thiết kế vật lý chọn chỉ mục, cách tổ chức tệp và các tối ưu theo workload thực tế. Denormalization hoặc thêm chỉ mục có thể giảm chi phí đọc nhưng đổi lại tăng lưu trữ, chi phí ghi và nguy cơ phải đồng bộ dữ liệu dư thừa.

Phân tích: View là quan hệ ảo được định nghĩa bằng truy vấn trên một hay nhiều bảng/view; thông thường nó lưu định nghĩa truy vấn chứ không lưu bản sao kết quả như materialized view. Do đó đáp án phải chỉ đối tượng được tạo từ truy vấn, không nhầm với bảng vật lý.

Đối chiếu phương án: A. Insert a new tuple into Studio; B. Delete a tuple from Studio; D. None of the others. Các lựa chọn này không đáp ứng định nghĩa/điều kiện được hỏi; khi ôn, hãy chỉ ra chính xác điểm lệch của từng lựa chọn thay vì học thuộc mỗi chữ cái.
\end{question}

\begin{question}{486}
% [ID gốc: 16.5]
% [Nguồn: DBI202 Order 16 Q5]
What is the main disadvantage of having too many indexes on a table?
\begin{choices}
\choice Improved query performance
\choice Increased storage space
\choice Enhanced data integrity
\choice Fast data modification operations
\end{choices}
\textbf{Đáp án:} B
\textbf{Giải thích:} Đáp án theo dữ liệu đề: B — B. Increased storage space.

Nền tảng kiến thức: Thiết kế vật lý chọn chỉ mục, cách tổ chức tệp và các tối ưu theo workload thực tế. Denormalization hoặc thêm chỉ mục có thể giảm chi phí đọc nhưng đổi lại tăng lưu trữ, chi phí ghi và nguy cơ phải đồng bộ dữ liệu dư thừa.

Phân tích: Đọc truy vấn theo các giai đoạn logic: FROM/JOIN tạo dữ liệu nguồn, WHERE lọc từng hàng, GROUP BY tạo nhóm, HAVING lọc nhóm, SELECT xác định biểu thức đầu ra và ORDER BY sắp kết quả. Với phép join, kiểm tra điều kiện ghép và các hàng không khớp trước khi tính kết quả.

Đối chiếu phương án: A. Improved query performance; C. Enhanced data integrity; D. Fast data modification operations. Các lựa chọn này không đáp ứng định nghĩa/điều kiện được hỏi; khi ôn, hãy chỉ ra chính xác điểm lệch của từng lựa chọn thay vì học thuộc mỗi chữ cái.
\end{question}

\begin{question}{487}
% [ID gốc: 16.8]
% [Nguồn: DBI202 Order 16 Q8]
In query optimization, why is it beneficial for the database engine to use indexes?
\begin{choices}
\choice To reduce table size
\choice To enforce foreign key constraints
\choice To speed up data retrieval
\end{choices}
\textbf{Đáp án:} C
\textbf{Giải thích:} Đáp án theo dữ liệu đề: C — C. To speed up data retrieval.

Lời giải gốc: Mục đích chính và quan trọng nhất của việc tạo Chỉ mục (Index) trong cơ sở dữ liệu là để giúp tối ưu hóa truy vấn, tăng tốc độ tìm kiếm và truy xuất dữ liệu (speed up data retrieval).

Nền tảng kiến thức: Thiết kế vật lý chọn chỉ mục, cách tổ chức tệp và các tối ưu theo workload thực tế. Denormalization hoặc thêm chỉ mục có thể giảm chi phí đọc nhưng đổi lại tăng lưu trữ, chi phí ghi và nguy cơ phải đồng bộ dữ liệu dư thừa.

Phân tích: Đọc truy vấn theo các giai đoạn logic: FROM/JOIN tạo dữ liệu nguồn, WHERE lọc từng hàng, GROUP BY tạo nhóm, HAVING lọc nhóm, SELECT xác định biểu thức đầu ra và ORDER BY sắp kết quả. Với phép join, kiểm tra điều kiện ghép và các hàng không khớp trước khi tính kết quả.

Đối chiếu phương án: A. To reduce table size; B. To enforce foreign key constraints. Các lựa chọn này không đáp ứng định nghĩa/điều kiện được hỏi; khi ôn, hãy chỉ ra chính xác điểm lệch của từng lựa chọn thay vì học thuộc mỗi chữ cái.
\end{question}

\begin{question}{488}
% [ID gốc: 16.9]
% [Nguồn: DBI202 Order 16 Q9]
In which scenario would denormalization be considered a good design choice?
\begin{choices}
\choice When maximizing disk usage is a priority
\choice When minimizing data redundancy is the main goal
\choice When performance of read-heavy queries is critical
\choice When multiple developers are working on schema design
\end{choices}
\textbf{Đáp án:} C
\textbf{Giải thích:} Đáp án theo dữ liệu đề: C — C. When performance of read-heavy queries is critical.

Lời giải gốc: Phi chuẩn hóa (Denormalization) là kỹ thuật "phá vỡ" các dạng chuẩn, cố tình gộp dữ liệu trùng lặp. Việc này giúp giảm số lượng lệnh JOIN, từ đó tăng tốc độ cho các hệ thống cần truy vấn (đọc) cực nhanh. Câu 38

Nền tảng kiến thức: Thiết kế vật lý chọn chỉ mục, cách tổ chức tệp và các tối ưu theo workload thực tế. Denormalization hoặc thêm chỉ mục có thể giảm chi phí đọc nhưng đổi lại tăng lưu trữ, chi phí ghi và nguy cơ phải đồng bộ dữ liệu dư thừa.

Phân tích: Denormalization chủ động chấp nhận một phần dư thừa hoặc lưu trước kết quả để giảm join/tính toán ở workload đọc. Nó có thể tăng tốc truy xuất, nhưng làm tăng dung lượng và buộc hệ thống cập nhật đồng bộ các bản sao; đây là quyết định theo workload, không phải quy tắc luôn có lợi.

Đối chiếu phương án: A. When maximizing disk usage is a priority; B. When minimizing data redundancy is the main goal; D. When multiple developers are working on schema design. Các lựa chọn này không đáp ứng định nghĩa/điều kiện được hỏi; khi ôn, hãy chỉ ra chính xác điểm lệch của từng lựa chọn thay vì học thuộc mỗi chữ cái.
\end{question}

\begin{question}{489}
% [ID gốc: 18.1]
% [Nguồn: DBI202 Order 18 Q1]
atomicity the all-or-nothing execution of transactions
consistency all databases have consistency constraints, or expectations about relationships among data elements (e.g., account balances may not be negative after a transaction finishes). Transactions are expected to preserve the consistency of the database.
isolation the fact that each transaction must appear to be executed as if no other transaction is executing at the same time.
durability the effect on the database of a transaction must never be lost, once the transaction has completed.
QN=40 (8560) A database transaction, by definition, must be ACID (atomic, consistent, isolated and durable). What does "Durable" mean?
\begin{choices}
\choice "Durable" means that: Transactions provide an "all-or-nothing" proposition, stating that each work-unit performed in a database must either complete in its entirety or have no effect whatsoever
\choice "Durable" means that: Transactions must not violate any integrity constraints during its execution
\choice "Durable" means that: Transactions that have committed will survive permanently
\choice All of the others
\end{choices}
\textbf{Đáp án:} C
\textbf{Giải thích:} Đáp án theo dữ liệu đề: C — C. "Durable" means that: Transactions that have committed will survive permanently.

Nền tảng kiến thức: Giao dịch là một đơn vị công việc cần thỏa ACID: atomicity (trọn vẹn hoặc hủy), consistency (giữ ràng buộc), isolation (hạn chế ảnh hưởng xen kẽ), durability (đã commit thì bền vững qua sự cố). COMMIT và ROLLBACK kết thúc theo hai hướng khác nhau.

Phân tích: Đối chiếu đúng thuộc tính ACID được hỏi: atomicity là tất-cả-hoặc-không; consistency giữ các ràng buộc; isolation kiểm soát ảnh hưởng giữa giao dịch đồng thời; durability bảo đảm giao dịch đã commit sống sót qua sự cố. COMMIT xác nhận, ROLLBACK hủy thay đổi chưa xác nhận.

Đối chiếu phương án: A. "Durable" means that: Transactions provide an "all-or-nothing" proposition, stating that each…; B. "Durable" means that: Transactions must not violate any integrity constraints during its execu…; D. All of the others. Các lựa chọn này không đáp ứng định nghĩa/điều kiện được hỏi; khi ôn, hãy chỉ ra chính xác điểm lệch của từng lựa chọn thay vì học thuộc mỗi chữ cái.
\end{question}

\begin{question}{490}
% [ID gốc: 18.2]
% [Nguồn: DBI202 Order 18 Q2]
A database transaction, by definition, must be ACID (atomic, consistent, isolated and durable). What does "Isolated" mean?
\begin{choices}
\choice "Isolated" means that: Transactions provide an "all-or-nothing" proposition, stating that each work-unit performed in a database must either complete in its entirety or have no effect whatsoever
\choice "Isolated" means that: Transactions must not violate any integrity constraints during its execution
\choice "Isolated" means that: Transactions that have committed will survive permanently
\choice "Isolated" means that: how/when the changes made by one operation in one transaction become visible to other concurrent operations in other transactions
\end{choices}
\textbf{Đáp án:} D
\textbf{Giải thích:} Đáp án theo dữ liệu đề: D — D. "Isolated" means that: how/when the changes made by one operation in one transaction become visible to other concurrent operations in other transactions.

Nền tảng kiến thức: Giao dịch là một đơn vị công việc cần thỏa ACID: atomicity (trọn vẹn hoặc hủy), consistency (giữ ràng buộc), isolation (hạn chế ảnh hưởng xen kẽ), durability (đã commit thì bền vững qua sự cố). COMMIT và ROLLBACK kết thúc theo hai hướng khác nhau.

Phân tích: Đối chiếu đúng thuộc tính ACID được hỏi: atomicity là tất-cả-hoặc-không; consistency giữ các ràng buộc; isolation kiểm soát ảnh hưởng giữa giao dịch đồng thời; durability bảo đảm giao dịch đã commit sống sót qua sự cố. COMMIT xác nhận, ROLLBACK hủy thay đổi chưa xác nhận.

Đối chiếu phương án: A. "Isolated" means that: Transactions provide an "all-or-nothing" proposition, stating that each…; B. "Isolated" means that: Transactions must not violate any integrity constraints during its exec…; C. "Isolated" means that: Transactions that have committed will survive permanently. Các lựa chọn này không đáp ứng định nghĩa/điều kiện được hỏi; khi ôn, hãy chỉ ra chính xác điểm lệch của từng lựa chọn thay vì học thuộc mỗi chữ cái.
\end{question}

\begin{question}{491}
% [ID gốc: 18.3]
% [Nguồn: DBI202 Order 18 Q3]
Regardless of whatever any other transaction is doing, a transaction must be able to continue with the exact same data set it started with.
The above describes which property of a transaction?
\begin{choices}
\choice Atomic
\choice Isolation
\choice Consistency
\choice Durability
\end{choices}
\textbf{Đáp án:} B
\textbf{Giải thích:} Đáp án theo dữ liệu đề: B — B. Isolation.

Nền tảng kiến thức: Giao dịch là một đơn vị công việc cần thỏa ACID: atomicity (trọn vẹn hoặc hủy), consistency (giữ ràng buộc), isolation (hạn chế ảnh hưởng xen kẽ), durability (đã commit thì bền vững qua sự cố). COMMIT và ROLLBACK kết thúc theo hai hướng khác nhau.

Phân tích: Đối chiếu đúng thuộc tính ACID được hỏi: atomicity là tất-cả-hoặc-không; consistency giữ các ràng buộc; isolation kiểm soát ảnh hưởng giữa giao dịch đồng thời; durability bảo đảm giao dịch đã commit sống sót qua sự cố. COMMIT xác nhận, ROLLBACK hủy thay đổi chưa xác nhận.

Đối chiếu phương án: A. Atomic; C. Consistency; D. Durability. Các lựa chọn này không đáp ứng định nghĩa/điều kiện được hỏi; khi ôn, hãy chỉ ra chính xác điểm lệch của từng lựa chọn thay vì học thuộc mỗi chữ cái.
\end{question}

\begin{question}{492}
% [ID gốc: 18.4]
% [Nguồn: DBI202 Order 18 Q4]
Choose the in-correct statement:
\begin{choices}
\choice In SQL Server, every DML operation is a transaction regardless of whether it has a BEGIN TRANSACTION or not
\choice COMMITTING a transaction makes permanent the changes resulting from all SQL statements in the transaction
\choice ROLLING back a transaction rejects any of the changes resulting from the SQL statements in the transaction
\choice When the sequence of tasks is complete, the SAVE POINT closes the transaction
\end{choices}
\textbf{Đáp án:} D
\textbf{Giải thích:} Đáp án theo dữ liệu đề: D — D. When the sequence of tasks is complete, the SAVE POINT closes the transaction.

Nền tảng kiến thức: Giao dịch là một đơn vị công việc cần thỏa ACID: atomicity (trọn vẹn hoặc hủy), consistency (giữ ràng buộc), isolation (hạn chế ảnh hưởng xen kẽ), durability (đã commit thì bền vững qua sự cố). COMMIT và ROLLBACK kết thúc theo hai hướng khác nhau.

Phân tích: DDL (Data Definition Language) khai báo/thay đổi cấu trúc như bảng và schema; DML truy vấn/thêm/sửa/xóa dữ liệu; DCL quản lý quyền. Chọn theo đối tượng bị tác động: cấu trúc, các hàng hay quyền truy cập.

Đối chiếu phương án: A. In SQL Server, every DML operation is a transaction regardless of whether it has a BEGIN TRANS…; B. COMMITTING a transaction makes permanent the changes resulting from all SQL statements in the…; C. ROLLING back a transaction rejects any of the changes resulting from the SQL statements in the…. Các lựa chọn này không đáp ứng định nghĩa/điều kiện được hỏi; khi ôn, hãy chỉ ra chính xác điểm lệch của từng lựa chọn thay vì học thuộc mỗi chữ cái.
\end{question}

\begin{question}{493}
% [ID gốc: 18.5]
% [Nguồn: DBI202 Order 18 Q5]
How many properties we must use to measure the quality of a transaction?
\begin{choices}
\choice 6
\choice 4
\choice 3
\choice 5
\end{choices}
\textbf{Đáp án:} B
\textbf{Giải thích:} A transaction is measured by the 4 ACID properties: Atomicity, Consistency, Isolation, Durability.
\end{question}

\begin{question}{494}
% [ID gốc: 18.6]
% [Nguồn: DBI202 Order 18 Q6]
Every row and value must agree with all constraints right after the transaction is complete. The above describes which property of a transaction?
\begin{choices}
\choice Isolation
\choice Consistency
\choice Atomic.
\choice Durability
\end{choices}
\textbf{Đáp án:} B
\textbf{Giải thích:} Đáp án theo dữ liệu đề: B — B. Consistency.

Nền tảng kiến thức: Giao dịch là một đơn vị công việc cần thỏa ACID: atomicity (trọn vẹn hoặc hủy), consistency (giữ ràng buộc), isolation (hạn chế ảnh hưởng xen kẽ), durability (đã commit thì bền vững qua sự cố). COMMIT và ROLLBACK kết thúc theo hai hướng khác nhau.

Phân tích: Đối chiếu đúng thuộc tính ACID được hỏi: atomicity là tất-cả-hoặc-không; consistency giữ các ràng buộc; isolation kiểm soát ảnh hưởng giữa giao dịch đồng thời; durability bảo đảm giao dịch đã commit sống sót qua sự cố. COMMIT xác nhận, ROLLBACK hủy thay đổi chưa xác nhận.

Đối chiếu phương án: A. Isolation; C. Atomic.; D. Durability. Các lựa chọn này không đáp ứng định nghĩa/điều kiện được hỏi; khi ôn, hãy chỉ ra chính xác điểm lệch của từng lựa chọn thay vì học thuộc mỗi chữ cái.
\end{question}

\begin{question}{495}
% [ID gốc: 18.7]
% [Nguồn: DBI202 Order 18 Q7]
What happens if an error occurs during a transaction and it is not explicitly handled or rolled back?
\begin{choices}
\choice The transaction is automatically committed
\choice The transaction is automatically rolled back
\choice The database becomes read-only
\choice The transaction remains in a pending state
\end{choices}
\textbf{Đáp án:} B
\textbf{Giải thích:} Đáp án theo dữ liệu đề: B — B. The transaction is automatically rolled back.

Nền tảng kiến thức: Giao dịch là một đơn vị công việc cần thỏa ACID: atomicity (trọn vẹn hoặc hủy), consistency (giữ ràng buộc), isolation (hạn chế ảnh hưởng xen kẽ), durability (đã commit thì bền vững qua sự cố). COMMIT và ROLLBACK kết thúc theo hai hướng khác nhau.

Phân tích: Đối chiếu đúng thuộc tính ACID được hỏi: atomicity là tất-cả-hoặc-không; consistency giữ các ràng buộc; isolation kiểm soát ảnh hưởng giữa giao dịch đồng thời; durability bảo đảm giao dịch đã commit sống sót qua sự cố. COMMIT xác nhận, ROLLBACK hủy thay đổi chưa xác nhận.

Đối chiếu phương án: A. The transaction is automatically committed; C. The database becomes read-only; D. The transaction remains in a pending state. Các lựa chọn này không đáp ứng định nghĩa/điều kiện được hỏi; khi ôn, hãy chỉ ra chính xác điểm lệch của từng lựa chọn thay vì học thuộc mỗi chữ cái.
\end{question}

\begin{question}{496}
% [ID gốc: 18.8]
% [Nguồn: DBI202 Order 18 Q8]
The transaction methods is used with the Connection object to save or cancel changes made to the data source.
\begin{choices}
\choice Begin Transaction, Rollback Transaction
\choice Begin Transaction, Commit Transaction
\choice Begin Transaction, Commit Transaction, Rollback Transaction
\choice Commit Transaction, Rollback Transaction
\end{choices}
\textbf{Đáp án:} C
\textbf{Giải thích:} Đáp án theo dữ liệu đề: C — C. Begin Transaction, Commit Transaction, Rollback Transaction.

Nền tảng kiến thức: Giao dịch là một đơn vị công việc cần thỏa ACID: atomicity (trọn vẹn hoặc hủy), consistency (giữ ràng buộc), isolation (hạn chế ảnh hưởng xen kẽ), durability (đã commit thì bền vững qua sự cố). COMMIT và ROLLBACK kết thúc theo hai hướng khác nhau.

Phân tích: Đối chiếu đúng thuộc tính ACID được hỏi: atomicity là tất-cả-hoặc-không; consistency giữ các ràng buộc; isolation kiểm soát ảnh hưởng giữa giao dịch đồng thời; durability bảo đảm giao dịch đã commit sống sót qua sự cố. COMMIT xác nhận, ROLLBACK hủy thay đổi chưa xác nhận.

Đối chiếu phương án: A. Begin Transaction, Rollback Transaction; B. Begin Transaction, Commit Transaction; D. Commit Transaction, Rollback Transaction. Các lựa chọn này không đáp ứng định nghĩa/điều kiện được hỏi; khi ôn, hãy chỉ ra chính xác điểm lệch của từng lựa chọn thay vì học thuộc mỗi chữ cái.
\end{question}

\begin{question}{497}
% [ID gốc: 18.9]
% [Nguồn: DBI202 Order 18 Q9]
What are the ACID properties of DBMS Transactions?
\begin{choices}
\choice Atomicity, Consistency, Isolation, Durability
\choice Atomicity, Consistency, Interpolate, Database
\choice Atomicity, Consistency, Inconsistent, Durability
\choice Automatic, Concurrent, Inconsistent, Deadlock-freedom
\end{choices}
\textbf{Đáp án:} A
\textbf{Giải thích:} Đáp án theo dữ liệu đề: A — A. Atomicity, Consistency, Isolation, Durability.

Nền tảng kiến thức: Giao dịch là một đơn vị công việc cần thỏa ACID: atomicity (trọn vẹn hoặc hủy), consistency (giữ ràng buộc), isolation (hạn chế ảnh hưởng xen kẽ), durability (đã commit thì bền vững qua sự cố). COMMIT và ROLLBACK kết thúc theo hai hướng khác nhau.

Phân tích: Đối chiếu đúng thuộc tính ACID được hỏi: atomicity là tất-cả-hoặc-không; consistency giữ các ràng buộc; isolation kiểm soát ảnh hưởng giữa giao dịch đồng thời; durability bảo đảm giao dịch đã commit sống sót qua sự cố. COMMIT xác nhận, ROLLBACK hủy thay đổi chưa xác nhận.

Đối chiếu phương án: B. Atomicity, Consistency, Interpolate, Database; C. Atomicity, Consistency, Inconsistent, Durability; D. Automatic, Concurrent, Inconsistent, Deadlock-freedom. Các lựa chọn này không đáp ứng định nghĩa/điều kiện được hỏi; khi ôn, hãy chỉ ra chính xác điểm lệch của từng lựa chọn thay vì học thuộc mỗi chữ cái.
\end{question}

\begin{question}{498}
% [ID gốc: 18.11]
% [Nguồn: DBI202 Order 18 Q11]
Which of the following makes the transaction permanent in the database?
\begin{choices}
\choice View
\choice Commit
\choice Rollback
\choice Flashback
\end{choices}
\textbf{Đáp án:} B
\textbf{Giải thích:} Đáp án theo dữ liệu đề: B — B. Commit.

Nền tảng kiến thức: Giao dịch là một đơn vị công việc cần thỏa ACID: atomicity (trọn vẹn hoặc hủy), consistency (giữ ràng buộc), isolation (hạn chế ảnh hưởng xen kẽ), durability (đã commit thì bền vững qua sự cố). COMMIT và ROLLBACK kết thúc theo hai hướng khác nhau.

Phân tích: Đối chiếu đúng thuộc tính ACID được hỏi: atomicity là tất-cả-hoặc-không; consistency giữ các ràng buộc; isolation kiểm soát ảnh hưởng giữa giao dịch đồng thời; durability bảo đảm giao dịch đã commit sống sót qua sự cố. COMMIT xác nhận, ROLLBACK hủy thay đổi chưa xác nhận.

Đối chiếu phương án: A. View; C. Rollback; D. Flashback. Các lựa chọn này không đáp ứng định nghĩa/điều kiện được hỏi; khi ôn, hãy chỉ ra chính xác điểm lệch của từng lựa chọn thay vì học thuộc mỗi chữ cái.
\end{question}

\begin{question}{499}
% [ID gốc: 18.12]
% [Nguồn: DBI202 Order 18 Q12]
In order to undo the work of transaction after last commit which one should be used?
\begin{choices}
\choice View
\choice Commit
\choice Rollback
\choice Flashback
\end{choices}
\textbf{Đáp án:} C
\textbf{Giải thích:} Đáp án theo dữ liệu đề: C — C. Rollback.

Nền tảng kiến thức: Giao dịch là một đơn vị công việc cần thỏa ACID: atomicity (trọn vẹn hoặc hủy), consistency (giữ ràng buộc), isolation (hạn chế ảnh hưởng xen kẽ), durability (đã commit thì bền vững qua sự cố). COMMIT và ROLLBACK kết thúc theo hai hướng khác nhau.

Phân tích: Đối chiếu đúng thuộc tính ACID được hỏi: atomicity là tất-cả-hoặc-không; consistency giữ các ràng buộc; isolation kiểm soát ảnh hưởng giữa giao dịch đồng thời; durability bảo đảm giao dịch đã commit sống sót qua sự cố. COMMIT xác nhận, ROLLBACK hủy thay đổi chưa xác nhận.

Đối chiếu phương án: A. View; B. Commit; D. Flashback. Các lựa chọn này không đáp ứng định nghĩa/điều kiện được hỏi; khi ôn, hãy chỉ ra chính xác điểm lệch của từng lựa chọn thay vì học thuộc mỗi chữ cái.
\end{question}

\begin{question}{500}
% [ID gốc: 18.13]
% [Nguồn: DBI202 Order 18 Q13]
The Isolation property in transactions ensures that:
\begin{choices}
\choice Each transaction is encrypted
\choice No two transactions run concurrently
\choice Transactions do not interfere with each other
\choice Only one user accesses the system
\end{choices}
\textbf{Đáp án:} C
\textbf{Giải thích:} Đáp án theo dữ liệu đề: C — C. Transactions do not interfere with each other.

Lời giải gốc: Tính cô lập (Isolation) trong thuộc tính ACID của giao dịch (transaction) đảm bảo rằng các giao dịch thực thi đồng thời không can thiệp hoặc ảnh hưởng lẫn nhau cho đến khi chúng hoàn tất.

Nền tảng kiến thức: Giao dịch là một đơn vị công việc cần thỏa ACID: atomicity (trọn vẹn hoặc hủy), consistency (giữ ràng buộc), isolation (hạn chế ảnh hưởng xen kẽ), durability (đã commit thì bền vững qua sự cố). COMMIT và ROLLBACK kết thúc theo hai hướng khác nhau.

Phân tích: Đối chiếu đúng thuộc tính ACID được hỏi: atomicity là tất-cả-hoặc-không; consistency giữ các ràng buộc; isolation kiểm soát ảnh hưởng giữa giao dịch đồng thời; durability bảo đảm giao dịch đã commit sống sót qua sự cố. COMMIT xác nhận, ROLLBACK hủy thay đổi chưa xác nhận.

Đối chiếu phương án: A. Each transaction is encrypted; B. No two transactions run concurrently; D. Only one user accesses the system. Các lựa chọn này không đáp ứng định nghĩa/điều kiện được hỏi; khi ôn, hãy chỉ ra chính xác điểm lệch của từng lựa chọn thay vì học thuộc mỗi chữ cái.
\end{question}

\begin{question}{501}
% [ID gốc: 18.14]
% [Nguồn: DBI202 Order 18 Q14]
Which SQL command is used to undo changes made by a transaction?
\begin{choices}
\choice COMMIT
\choice ROLLBACK
\choice SAVEPOINT
\choice UNDO
\end{choices}
\textbf{Đáp án:} B
\textbf{Giải thích:} Đáp án theo dữ liệu đề: B — B. ROLLBACK.

Lời giải gốc: Lệnh ROLLBACK dùng để hủy bỏ các thay đổi chưa được lưu của một giao dịch và đưa cơ sở dữ liệu về trạng thái ban đầu.

Nền tảng kiến thức: Giao dịch là một đơn vị công việc cần thỏa ACID: atomicity (trọn vẹn hoặc hủy), consistency (giữ ràng buộc), isolation (hạn chế ảnh hưởng xen kẽ), durability (đã commit thì bền vững qua sự cố). COMMIT và ROLLBACK kết thúc theo hai hướng khác nhau.

Phân tích: Đối chiếu đúng thuộc tính ACID được hỏi: atomicity là tất-cả-hoặc-không; consistency giữ các ràng buộc; isolation kiểm soát ảnh hưởng giữa giao dịch đồng thời; durability bảo đảm giao dịch đã commit sống sót qua sự cố. COMMIT xác nhận, ROLLBACK hủy thay đổi chưa xác nhận.

Đối chiếu phương án: A. COMMIT; C. SAVEPOINT; D. UNDO. Các lựa chọn này không đáp ứng định nghĩa/điều kiện được hỏi; khi ôn, hãy chỉ ra chính xác điểm lệch của từng lựa chọn thay vì học thuộc mỗi chữ cái.
\end{question}

\begin{question}{502}
% [ID gốc: 18.17]
% [Nguồn: DBI202 Order 18 Q17]
Which property of ACID transactions ensures that a transaction, once committed, remains committed even in the event of a system failure?
\begin{choices}
\choice Atomicity
\choice Consistency
\choice Isolation
\choice Durability
\end{choices}
\textbf{Đáp án:} D
\textbf{Giải thích:} Đáp án theo dữ liệu đề: D — D. Durability.

Lời giải gốc: Tính bền vững (Durability) đảm bảo rằng dữ liệu của các giao dịch đã commit sẽ được lưu trữ an toàn ngay cả khi xảy ra sự cố hệ thống.

Nền tảng kiến thức: Giao dịch là một đơn vị công việc cần thỏa ACID: atomicity (trọn vẹn hoặc hủy), consistency (giữ ràng buộc), isolation (hạn chế ảnh hưởng xen kẽ), durability (đã commit thì bền vững qua sự cố). COMMIT và ROLLBACK kết thúc theo hai hướng khác nhau.

Phân tích: Đối chiếu đúng thuộc tính ACID được hỏi: atomicity là tất-cả-hoặc-không; consistency giữ các ràng buộc; isolation kiểm soát ảnh hưởng giữa giao dịch đồng thời; durability bảo đảm giao dịch đã commit sống sót qua sự cố. COMMIT xác nhận, ROLLBACK hủy thay đổi chưa xác nhận.

Đối chiếu phương án: A. Atomicity; B. Consistency; C. Isolation. Các lựa chọn này không đáp ứng định nghĩa/điều kiện được hỏi; khi ôn, hãy chỉ ra chính xác điểm lệch của từng lựa chọn thay vì học thuộc mỗi chữ cái.
\end{question}

\begin{question}{503}
% [ID gốc: 18.18]
% [Nguồn: DBI202 Order 18 Q18]
Within the context of ACID properties, which property ensures that even if the system crashes during a transaction, the database will eventually reflect either the full completion or the complete rollback of that transaction?
\begin{choices}
\choice Atomicity
\choice Consistency
\choice Isolation
\choice Durability
\end{choices}
\textbf{Đáp án:} A
\textbf{Giải thích:} Đáp án theo dữ liệu đề: A — A. Atomicity.

Lời giải gốc: Tính nguyên tử (Atomicity) đảm bảo giao dịch thực thi toàn bộ hoặc không thực hiện gì cả, không có trạng thái nửa vời.

Nền tảng kiến thức: Giao dịch là một đơn vị công việc cần thỏa ACID: atomicity (trọn vẹn hoặc hủy), consistency (giữ ràng buộc), isolation (hạn chế ảnh hưởng xen kẽ), durability (đã commit thì bền vững qua sự cố). COMMIT và ROLLBACK kết thúc theo hai hướng khác nhau.

Phân tích: Đối chiếu đúng thuộc tính ACID được hỏi: atomicity là tất-cả-hoặc-không; consistency giữ các ràng buộc; isolation kiểm soát ảnh hưởng giữa giao dịch đồng thời; durability bảo đảm giao dịch đã commit sống sót qua sự cố. COMMIT xác nhận, ROLLBACK hủy thay đổi chưa xác nhận.

Đối chiếu phương án: B. Consistency; C. Isolation; D. Durability. Các lựa chọn này không đáp ứng định nghĩa/điều kiện được hỏi; khi ôn, hãy chỉ ra chính xác điểm lệch của từng lựa chọn thay vì học thuộc mỗi chữ cái.
\end{question}

\begin{question}{504}
% [ID gốc: 18.19]
% [Nguồn: DBI202 Order 18 Q19]
The transaction methods are used with the Connection object to save or cancel changes made to the data source.
\begin{choices}
\choice Begin Transaction, Rollback Transaction
\choice Begin Transaction, Commit Transaction
\choice Begin Transaction, Commit Transaction, Rollback Transaction
\choice Commit Transaction, Rollback Transaction
\end{choices}
\textbf{Đáp án:} C
\textbf{Giải thích:} Đáp án theo dữ liệu đề: D — D. Commit Transaction, Rollback Transaction.

Lời giải gốc: Các phương thức dùng để xác nhận lưu hoặc hủy bỏ các thay đổi dữ liệu trong giao dịch là Commit Transaction và Rollback Transaction. ---

Nền tảng kiến thức: Giao dịch là một đơn vị công việc cần thỏa ACID: atomicity (trọn vẹn hoặc hủy), consistency (giữ ràng buộc), isolation (hạn chế ảnh hưởng xen kẽ), durability (đã commit thì bền vững qua sự cố). COMMIT và ROLLBACK kết thúc theo hai hướng khác nhau.

Phân tích: Đối chiếu đúng thuộc tính ACID được hỏi: atomicity là tất-cả-hoặc-không; consistency giữ các ràng buộc; isolation kiểm soát ảnh hưởng giữa giao dịch đồng thời; durability bảo đảm giao dịch đã commit sống sót qua sự cố. COMMIT xác nhận, ROLLBACK hủy thay đổi chưa xác nhận.

Đối chiếu phương án: A. Begin Transaction, Rollback Transaction; B. Begin Transaction, Commit Transaction; C. Begin Transaction, Commit Transaction, Rollback Transaction. Các lựa chọn này không đáp ứng định nghĩa/điều kiện được hỏi; khi ôn, hãy chỉ ra chính xác điểm lệch của từng lựa chọn thay vì học thuộc mỗi chữ cái.
\end{question}

\begin{question}{505}
% [ID gốc: 18.20]
% [Nguồn: DBI202 Order 18 Q20]
Which of the following best describes Atomicity in transactions?
\begin{choices}
\choice Transactions run in parallel
\choice Transactions are executed step by step
\choice All operations in a transaction succeed or none are applied
\choice Data is stored permanently
\end{choices}
\textbf{Đáp án:} C
\textbf{Giải thích:} Đáp án theo dữ liệu đề: C — C. All operations in a transaction succeed or none are applied.

Lời giải gốc: Tính nguyên tố (Atomicity - thuộc tính 'A' trong ACID) đảm bảo rằng tất cả các thao tác trong một giao dịch phải được thực hiện thành công toàn bộ hoặc không có thao tác nào được áp dụng (All or Nothing).

Nền tảng kiến thức: Giao dịch là một đơn vị công việc cần thỏa ACID: atomicity (trọn vẹn hoặc hủy), consistency (giữ ràng buộc), isolation (hạn chế ảnh hưởng xen kẽ), durability (đã commit thì bền vững qua sự cố). COMMIT và ROLLBACK kết thúc theo hai hướng khác nhau.

Phân tích: Đối chiếu đúng thuộc tính ACID được hỏi: atomicity là tất-cả-hoặc-không; consistency giữ các ràng buộc; isolation kiểm soát ảnh hưởng giữa giao dịch đồng thời; durability bảo đảm giao dịch đã commit sống sót qua sự cố. COMMIT xác nhận, ROLLBACK hủy thay đổi chưa xác nhận.

Đối chiếu phương án: A. Transactions run in parallel; B. Transactions are executed step by step; D. Data is stored permanently. Các lựa chọn này không đáp ứng định nghĩa/điều kiện được hỏi; khi ôn, hãy chỉ ra chính xác điểm lệch của từng lựa chọn thay vì học thuộc mỗi chữ cái.
\end{question}

\begin{question}{506}
% [ID gốc: 18.23]
% [Nguồn: DBI202 Order 18 Q23]
What is the default behavior if autocommit is ON?
\begin{choices}
\choice Manual COMMIT needed
\choice Every command is saved
\choice No command is saved
\choice Only SELECT works
\end{choices}
\textbf{Đáp án:} B
\textbf{Giải thích:} Đáp án theo dữ liệu đề: B — B. Every command is saved.

Lời giải gốc: Khi chế độ autocommit ở trạng thái ON, mỗi câu lệnh SQL thay đổi dữ liệu được tự động thực thi và lưu (commit) ngay lập tức vào cơ sở dữ liệu.

Nền tảng kiến thức: Giao dịch là một đơn vị công việc cần thỏa ACID: atomicity (trọn vẹn hoặc hủy), consistency (giữ ràng buộc), isolation (hạn chế ảnh hưởng xen kẽ), durability (đã commit thì bền vững qua sự cố). COMMIT và ROLLBACK kết thúc theo hai hướng khác nhau.

Phân tích: Đọc truy vấn theo các giai đoạn logic: FROM/JOIN tạo dữ liệu nguồn, WHERE lọc từng hàng, GROUP BY tạo nhóm, HAVING lọc nhóm, SELECT xác định biểu thức đầu ra và ORDER BY sắp kết quả. Với phép join, kiểm tra điều kiện ghép và các hàng không khớp trước khi tính kết quả.

Đối chiếu phương án: A. Manual COMMIT needed; C. No command is saved; D. Only SELECT works. Các lựa chọn này không đáp ứng định nghĩa/điều kiện được hỏi; khi ôn, hãy chỉ ra chính xác điểm lệch của từng lựa chọn thay vì học thuộc mỗi chữ cái.
\end{question}

\begin{question}{507}
% [ID gốc: 18.24]
% [Nguồn: DBI202 Order 18 Q24]
DB1202_SP26_RE | Multiple Choice Question 34 (Choose 1 answer) The transaction methods is used with the Connection object to save or cancel changes made to the data source.
\begin{choices}
\choice Begin Transaction, Rollback Transaction
\choice Begin Transaction, Commit Transaction
\choice Begin Transaction, Commit Transaction, Rollback Transaction
\choice Commit Transaction, Rollback Transaction
\end{choices}
\textbf{Đáp án:} C
\textbf{Giải thích:} Đáp án theo dữ liệu đề: C — C. Begin Transaction, Commit Transaction, Rollback Transaction.

Lời giải gốc: Vòng đời của một giao dịch (Transaction) có 3 lệnh chính: Begin (bắt đầu), Commit (lưu kết quả vĩnh viễn), và Rollback (hủy bỏ, quay về trạng thái lúc mới Begin khi có lỗi).

Nền tảng kiến thức: Giao dịch là một đơn vị công việc cần thỏa ACID: atomicity (trọn vẹn hoặc hủy), consistency (giữ ràng buộc), isolation (hạn chế ảnh hưởng xen kẽ), durability (đã commit thì bền vững qua sự cố). COMMIT và ROLLBACK kết thúc theo hai hướng khác nhau.

Phân tích: Đối chiếu đúng thuộc tính ACID được hỏi: atomicity là tất-cả-hoặc-không; consistency giữ các ràng buộc; isolation kiểm soát ảnh hưởng giữa giao dịch đồng thời; durability bảo đảm giao dịch đã commit sống sót qua sự cố. COMMIT xác nhận, ROLLBACK hủy thay đổi chưa xác nhận.

Đối chiếu phương án: A. Begin Transaction, Rollback Transaction; B. Begin Transaction, Commit Transaction; D. Commit Transaction, Rollback Transaction. Các lựa chọn này không đáp ứng định nghĩa/điều kiện được hỏi; khi ôn, hãy chỉ ra chính xác điểm lệch của từng lựa chọn thay vì học thuộc mỗi chữ cái.
\end{question}

\begin{question}{508}
% [ID gốc: 18.26]
% [Nguồn: DBI202 Order 18 Q26]
Which transaction methods are used with the Connection object to save or cancel changes made to the data source?
\begin{choices}
\choice Begin Transaction, Rollback Transaction
\choice Begin Transaction, Commit Transaction
\choice Begin Transaction, Commit Transaction, Rollback Transaction
\choice Commit Transaction, Rollback Transaction
\end{choices}
\textbf{Đáp án:} C
\textbf{Giải thích:} Đáp án theo dữ liệu đề: C — C. Begin Transaction, Commit Transaction, Rollback Transaction.

Lời giải gốc: Để quản lý một giao dịch hoàn chỉnh (lưu hoặc hủy thay đổi), các phương thức kết hợp gồm bắt đầu (Begin), xác nhận lưu (Commit) và hủy bỏ (Rollback).

Nền tảng kiến thức: Giao dịch là một đơn vị công việc cần thỏa ACID: atomicity (trọn vẹn hoặc hủy), consistency (giữ ràng buộc), isolation (hạn chế ảnh hưởng xen kẽ), durability (đã commit thì bền vững qua sự cố). COMMIT và ROLLBACK kết thúc theo hai hướng khác nhau.

Phân tích: Đối chiếu đúng thuộc tính ACID được hỏi: atomicity là tất-cả-hoặc-không; consistency giữ các ràng buộc; isolation kiểm soát ảnh hưởng giữa giao dịch đồng thời; durability bảo đảm giao dịch đã commit sống sót qua sự cố. COMMIT xác nhận, ROLLBACK hủy thay đổi chưa xác nhận.

Đối chiếu phương án: A. Begin Transaction, Rollback Transaction; B. Begin Transaction, Commit Transaction; D. Commit Transaction, Rollback Transaction. Các lựa chọn này không đáp ứng định nghĩa/điều kiện được hỏi; khi ôn, hãy chỉ ra chính xác điểm lệch của từng lựa chọn thay vì học thuộc mỗi chữ cái.
\end{question}

\begin{question}{509}
% [ID gốc: 18.27]
% [Nguồn: DBI202 Order 18 Q27]
In the context of a relational database, what is the correct term for "A set of changes that must all be made together"?
\begin{choices}
\choice Consistency
\choice Concurrency
\choice Transaction
\choice Atom
\end{choices}
\textbf{Đáp án:} C
\textbf{Giải thích:} Đáp án theo dữ liệu đề: C — C. Transaction.

Lời giải gốc: Giao dịch (Transaction) là một tập hợp các thay đổi thao tác dữ liệu phải được thực hiện đồng thời trọn vẹn.

Nền tảng kiến thức: Giao dịch là một đơn vị công việc cần thỏa ACID: atomicity (trọn vẹn hoặc hủy), consistency (giữ ràng buộc), isolation (hạn chế ảnh hưởng xen kẽ), durability (đã commit thì bền vững qua sự cố). COMMIT và ROLLBACK kết thúc theo hai hướng khác nhau.

Phân tích: Đối chiếu đúng thuộc tính ACID được hỏi: atomicity là tất-cả-hoặc-không; consistency giữ các ràng buộc; isolation kiểm soát ảnh hưởng giữa giao dịch đồng thời; durability bảo đảm giao dịch đã commit sống sót qua sự cố. COMMIT xác nhận, ROLLBACK hủy thay đổi chưa xác nhận.

Đối chiếu phương án: A. Consistency; B. Concurrency; D. Atom. Các lựa chọn này không đáp ứng định nghĩa/điều kiện được hỏi; khi ôn, hãy chỉ ra chính xác điểm lệch của từng lựa chọn thay vì học thuộc mỗi chữ cái.
\end{question}

\begin{question}{510}
% [ID gốc: 18.28]
% [Nguồn: DBI202 Order 18 Q28]
In SQL, what does the acronym ACID stand for concerning transactions?
\begin{choices}
\choice Atomicity, Consistency, Isolation, Durability
\choice Association, Concurrency, Integrity, Database
\choice Aggregate, Criteria, Index, Delete
\choice All Columns, Insertion, Deletion
\end{choices}
\textbf{Đáp án:} A
\textbf{Giải thích:} Đáp án theo dữ liệu đề: A — A. Atomicity, Consistency, Isolation, Durability.

Lời giải gốc: ACID viết tắt cho tính nguyên tử (Atomicity), tính nhất quán (Consistency), tính cô lập (Isolation) và tính bền vững (Durability) của giao dịch.

Nền tảng kiến thức: Giao dịch là một đơn vị công việc cần thỏa ACID: atomicity (trọn vẹn hoặc hủy), consistency (giữ ràng buộc), isolation (hạn chế ảnh hưởng xen kẽ), durability (đã commit thì bền vững qua sự cố). COMMIT và ROLLBACK kết thúc theo hai hướng khác nhau.

Phân tích: Đối chiếu đúng thuộc tính ACID được hỏi: atomicity là tất-cả-hoặc-không; consistency giữ các ràng buộc; isolation kiểm soát ảnh hưởng giữa giao dịch đồng thời; durability bảo đảm giao dịch đã commit sống sót qua sự cố. COMMIT xác nhận, ROLLBACK hủy thay đổi chưa xác nhận.

Đối chiếu phương án: B. Association, Concurrency, Integrity, Database; C. Aggregate, Criteria, Index, Delete; D. All Columns, Insertion, Deletion. Các lựa chọn này không đáp ứng định nghĩa/điều kiện được hỏi; khi ôn, hãy chỉ ra chính xác điểm lệch của từng lựa chọn thay vì học thuộc mỗi chữ cái.
\end{question}

\begin{question}{511}
% [ID gốc: 18.29]
% [Nguồn: DBI202 Order 18 Q29]
In the context of a banking application, which of the following best represents a transaction that violates atomicity? (Choose two answers)
\begin{choices}
\choice A money transfer where funds are withdrawn from one account but not credited to another due to a crash
\choice A funds transfer request fails before any account is updated, so neither the debit nor the credit happens
\choice Both the debit from Account A and the credit to Account B are completed successfully.
\choice Create a new account successfully but cannot link it to the customer profile.
\end{choices}
\textbf{Đáp án:} AD
\textbf{Giải thích:} Đáp án theo dữ liệu đề: A, D — A. A money transfer where funds are withdrawn from one account but not credited to another due to a crash; D. Create a new account successfully but cannot link it to the customer profile..

Lời giải gốc: Tính nguyên tử (Atomicity) yêu cầu một giao dịch phải thực hiện trọn vẹn "Tất cả hoặc không có gì" (All or Nothing). Việc tiền đã trừ ở tài khoản gửi nhưng chưa cộng ở tài khoản nhận do crash (A) hay việc tạo tài khoản thành công nhưng không thể liên kết profile (D) đều là trạng thái dở dang, vi phạm nghiêm trọng tính nguyên tử.

Nền tảng kiến thức: Giao dịch là một đơn vị công việc cần thỏa ACID: atomicity (trọn vẹn hoặc hủy), consistency (giữ ràng buộc), isolation (hạn chế ảnh hưởng xen kẽ), durability (đã commit thì bền vững qua sự cố). COMMIT và ROLLBACK kết thúc theo hai hướng khác nhau.

Phân tích: Đối chiếu đúng thuộc tính ACID được hỏi: atomicity là tất-cả-hoặc-không; consistency giữ các ràng buộc; isolation kiểm soát ảnh hưởng giữa giao dịch đồng thời; durability bảo đảm giao dịch đã commit sống sót qua sự cố. COMMIT xác nhận, ROLLBACK hủy thay đổi chưa xác nhận.

Đối chiếu phương án: B. A funds transfer request fails before any account is updated, so neither the debit nor the cre…; C. Both the debit from Account A and the credit to Account B are completed successfully.. Các lựa chọn này không đáp ứng định nghĩa/điều kiện được hỏi; khi ôn, hãy chỉ ra chính xác điểm lệch của từng lựa chọn thay vì học thuộc mỗi chữ cái.
\end{question}

\begin{question}{512}
% [ID gốc: 18.30]
% [Nguồn: DBI202 Order 18 Q30]
Consider the statement: "Once money is transferred, it won't disappear due to a power outage." In database transactions, which property does this statement exemplify?
\begin{choices}
\choice Durability
\choice Isolation
\choice Consistency
\choice Atomicity
\end{choices}
\textbf{Đáp án:} A
\textbf{Giải thích:} Đáp án theo dữ liệu đề: A — A. Durability.

Lời giải gốc: Tính bền vững (Durability) đảm bảo rằng một khi giao dịch đã hoàn tất thành công, dữ liệu được lưu trữ vĩnh viễn và không bị mất ngay cả khi có sự cố mất điện. ---

Nền tảng kiến thức: Giao dịch là một đơn vị công việc cần thỏa ACID: atomicity (trọn vẹn hoặc hủy), consistency (giữ ràng buộc), isolation (hạn chế ảnh hưởng xen kẽ), durability (đã commit thì bền vững qua sự cố). COMMIT và ROLLBACK kết thúc theo hai hướng khác nhau.

Phân tích: Đối chiếu đúng thuộc tính ACID được hỏi: atomicity là tất-cả-hoặc-không; consistency giữ các ràng buộc; isolation kiểm soát ảnh hưởng giữa giao dịch đồng thời; durability bảo đảm giao dịch đã commit sống sót qua sự cố. COMMIT xác nhận, ROLLBACK hủy thay đổi chưa xác nhận.

Đối chiếu phương án: B. Isolation; C. Consistency; D. Atomicity. Các lựa chọn này không đáp ứng định nghĩa/điều kiện được hỏi; khi ôn, hãy chỉ ra chính xác điểm lệch của từng lựa chọn thay vì học thuộc mỗi chữ cái.
\end{question}

\begin{question}{513}
% [ID gốc: 18.31]
% [Nguồn: DBI202 Order 18 Q31]
Which of the following accurately represents the purpose of a transaction in a database system? Choose two answers.
\begin{choices}
\choice To maintain data integrity and consistency.
\choice To define access permissions
\choice To ensure a sequence of operations is completed fully or not at all
\choice To allow concurrent modifications without restrictions
\end{choices}
\textbf{Đáp án:} AC
\textbf{Giải thích:} Đáp án theo dữ liệu đề: A, C — A. To maintain data integrity and consistency.; C. To ensure a sequence of operations is completed fully or not at all.

Lời giải gốc: Mục tiêu của giao dịch (Transaction) là đảm bảo tính an toàn của CSDL (A - Toàn vẹn, nhất quán). Ngoài ra, nó thực thi nguyên lý "All or Nothing" (C - Tính nguyên tử/Atomicity), tức là mọi lệnh trong giao dịch phải cùng thành công, hoặc không có lệnh nào được lưu nếu xảy ra lỗi.

Nền tảng kiến thức: Giao dịch là một đơn vị công việc cần thỏa ACID: atomicity (trọn vẹn hoặc hủy), consistency (giữ ràng buộc), isolation (hạn chế ảnh hưởng xen kẽ), durability (đã commit thì bền vững qua sự cố). COMMIT và ROLLBACK kết thúc theo hai hướng khác nhau.

Phân tích: Đối chiếu đúng thuộc tính ACID được hỏi: atomicity là tất-cả-hoặc-không; consistency giữ các ràng buộc; isolation kiểm soát ảnh hưởng giữa giao dịch đồng thời; durability bảo đảm giao dịch đã commit sống sót qua sự cố. COMMIT xác nhận, ROLLBACK hủy thay đổi chưa xác nhận.

Đối chiếu phương án: B. To define access permissions; C. To ensure a sequence of operations is completed fully or not at all; D. To allow concurrent modifications without restrictions. Các lựa chọn này không đáp ứng định nghĩa/điều kiện được hỏi; khi ôn, hãy chỉ ra chính xác điểm lệch của từng lựa chọn thay vì học thuộc mỗi chữ cái.
\end{question}

\begin{question}{514}
% [ID gốc: 19.1]
% [Nguồn: DBI202 Order 19 Q1]
A lock in DBMS is a mechanism used to control concurrent access to data. It is a variable associated with a data item that describes the status of the item with respect to possible operations that can be applied to it12. Locks are used as a means of synchronizing access by concurrent transactions to database items2.
QN=39 (8558) Three properties of a lock in DBMS are:
a)Granularity: the size of the lock
\begin{choices}
\choice Granularity: the type of the lock
\choice Mode: the type of the lock
\choice Duration: the time in seconds that the DBMS waits for a lock to be released
\choice (a) and (c) and (d) are true
\choice (a) and (b) and (d) are true
\choice (b) and (c) and (d) are true
\choice (a) and (b) and (c) are true
\end{choices}
\textbf{Đáp án:} G
\textbf{Giải thích:} Đáp án theo dữ liệu đề: A, C, D — A. (a) and (c) and (d) are true; C. (b) and (c) and (d) are true; D. (a) and (b) and (c) are true.

Nền tảng kiến thức: Điều khiển đồng thời giữ lịch thực thi nhiều giao dịch tương đương một thứ tự an toàn. Khóa điều phối đọc/ghi; protocol như two-phase locking hỗ trợ serializability, còn phát hiện deadlock và xử lý các anomaly như lost update/dirty read là phần riêng.

Phân tích: Concurrency control điều phối lịch xen kẽ để tránh anomaly. Shared lock thường cho phép nhiều reader; exclusive lock bảo vệ cập nhật; hai giao dịch chờ tài nguyên vòng tròn có thể deadlock và DBMS thường phá vòng chờ bằng cách hủy/rollback một giao dịch.

Đối chiếu phương án: B. Granularity: the type of the lock; B. (a) and (b) and (d) are true. Các lựa chọn này không đáp ứng định nghĩa/điều kiện được hỏi; khi ôn, hãy chỉ ra chính xác điểm lệch của từng lựa chọn thay vì học thuộc mỗi chữ cái.
\end{question}

\begin{question}{515}
% [ID gốc: 19.2]
% [Nguồn: DBI202 Order 19 Q2]
Two concurrent executing transactions T1 and T2 are allowed to update the same record in an uncontrolled manner. In such a scenario. The following problem may not occur if database system has no concurrency module and allows concurrent execution of above two transactions.Which of the following is correct in this case?
\begin{choices}
\choice Transaction failure
\choice Dirty read problem
\choice Lost update problem
\choice Inconsistent database state
\end{choices}
\textbf{Đáp án:} A
\textbf{Giải thích:} Đáp án theo dữ liệu đề: C — C. Lost update problem.

Nền tảng kiến thức: Điều khiển đồng thời giữ lịch thực thi nhiều giao dịch tương đương một thứ tự an toàn. Khóa điều phối đọc/ghi; protocol như two-phase locking hỗ trợ serializability, còn phát hiện deadlock và xử lý các anomaly như lost update/dirty read là phần riêng.

Phân tích: Concurrency control điều phối lịch xen kẽ để tránh anomaly. Shared lock thường cho phép nhiều reader; exclusive lock bảo vệ cập nhật; hai giao dịch chờ tài nguyên vòng tròn có thể deadlock và DBMS thường phá vòng chờ bằng cách hủy/rollback một giao dịch.

Đối chiếu phương án: A. Transaction failure; B. Dirty read problem; D. Inconsistent database state. Các lựa chọn này không đáp ứng định nghĩa/điều kiện được hỏi; khi ôn, hãy chỉ ra chính xác điểm lệch của từng lựa chọn thay vì học thuộc mỗi chữ cái.
\end{question}

\begin{question}{516}
% [ID gốc: 19.3]
% [Nguồn: DBI202 Order 19 Q3]
According to the technology deployed by the database management system, which of the following is correct?
\begin{choices}
\choice Locks are used to maintain transactional integrity and consistency.
\choice Cursors are used to maintain transactional integrity and consistency.
\choice Procedures are used to maintain transactional integrity and consistency.
\choice Functions are used to maintain transactional integrity and consistency.
\end{choices}
\textbf{Đáp án:} A
\textbf{Giải thích:} Đáp án theo dữ liệu đề: A — A. Locks are used to maintain transactional integrity and consistency..

Nền tảng kiến thức: Điều khiển đồng thời giữ lịch thực thi nhiều giao dịch tương đương một thứ tự an toàn. Khóa điều phối đọc/ghi; protocol như two-phase locking hỗ trợ serializability, còn phát hiện deadlock và xử lý các anomaly như lost update/dirty read là phần riêng.

Phân tích: Đối chiếu đúng thuộc tính ACID được hỏi: atomicity là tất-cả-hoặc-không; consistency giữ các ràng buộc; isolation kiểm soát ảnh hưởng giữa giao dịch đồng thời; durability bảo đảm giao dịch đã commit sống sót qua sự cố. COMMIT xác nhận, ROLLBACK hủy thay đổi chưa xác nhận.

Đối chiếu phương án: B. Cursors are used to maintain transactional integrity and consistency.; C. Procedures are used to maintain transactional integrity and consistency.; D. Functions are used to maintain transactional integrity and consistency.. Các lựa chọn này không đáp ứng định nghĩa/điều kiện được hỏi; khi ôn, hãy chỉ ra chính xác điểm lệch của từng lựa chọn thay vì học thuộc mỗi chữ cái.
\end{question}

\begin{question}{517}
% [ID gốc: 19.4]
% [Nguồn: DBI202 Order 19 Q4]
The deadlock state can be changed back to stable state by using
statement.
\begin{choices}
\choice Delete
\choice Deadlock
\choice Commit
\choice Rollback
\end{choices}
\textbf{Đáp án:} D
\textbf{Giải thích:} Đáp án theo dữ liệu đề: D — D. Rollback.

Nền tảng kiến thức: Điều khiển đồng thời giữ lịch thực thi nhiều giao dịch tương đương một thứ tự an toàn. Khóa điều phối đọc/ghi; protocol như two-phase locking hỗ trợ serializability, còn phát hiện deadlock và xử lý các anomaly như lost update/dirty read là phần riêng.

Phân tích: Concurrency control điều phối lịch xen kẽ để tránh anomaly. Shared lock thường cho phép nhiều reader; exclusive lock bảo vệ cập nhật; hai giao dịch chờ tài nguyên vòng tròn có thể deadlock và DBMS thường phá vòng chờ bằng cách hủy/rollback một giao dịch.

Đối chiếu phương án: A. Delete; B. Deadlock; C. Commit. Các lựa chọn này không đáp ứng định nghĩa/điều kiện được hỏi; khi ôn, hãy chỉ ra chính xác điểm lệch của từng lựa chọn thay vì học thuộc mỗi chữ cái.
\end{question}

\begin{question}{518}
% [ID gốc: 19.5]
% [Nguồn: DBI202 Order 19 Q5]
Two concurrently executing transactions T1 and T2 are allowed to update the same record in an uncontrolled manner. If the database system has no concurrency module, which of the following problems may occur?
\begin{choices}
\choice Transaction failure
\choice Dirty read problem
\choice Lost update problem
\choice Inconsistent database state
\end{choices}
\textbf{Đáp án:} C
\textbf{Giải thích:} Đáp án theo dữ liệu đề: C — C. Lost update problem.

Lời giải gốc: Khi hai giao dịch cùng cập nhật một bản ghi mà không được kiểm soát đồng thời, giao dịch ghi sau sẽ ghi đè lên kết quả của giao dịch ghi trước, dẫn đến vấn đề mất cập nhật (Lost update problem).

Nền tảng kiến thức: Điều khiển đồng thời giữ lịch thực thi nhiều giao dịch tương đương một thứ tự an toàn. Khóa điều phối đọc/ghi; protocol như two-phase locking hỗ trợ serializability, còn phát hiện deadlock và xử lý các anomaly như lost update/dirty read là phần riêng.

Phân tích: Concurrency control điều phối lịch xen kẽ để tránh anomaly. Shared lock thường cho phép nhiều reader; exclusive lock bảo vệ cập nhật; hai giao dịch chờ tài nguyên vòng tròn có thể deadlock và DBMS thường phá vòng chờ bằng cách hủy/rollback một giao dịch.

Đối chiếu phương án: A. Transaction failure; B. Dirty read problem; D. Inconsistent database state. Các lựa chọn này không đáp ứng định nghĩa/điều kiện được hỏi; khi ôn, hãy chỉ ra chính xác điểm lệch của từng lựa chọn thay vì học thuộc mỗi chữ cái.
\end{question}

\begin{question}{519}
% [ID gốc: 19.6]
% [Nguồn: DBI202 Order 19 Q6]
The deadlock state can be changed back to a stable state by using which statement?
\begin{choices}
\choice DELETE
\choice DEADLOCK
\choice COMMIT
\choice ROLLBACK
\end{choices}
\textbf{Đáp án:} D
\textbf{Giải thích:} Đáp án theo dữ liệu đề: D — D. ROLLBACK.

Lời giải gốc: Khi xảy ra bế tắc (deadlock), hệ quản trị cơ sở dữ liệu thường hủy bỏ một giao dịch (qua ROLLBACK hoặc hủy tiến trình) để giải phóng tài nguyên và đưa hệ thống về trạng thái ổn định.

Nền tảng kiến thức: Điều khiển đồng thời giữ lịch thực thi nhiều giao dịch tương đương một thứ tự an toàn. Khóa điều phối đọc/ghi; protocol như two-phase locking hỗ trợ serializability, còn phát hiện deadlock và xử lý các anomaly như lost update/dirty read là phần riêng.

Phân tích: Concurrency control điều phối lịch xen kẽ để tránh anomaly. Shared lock thường cho phép nhiều reader; exclusive lock bảo vệ cập nhật; hai giao dịch chờ tài nguyên vòng tròn có thể deadlock và DBMS thường phá vòng chờ bằng cách hủy/rollback một giao dịch.

Đối chiếu phương án: A. DELETE; B. DEADLOCK; C. COMMIT. Các lựa chọn này không đáp ứng định nghĩa/điều kiện được hỏi; khi ôn, hãy chỉ ra chính xác điểm lệch của từng lựa chọn thay vì học thuộc mỗi chữ cái.
\end{question}

\begin{question}{520}
% [ID gốc: 19.9]
% [Nguồn: DBI202 Order 19 Q9]
In database transactions, which isolation level is most restrictive?
\begin{choices}
\choice READ UNCOMMITTED
\choice READ COMMITTED
\choice REPEATABLE READ
\choice SERIALIZABLE
\end{choices}
\textbf{Đáp án:} D
\textbf{Giải thích:} Đáp án theo dữ liệu đề: D — D. SERIALIZABLE.

Lời giải gốc: Cấp độ cô lập SERIALIZABLE là nghiêm ngặt, hạn chế và an toàn nhất. Nó ép buộc các giao dịch phải thực thi tuần tự tuần tự như thể không có sự đồng thời, giúp ngăn chặn toàn bộ các lỗi đồng thời (Dirty read, Non-repeatable read, Phantom read) nhưng làm giảm hiệu năng hệ thống.

Nền tảng kiến thức: Điều khiển đồng thời giữ lịch thực thi nhiều giao dịch tương đương một thứ tự an toàn. Khóa điều phối đọc/ghi; protocol như two-phase locking hỗ trợ serializability, còn phát hiện deadlock và xử lý các anomaly như lost update/dirty read là phần riêng.

Phân tích: Đối chiếu đúng thuộc tính ACID được hỏi: atomicity là tất-cả-hoặc-không; consistency giữ các ràng buộc; isolation kiểm soát ảnh hưởng giữa giao dịch đồng thời; durability bảo đảm giao dịch đã commit sống sót qua sự cố. COMMIT xác nhận, ROLLBACK hủy thay đổi chưa xác nhận.

Đối chiếu phương án: A. READ UNCOMMITTED; B. READ COMMITTED; C. REPEATABLE READ. Các lựa chọn này không đáp ứng định nghĩa/điều kiện được hỏi; khi ôn, hãy chỉ ra chính xác điểm lệch của từng lựa chọn thay vì học thuộc mỗi chữ cái.
\end{question}

\begin{question}{521}
% [ID gốc: 19.10]
% [Nguồn: DBI202 Order 19 Q10]
The deadlock state can be changed back to a stable state by using _ _ _ _ _ statement.
\begin{choices}
\choice Delete
\choice Deadlock
\choice Commit
\choice Rollback
\end{choices}
\textbf{Đáp án:} D
\textbf{Giải thích:} Đáp án theo dữ liệu đề: D — D. Rollback.

Lời giải gốc: Khi xảy ra deadlock (bế tắc), hệ thống sẽ hủy bỏ (Rollback) một giao dịch "nạn nhân" để đưa hệ thống trở lại trạng thái an toàn/ổn định.

Nền tảng kiến thức: Điều khiển đồng thời giữ lịch thực thi nhiều giao dịch tương đương một thứ tự an toàn. Khóa điều phối đọc/ghi; protocol như two-phase locking hỗ trợ serializability, còn phát hiện deadlock và xử lý các anomaly như lost update/dirty read là phần riêng.

Phân tích: Concurrency control điều phối lịch xen kẽ để tránh anomaly. Shared lock thường cho phép nhiều reader; exclusive lock bảo vệ cập nhật; hai giao dịch chờ tài nguyên vòng tròn có thể deadlock và DBMS thường phá vòng chờ bằng cách hủy/rollback một giao dịch.

Đối chiếu phương án: A. Delete; B. Deadlock; C. Commit. Các lựa chọn này không đáp ứng định nghĩa/điều kiện được hỏi; khi ôn, hãy chỉ ra chính xác điểm lệch của từng lựa chọn thay vì học thuộc mỗi chữ cái.
\end{question}

\begin{question}{522}
% [ID gốc: 23.1]
% [Nguồn: DBI202 Order 23 Q1]
Transaction processing is associated with everything below except.
\begin{choices}
\choice Producing detail, summery, or exception reports
\choice Recording a business activity
\choice Confirming an action or triggering a response
\choice Maintaining data
\end{choices}
\textbf{Đáp án:} A
\textbf{Giải thích:} Đáp án theo dữ liệu đề: A — A. Producing detail, summery, or exception reports.

Nền tảng kiến thức: Decision support và OLAP phục vụ phân tích, tổng hợp và báo cáo trên khối lượng dữ liệu lớn; chúng khác với OLTP vốn ưu tiên ghi nhận các giao dịch ngắn, thường xuyên. Kho dữ liệu, chiều phân tích và phép tổng hợp hỗ trợ câu hỏi quản trị.

Phân tích: Đối chiếu đúng thuộc tính ACID được hỏi: atomicity là tất-cả-hoặc-không; consistency giữ các ràng buộc; isolation kiểm soát ảnh hưởng giữa giao dịch đồng thời; durability bảo đảm giao dịch đã commit sống sót qua sự cố. COMMIT xác nhận, ROLLBACK hủy thay đổi chưa xác nhận.

Đối chiếu phương án: B. Recording a business activity; C. Confirming an action or triggering a response; D. Maintaining data. Các lựa chọn này không đáp ứng định nghĩa/điều kiện được hỏi; khi ôn, hãy chỉ ra chính xác điểm lệch của từng lựa chọn thay vì học thuộc mỗi chữ cái.
\end{question}


\newpage
\section*{Chương 8: Views and Indexes}

\begin{question}{523}
% [ID gốc: 1.11]
% [Nguồn: DBI202 Order 1 Q11]
In the three-tier architecture, the database tier's function is to
\begin{choices}
\choice All of the others.
\choice Execute the business logic of the organization operating the database.
\choice Manage the interactions with the user.
\choice Execute queries that are requested from the application tier.
\end{choices}
\textbf{Đáp án:} D
\textbf{Giải thích:} Đáp án theo dữ liệu đề: D — D. Execute queries that are requested from the application tier..

Nền tảng kiến thức: Chương này phân biệt dữ liệu, cơ sở dữ liệu, hệ quản trị CSDL (DBMS), mô hình dữ liệu, lược đồ và trạng thái dữ liệu. Cơ sở dữ liệu lưu dữ liệu có liên hệ; DBMS là phần mềm định nghĩa, truy vấn, cập nhật và bảo vệ dữ liệu đó.

Phân tích: Kiến trúc ba tầng tách giao diện (presentation), xử lý nghiệp vụ (application) và lưu trữ/truy vấn dữ liệu (database). Xác định tầng được hỏi đang nhận yêu cầu từ tầng nào và thực hiện trách nhiệm nào; không gán xử lý giao diện hay nghiệp vụ cho tầng dữ liệu.

Đối chiếu phương án: A. All of the others.; B. Execute the business logic of the organization operating the database.; C. Manage the interactions with the user.. Các lựa chọn này không đáp ứng định nghĩa/điều kiện được hỏi; khi ôn, hãy chỉ ra chính xác điểm lệch của từng lựa chọn thay vì học thuộc mỗi chữ cái.
\end{question}

\begin{question}{524}
% [ID gốc: 5.19]
% [Nguồn: DBI202 Order 5 Q19]
Choose the correct statement.
\begin{choices}
\choice You can remove a trigger by dropping it or by dropping the trigger table.
\choice The syntax to remove a trigger is: DROP TRIGGER
\choice Use ALTER TRIGGER to change the definition of a trigger
\choice All of the others.
\end{choices}
\textbf{Đáp án:} D
\textbf{Giải thích:} Đáp án theo dữ liệu đề: D — D. All of the others..

Nền tảng kiến thức: SQL tách truy vấn dữ liệu (DML) khỏi định nghĩa cấu trúc (DDL) và quyền truy cập (DCL). Với SELECT, cần xác định bảng nguồn, điều kiện lọc, phép kết hợp, nhóm và điều kiện nhóm; các ràng buộc, view, thủ tục và trigger đảm nhiệm các vai trò khác nhau.

Phân tích: Với lựa chọn tổng hợp, không chọn theo cảm tính: kiểm tra từng phát biểu độc lập. “All of the others/above” chỉ đúng khi mọi phát biểu còn lại đúng; “None” chỉ đúng khi không phát biểu nào trong nhóm thỏa yêu cầu. Đáp án lưu cần được đọc cùng các mệnh đề thành phần.

Đối chiếu phương án: A. You can remove a trigger by dropping it or by dropping the trigger table.; B. The syntax to remove a trigger is: DROP TRIGGER; C. Use ALTER TRIGGER to change the definition of a trigger. Các lựa chọn này không đáp ứng định nghĩa/điều kiện được hỏi; khi ôn, hãy chỉ ra chính xác điểm lệch của từng lựa chọn thay vì học thuộc mỗi chữ cái.
\end{question}

\begin{question}{525}
% [ID gốc: 5.36]
% [Nguồn: DBI202 Order 5 Q36]
In PSM (Persistent Stored Modules), the difference between a stored procedure and a function is that:
\begin{choices}
\choice A function has the return statement.
\choice We can declare local variables in a function.
\choice Loops are not allowed in a function.
\choice All of the others.
\end{choices}
\textbf{Đáp án:} A
\textbf{Giải thích:} Đáp án theo dữ liệu đề: A — A. A function has the return statement..

Nền tảng kiến thức: SQL tách truy vấn dữ liệu (DML) khỏi định nghĩa cấu trúc (DDL) và quyền truy cập (DCL). Với SELECT, cần xác định bảng nguồn, điều kiện lọc, phép kết hợp, nhóm và điều kiện nhóm; các ràng buộc, view, thủ tục và trigger đảm nhiệm các vai trò khác nhau.

Phân tích: Phép tập hợp trên quan hệ cần so sánh các tuple theo thuộc tính tương ứng; union lấy hợp, intersection lấy phần chung và difference lấy phần có ở vế trái mà không có ở vế phải. Trước tiên kiểm tra tính tương thích schema, sau đó áp phép theo đúng hướng/toán tử.

Đối chiếu phương án: B. We can declare local variables in a function.; C. Loops are not allowed in a function.; D. All of the others.. Các lựa chọn này không đáp ứng định nghĩa/điều kiện được hỏi; khi ôn, hãy chỉ ra chính xác điểm lệch của từng lựa chọn thay vì học thuộc mỗi chữ cái.
\end{question}

\begin{question}{526}
% [ID gốc: 5.44]
% [Nguồn: DBI202 Order 5 Q44]
Why we should use stored procedures?
\begin{choices}
\choice They allow faster execution
\choice They can reduce network traffic
\choice They allow modular programming
\choice They can be used as a security mechanism
\choice All of the others
\end{choices}
\textbf{Đáp án:} E
\textbf{Giải thích:} Đáp án theo dữ liệu đề: E — E. All of the others.

Nền tảng kiến thức: SQL tách truy vấn dữ liệu (DML) khỏi định nghĩa cấu trúc (DDL) và quyền truy cập (DCL). Với SELECT, cần xác định bảng nguồn, điều kiện lọc, phép kết hợp, nhóm và điều kiện nhóm; các ràng buộc, view, thủ tục và trigger đảm nhiệm các vai trò khác nhau.

Phân tích: Stored procedure là tập lệnh SQL có tên lưu trong CSDL, được gọi để thực thi một thủ tục; nó có thể nhận tham số và chứa nhiều câu lệnh. Đừng nhầm với trigger tự chạy khi có sự kiện hoặc function có quy tắc trả về giá trị.

Đối chiếu phương án: A. They allow faster execution; B. They can reduce network traffic; C. They allow modular programming; D. They can be used as a security mechanism. Các lựa chọn này không đáp ứng định nghĩa/điều kiện được hỏi; khi ôn, hãy chỉ ra chính xác điểm lệch của từng lựa chọn thay vì học thuộc mỗi chữ cái.
\end{question}

\begin{question}{527}
% [ID gốc: 5.58]
% [Nguồn: DBI202 Order 5 Q58]
Exception handler in PSM is defined as follows:
DECLARE  HANDLER FOR condition list> 
The  clause can be:
\begin{choices}
\choice UNDO
\choice EXIT
\choice All of the others
\choice CONTINUE
\end{choices}
\textbf{Đáp án:} C
\textbf{Giải thích:} Đáp án theo dữ liệu đề: C — C. All of the others.

Nền tảng kiến thức: SQL tách truy vấn dữ liệu (DML) khỏi định nghĩa cấu trúc (DDL) và quyền truy cập (DCL). Với SELECT, cần xác định bảng nguồn, điều kiện lọc, phép kết hợp, nhóm và điều kiện nhóm; các ràng buộc, view, thủ tục và trigger đảm nhiệm các vai trò khác nhau.

Phân tích: Với lựa chọn tổng hợp, không chọn theo cảm tính: kiểm tra từng phát biểu độc lập. “All of the others/above” chỉ đúng khi mọi phát biểu còn lại đúng; “None” chỉ đúng khi không phát biểu nào trong nhóm thỏa yêu cầu. Đáp án lưu cần được đọc cùng các mệnh đề thành phần.

Đối chiếu phương án: A. UNDO; B. EXIT; D. CONTINUE. Các lựa chọn này không đáp ứng định nghĩa/điều kiện được hỏi; khi ôn, hãy chỉ ra chính xác điểm lệch của từng lựa chọn thay vì học thuộc mỗi chữ cái.
\end{question}

\begin{question}{528}
% [ID gốc: 5.65]
% [Nguồn: DBI202 Order 5 Q65]
In SQL, what is the primary advantage of using stored procedures?
\begin{choices}
\choice They allow for the creation of tables
\choice They enhance security by restricting access to data
\choice They provide a way to encapsulate and reuse a set of SQL statements
\choice They are more efficient than regular SQL queries
\end{choices}
\textbf{Đáp án:} C
\textbf{Giải thích:} Đáp án theo dữ liệu đề: C — C. They provide a way to encapsulate and reuse a set of SQL statements.

Nền tảng kiến thức: SQL tách truy vấn dữ liệu (DML) khỏi định nghĩa cấu trúc (DDL) và quyền truy cập (DCL). Với SELECT, cần xác định bảng nguồn, điều kiện lọc, phép kết hợp, nhóm và điều kiện nhóm; các ràng buộc, view, thủ tục và trigger đảm nhiệm các vai trò khác nhau.

Phân tích: Stored procedure là tập lệnh SQL có tên lưu trong CSDL, được gọi để thực thi một thủ tục; nó có thể nhận tham số và chứa nhiều câu lệnh. Đừng nhầm với trigger tự chạy khi có sự kiện hoặc function có quy tắc trả về giá trị.

Đối chiếu phương án: A. They allow for the creation of tables; B. They enhance security by restricting access to data; D. They are more efficient than regular SQL queries. Các lựa chọn này không đáp ứng định nghĩa/điều kiện được hỏi; khi ôn, hãy chỉ ra chính xác điểm lệch của từng lựa chọn thay vì học thuộc mỗi chữ cái.
\end{question}

\begin{question}{529}
% [ID gốc: 5.66]
% [Nguồn: DBI202 Order 5 Q66]
Which type of user-defined function in SQL returns a single value?
\begin{choices}
\choice Scalar function
\choice Table-valued function
\choice Inline function
\choice Aggregate function
\end{choices}
\textbf{Đáp án:} A
\textbf{Giải thích:} Đáp án theo dữ liệu đề: A — A. Scalar function.

Nền tảng kiến thức: SQL tách truy vấn dữ liệu (DML) khỏi định nghĩa cấu trúc (DDL) và quyền truy cập (DCL). Với SELECT, cần xác định bảng nguồn, điều kiện lọc, phép kết hợp, nhóm và điều kiện nhóm; các ràng buộc, view, thủ tục và trigger đảm nhiệm các vai trò khác nhau.

Phân tích: Aggregate function tính một kết quả từ một tập nhiều hàng/giá trị, ví dụ COUNT, SUM, AVG, MIN và MAX. Trong truy vấn nhóm, GROUP BY xác định từng tập đầu vào; HAVING lọc các nhóm sau khi tính tổng hợp.

Đối chiếu phương án: B. Table-valued function; C. Inline function; D. Aggregate function. Các lựa chọn này không đáp ứng định nghĩa/điều kiện được hỏi; khi ôn, hãy chỉ ra chính xác điểm lệch của từng lựa chọn thay vì học thuộc mỗi chữ cái.
\end{question}

\begin{question}{530}
% [ID gốc: 5.70]
% [Nguồn: DBI202 Order 5 Q70]
Which of the following is true about stored procedure in MS-SQL Server?
\begin{choices}
\choice They include procedural and SQL statements.
\choice They are stored in Rules folder
\choice They do not need to have a unique name
\choice They are same as void function
\end{choices}
\textbf{Đáp án:} A
\textbf{Giải thích:} Đáp án theo dữ liệu đề: A — A. They include procedural and SQL statements..

Nền tảng kiến thức: SQL tách truy vấn dữ liệu (DML) khỏi định nghĩa cấu trúc (DDL) và quyền truy cập (DCL). Với SELECT, cần xác định bảng nguồn, điều kiện lọc, phép kết hợp, nhóm và điều kiện nhóm; các ràng buộc, view, thủ tục và trigger đảm nhiệm các vai trò khác nhau.

Phân tích: Stored procedure là tập lệnh SQL có tên lưu trong CSDL, được gọi để thực thi một thủ tục; nó có thể nhận tham số và chứa nhiều câu lệnh. Đừng nhầm với trigger tự chạy khi có sự kiện hoặc function có quy tắc trả về giá trị.

Đối chiếu phương án: B. They are stored in Rules folder; C. They do not need to have a unique name; D. They are same as void function. Các lựa chọn này không đáp ứng định nghĩa/điều kiện được hỏi; khi ôn, hãy chỉ ra chính xác điểm lệch của từng lựa chọn thay vì học thuộc mỗi chữ cái.
\end{question}

\begin{question}{531}
% [ID gốc: 5.71]
% [Nguồn: DBI202 Order 5 Q71]
______ A variable is local, and its value is not preserved by the DBMS after a run-ning of the function or
procedure.
\begin{choices}
\choice DECLARE
\choice RETURN
\choice CALL
\choice EXEC
\end{choices}
\textbf{Đáp án:} A
\textbf{Giải thích:} Đáp án theo dữ liệu đề: A — A. DECLARE.

Nền tảng kiến thức: SQL tách truy vấn dữ liệu (DML) khỏi định nghĩa cấu trúc (DDL) và quyền truy cập (DCL). Với SELECT, cần xác định bảng nguồn, điều kiện lọc, phép kết hợp, nhóm và điều kiện nhóm; các ràng buộc, view, thủ tục và trigger đảm nhiệm các vai trò khác nhau.

Phân tích: Stored procedure là tập lệnh SQL có tên lưu trong CSDL, được gọi để thực thi một thủ tục; nó có thể nhận tham số và chứa nhiều câu lệnh. Đừng nhầm với trigger tự chạy khi có sự kiện hoặc function có quy tắc trả về giá trị.

Đối chiếu phương án: B. RETURN; C. CALL; D. EXEC. Các lựa chọn này không đáp ứng định nghĩa/điều kiện được hỏi; khi ôn, hãy chỉ ra chính xác điểm lệch của từng lựa chọn thay vì học thuộc mỗi chữ cái.
\end{question}

\begin{question}{532}
% [ID gốc: 5.72]
% [Nguồn: DBI202 Order 5 Q72]
What is NOT true about the cursor in SQL?
\begin{choices}
\choice It is considered the best practice to use when we want insert, update, or delete data in a table of SQL Server.
\choice The cursor allows users to process data from a result set, one row at a time.
\choice Cursors are an alternative to commands, which operate on all rows in a result set at the same time.
\choice Unlike commands, cursors can be used to update data on a row-by-row basis.
\end{choices}
\textbf{Đáp án:} A
\textbf{Giải thích:} Đáp án theo dữ liệu đề: A — A. It is considered the best practice to use when we want insert, update, or delete data in a table of SQL Server..

Nền tảng kiến thức: SQL tách truy vấn dữ liệu (DML) khỏi định nghĩa cấu trúc (DDL) và quyền truy cập (DCL). Với SELECT, cần xác định bảng nguồn, điều kiện lọc, phép kết hợp, nhóm và điều kiện nhóm; các ràng buộc, view, thủ tục và trigger đảm nhiệm các vai trò khác nhau.

Phân tích: Cursor cho phép duyệt kết quả theo từng hàng và chỉ nên dùng khi xử lý từng bản ghi là cần thiết hoặc không thể biểu diễn gọn bằng thao tác tập hợp. Với INSERT/UPDATE/DELETE hàng loạt, truy vấn set-based thường đơn giản và hiệu quả hơn.

Đối chiếu phương án: B. The cursor allows users to process data from a result set, one row at a time.; C. Cursors are an alternative to commands, which operate on all rows in a result set at the same…; D. Unlike commands, cursors can be used to update data on a row-by-row basis.. Các lựa chọn này không đáp ứng định nghĩa/điều kiện được hỏi; khi ôn, hãy chỉ ra chính xác điểm lệch của từng lựa chọn thay vì học thuộc mỗi chữ cái.
\end{question}

\begin{question}{533}
% [ID gốc: 5.73]
% [Nguồn: DBI202 Order 5 Q73]
How can you delete a user-defined function in SQL?
\begin{choices}
\choice DELETE FUNCTION function_name
\choice REMOVE FUNCTION function_name
\choice DROP FUNCTION function_name
\choice ERASE FUNCTION function_name
\end{choices}
\textbf{Đáp án:} C
\textbf{Giải thích:} Đáp án theo dữ liệu đề: C — C. DROP FUNCTION function_name.

Nền tảng kiến thức: SQL tách truy vấn dữ liệu (DML) khỏi định nghĩa cấu trúc (DDL) và quyền truy cập (DCL). Với SELECT, cần xác định bảng nguồn, điều kiện lọc, phép kết hợp, nhóm và điều kiện nhóm; các ràng buộc, view, thủ tục và trigger đảm nhiệm các vai trò khác nhau.

Phân tích: DROP FUNCTION xóa đối tượng hàm do người dùng định nghĩa; tên hàm/schema cần chỉ đúng đối tượng. ALTER FUNCTION sửa định nghĩa còn EXECUTE gọi hàm/thủ tục, nên không dùng để loại bỏ hàm.

Đối chiếu phương án: A. DELETE FUNCTION function_name; B. REMOVE FUNCTION function_name; D. ERASE FUNCTION function_name. Các lựa chọn này không đáp ứng định nghĩa/điều kiện được hỏi; khi ôn, hãy chỉ ra chính xác điểm lệch của từng lựa chọn thay vì học thuộc mỗi chữ cái.
\end{question}

\begin{question}{534}
% [ID gốc: 5.98]
% [Nguồn: DBI202 Order 5 Q98]
Triggers are stored blocks of code that have to be called in order to operate.
\begin{choices}
\choice True
\choice False
\end{choices}
\textbf{Đáp án:} B
\textbf{Giải thích:} Đáp án theo dữ liệu đề: B — B. False.

Nền tảng kiến thức: SQL tách truy vấn dữ liệu (DML) khỏi định nghĩa cấu trúc (DDL) và quyền truy cập (DCL). Với SELECT, cần xác định bảng nguồn, điều kiện lọc, phép kết hợp, nhóm và điều kiện nhóm; các ràng buộc, view, thủ tục và trigger đảm nhiệm các vai trò khác nhau.

Phân tích: Phân biệt tệp dữ liệu chứa trang/bảng với tệp transaction log ghi lại thay đổi phục vụ commit và recovery. Một trang/khối là đơn vị I/O cơ bản; lựa chọn tổ chức vật lý ảnh hưởng số lần đọc ghi và cách tìm bản ghi.

Đối chiếu phương án: A. True. Các lựa chọn này không đáp ứng định nghĩa/điều kiện được hỏi; khi ôn, hãy chỉ ra chính xác điểm lệch của từng lựa chọn thay vì học thuộc mỗi chữ cái.
\end{question}

\begin{question}{535}
% [ID gốc: 5.99]
% [Nguồn: DBI202 Order 5 Q99]
In SQL, when does a BEFORE INSERT trigger fire
\begin{choices}
\choice Before the execution of a SELECT statement
\choice Before a new record is inserted into a table
\choice Before a transaction is committed
\choice Before a table is dropped
\end{choices}
\textbf{Đáp án:} B
\textbf{Giải thích:} Đáp án theo dữ liệu đề: B — B. Before a new record is inserted into a table.

Nền tảng kiến thức: SQL tách truy vấn dữ liệu (DML) khỏi định nghĩa cấu trúc (DDL) và quyền truy cập (DCL). Với SELECT, cần xác định bảng nguồn, điều kiện lọc, phép kết hợp, nhóm và điều kiện nhóm; các ràng buộc, view, thủ tục và trigger đảm nhiệm các vai trò khác nhau.

Phân tích: Trigger là hành động do DBMS tự kích hoạt khi một sự kiện dữ liệu như INSERT, UPDATE hoặc DELETE xảy ra trên đối tượng đã khai báo. Nó khác stored procedure vốn thường được gọi tường minh; đáp án đúng cần thể hiện đúng sự kiện kích hoạt.

Đối chiếu phương án: A. Before the execution of a SELECT statement; C. Before a transaction is committed; D. Before a table is dropped. Các lựa chọn này không đáp ứng định nghĩa/điều kiện được hỏi; khi ôn, hãy chỉ ra chính xác điểm lệch của từng lựa chọn thay vì học thuộc mỗi chữ cái.
\end{question}

\begin{question}{536}
% [ID gốc: 5.100]
% [Nguồn: DBI202 Order 5 Q100]
In SQL, what is the purpose of the CLOSE statement in relation to cursors?
\begin{choices}
\choice To close the database connection
\choice To close the cursor and release associated resources
\choice To close the transaction
\choice To close the result set
\end{choices}
\textbf{Đáp án:} B
\textbf{Giải thích:} Đáp án theo dữ liệu đề: B — B. To close the cursor and release associated resources.

Nền tảng kiến thức: SQL tách truy vấn dữ liệu (DML) khỏi định nghĩa cấu trúc (DDL) và quyền truy cập (DCL). Với SELECT, cần xác định bảng nguồn, điều kiện lọc, phép kết hợp, nhóm và điều kiện nhóm; các ràng buộc, view, thủ tục và trigger đảm nhiệm các vai trò khác nhau.

Phân tích: Cursor cho phép duyệt kết quả theo từng hàng và chỉ nên dùng khi xử lý từng bản ghi là cần thiết hoặc không thể biểu diễn gọn bằng thao tác tập hợp. Với INSERT/UPDATE/DELETE hàng loạt, truy vấn set-based thường đơn giản và hiệu quả hơn.

Đối chiếu phương án: A. To close the database connection; C. To close the transaction; D. To close the result set. Các lựa chọn này không đáp ứng định nghĩa/điều kiện được hỏi; khi ôn, hãy chỉ ra chính xác điểm lệch của từng lựa chọn thay vì học thuộc mỗi chữ cái.
\end{question}

\begin{question}{537}
% [ID gốc: 5.101]
% [Nguồn: DBI202 Order 5 Q101]
Which SQL statement is used to declare a cursor?
\begin{choices}
\choice OPEN
\choice DECLARE
\choice CREATE CURSOR
\choice BEGIN CURSOR
\end{choices}
\textbf{Đáp án:} B
\textbf{Giải thích:} Đáp án theo dữ liệu đề: B — B. DECLARE.

Nền tảng kiến thức: SQL tách truy vấn dữ liệu (DML) khỏi định nghĩa cấu trúc (DDL) và quyền truy cập (DCL). Với SELECT, cần xác định bảng nguồn, điều kiện lọc, phép kết hợp, nhóm và điều kiện nhóm; các ràng buộc, view, thủ tục và trigger đảm nhiệm các vai trò khác nhau.

Phân tích: Cursor cho phép duyệt kết quả theo từng hàng và chỉ nên dùng khi xử lý từng bản ghi là cần thiết hoặc không thể biểu diễn gọn bằng thao tác tập hợp. Với INSERT/UPDATE/DELETE hàng loạt, truy vấn set-based thường đơn giản và hiệu quả hơn.

Đối chiếu phương án: A. OPEN; C. CREATE CURSOR; D. BEGIN CURSOR. Các lựa chọn này không đáp ứng định nghĩa/điều kiện được hỏi; khi ôn, hãy chỉ ra chính xác điểm lệch của từng lựa chọn thay vì học thuộc mỗi chữ cái.
\end{question}

\begin{question}{538}
% [ID gốc: 5.102]
% [Nguồn: DBI202 Order 5 Q102]
How is User-defined function different from Store Procedure?
\begin{choices}
\choice User-defined functions do not have input parameters, Store Procedures do
\choice User-defined functions cannot contain transactions, Store Procedures can
\choice Store Procedure does not have return results, User-defined function does
\choice All of the mentioned answers are incorrect
\end{choices}
\textbf{Đáp án:} B
\textbf{Giải thích:} Đáp án theo dữ liệu đề: B — B. User-defined functions cannot contain transactions, Store Procedures can.

Nền tảng kiến thức: SQL tách truy vấn dữ liệu (DML) khỏi định nghĩa cấu trúc (DDL) và quyền truy cập (DCL). Với SELECT, cần xác định bảng nguồn, điều kiện lọc, phép kết hợp, nhóm và điều kiện nhóm; các ràng buộc, view, thủ tục và trigger đảm nhiệm các vai trò khác nhau.

Phân tích: Stored procedure là tập lệnh SQL có tên lưu trong CSDL, được gọi để thực thi một thủ tục; nó có thể nhận tham số và chứa nhiều câu lệnh. Đừng nhầm với trigger tự chạy khi có sự kiện hoặc function có quy tắc trả về giá trị.

Đối chiếu phương án: A. User-defined functions do not have input parameters, Store Procedures do; C. Store Procedure does not have return results, User-defined function does; D. All of the mentioned answers are incorrect. Các lựa chọn này không đáp ứng định nghĩa/điều kiện được hỏi; khi ôn, hãy chỉ ra chính xác điểm lệch của từng lựa chọn thay vì học thuộc mỗi chữ cái.
\end{question}

\begin{question}{539}
% [ID gốc: 5.103]
% [Nguồn: DBI202 Order 5 Q103]
How to execute a procedure
\begin{choices}
\choice execute name_sp
\choice exe name_sp
\choice call name_sp
\choice All the answers above
\choice a and b are correct
\end{choices}
\textbf{Đáp án:} A
\textbf{Giải thích:} Đáp án theo dữ liệu đề: A — A. execute name_sp.

Nền tảng kiến thức: SQL tách truy vấn dữ liệu (DML) khỏi định nghĩa cấu trúc (DDL) và quyền truy cập (DCL). Với SELECT, cần xác định bảng nguồn, điều kiện lọc, phép kết hợp, nhóm và điều kiện nhóm; các ràng buộc, view, thủ tục và trigger đảm nhiệm các vai trò khác nhau.

Phân tích: Stored procedure là tập lệnh SQL có tên lưu trong CSDL, được gọi để thực thi một thủ tục; nó có thể nhận tham số và chứa nhiều câu lệnh. Đừng nhầm với trigger tự chạy khi có sự kiện hoặc function có quy tắc trả về giá trị.

Đối chiếu phương án: B. exe name_sp; C. call name_sp; D. All the answers above; E. a and b are correct. Các lựa chọn này không đáp ứng định nghĩa/điều kiện được hỏi; khi ôn, hãy chỉ ra chính xác điểm lệch của từng lựa chọn thay vì học thuộc mỗi chữ cái.
\end{question}

\begin{question}{540}
% [ID gốc: 5.104]
% [Nguồn: DBI202 Order 5 Q104]
A stored procedure in SQL is a
\begin{choices}
\choice Block of functions
\choice Group of Transact-SQL statements compiled into a single execution plan
\choice Group of distinct SQL statements
\choice None of the mentioned
\end{choices}
\textbf{Đáp án:} B
\textbf{Giải thích:} Đáp án theo dữ liệu đề: B — B. Group of Transact-SQL statements compiled into a single execution plan.

Nền tảng kiến thức: SQL tách truy vấn dữ liệu (DML) khỏi định nghĩa cấu trúc (DDL) và quyền truy cập (DCL). Với SELECT, cần xác định bảng nguồn, điều kiện lọc, phép kết hợp, nhóm và điều kiện nhóm; các ràng buộc, view, thủ tục và trigger đảm nhiệm các vai trò khác nhau.

Phân tích: Stored procedure là tập lệnh SQL có tên lưu trong CSDL, được gọi để thực thi một thủ tục; nó có thể nhận tham số và chứa nhiều câu lệnh. Đừng nhầm với trigger tự chạy khi có sự kiện hoặc function có quy tắc trả về giá trị.

Đối chiếu phương án: A. Block of functions; C. Group of distinct SQL statements; D. None of the mentioned. Các lựa chọn này không đáp ứng định nghĩa/điều kiện được hỏi; khi ôn, hãy chỉ ra chính xác điểm lệch của từng lựa chọn thay vì học thuộc mỗi chữ cái.
\end{question}

\begin{question}{541}
% [ID gốc: 5.105]
% [Nguồn: DBI202 Order 5 Q105]
When is it recommended to use a cursor in SQL?
\begin{choices}
\choice Always, as cursors are the most efficient way to process data
\choice When processing data row by row is necessary or unavoidable
\choice Only when creating complex joins between tables
\choice Never, as set-based operations are always more efficient
\end{choices}
\textbf{Đáp án:} B
\textbf{Giải thích:} Đáp án theo dữ liệu đề: B — B. When processing data row by row is necessary or unavoidable.

Nền tảng kiến thức: SQL tách truy vấn dữ liệu (DML) khỏi định nghĩa cấu trúc (DDL) và quyền truy cập (DCL). Với SELECT, cần xác định bảng nguồn, điều kiện lọc, phép kết hợp, nhóm và điều kiện nhóm; các ràng buộc, view, thủ tục và trigger đảm nhiệm các vai trò khác nhau.

Phân tích: Cursor cho phép duyệt kết quả theo từng hàng và chỉ nên dùng khi xử lý từng bản ghi là cần thiết hoặc không thể biểu diễn gọn bằng thao tác tập hợp. Với INSERT/UPDATE/DELETE hàng loạt, truy vấn set-based thường đơn giản và hiệu quả hơn.

Đối chiếu phương án: A. Always, as cursors are the most efficient way to process data; C. Only when creating complex joins between tables; D. Never, as set-based operations are always more efficient. Các lựa chọn này không đáp ứng định nghĩa/điều kiện được hỏi; khi ôn, hãy chỉ ra chính xác điểm lệch của từng lựa chọn thay vì học thuộc mỗi chữ cái.
\end{question}

\begin{question}{542}
% [ID gốc: 5.106]
% [Nguồn: DBI202 Order 5 Q106]
How can you invoke a user-defined function in a SQL query?
\begin{choices}
\choice Using the EXECUTE statement
\choice Using the CALL statement
\choice Using the SELECT statement
\choice Using the RUN statement
\end{choices}
\textbf{Đáp án:} C
\textbf{Giải thích:} Đáp án theo dữ liệu đề: C — C. Using the SELECT statement.

Nền tảng kiến thức: SQL tách truy vấn dữ liệu (DML) khỏi định nghĩa cấu trúc (DDL) và quyền truy cập (DCL). Với SELECT, cần xác định bảng nguồn, điều kiện lọc, phép kết hợp, nhóm và điều kiện nhóm; các ràng buộc, view, thủ tục và trigger đảm nhiệm các vai trò khác nhau.

Phân tích: Đọc truy vấn theo các giai đoạn logic: FROM/JOIN tạo dữ liệu nguồn, WHERE lọc từng hàng, GROUP BY tạo nhóm, HAVING lọc nhóm, SELECT xác định biểu thức đầu ra và ORDER BY sắp kết quả. Với phép join, kiểm tra điều kiện ghép và các hàng không khớp trước khi tính kết quả.

Đối chiếu phương án: A. Using the EXECUTE statement; B. Using the CALL statement; D. Using the RUN statement. Các lựa chọn này không đáp ứng định nghĩa/điều kiện được hỏi; khi ôn, hãy chỉ ra chính xác điểm lệch của từng lựa chọn thay vì học thuộc mỗi chữ cái.
\end{question}

\begin{question}{543}
% [ID gốc: 5.116]
% [Nguồn: DBI202 Order 5 Q116]
Which of the following is not mentioned how does a local variable is defined using T-SQL?
\begin{choices}
\choice The name of the local variable should begin with "@@" sign as the first character of its name.
\choice A local variable is defined using T-SQL by using statement "DECLARE".
\choice When a variable is first declared, its value is set to NULL. To assign a value to a variable, use the SET statement.
\choice The data type of local variable is declared after its name and assigned a system-supplied or user-defined data type and a length
\end{choices}
\textbf{Đáp án:} A
\textbf{Giải thích:} Đáp án theo dữ liệu đề: A — A. The name of the local variable should begin with "@@" sign as the first character of its name..

Nền tảng kiến thức: SQL tách truy vấn dữ liệu (DML) khỏi định nghĩa cấu trúc (DDL) và quyền truy cập (DCL). Với SELECT, cần xác định bảng nguồn, điều kiện lọc, phép kết hợp, nhóm và điều kiện nhóm; các ràng buộc, view, thủ tục và trigger đảm nhiệm các vai trò khác nhau.

Phân tích: Trong T-SQL, biến cục bộ do người dùng khai báo bắt đầu bằng một dấu `@` (ví dụ `@count`); tên bắt đầu `@@` được dùng cho các hàm/biến hệ thống. Đừng nhầm tiền tố biến cục bộ với các biến hệ thống như `@@ROWCOUNT`.

Đối chiếu phương án: B. A local variable is defined using T-SQL by using statement "DECLARE".; C. When a variable is first declared, its value is set to NULL. To assign a value to a variable,…; D. The data type of local variable is declared after its name and assigned a system-supplied or u…. Các lựa chọn này không đáp ứng định nghĩa/điều kiện được hỏi; khi ôn, hãy chỉ ra chính xác điểm lệch của từng lựa chọn thay vì học thuộc mỗi chữ cái.
\end{question}

\begin{question}{544}
% [ID gốc: 5.117]
% [Nguồn: DBI202 Order 5 Q117]
A/An ___ is a statement gives the user a number of different options in the event, condition, and action parts
\begin{choices}
\choice Trigger
\choice Stored-function
\choice Index
\choice Stored-procedure
\end{choices}
\textbf{Đáp án:} A
\textbf{Giải thích:} Đáp án theo dữ liệu đề: A — A. Trigger.

Nền tảng kiến thức: SQL tách truy vấn dữ liệu (DML) khỏi định nghĩa cấu trúc (DDL) và quyền truy cập (DCL). Với SELECT, cần xác định bảng nguồn, điều kiện lọc, phép kết hợp, nhóm và điều kiện nhóm; các ràng buộc, view, thủ tục và trigger đảm nhiệm các vai trò khác nhau.

Phân tích: Trigger là hành động do DBMS tự kích hoạt khi một sự kiện dữ liệu như INSERT, UPDATE hoặc DELETE xảy ra trên đối tượng đã khai báo. Nó khác stored procedure vốn thường được gọi tường minh; đáp án đúng cần thể hiện đúng sự kiện kích hoạt.

Đối chiếu phương án: B. Stored-function; C. Index; D. Stored-procedure. Các lựa chọn này không đáp ứng định nghĩa/điều kiện được hỏi; khi ôn, hãy chỉ ra chính xác điểm lệch của từng lựa chọn thay vì học thuộc mỗi chữ cái.
\end{question}

\begin{question}{545}
% [ID gốc: 5.118]
% [Nguồn: DBI202 Order 5 Q118]
Which of the following describes an SQL cursor?
\begin{choices}
\choice It allows you to traverse the result set of a query and process each row individually
\choice A field within a database
\choice is a precompiled collection of one or more SQL statements that can be executed as a single unit
\choice Cursors are not related to SQL
\end{choices}
\textbf{Đáp án:} A
\textbf{Giải thích:} Đáp án theo dữ liệu đề: A — A. It allows you to traverse the result set of a query and process each row individually.

Nền tảng kiến thức: SQL tách truy vấn dữ liệu (DML) khỏi định nghĩa cấu trúc (DDL) và quyền truy cập (DCL). Với SELECT, cần xác định bảng nguồn, điều kiện lọc, phép kết hợp, nhóm và điều kiện nhóm; các ràng buộc, view, thủ tục và trigger đảm nhiệm các vai trò khác nhau.

Phân tích: Cursor cho phép duyệt kết quả theo từng hàng và chỉ nên dùng khi xử lý từng bản ghi là cần thiết hoặc không thể biểu diễn gọn bằng thao tác tập hợp. Với INSERT/UPDATE/DELETE hàng loạt, truy vấn set-based thường đơn giản và hiệu quả hơn.

Đối chiếu phương án: C. A field within a database; B. is a precompiled collection of one or more SQL statements that can be executed as a single unit; D. Cursors are not related to SQL. Các lựa chọn này không đáp ứng định nghĩa/điều kiện được hỏi; khi ôn, hãy chỉ ra chính xác điểm lệch của từng lựa chọn thay vì học thuộc mỗi chữ cái.
\end{question}

\begin{question}{546}
% [ID gốc: 5.119]
% [Nguồn: DBI202 Order 5 Q119]
How many types of variables are there in Stored Procedure
\begin{choices}
\choice Default variable and input variable
\choice Default variable and output variable
\choice Input variable and output variable
\choice Local variables and global variables
\end{choices}
\textbf{Đáp án:} C
\textbf{Giải thích:} Đáp án theo dữ liệu đề: C — C. Input variable and output variable.

Nền tảng kiến thức: SQL tách truy vấn dữ liệu (DML) khỏi định nghĩa cấu trúc (DDL) và quyền truy cập (DCL). Với SELECT, cần xác định bảng nguồn, điều kiện lọc, phép kết hợp, nhóm và điều kiện nhóm; các ràng buộc, view, thủ tục và trigger đảm nhiệm các vai trò khác nhau.

Phân tích: Stored procedure là tập lệnh SQL có tên lưu trong CSDL, được gọi để thực thi một thủ tục; nó có thể nhận tham số và chứa nhiều câu lệnh. Đừng nhầm với trigger tự chạy khi có sự kiện hoặc function có quy tắc trả về giá trị.

Đối chiếu phương án: A. Default variable and input variable; B. Default variable and output variable; D. Local variables and global variables. Các lựa chọn này không đáp ứng định nghĩa/điều kiện được hỏi; khi ôn, hãy chỉ ra chính xác điểm lệch của từng lựa chọn thay vì học thuộc mỗi chữ cái.
\end{question}

\begin{question}{547}
% [ID gốc: 5.127]
% [Nguồn: DBI202 Order 5 Q127]
What are the different events in Triggers?
\begin{choices}
\choice Select, Insert
\choice Select, Update
\choice Select, Commit
\choice Insert, Update, Delete
\end{choices}
\textbf{Đáp án:} D
\textbf{Giải thích:} Đáp án theo dữ liệu đề: D — D. Insert, Update, Delete.

Nền tảng kiến thức: SQL tách truy vấn dữ liệu (DML) khỏi định nghĩa cấu trúc (DDL) và quyền truy cập (DCL). Với SELECT, cần xác định bảng nguồn, điều kiện lọc, phép kết hợp, nhóm và điều kiện nhóm; các ràng buộc, view, thủ tục và trigger đảm nhiệm các vai trò khác nhau.

Phân tích: Đọc truy vấn theo các giai đoạn logic: FROM/JOIN tạo dữ liệu nguồn, WHERE lọc từng hàng, GROUP BY tạo nhóm, HAVING lọc nhóm, SELECT xác định biểu thức đầu ra và ORDER BY sắp kết quả. Với phép join, kiểm tra điều kiện ghép và các hàng không khớp trước khi tính kết quả.

Đối chiếu phương án: A. Select, Insert; B. Select, Update; C. Select, Commit. Các lựa chọn này không đáp ứng định nghĩa/điều kiện được hỏi; khi ôn, hãy chỉ ra chính xác điểm lệch của từng lựa chọn thay vì học thuộc mỗi chữ cái.
\end{question}

\begin{question}{548}
% [ID gốc: 5.128]
% [Nguồn: DBI202 Order 5 Q128]
Which of the following methods fetch data from cursor to access data?
\begin{choices}
\choice First, last, next, prior, absolute, relative
\choice First, last, next
\choice First, last, next, prior
\choice First, absolute, relative
\end{choices}
\textbf{Đáp án:} A
\textbf{Giải thích:} Đáp án theo dữ liệu đề: A — A. First, last, next, prior, absolute, relative.

Nền tảng kiến thức: SQL tách truy vấn dữ liệu (DML) khỏi định nghĩa cấu trúc (DDL) và quyền truy cập (DCL). Với SELECT, cần xác định bảng nguồn, điều kiện lọc, phép kết hợp, nhóm và điều kiện nhóm; các ràng buộc, view, thủ tục và trigger đảm nhiệm các vai trò khác nhau.

Phân tích: Cursor cho phép duyệt kết quả theo từng hàng và chỉ nên dùng khi xử lý từng bản ghi là cần thiết hoặc không thể biểu diễn gọn bằng thao tác tập hợp. Với INSERT/UPDATE/DELETE hàng loạt, truy vấn set-based thường đơn giản và hiệu quả hơn.

Đối chiếu phương án: B. First, last, next; C. First, last, next, prior; D. First, absolute, relative. Các lựa chọn này không đáp ứng định nghĩa/điều kiện được hỏi; khi ôn, hãy chỉ ra chính xác điểm lệch của từng lựa chọn thay vì học thuộc mỗi chữ cái.
\end{question}

\begin{question}{549}
% [ID gốc: 5.129]
% [Nguồn: DBI202 Order 5 Q129]
When does the Deleted temporary trigger table become populated?
\begin{choices}
\choice With DELETE and INSERT triggers
\choice With DELETE triggers only
\choice With DELETE, UPDATE, and INSERT triggers
\choice With DELETE and UPDATE triggers
\end{choices}
\textbf{Đáp án:} D
\textbf{Giải thích:} Đáp án theo dữ liệu đề: D — D. With DELETE and UPDATE triggers.

Nền tảng kiến thức: SQL tách truy vấn dữ liệu (DML) khỏi định nghĩa cấu trúc (DDL) và quyền truy cập (DCL). Với SELECT, cần xác định bảng nguồn, điều kiện lọc, phép kết hợp, nhóm và điều kiện nhóm; các ràng buộc, view, thủ tục và trigger đảm nhiệm các vai trò khác nhau.

Phân tích: Trigger là hành động do DBMS tự kích hoạt khi một sự kiện dữ liệu như INSERT, UPDATE hoặc DELETE xảy ra trên đối tượng đã khai báo. Nó khác stored procedure vốn thường được gọi tường minh; đáp án đúng cần thể hiện đúng sự kiện kích hoạt.

Đối chiếu phương án: A. With DELETE and INSERT triggers; B. With DELETE triggers only; C. With DELETE, UPDATE, and INSERT triggers. Các lựa chọn này không đáp ứng định nghĩa/điều kiện được hỏi; khi ôn, hãy chỉ ra chính xác điểm lệch của từng lựa chọn thay vì học thuộc mỗi chữ cái.
\end{question}

\begin{question}{550}
% [ID gốc: 5.130]
% [Nguồn: DBI202 Order 5 Q130]
Which of the following is used to define code that is fired when certain actions or event occur?
\begin{choices}
\choice You use an application Trigger to fire when certain actions or event occur.
\choice You use an application Cursor to fire when certain actions or event occur.
\choice You use an application Replace to fire when certain actions or event occur.
\choice You use an application Keyword to fire when certain actions or event occur.
\end{choices}
\textbf{Đáp án:} A
\textbf{Giải thích:} Đáp án theo dữ liệu đề: A — A. You use an application Trigger to fire when certain actions or event occur..

Nền tảng kiến thức: SQL tách truy vấn dữ liệu (DML) khỏi định nghĩa cấu trúc (DDL) và quyền truy cập (DCL). Với SELECT, cần xác định bảng nguồn, điều kiện lọc, phép kết hợp, nhóm và điều kiện nhóm; các ràng buộc, view, thủ tục và trigger đảm nhiệm các vai trò khác nhau.

Phân tích: Cursor cho phép duyệt kết quả theo từng hàng và chỉ nên dùng khi xử lý từng bản ghi là cần thiết hoặc không thể biểu diễn gọn bằng thao tác tập hợp. Với INSERT/UPDATE/DELETE hàng loạt, truy vấn set-based thường đơn giản và hiệu quả hơn.

Đối chiếu phương án: B. You use an application Cursor to fire when certain actions or event occur.; C. You use an application Replace to fire when certain actions or event occur.; D. You use an application Keyword to fire when certain actions or event occur.. Các lựa chọn này không đáp ứng định nghĩa/điều kiện được hỏi; khi ôn, hãy chỉ ra chính xác điểm lệch của từng lựa chọn thay vì học thuộc mỗi chữ cái.
\end{question}

\begin{question}{551}
% [ID gốc: 5.133]
% [Nguồn: DBI202 Order 5 Q133]
Which of the following is not a typical trigger action?
\begin{choices}
\choice Insert
\choice Select
\choice Update
\choice Delete
\end{choices}
\textbf{Đáp án:} B
\textbf{Giải thích:} Đáp án theo dữ liệu đề: B — B. Select.

Lời giải gốc: Trigger trong SQL được kích hoạt tự động khi có sự thay đổi dữ liệu thông qua các lệnh DML bao gồm INSERT, UPDATE, hoặc DELETE. Lệnh SELECT chỉ dùng để truy vấn, không làm thay đổi dữ liệu nên không kích hoạt trigger.

Nền tảng kiến thức: SQL tách truy vấn dữ liệu (DML) khỏi định nghĩa cấu trúc (DDL) và quyền truy cập (DCL). Với SELECT, cần xác định bảng nguồn, điều kiện lọc, phép kết hợp, nhóm và điều kiện nhóm; các ràng buộc, view, thủ tục và trigger đảm nhiệm các vai trò khác nhau.

Phân tích: Trigger là hành động do DBMS tự kích hoạt khi một sự kiện dữ liệu như INSERT, UPDATE hoặc DELETE xảy ra trên đối tượng đã khai báo. Nó khác stored procedure vốn thường được gọi tường minh; đáp án đúng cần thể hiện đúng sự kiện kích hoạt.

Đối chiếu phương án: A. Insert; C. Update; D. Delete. Các lựa chọn này không đáp ứng định nghĩa/điều kiện được hỏi; khi ôn, hãy chỉ ra chính xác điểm lệch của từng lựa chọn thay vì học thuộc mỗi chữ cái.
\end{question}

\begin{question}{552}
% [ID gốc: 5.139]
% [Nguồn: DBI202 Order 5 Q139]
Which of the following is NOT mentioned on how a local variable is defined using T-SQL?
\begin{choices}
\choice The name of the local variable should begin with "@" sign as the first character of its name.
\choice A local variable is defined using T-SQL by using statement "DECLARE".
\choice When a variable is first declared, its value is set to NULL. To assign a value to a variable, use the SET statement.
\choice The data type of local variable is declared after its name and assigned a system-supplied or user-defined data type and a length.
\end{choices}
\textbf{Đáp án:} A
\textbf{Giải thích:} Đáp án theo dữ liệu đề: A — A. The name of the local variable should begin with "@" sign as the first character of its name..

Lời giải gốc: Biến cục bộ trong T-SQL chỉ bắt đầu bằng một ký tự @ (ví dụ: @VarName). Phương án A ghi sai thành ký tự @@ (vì @@ trong T-SQL được dùng cho các biến hệ thống/biến toàn cục).

Nền tảng kiến thức: SQL tách truy vấn dữ liệu (DML) khỏi định nghĩa cấu trúc (DDL) và quyền truy cập (DCL). Với SELECT, cần xác định bảng nguồn, điều kiện lọc, phép kết hợp, nhóm và điều kiện nhóm; các ràng buộc, view, thủ tục và trigger đảm nhiệm các vai trò khác nhau.

Phân tích: Trong T-SQL, biến cục bộ do người dùng khai báo bắt đầu bằng một dấu `@` (ví dụ `@count`); tên bắt đầu `@@` được dùng cho các hàm/biến hệ thống. Đừng nhầm tiền tố biến cục bộ với các biến hệ thống như `@@ROWCOUNT`.

Đối chiếu phương án: B. A local variable is defined using T-SQL by using statement "DECLARE".; C. When a variable is first declared, its value is set to NULL. To assign a value to a variable,…; D. The data type of local variable is declared after its name and assigned a system-supplied or u…. Các lựa chọn này không đáp ứng định nghĩa/điều kiện được hỏi; khi ôn, hãy chỉ ra chính xác điểm lệch của từng lựa chọn thay vì học thuộc mỗi chữ cái.
\end{question}

\begin{question}{553}
% [ID gốc: 5.149]
% [Nguồn: DBI202 Order 5 Q149]
When does the Inserted temporary trigger table become populated?
\begin{choices}
\choice With INSERT and UPDATE triggers
\choice With DELETE triggers only
\choice With INSERT, UPDATE, and DELETE triggers
\choice With UPDATE and DELETE triggers
\end{choices}
\textbf{Đáp án:} A
\textbf{Giải thích:} Đáp án theo dữ liệu đề: A — A. With INSERT and UPDATE triggers.

Lời giải gốc: Bảng tạm Inserted lưu trữ các bản ghi mới được thêm vào hoặc giá trị sau khi cập nhật, do đó nó được điền dữ liệu trong các trigger INSERT và UPDATE.

Nền tảng kiến thức: SQL tách truy vấn dữ liệu (DML) khỏi định nghĩa cấu trúc (DDL) và quyền truy cập (DCL). Với SELECT, cần xác định bảng nguồn, điều kiện lọc, phép kết hợp, nhóm và điều kiện nhóm; các ràng buộc, view, thủ tục và trigger đảm nhiệm các vai trò khác nhau.

Phân tích: Trigger là hành động do DBMS tự kích hoạt khi một sự kiện dữ liệu như INSERT, UPDATE hoặc DELETE xảy ra trên đối tượng đã khai báo. Nó khác stored procedure vốn thường được gọi tường minh; đáp án đúng cần thể hiện đúng sự kiện kích hoạt.

Đối chiếu phương án: B. With DELETE triggers only; C. With INSERT, UPDATE, and DELETE triggers; D. With UPDATE and DELETE triggers. Các lựa chọn này không đáp ứng định nghĩa/điều kiện được hỏi; khi ôn, hãy chỉ ra chính xác điểm lệch của từng lựa chọn thay vì học thuộc mỗi chữ cái.
\end{question}

\begin{question}{554}
% [ID gốc: 5.163]
% [Nguồn: DBI202 Order 5 Q163]
Which of the following is FALSE about database triggers?
\begin{choices}
\choice It automatically runs when an event occurs in the database server
\choice DML triggers are frequently used for enforcing business rules and data integrity
\choice The same trigger action cannot be defined for more than one user action (for example, INSERT and UPDATE) within a single CREATE TRIGGER statement
\choice When a trigger fires, results are returned to the calling application, just like with stored procedures.
\end{choices}
\textbf{Đáp án:} D
\textbf{Giải thích:} Đáp án theo dữ liệu đề: D — D. When a trigger fires, results are returned to the calling application, just like with stored procedures..

Lời giải gốc: Trigger hoạt động ngầm bên trong cơ sở dữ liệu để đáp ứng các sự kiện thay đổi dữ liệu và thường không trả về tập kết quả trực tiếp cho ứng dụng gọi giống như stored procedure.

Nền tảng kiến thức: SQL tách truy vấn dữ liệu (DML) khỏi định nghĩa cấu trúc (DDL) và quyền truy cập (DCL). Với SELECT, cần xác định bảng nguồn, điều kiện lọc, phép kết hợp, nhóm và điều kiện nhóm; các ràng buộc, view, thủ tục và trigger đảm nhiệm các vai trò khác nhau.

Phân tích: DDL (Data Definition Language) khai báo/thay đổi cấu trúc như bảng và schema; DML truy vấn/thêm/sửa/xóa dữ liệu; DCL quản lý quyền. Chọn theo đối tượng bị tác động: cấu trúc, các hàng hay quyền truy cập.

Đối chiếu phương án: A. It automatically runs when an event occurs in the database server; B. DML triggers are frequently used for enforcing business rules and data integrity; C. The same trigger action cannot be defined for more than one user action (for example, INSERT a…. Các lựa chọn này không đáp ứng định nghĩa/điều kiện được hỏi; khi ôn, hãy chỉ ra chính xác điểm lệch của từng lựa chọn thay vì học thuộc mỗi chữ cái.
\end{question}

\begin{question}{555}
% [ID gốc: 5.164]
% [Nguồn: DBI202 Order 5 Q164]
Which of the following is used to define code that is fired when certain actions or events occur?
\begin{choices}
\choice You use an application Trigger to fire when certain actions or events occur.
\choice You use an application Cursor to fire when certain actions or events occur.
\choice You use an application Replace to fire when certain actions or events occur.
\choice You use an application Keyword to fire when certain actions or events occur.
\end{choices}
\textbf{Đáp án:} A
\textbf{Giải thích:} Đáp án theo dữ liệu đề: A — A. You use an application Trigger to fire when certain actions or events occur..

Lời giải gốc: Trigger là một đoạn mã tự động thực thi (fire) khi có một sự kiện thao tác dữ liệu cụ thể (INSERT, UPDATE, DELETE) xảy ra trên bảng.

Nền tảng kiến thức: SQL tách truy vấn dữ liệu (DML) khỏi định nghĩa cấu trúc (DDL) và quyền truy cập (DCL). Với SELECT, cần xác định bảng nguồn, điều kiện lọc, phép kết hợp, nhóm và điều kiện nhóm; các ràng buộc, view, thủ tục và trigger đảm nhiệm các vai trò khác nhau.

Phân tích: Cursor cho phép duyệt kết quả theo từng hàng và chỉ nên dùng khi xử lý từng bản ghi là cần thiết hoặc không thể biểu diễn gọn bằng thao tác tập hợp. Với INSERT/UPDATE/DELETE hàng loạt, truy vấn set-based thường đơn giản và hiệu quả hơn.

Đối chiếu phương án: B. You use an application Cursor to fire when certain actions or events occur.; C. You use an application Replace to fire when certain actions or events occur.; D. You use an application Keyword to fire when certain actions or events occur.. Các lựa chọn này không đáp ứng định nghĩa/điều kiện được hỏi; khi ôn, hãy chỉ ra chính xác điểm lệch của từng lựa chọn thay vì học thuộc mỗi chữ cái.
\end{question}

\begin{question}{556}
% [ID gốc: 5.176]
% [Nguồn: DBI202 Order 5 Q176]
In which line of the following code is the error? CREATE TRIGGER trigger1 --1 ON dbo.EMPLOYEE -- 2 AFTER FOR INSERT AS -- 3 BEGIN -- 4 INSERT INTO EMPLOYEE BACKUP -- 5 SELECT * FROM INSERTED -- 6 END -- 7
\begin{choices}
\choice Line 1
\choice Line 2
\choice Line 3
\choice Line 5
\choice Line 6
F. There are no lines in error
\end{choices}
\textbf{Đáp án:} C
\textbf{Giải thích:} Đáp án theo dữ liệu đề: C — C. Line 3.

Lời giải gốc: Dòng 3 bị thừa từ khóa. Cú pháp đúng phải là AFTER INSERT hoặc FOR INSERT, không dùng đồng thời cả AFTER FOR INSERT.

Nền tảng kiến thức: SQL tách truy vấn dữ liệu (DML) khỏi định nghĩa cấu trúc (DDL) và quyền truy cập (DCL). Với SELECT, cần xác định bảng nguồn, điều kiện lọc, phép kết hợp, nhóm và điều kiện nhóm; các ràng buộc, view, thủ tục và trigger đảm nhiệm các vai trò khác nhau.

Phân tích: Đọc truy vấn theo các giai đoạn logic: FROM/JOIN tạo dữ liệu nguồn, WHERE lọc từng hàng, GROUP BY tạo nhóm, HAVING lọc nhóm, SELECT xác định biểu thức đầu ra và ORDER BY sắp kết quả. Với phép join, kiểm tra điều kiện ghép và các hàng không khớp trước khi tính kết quả.

Đối chiếu phương án: A. Line 1; B. Line 2; D. Line 5; E. Line 6 F. There are no lines in error. Các lựa chọn này không đáp ứng định nghĩa/điều kiện được hỏi; khi ôn, hãy chỉ ra chính xác điểm lệch của từng lựa chọn thay vì học thuộc mỗi chữ cái.
\end{question}

\begin{question}{557}
% [ID gốc: 5.182]
% [Nguồn: DBI202 Order 5 Q182]
The CREATE TRIGGER statement is used to create a trigger. The ____ clause specifies the table name on which the trigger is to be attached. The ____ specifies that this is an AFTER INSERT trigger.
\begin{choices}
\choice for insert, on
\choice on, for insert
\choice for, insert
\choice None of the mentioned
\end{choices}
\textbf{Đáp án:} B
\textbf{Giải thích:} Đáp án theo dữ liệu đề: B — B. on, for insert.

Lời giải gốc: Cú pháp tạo trigger trong SQL Server dùng mệnh đề ON để chỉ định bảng và tùy chọn thời điểm thực thi như AFTER INSERT.

Nền tảng kiến thức: SQL tách truy vấn dữ liệu (DML) khỏi định nghĩa cấu trúc (DDL) và quyền truy cập (DCL). Với SELECT, cần xác định bảng nguồn, điều kiện lọc, phép kết hợp, nhóm và điều kiện nhóm; các ràng buộc, view, thủ tục và trigger đảm nhiệm các vai trò khác nhau.

Phân tích: Trigger là hành động do DBMS tự kích hoạt khi một sự kiện dữ liệu như INSERT, UPDATE hoặc DELETE xảy ra trên đối tượng đã khai báo. Nó khác stored procedure vốn thường được gọi tường minh; đáp án đúng cần thể hiện đúng sự kiện kích hoạt.

Đối chiếu phương án: A. for insert, on; C. for, insert; D. None of the mentioned. Các lựa chọn này không đáp ứng định nghĩa/điều kiện được hỏi; khi ôn, hãy chỉ ra chính xác điểm lệch của từng lựa chọn thay vì học thuộc mỗi chữ cái.
\end{question}

\begin{question}{558}
% [ID gốc: 5.186]
% [Nguồn: DBI202 Order 5 Q186]
How can you pass parameters to a stored procedure in SQL Server?
\begin{choices}
\choice Using @parameter_name = value
\choice Using ?parameter_name = value
\choice Using :parameter_name = value
\choice Using param_name = value
\end{choices}
\textbf{Đáp án:} A
\textbf{Giải thích:} Đáp án theo dữ liệu đề: A — A. Using @parameter_name = value.

Lời giải gốc: Trong SQL Server, các biến và tham số truyền vào Stored Procedure luôn bắt đầu bằng ký tự @ (ví dụ: @parameter_name = value).

Nền tảng kiến thức: SQL tách truy vấn dữ liệu (DML) khỏi định nghĩa cấu trúc (DDL) và quyền truy cập (DCL). Với SELECT, cần xác định bảng nguồn, điều kiện lọc, phép kết hợp, nhóm và điều kiện nhóm; các ràng buộc, view, thủ tục và trigger đảm nhiệm các vai trò khác nhau.

Phân tích: Stored procedure là tập lệnh SQL có tên lưu trong CSDL, được gọi để thực thi một thủ tục; nó có thể nhận tham số và chứa nhiều câu lệnh. Đừng nhầm với trigger tự chạy khi có sự kiện hoặc function có quy tắc trả về giá trị.

Đối chiếu phương án: B. Using ?parameter_name = value; C. Using :parameter_name = value; D. Using param_name = value. Các lựa chọn này không đáp ứng định nghĩa/điều kiện được hỏi; khi ôn, hãy chỉ ra chính xác điểm lệch của từng lựa chọn thay vì học thuộc mỗi chữ cái.
\end{question}

\begin{question}{559}
% [ID gốc: 5.192]
% [Nguồn: DBI202 Order 5 Q192]
Consider a recursive CTE with OPTION (MAXRECURSION 100) that would generate 101 levels. What is the result of executing this code?
\begin{choices}
\choice 101 rows will be returned with no error.
\choice 10,001 rows will be returned with no error.
\choice 101 rows will be returned with a maximum recursion error.
\choice 10,001 rows will be returned with a maximum recursion error.
\end{choices}
\textbf{Đáp án:} C
\textbf{Giải thích:} Đáp án theo dữ liệu đề: C — C. 101 rows will be returned with a maximum recursion error..

Lời giải gốc: Mệnh đề OPTION (MAXRECURSION 100) giới hạn số lần đệ quy tối đa là 100; khi vượt quá giới hạn này, hệ thống dừng lại và phát sinh lỗi tràn giới hạn đệ quy (maximum recursion error).

Nền tảng kiến thức: SQL tách truy vấn dữ liệu (DML) khỏi định nghĩa cấu trúc (DDL) và quyền truy cập (DCL). Với SELECT, cần xác định bảng nguồn, điều kiện lọc, phép kết hợp, nhóm và điều kiện nhóm; các ràng buộc, view, thủ tục và trigger đảm nhiệm các vai trò khác nhau.

Phân tích: Trong SQL Server, MAXRECURSION đặt giới hạn mức đệ quy của CTE để ngăn vòng lặp vô hạn. Nếu truy vấn cần vượt giới hạn cấu hình, nó phát sinh lỗi/không trả về kết quả thành công đầy đủ; không nên suy luận rằng mọi mức vượt giới hạn vẫn được trả như một kết quả hoàn chỉnh.

Đối chiếu phương án: A. 101 rows will be returned with no error.; B. 10,001 rows will be returned with no error.; D. 10,001 rows will be returned with a maximum recursion error.. Các lựa chọn này không đáp ứng định nghĩa/điều kiện được hỏi; khi ôn, hãy chỉ ra chính xác điểm lệch của từng lựa chọn thay vì học thuộc mỗi chữ cái.
\end{question}

\begin{question}{560}
% [ID gốc: 5.202]
% [Nguồn: DBI202 Order 5 Q202]
Which of the following is NOT correct about how a local variable is defined using T-SQL?
\begin{choices}
\choice The name of the local variable should begin with the "@@" sign as the first character of its name.
\choice A local variable is defined using T-SQL with the statement DECLARE.
\choice When a variable is first declared, its value is set to NULL. To assign a value, use the SET statement.
\choice The data type of a local variable is declared after its name, assigned a system-supplied or user-defined data type and a length.
\end{choices}
\textbf{Đáp án:} A
\textbf{Giải thích:} Đáp án theo dữ liệu đề: A — A. The name of the local variable should begin with the "@@" sign as the first character of its name..

Lời giải gốc: Biến cục bộ trong T-SQL phải bắt đầu bằng một ký tự @ (một dấu a-còng); tiền tố @@ được dành riêng cho các biến toàn cục hệ thống. Do đó phát biểu A là sai.

Nền tảng kiến thức: SQL tách truy vấn dữ liệu (DML) khỏi định nghĩa cấu trúc (DDL) và quyền truy cập (DCL). Với SELECT, cần xác định bảng nguồn, điều kiện lọc, phép kết hợp, nhóm và điều kiện nhóm; các ràng buộc, view, thủ tục và trigger đảm nhiệm các vai trò khác nhau.

Phân tích: Trong T-SQL, biến cục bộ do người dùng khai báo bắt đầu bằng một dấu `@` (ví dụ `@count`); tên bắt đầu `@@` được dùng cho các hàm/biến hệ thống. Đừng nhầm tiền tố biến cục bộ với các biến hệ thống như `@@ROWCOUNT`.

Đối chiếu phương án: B. A local variable is defined using T-SQL with the statement DECLARE.; C. When a variable is first declared, its value is set to NULL. To assign a value, use the SET st…; D. The data type of a local variable is declared after its name, assigned a system-supplied or us…. Các lựa chọn này không đáp ứng định nghĩa/điều kiện được hỏi; khi ôn, hãy chỉ ra chính xác điểm lệch của từng lựa chọn thay vì học thuộc mỗi chữ cái.
\end{question}

\begin{question}{561}
% [ID gốc: 5.209]
% [Nguồn: DBI202 Order 5 Q209]
Which of the following is used to invoke a procedure in SQL?
\begin{choices}
\choice EXECUTE
\choice SELECT
\choice PERFORM
\choice SHOW
\end{choices}
\textbf{Đáp án:} A
\textbf{Giải thích:} Đáp án theo dữ liệu đề: A — A. EXECUTE.

Nền tảng kiến thức: SQL tách truy vấn dữ liệu (DML) khỏi định nghĩa cấu trúc (DDL) và quyền truy cập (DCL). Với SELECT, cần xác định bảng nguồn, điều kiện lọc, phép kết hợp, nhóm và điều kiện nhóm; các ràng buộc, view, thủ tục và trigger đảm nhiệm các vai trò khác nhau.

Phân tích: Stored procedure là tập lệnh SQL có tên lưu trong CSDL, được gọi để thực thi một thủ tục; nó có thể nhận tham số và chứa nhiều câu lệnh. Đừng nhầm với trigger tự chạy khi có sự kiện hoặc function có quy tắc trả về giá trị.

Đối chiếu phương án: B. SELECT; C. PERFORM; D. SHOW. Các lựa chọn này không đáp ứng định nghĩa/điều kiện được hỏi; khi ôn, hãy chỉ ra chính xác điểm lệch của từng lựa chọn thay vì học thuộc mỗi chữ cái.
\end{question}

\begin{question}{562}
% [ID gốc: 5.213]
% [Nguồn: DBI202 Order 5 Q213]
Which of the following is the correct way to execute a stored procedure?
\begin{choices}
\choice EXEC SelectAllCustomers @City = 'London', @PostalCode = 'WA11DP';
\choice EXECUTE SelectAllCustomers 'London', 'WA11DP' OUTPUT
\choice EXEC @City = 'London', @PostalCode = 'WA11DP'
\choice EXECUTE SelectAllCustomers
\end{choices}
\textbf{Đáp án:} A
\textbf{Giải thích:} Đáp án theo dữ liệu đề: A — A. EXEC SelectAllCustomers @City = 'London', @PostalCode = 'WA11DP';.

Lời giải gốc: Cú pháp chuẩn để thực thi một stored procedure có tham số trong SQL Server dùng từ khóa EXEC hoặc EXECUTE kèm theo tên tham số và giá trị tương ứng.

Nền tảng kiến thức: SQL tách truy vấn dữ liệu (DML) khỏi định nghĩa cấu trúc (DDL) và quyền truy cập (DCL). Với SELECT, cần xác định bảng nguồn, điều kiện lọc, phép kết hợp, nhóm và điều kiện nhóm; các ràng buộc, view, thủ tục và trigger đảm nhiệm các vai trò khác nhau.

Phân tích: Stored procedure là tập lệnh SQL có tên lưu trong CSDL, được gọi để thực thi một thủ tục; nó có thể nhận tham số và chứa nhiều câu lệnh. Đừng nhầm với trigger tự chạy khi có sự kiện hoặc function có quy tắc trả về giá trị.

Đối chiếu phương án: B. EXECUTE SelectAllCustomers 'London', 'WA11DP' OUTPUT; C. EXEC @City = 'London', @PostalCode = 'WA11DP'; D. EXECUTE SelectAllCustomers. Các lựa chọn này không đáp ứng định nghĩa/điều kiện được hỏi; khi ôn, hãy chỉ ra chính xác điểm lệch của từng lựa chọn thay vì học thuộc mỗi chữ cái.
\end{question}

\begin{question}{563}
% [ID gốc: 5.215]
% [Nguồn: DBI202 Order 5 Q215]
What makes stored procedures faster than sending a Transact-SQL statement to SQL Server?
\begin{choices}
\choice They are not faster.
\choice The query plan is compiled and cached on the server.
\choice The Transact-SQL statements are shorter.
\choice There are less commands available, so parsing time is shorter.
\end{choices}
\textbf{Đáp án:} B
\textbf{Giải thích:} Đáp án theo dữ liệu đề: B — B. The query plan is compiled and cached on the server..

Lời giải gốc: Stored Procedure được biên dịch sẵn và kế hoạch thực thi (Query execution plan) được lưu vào bộ nhớ đệm (cache) trên server, giúp giảm thời gian biên dịch lại ở các lần gọi sau.

Nền tảng kiến thức: SQL tách truy vấn dữ liệu (DML) khỏi định nghĩa cấu trúc (DDL) và quyền truy cập (DCL). Với SELECT, cần xác định bảng nguồn, điều kiện lọc, phép kết hợp, nhóm và điều kiện nhóm; các ràng buộc, view, thủ tục và trigger đảm nhiệm các vai trò khác nhau.

Phân tích: Trong T-SQL, biến cục bộ do người dùng khai báo bắt đầu bằng một dấu `@` (ví dụ `@count`); tên bắt đầu `@@` được dùng cho các hàm/biến hệ thống. Đừng nhầm tiền tố biến cục bộ với các biến hệ thống như `@@ROWCOUNT`.

Đối chiếu phương án: A. They are not faster.; C. The Transact-SQL statements are shorter.; D. There are less commands available, so parsing time is shorter.. Các lựa chọn này không đáp ứng định nghĩa/điều kiện được hỏi; khi ôn, hãy chỉ ra chính xác điểm lệch của từng lựa chọn thay vì học thuộc mỗi chữ cái.
\end{question}

\begin{question}{564}
% [ID gốc: 5.220]
% [Nguồn: DBI202 Order 5 Q220]
Which of the following statements is true about INSTEAD OF DELETE triggers?
\begin{choices}
\choice They execute in place of the standard DELETE operation.
\choice They execute before the standard DELETE operation.
\choice They cannot modify the DELETED table.
\choice They cannot be used on tables.
\end{choices}
\textbf{Đáp án:} A
\textbf{Giải thích:} Đáp án theo dữ liệu đề: A — A. They execute in place of the standard DELETE operation..

Lời giải gốc: Trigger INSTEAD OF DELETE thực thi thay cho (in place of) thao tác xóa tiêu chuẩn trên bảng hoặc view. ---

Nền tảng kiến thức: SQL tách truy vấn dữ liệu (DML) khỏi định nghĩa cấu trúc (DDL) và quyền truy cập (DCL). Với SELECT, cần xác định bảng nguồn, điều kiện lọc, phép kết hợp, nhóm và điều kiện nhóm; các ràng buộc, view, thủ tục và trigger đảm nhiệm các vai trò khác nhau.

Phân tích: INSTEAD OF DELETE trigger chạy thay cho thao tác DELETE thông thường trên đối tượng được gắn trigger. Trigger có thể kiểm tra hoặc thay đổi hành động; nó khác AFTER trigger, vốn chạy sau khi thao tác gốc được thực hiện.

Đối chiếu phương án: B. They execute before the standard DELETE operation.; C. They cannot modify the DELETED table.; D. They cannot be used on tables.. Các lựa chọn này không đáp ứng định nghĩa/điều kiện được hỏi; khi ôn, hãy chỉ ra chính xác điểm lệch của từng lựa chọn thay vì học thuộc mỗi chữ cái.
\end{question}

\begin{question}{565}
% [ID gốc: 5.238]
% [Nguồn: DBI202 Order 5 Q238]
How do you execute a stored procedure named name_sp?
\begin{choices}
\choice EXECUTE name_sp
\choice EXEC name_sp
\choice CALL name_sp
\choice All of the answers above
\choice A and B are correct
\end{choices}
\textbf{Đáp án:} E
\textbf{Giải thích:} Đáp án theo dữ liệu đề: E — E. A and B are correct.

Lời giải gốc: Trong SQL Server, cú pháp thực thi procedure hỗ trợ từ khóa EXECUTE hoặc viết tắt là EXEC. Cú pháp CALL được dùng ở một số hệ quản trị khác (MySQL, Oracle) chứ không phải SQL Server.

Nền tảng kiến thức: SQL tách truy vấn dữ liệu (DML) khỏi định nghĩa cấu trúc (DDL) và quyền truy cập (DCL). Với SELECT, cần xác định bảng nguồn, điều kiện lọc, phép kết hợp, nhóm và điều kiện nhóm; các ràng buộc, view, thủ tục và trigger đảm nhiệm các vai trò khác nhau.

Phân tích: Stored procedure là tập lệnh SQL có tên lưu trong CSDL, được gọi để thực thi một thủ tục; nó có thể nhận tham số và chứa nhiều câu lệnh. Đừng nhầm với trigger tự chạy khi có sự kiện hoặc function có quy tắc trả về giá trị.

Đối chiếu phương án: A. EXECUTE name_sp; B. EXEC name_sp; C. CALL name_sp; D. All of the answers above. Các lựa chọn này không đáp ứng định nghĩa/điều kiện được hỏi; khi ôn, hãy chỉ ra chính xác điểm lệch của từng lựa chọn thay vì học thuộc mỗi chữ cái.
\end{question}

\begin{question}{566}
% [ID gốc: 5.241]
% [Nguồn: DBI202 Order 5 Q241]
Can a stored procedure call another stored procedure?
\begin{choices}
\choice Yes, a stored procedure can call another stored procedure
\choice No, stored procedures cannot call each other
\choice Only if both procedures are in the same schema
\choice Only if one procedure is a system procedure
\end{choices}
\textbf{Đáp án:} A
\textbf{Giải thích:} Đáp án theo dữ liệu đề: A — A. Yes, a stored procedure can call another stored procedure.

Lời giải gốc: Một stored procedure hoàn toàn có thể gọi đến một stored procedure khác bên trong nội dung của nó (nested stored procedure).

Nền tảng kiến thức: SQL tách truy vấn dữ liệu (DML) khỏi định nghĩa cấu trúc (DDL) và quyền truy cập (DCL). Với SELECT, cần xác định bảng nguồn, điều kiện lọc, phép kết hợp, nhóm và điều kiện nhóm; các ràng buộc, view, thủ tục và trigger đảm nhiệm các vai trò khác nhau.

Phân tích: Stored procedure là tập lệnh SQL có tên lưu trong CSDL, được gọi để thực thi một thủ tục; nó có thể nhận tham số và chứa nhiều câu lệnh. Đừng nhầm với trigger tự chạy khi có sự kiện hoặc function có quy tắc trả về giá trị.

Đối chiếu phương án: B. No, stored procedures cannot call each other; C. Only if both procedures are in the same schema; D. Only if one procedure is a system procedure. Các lựa chọn này không đáp ứng định nghĩa/điều kiện được hỏi; khi ôn, hãy chỉ ra chính xác điểm lệch của từng lựa chọn thay vì học thuộc mỗi chữ cái.
\end{question}

\begin{question}{567}
% [ID gốc: 5.249]
% [Nguồn: DBI202 Order 5 Q249]
Which two statements correctly describe stored procedures? (Choose two correct answers)
\begin{choices}
\choice They can return multiple result sets
\choice They are compiled each time they run
\choice They support conditional logic like IF...ELSE
\choice They replace the need for triggers
\end{choices}
\textbf{Đáp án:} AC
\textbf{Giải thích:} Đáp án theo dữ liệu đề: A, C — A. They can return multiple result sets; C. They support conditional logic like IF...ELSE.

Lời giải gốc: Thủ tục lưu trữ (Stored Procedure) có khả năng trả về nhiều tập kết quả khác nhau từ các câu lệnh SELECT bên trong nó (A), hỗ trợ đầy đủ cấu trúc điều khiển logic lập trình như IF...ELSE, vòng lặp (C). Chúng được biên dịch sẵn một lần chứ không phải biên dịch lại mỗi lần chạy (B sai) và không thay thế hoàn toàn được trigger (D sai).

Nền tảng kiến thức: SQL tách truy vấn dữ liệu (DML) khỏi định nghĩa cấu trúc (DDL) và quyền truy cập (DCL). Với SELECT, cần xác định bảng nguồn, điều kiện lọc, phép kết hợp, nhóm và điều kiện nhóm; các ràng buộc, view, thủ tục và trigger đảm nhiệm các vai trò khác nhau.

Phân tích: Stored procedure là tập lệnh SQL có tên lưu trong CSDL, được gọi để thực thi một thủ tục; nó có thể nhận tham số và chứa nhiều câu lệnh. Đừng nhầm với trigger tự chạy khi có sự kiện hoặc function có quy tắc trả về giá trị.

Đối chiếu phương án: B. They are compiled each time they run; D. They replace the need for triggers. Các lựa chọn này không đáp ứng định nghĩa/điều kiện được hỏi; khi ôn, hãy chỉ ra chính xác điểm lệch của từng lựa chọn thay vì học thuộc mỗi chữ cái.
\end{question}

\begin{question}{568}
% [ID gốc: 5.254]
% [Nguồn: DBI202 Order 5 Q254]
Which is the correct statement about Stored Procedures?
\begin{choices}
\choice It is a set of SQL statements used to perform a certain task and works like a function in other programming languages
\choice It is a single SQL statement used to perform a certain task and works like a function in other programming languages
\choice It is a set of SQL statements used to perform many different tasks and works like a function in other programming languages
\choice a and b are correct
\end{choices}
\textbf{Đáp án:} A
\textbf{Giải thích:} Đáp án theo dữ liệu đề: A — A. It is a set of SQL statements used to perform a certain task and works like a function in other programming languages.

Lời giải gốc: Stored Procedure là một tập hợp các câu lệnh SQL được lưu trữ để thực hiện một tác vụ cụ thể, hoạt động tương tự như một hàm trong ngôn ngữ lập trình.

Nền tảng kiến thức: SQL tách truy vấn dữ liệu (DML) khỏi định nghĩa cấu trúc (DDL) và quyền truy cập (DCL). Với SELECT, cần xác định bảng nguồn, điều kiện lọc, phép kết hợp, nhóm và điều kiện nhóm; các ràng buộc, view, thủ tục và trigger đảm nhiệm các vai trò khác nhau.

Phân tích: Stored procedure là tập lệnh SQL có tên lưu trong CSDL, được gọi để thực thi một thủ tục; nó có thể nhận tham số và chứa nhiều câu lệnh. Đừng nhầm với trigger tự chạy khi có sự kiện hoặc function có quy tắc trả về giá trị.

Đối chiếu phương án: B. It is a single SQL statement used to perform a certain task and works like a function in other…; C. It is a set of SQL statements used to perform many different tasks and works like a function i…; D. a and b are correct. Các lựa chọn này không đáp ứng định nghĩa/điều kiện được hỏi; khi ôn, hãy chỉ ra chính xác điểm lệch của từng lựa chọn thay vì học thuộc mỗi chữ cái.
\end{question}

\begin{question}{569}
% [ID gốc: 5.258]
% [Nguồn: DBI202 Order 5 Q258]
Given a relation: Laptop(model, speed, ram, hd, price) What will this function return when called with dbo.GetLaptopPrice(1)? (Đoạn code định nghĩa hàm chọn price của model = @model)
\begin{choices}
\choice The average laptop price with model = 1
\choice The price of all laptops
\choice The price of laptop with model = 1
\choice The model of laptop with price = 1
\end{choices}
\textbf{Đáp án:} C
\textbf{Giải thích:} Đáp án theo dữ liệu đề: C — C. The price of laptop with model = 1.

Lời giải gốc: Nhìn vào đoạn code: SELECT @price = price FROM Laptop WHERE model = @model. Khi truyền giá trị tham số đầu vào @model = 1, câu lệnh sẽ tìm hàng có model = 1 và lấy ra giá trị cột giá (price) của máy đó để trả về.

Nền tảng kiến thức: SQL tách truy vấn dữ liệu (DML) khỏi định nghĩa cấu trúc (DDL) và quyền truy cập (DCL). Với SELECT, cần xác định bảng nguồn, điều kiện lọc, phép kết hợp, nhóm và điều kiện nhóm; các ràng buộc, view, thủ tục và trigger đảm nhiệm các vai trò khác nhau.

Phân tích: Hàm vô hướng `dbo.GetLaptopPrice(1)` nhận mã model 1 rồi trả về cột `price` của hàng Laptop có model đó. Đối số truyền vào là điều kiện tìm model, không phải tốc độ/RAM/ổ cứng; giá trị số cụ thể phụ thuộc dữ liệu bảng.

Đối chiếu phương án: A. The average laptop price with model = 1; B. The price of all laptops; D. The model of laptop with price = 1. Các lựa chọn này không đáp ứng định nghĩa/điều kiện được hỏi; khi ôn, hãy chỉ ra chính xác điểm lệch của từng lựa chọn thay vì học thuộc mỗi chữ cái.
\end{question}

\begin{question}{570}
% [ID gốc: 5.263]
% [Nguồn: DBI202 Order 5 Q263]
The following SQL statement creates a stored procedure that selects Customers from a particular City with a particular PostalCode from the "Customers" table: CREATE PROCEDURE SelectAllCustomers @City nvarchar(30), @PostalCode nvarchar(10) AS BEGIN SELECT * FROM Customers WHERE City = @City AND PostalCode = @PostalCode; END; GO Which of the following is the way to execute a procedure?
\begin{choices}
\choice EXEC SelectAllCustomers @City = 'London', @PostalCode = 'WA11DP';
\choice EXECUTE SelectAllCustomers 'London', 'WA11DP' OUTPUT
\choice EXEC @City = 'London', @PostalCode = 'WA11DP';
\choice EXECUTE SelectAllCustomers
\end{choices}
\textbf{Đáp án:} A
\textbf{Giải thích:} Đáp án theo dữ liệu đề: A — A. EXEC SelectAllCustomers @City = 'London', @PostalCode = 'WA11DP';.

Lời giải gốc: Cú pháp gọi Stored Procedure có tham số truyền vào đúng tên là EXEC ProcedureName @Param1 = Value1, @Param2 = Value2.

Nền tảng kiến thức: SQL tách truy vấn dữ liệu (DML) khỏi định nghĩa cấu trúc (DDL) và quyền truy cập (DCL). Với SELECT, cần xác định bảng nguồn, điều kiện lọc, phép kết hợp, nhóm và điều kiện nhóm; các ràng buộc, view, thủ tục và trigger đảm nhiệm các vai trò khác nhau.

Phân tích: NVARCHAR lưu văn bản Unicode nên phù hợp với tiếng Việt có dấu; VARCHAR không Unicode có thể phụ thuộc code page. Kích thước trong ngoặc là giới hạn số đơn vị ký tự theo SQL Server, cần chọn dựa trên độ dài dữ liệu thực tế.

Đối chiếu phương án: B. EXECUTE SelectAllCustomers 'London', 'WA11DP' OUTPUT; C. EXEC @City = 'London', @PostalCode = 'WA11DP';; D. EXECUTE SelectAllCustomers. Các lựa chọn này không đáp ứng định nghĩa/điều kiện được hỏi; khi ôn, hãy chỉ ra chính xác điểm lệch của từng lựa chọn thay vì học thuộc mỗi chữ cái.
\end{question}

\begin{question}{571}
% [ID gốc: 5.265]
% [Nguồn: DBI202 Order 5 Q265]
The ____ object is used to execute SQL queries and update statements in JDBC.
\begin{choices}
\choice ResultSet
\choice Connection
\choice Statement
\choice CallableStatement
\end{choices}
\textbf{Đáp án:} C
\textbf{Giải thích:} Đáp án theo dữ liệu đề: C — C. Statement.

Lời giải gốc: Trong JDBC, đối tượng Statement được dùng để gửi các câu lệnh truy vấn và cập nhật SQL tới cơ sở dữ liệu.

Nền tảng kiến thức: SQL tách truy vấn dữ liệu (DML) khỏi định nghĩa cấu trúc (DDL) và quyền truy cập (DCL). Với SELECT, cần xác định bảng nguồn, điều kiện lọc, phép kết hợp, nhóm và điều kiện nhóm; các ràng buộc, view, thủ tục và trigger đảm nhiệm các vai trò khác nhau.

Phân tích: Trong JDBC, Statement đại diện câu lệnh SQL gửi tới DBMS; PreparedStatement phù hợp khi truy vấn có tham số và giúp tách giá trị khỏi câu lệnh. ResultSet thường chứa các hàng trả về từ truy vấn SELECT.

Đối chiếu phương án: A. ResultSet; B. Connection; D. CallableStatement. Các lựa chọn này không đáp ứng định nghĩa/điều kiện được hỏi; khi ôn, hãy chỉ ra chính xác điểm lệch của từng lựa chọn thay vì học thuộc mỗi chữ cái.
\end{question}

\begin{question}{572}
% [ID gốc: 18.21]
% [Nguồn: DBI202 Order 18 Q21]
In which objects can a Transaction be used?
\begin{choices}
\choice View, User-defined function
\choice User-defined function, Stored Procedure
\choice Stored Procedure, Trigger, View
\choice User-defined function, Stored Procedure, Trigger
\choice All of the mentioned answers are wrong
\end{choices}
\textbf{Đáp án:} E
\textbf{Giải thích:} Đáp án theo dữ liệu đề: E — E. All of the mentioned answers are wrong.

Lời giải gốc: Giao dịch (Transaction) thường được đặt trong Stored Procedures và Triggers, nhưng không thể dùng bên trong User-defined functions hoặc View. Do đó các lựa chọn A, B, C, D đều chứa đối tượng không hợp lệ.

Nền tảng kiến thức: Giao dịch là một đơn vị công việc cần thỏa ACID: atomicity (trọn vẹn hoặc hủy), consistency (giữ ràng buộc), isolation (hạn chế ảnh hưởng xen kẽ), durability (đã commit thì bền vững qua sự cố). COMMIT và ROLLBACK kết thúc theo hai hướng khác nhau.

Phân tích: Đối chiếu đúng thuộc tính ACID được hỏi: atomicity là tất-cả-hoặc-không; consistency giữ các ràng buộc; isolation kiểm soát ảnh hưởng giữa giao dịch đồng thời; durability bảo đảm giao dịch đã commit sống sót qua sự cố. COMMIT xác nhận, ROLLBACK hủy thay đổi chưa xác nhận.

Đối chiếu phương án: A. View, User-defined function; B. User-defined function, Stored Procedure; C. Stored Procedure, Trigger, View; D. User-defined function, Stored Procedure, Trigger. Các lựa chọn này không đáp ứng định nghĩa/điều kiện được hỏi; khi ôn, hãy chỉ ra chính xác điểm lệch của từng lựa chọn thay vì học thuộc mỗi chữ cái.
\end{question}

\begin{question}{573}
% [ID gốc: 18.25]
% [Nguồn: DBI202 Order 18 Q25]
In which objects can Transaction be used?
\begin{choices}
\choice View, User-defined function
\choice User-defined function, Store Procedure
\choice Store Procedure, Trigger, View
\choice User-defined function, Store Procedure, Trigger
\choice All of the mentioned answers are wrong
\end{choices}
\textbf{Đáp án:} E
\textbf{Giải thích:} Đáp án theo dữ liệu đề: E — E. All of the mentioned answers are wrong.

Lời giải gốc: Giao dịch (Transaction) không được phép sử dụng trong User-defined functions (UDF). Do đó các phương án chứa UDF đều sai. Phương án E là chính xác.

Nền tảng kiến thức: Giao dịch là một đơn vị công việc cần thỏa ACID: atomicity (trọn vẹn hoặc hủy), consistency (giữ ràng buộc), isolation (hạn chế ảnh hưởng xen kẽ), durability (đã commit thì bền vững qua sự cố). COMMIT và ROLLBACK kết thúc theo hai hướng khác nhau.

Phân tích: Đối chiếu đúng thuộc tính ACID được hỏi: atomicity là tất-cả-hoặc-không; consistency giữ các ràng buộc; isolation kiểm soát ảnh hưởng giữa giao dịch đồng thời; durability bảo đảm giao dịch đã commit sống sót qua sự cố. COMMIT xác nhận, ROLLBACK hủy thay đổi chưa xác nhận.

Đối chiếu phương án: A. View, User-defined function; B. User-defined function, Store Procedure; C. Store Procedure, Trigger, View; D. User-defined function, Store Procedure, Trigger. Các lựa chọn này không đáp ứng định nghĩa/điều kiện được hỏi; khi ôn, hãy chỉ ra chính xác điểm lệch của từng lựa chọn thay vì học thuộc mỗi chữ cái.
\end{question}

\end{document}
